%==============================================================================
% Auto-generated by autoinformalization
% Library: balanced_product_codes
% Generated: 2026-03-13 18:27:56
%==============================================================================

%--- Rem_1: BaseField ---
\chapter{Rem 1: Base Field}

For the remainder of this manuscript we consider only vector spaces over the field $\mathbb{F}_2$. All tensor products, linear maps, and homology computations are taken over $\mathbb{F}_2$ unless stated otherwise.

\begin{definition}[$\mathbb{F}_2$]
\label{def:F2}
\lean{𝔽₂}
\leanok

The \emph{base field} $\mathbb{F}_2$ is the field with two elements, defined as $\mathbb{F}_2 := \mathbb{Z}/2\mathbb{Z}$ (i.e., \texttt{ZMod 2}). All vector spaces, tensor products, linear maps, and homology computations in this paper are over $\mathbb{F}_2$.
\end{definition}

\begin{theorem}[Cardinality of $\mathbb{F}_2$]
\label{thm:F2.card_eq}
\lean{𝔽₂.card_eq}
\leanok
\uses{def:F2}
The cardinality of $\mathbb{F}_2$ is $2$, i.e., $|\mathbb{F}_2| = 2$.
\end{theorem}

\begin{proof}
\leanok
\uses{def:F2}
This follows directly from \texttt{ZMod.card 2}, which computes the cardinality of $\mathbb{Z}/2\mathbb{Z}$ to be $2$.
\end{proof}

\begin{theorem}[Negation is Identity in $\mathbb{F}_2$]
\label{thm:F2.neg_eq_self}
\lean{𝔽₂.neg_eq_self}
\leanok
\uses{def:F2}
For every element $a \in \mathbb{F}_2$, we have $-a = a$.
\end{theorem}

\begin{proof}
\leanok
\uses{def:F2}
Since $\mathbb{F}_2$ has characteristic $2$, the result follows from the general fact \texttt{CharTwo.neg\_eq}, which states that in any ring of characteristic $2$, negation is the identity map. Applying this to an arbitrary $a \in \mathbb{F}_2$ gives $-a = a$.
\end{proof}

\begin{theorem}[$a + a = 0$ in $\mathbb{F}_2$]
\label{thm:F2.add_self_eq_zero}
\lean{𝔽₂.add_self_eq_zero}
\leanok
\uses{def:F2}
For every element $a \in \mathbb{F}_2$, we have $a + a = 0$.
\end{theorem}

\begin{proof}
\leanok
\uses{def:F2}
Since $\mathbb{F}_2$ has characteristic $2$, this follows from \texttt{CharTwo.add\_self\_eq\_zero}, which asserts that in any ring of characteristic $2$, every element $a$ satisfies $a + a = 0$. Applying this to arbitrary $a \in \mathbb{F}_2$ yields the result.
\end{proof}

\begin{theorem}[Subtraction Equals Addition in $\mathbb{F}_2$]
\label{thm:F2.sub_eq_add}
\lean{𝔽₂.sub_eq_add}
\leanok
\uses{def:F2, thm:F2.neg_eq_self}
For all $a, b \in \mathbb{F}_2$, we have $a - b = a + b$.
\end{theorem}

\begin{proof}
\leanok
\uses{def:F2, thm:F2.neg_eq_self}
Since $\mathbb{F}_2$ has characteristic $2$, we apply \texttt{CharTwo.sub\_eq\_add}, which states that in any ring of characteristic $2$, subtraction coincides with addition (as $-b = b$ for all $b$). Applying this to arbitrary $a, b \in \mathbb{F}_2$ gives $a - b = a + b$.
\end{proof}

%--- Rem_2: NotationConventions ---
\chapter{Rem 2: Notation Conventions}

This chapter establishes the notation conventions used throughout the formalization, documenting the precise correspondence between paper notation and Mathlib definitions.

\begin{definition}[Chain Complex over $\mathbb{F}_2$]
\label{def:ChainComplexF2}
\lean{ChainComplex𝔽₂}
\leanok
\uses{def:F2}
A \emph{chain complex of $\mathbb{F}_2$-vector spaces} indexed by $\mathbb{Z}$ is defined as
\[
  \operatorname{ChainComplex}(\operatorname{ModuleCat}\,\mathbb{F}_2,\,\mathbb{Z}).
\]
The differential $\partial_i : C_i \to C_{i-1}$ is given by $\mathtt{d}\,i\,j$ in Mathlib where the complex shape satisfies $j + 1 = i$ (i.e.\ $j = i - 1$). This corresponds to the paper's $(C_\bullet, \partial)$.
\end{definition}

\begin{definition}[Cochain Complex over $\mathbb{F}_2$]
\label{def:CochainComplexF2}
\lean{CochainComplex𝔽₂}
\leanok
\uses{def:F2}
A \emph{cochain complex of $\mathbb{F}_2$-vector spaces} indexed by $\mathbb{Z}$ is defined as
\[
  \operatorname{CochainComplex}(\operatorname{ModuleCat}\,\mathbb{F}_2,\,\mathbb{Z}).
\]
The differential $\delta^i : C^i \to C^{i+1}$ is given by $\mathtt{d}\,i\,j$ in Mathlib where $i + 1 = j$. In the paper, the coboundary $\delta^i$ is the transpose of $\partial_{i+1}$.
\end{definition}

\begin{definition}[Boundary Map]
\label{def:NotationConventions.boundary}
\lean{NotationConventions.boundary}
\leanok
\uses{def:ChainComplexF2}
For a chain complex $C \in \operatorname{ChainComplex}_{\mathbb{F}_2}$ and indices $i, j \in \mathbb{Z}$, the \emph{boundary map} $\partial_i : C_i \to C_{i-1}$ is defined as
\[
  \operatorname{boundary}(C, i, j) \;:=\; C.\mathtt{d}\,i\,j.
\]
This is the paper's differential $\partial$ accessed via the Mathlib API \texttt{HomologicalComplex.d}.
\end{definition}

\begin{definition}[Coboundary Map]
\label{def:NotationConventions.coboundary}
\lean{NotationConventions.coboundary}
\leanok
\uses{def:CochainComplexF2}
For a cochain complex $C \in \operatorname{CochainComplex}_{\mathbb{F}_2}$ and indices $i, j \in \mathbb{Z}$, the \emph{coboundary map} $\delta^i : C^i \to C^{i+1}$ is defined as
\[
  \operatorname{coboundary}(C, i, j) \;:=\; C.\mathtt{d}\,i\,j.
\]
This corresponds to the paper's coboundary $\delta$, accessed via \texttt{HomologicalComplex.d} on cochain complexes.
\end{definition}

\begin{definition}[Image of a Linear Map]
\label{def:NotationConventions.im}
\lean{NotationConventions.im}
\leanok
\uses{def:F2}
For $\mathbb{F}_2$-modules $M$ and $N$ and a linear map $f : M \to_{\mathbb{F}_2} N$, the \emph{image} of $f$, denoted $\operatorname{im}(f)$ in the paper, is defined as
\[
  \operatorname{im}(f) \;:=\; \operatorname{LinearMap.range}(f).
\]
\end{definition}

\begin{definition}[Identity Map]
\label{def:NotationConventions.idMap}
\lean{NotationConventions.idMap}
\leanok
\uses{def:F2}
For an $\mathbb{F}_2$-module $M$, the \emph{identity linear map}, denoted $\operatorname{id}$ in the paper, is defined as
\[
  \operatorname{idMap}(M) \;:=\; \operatorname{LinearMap.id}_{\mathbb{F}_2,\,M}.
\]
\end{definition}

\begin{definition}[Transpose of a Linear Map]
\label{def:NotationConventions.transpose}
\lean{NotationConventions.transpose}
\leanok
\uses{def:F2}
For $\mathbb{F}_2$-modules $M$ and $N$ and a linear map $f : M \to_{\mathbb{F}_2} N$, the \emph{transpose} (dual) of $f$ is defined as
\[
  \operatorname{transpose}(f) \;:=\; \operatorname{Module.Dual.transpose}(f) : N^* \to M^*.
\]
In the paper, the coboundary satisfies $\delta = \partial^\top$.
\end{definition}

\begin{theorem}[Boundary Squares to Zero]
\label{thm:NotationConventions.boundary_comp_boundary}
\lean{NotationConventions.boundary_comp_boundary}
\leanok
\uses{def:ChainComplexF2, def:NotationConventions.boundary}
For any chain complex $C \in \operatorname{ChainComplex}_{\mathbb{F}_2}$ and indices $i, j, k \in \mathbb{Z}$,
\[
  \partial_j \circ \partial_i \;=\; 0,
\]
i.e.\ $C.\mathtt{d}\,i\,j \,\gg\, C.\mathtt{d}\,j\,k = 0$.
\end{theorem}

\begin{proof}
\leanok
\uses{def:ChainComplexF2}
This follows directly from the Mathlib lemma \texttt{HomologicalComplex.d\_comp\_d}, applied to $C$ at indices $i$, $j$, $k$. Specifically, we have $C.\mathtt{d\_comp\_d}\,i\,j\,k = 0$ by the axioms of a homological complex, which states that consecutive differentials compose to zero.
\end{proof}

\begin{theorem}[Coboundary Squares to Zero]
\label{thm:NotationConventions.coboundary_comp_coboundary}
\lean{NotationConventions.coboundary_comp_coboundary}
\leanok
\uses{def:CochainComplexF2, def:NotationConventions.coboundary}
For any cochain complex $C \in \operatorname{CochainComplex}_{\mathbb{F}_2}$ and indices $i, j, k \in \mathbb{Z}$,
\[
  \delta^j \circ \delta^i \;=\; 0,
\]
i.e.\ $C.\mathtt{d}\,i\,j \,\gg\, C.\mathtt{d}\,j\,k = 0$.
\end{theorem}

\begin{proof}
\leanok
\uses{def:CochainComplexF2}
This follows directly from the Mathlib lemma \texttt{HomologicalComplex.d\_comp\_d}, applied to the cochain complex $C$ at indices $i$, $j$, $k$. Consecutive coboundary maps compose to zero by the axioms of a homological complex.
\end{proof}

\begin{theorem}[Image Equals Range]
\label{thm:NotationConventions.im_eq_range}
\lean{NotationConventions.im_eq_range}
\leanok
\uses{def:NotationConventions.im}
For $\mathbb{F}_2$-modules $M$, $N$ and a linear map $f : M \to_{\mathbb{F}_2} N$,
\[
  \operatorname{im}(f) \;=\; \operatorname{LinearMap.range}(f).
\]
\end{theorem}

\begin{proof}
\leanok
\uses{def:NotationConventions.im}
This holds by reflexivity, since $\operatorname{im}(f)$ is defined to be $\operatorname{LinearMap.range}(f)$ by definition.
\end{proof}

\begin{theorem}[Identity Map Equals \texttt{LinearMap.id}]
\label{thm:NotationConventions.idMap_eq_id}
\lean{NotationConventions.idMap_eq_id}
\leanok
\uses{def:NotationConventions.idMap}
For any $\mathbb{F}_2$-module $M$,
\[
  \operatorname{idMap}(M) \;=\; \operatorname{LinearMap.id}.
\]
\end{theorem}

\begin{proof}
\leanok
\uses{def:NotationConventions.idMap}
This holds by reflexivity, since $\operatorname{idMap}(M)$ is defined to be $\operatorname{LinearMap.id}_{\mathbb{F}_2,\,M}$ by definition.
\end{proof}

\begin{remark}[Mathlib Correspondence for Homological Algebra]
The following Mathlib definitions correspond to the paper's notation:
\begin{itemize}
  \item $\partial$ or $\delta$ (differential): \texttt{HomologicalComplex.d\,i\,j}
  \item $Z_i(C)$ (cycles $= \ker \partial_i$): \texttt{HomologicalComplex.cycles\,C\,i}
  \item $B_i(C)$ (boundaries $= \operatorname{im}\,\partial_{i+1}$): image of the incoming differential
  \item $H_i(C)$ (homology $= Z_i(C)/B_i(C)$): \texttt{HomologicalComplex.homology\,C\,i}
  \item $\operatorname{im}(f)$ (image): \texttt{LinearMap.range\,f}
  \item $\delta = \partial^\top$ (coboundary as transpose): \texttt{Module.Dual.transpose}
  \item $\operatorname{id}$ (identity): \texttt{LinearMap.id}
  \item $\operatorname{Tot}$ (total complex functor): \texttt{HomologicalComplex.total}
\end{itemize}
The same Mathlib API applies to cochains ($Z^i$, $B^i$, $H^i$) via \texttt{CochainComplex}.
\end{remark}

%--- Rem_3: ExpandingMatrixDefinition ---
\chapter{Rem 3: Expanding Matrix Definition}

\begin{definition}[Hamming Weight]
\label{def:hammingWeight}
\lean{hammingWeight}
\leanok

The \emph{Hamming weight} of a vector $x : \operatorname{Fin} n \to \mathbb{F}_2$, denoted $|x|$, is the number of nonzero entries of $x$. This is the specialization of \texttt{hammingNorm} from Mathlib to vectors over $\mathbb{F}_2$.
\end{definition}

\begin{definition}[Expanding Matrix]
\label{def:IsExpanding}
\lean{IsExpanding}
\leanok
\uses{def:hammingWeight, def:F2}
A matrix $A \in \mathbb{F}_2^{m \times n}$ is called \emph{$(\alpha, \beta)$-expanding}, where $\alpha, \beta \in \mathbb{R}$, if the following conditions hold:
\begin{enumerate}
  \item $0 < \alpha \leq 1$,
  \item $0 < \beta$,
  \item for every vector $x \in \mathbb{F}_2^n$ satisfying $|x| \leq \alpha \cdot n$, we have $|Ax| \geq \beta \cdot |x|$,
\end{enumerate}
where $|\cdot|$ denotes the Hamming weight, and $Ax$ denotes the matrix-vector product $A \cdot x$ over $\mathbb{F}_2$.

Formally, $\operatorname{IsExpanding}(A, \alpha, \beta)$ holds if and only if
\[
  0 < \alpha \;\land\; \alpha \leq 1 \;\land\; 0 < \beta \;\land\; \forall\, x : \operatorname{Fin} n \to \mathbb{F}_2,\quad (|x| : \mathbb{R}) \leq \alpha \cdot n \;\Rightarrow\; (|Ax| : \mathbb{R}) \geq \beta \cdot (|x| : \mathbb{R}).
\]
\end{definition}

%--- Def_1: ChainComplex ---
\chapter{Def 1: Chain Complex}

\begin{definition}[Differential]
\label{def:ChainComplexF2.differential}
\lean{ChainComplex𝔽₂.differential}
\leanok
\uses{def:ChainComplexF2}
Let $C$ be a chain complex of $\mathbb{F}_2$-vector spaces. The \emph{differential}
$\partial_i : C_i \to C_{i-1}$ is the $\mathbb{F}_2$-linear map obtained from the Mathlib
differential $C.d\, i\, (i-1)$ (with respect to the descending complex shape on $\mathbb{Z}$)
by extracting its underlying linear map via the forgetful functor from $\mathbf{Mod}_{\mathbb{F}_2}$
to $\mathbb{F}_2$-linear maps.
\end{definition}

\begin{definition}[Incoming Differential]
\label{def:ChainComplexF2.incomingDifferential}
\lean{ChainComplex𝔽₂.incomingDifferential}
\leanok
\uses{def:ChainComplexF2}
Let $C$ be a chain complex of $\mathbb{F}_2$-vector spaces. The \emph{incoming differential}
at degree $i$ is the $\mathbb{F}_2$-linear map
$\partial_{i+1} : C_{i+1} \to C_i$
extracted from the Mathlib differential $C.d\,(i+1)\,i$. This formulation avoids
definitional issues with $(i+1)-1$ vs.\ $i$ in integer arithmetic.
\end{definition}

\begin{definition}[Cycles]
\label{def:ChainComplexF2.cycles}
\lean{ChainComplex𝔽₂.cycles}
\leanok
\uses{def:ChainComplexF2, def:ChainComplexF2.differential}
Let $C$ be a chain complex of $\mathbb{F}_2$-vector spaces. The \emph{cycles} at degree $i$ are
\[
  Z_i(C) \;=\; \ker \partial_i \;\subseteq\; C_i,
\]
defined as the kernel submodule of the differential $\partial_i : C_i \to C_{i-1}$.
\end{definition}

\begin{definition}[Boundaries]
\label{def:ChainComplexF2.boundaries}
\lean{ChainComplex𝔽₂.boundaries}
\leanok
\uses{def:ChainComplexF2, def:ChainComplexF2.incomingDifferential}
Let $C$ be a chain complex of $\mathbb{F}_2$-vector spaces. The \emph{boundaries} at degree $i$ are
\[
  B_i(C) \;=\; \operatorname{im} \partial_{i+1} \;\subseteq\; C_i,
\]
defined as the image (range) submodule of the incoming differential $\partial_{i+1} : C_{i+1} \to C_i$.
\end{definition}

\begin{theorem}[Chain Condition: Differential Composed with Incoming Differential is Zero]
\label{thm:ChainComplexF2.differential_comp_incomingDifferential}
\lean{ChainComplex𝔽₂.differential_comp_incomingDifferential}
\leanok
\uses{def:ChainComplexF2, def:ChainComplexF2.differential, def:ChainComplexF2.incomingDifferential}
Let $C$ be a chain complex of $\mathbb{F}_2$-vector spaces. For every $i \in \mathbb{Z}$,
\[
  \partial_i \circ \partial_{i+1} \;=\; 0,
\]
as a composition of $\mathbb{F}_2$-linear maps $C_{i+1} \to C_{i-1}$.
\end{theorem}

\begin{proof}
\leanok
\uses{def:ChainComplexF2.differential, def:ChainComplexF2.incomingDifferential}
By extensionality, it suffices to show that $(\partial_i \circ \partial_{i+1})(x) = 0$ for an
arbitrary element $x \in C_{i+1}$.
Unfolding the definitions of $\partial_i$ (i.e., $\mathtt{differential}$) and $\partial_{i+1}$
(i.e., $\mathtt{incomingDifferential}$) and simplifying the linear map composition and zero map,
we reduce to showing that $(C.d\,(i+1)\,i \gg C.d\,i\,(i-1))(x) = 0$.
By the Mathlib chain complex axiom $C.\mathtt{d\_comp\_d}\,(i+1)\,i\,(i-1)$, the composition
$C.d\,(i+1)\,i \gg C.d\,i\,(i-1) = 0$, so after simplification the expression vanishes.
Rewriting using the $\mathbf{Mod}_{\mathbb{F}_2}$ composition application lemma completes the proof.
\end{proof}

\begin{theorem}[Boundaries are Contained in Cycles]
\label{thm:ChainComplexF2.boundaries_le_cycles}
\lean{ChainComplex𝔽₂.boundaries_le_cycles}
\leanok
\uses{def:ChainComplexF2, def:ChainComplexF2.boundaries, def:ChainComplexF2.cycles, thm:ChainComplexF2.differential_comp_incomingDifferential}
Let $C$ be a chain complex of $\mathbb{F}_2$-vector spaces. For every $i \in \mathbb{Z}$,
\[
  B_i(C) \;\subseteq\; Z_i(C).
\]
That is, every boundary is a cycle.
\end{theorem}

\begin{proof}
\leanok
\uses{def:ChainComplexF2.boundaries, def:ChainComplexF2.cycles, thm:ChainComplexF2.differential_comp_incomingDifferential}
Unfolding the definitions of $B_i(C)$ (i.e., $\mathtt{boundaries}$) and $Z_i(C)$
(i.e., $\mathtt{cycles}$), the inclusion $\operatorname{im}\partial_{i+1} \subseteq \ker \partial_i$
is equivalent (by the submodule lemma $\mathtt{LinearMap.range\_le\_ker\_iff}$) to the condition
$\partial_i \circ \partial_{i+1} = 0$. This follows immediately from
$\mathtt{differential\_comp\_incomingDifferential}$.
\end{proof}

\begin{definition}[Boundaries as a Submodule of Cycles]
\label{def:ChainComplexF2.boundariesSubCycles}
\lean{ChainComplex𝔽₂.boundariesSubCycles}
\leanok
\uses{def:ChainComplexF2, def:ChainComplexF2.boundaries, def:ChainComplexF2.cycles}
Let $C$ be a chain complex of $\mathbb{F}_2$-vector spaces. The boundaries $B_i(C)$, viewed as a
submodule of the cycles $Z_i(C)$, are defined as the preimage of $B_i(C) \subseteq C_i$ under
the canonical subtype embedding $Z_i(C) \hookrightarrow C_i$. Concretely,
\[
  B_i^{Z}(C) \;=\; \bigl\{\, z \in Z_i(C) \;\bigm|\; z \in B_i(C) \,\bigr\}
  \;\subseteq\; Z_i(C).
\]
\end{definition}

\begin{definition}[Homology]
\label{def:ChainComplexF2.homologyPrime}
\lean{ChainComplex𝔽₂.homology'}
\leanok
\uses{def:ChainComplexF2, def:ChainComplexF2.cycles, def:ChainComplexF2.boundariesSubCycles}
Let $C$ be a chain complex of $\mathbb{F}_2$-vector spaces. The \emph{$i$-th homology} of $C$ is
the quotient $\mathbb{F}_2$-module
\[
  H_i(C) \;=\; Z_i(C) \,/\, B_i^{Z}(C),
\]
where $B_i^{Z}(C)$ is the boundaries viewed as a submodule of the cycles $Z_i(C)$.
\end{definition}

%--- Def_2: CochainsCohomology ---
\chapter{Def 2: Cochains and Cohomology}

\begin{definition}[Coboundary Map]
\label{def:ChainComplexF2.coboundaryMap}
\lean{ChainComplex𝔽₂.coboundaryMap}
\leanok
\uses{def:ChainComplexF2, def:ChainComplexF2.incomingDifferential}
Let $C = (C_\bullet, \partial)$ be a chain complex over $\mathbb{F}_2$. For each $i \in \mathbb{Z}$, the \emph{coboundary map} $\delta^i : \operatorname{Dual}(C_i) \to \operatorname{Dual}(C_{i+1})$ is defined as the dual (transpose) of the incoming differential $\partial_{i+1} : C_{i+1} \to C_i$:
\[
  \delta^i \;=\; (\partial_{i+1})^{\mathrm{tr}} \;:=\; \operatorname{dualMap}(C.\operatorname{incomingDifferential}(i)).
\]
\end{definition}

\begin{definition}[Incoming Coboundary Map]
\label{def:ChainComplexF2.incomingCoboundaryMap}
\lean{ChainComplex𝔽₂.incomingCoboundaryMap}
\leanok
\uses{def:ChainComplexF2, def:ChainComplexF2.differential}
Let $C = (C_\bullet, \partial)$ be a chain complex over $\mathbb{F}_2$. For each $i \in \mathbb{Z}$, the \emph{incoming coboundary map} $\delta^{i-1} : \operatorname{Dual}(C_{i-1}) \to \operatorname{Dual}(C_i)$ is defined as the dual (transpose) of the differential $\partial_i : C_i \to C_{i-1}$:
\[
  \delta^{i-1} \;=\; (\partial_i)^{\mathrm{tr}} \;:=\; \operatorname{dualMap}(C.\operatorname{differential}(i)).
\]
This formulation avoids index-arithmetic issues by using $C.\operatorname{differential}(i)$ directly.
\end{definition}

\begin{definition}[Cocycles]
\label{def:ChainComplexF2.cocycles}
\lean{ChainComplex𝔽₂.cocycles}
\leanok
\uses{def:ChainComplexF2.coboundaryMap}
Let $C$ be a chain complex over $\mathbb{F}_2$. The \emph{cocycles} in degree $i$ are the kernel of the coboundary map $\delta^i$:
\[
  Z^i(C) \;=\; \ker \delta^i \;\subseteq\; \operatorname{Dual}(C_i),
\]
viewed as a submodule of $\operatorname{Dual}(C_i)$ over $\mathbb{F}_2$.
\end{definition}

\begin{definition}[Coboundaries]
\label{def:ChainComplexF2.coboundaries}
\lean{ChainComplex𝔽₂.coboundaries}
\leanok
\uses{def:ChainComplexF2.incomingCoboundaryMap}
Let $C$ be a chain complex over $\mathbb{F}_2$. The \emph{coboundaries} in degree $i$ are the image (range) of the incoming coboundary map $\delta^{i-1}$:
\[
  B^i(C) \;=\; \operatorname{im}\,\delta^{i-1} \;\subseteq\; \operatorname{Dual}(C_i),
\]
viewed as a submodule of $\operatorname{Dual}(C_i)$ over $\mathbb{F}_2$.
\end{definition}

\begin{theorem}[Coboundary Composition is Zero]
\label{thm:ChainComplexF2.coboundaryMap_comp_incomingCoboundaryMap}
\lean{ChainComplex𝔽₂.coboundaryMap_comp_incomingCoboundaryMap}
\leanok
\uses{def:ChainComplexF2.coboundaryMap, def:ChainComplexF2.incomingCoboundaryMap, thm:ChainComplexF2.differential_comp_incomingDifferential}
Let $C$ be a chain complex over $\mathbb{F}_2$. For each $i \in \mathbb{Z}$, the composition of coboundary maps satisfies:
\[
  \delta^i \circ \delta^{i-1} \;=\; 0 \;:\; \operatorname{Dual}(C_{i-1}) \;\to\; \operatorname{Dual}(C_{i+1}).
\]
\end{theorem}

\begin{proof}
\leanok
\uses{def:ChainComplexF2.coboundaryMap, def:ChainComplexF2.incomingCoboundaryMap, thm:ChainComplexF2.differential_comp_incomingDifferential}
Unfolding the definitions of $\delta^i = \operatorname{dualMap}(\partial_{i+1})$ and $\delta^{i-1} = \operatorname{dualMap}(\partial_i)$, we apply the identity $\operatorname{dualMap}(f) \circ \operatorname{dualMap}(g) = \operatorname{dualMap}(g \circ f)$ (i.e., $\operatorname{dualMap}$ is contravariant). Thus
\[
  \delta^i \circ \delta^{i-1} \;=\; \operatorname{dualMap}(\partial_i \circ \partial_{i+1}).
\]
Rewriting using the chain complex relation $\partial_i \circ \partial_{i+1} = 0$ (which holds by \texttt{differential\_comp\_incomingDifferential}), we obtain $\operatorname{dualMap}(0) = 0$, which follows from the fact that any linear map sends zero to zero.
\end{proof}

\begin{theorem}[Coboundaries are Contained in Cocycles]
\label{thm:ChainComplexF2.coboundaries_le_cocycles}
\lean{ChainComplex𝔽₂.coboundaries_le_cocycles}
\leanok
\uses{def:ChainComplexF2.coboundaries, def:ChainComplexF2.cocycles, thm:ChainComplexF2.coboundaryMap_comp_incomingCoboundaryMap}
Let $C$ be a chain complex over $\mathbb{F}_2$. For each $i \in \mathbb{Z}$:
\[
  B^i(C) \;\subseteq\; Z^i(C).
\]
Every coboundary is a cocycle.
\end{theorem}

\begin{proof}
\leanok
\uses{def:ChainComplexF2.coboundaries, def:ChainComplexF2.cocycles, thm:ChainComplexF2.coboundaryMap_comp_incomingCoboundaryMap}
Unfolding the definitions of $B^i(C) = \operatorname{im}\,\delta^{i-1}$ and $Z^i(C) = \ker \delta^i$, we apply the standard criterion: $\operatorname{im}(f) \subseteq \ker(g)$ if and only if $g \circ f = 0$. Rewriting using the range--kernel inclusion lemma, it suffices to show $\delta^i \circ \delta^{i-1} = 0$, which is exactly the content of \texttt{coboundaryMap\_comp\_incomingCoboundaryMap}.
\end{proof}

\begin{definition}[Coboundaries as Submodule of Cocycles]
\label{def:ChainComplexF2.coboundariesSubCocycles}
\lean{ChainComplex𝔽₂.coboundariesSubCocycles}
\leanok
\uses{def:ChainComplexF2.coboundaries, def:ChainComplexF2.cocycles, thm:ChainComplexF2.coboundaries_le_cocycles}
Since $B^i(C) \subseteq Z^i(C)$, we define $B^i(C)$ as a submodule of $Z^i(C)$ by taking the preimage of $B^i(C)$ under the subtype embedding $Z^i(C) \hookrightarrow \operatorname{Dual}(C_i)$:
\[
  B^i(C)|_{Z^i} \;:=\; \operatorname{comap}\bigl(\iota_{Z^i}\bigr)\bigl(B^i(C)\bigr) \;\subseteq\; Z^i(C),
\]
where $\iota_{Z^i} : Z^i(C) \to \operatorname{Dual}(C_i)$ is the canonical inclusion.
\end{definition}

\begin{definition}[Cohomology]
\label{def:ChainComplexF2.cohomology}
\lean{ChainComplex𝔽₂.cohomology}
\leanok
\uses{def:ChainComplexF2.cocycles, def:ChainComplexF2.coboundariesSubCocycles}
Let $C$ be a chain complex over $\mathbb{F}_2$. The \emph{$i$-th cohomology} of $C$ is the quotient $\mathbb{F}_2$-vector space:
\[
  H^i(C) \;=\; Z^i(C) \,/\, B^i(C).
\]
More precisely, it is the quotient of $Z^i(C)$ by the submodule $B^i(C)|_{Z^i}$ of $Z^i(C)$.
\end{definition}

\begin{theorem}[Cocycles are the Dual Annihilator of Boundaries]
\label{thm:ChainComplexF2.cocycles_eq_dualAnnihilator_boundaries}
\lean{ChainComplex𝔽₂.cocycles_eq_dualAnnihilator_boundaries}
\leanok
\uses{def:ChainComplexF2.cocycles, def:ChainComplexF2.boundaries}
Let $C$ be a chain complex over $\mathbb{F}_2$. For each $i \in \mathbb{Z}$:
\[
  Z^i(C) \;=\; B_i(C)^{\mathrm{ann}},
\]
where $B_i(C)^{\mathrm{ann}} \subseteq \operatorname{Dual}(C_i)$ denotes the dual annihilator of the homological boundaries $B_i(C) = \operatorname{im}\,\partial_{i+1}$.
\end{theorem}

\begin{proof}
\leanok
\uses{def:ChainComplexF2.cocycles, def:ChainComplexF2.coboundaryMap, def:ChainComplexF2.boundaries}
Unfolding the definitions $Z^i(C) = \ker \delta^i = \ker\bigl(\operatorname{dualMap}(\partial_{i+1})\bigr)$ and $B_i(C) = \operatorname{im}\,\partial_{i+1}$, the result follows directly from the identity $\ker(f^{\mathrm{tr}}) = (\operatorname{im}\,f)^{\mathrm{ann}}$, which holds over any commutative ring. Specifically, we apply \texttt{LinearMap.ker\_dualMap\_eq\_dualAnnihilator\_range} to the incoming differential $\partial_{i+1} = C.\operatorname{incomingDifferential}(i)$.
\end{proof}

\begin{theorem}[Coboundaries are the Dual Annihilator of Cycles]
\label{thm:ChainComplexF2.coboundaries_eq_dualAnnihilator_cycles}
\lean{ChainComplex𝔽₂.coboundaries_eq_dualAnnihilator_cycles}
\leanok
\uses{def:ChainComplexF2.coboundaries, def:ChainComplexF2.cycles}
Let $C$ be a chain complex over $\mathbb{F}_2$, and assume $C_i$ and $C_{i-1}$ are finite-dimensional over $\mathbb{F}_2$. For each $i \in \mathbb{Z}$:
\[
  B^i(C) \;=\; Z_i(C)^{\mathrm{ann}},
\]
where $Z_i(C)^{\mathrm{ann}} \subseteq \operatorname{Dual}(C_i)$ is the dual annihilator of the homological cycles $Z_i(C) = \ker\,\partial_i$.
\end{theorem}

\begin{proof}
\leanok
\uses{def:ChainComplexF2.coboundaries, def:ChainComplexF2.incomingCoboundaryMap, def:ChainComplexF2.cycles, def:ChainComplexF2.differential}
Unfolding the definitions $B^i(C) = \operatorname{im}\,\delta^{i-1} = \operatorname{im}\bigl(\operatorname{dualMap}(\partial_i)\bigr)$ and $Z_i(C) = \ker\,\partial_i$, the result follows from the identity $\operatorname{im}(f^{\mathrm{tr}}) = (\ker f)^{\mathrm{ann}}$ for linear maps between finite-dimensional vector spaces over a field. Specifically, we apply \texttt{LinearMap.range\_dualMap\_eq\_dualAnnihilator\_ker} to the differential $\partial_i = C.\operatorname{differential}(i)$. Finite-dimensionality of $C_i$ and $C_{i-1}$ is used to ensure the field-theoretic transpose identity holds.
\end{proof}

\begin{lemma}[Finrank of Boundaries Sub-Cycles Equals Finrank of Boundaries]
\label{lem:ChainComplexF2.finrank_boundariesSubCycles_eq}
\lean{ChainComplex𝔽₂.finrank_boundariesSubCycles_eq}
\leanok
\uses{def:ChainComplexF2.boundariesSubCycles, def:ChainComplexF2.boundaries, thm:ChainComplexF2.boundaries_le_cycles}
Let $C$ be a chain complex over $\mathbb{F}_2$ with $C_i$ and $C_{i+1}$ finite-dimensional. Then:
\[
  \operatorname{finrank}_{\mathbb{F}_2}\bigl(B_i(C)|_{Z_i}\bigr) \;=\; \operatorname{finrank}_{\mathbb{F}_2}\bigl(B_i(C)\bigr).
\]
\end{lemma}

\begin{proof}
\leanok
\uses{def:ChainComplexF2.boundariesSubCycles, def:ChainComplexF2.boundaries, thm:ChainComplexF2.boundaries_le_cycles}
Let $h_{\le}$ denote the inclusion $B_i(C) \subseteq Z_i(C)$ from \texttt{boundaries\_le\_cycles}. Unfolding $B_i(C)|_{Z_i} = \operatorname{comap}(\iota_{Z_i})(B_i(C))$, we apply the identity
\[
  \operatorname{finrank}\bigl(\operatorname{comap}(\iota_{Z_i})(B_i)\bigr) \;=\; \operatorname{finrank}\bigl(\operatorname{map}(\iota_{Z_i})\bigl(\operatorname{comap}(\iota_{Z_i})(B_i)\bigr)\bigr)
\]
using \texttt{Submodule.finrank\_map\_subtype\_eq}. Next, we apply \texttt{Submodule.map\_comap\_subtype} to obtain $\operatorname{map}(\iota_{Z_i})(\operatorname{comap}(\iota_{Z_i})(B_i)) = Z_i \sqcap B_i$. Since $B_i \subseteq Z_i$ (i.e., $h_{\le}$ holds), we conclude $Z_i \sqcap B_i = B_i$ by \texttt{inf\_eq\_right}, completing the proof.
\end{proof}

\begin{lemma}[Finrank of Coboundaries Sub-Cocycles Equals Finrank of Coboundaries]
\label{lem:ChainComplexF2.finrank_coboundariesSubCocycles_eq}
\lean{ChainComplex𝔽₂.finrank_coboundariesSubCocycles_eq}
\leanok
\uses{def:ChainComplexF2.coboundariesSubCocycles, def:ChainComplexF2.coboundaries, thm:ChainComplexF2.coboundaries_le_cocycles}
Let $C$ be a chain complex over $\mathbb{F}_2$ with $C_i$, $C_{i-1}$, and $C_{i+1}$ finite-dimensional. Then:
\[
  \operatorname{finrank}_{\mathbb{F}_2}\bigl(B^i(C)|_{Z^i}\bigr) \;=\; \operatorname{finrank}_{\mathbb{F}_2}\bigl(B^i(C)\bigr).
\]
\end{lemma}

\begin{proof}
\leanok
\uses{def:ChainComplexF2.coboundariesSubCocycles, def:ChainComplexF2.coboundaries, thm:ChainComplexF2.coboundaries_le_cocycles}
Let $h_{\le}$ denote the inclusion $B^i(C) \subseteq Z^i(C)$ from \texttt{coboundaries\_le\_cocycles}. Unfolding $B^i(C)|_{Z^i} = \operatorname{comap}(\iota_{Z^i})(B^i(C))$, we apply \texttt{Submodule.finrank\_map\_subtype\_eq} to obtain
\[
  \operatorname{finrank}\bigl(\operatorname{comap}(\iota_{Z^i})(B^i)\bigr) \;=\; \operatorname{finrank}\bigl(\operatorname{map}(\iota_{Z^i})(\operatorname{comap}(\iota_{Z^i})(B^i))\bigr).
\]
By \texttt{Submodule.map\_comap\_subtype}, this equals $\operatorname{finrank}(Z^i \sqcap B^i)$. Since $B^i \subseteq Z^i$, we have $Z^i \sqcap B^i = B^i$ by \texttt{inf\_eq\_right}, completing the proof.
\end{proof}

\begin{theorem}[Homology and Cohomology Have Equal Finrank]
\label{thm:ChainComplexF2.finrank_homology_eq_finrank_cohomology}
\lean{ChainComplex𝔽₂.finrank_homology_eq_finrank_cohomology}
\leanok
\uses{def:ChainComplexF2.boundariesSubCycles, def:ChainComplexF2.coboundariesSubCocycles, lem:ChainComplexF2.finrank_boundariesSubCycles_eq, lem:ChainComplexF2.finrank_coboundariesSubCocycles_eq, thm:ChainComplexF2.coboundaries_eq_dualAnnihilator_cycles, thm:ChainComplexF2.cocycles_eq_dualAnnihilator_boundaries}
Let $C$ be a chain complex over $\mathbb{F}_2$ with $C_i$, $C_{i-1}$, and $C_{i+1}$ finite-dimensional. Then:
\[
  \operatorname{finrank}_{\mathbb{F}_2}\bigl(H_i(C)\bigr) \;=\; \operatorname{finrank}_{\mathbb{F}_2}\bigl(H^i(C)\bigr),
\]
where $H_i(C) = Z_i(C)/B_i(C)$ is homology and $H^i(C) = Z^i(C)/B^i(C)$ is cohomology.
\end{theorem}

\begin{proof}
\leanok
\uses{lem:ChainComplexF2.finrank_boundariesSubCycles_eq, lem:ChainComplexF2.finrank_coboundariesSubCocycles_eq, thm:ChainComplexF2.coboundaries_eq_dualAnnihilator_cycles, thm:ChainComplexF2.cocycles_eq_dualAnnihilator_boundaries}
We apply the dimension formula $\operatorname{finrank}(M/N) + \operatorname{finrank}(N) = \operatorname{finrank}(M)$ twice:
\begin{align*}
  \operatorname{finrank}(H_i) + \operatorname{finrank}(B_i|_{Z_i}) &= \operatorname{finrank}(Z_i), \\
  \operatorname{finrank}(H^i) + \operatorname{finrank}(B^i|_{Z^i}) &= \operatorname{finrank}(Z^i).
\end{align*}
By \texttt{finrank\_boundariesSubCycles\_eq}, $\operatorname{finrank}(B_i|_{Z_i}) = \operatorname{finrank}(B_i)$. By \texttt{finrank\_coboundariesSubCocycles\_eq}, $\operatorname{finrank}(B^i|_{Z^i}) = \operatorname{finrank}(B^i)$.

We then apply the dual annihilator dimension formula $\operatorname{finrank}(W) + \operatorname{finrank}(W^{\mathrm{ann}}) = \operatorname{finrank}(V)$ (for $W \subseteq V$ finite-dimensional) to both $B_i(C)$ and $Z_i(C)$:
\begin{align*}
  \operatorname{finrank}(B_i) + \operatorname{finrank}(B_i^{\mathrm{ann}}) &= \operatorname{finrank}(C_i), \\
  \operatorname{finrank}(Z_i) + \operatorname{finrank}(Z_i^{\mathrm{ann}}) &= \operatorname{finrank}(C_i).
\end{align*}
By \texttt{cocycles\_eq\_dualAnnihilator\_boundaries}, $\operatorname{finrank}(Z^i) = \operatorname{finrank}(B_i^{\mathrm{ann}})$. By \texttt{coboundaries\_eq\_dualAnnihilator\_cycles}, $\operatorname{finrank}(B^i) = \operatorname{finrank}(Z_i^{\mathrm{ann}})$. Combining all four equations, the equality $\operatorname{finrank}(H_i) = \operatorname{finrank}(H^i)$ follows by integer arithmetic (\texttt{omega}).
\end{proof}

\begin{definition}[Homology--Cohomology Isomorphism]
\label{def:ChainComplexF2.homologyCohomologyEquiv}
\lean{ChainComplex𝔽₂.homologyCohomologyEquiv}
\leanok
\uses{thm:ChainComplexF2.finrank_homology_eq_finrank_cohomology}
Let $C$ be a chain complex over $\mathbb{F}_2$ with $C_i$, $C_{i-1}$, and $C_{i+1}$ finite-dimensional. There is a linear isomorphism of $\mathbb{F}_2$-vector spaces:
\[
  H_i(C) \;\simeq_{\mathbb{F}_2}\; H^i(C).
\]
This isomorphism is constructed via \texttt{LinearEquiv.ofFinrankEq}, using the fact that both spaces are finite-dimensional $\mathbb{F}_2$-vector spaces of equal dimension, as established by \texttt{finrank\_homology\_eq\_finrank\_cohomology}.
\end{definition}

%--- Def_3: ClassicalCode ---
\chapter{Def 3: Classical Code}

\begin{definition}[Classical Linear Binary Code]
\label{def:ClassicalCode}
\lean{ClassicalCode}
\leanok
\uses{def:F2}
A \emph{classical linear binary code} of block length $n$ is a subspace $\mathcal{C} \subseteq \mathbb{F}_2^n$. Formally, it is a structure wrapping a submodule
\[
  \mathcal{C} \leq \mathbb{F}_2^n,
\]
where $\mathbb{F}_2^n$ is viewed as the $\mathbb{F}_2$-module of functions $\operatorname{Fin}(n) \to \mathbb{F}_2$.
\end{definition}

\begin{definition}[Code Dimension]
\label{def:ClassicalCode.dimension}
\lean{ClassicalCode.dimension}
\leanok
\uses{def:ClassicalCode, def:F2}
Let $\mathcal{C}$ be a classical code of block length $n$. The \emph{dimension} of $\mathcal{C}$ is
\[
  k \;=\; \dim_{\mathbb{F}_2} \mathcal{C} \;=\; \operatorname{finrank}_{\mathbb{F}_2}(\mathcal{C}),
\]
the finite rank of $\mathcal{C}$ as an $\mathbb{F}_2$-module. This equals the number of encodable bits.
\end{definition}

\begin{definition}[Code Distance]
\label{def:ClassicalCode.distance}
\lean{ClassicalCode.distance}
\leanok
\uses{def:ClassicalCode, def:hammingWeight}
Let $\mathcal{C}$ be a classical code of block length $n$. The \emph{distance} of $\mathcal{C}$ is
\[
  d \;=\; \inf \bigl\{\, w \in \mathbb{N} \;\big|\; \exists\, x \in \mathcal{C},\; x \neq 0,\; \operatorname{wt}(x) = w \,\bigr\},
\]
where $\operatorname{wt}(x)$ denotes the Hamming weight of $x$. Here $\inf$ is taken over $\mathbb{N}$, and by convention $\inf \emptyset = 0$, which gives $d = 0$ for the trivial code $\mathcal{C} = \{0\}$.
\end{definition}

\begin{definition}[{$[n,k,d]$-Code}]
\label{def:ClassicalCode.IsNKDCode}
\lean{ClassicalCode.IsNKDCode}
\leanok
\uses{def:ClassicalCode, def:ClassicalCode.dimension, def:ClassicalCode.distance}
A classical code $\mathcal{C}$ of block length $n$ is called an $[n, k, d]$\emph{-code} if its dimension equals $k$ and its distance equals $d$:
\[
  \mathcal{C}.\operatorname{dimension} = k \;\wedge\; \mathcal{C}.\operatorname{distance} = d.
\]
\end{definition}

\begin{definition}[Code from Parity Check]
\label{def:ClassicalCode.ofParityCheck}
\lean{ClassicalCode.ofParityCheck}
\leanok
\uses{def:ClassicalCode, def:F2}
Given an $\mathbb{F}_2$-linear map $\partial^C : \mathbb{F}_2^n \to \mathbb{F}_2^m$ (the \emph{parity check map}), we construct a classical code of block length $n$ by taking the kernel:
\[
  \mathcal{C} \;=\; \ker(\partial^C) \;\leq\; \mathbb{F}_2^n.
\]
\end{definition}

\begin{definition}[Parity Check Chain Complex]
\label{def:ClassicalCode.parityCheckComplex}
\lean{ClassicalCode.parityCheckComplex}
\leanok
\uses{def:ClassicalCode.ofParityCheck, def:ChainComplexF2, def:F2}
Given an $\mathbb{F}_2$-linear parity check map $\partial^C : \mathbb{F}_2^n \to \mathbb{F}_2^m$, the associated \emph{two-term chain complex} is
\[
  C \;=\; \bigl(C_1 \xrightarrow{\partial^C} C_0\bigr),
\]
where $C_1 = \mathbb{F}_2^n$ sits in degree $1$ and $C_0 = \mathbb{F}_2^m$ sits in degree $0$. This is constructed via $\operatorname{HomologicalComplex.double}$ applied to the morphism $\partial^C$ in the category of $\mathbb{F}_2$-modules, using the complex shape relation $1 \to 0$.
\end{definition}

\begin{definition}[Parity Check Complex Isomorphism at Degree 1]
\label{def:ClassicalCode.parityCheckComplexXIso1}
\lean{ClassicalCode.parityCheckComplexXIso₁}
\leanok
\uses{def:ClassicalCode.parityCheckComplex, def:F2}
For a parity check map $\partial^C : \mathbb{F}_2^n \to \mathbb{F}_2^m$, there is a canonical isomorphism
\[
  \varphi \;:\; C(1) \;\xrightarrow{\;\sim\;}\; \mathbb{F}_2^n
\]
between the degree-$1$ object of the parity check complex $C$ and $\mathbb{F}_2^n$, given by $\operatorname{HomologicalComplex.doubleXIso_0}$.
\end{definition}

\begin{theorem}[Code Equals Kernel]
\label{thm:ClassicalCode.code_eq_ker}
\lean{ClassicalCode.code_eq_ker}
\leanok
\uses{def:ClassicalCode.ofParityCheck}
For any $\mathbb{F}_2$-linear parity check map $\partial^C : \mathbb{F}_2^n \to \mathbb{F}_2^m$, the underlying submodule of the code constructed from $\partial^C$ equals the kernel of $\partial^C$:
\[
  (\operatorname{ofParityCheck}(\partial^C)).\operatorname{code} \;=\; \ker(\partial^C).
\]
\end{theorem}

\begin{proof}
\leanok
\uses{def:ClassicalCode.ofParityCheck}
This holds by reflexivity: $\operatorname{ofParityCheck}(\partial^C).\operatorname{code}$ is defined to be $\ker(\partial^C)$, so the equality is definitional.
\end{proof}

\begin{theorem}[Code Equals Cycles]
\label{thm:ClassicalCode.code_eq_cycles_map}
\lean{ClassicalCode.code_eq_cycles_map}
\leanok
\uses{def:ClassicalCode.ofParityCheck, def:ClassicalCode.parityCheckComplex, def:ClassicalCode.parityCheckComplexXIso1, def:ChainComplexF2.cycles, def:ChainComplexF2.differential}
For any $\mathbb{F}_2$-linear parity check map $\partial^C : \mathbb{F}_2^n \to \mathbb{F}_2^m$, the code $\mathcal{C} = \ker(\partial^C)$ equals the image of the degree-$1$ cycles of the parity check complex $C$ under the canonical isomorphism $\varphi : C(1) \xrightarrow{\sim} \mathbb{F}_2^n$:
\[
  \mathcal{C} \;=\; \varphi\bigl(Z_1(C)\bigr),
\]
where $Z_1(C) = \ker(\partial^C : C(1) \to C(0))$ is the submodule of $1$-cycles.
\end{theorem}

\begin{proof}
\leanok
\uses{def:ClassicalCode.ofParityCheck, def:ClassicalCode.parityCheckComplex, def:ClassicalCode.parityCheckComplexXIso1, def:ChainComplexF2.cycles, def:ChainComplexF2.differential}
Let $C = \operatorname{parityCheckComplex}(\partial^C)$, let $\varphi = \operatorname{parityCheckComplexXIso_1}(\partial^C) : C(1) \xrightarrow{\sim} \mathbb{F}_2^n$, and let $\psi = \operatorname{HomologicalComplex.doubleXIso_1}(\partial^C) : C(0) \xrightarrow{\sim} \mathbb{F}_2^m$. We note the factorization
\[
  \partial^C_{\,C(1) \to C(0)} \;=\; \varphi^{-1} \circ \partial^C \circ \psi^{-1},
\]
established by $\operatorname{HomologicalComplex.double_d}$.

By extensionality, it suffices to show that for each $x \in \mathbb{F}_2^n$:
\[
  x \in \ker(\partial^C) \;\Longleftrightarrow\; \exists\, y \in Z_1(C),\; \varphi(y) = x.
\]

$(\Rightarrow)$: Assume $\partial^C(x) = 0$. We produce the preimage $y = \varphi^{-1}(x) \in C(1)$. Using the factorization of the differential, we compute:
\[
  d(y) = \psi^{-1}(\partial^C(\varphi(\varphi^{-1}(x)))) = \psi^{-1}(\partial^C(x)) = \psi^{-1}(0) = 0,
\]
where we used $\varphi \circ \varphi^{-1} = \operatorname{id}$ (i.e., $(\varphi^{-1} \circ \varphi)(x) = x$ by the isomorphism identity $\varphi^{-1} \circ \varphi = \operatorname{id}$, and $\psi^{-1}(0) = 0$ since $\psi^{-1}$ is linear). Thus $y \in Z_1(C)$, and $\varphi(y) = \varphi(\varphi^{-1}(x)) = x$ by the isomorphism identity.

$(\Leftarrow)$: Assume $y \in Z_1(C)$ with $x = \varphi(y)$. By simplification using the definition of cycles and the differential, $y \in \ker(d_{1,0})$, meaning $d_{1,0}(y) = 0$. Using the factorization $d_{1,0} = \varphi^{-1} \circ \partial^C \circ \psi^{-1}$ and substituting, we obtain
\[
  \psi^{-1}(\partial^C(\varphi(y))) = 0.
\]
We show $\psi^{-1}$ is injective: given $a, b$ with $\psi^{-1}(a) = \psi^{-1}(b)$, applying $\psi$ to both sides and using $\psi \circ \psi^{-1} = \operatorname{id}$ yields $a = b$. By injectivity of $\psi^{-1}$, from $\psi^{-1}(\partial^C(\varphi(y))) = \psi^{-1}(0)$ we conclude $\partial^C(\varphi(y)) = 0$, i.e., $\partial^C(x) = 0$. Thus $x \in \ker(\partial^C) = \mathcal{C}$.
\end{proof}

%--- Def_4: CSSCode ---
\chapter{Def 4: CSS Code}

\begin{definition}[CSS Code]
\label{def:CSSCode}
\lean{CSSCode}
\leanok
\uses{def:F2}
A \emph{CSS (Calderbank--Shor--Steane) quantum code} of parameters $(n, r_X, r_Z)$ consists of:
\begin{itemize}
  \item An \emph{X-type parity check matrix} $H_X : \mathbb{F}_2^n \to \mathbb{F}_2^{r_X}$, a linear map over $\mathbb{F}_2$.
  \item A \emph{Z-type parity check matrix transpose} $H_Z^T : \mathbb{F}_2^{r_Z} \to \mathbb{F}_2^n$, a linear map over $\mathbb{F}_2$.
  \item The \emph{CSS condition}: $H_X \circ H_Z^T = 0$.
\end{itemize}
The map $H_Z^T$ serves as the differential from degree $1$ to degree $0$ in the associated chain complex $C_1 \xrightarrow{H_Z^T} C_0 \xrightarrow{H_X} C_{-1}$.
\end{definition}

\begin{definition}[Three-term chain complex of a CSS code]
\label{def:CSSCode.complex}
\lean{CSSCode.complex}
\leanok
\uses{def:CSSCode, def:ChainComplexF2}
Let $Q = (H_X, H_Z^T)$ be a CSS code with parameters $(n, r_X, r_Z)$. The \emph{associated chain complex} $\operatorname{complex}(Q)$ is the three-term chain complex in $\operatorname{Mod}_{\mathbb{F}_2}$ of the form
\[
  C_1 \xrightarrow{\,H_Z^T\,} C_0 \xrightarrow{\,H_X\,} C_{-1},
\]
where the objects are
\[
  C_i = \begin{cases} \mathbb{F}_2^{r_Z} & \text{if } i = 1, \\ \mathbb{F}_2^n & \text{if } i = 0, \\ \mathbb{F}_2^{r_X} & \text{if } i = -1, \\ 0 & \text{otherwise,} \end{cases}
\]
and the differentials are $d_{1,0} = H_Z^T$ and $d_{0,-1} = H_X$, with all other differentials zero. This is well-defined as a chain complex because $\partial^2 = 0$, which follows from the CSS condition $H_X \circ H_Z^T = 0$.
\end{definition}

\begin{proof}
\leanok
\uses{def:CSSCode, def:ChainComplexF2}
We must verify that $d \circ d = 0$, i.e., for all integers $i, j, k$ with $j + 1 = i$ and $k + 1 = j$, the composition $d_{i,j} \circ d_{j,k} = 0$.

For degrees $i, j, k$ with $j + 1 = i$ and $k + 1 = j$, the only non-trivial case occurs when $i = 1$, $j = 0$, and $k = -1$ (since all other differentials are zero by definition). In this case, we must show
\[
  d_{1,0} \circ d_{0,-1} = H_Z^T \circ H_X = 0.
\]

We expand the composition: rewriting using $i = 1$, the first differential is $d_{1,0} = \operatorname{eqToHom} \circ H_Z^T \circ \operatorname{eqToHom}$, and the second is $d_{0,-1} = \operatorname{eqToHom} \circ H_X \circ \operatorname{eqToHom}$. Applying associativity of composition and composing the adjacent $\operatorname{eqToHom}$ morphisms (which cancel to the identity since they are inverse transport maps), we reduce to showing
\[
  H_Z^T \circ H_X = 0.
\]
We have $\operatorname{ofHom}(H_Z^T) \circ \operatorname{ofHom}(H_X) = \operatorname{ofHom}(H_Z^T \circ H_X)$, and by the CSS condition $Q.\mathtt{css\_condition}$, we have $H_X \circ H_Z^T = 0$, which gives $H_Z^T \circ H_X = 0$ after noting the composition order. Thus the composition is $0$.

For all other combinations of $i, j, k$: if $i = 1$ but $j \neq 0$, this contradicts $j + 1 = 1$, hence $j = 0$ by integer arithmetic. If $i \neq 1$ and $i = 0$, $j = -1$, then $k = -2$, and $d_{-1,-2} = 0$ by the shape condition (neither $-1 \to 0$ nor $-1 \to -2$ are non-trivial differentials), so the composition is $0$. In all remaining cases, $d_{i,j} = 0$, making the composition immediately $0$.
\end{proof}

\begin{definition}[Number of physical qubits]
\label{def:CSSCode.numQubits}
\lean{CSSCode.numQubits}
\leanok
\uses{def:CSSCode}
Let $Q$ be a CSS code with parameters $(n, r_X, r_Z)$. The \emph{number of physical qubits} is
\[
  n(Q) := n = \dim C_0.
\]
\end{definition}

\begin{definition}[Number of logical qubits]
\label{def:CSSCode.logicalQubits}
\lean{CSSCode.logicalQubits}
\leanok
\uses{def:CSSCode, def:CSSCode.complex, def:ChainComplexF2}
Let $Q = (H_X, H_Z^T)$ be a CSS code. The \emph{number of logical qubits} is the dimension of the degree-$0$ homology of the chain complex:
\[
  k(Q) := \dim_{\mathbb{F}_2} H_0(Q) = \dim_{\mathbb{F}_2}\!\left(\ker H_X \,\big/\, \operatorname{im}(H_Z^T)\big|_{\ker H_X}\right).
\]
More precisely, $k(Q) = \operatorname{finrank}_{\mathbb{F}_2}\bigl(\ker H_X \,/\, (\operatorname{range}(H_Z^T))^{-1}(\ker H_X)\bigr)$, where the quotient is taken of the subspace of $\ker H_X$ by the image of $H_Z^T$ within $\ker H_X$.
\end{definition}

\begin{definition}[X-distance]
\label{def:CSSCode.dX}
\lean{CSSCode.dX}
\leanok
\uses{def:CSSCode, def:hammingWeight}
Let $Q = (H_X, H_Z^T)$ be a CSS code. The \emph{X-distance} is the minimum Hamming weight of a vector in $\ker H_X$ that does not lie in $\operatorname{im}(H_Z^T)$:
\[
  d_X(Q) := \inf\bigl\{\, \operatorname{wt}(x) \mid x \in \ker H_X,\; x \notin \operatorname{im}(H_Z^T) \,\bigr\},
\]
where $\operatorname{wt}(x) = |\operatorname{supp}(x)|$ is the Hamming weight. By convention, $d_X = 0$ when $\ker H_X = \operatorname{im}(H_Z^T)$ (trivial homology), since the infimum of the empty set is $0$.
\end{definition}

\begin{definition}[Z-distance]
\label{def:CSSCode.dZ}
\lean{CSSCode.dZ}
\leanok
\uses{def:CSSCode, def:hammingWeight}
Let $Q = (H_X, H_Z^T)$ be a CSS code. The \emph{Z-distance} is defined using the dual maps. Since $H_Z = (H_Z^T)^T$, the kernel of $H_Z$ corresponds to the kernel of the dual map $(H_Z^T)^\vee : ({\mathbb{F}_2^n})^\vee \to ({\mathbb{F}_2^{r_Z}})^\vee$, and the image of $H_X^T$ corresponds to the range of the dual map $H_X^\vee : ({\mathbb{F}_2^{r_X}})^\vee \to ({\mathbb{F}_2^n})^\vee$. Over $\mathbb{F}_2$ with canonical bases, $(\mathbb{F}_2^n)^\vee \cong \mathbb{F}_2^n$ via the dot product equivalence. Using this identification, the \emph{Z-distance} is
\[
  d_Z(Q) := \inf\bigl\{\, \operatorname{wt}(\varphi) \mid \varphi \in \ker\bigl((H_Z^T)^\vee\bigr),\; \varphi \notin \operatorname{im}(H_X^\vee) \,\bigr\},
\]
where $\operatorname{wt}(\varphi)$ is the Hamming weight of $\varphi$ viewed as a vector in $\mathbb{F}_2^n$ via the dot product isomorphism. By convention, $d_Z = 0$ when the cohomology is trivial.
\end{definition}

\begin{definition}[CSS code distance]
\label{def:CSSCode.distance}
\lean{CSSCode.distance}
\leanok
\uses{def:CSSCode, def:CSSCode.dX, def:CSSCode.dZ}
Let $Q$ be a CSS code. The \emph{overall distance} is
\[
  d(Q) := \min\bigl(d_X(Q),\, d_Z(Q)\bigr).
\]
\end{definition}

\begin{definition}[{$[\![n,k,d]\!]$-code predicate}]
\label{def:CSSCode.IsNKDCode}
\lean{CSSCode.IsNKDCode}
\leanok
\uses{def:CSSCode, def:CSSCode.logicalQubits, def:CSSCode.distance}
Let $Q$ be a CSS code and let $k, d \in \mathbb{N}$. We say $Q$ is an \emph{$[\![n,k,d]\!]$-code} if
\[
  k(Q) = k \quad \text{and} \quad d(Q) = d,
\]
i.e., $Q$ encodes exactly $k$ logical qubits with overall distance $d$.
\end{definition}

\begin{definition}[{$[\![n,k,d_X,d_Z]\!]$-code predicate}]
\label{def:CSSCode.IsNKDXZCode}
\lean{CSSCode.IsNKDXZCode}
\leanok
\uses{def:CSSCode, def:CSSCode.logicalQubits, def:CSSCode.dX, def:CSSCode.dZ}
Let $Q$ be a CSS code and let $k, d_X, d_Z \in \mathbb{N}$. We say $Q$ is an \emph{$[\![n,k,d_X,d_Z]\!]$-code} if
\[
  k(Q) = k, \quad d_X(Q) = d_X, \quad \text{and} \quad d_Z(Q) = d_Z.
\]
This predicate is used when the X- and Z-distances are distinct and one wishes to record both separately.
\end{definition}

%--- Def_5: LDPCCode ---
\chapter{Def 5: LDPC Code}

\begin{definition}[Row Weight]
\label{def:rowWeight}
\lean{rowWeight}
\leanok
\uses{def:hammingWeight, def:F2}
For a linear map $f : (\operatorname{Fin} n \to \mathbb{F}_2) \to_{\mathbb{F}_2} (\operatorname{Fin} m \to \mathbb{F}_2)$ and an index $i : \operatorname{Fin} m$, the \emph{row weight} of row $i$ of the matrix representation of $f$ is defined as
\[
\operatorname{rowWeight}(f, i) := \operatorname{hammingWeight}\!\left(\lambda\, j : \operatorname{Fin} n,\ f(\operatorname{Pi.single}\ j\ 1)\ i\right),
\]
i.e., the number of column indices $j$ such that the $(i,j)$-entry of the matrix is nonzero.
\end{definition}

\begin{definition}[Column Weight]
\label{def:colWeight}
\lean{colWeight}
\leanok
\uses{def:hammingWeight, def:F2}
For a linear map $f : (\operatorname{Fin} n \to \mathbb{F}_2) \to_{\mathbb{F}_2} (\operatorname{Fin} m \to \mathbb{F}_2)$ and an index $j : \operatorname{Fin} n$, the \emph{column weight} of column $j$ of the matrix representation of $f$ is defined as
\[
\operatorname{colWeight}(f, j) := \operatorname{hammingWeight}\!\left(f(\operatorname{Pi.single}\ j\ 1)\right),
\]
i.e., the number of row indices $i$ such that the $(i,j)$-entry of the matrix is nonzero.
\end{definition}

\begin{definition}[Bounded Weight]
\label{def:HasBoundedWeight}
\lean{HasBoundedWeight}
\leanok
\uses{def:rowWeight, def:colWeight, def:F2}
A linear map $f : (\operatorname{Fin} n \to \mathbb{F}_2) \to_{\mathbb{F}_2} (\operatorname{Fin} m \to \mathbb{F}_2)$ is said to have \emph{bounded weight $w$} (for some $w \in \mathbb{N}$) if every row and every column of its matrix representation has Hamming weight at most $w$:
\[
\operatorname{HasBoundedWeight}(f, w) \;:=\; \left(\forall\, i : \operatorname{Fin} m,\ \operatorname{rowWeight}(f, i) \leq w\right) \;\land\; \left(\forall\, j : \operatorname{Fin} n,\ \operatorname{colWeight}(f, j) \leq w\right).
\]
\end{definition}

\begin{definition}[LDPC Code Family]
\label{def:CSSCode.IsLDPC}
\lean{CSSCode.IsLDPC}
\leanok
\uses{def:CSSCode, def:HasBoundedWeight}
A family of CSS codes indexed by a type $\iota$, say $\{C_\alpha\}_{\alpha \in \iota}$ where each $C_\alpha$ is a CSS code with parameters $(n(\alpha),\, r_X(\alpha),\, r_Z(\alpha))$, is called \emph{low-density parity-check (LDPC)} if there exists a uniform bound $w \in \mathbb{N}$ such that for every $\alpha \in \iota$, both the parity-check matrix $H_X$ and the transposed parity-check matrix $H_Z^T$ (as stored in the CSS code structure) have bounded weight $w$:
\[
\operatorname{IsLDPC}(\{C_\alpha\}) \;:=\; \exists\, w \in \mathbb{N},\ \forall\, \alpha \in \iota,\ \operatorname{HasBoundedWeight}((C_\alpha).H_X,\, w) \;\land\; \operatorname{HasBoundedWeight}((C_\alpha).H_Z^T,\, w).
\]

\textit{Note: Since the CSS code structure stores $H_Z^T$ rather than $H_Z$ directly, and since the weight bound $w$ applies symmetrically to both rows and columns, the condition $\operatorname{HasBoundedWeight}(H_Z^T, w)$ is equivalent to requiring that $H_Z$ itself has all row weights and column weights bounded by $w$.}
\end{definition}

%--- Def_6: SubsystemCSSCode ---
\chapter{Def 6: Subsystem CSS Code}

\begin{definition}[Boundaries in Cycles]
\label{def:CSSCode.boundariesInCycles}
\lean{CSSCode.boundariesInCycles}
\leanok
\uses{def:CSSCode}
Let $Q$ be a CSS code with parity check matrices $H_X$ and $H_Z^T$. The submodule of \emph{boundaries inside cycles} is
\[
  \operatorname{im}(H_Z^T) \cap \ker(H_X),
\]
viewed as a submodule of $\ker(H_X)$ via the comap of the subtype inclusion. Formally, it is the $\mathbb{F}_2$-submodule
\[
  (\operatorname{im}(H_Z^T)).\operatorname{comap}(\ker(H_X).\operatorname{subtype}) \;\subseteq\; \ker(H_X).
\]
\end{definition}

\begin{definition}[Homology Type $H_0$]
\label{def:CSSCode.Homology0}
\lean{CSSCode.Homology₀}
\leanok
\uses{def:CSSCode, def:CSSCode.boundariesInCycles}
The \emph{homology type} of a CSS code $Q$ is the quotient $\mathbb{F}_2$-vector space
\[
  H_0 \;=\; \ker(H_X) \,/\, \operatorname{im}(H_Z^T),
\]
i.e., the quotient of $\ker(H_X)$ by the submodule of boundaries in cycles.
\end{definition}

\begin{definition}[Coboundaries in Cocycles]
\label{def:CSSCode.coboundariesInCocycles}
\lean{CSSCode.coboundariesInCocycles}
\leanok
\uses{def:CSSCode}
The submodule of \emph{coboundaries inside cocycles} is
\[
  \operatorname{im}(H_X^T) \cap \ker(H_Z),
\]
where $H_Z = (H_Z^T)^T = \operatorname{dualMap}(H_Z^T)$ and $H_X^T = \operatorname{dualMap}(H_X)$, viewed as a submodule of $\ker(H_Z)$. Formally, it is
\[
  (\operatorname{im}(\operatorname{dualMap}(H_X))).\operatorname{comap}(\ker(\operatorname{dualMap}(H_Z^T)).\operatorname{subtype}).
\]
\end{definition}

\begin{definition}[Cohomology Type $H^0$]
\label{def:CSSCode.Cohomology0}
\lean{CSSCode.Cohomology₀}
\leanok
\uses{def:CSSCode, def:CSSCode.coboundariesInCocycles}
The \emph{cohomology type} of a CSS code $Q$ is the quotient $\mathbb{F}_2$-vector space
\[
  H^0 \;=\; \ker(H_Z) \,/\, \operatorname{im}(H_X^T),
\]
i.e., the quotient of $\ker(\operatorname{dualMap}(H_Z^T))$ by the submodule of coboundaries in cocycles.
\end{definition}

\begin{definition}[Homology Quotient Map]
\label{def:CSSCode.homologyMkQ}
\lean{CSSCode.homologyMkQ}
\leanok
\uses{def:CSSCode.Homology0, def:CSSCode.boundariesInCycles}
The \emph{homology quotient map} is the canonical $\mathbb{F}_2$-linear surjection
\[
  \pi_H \;:\; \ker(H_X) \;\longrightarrow\; H_0,
\]
defined as the module quotient map $\operatorname{mkQ}$ for the submodule of boundaries in cycles.
\end{definition}

\begin{definition}[Cohomology Quotient Map]
\label{def:CSSCode.cohomologyMkQ}
\lean{CSSCode.cohomologyMkQ}
\leanok
\uses{def:CSSCode.Cohomology0, def:CSSCode.coboundariesInCocycles}
The \emph{cohomology quotient map} is the canonical $\mathbb{F}_2$-linear surjection
\[
  \pi_{H^0} \;:\; \ker(H_Z) \;\longrightarrow\; H^0,
\]
defined as the module quotient map $\operatorname{mkQ}$ for the submodule of coboundaries in cocycles.
\end{definition}

\begin{definition}[Subsystem CSS Code]
\label{def:SubsystemCSSCode}
\lean{SubsystemCSSCode}
\leanok
\uses{def:CSSCode}
A \emph{subsystem CSS code} with parameters $n, r_X, r_Z \in \mathbb{N}$ is a CSS code $Q$ (Definition~\ref{def:CSSCode}) together with the following data:
\begin{itemize}
  \item A \emph{logical subspace} $H_0^L \subseteq H_0$: an $\mathbb{F}_2$-submodule of the homology $H_0 = \ker(H_X)/\operatorname{im}(H_Z^T)$.
  \item A \emph{gauge subspace} $H_0^G \subseteq H_0$: an $\mathbb{F}_2$-submodule of the homology.
  \item A \emph{direct sum decomposition} $H_0 = H_0^L \oplus H_0^G$: the pair $(H_0^L, H_0^G)$ satisfies $\operatorname{IsCompl}(H_0^L, H_0^G)$, meaning $H_0^L \cap H_0^G = 0$ and $H_0^L + H_0^G = H_0$.
  \item A \emph{linear equivalence} $\phi : H_0 \xrightarrow{\;\sim\;} H^0$ between the homology and cohomology, encoding the nondegenerate pairing. The cohomology splitting is induced by $\phi$: $H^0_L = \phi(H_0^L)$ and $H^0_G = \phi(H_0^G)$, ensuring that $H_0^L$ pairs nondegenerately with $H^0_L$, $H_0^G$ pairs nondegenerately with $H^0_G$, and the cross-pairings vanish.
\end{itemize}
\end{definition}

\begin{definition}[Logical Cohomology Subspace]
\label{def:SubsystemCSSCode.coHL}
\lean{SubsystemCSSCode.coHL}
\leanok
\uses{def:SubsystemCSSCode}
Let $S$ be a subsystem CSS code with equivalence $\phi : H_0 \xrightarrow{\;\sim\;} H^0$. The \emph{logical cohomology subspace} is
\[
  H^0_L \;:=\; \phi(H_0^L) \;=\; \operatorname{Submodule.map}(\phi, H_0^L) \;\subseteq\; H^0.
\]
\end{definition}

\begin{definition}[Gauge Cohomology Subspace]
\label{def:SubsystemCSSCode.coHG}
\lean{SubsystemCSSCode.coHG}
\leanok
\uses{def:SubsystemCSSCode}
Let $S$ be a subsystem CSS code with equivalence $\phi : H_0 \xrightarrow{\;\sim\;} H^0$. The \emph{gauge cohomology subspace} is
\[
  H^0_G \;:=\; \phi(H_0^G) \;=\; \operatorname{Submodule.map}(\phi, H_0^G) \;\subseteq\; H^0.
\]
\end{definition}

\begin{theorem}[Cohomology Direct Sum Decomposition]
\label{thm:SubsystemCSSCode.cohcompl}
\lean{SubsystemCSSCode.cohcompl}
\leanok
\uses{def:SubsystemCSSCode, def:SubsystemCSSCode.coHL, def:SubsystemCSSCode.coHG}
Let $S$ be a subsystem CSS code. The cohomology splitting $H^0 = H^0_L \oplus H^0_G$ holds, i.e., $\operatorname{IsCompl}(H^0_L, H^0_G)$.
\end{theorem}

\begin{proof}
\leanok
\uses{def:SubsystemCSSCode, def:SubsystemCSSCode.coHL, def:SubsystemCSSCode.coHG}
We prove both conditions of $\operatorname{IsCompl}(H^0_L, H^0_G)$ separately.

\textbf{Disjointness} ($H^0_L \cap H^0_G = 0$): We use the definition of $\operatorname{Submodule.disjoint_def}$. Let $x \in H^0_L \cap H^0_G$. Since $x \in H^0_L = \phi(H_0^L)$, there exists $a \in H_0^L$ with $\phi(a) = x$. Since $x \in H^0_G = \phi(H_0^G)$, there exists $b \in H_0^G$ with $\phi(b) = x$. By injectivity of $\phi$, we have $a = b$. Therefore $a \in H_0^L \cap H_0^G$. Since $H_0 = H_0^L \oplus H_0^G$ (i.e., $\operatorname{IsCompl}(H_0^L, H_0^G)$), we have $H_0^L \cap H_0^G = 0$, so $a = 0$. Hence $x = \phi(a) = \phi(0) = 0$ by the linear map property, and $x = 0$ follows by simplification.

\textbf{Codisjointness} ($H^0_L + H^0_G = H^0$): Unfolding $H^0_L = \phi(H_0^L)$ and $H^0_G = \phi(H_0^G)$, we rewrite using codisjointness and compute
\[
  H^0_L + H^0_G = \phi(H_0^L) + \phi(H_0^G) = \phi(H_0^L + H_0^G).
\]
Since $H_0 = H_0^L \oplus H_0^G$, we have $H_0^L + H_0^G = H_0$ (i.e., $\operatorname{sup_eq_top}$). Therefore $\phi(H_0) = \operatorname{Submodule.map_top}(\phi)$. By surjectivity of $\phi$, $\operatorname{LinearMap.range_eq_top}$ gives $\operatorname{range}(\phi) = H^0$, so $H^0_L + H^0_G = H^0$.
\end{proof}

\begin{definition}[Logical Projection]
\label{def:SubsystemCSSCode.projL}
\lean{SubsystemCSSCode.projL}
\leanok
\uses{def:SubsystemCSSCode}
The \emph{logical projection} is the $\mathbb{F}_2$-linear map
\[
  \pi_L \;:\; H_0 \;\longrightarrow\; H_0^L,
\]
defined as the linear projection onto $H_0^L$ along $H_0^G$, using the direct sum decomposition $H_0 = H_0^L \oplus H_0^G$ (via $\operatorname{Submodule.linearProjOfIsCompl}$).
\end{definition}

\begin{definition}[Logical Cohomology Projection]
\label{def:SubsystemCSSCode.coprojL}
\lean{SubsystemCSSCode.coprojL}
\leanok
\uses{def:SubsystemCSSCode, def:SubsystemCSSCode.coHL, thm:SubsystemCSSCode.cohcompl}
The \emph{logical cohomology projection} is the $\mathbb{F}_2$-linear map
\[
  \pi_{L}^* \;:\; H^0 \;\longrightarrow\; H^0_L,
\]
defined as the linear projection onto $H^0_L$ along $H^0_G$, using the induced cohomology direct sum decomposition $H^0 = H^0_L \oplus H^0_G$ (via $\operatorname{Submodule.linearProjOfIsCompl}$ and $\operatorname{cohcompl}$).
\end{definition}

\begin{definition}[Number of Logical Qubits]
\label{def:SubsystemCSSCode.logicalQubits}
\lean{SubsystemCSSCode.logicalQubits}
\leanok
\uses{def:SubsystemCSSCode}
The \emph{number of logical qubits} of a subsystem CSS code $S$ is
\[
  k \;:=\; \dim_{\mathbb{F}_2}(H_0^L),
\]
i.e., the $\mathbb{F}_2$-dimension (finrank) of the logical subspace $H_0^L$. In a subsystem code, only $H_0^L$ encodes logical information, while $H_0^G$ corresponds to gauge degrees of freedom.
\end{definition}

\begin{definition}[Z-Distance]
\label{def:SubsystemCSSCode.dZ}
\lean{SubsystemCSSCode.dZ}
\leanok
\uses{def:SubsystemCSSCode, def:SubsystemCSSCode.projL, def:CSSCode.homologyMkQ, def:hammingWeight}
The \emph{Z-distance} $d_Z$ of a subsystem CSS code $S$ is the minimum Hamming weight of a representative $z \in \ker(H_X)$ of a homology class $[z] \in H_0$ whose projection onto the logical subspace is nonzero:
\[
  d_Z \;:=\; \inf\bigl\{\, \operatorname{wt}(z) \;\bigm|\; z \in \ker(H_X),\; \pi_L([z]) \neq 0 \,\bigr\},
\]
where $\operatorname{wt}(z) = \operatorname{hammingWeight}(z)$ denotes the Hamming weight of $z \in \mathbb{F}_2^n$, $[z] = \pi_H(z)$ is the image under the homology quotient map, and $\pi_L : H_0 \to H_0^L$ is the logical projection. By convention, $d_Z = 0$ when no such class exists (i.e., $H_0^L = 0$).
\end{definition}

\begin{definition}[X-Distance]
\label{def:SubsystemCSSCode.dX}
\lean{SubsystemCSSCode.dX}
\leanok
\uses{def:SubsystemCSSCode, def:SubsystemCSSCode.coprojL, def:CSSCode.cohomologyMkQ, def:hammingWeight}
The \emph{X-distance} $d_X$ of a subsystem CSS code $S$ is the minimum Hamming weight (under the identification $\operatorname{dotProductEquiv} : (\mathbb{F}_2^n)^* \cong \mathbb{F}_2^n$) of a representative $\zeta \in \ker(H_Z)$ of a cohomology class $[\zeta] \in H^0$ whose projection onto $H^0_L$ is nonzero:
\[
  d_X \;:=\; \inf\bigl\{\, \operatorname{wt}(\phi^{-1}(\zeta)) \;\bigm|\; \zeta \in \ker(H_Z),\; \pi_L^*([\zeta]) \neq 0 \,\bigr\},
\]
where $\phi : \mathbb{F}_2^n \xrightarrow{\;\sim\;} (\mathbb{F}_2^n)^*$ is the dot-product equivalence, $[\zeta] = \pi_{H^0}(\zeta)$ is the image under the cohomology quotient map, and $\pi_L^* : H^0 \to H^0_L$ is the logical cohomology projection. By convention, $d_X = 0$ when $H^0_L = 0$.
\end{definition}

\begin{definition}[Distance]
\label{def:SubsystemCSSCode.distance}
\lean{SubsystemCSSCode.distance}
\leanok
\uses{def:SubsystemCSSCode.dX, def:SubsystemCSSCode.dZ}
The \emph{distance} of a subsystem CSS code $S$ is
\[
  d \;:=\; \min(d_X,\, d_Z).
\]
\end{definition}

\begin{definition}[{$[\![n,k,d]\!]$-Code Predicate}]
\label{def:SubsystemCSSCode.IsNKDCode}
\lean{SubsystemCSSCode.IsNKDCode}
\leanok
\uses{def:SubsystemCSSCode.logicalQubits, def:SubsystemCSSCode.distance}
A subsystem CSS code $S$ is an $[\![n,k,d]\!]$-code (written $\operatorname{IsNKDCode}(k, d)$) if
\[
  \dim_{\mathbb{F}_2}(H_0^L) = k \quad \text{and} \quad d_S = d,
\]
i.e., it encodes $k$ logical qubits with overall distance $d$.
\end{definition}

\begin{definition}[{$[\![n,k,d_X,d_Z]\!]$-Code Predicate}]
\label{def:SubsystemCSSCode.IsNKDXZCode}
\lean{SubsystemCSSCode.IsNKDXZCode}
\leanok
\uses{def:SubsystemCSSCode.logicalQubits, def:SubsystemCSSCode.dX, def:SubsystemCSSCode.dZ}
A subsystem CSS code $S$ is an $[\![n,k,d_X,d_Z]\!]$-code (written $\operatorname{IsNKDXZCode}(k, d_X, d_Z)$) if
\[
  \dim_{\mathbb{F}_2}(H_0^L) = k, \quad d_X^S = d_X, \quad \text{and} \quad d_Z^S = d_Z,
\]
i.e., it encodes $k$ logical qubits with X-distance $d_X$ and Z-distance $d_Z$ separately specified.
\end{definition}

%--- Def_7: CellComplex ---
\chapter{Def 7: Cell Complex}

\begin{definition}[Cell Complex]
\label{def:CellComplex}
\lean{CellComplex}
\leanok

A \textbf{regular cell complex} $X$ consists of:
\begin{itemize}
  \item A family of types $X_n$ (for each $n \in \mathbb{Z}$), whose elements are called $n$-cells, each equipped with a \texttt{Fintype} instance and a \texttt{DecidableEq} instance.
  \item A boundary map $\partial$: for each $n$-cell $\sigma \in X_n$, a finite set $\partial\sigma \subseteq X_{n-1}$ of $(n-1)$-cells in the boundary of $\sigma$.
  \item The \textbf{chain complex condition}: for every $n \in \mathbb{Z}$, every $n$-cell $\sigma \in X_n$, and every $(n-2)$-cell $\rho \in X_{n-2}$, the cardinality
  \[
    \left|\{\, \tau \in \partial\sigma \mid \rho \in \partial\tau \,\}\right|
  \]
  is even (i.e., $\equiv 0 \pmod{2}$).
\end{itemize}
\end{definition}

\begin{definition}[Differential Map]
\label{def:CellComplex.differentialMap}
\lean{CellComplex.differentialMap}
\leanok
\uses{def:CellComplex}
Let $X$ be a cell complex and let $C_i(X) = (X_i \to \mathbb{F}_2)$ denote the free $\mathbb{F}_2$-vector space on the $i$-cells (i.e., functions from $X_i$ to $\mathbb{F}_2$). The \textbf{differential}
\[
  \partial_n : C_n(X) \longrightarrow C_{n-1}(X)
\]
is the $\mathbb{F}_2$-linear map defined for $f \in C_n(X)$ and $\tau \in X_{n-1}$ by
\[
  \partial_n(f)(\tau) \;=\; \sum_{\substack{\sigma \in X_n \\ \tau \,\in\, \partial\sigma}} f(\sigma).
\]
Equivalently, on a basis element $e_\sigma$ (the indicator function of $\sigma$), we have $\partial_n(e_\sigma) = \sum_{\tau \in \partial\sigma} e_\tau$, extended linearly.
\end{definition}

\begin{theorem}[Chain Complex Condition: $\partial_{n-1} \circ \partial_n = 0$]
\label{thm:CellComplex.differentialMap_comp}
\lean{CellComplex.differentialMap_comp}
\leanok
\uses{def:CellComplex, def:CellComplex.differentialMap}
Let $X$ be a cell complex. For every $n \in \mathbb{Z}$,
\[
  \partial_{n-1} \circ \partial_n = 0 : C_n(X) \longrightarrow C_{n-2}(X).
\]
\end{theorem}

\begin{proof}
\leanok
\uses{def:CellComplex, def:CellComplex.differentialMap}
Let $f \in C_n(X)$ and $\rho \in X_{n-2}$ be arbitrary. By extensionality, it suffices to show
\[
  (\partial_{n-1} \circ \partial_n)(f)(\rho) = 0.
\]
Simplifying using the definition of $\partial_{n-1}$, $\partial_n$, and the zero map, the left-hand side becomes
\[
  \sum_{\tau \;\in\; X_{n-1},\; \rho \,\in\, \partial\tau} \left(\sum_{\sigma \;\in\; X_n,\; \tau \,\in\, \partial\sigma} f(\sigma)\right).
\]
We rewrite the inner filtered sum using the identity $\sum_{\sigma \in S,\, P(\sigma)} f(\sigma) = \sum_{\sigma \in S} \mathbf{1}_{P(\sigma)}\, f(\sigma)$, converting to a sum over all of $X_n$ with an indicator. Then we swap the order of summation (both index sets are finite), obtaining
\[
  \sum_{\sigma \in X_n} \left(\sum_{\tau \;\in\; X_{n-1},\; \rho \,\in\, \partial\tau} \mathbf{1}_{\tau \,\in\, \partial\sigma}\, f(\sigma)\right).
\]
It suffices to show each inner summand (over $\sigma$) is zero. For fixed $\sigma \in X_n$, we apply the identity for sums of if-then-else expressions: the sum over $\tau$ splits as
\[
  \sum_{\substack{\tau \in X_{n-1} \\ \rho \,\in\, \partial\tau}} \mathbf{1}_{\tau \,\in\, \partial\sigma}\, f(\sigma)
  \;=\; f(\sigma) \cdot \left|\bigl\{\, \tau \in X_{n-1} \;\big|\; \rho \in \partial\tau \;\text{ and }\; \tau \in \partial\sigma \,\bigr\}\right|_{\mathbb{F}_2}.
\]
Combining the two filter conditions with \texttt{filter\_filter}, we identify the counting set as
\[
  \bigl\{\, \tau \in \partial\sigma \;\big|\; \rho \in \partial\tau \,\bigr\},
\]
whose cardinality equals $\left|\{\tau \in \partial\sigma \mid \rho \in \partial\tau\}\right|$. By the chain complex condition (\texttt{bdry\_bdry}) applied to $\sigma$ and $\rho$, this cardinality is even. Since the cardinality of an even set, cast to $\mathbb{F}_2$, equals $0$ (by \texttt{Even.natCast\_zmod\_two}), we have $f(\sigma) \cdot 0 = 0$. Hence the entire sum is zero.
\end{proof}

\begin{definition}[Associated Chain Complex $C(X)$]
\label{def:CellComplex.chainComplex}
\lean{CellComplex.chainComplex}
\leanok
\uses{def:CellComplex, def:CellComplex.differentialMap, thm:CellComplex.differentialMap_comp, def:ChainComplexF2}
Let $X$ be a cell complex. The \textbf{associated chain complex} $C(X) = (C_\bullet(X), \partial)$ is the $\mathbb{F}_2$-chain complex defined as follows:
\begin{itemize}
  \item In degree $i \in \mathbb{Z}$, the chain module is $C_i(X) = (X_i \to \mathbb{F}_2)$, the free $\mathbb{F}_2$-vector space on the $i$-cells.
  \item The differential $d_{i,j} : C_i(X) \to C_j(X)$ is the linear map $\partial_{j+1}$ when $j + 1 = i$, and zero otherwise.
\end{itemize}
The chain complex condition $d_{j} \circ d_{j+1} = 0$ is verified by composing the differentials and applying \texttt{differentialMap\_comp}, after cancelling the coherence isomorphisms (\texttt{eqToHom}) that arise from the \texttt{ModuleCat} packaging.
\end{definition}

\begin{definition}[Cell Homology]
\label{def:CellComplex.cellHomology}
\lean{CellComplex.cellHomology}
\leanok
\uses{def:CellComplex, def:CellComplex.chainComplex, def:ChainComplexF2.homologyPrime}
Let $X$ be a cell complex. The \textbf{homology} of $X$ in degree $i \in \mathbb{Z}$ is
\[
  H_i(X) \;:=\; H_i(C(X)),
\]
where $H_i(C(X))$ is the $i$-th homology of the associated chain complex $C(X)$ (as defined in the chain complex homology construction of Definition~1).
\end{definition}

%--- Def_8: CycleGraph ---
\chapter{Def 8: Cycle Graph}

\begin{definition}[Cycle Graph Differential]
\label{def:CycleGraph.differential}
\lean{CycleGraph.differential}
\leanok
\uses{def:F2}
Let $\ell \geq 1$ be a natural number. The \emph{differential} $\partial_1 : \mathbb{F}_2^\ell \to \mathbb{F}_2^\ell$ of the cycle graph is the linear map defined by
\[
(\partial_1 f)(i) = f(i) + f\!\left(\lfloor(i + \ell - 1) \bmod \ell\rfloor\right)
\]
for $f : \operatorname{Fin}(\ell) \to \mathbb{F}_2$ and $i \in \operatorname{Fin}(\ell)$. On basis vectors $e_j$, one has $\partial_1(e_j) = e_j + e_{(j+1) \bmod \ell}$, so column $j$ of the corresponding matrix has $1$-entries in rows $j$ and $(j+1) \bmod \ell$, encoding the boundary relation $\partial\sigma_j = \{\tau_j, \tau_{(j+1) \bmod \ell}\}$.
\end{definition}

\begin{definition}[All-Ones Vector]
\label{def:CycleGraph.onesVec}
\lean{CycleGraph.onesVec}
\leanok
\uses{def:F2}
The \emph{all-ones vector} $\mathbf{1} \in \mathbb{F}_2^\ell$ is defined by $\mathbf{1}(i) = 1$ for all $i \in \operatorname{Fin}(\ell)$.
\end{definition}

\begin{lemma}[Differential Annihilates All-Ones Vector]
\label{lem:CycleGraph.differential_onesVec}
\lean{CycleGraph.differential_onesVec}
\leanok
\uses{def:CycleGraph.differential, def:CycleGraph.onesVec, thm:F2.add_self_eq_zero}
$\partial_1(\mathbf{1}) = 0$ in $\mathbb{F}_2^\ell$.
\end{lemma}

\begin{proof}
\leanok
\uses{def:CycleGraph.differential, def:CycleGraph.onesVec, thm:F2.add_self_eq_zero}
By extensionality, it suffices to show that $(\partial_1 \mathbf{1})(i) = 0$ for each $i \in \operatorname{Fin}(\ell)$. Unfolding the definitions, $(\partial_1 \mathbf{1})(i) = \mathbf{1}(i) + \mathbf{1}((i+\ell-1)\bmod\ell) = 1 + 1 = 0$ in $\mathbb{F}_2$, where the last equality holds by $x + x = 0$ in $\mathbb{F}_2$.
\end{proof}

\begin{lemma}[Kernel Condition Implies Adjacent Equality]
\label{lem:CycleGraph.ker_differential_succ}
\lean{CycleGraph.ker_differential_succ}
\leanok
\uses{def:CycleGraph.differential, thm:F2.sub_eq_add}
Let $f : \operatorname{Fin}(\ell) \to \mathbb{F}_2$ satisfy $\partial_1 f = 0$. For any $i \in \mathbb{N}$ with $i + 1 < \ell$, we have $f(i+1) = f(i)$.
\end{lemma}

\begin{proof}
\leanok
\uses{def:CycleGraph.differential, thm:F2.sub_eq_add}
Let $h$ denote the hypothesis $\partial_1 f = 0$. Evaluating at position $i+1$, we obtain $h_{i+1} : (\partial_1 f)(i+1) = 0$, which by the definition of the differential gives $f(i+1) + f\!\left((i+1+\ell-1)\bmod\ell\right) = 0$. We compute $(i+1+\ell-1)\bmod\ell = (i+\ell)\bmod\ell = i\bmod\ell = i$ (since $i < \ell$), so after rewriting we get $f(i+1) + f(i) = 0$. It follows that $f(i+1) - f(i) = 0$, and since subtraction equals addition in $\mathbb{F}_2$ (using $\mathbb{F}_2.\mathtt{sub\_eq\_add}$), this gives $f(i+1) = f(i)$.
\end{proof}

\begin{lemma}[Kernel Elements Are Constant]
\label{lem:CycleGraph.ker_differential_const}
\lean{CycleGraph.ker_differential_const}
\leanok
\uses{def:CycleGraph.differential, lem:CycleGraph.ker_differential_succ}
If $\partial_1 f = 0$, then $f(i) = f(0)$ for all $i \in \operatorname{Fin}(\ell)$.
\end{lemma}

\begin{proof}
\leanok
\uses{lem:CycleGraph.ker_differential_succ}
Write $i = \langle n, h_n \rangle$. We proceed by induction on $n$. In the base case $n = 0$, the equality $f(0) = f(0)$ holds by reflexivity. In the inductive step, assuming $f(k) = f(0)$ for $k$ with $k < \ell$, we apply the lemma $\mathtt{ker\_differential\_succ}$ to obtain $f(k+1) = f(k)$, and the result follows by the inductive hypothesis.
\end{proof}

\begin{theorem}[Kernel Equals Span of All-Ones Vector]
\label{thm:CycleGraph.ker_differential_eq_span}
\lean{CycleGraph.ker_differential_eq_span}
\leanok
\uses{def:CycleGraph.differential, def:CycleGraph.onesVec, lem:CycleGraph.ker_differential_const, thm:F2.add_self_eq_zero}
$\ker(\partial_1) = \operatorname{span}_{\mathbb{F}_2}\{\mathbf{1}\}$ as submodules of $\mathbb{F}_2^\ell$.
\end{theorem}

\begin{proof}
\leanok
\uses{def:CycleGraph.differential, def:CycleGraph.onesVec, lem:CycleGraph.ker_differential_const, thm:F2.add_self_eq_zero}
By extensionality, it suffices to show for all $f$ that $f \in \ker(\partial_1)$ if and only if $f \in \operatorname{span}\{\mathbf{1}\}$, i.e., $\partial_1 f = 0 \Leftrightarrow \exists\, c \in \mathbb{F}_2,\ f = c \cdot \mathbf{1}$.

$(\Rightarrow)$: Suppose $\partial_1 f = 0$. We claim $f = f(0) \cdot \mathbf{1}$. By extensionality at $i$, we need $f(i) = f(0) \cdot \mathbf{1}(i) = f(0) \cdot 1 = f(0)$. This follows from $\mathtt{ker\_differential\_const}$, applying the symmetry of the equality.

$(\Leftarrow)$: Suppose $f = c \cdot \mathbf{1}$ for some $c \in \mathbb{F}_2$. Then for each $i$, $(\partial_1 f)(i) = (\partial_1(c \cdot \mathbf{1}))(i) = c \cdot (\partial_1 \mathbf{1})(i) = c \cdot (c + c) = c \cdot 0$. More precisely, unfolding the differential, $(\partial_1(c\cdot\mathbf{1}))(i) = c\cdot\mathbf{1}(i) + c\cdot\mathbf{1}((i+\ell-1)\bmod\ell) = c + c = 0$ by $x + x = 0$ in $\mathbb{F}_2$.
\end{proof}

\begin{lemma}[All-Ones Vector Is Nonzero]
\label{lem:CycleGraph.onesVec_ne_zero}
\lean{CycleGraph.onesVec_ne_zero}
\leanok
\uses{def:CycleGraph.onesVec}
If $\ell \geq 1$, then $\mathbf{1} \neq 0$ in $\mathbb{F}_2^\ell$.
\end{lemma}

\begin{proof}
\leanok
\uses{def:CycleGraph.onesVec}
Suppose for contradiction that $\mathbf{1} = 0$. Evaluating both sides at $\langle 0, h_\ell \rangle \in \operatorname{Fin}(\ell)$ (which exists since $\ell \geq 1$), we obtain $\mathbf{1}(0) = 1 = 0$, which is a contradiction since $1 \neq 0$ in $\mathbb{F}_2$.
\end{proof}

\begin{theorem}[Kernel Dimension Is One]
\label{thm:CycleGraph.finrank_ker_differential}
\lean{CycleGraph.finrank_ker_differential}
\leanok
\uses{thm:CycleGraph.ker_differential_eq_span, lem:CycleGraph.onesVec_ne_zero}
If $\ell \geq 1$, then $\dim_{\mathbb{F}_2} \ker(\partial_1) = 1$.
\end{theorem}

\begin{proof}
\leanok
\uses{thm:CycleGraph.ker_differential_eq_span, lem:CycleGraph.onesVec_ne_zero}
We rewrite using $\mathtt{ker\_differential\_eq\_span}$, which gives $\ker(\partial_1) = \operatorname{span}\{\mathbf{1}\}$. Since $\mathbf{1} \neq 0$ by $\mathtt{onesVec\_ne\_zero}$ (using $\ell \geq 1$), the span of a single nonzero vector in a vector space over a field has dimension $1$ (by $\mathtt{finrank\_span\_singleton}$).
\end{proof}

\begin{definition}[Cell Type of the Cycle Graph]
\label{def:CycleGraph.cellType}
\lean{CycleGraph.cellType}
\leanok

The \emph{cell type} function for the cycle graph assigns to each integer dimension $n$ the type of $n$-cells:
\[
\operatorname{cellType}(\ell, n) = \begin{cases} \operatorname{Fin}(\ell) & \text{if } n = 0,\\ \operatorname{Fin}(\ell) & \text{if } n = 1,\\ \varnothing & \text{otherwise.}\end{cases}
\]
Thus the cycle graph has $\ell$ vertices (0-cells) and $\ell$ edges (1-cells), and no cells in any other dimension.
\end{definition}

\begin{definition}[Boundary Map of the Cycle Graph]
\label{def:CycleGraph.cellBdry}
\lean{CycleGraph.cellBdry}
\leanok
\uses{def:CycleGraph.cellType}
The \emph{boundary map} for the cycle graph $C_\ell$ assigns to each cell its boundary:
\begin{itemize}
\item For a 1-cell $\sigma_i$ (with $i \in \operatorname{Fin}(\ell)$): $\partial\sigma_i = \{\tau_i,\, \tau_{(i+1)\bmod\ell}\} \subseteq \operatorname{Fin}(\ell)$, encoding the incidence relation of each edge with its two endpoint vertices.
\item For a 0-cell: $\partial\tau_i = \varnothing$ (the boundary of a vertex is empty, valued in $\operatorname{cellType}(\ell,-1) = \varnothing$).
\item For cells in all other dimensions: the boundary is vacuously empty (the cell type is $\varnothing$).
\end{itemize}
\end{definition}

\begin{theorem}[Boundary-of-Boundary Vanishes]
\label{thm:CycleGraph.cellBdry_bdry}
\lean{CycleGraph.cellBdry_bdry}
\leanok
\uses{def:CycleGraph.cellType, def:CycleGraph.cellBdry}
For all $n \in \mathbb{Z}$, all $n$-cells $\sigma$, and all $(n-2)$-cells $\rho$, the number of $(n-1)$-cells in $\partial\sigma$ whose boundary contains $\rho$ is even.
\end{theorem}

\begin{proof}
\leanok
\uses{def:CycleGraph.cellType, def:CycleGraph.cellBdry}
We verify by cases on $n$. For $n = 1$: the target type $\operatorname{cellType}(\ell, -1) = \varnothing$, so $\rho : \varnothing$ is absurd; the statement holds vacuously. For $n = 0$: similarly, $\operatorname{cellType}(\ell, -2) = \varnothing$ gives an absurd $\rho$. For $n = k+2$ with $k \geq 0$: $\operatorname{cellType}(\ell, k+2) = \varnothing$, so $\sigma : \varnothing$ is absurd. For $n < 0$: similarly $\operatorname{cellType}(\ell, n) = \varnothing$ gives an absurd $\sigma$. In all cases the condition is vacuously satisfied.
\end{proof}

\begin{definition}[Cycle Graph as a Cell Complex]
\label{def:CycleGraph.cycleGraph}
\lean{CycleGraph.cycleGraph}
\leanok
\uses{def:CellComplex, def:CycleGraph.cellType, def:CycleGraph.cellBdry, thm:CycleGraph.cellBdry_bdry}
The \emph{cycle graph} $C_\ell$ is the cell complex (Definition~\ref{def:CellComplex}) with:
\begin{itemize}
\item cells given by $\operatorname{cellType}(\ell)$: $\ell$ vertices (0-cells) and $\ell$ edges (1-cells), all indexed by $\operatorname{Fin}(\ell)$;
\item boundary map given by $\operatorname{cellBdry}(\ell)$: $\partial\sigma_i = \{\tau_i, \tau_{(i+1)\bmod\ell}\}$;
\item the chain complex condition $\partial\circ\partial = 0$ verified by $\mathtt{cellBdry\_bdry}$.
\end{itemize}
\end{definition}

\begin{theorem}[1-Cells of Cycle Graph]
\label{thm:CycleGraph.cycleGraph_cells_one}
\lean{CycleGraph.cycleGraph_cells_one}
\leanok
\uses{def:CycleGraph.cycleGraph, def:CycleGraph.cellType}
The type of 1-cells of $C_\ell$ equals $\operatorname{Fin}(\ell)$, i.e., $(C_\ell).\mathtt{cells}(1) = \operatorname{Fin}(\ell)$.
\end{theorem}

\begin{proof}
\leanok
\uses{def:CycleGraph.cycleGraph, def:CycleGraph.cellType}
This holds by reflexivity, since $\operatorname{cellType}(\ell, 1) = \operatorname{Fin}(\ell)$ by definition.
\end{proof}

\begin{theorem}[0-Cells of Cycle Graph]
\label{thm:CycleGraph.cycleGraph_cells_zero}
\lean{CycleGraph.cycleGraph_cells_zero}
\leanok
\uses{def:CycleGraph.cycleGraph, def:CycleGraph.cellType}
The type of 0-cells of $C_\ell$ equals $\operatorname{Fin}(\ell)$, i.e., $(C_\ell).\mathtt{cells}(0) = \operatorname{Fin}(\ell)$.
\end{theorem}

\begin{proof}
\leanok
\uses{def:CycleGraph.cycleGraph, def:CycleGraph.cellType}
This holds by reflexivity, since $\operatorname{cellType}(\ell, 0) = \operatorname{Fin}(\ell)$ by definition.
\end{proof}

\begin{lemma}[Predecessor-Successor Modular Identity]
\label{lem:CycleGraph.pred_mod_succ}
\lean{CycleGraph.pred_mod_succ}
\leanok

For $\ell \geq 1$ and $i \in \operatorname{Fin}(\ell)$, we have $\left((i + \ell - 1)\bmod\ell + 1\right)\bmod\ell = i$.
\end{lemma}

\begin{proof}
\leanok

We rewrite using $\mathtt{Nat.add\_mod}$ and $\mathtt{Nat.mod\_mod}$. Observe that $(i + \ell - 1) + 1 = i + 1 \cdot \ell$, so $(i + \ell - 1 + 1) \bmod \ell = (i + \ell)\bmod\ell = i\bmod\ell = i$ (since $i < \ell$). By $\mathtt{Nat.add\_mul\_mod\_self\_right}$ and $\mathtt{Nat.mod\_eq\_of\_lt}$, the claim follows.
\end{proof}

\begin{lemma}[No Element Equals Its Successor Modulo $\ell$]
\label{lem:CycleGraph.fin_ne_succ_mod}
\lean{CycleGraph.fin_ne_succ_mod}
\leanok

For $\ell \geq 2$ and $x \in \operatorname{Fin}(\ell)$, we have $x \neq \langle (x + 1)\bmod\ell,\, \cdot\rangle$ as elements of $\operatorname{Fin}(\ell)$.
\end{lemma}

\begin{proof}
\leanok

Suppose for contradiction that $x = \langle (x+1)\bmod\ell, \cdot\rangle$, i.e., $x = (x+1)\bmod\ell$. We split on whether $x + 1 < \ell$ or $x + 1 \geq \ell$. If $x + 1 < \ell$, then $(x+1)\bmod\ell = x+1$, giving $x = x+1$, a contradiction by integer arithmetic. If $x + 1 \geq \ell$, then since $x < \ell$ we have $x + 1 = \ell$, so $(x+1)\bmod\ell = 0$, giving $x = 0$ and thus $\ell = 1$, contradicting $\ell \geq 2$.
\end{proof}

\begin{lemma}[No Element Equals Its Predecessor Modulo $\ell$]
\label{lem:CycleGraph.fin_ne_pred_mod}
\lean{CycleGraph.fin_ne_pred_mod}
\leanok

For $\ell \geq 2$ and $i \in \operatorname{Fin}(\ell)$, we have $i \neq \langle (i + \ell - 1)\bmod\ell,\, \cdot\rangle$ as elements of $\operatorname{Fin}(\ell)$.
\end{lemma}

\begin{proof}
\leanok

Suppose for contradiction that $i = (i + \ell - 1)\bmod\ell$. We split on $i = 0$ and $i \geq 1$. If $i = 0$: $(0 + \ell - 1)\bmod\ell = (\ell - 1)\bmod\ell = \ell - 1$ (since $\ell - 1 < \ell$), giving $0 = \ell - 1$, i.e., $\ell = 1$, contradicting $\ell \geq 2$. If $i \geq 1$: $i + \ell - 1 = (i-1) + 1\cdot\ell$, so by $\mathtt{Nat.add\_mul\_mod\_self\_right}$ we get $(i+\ell-1)\bmod\ell = (i-1)\bmod\ell = i-1$ (since $i-1 < \ell$), giving $i = i-1$, a contradiction.
\end{proof}

\begin{lemma}[Filter Set for Differential Equals Two-Element Set]
\label{lem:CycleGraph.differentialMap_filter_eq}
\lean{CycleGraph.differentialMap_filter_eq}
\leanok
\uses{def:CycleGraph.cycleGraph, lem:CycleGraph.pred_mod_succ, lem:CycleGraph.fin_ne_succ_mod}
For $\ell \geq 2$ and $i \in \operatorname{Fin}(\ell)$, letting $j = \langle (i + \ell - 1)\bmod\ell,\, \cdot\rangle$, the set of 1-cells $\sigma$ of $C_\ell$ whose boundary contains $i$ equals $\{i, j\}$:
\[
\Bigl\{\sigma \in \operatorname{Fin}(\ell) \;\Big|\; i \in \partial\sigma\Bigr\} = \{i,\, j\}.
\]
\end{lemma}

\begin{proof}
\leanok
\uses{def:CycleGraph.cycleGraph, lem:CycleGraph.pred_mod_succ, lem:CycleGraph.fin_ne_succ_mod}
By extensionality, for any $\sigma \in \operatorname{Fin}(\ell)$ we show $i \in \partial\sigma \Leftrightarrow \sigma \in \{i, j\}$.

Recall $\partial\sigma = \{\sigma, \langle (\sigma+1)\bmod\ell,\cdot\rangle\}$, so $i \in \partial\sigma$ iff $i = \sigma$ or $i = \langle (\sigma+1)\bmod\ell,\cdot\rangle$.

$(\Rightarrow)$: If $i = \sigma$, then $\sigma \in \{i,j\}$ directly (as the left member). If $i = \langle(\sigma+1)\bmod\ell, \cdot\rangle$, then $i = (\sigma+1)\bmod\ell$, i.e., $\sigma = (i+\ell-1)\bmod\ell = j$ by the modular identity $\mathtt{pred\_mod\_succ}$, so $\sigma \in \{i,j\}$ as the right member.

$(\Leftarrow)$: If $\sigma = i$, then $i = \sigma \in \partial\sigma$ directly. If $\sigma = j = \langle(i+\ell-1)\bmod\ell,\cdot\rangle$, then $\langle(\sigma+1)\bmod\ell,\cdot\rangle = \langle(j+1)\bmod\ell,\cdot\rangle = i$ by $\mathtt{pred\_mod\_succ}$, so $i \in \partial j$.
\end{proof}

\begin{theorem}[Cell Complex Differential Agrees with Algebraic Differential]
\label{thm:CycleGraph.differentialMap_eq_differential}
\lean{CycleGraph.differentialMap_eq_differential}
\leanok
\uses{def:CycleGraph.cycleGraph, def:CycleGraph.differential, def:CellComplex, lem:CycleGraph.differentialMap_filter_eq, lem:CycleGraph.fin_ne_pred_mod}
For $\ell \geq 2$, the cell complex differential $\partial_1^{\mathrm{cell}} : \mathbb{F}_2^{C_1} \to \mathbb{F}_2^{C_0}$ of $C_\ell$ (as defined in Definition~\ref{def:CellComplex}) equals the algebraic map $\partial_1$ of Definition~\ref{def:CycleGraph.differential}:
\[
(C_\ell).\mathtt{differentialMap}(1) = \partial_1.
\]
\end{theorem}

\begin{proof}
\leanok
\uses{def:CycleGraph.cycleGraph, def:CycleGraph.differential, def:CellComplex, lem:CycleGraph.differentialMap_filter_eq, lem:CycleGraph.fin_ne_pred_mod}
By extensionality of linear maps and of functions, it suffices to show that for every $f \in \mathbb{F}_2^\ell$ and every $i \in \operatorname{Fin}(\ell)$, the two maps agree at $(f, i)$. We apply $\mathtt{CellComplex.differentialMap\_apply}$ to rewrite the left side as a sum:
\[
\bigl((C_\ell).\mathtt{differentialMap}(1)\, f\bigr)(i) = \sum_{\sigma \in \operatorname{univ},\, i \in \partial\sigma} f(\sigma).
\]
Setting $j = \langle(i+\ell-1)\bmod\ell,\cdot\rangle$ and applying $\mathtt{differentialMap\_filter\_eq}$, the filter set equals $\{i, j\}$. Since $i \neq j$ by $\mathtt{fin\_ne\_pred\_mod}$ (using $\ell \geq 2$), the sum over the pair $\{i, j\}$ equals $f(i) + f(j)$ by $\mathtt{Finset.sum\_pair}$. This is precisely $(\partial_1 f)(i) = f(i) + f\!\left((i+\ell-1)\bmod\ell\right)$, completing the proof.
\end{proof}

\begin{lemma}[Dimension of $\mathbb{F}_2^\ell$]
\label{lem:CycleGraph.finrank_F2_fin}
\lean{CycleGraph.finrank_F2_fin}
\leanok
\uses{def:F2}
$\dim_{\mathbb{F}_2}(\operatorname{Fin}(\ell) \to \mathbb{F}_2) = \ell$.
\end{lemma}

\begin{proof}
\leanok
\uses{def:F2}
This follows directly from $\mathtt{Module.finrank\_fin\_fun}$, which gives $\dim_F(F^n) = n$ for any field $F$.
\end{proof}

\begin{theorem}[Image Dimension Is $\ell - 1$]
\label{thm:CycleGraph.finrank_range_differential}
\lean{CycleGraph.finrank_range_differential}
\leanok
\uses{def:CycleGraph.differential, thm:CycleGraph.finrank_ker_differential, lem:CycleGraph.finrank_F2_fin}
If $\ell \geq 1$, then $\dim_{\mathbb{F}_2}\operatorname{im}(\partial_1) = \ell - 1$.
\end{theorem}

\begin{proof}
\leanok
\uses{def:CycleGraph.differential, thm:CycleGraph.finrank_ker_differential, lem:CycleGraph.finrank_F2_fin}
By the rank--nullity theorem ($\mathtt{LinearMap.finrank\_range\_add\_finrank\_ker}$):
\[
\dim\operatorname{im}(\partial_1) + \dim\ker(\partial_1) = \dim\mathbb{F}_2^\ell.
\]
Substituting $\dim\mathbb{F}_2^\ell = \ell$ (by $\mathtt{finrank\_F2\_fin}$) and $\dim\ker(\partial_1) = 1$ (by $\mathtt{finrank\_ker\_differential}$, using $\ell \geq 1$), we obtain $\dim\operatorname{im}(\partial_1) + 1 = \ell$, hence $\dim\operatorname{im}(\partial_1) = \ell - 1$ by integer arithmetic.
\end{proof}

\begin{theorem}[First Homology Dimension of the Cycle Graph]
\label{thm:CycleGraph.dim_H1_cycleGraph}
\lean{CycleGraph.dim_H1_cycleGraph}
\leanok
\uses{thm:CycleGraph.finrank_ker_differential}
If $\ell \geq 1$, then $\dim_{\mathbb{F}_2} H_1(C_\ell) = 1$.
\end{theorem}

\begin{proof}
\leanok
\uses{thm:CycleGraph.finrank_ker_differential}
Since $C_\ell$ has no 2-cells, the image of $\partial_2$ is zero, so $H_1(C_\ell) = \ker(\partial_1) / \operatorname{im}(\partial_2) \cong \ker(\partial_1)$. Thus $\dim H_1(C_\ell) = \dim\ker(\partial_1) = 1$ by $\mathtt{finrank\_ker\_differential}$ (using $\ell \geq 1$).
\end{proof}

\begin{theorem}[Zeroth Homology Dimension of the Cycle Graph]
\label{thm:CycleGraph.dim_H0_cycleGraph}
\lean{CycleGraph.dim_H0_cycleGraph}
\leanok
\uses{thm:CycleGraph.finrank_range_differential, lem:CycleGraph.finrank_F2_fin}
If $\ell \geq 1$, then $\dim_{\mathbb{F}_2}\left(\mathbb{F}_2^\ell \big/ \operatorname{im}(\partial_1)\right) = 1$.
\end{theorem}

\begin{proof}
\leanok
\uses{thm:CycleGraph.finrank_range_differential, lem:CycleGraph.finrank_F2_fin}
By $\mathtt{Submodule.finrank\_quotient\_add\_finrank}$:
\[
\dim\left(\mathbb{F}_2^\ell / \operatorname{im}(\partial_1)\right) + \dim\operatorname{im}(\partial_1) = \dim\mathbb{F}_2^\ell.
\]
Substituting $\dim\mathbb{F}_2^\ell = \ell$ (by $\mathtt{finrank\_F2\_fin}$) and $\dim\operatorname{im}(\partial_1) = \ell - 1$ (by $\mathtt{finrank\_range\_differential}$, using $\ell \geq 1$), we get $\dim H_0 + (\ell - 1) = \ell$, so $\dim H_0 = 1$ by integer arithmetic.
\end{proof}

\begin{definition}[Repetition Code]
\label{def:CycleGraph.repetitionCode}
\lean{CycleGraph.repetitionCode}
\leanok
\uses{def:ClassicalCode, def:ClassicalCode.ofParityCheck, def:CycleGraph.differential}
The \emph{repetition code} of length $\ell$ is the classical code (Definition~\ref{def:ClassicalCode}) defined by
\[
\operatorname{repetitionCode}(\ell) = \mathtt{ClassicalCode.ofParityCheck}(\partial_1),
\]
i.e., it is the classical code whose codeword space is $\ker(\partial_1) = \operatorname{span}\{\mathbf{1}\} \subseteq \mathbb{F}_2^\ell$.
\end{definition}

\begin{theorem}[Repetition Code Equals Kernel]
\label{thm:CycleGraph.repetitionCode_eq_ker}
\lean{CycleGraph.repetitionCode_eq_ker}
\leanok
\uses{def:CycleGraph.repetitionCode, def:CycleGraph.differential}
The codeword space of the repetition code equals the kernel of $\partial_1$:
\[
\operatorname{repetitionCode}(\ell).\mathtt{code} = \ker(\partial_1).
\]
\end{theorem}

\begin{proof}
\leanok
\uses{def:CycleGraph.repetitionCode, def:CycleGraph.differential}
This holds by reflexivity, as $\mathtt{ClassicalCode.ofParityCheck}(\partial_1).\mathtt{code}$ is defined to be $\ker(\partial_1)$.
\end{proof}

\begin{theorem}[Repetition Code Has Dimension One]
\label{thm:CycleGraph.repetitionCode_dimension}
\lean{CycleGraph.repetitionCode_dimension}
\leanok
\uses{def:CycleGraph.repetitionCode, def:ClassicalCode.dimension, thm:CycleGraph.finrank_ker_differential, thm:CycleGraph.repetitionCode_eq_ker}
If $\ell \geq 1$, then $\operatorname{repetitionCode}(\ell).\mathtt{dimension} = 1$.
\end{theorem}

\begin{proof}
\leanok
\uses{def:CycleGraph.repetitionCode, def:ClassicalCode.dimension, thm:CycleGraph.finrank_ker_differential, thm:CycleGraph.repetitionCode_eq_ker}
Unfolding $\mathtt{ClassicalCode.dimension}$, the dimension is defined as $\dim_{\mathbb{F}_2}(\mathtt{code})$. Rewriting by $\mathtt{repetitionCode\_eq\_ker}$, this becomes $\dim_{\mathbb{F}_2}\ker(\partial_1)$, which equals $1$ by $\mathtt{finrank\_ker\_differential}$ using $\ell \geq 1$.
\end{proof}

%--- Def_9: DoubleComplex ---
\chapter{Def 9: Double Complex}

\begin{definition}[Double Complex over $\mathbb{F}_2$]
\label{def:DoubleComplexF2}
\lean{DoubleComplex𝔽₂}
\leanok

A \emph{double complex} of $\mathbb{F}_2$-vector spaces is an array of $\mathbb{F}_2$-vector spaces $E_{p,q}$ indexed by $(p,q) \in \mathbb{Z} \times \mathbb{Z}$, equipped with:
\begin{itemize}
    \item a \emph{vertical differential} $\partial^v_{p,q} : E_{p,q} \to E_{p,q-1}$ (decreasing $q$),
    \item a \emph{horizontal differential} $\partial^h_{p,q} : E_{p,q} \to E_{p-1,q}$ (decreasing $p$),
\end{itemize}
satisfying:
\begin{enumerate}
    \item $\partial^v_{p,q-1} \circ \partial^v_{p,q} = 0$ for all $p, q$,
    \item $\partial^h_{p-1,q} \circ \partial^h_{p,q} = 0$ for all $p, q$,
    \item $\partial^v_{p-1,q} \circ \partial^h_{p,q} = \partial^h_{p,q-1} \circ \partial^v_{p,q}$ for all $p, q$ (commutativity).
\end{enumerate}
Over $\mathbb{F}_2$, commutativity and anticommutativity coincide since $-1 = 1$ in characteristic $2$.

Formally, $\operatorname{DoubleComplex}_{\mathbb{F}_2}$ is defined as $\operatorname{HomologicalComplex}_2(\operatorname{ModuleCat}\ \mathbb{F}_2,\ \operatorname{ComplexShape.down}\ \mathbb{Z},\ \operatorname{ComplexShape.down}\ \mathbb{Z})$, where the first $\operatorname{ComplexShape.down}\ \mathbb{Z}$ is the horizontal direction and the second is the vertical direction.
\end{definition}

\begin{definition}[Component Vector Space]
\label{def:DoubleComplexF2.obj}
\lean{DoubleComplex𝔽₂.obj}
\leanok
\uses{def:DoubleComplexF2}
Let $E : \operatorname{DoubleComplex}_{\mathbb{F}_2}$. The $\mathbb{F}_2$-vector space at position $(p, q) \in \mathbb{Z} \times \mathbb{Z}$ is defined by
\[
  E_{p,q} := (E_p)_q,
\]
where $E_p$ denotes the $p$-th column complex and $(E_p)_q$ its $q$-th component.
\end{definition}

\begin{definition}[Vertical Differential]
\label{def:DoubleComplexF2.dv}
\lean{DoubleComplex𝔽₂.dv}
\leanok
\uses{def:DoubleComplexF2.obj}
Let $E : \operatorname{DoubleComplex}_{\mathbb{F}_2}$. The \emph{vertical differential} at position $(p,q)$ is the $\mathbb{F}_2$-linear map
\[
  \partial^v_{p,q} : E_{p,q} \to E_{p,q-1}
\]
defined as the underlying linear map of the chain complex differential $(E_p).d_{q,q-1}$.
\end{definition}

\begin{definition}[Horizontal Differential]
\label{def:DoubleComplexF2.dh}
\lean{DoubleComplex𝔽₂.dh}
\leanok
\uses{def:DoubleComplexF2.obj}
Let $E : \operatorname{DoubleComplex}_{\mathbb{F}_2}$. The \emph{horizontal differential} at position $(p,q)$ is the $\mathbb{F}_2$-linear map
\[
  \partial^h_{p,q} : E_{p,q} \to E_{p-1,q}
\]
defined as the underlying linear map of the $q$-th component of the horizontal chain map $E.d_{p,p-1}$.
\end{definition}

\begin{theorem}[Vertical Differential Squares to Zero]
\label{thm:DoubleComplexF2.dv_comp_dv}
\lean{DoubleComplex𝔽₂.dv_comp_dv}
\leanok
\uses{def:DoubleComplexF2.dv}
Let $E : \operatorname{DoubleComplex}_{\mathbb{F}_2}$. For all $p, q \in \mathbb{Z}$,
\[
  \partial^v_{p,q-1} \circ \partial^v_{p,q} = 0.
\]
\end{theorem}

\begin{proof}
\leanok
\uses{def:DoubleComplexF2.dv}
Let $p, q \in \mathbb{Z}$ be arbitrary. By unfolding the definition of $\partial^v$, it suffices to show that for all $x \in E_{p,q}$,
\[
  ((E_p).d_{q-1,q-2} \circ (E_p).d_{q,q-1})(x) = 0.
\]
By the chain complex axiom $d \circ d = 0$ applied to $E_p$ (i.e., $(E_p).d_{q,q-1} \gg (E_p).d_{q-1,q-2} = 0$), and simplifying, the composition is the zero map. Rewriting the categorical composition as a composition of module maps via $\operatorname{ModuleCat.comp_apply}$, we conclude that the value at $x$ is $0$.
\end{proof}

\begin{theorem}[Horizontal Differential Squares to Zero]
\label{thm:DoubleComplexF2.dh_comp_dh}
\lean{DoubleComplex𝔽₂.dh_comp_dh}
\leanok
\uses{def:DoubleComplexF2.dh}
Let $E : \operatorname{DoubleComplex}_{\mathbb{F}_2}$. For all $p, q \in \mathbb{Z}$,
\[
  \partial^h_{p-1,q} \circ \partial^h_{p,q} = 0.
\]
\end{theorem}

\begin{proof}
\leanok
\uses{def:DoubleComplexF2.dh}
Let $p, q \in \mathbb{Z}$ be arbitrary. By unfolding the definition of $\partial^h$, it suffices to show that for all $x \in E_{p,q}$,
\[
  (((E.d_{p-1,p-2}).f_q) \circ ((E.d_{p,p-1}).f_q))(x) = 0.
\]
From the chain complex axiom $d \circ d = 0$ applied to $E$ (i.e., $E.d_{p,p-1} \gg E.d_{p-1,p-2} = 0$), we obtain that $((E.d_{p,p-1} \gg E.d_{p-1,p-2}).f_q)(x) = 0$ after simplification. Rewriting using $\operatorname{HomologicalComplex.comp_f}$ and $\operatorname{ModuleCat.comp_apply}$, the result follows.
\end{proof}

\begin{theorem}[Commutativity of Differentials]
\label{thm:DoubleComplexF2.comm}
\lean{DoubleComplex𝔽₂.comm}
\leanok
\uses{def:DoubleComplexF2.dv, def:DoubleComplexF2.dh}
Let $E : \operatorname{DoubleComplex}_{\mathbb{F}_2}$. For all $p, q \in \mathbb{Z}$, the vertical and horizontal differentials commute:
\[
  \partial^v_{p-1,q} \circ \partial^h_{p,q} = \partial^h_{p,q-1} \circ \partial^v_{p,q}.
\]
\end{theorem}

\begin{proof}
\leanok
\uses{def:DoubleComplexF2.dv, def:DoubleComplexF2.dh}
Let $p, q \in \mathbb{Z}$ be arbitrary. By unfolding the definitions of $\partial^v$ and $\partial^h$, it suffices to show that for all $x \in E_{p,q}$,
\[
  ((E.d_{p,p-1}).f_q \text{ composed vertically then horizontally})(x) = ((E.d_{p,p-1}).f_q \text{ composed horizontally then vertically})(x).
\]
We apply the Mathlib lemma $\operatorname{HomologicalComplex}_2\texttt{.d\_comm}$ to $E$ at indices $p, p-1, q, q-1$, which gives the required commutativity identity as a morphism in $\operatorname{ModuleCat}\ \mathbb{F}_2$. Taking the underlying linear map via $\operatorname{ModuleCat.Hom.hom}$ and evaluating at $x$ (using $\operatorname{congr_arg}$), the result follows by naturality of the bicomplex structure.
\end{proof}

\begin{theorem}[Anticommutativity Equals Commutativity over $\mathbb{F}_2$]
\label{thm:DoubleComplexF2.anticomm_eq_comm}
\lean{DoubleComplex𝔽₂.anticomm_eq_comm}
\leanok
\uses{def:DoubleComplexF2.dv, def:DoubleComplexF2.dh, thm:DoubleComplexF2.comm}
Let $E : \operatorname{DoubleComplex}_{\mathbb{F}_2}$. For all $p, q \in \mathbb{Z}$, the anticommutativity condition holds:
\[
  \partial^v_{p-1,q} \circ \partial^h_{p,q} + \partial^h_{p,q-1} \circ \partial^v_{p,q} = 0.
\]
That is, over $\mathbb{F}_2$, anticommutativity and commutativity of the differentials coincide, since $-1 = 1$ in characteristic $2$.
\end{theorem}

\begin{proof}
\leanok
\uses{def:DoubleComplexF2.dv, def:DoubleComplexF2.dh, thm:DoubleComplexF2.comm}
Rewriting the left-hand side using the commutativity result $\partial^v_{p-1,q} \circ \partial^h_{p,q} = \partial^h_{p,q-1} \circ \partial^v_{p,q}$ (Theorem~\ref{thm:DoubleComplexF2.comm}), it suffices to show that for all $x \in E_{p,q}$,
\[
  (\partial^h_{p,q-1} \circ \partial^v_{p,q})(x) + (\partial^h_{p,q-1} \circ \partial^v_{p,q})(x) = 0.
\]
We use the fact that $2 = 0$ in $\mathbb{F}_2$, which is verified by computation ($\mathtt{decide}$). Then $2 \cdot v = 0$ for any $v \in \mathbb{F}_2$. Expressing the sum as $2 \cdot (\partial^h_{p,q-1} \circ \partial^v_{p,q})(x)$ via $\mathtt{two\_smul}$ and substituting $2 = 0$ with $\mathtt{zero\_smul}$, we obtain $0$, completing the proof.
\end{proof}

%--- Def_10: TotalComplex ---
\chapter{Def 10: Total Complex}

\begin{definition}[Total Complex of a Double Complex]
\label{def:DoubleComplexF2.totalComplex}
\lean{DoubleComplex𝔽₂.totalComplex}
\leanok
\uses{def:DoubleComplexF2}

Let $E = (E_{\bullet,\bullet}, \partial^v, \partial^h)$ be a double complex of $\mathbb{F}_2$-vector spaces (see Definition~\ref{def:DoubleComplexF2}). The \emph{total complex} $\operatorname{Tot}(E)$ is the chain complex over $\mathbb{F}_2$ indexed by $\mathbb{Z}$ defined as follows:
\begin{itemize}
  \item The $n$-th object is the direct sum
    \[
      \operatorname{Tot}(E)_n \;=\; \bigoplus_{p+q=n} E_{p,q},
    \]
    taken over all pairs of integers $(p,q)$ with $p + q = n$.
  \item The differential $\partial_n \colon \operatorname{Tot}(E)_n \to \operatorname{Tot}(E)_{n-1}$ restricts to each summand $E_{p,q}$ as
    \[
      \partial_n\big|_{E_{p,q}} \;=\; \partial^v_{p,q} + \partial^h_{p,q},
    \]
    where $\partial^v_{p,q} \colon E_{p,q} \to E_{p,q-1}$ maps into the $(p, q-1)$-summand and $\partial^h_{p,q} \colon E_{p,q} \to E_{p-1,q}$ maps into the $(p-1,q)$-summand of $\operatorname{Tot}(E)_{n-1}$.
\end{itemize}

Since $\operatorname{ModuleCat}(\mathbb{F}_2)$ has all coproducts, the total complex exists. This construction is realized via Mathlib's \texttt{HomologicalComplex\textsubscript{2}.total} with the descending complex shape on $\mathbb{Z}$.

\textit{Sign convention:} Mathlib uses sign maps $\varepsilon(n) = (-1)^n$ (implemented by \texttt{Int.negOnePow}) for the \texttt{TensorSigns} instance on \texttt{ComplexShape.down}~$\mathbb{Z}$, satisfying: for $p + q + 1 = n$, the sign condition $\varepsilon'_{\mathrm{succ}}$ gives $\varepsilon(p+1) = -\varepsilon(p)$, verified by simplification using \texttt{Int.negOnePow\_succ} and $-(-x)=x$. Over $\mathbb{F}_2$ these signs are trivial (since $-1 = 1$), so Mathlib's construction agrees with the sign-free definition in the paper.
\end{definition}

\begin{definition}[Inclusion into Total Complex]
\label{def:DoubleComplexF2.iotaTotalComplex}
\lean{DoubleComplex𝔽₂.ιTotalComplex}
\leanok
\uses{def:DoubleComplexF2.totalComplex}

Let $E$ be a double complex of $\mathbb{F}_2$-vector spaces. For integers $p, q, n$ with $p + q = n$, the \emph{inclusion map}
\[
  \iota_{p,q}^n \colon E_{p,q} \;\longrightarrow\; \operatorname{Tot}(E)_n
\]
is the canonical injection of the summand $E_{p,q}$ into the coproduct $\operatorname{Tot}(E)_n = \bigoplus_{p'+q'=n} E_{p',q'}$.
\end{definition}

\begin{theorem}[Differential of Total Complex Squares to Zero]
\label{thm:DoubleComplexF2.totalComplex_d_sq_zero}
\lean{DoubleComplex𝔽₂.totalComplex_d_sq_zero}
\leanok
\uses{def:DoubleComplexF2.totalComplex, def:DoubleComplexF2}

Let $E$ be a double complex of $\mathbb{F}_2$-vector spaces. For all integers $n, m, k$, the composition of consecutive differentials of $\operatorname{Tot}(E)$ vanishes:
\[
  \partial_m \circ \partial_n \;=\; 0 \;\colon\; \operatorname{Tot}(E)_n \;\longrightarrow\; \operatorname{Tot}(E)_k.
\]
\end{theorem}

\begin{proof}
\leanok
\uses{def:DoubleComplexF2.totalComplex}

This follows directly from the general fact that in any homological complex, the composition of two consecutive differentials is zero, applied to $\operatorname{Tot}(E)$ via \texttt{d\_comp\_d}.

Conceptually, $\partial^2 = (\partial^v + \partial^h)^2 = (\partial^v)^2 + (\partial^v \partial^h + \partial^h \partial^v) + (\partial^h)^2$. Each term vanishes: $(\partial^v)^2 = 0$ and $(\partial^h)^2 = 0$ by the double complex axioms, and $\partial^v \partial^h + \partial^h \partial^v = 0$ over $\mathbb{F}_2$ because $\partial^v \partial^h = \partial^h \partial^v$ (commutativity up to sign, and $-1 = 1$ in $\mathbb{F}_2$).
\end{proof}

\begin{theorem}[Objects of Total Complex as Coproducts]
\label{thm:DoubleComplexF2.totalComplex_obj_eq}
\lean{DoubleComplex𝔽₂.totalComplex_obj_eq}
\leanok
\uses{def:DoubleComplexF2.totalComplex, def:DoubleComplexF2}

Let $E$ be a double complex of $\mathbb{F}_2$-vector spaces. For each integer $n$, the $n$-th object of the total complex equals the coproduct
\[
  \operatorname{Tot}(E)_n \;=\; \bigl(E^{\mathrm{gr}}\bigr)^{\operatorname{map}}_n,
\]
where $E^{\mathrm{gr}}$ denotes the underlying graded object of $E$ and the superscript denotes the graded-object mapping along the complex-shape projection $\pi \colon \mathbb{Z} \times \mathbb{Z} \to \mathbb{Z}$, $(p,q) \mapsto p+q$.
\end{theorem}

\begin{proof}
\leanok
\uses{def:DoubleComplexF2.totalComplex}

This holds by reflexivity: the definition of \texttt{totalComplex} as \texttt{E.total (ComplexShape.down}~$\mathbb{Z}$\texttt{)} and the definition of the total complex object in Mathlib make both sides definitionally equal.
\end{proof}

\begin{theorem}[Extensionality for Morphisms from Total Complex]
\label{thm:DoubleComplexF2.totalComplex_hom_ext}
\lean{DoubleComplex𝔽₂.totalComplex_hom_ext}
\leanok
\uses{def:DoubleComplexF2.totalComplex, def:DoubleComplexF2.iotaTotalComplex}

Let $E$ be a double complex of $\mathbb{F}_2$-vector spaces, let $A$ be an $\mathbb{F}_2$-module, and let $n \in \mathbb{Z}$. Suppose $f, g \colon \operatorname{Tot}(E)_n \to A$ are two morphisms such that for all integers $p, q$ with $p + q = n$,
\[
  f \circ \iota_{p,q}^n \;=\; g \circ \iota_{p,q}^n.
\]
Then $f = g$.
\end{theorem}

\begin{proof}
\leanok
\uses{def:DoubleComplexF2.iotaTotalComplex}

This follows from the universal property of coproducts: since $\operatorname{Tot}(E)_n = \bigoplus_{p+q=n} E_{p,q}$, a morphism out of it is uniquely determined by its restrictions to each summand. Formally, this is an application of \texttt{HomologicalComplex\textsubscript{2}.total.hom\_ext} with the descending complex shape on $\mathbb{Z}$, using the hypothesis $h$ that $\iota_{p,q}^n \gg f = \iota_{p,q}^n \gg g$ for all $p, q$ with $p + q = n$.
\end{proof}

%--- Def_11: TensorProductDoubleComplex ---
\chapter{Def 11: Tensor Product Double Complex}

\begin{definition}[Tensor Object]
\label{def:ChainComplexF2.tensorObj'}
\lean{ChainComplex𝔽₂.tensorObj'}
\leanok
\uses{def:ChainComplexF2}
Let $C$ and $D$ be chain complexes over $\mathbb{F}_2$. The graded object underlying the tensor
product double complex is defined, for integers $p, q \in \mathbb{Z}$, by
\[
  (C \boxtimes D)_{p,q} \;=\; C_p \otimes_{\mathbb{F}_2} D_q,
\]
where $\otimes_{\mathbb{F}_2}$ denotes the tensor product of $\mathbb{F}_2$-modules.
\end{definition}

\begin{definition}[Horizontal Differential]
\label{def:ChainComplexF2.tensorDh}
\lean{ChainComplex𝔽₂.tensorDh}
\leanok
\uses{def:ChainComplexF2, def:ChainComplexF2.tensorObj'}
Let $C$ and $D$ be chain complexes over $\mathbb{F}_2$. The \emph{horizontal differential}
\[
  \partial^h_{p,p',q} \;:\; C_p \otimes D_q \;\longrightarrow\; C_{p'} \otimes D_q
\]
is defined by
\[
  \partial^h_{p,p',q} \;=\; \partial^C_{p \to p'} \otimes \operatorname{id}_{D_q},
\]
that is, it applies the differential $\partial^C$ of $C$ in the first (horizontal, $p$-)direction
and the identity on $D_q$.
\end{definition}

\begin{definition}[Vertical Differential]
\label{def:ChainComplexF2.tensorDv}
\lean{ChainComplex𝔽₂.tensorDv}
\leanok
\uses{def:ChainComplexF2, def:ChainComplexF2.tensorObj'}
Let $C$ and $D$ be chain complexes over $\mathbb{F}_2$. The \emph{vertical differential}
\[
  \partial^v_{p,q,q'} \;:\; C_p \otimes D_q \;\longrightarrow\; C_p \otimes D_{q'}
\]
is defined by
\[
  \partial^v_{p,q,q'} \;=\; \operatorname{id}_{C_p} \otimes \partial^D_{q \to q'},
\]
that is, it applies the identity on $C_p$ and the differential $\partial^D$ of $D$ in the
second (vertical, $q$-)direction.
\end{definition}

\begin{definition}[Tensor Product Double Complex]
\label{def:ChainComplexF2.tensorDoubleComplex}
\lean{ChainComplex𝔽₂.tensorDoubleComplex}
\leanok
\uses{def:ChainComplexF2, def:DoubleComplexF2, def:ChainComplexF2.tensorObj', def:ChainComplexF2.tensorDh, def:ChainComplexF2.tensorDv}
Let $C$ and $D$ be chain complexes over $\mathbb{F}_2$. The \emph{tensor product double complex}
$C \boxtimes D$ is the double complex (in the sense of $\operatorname{DoubleComplex}_{\mathbb{F}_2}$)
defined as follows:
\begin{itemize}
  \item \textbf{Objects:} $(C \boxtimes D)_{p,q} = C_p \otimes_{\mathbb{F}_2} D_q$ for all
    $p, q \in \mathbb{Z}$.
  \item \textbf{Horizontal differential:}
    $\partial^h_{p,q} = \partial^C_p \otimes \operatorname{id}_{D_q} : C_p \otimes D_q \to C_{p-1} \otimes D_q$.
  \item \textbf{Vertical differential:}
    $\partial^v_{p,q} = \operatorname{id}_{C_p} \otimes \partial^D_q : C_p \otimes D_q \to C_p \otimes D_{q-1}$.
\end{itemize}
The double complex conditions are verified as follows. When the complex shape relation
$\operatorname{Rel}(p, p')$ (respectively $\operatorname{Rel}(q, q')$) does not hold, the
corresponding differential vanishes. The horizontal and vertical differentials each square to zero:
$\partial^h \circ \partial^h = 0$ and $\partial^v \circ \partial^v = 0$. Moreover, the horizontal and
vertical differentials commute: $\partial^h \circ \partial^v = \partial^v \circ \partial^h$, which
over $\mathbb{F}_2$ replaces the usual anticommutativity requirement. This is a consequence of the
whisker exchange identity $(\partial^C \otimes \operatorname{id}) \circ (\operatorname{id} \otimes \partial^D)
= (\operatorname{id} \otimes \partial^D) \circ (\partial^C \otimes \operatorname{id})$ in any monoidal
category.

The double complex is assembled via \texttt{HomologicalComplex\textsubscript{2}.ofGradedObject}
using the descending complex shape on $\mathbb{Z}$ in both directions.
\end{definition}

\begin{theorem}[Object of Tensor Double Complex]
\label{thm:ChainComplexF2.tensorDoubleComplex_obj}
\lean{ChainComplex𝔽₂.tensorDoubleComplex_obj}
\leanok
\uses{def:ChainComplexF2.tensorDoubleComplex, def:ChainComplexF2.tensorObj'}
Let $C, D$ be chain complexes over $\mathbb{F}_2$ and let $p, q \in \mathbb{Z}$. Then
\[
  (C \boxtimes D)_{p,q} \;=\; C_p \otimes_{\mathbb{F}_2} D_q.
\]
\end{theorem}

\begin{proof}
\leanok
\uses{def:ChainComplexF2.tensorDoubleComplex, def:ChainComplexF2.tensorObj'}
This holds by reflexivity: the object at position $(p,q)$ in the tensor double complex is
definitionally equal to $C_p \otimes D_q$ by construction in the definition of
$\operatorname{tensorDoubleComplex}$.
\end{proof}

\begin{theorem}[Vertical Differential of Tensor Double Complex]
\label{thm:ChainComplexF2.tensorDoubleComplex_dv_eq}
\lean{ChainComplex𝔽₂.tensorDoubleComplex_dv_eq}
\leanok
\uses{def:ChainComplexF2.tensorDoubleComplex, def:ChainComplexF2.tensorDv}
Let $C, D$ be chain complexes over $\mathbb{F}_2$ and let $p, q \in \mathbb{Z}$. The vertical
differential of $C \boxtimes D$ at position $(p,q)$ is
\[
  \bigl((C \boxtimes D)_p\bigr)_{q \to q-1} \;=\; \partial^v_{p,q,q-1}
  \;=\; \operatorname{id}_{C_p} \otimes \partial^D_{q \to q-1}.
\]
\end{theorem}

\begin{proof}
\leanok
\uses{def:ChainComplexF2.tensorDoubleComplex, def:ChainComplexF2.tensorDv}
This holds by reflexivity: the vertical differential at $(p, q)$ is definitionally equal to
$\operatorname{tensorDv}_{p,q,q-1}$ by construction of $\operatorname{tensorDoubleComplex}$.
\end{proof}

\begin{theorem}[Horizontal Differential of Tensor Double Complex]
\label{thm:ChainComplexF2.tensorDoubleComplex_dh_eq}
\lean{ChainComplex𝔽₂.tensorDoubleComplex_dh_eq}
\leanok
\uses{def:ChainComplexF2.tensorDoubleComplex, def:ChainComplexF2.tensorDh}
Let $C, D$ be chain complexes over $\mathbb{F}_2$ and let $p, q \in \mathbb{Z}$. The horizontal
differential of $C \boxtimes D$ from column $p$ to column $p-1$ at row $q$ is
\[
  \bigl((C \boxtimes D)_{p \to p-1}\bigr)_q \;=\; \partial^h_{p,p-1,q}
  \;=\; \partial^C_{p \to p-1} \otimes \operatorname{id}_{D_q}.
\]
\end{theorem}

\begin{proof}
\leanok
\uses{def:ChainComplexF2.tensorDoubleComplex, def:ChainComplexF2.tensorDh}
This holds by reflexivity: the horizontal differential at $(p,q)$ is definitionally equal to
$\operatorname{tensorDh}_{p,p-1,q}$ by construction of $\operatorname{tensorDoubleComplex}$.
\end{proof}

\begin{definition}[Tensor Product Complex]
\label{def:ChainComplexF2.tensorComplex}
\lean{ChainComplex𝔽₂.tensorComplex}
\leanok
\uses{def:ChainComplexF2.tensorDoubleComplex, def:DoubleComplexF2.totalComplex}
Let $C$ and $D$ be chain complexes over $\mathbb{F}_2$. The \emph{tensor product complex}
$C \otimes D$ is defined as the total complex of the tensor product double complex:
\[
  C \otimes D \;:=\; \operatorname{Tot}(C \boxtimes D).
\]
Concretely, its degree-$n$ component is
\[
  (C \otimes D)_n \;=\; \bigoplus_{p + q = n} C_p \otimes_{\mathbb{F}_2} D_q,
\]
and its differential is $\partial = \partial^C \otimes \operatorname{id} + \operatorname{id} \otimes \partial^D$.
Over $\mathbb{F}_2$, the usual sign factor $(-1)^p$ is trivial since $-1 = 1$.
\end{definition}

\begin{definition}[Inclusion into Tensor Complex]
\label{def:ChainComplexF2.iotaTensorComplex}
\lean{ChainComplex𝔽₂.ιTensorComplex}
\leanok
\uses{def:ChainComplexF2.tensorDoubleComplex, def:ChainComplexF2.tensorComplex, def:DoubleComplexF2.iotaTotalComplex}
Let $C$ and $D$ be chain complexes over $\mathbb{F}_2$, and let $p, q, n \in \mathbb{Z}$ with
$p + q = n$. The canonical inclusion of the summand $C_p \otimes D_q$ into $(C \otimes D)_n$
is the map
\[
  \iota_{p,q}^n \;:\; C_p \otimes D_q \;\longrightarrow\; (C \otimes D)_n,
\]
defined as the total complex inclusion $\iota_{p,q,n}$ applied to the tensor product double
complex $C \boxtimes D$.
\end{definition}

\begin{definition}[Hypergraph Product Code]
\label{def:ChainComplexF2.hypergraphProductCode}
\lean{ChainComplex𝔽₂.hypergraphProductCode}
\leanok
\uses{def:ChainComplexF2.tensorComplex, def:ClassicalCode.parityCheckComplex}
Let $n_1, m_1, n_2, m_2 \in \mathbb{N}$, and let
\[
  \partial^C : \mathbb{F}_2^{n_1} \to \mathbb{F}_2^{m_1}, \qquad
  \partial^D : \mathbb{F}_2^{n_2} \to \mathbb{F}_2^{m_2}
\]
be linear maps over $\mathbb{F}_2$ (parity-check matrices of two classical codes). The
\emph{hypergraph product code} is the chain complex
\[
  \operatorname{HGP}(\partial^C, \partial^D)
  \;:=\; \operatorname{parityCheckComplex}(\partial^C)
    \otimes \operatorname{parityCheckComplex}(\partial^D),
\]
i.e., the tensor product complex of the parity-check complexes associated to $\partial^C$ and
$\partial^D$. Concretely, the underlying 1-complexes are
\[
  C: \quad \mathbb{F}_2^{n_1} \xrightarrow{\partial^C} \mathbb{F}_2^{m_1}, \qquad
  D: \quad \mathbb{F}_2^{n_2} \xrightarrow{\partial^D} \mathbb{F}_2^{m_2},
\]
and their tensor product double complex has three nontrivial degrees:
\begin{itemize}
  \item Degree $2$: $\mathbb{F}_2^{n_1} \otimes \mathbb{F}_2^{n_2}$,
  \item Degree $1$: $(\mathbb{F}_2^{n_1} \otimes \mathbb{F}_2^{m_2}) \oplus (\mathbb{F}_2^{m_1} \otimes \mathbb{F}_2^{n_2})$,
  \item Degree $0$: $\mathbb{F}_2^{m_1} \otimes \mathbb{F}_2^{m_2}$.
\end{itemize}
The differential from degree $2$ to degree $1$ plays the role of $H_Z^\top$, and the differential
from degree $1$ to degree $0$ plays the role of $H_X$ in the CSS code convention. The CSS code
structure is read off from the total complex shifted by $1$ to match the convention
$C_1 \xrightarrow{H_Z^\top} C_0 \xrightarrow{H_X} C_{-1}$.
\end{definition}

%--- Thm_2: SmallDoubleComplexHomology ---
\chapter{Thm 2: Small Double Complex Homology}

\begin{definition}[Small $2 \times 2$ Double Complex]
\label{def:SmallDC}
\lean{SmallDC}
\leanok

A \emph{small double complex} $E$ over $\mathbb{F}_2$ is a $2 \times 2$ double complex consisting of four $\mathbb{F}_2$-modules $A_{0,0}$, $A_{0,1}$, $A_{1,0}$, $A_{1,1}$ together with vertical linear maps
\[
  \partial^v_0 : A_{0,1} \to A_{0,0}, \quad \partial^v_1 : A_{1,1} \to A_{1,0}
\]
and horizontal linear maps
\[
  \partial^h_0 : A_{1,0} \to A_{0,0}, \quad \partial^h_1 : A_{1,1} \to A_{0,1}
\]
satisfying the commutativity condition
\[
  \partial^v_0 \circ \partial^h_1 = \partial^h_0 \circ \partial^v_1.
\]
This is the data of a double complex supported in bi-degrees $(p,q) \in \{0,1\}^2$.
\end{definition}

\begin{definition}[Total Complex Differential $\partial_2$]
\label{def:SmallDC.d2}
\lean{SmallDC.d₂}
\leanok
\uses{def:SmallDC}
Given a small double complex $E$, define the degree-2 total complex differential
\[
  \partial_2 : A_{1,1} \to A_{1,0} \times A_{0,1}, \quad x \mapsto (\partial^v_1(x),\, \partial^h_1(x)).
\]
This is the product map $\partial^v_1 \times \partial^h_1$.
\end{definition}

\begin{definition}[Total Complex Differential $\partial_1$]
\label{def:SmallDC.d1}
\lean{SmallDC.d₁}
\leanok
\uses{def:SmallDC}
Given a small double complex $E$, define the degree-1 total complex differential
\[
  \partial_1 : A_{1,0} \times A_{0,1} \to A_{0,0}, \quad (a,b) \mapsto \partial^h_0(a) + \partial^v_0(b).
\]
This is the coproduct (direct sum) map $[\partial^h_0,\, \partial^v_0]$.
\end{definition}

\begin{lemma}[First Component of $\partial_2$]
\label{lem:SmallDC.d2_fst}
\lean{SmallDC.d₂_fst}
\leanok
\uses{def:SmallDC.d2}
For all $x \in A_{1,1}$, the first component of $\partial_2(x)$ equals $\partial^v_1(x)$:
\[
  (\partial_2(x))_1 = \partial^v_1(x).
\]
\end{lemma}

\begin{proof}
\leanok
\uses{def:SmallDC.d2}
This holds by applying the identity $\operatorname{fst} \circ (\partial^v_1 \times \partial^h_1) = \partial^v_1$, which follows from the definition of the product map and the projection axiom $\operatorname{fst}(\partial^v_1(x), \partial^h_1(x)) = \partial^v_1(x)$. Verified by simplification using the component projections of the product linear map.
\end{proof}

\begin{lemma}[Second Component of $\partial_2$]
\label{lem:SmallDC.d2_snd}
\lean{SmallDC.d₂_snd}
\leanok
\uses{def:SmallDC.d2}
For all $x \in A_{1,1}$, the second component of $\partial_2(x)$ equals $\partial^h_1(x)$:
\[
  (\partial_2(x))_2 = \partial^h_1(x).
\]
\end{lemma}

\begin{proof}
\leanok
\uses{def:SmallDC.d2}
This holds by applying the identity $\operatorname{snd} \circ (\partial^v_1 \times \partial^h_1) = \partial^h_1$, verified by simplification using the second projection of the product linear map.
\end{proof}

\begin{lemma}[Total Complex is a Complex]
\label{lem:SmallDC.d1_comp_d2}
\lean{SmallDC.d₁_comp_d₂}
\leanok
\uses{def:SmallDC.d1, def:SmallDC.d2, lem:SmallDC.d2_fst, lem:SmallDC.d2_snd}
The composition $\partial_1 \circ \partial_2 = 0$.
\end{lemma}

\begin{proof}
\leanok
\uses{def:SmallDC.d1, def:SmallDC.d2, lem:SmallDC.d2_fst, lem:SmallDC.d2_snd}
By extensionality, let $x \in A_{1,1}$. We compute:
\[
  (\partial_1 \circ \partial_2)(x) = \partial_1(\partial_2(x)) = \partial^h_0((\partial_2(x))_1) + \partial^v_0((\partial_2(x))_2).
\]
Rewriting using the component lemmas $(\partial_2(x))_1 = \partial^v_1(x)$ and $(\partial_2(x))_2 = \partial^h_1(x)$, we get
\[
  \partial^h_0(\partial^v_1(x)) + \partial^v_0(\partial^h_1(x)).
\]
By the commutativity condition $\partial^v_0 \circ \partial^h_1 = \partial^h_0 \circ \partial^v_1$, this equals
\[
  \partial^h_0(\partial^v_1(x)) + \partial^h_0(\partial^v_1(x)).
\]
Since we are over $\mathbb{F}_2$, any element added to itself is zero, so this sum is $0$.
\end{proof}

\begin{definition}[Degree-0 Total Homology]
\label{def:SmallDC.totH0}
\lean{SmallDC.totH₀}
\leanok
\uses{def:SmallDC.d1}
The degree-0 homology of the total complex is defined as the quotient
\[
  H_0(\operatorname{Tot}(E)) := A_{0,0} / \operatorname{im}(\partial_1).
\]
\end{definition}

\begin{definition}[Degree-1 Total Cycles]
\label{def:SmallDC.totZ1}
\lean{SmallDC.totZ₁}
\leanok
\uses{def:SmallDC.d1}
The degree-1 total cycles are
\[
  Z_1 := \ker(\partial_1) \subseteq A_{1,0} \times A_{0,1}.
\]
\end{definition}

\begin{definition}[Degree-1 Total Boundaries]
\label{def:SmallDC.totB1}
\lean{SmallDC.totB₁}
\leanok
\uses{def:SmallDC.d2}
The degree-1 total boundaries are
\[
  B_1 := \operatorname{im}(\partial_2) \subseteq A_{1,0} \times A_{0,1}.
\]
\end{definition}

\begin{lemma}[Boundaries Contained in Cycles]
\label{lem:SmallDC.totB1_le_totZ1}
\lean{SmallDC.totB₁_le_totZ₁}
\leanok
\uses{def:SmallDC.totB1, def:SmallDC.totZ1, lem:SmallDC.d1_comp_d2}
$B_1 \leq Z_1$, i.e., every boundary is a cycle.
\end{lemma}

\begin{proof}
\leanok
\uses{def:SmallDC.totB1, def:SmallDC.totZ1, lem:SmallDC.d1_comp_d2}
Let $x \in B_1$. We must show $x \in Z_1 = \ker(\partial_1)$. Since $x \in \operatorname{im}(\partial_2)$, obtain $y \in A_{1,1}$ such that $x = \partial_2(y)$. By the identity $\partial_1 \circ \partial_2 = 0$ (Lemma~\ref{lem:SmallDC.d1_comp_d2}), we have $\partial_1(x) = \partial_1(\partial_2(y)) = 0$, so $x \in \ker(\partial_1) = Z_1$.
\end{proof}

\begin{definition}[Boundaries as Submodule of Cycles]
\label{def:SmallDC.totB1InZ1}
\lean{SmallDC.totB₁InZ₁}
\leanok
\uses{def:SmallDC.totZ1, def:SmallDC.totB1}
We regard $B_1$ as a submodule of $Z_1$ via the comap of the inclusion $Z_1 \hookrightarrow A_{1,0} \times A_{0,1}$:
\[
  B_1^{Z} := B_1 \cdot_{\iota} Z_1 = B_1 \cap Z_1 \text{ (as a submodule of } Z_1\text{)}.
\]
\end{definition}

\begin{definition}[Degree-1 Total Homology]
\label{def:SmallDC.totH1}
\lean{SmallDC.totH₁}
\leanok
\uses{def:SmallDC.totZ1, def:SmallDC.totB1InZ1}
The degree-1 homology of the total complex is
\[
  H_1(\operatorname{Tot}(E)) := Z_1 / B_1^{Z}.
\]
\end{definition}

\begin{definition}[Degree-2 Total Homology]
\label{def:SmallDC.totH2}
\lean{SmallDC.totH₂}
\leanok
\uses{def:SmallDC.d2}
The degree-2 homology of the total complex is
\[
  H_2(\operatorname{Tot}(E)) := \ker(\partial_2),
\]
since there are no boundaries in degree 2 (as there are no non-zero modules in bi-degree $(p,q)$ with $p+q=3$).
\end{definition}

\begin{definition}[Vertical Homology Groups]
\label{def:SmallDC.vH00}
\lean{SmallDC.vH₀₀}
\leanok
\uses{def:SmallDC}
The vertical homology groups of $E$ are:
\begin{align*}
  H^v_{0,0} &:= A_{0,0} / \operatorname{im}(\partial^v_0), \\
  H^v_{1,0} &:= A_{1,0} / \operatorname{im}(\partial^v_1), \\
  H^v_{0,1} &:= \ker(\partial^v_0), \\
  H^v_{1,1} &:= \ker(\partial^v_1).
\end{align*}
Here $H^v_{p,0}$ is the degree-0 vertical homology in column $p$, and $H^v_{p,1}$ is the degree-1 vertical homology in column $p$.
\end{definition}

\begin{lemma}[Horizontal Map Preserves Vertical Boundaries]
\label{lem:SmallDC.dh0_maps_range}
\lean{SmallDC.dh₀_maps_range}
\leanok
\uses{def:SmallDC}
The map $\partial^h_0$ sends $\operatorname{im}(\partial^v_1)$ into $\operatorname{im}(\partial^v_0)$:
\[
  \operatorname{im}(\partial^v_1) \leq (\operatorname{im}(\partial^v_0))^{(\partial^h_0)},
\]
i.e., for every $x = \partial^v_1(z)$ in the image of $\partial^v_1$, we have $\partial^h_0(x) \in \operatorname{im}(\partial^v_0)$.
\end{lemma}

\begin{proof}
\leanok
\uses{def:SmallDC}
Let $x \in \operatorname{im}(\partial^v_1)$, so $x = \partial^v_1(z)$ for some $z \in A_{1,1}$. By the commutativity condition $\partial^v_0 \circ \partial^h_1 = \partial^h_0 \circ \partial^v_1$, we have
\[
  \partial^h_0(x) = \partial^h_0(\partial^v_1(z)) = \partial^v_0(\partial^h_1(z)) \in \operatorname{im}(\partial^v_0).
\]
\end{proof}

\begin{definition}[Induced Horizontal Map $\bar{\partial}^h_0$]
\label{def:SmallDC.barDh0}
\lean{SmallDC.barDh₀}
\leanok
\uses{def:SmallDC.vH00, lem:SmallDC.dh0_maps_range}
Since $\partial^h_0$ maps $\operatorname{im}(\partial^v_1)$ into $\operatorname{im}(\partial^v_0)$ (Lemma~\ref{lem:SmallDC.dh0_maps_range}), it descends to a well-defined $\mathbb{F}_2$-linear map on vertical homology:
\[
  \bar{\partial}^h_0 : H^v_{1,0} \to H^v_{0,0}, \quad [a] \mapsto [\partial^h_0(a)].
\]
This is constructed as a quotient map $\operatorname{mapQ}$ applied to $\partial^h_0$.
\end{definition}

\begin{lemma}[Horizontal Map Preserves Vertical Cycles]
\label{lem:SmallDC.dh1_maps_ker}
\lean{SmallDC.dh₁_maps_ker}
\leanok
\uses{def:SmallDC}
For every $x \in \ker(\partial^v_1)$, we have $\partial^h_1(x) \in \ker(\partial^v_0)$.
\end{lemma}

\begin{proof}
\leanok
\uses{def:SmallDC}
Let $x \in \ker(\partial^v_1)$, so $\partial^v_1(x) = 0$. We must show $\partial^v_0(\partial^h_1(x)) = 0$. By the commutativity condition $\partial^v_0 \circ \partial^h_1 = \partial^h_0 \circ \partial^v_1$:
\[
  \partial^v_0(\partial^h_1(x)) = \partial^h_0(\partial^v_1(x)) = \partial^h_0(0) = 0.
\]
Rewriting using the membership $x \in \ker(\partial^v_1)$ and applying linearity, we conclude $\partial^h_1(x) \in \ker(\partial^v_0)$.
\end{proof}

\begin{definition}[Induced Horizontal Map $\bar{\partial}^h_1$]
\label{def:SmallDC.barDh1}
\lean{SmallDC.barDh₁}
\leanok
\uses{def:SmallDC.vH00, lem:SmallDC.dh1_maps_ker}
Since $\partial^h_1$ maps $\ker(\partial^v_1)$ into $\ker(\partial^v_0)$ (Lemma~\ref{lem:SmallDC.dh1_maps_ker}), it restricts to a well-defined $\mathbb{F}_2$-linear map on vertical homology:
\[
  \bar{\partial}^h_1 : H^v_{1,1} \to H^v_{0,1}, \quad \langle x, hx \rangle \mapsto \langle \partial^h_1(x), \text{(proof)} \rangle.
\]
This is constructed by codomain restriction of the domain-restricted map $\partial^h_1|_{\ker(\partial^v_1)}$.
\end{definition}

\begin{definition}[Iterated Homology at Degree 0]
\label{def:SmallDC.itH0}
\lean{SmallDC.itH₀}
\leanok
\uses{def:SmallDC.vH00, def:SmallDC.barDh0}
The iterated (horizontal) homology at total degree 0 is
\[
  \mathcal{H}_0 := H^v_{0,0} / \operatorname{im}(\bar{\partial}^h_0) = \bigoplus_{p+q=0} H_p(H_q(E_{\bullet,\bullet}, \partial^v), \bar{\partial}^h).
\]
\end{definition}

\begin{definition}[Iterated Homology at Degree 2]
\label{def:SmallDC.itH2}
\lean{SmallDC.itH₂}
\leanok
\uses{def:SmallDC.vH00, def:SmallDC.barDh1}
The iterated (horizontal) homology at total degree 2 is
\[
  \mathcal{H}_2 := \ker(\bar{\partial}^h_1) = \bigoplus_{p+q=2} H_p(H_q(E_{\bullet,\bullet}, \partial^v), \bar{\partial}^h).
\]
\end{definition}

\begin{definition}[First Factor of Iterated Homology at Degree 1]
\label{def:SmallDC.itH1_fst}
\lean{SmallDC.itH₁_fst}
\leanok
\uses{def:SmallDC.vH00, def:SmallDC.barDh0}
The first summand of the iterated homology at total degree 1, corresponding to bi-degree $(1,0)$, is
\[
  \mathcal{H}_{1,\text{fst}} := \ker(\bar{\partial}^h_0) \subseteq H^v_{1,0}.
\]
\end{definition}

\begin{definition}[Second Factor of Iterated Homology at Degree 1]
\label{def:SmallDC.itH1_snd}
\lean{SmallDC.itH₁_snd}
\leanok
\uses{def:SmallDC.vH00, def:SmallDC.barDh1}
The second summand of the iterated homology at total degree 1, corresponding to bi-degree $(0,1)$, is
\[
  \mathcal{H}_{1,\text{snd}} := H^v_{0,1} / \operatorname{im}(\bar{\partial}^h_1).
\]
\end{definition}

\begin{lemma}[Kernel of $\partial_2$ Decomposes]
\label{lem:SmallDC.ker_d2_eq}
\lean{SmallDC.ker_d₂_eq}
\leanok
\uses{def:SmallDC.d2}
The kernel of $\partial_2 = \partial^v_1 \times \partial^h_1$ is the intersection:
\[
  \ker(\partial_2) = \ker(\partial^v_1) \cap \ker(\partial^h_1).
\]
\end{lemma}

\begin{proof}
\leanok
\uses{def:SmallDC.d2}
This follows directly from the standard identity $\ker(f \times g) = \ker(f) \cap \ker(g)$ for product linear maps, applied to $\partial_2 = \partial^v_1 \times \partial^h_1$: an element $x$ satisfies $\partial_2(x) = 0$ if and only if $(\partial^v_1(x), \partial^h_1(x)) = (0, 0)$, which holds if and only if $x \in \ker(\partial^v_1) \cap \ker(\partial^h_1)$.
\end{proof}

\begin{lemma}[Value of $\bar{\partial}^h_1$]
\label{lem:SmallDC.barDh1_val}
\lean{SmallDC.barDh₁_val}
\leanok
\uses{def:SmallDC.barDh1}
For $x \in H^v_{1,1} = \ker(\partial^v_1)$, the underlying value of $\bar{\partial}^h_1(x)$ is:
\[
  (\bar{\partial}^h_1(x))_{\text{val}} = \partial^h_1(x_{\text{val}}).
\]
\end{lemma}

\begin{proof}
\leanok
\uses{def:SmallDC.barDh1}
This holds by reflexivity from the definition of $\bar{\partial}^h_1$ as a codomain restriction: $\bar{\partial}^h_1\langle x, hx \rangle = \langle \partial^h_1(x), \ldots \rangle$, so its underlying value is $\partial^h_1(x)$.
\end{proof}

\begin{lemma}[Kernel of $\bar{\partial}^h_1$]
\label{lem:SmallDC.barDh1_ker_iff}
\lean{SmallDC.barDh₁_ker_iff}
\leanok
\uses{def:SmallDC.barDh1, lem:SmallDC.barDh1_val}
For $x \in H^v_{1,1}$:
\[
  x \in \ker(\bar{\partial}^h_1) \iff x_{\text{val}} \in \ker(\partial^h_1).
\]
\end{lemma}

\begin{proof}
\leanok
\uses{def:SmallDC.barDh1, lem:SmallDC.barDh1_val}
We use the kernel membership criterion $x \in \ker(\bar{\partial}^h_1) \Leftrightarrow \bar{\partial}^h_1(x) = 0$. Since $H^v_{0,1}$ is a subtype, $\bar{\partial}^h_1(x) = 0$ holds in the subtype if and only if its underlying value equals 0 in $A_{0,1}$.
$(\Rightarrow)$: Suppose $\bar{\partial}^h_1(x) = 0$. Taking the underlying value of both sides: $(\bar{\partial}^h_1(x))_{\text{val}} = 0_{\text{val}} = 0$. By Lemma~\ref{lem:SmallDC.barDh1_val}, $\partial^h_1(x_{\text{val}}) = 0$, so $x_{\text{val}} \in \ker(\partial^h_1)$.
$(\Leftarrow)$: Suppose $x_{\text{val}} \in \ker(\partial^h_1)$, i.e., $\partial^h_1(x_{\text{val}}) = 0$. By subtype injectivity and Lemma~\ref{lem:SmallDC.barDh1_val}, the underlying value of $\bar{\partial}^h_1(x)$ is 0, so $\bar{\partial}^h_1(x) = 0$.
\end{proof}

\begin{theorem}[Isomorphism at Degree 2]
\label{thm:SmallDC.homologyEquiv_n2}
\lean{SmallDC.homologyEquiv_n2}
\leanok
\uses{def:SmallDC.totH2, def:SmallDC.itH2, lem:SmallDC.ker_d2_eq, lem:SmallDC.barDh1_ker_iff}
There is a canonical $\mathbb{F}_2$-linear isomorphism
\[
  H_2(\operatorname{Tot}(E)) \cong \mathcal{H}_2,
\]
i.e., $\ker(\partial_2) \cong \ker(\bar{\partial}^h_1)$.
\end{theorem}

\begin{proof}
\leanok
\uses{def:SmallDC.totH2, def:SmallDC.itH2, lem:SmallDC.ker_d2_eq, lem:SmallDC.barDh1_ker_iff}
We construct an explicit $\mathbb{F}_2$-linear equivalence. By Lemma~\ref{lem:SmallDC.ker_d2_eq}, $\ker(\partial_2) = \ker(\partial^v_1) \cap \ker(\partial^h_1)$.

\textbf{Forward map}: Given $\langle x, hx \rangle \in \ker(\partial_2)$, we have $x \in \ker(\partial^v_1) \cap \ker(\partial^h_1)$. From the intersection decomposition, extract $hx_1 : x \in \ker(\partial^v_1)$ and $hx_2 : x \in \ker(\partial^h_1)$. Define the image as $\langle \langle x, hx_1 \rangle, (\bar{\partial}^h_1\text{-kernel membership obtained from } hx_2 \text{ via Lemma~\ref{lem:SmallDC.barDh1_ker_iff}}) \rangle \in \ker(\bar{\partial}^h_1)$.

\textbf{Backward map}: Given $\langle \langle x, hx_v \rangle, hx_h \rangle$ where $\langle x, hx_v \rangle \in H^v_{1,1}$ and the pair lies in $\ker(\bar{\partial}^h_1)$, by Lemma~\ref{lem:SmallDC.barDh1_ker_iff} we have $x \in \ker(\partial^h_1)$. Combined with $hx_v : x \in \ker(\partial^v_1)$, Lemma~\ref{lem:SmallDC.ker_d2_eq} gives $x \in \ker(\partial_2)$.

\textbf{Linearity}: Both $\operatorname{map_add'}$ and $\operatorname{map_smul'}$ hold by subtype extensionality and reflexivity.

\textbf{Inverse laws}: Both left and right inverses hold by subtype extensionality and reflexivity.
\end{proof}

\begin{lemma}[Image of $\partial_1$ Decomposes]
\label{lem:SmallDC.range_d1_eq}
\lean{SmallDC.range_d₁_eq}
\leanok
\uses{def:SmallDC.d1}
The image of $\partial_1$ equals the join of the images of $\partial^h_0$ and $\partial^v_0$:
\[
  \operatorname{im}(\partial_1) = \operatorname{im}(\partial^h_0) + \operatorname{im}(\partial^v_0).
\]
\end{lemma}

\begin{proof}
\leanok
\uses{def:SmallDC.d1}
By extensionality in $A_{0,0}$, using the definition $\partial_1(a,b) = \partial^h_0(a) + \partial^v_0(b)$.
$(\subseteq)$: If $y \in \operatorname{im}(\partial_1)$, obtain $(a,b) \in A_{1,0} \times A_{0,1}$ with $\partial_1(a,b) = y$, i.e., $y = \partial^h_0(a) + \partial^v_0(b)$. Then $\partial^h_0(a) \in \operatorname{im}(\partial^h_0)$ and $\partial^v_0(b) \in \operatorname{im}(\partial^v_0)$, so $y \in \operatorname{im}(\partial^h_0) + \operatorname{im}(\partial^v_0)$.
$(\supseteq)$: If $y = u + v$ with $u = \partial^h_0(a) \in \operatorname{im}(\partial^h_0)$ and $v = \partial^v_0(b) \in \operatorname{im}(\partial^v_0)$, then $y = \partial_1(a,b) \in \operatorname{im}(\partial_1)$.
\end{proof}

\begin{lemma}[Range of $\bar{\partial}^h_0$ in Terms of Sum]
\label{lem:SmallDC.barDh0_range_eq}
\lean{SmallDC.barDh₀_range_eq}
\leanok
\uses{def:SmallDC.barDh0, def:SmallDC.vH00}
The image of $\bar{\partial}^h_0$ in $H^v_{0,0} = A_{0,0}/\operatorname{im}(\partial^v_0)$ equals the image of $\operatorname{im}(\partial^h_0) + \operatorname{im}(\partial^v_0)$ under the quotient map:
\[
  \operatorname{im}(\bar{\partial}^h_0) = (\operatorname{im}(\partial^h_0) + \operatorname{im}(\partial^v_0)) / \operatorname{im}(\partial^v_0).
\]
\end{lemma}

\begin{proof}
\leanok
\uses{def:SmallDC.barDh0, def:SmallDC.vH00}
By extensionality in $H^v_{0,0}$.
$(\subseteq)$: Given $y \in \operatorname{im}(\bar{\partial}^h_0)$, obtain $[a] \in H^v_{1,0}$ with $\bar{\partial}^h_0([a]) = y$. By induction on the quotient representative $a \in A_{1,0}$, the element $\partial^h_0(a) \in \operatorname{im}(\partial^h_0) \subseteq \operatorname{im}(\partial^h_0) + \operatorname{im}(\partial^v_0)$, and its image under $\pi = \operatorname{mkQ}(\operatorname{im}(\partial^v_0))$ equals $y$.
$(\supseteq)$: Given $y = \pi(z)$ with $z \in \operatorname{im}(\partial^h_0) + \operatorname{im}(\partial^v_0)$, write $z = \partial^h_0(a) + v$ with $a \in A_{1,0}$ and $v \in \operatorname{im}(\partial^v_0)$. Then $\pi(v) = 0$ in the quotient, so $\pi(z) = \pi(\partial^h_0(a)) + 0 = \bar{\partial}^h_0([a]) \in \operatorname{im}(\bar{\partial}^h_0)$.
\end{proof}

\begin{theorem}[Isomorphism at Degree 0]
\label{thm:SmallDC.homologyEquiv_n0}
\lean{SmallDC.homologyEquiv_n0}
\leanok
\uses{def:SmallDC.totH0, def:SmallDC.itH0, lem:SmallDC.range_d1_eq, lem:SmallDC.barDh0_range_eq}
There is a canonical $\mathbb{F}_2$-linear isomorphism
\[
  H_0(\operatorname{Tot}(E)) \cong \mathcal{H}_0,
\]
i.e., $A_{0,0}/\operatorname{im}(\partial_1) \cong H^v_{0,0}/\operatorname{im}(\bar{\partial}^h_0)$.
\end{theorem}

\begin{proof}
\leanok
\uses{def:SmallDC.totH0, def:SmallDC.itH0, lem:SmallDC.range_d1_eq, lem:SmallDC.barDh0_range_eq}
The isomorphism is constructed as a composition of three quotient equivalences. By Lemma~\ref{lem:SmallDC.range_d1_eq}, $\operatorname{im}(\partial_1) = \operatorname{im}(\partial^h_0) + \operatorname{im}(\partial^v_0)$, giving
\[
  A_{0,0}/\operatorname{im}(\partial_1) \cong A_{0,0}/(\operatorname{im}(\partial^h_0) + \operatorname{im}(\partial^v_0)).
\]
By the second isomorphism theorem for modules (quotient-quotient equivalence), since $\operatorname{im}(\partial^v_0) \leq \operatorname{im}(\partial^h_0) + \operatorname{im}(\partial^v_0)$:
\[
  A_{0,0}/(\operatorname{im}(\partial^h_0) + \operatorname{im}(\partial^v_0)) \cong (A_{0,0}/\operatorname{im}(\partial^v_0)) / ((\operatorname{im}(\partial^h_0) + \operatorname{im}(\partial^v_0))/\operatorname{im}(\partial^v_0)).
\]
By Lemma~\ref{lem:SmallDC.barDh0_range_eq}, $(\operatorname{im}(\partial^h_0) + \operatorname{im}(\partial^v_0))/\operatorname{im}(\partial^v_0) = \operatorname{im}(\bar{\partial}^h_0)$, so the result is $H^v_{0,0}/\operatorname{im}(\bar{\partial}^h_0) = \mathcal{H}_0$.
\end{proof}

\begin{definition}[Projection Map $\pi$]
\label{def:SmallDC.piMap}
\lean{SmallDC.piMap}
\leanok
\uses{def:SmallDC.totZ1, def:SmallDC.vH00}
Define the projection $\pi : Z_1 \to H^v_{1,0}$ by composing the inclusion $Z_1 \hookrightarrow A_{1,0} \times A_{0,1}$, the first projection $\operatorname{fst} : A_{1,0} \times A_{0,1} \to A_{1,0}$, and the quotient map $\pi_0 : A_{1,0} \to H^v_{1,0}$:
\[
  \pi : Z_1 \hookrightarrow A_{1,0} \times A_{0,1} \xrightarrow{\operatorname{pr}_1} A_{1,0} \twoheadrightarrow H^v_{1,0}, \quad (a,b) \mapsto [a].
\]
\end{definition}

\begin{lemma}[$\pi$ Lands in $\ker(\bar{\partial}^h_0)$]
\label{lem:SmallDC.piMap_mem_ker}
\lean{SmallDC.piMap_mem_ker}
\leanok
\uses{def:SmallDC.piMap, def:SmallDC.barDh0}
For every $z \in Z_1$, $\pi(z) \in \ker(\bar{\partial}^h_0)$.
\end{lemma}

\begin{proof}
\leanok
\uses{def:SmallDC.piMap, def:SmallDC.barDh0}
Let $z = (a,b) \in Z_1 = \ker(\partial_1)$. We must show $\bar{\partial}^h_0([a]) = 0$ in $H^v_{0,0}$. By definition, $\bar{\partial}^h_0([a]) = [\partial^h_0(a)]$ in $H^v_{0,0} = A_{0,0}/\operatorname{im}(\partial^v_0)$, so we need $\partial^h_0(a) \in \operatorname{im}(\partial^v_0)$. Since $(a,b) \in \ker(\partial_1)$, we have $\partial^h_0(a) + \partial^v_0(b) = 0$, hence $\partial^h_0(a) = -\partial^v_0(b)$. Over $\mathbb{F}_2$, negation is the identity, so $\partial^h_0(a) = \partial^v_0(b) \in \operatorname{im}(\partial^v_0)$.
\end{proof}

\begin{definition}[Lifted Projection $\pi'$]
\label{def:SmallDC.piMap'}
\lean{SmallDC.piMap'}
\leanok
\uses{def:SmallDC.piMap, lem:SmallDC.piMap_mem_ker}
By Lemma~\ref{lem:SmallDC.piMap_mem_ker}, $\pi$ lands in $\ker(\bar{\partial}^h_0) = \mathcal{H}_{1,\text{fst}}$. Define the codomain-restricted map:
\[
  \pi' : Z_1 \to \ker(\bar{\partial}^h_0), \quad z \mapsto (\pi(z), \text{proof that } \pi(z) \in \ker(\bar{\partial}^h_0)).
\]
\end{definition}

\begin{lemma}[$\pi'$ Vanishes on $B_1^Z$]
\label{lem:SmallDC.piMap'_vanishes_on_B1}
\lean{SmallDC.piMap'_vanishes_on_B₁}
\leanok
\uses{def:SmallDC.piMap', def:SmallDC.totB1InZ1, lem:SmallDC.d2_fst}
$B_1^Z \leq \ker(\pi')$, i.e., $\pi'$ vanishes on boundaries.
\end{lemma}

\begin{proof}
\leanok
\uses{def:SmallDC.piMap', def:SmallDC.totB1InZ1, lem:SmallDC.d2_fst}
Let $v = (a',b') \in B_1^Z \subseteq Z_1$. We must show $\pi'(v) = 0$, i.e., the image of $\pi'$ at $v$ is zero in $\ker(\bar{\partial}^h_0)$. Since $v \in B_1^Z$, we have $(a',b') \in \operatorname{im}(\partial_2)$; obtain $x \in A_{1,1}$ with $\partial_2(x) = (a',b')$. By Lemma~\ref{lem:SmallDC.d2_fst}, $a' = (\partial_2(x))_1 = \partial^v_1(x) \in \operatorname{im}(\partial^v_1)$. Therefore $[a'] = 0$ in $H^v_{1,0} = A_{1,0}/\operatorname{im}(\partial^v_1)$, so $\pi'(v) = 0$.
\end{proof}

\begin{definition}[Descended Projection $\bar{\pi}$]
\label{def:SmallDC.piBar}
\lean{SmallDC.piBar}
\leanok
\uses{def:SmallDC.totH1, def:SmallDC.itH1_fst, def:SmallDC.piMap', lem:SmallDC.piMap'_vanishes_on_B1}
Since $\pi'$ vanishes on $B_1^Z$ (Lemma~\ref{lem:SmallDC.piMap'_vanishes_on_B1}), it descends to a well-defined map on the degree-1 homology:
\[
  \bar{\pi} : H_1(\operatorname{Tot}(E)) \to \mathcal{H}_{1,\text{fst}}, \quad [(a,b)] \mapsto [a].
\]
\end{definition}

\begin{definition}[Inclusion Map $\iota$]
\label{def:SmallDC.iotaMap}
\lean{SmallDC.iotaMap}
\leanok
\uses{def:SmallDC.totZ1}
Define the inclusion $\iota : H^v_{0,1} = \ker(\partial^v_0) \to Z_1$ by sending $b \mapsto (0, b)$:
\[
  \iota : \ker(\partial^v_0) \to Z_1, \quad \langle b, hb \rangle \mapsto \langle (0,b), \text{proof}\rangle.
\]
This is well-defined because $\partial_1(0,b) = \partial^h_0(0) + \partial^v_0(b) = 0 + 0 = 0$ for $b \in \ker(\partial^v_0)$.
\end{definition}

\begin{lemma}[$\iota$ Vanishes on $\operatorname{im}(\bar{\partial}^h_1)$]
\label{lem:SmallDC.iotaMap_vanishes_on_barDh1}
\lean{SmallDC.iotaMap_vanishes_on_barDh₁}
\leanok
\uses{def:SmallDC.iotaMap, def:SmallDC.barDh1, def:SmallDC.totB1InZ1, lem:SmallDC.d2_fst, lem:SmallDC.d2_snd}
$\operatorname{im}(\bar{\partial}^h_1) \leq \ker(q \circ \iota)$, where $q : Z_1 \to H_1(\operatorname{Tot}(E))$ is the quotient map.
\end{lemma}

\begin{proof}
\leanok
\uses{def:SmallDC.iotaMap, def:SmallDC.barDh1, def:SmallDC.totB1InZ1, lem:SmallDC.d2_fst, lem:SmallDC.d2_snd}
Let $v \in \operatorname{im}(\bar{\partial}^h_1)$, so $v = \bar{\partial}^h_1(\langle w, hw \rangle)$ for some $w \in A_{1,1}$ with $hw : w \in \ker(\partial^v_1)$. We must show $q(\iota(v)) = 0$ in $H_1(\operatorname{Tot}(E))$, i.e., $\iota(v) \in B_1^Z = \operatorname{im}(\partial_2) \cap Z_1$.

We have $\iota(v) = \iota(\bar{\partial}^h_1\langle w, hw \rangle) = (0, \partial^h_1(w)) \in Z_1$. It suffices to show $(0, \partial^h_1(w)) \in \operatorname{im}(\partial_2)$. We claim $\partial_2(w) = (0, \partial^h_1(w))$:
\begin{itemize}
\item First component: $(\partial_2(w))_1 = \partial^v_1(w) = 0$ by Lemma~\ref{lem:SmallDC.d2_fst} since $w \in \ker(\partial^v_1)$.
\item Second component: $(\partial_2(w))_2 = \partial^h_1(w)$ by Lemma~\ref{lem:SmallDC.d2_snd}.
\end{itemize}
Thus $(0, \partial^h_1(w)) = \partial_2(w) \in \operatorname{im}(\partial_2)$.
\end{proof}

\begin{definition}[Descended Inclusion $\bar{\iota}$]
\label{def:SmallDC.iotaBar}
\lean{SmallDC.iotaBar}
\leanok
\uses{def:SmallDC.itH1_snd, def:SmallDC.totH1, def:SmallDC.iotaMap, lem:SmallDC.iotaMap_vanishes_on_barDh1}
Since $q \circ \iota$ vanishes on $\operatorname{im}(\bar{\partial}^h_1)$ (Lemma~\ref{lem:SmallDC.iotaMap_vanishes_on_barDh1}), the map descends to
\[
  \bar{\iota} : \mathcal{H}_{1,\text{snd}} = H^v_{0,1}/\operatorname{im}(\bar{\partial}^h_1) \to H_1(\operatorname{Tot}(E)), \quad [b] \mapsto [(0,b)].
\]
\end{definition}

\begin{lemma}[Exactness: $\bar{\pi} \circ \bar{\iota} = 0$]
\label{lem:SmallDC.piBar_comp_iotaBar}
\lean{SmallDC.piBar_comp_iotaBar}
\leanok
\uses{def:SmallDC.piBar, def:SmallDC.iotaBar}
The composition $\bar{\pi} \circ \bar{\iota} = 0$.
\end{lemma}

\begin{proof}
\leanok
\uses{def:SmallDC.piBar, def:SmallDC.iotaBar}
By induction on representatives, it suffices to show that for each $\langle b, hb \rangle \in H^v_{0,1}$:
\[
  \bar{\pi}(\bar{\iota}([\langle b, hb \rangle])) = 0.
\]
\textbf{Step 1}: By definition of $\bar{\iota}$, $\bar{\iota}([\langle b, hb \rangle]) = q(\iota(\langle b, hb \rangle))$, where $q$ is the quotient map $Z_1 \to H_1(\operatorname{Tot}(E))$.

\textbf{Step 2}: By definition of $\bar{\pi}$, $\bar{\pi}(q(\iota(\langle b, hb \rangle))) = \pi'(\iota(\langle b, hb \rangle))$.

\textbf{Step 3}: We compute $\pi(\iota(\langle b, hb \rangle)) = [(\iota\langle b, hb \rangle)_1] = [(0, b)_1] = [0] = 0$ in $H^v_{1,0}$, since $\iota$ maps to $(0, b)$ and the first component of $(0,b)$ is $0 \in \operatorname{im}(\partial^v_1)$. Therefore $\pi'(\iota(\langle b, hb \rangle)) = 0$.
\end{proof}

\begin{lemma}[Helper: $(0,b) \in \operatorname{im}(\partial_2)$ Implies $b \in \operatorname{im}(\bar{\partial}^h_1)$]
\label{lem:SmallDC.mem_range_barDh1_of_inr_mem_range_d2}
\lean{SmallDC.mem_range_barDh₁_of_inr_mem_range_d₂}
\leanok
\uses{def:SmallDC.barDh1, lem:SmallDC.d2_fst, lem:SmallDC.d2_snd}
Let $b \in \ker(\partial^v_0)$. If $(0, b) \in \operatorname{im}(\partial_2)$, then $\langle b, hb \rangle \in \operatorname{im}(\bar{\partial}^h_1)$.
\end{lemma}

\begin{proof}
\leanok
\uses{def:SmallDC.barDh1, lem:SmallDC.d2_fst, lem:SmallDC.d2_snd}
Suppose $(0, b) \in \operatorname{im}(\partial_2)$; obtain $w \in A_{1,1}$ with $\partial_2(w) = (0, b)$.

By Lemma~\ref{lem:SmallDC.d2_fst}: $\partial^v_1(w) = (\partial_2(w))_1 = 0$, so $w \in \ker(\partial^v_1) = H^v_{1,1}$.

By Lemma~\ref{lem:SmallDC.d2_snd}: $\partial^h_1(w) = (\partial_2(w))_2 = b$.

Therefore $\langle w, \partial^v_1(w) = 0 \rangle \in H^v_{1,1}$ and $\bar{\partial}^h_1(\langle w, \cdot \rangle) = \langle \partial^h_1(w), \cdot \rangle = \langle b, hb \rangle$, showing $\langle b, hb \rangle \in \operatorname{im}(\bar{\partial}^h_1)$.
\end{proof}

\begin{lemma}[$\bar{\iota}$ is Injective]
\label{lem:SmallDC.iotaBar_injective}
\lean{SmallDC.iotaBar_injective}
\leanok
\uses{def:SmallDC.iotaBar, lem:SmallDC.mem_range_barDh1_of_inr_mem_range_d2}
The map $\bar{\iota} : \mathcal{H}_{1,\text{snd}} \to H_1(\operatorname{Tot}(E))$ is injective.
\end{lemma}

\begin{proof}
\leanok
\uses{def:SmallDC.iotaBar, lem:SmallDC.mem_range_barDh1_of_inr_mem_range_d2}
It suffices to show $\ker(\bar{\iota}) = \{0\}$. Let $y \in \ker(\bar{\iota})$ and lift to a representative $\langle b, hb \rangle \in H^v_{0,1}$. We must show $[\langle b, hb \rangle] = 0$ in $\mathcal{H}_{1,\text{snd}}$, i.e., $\langle b, hb \rangle \in \operatorname{im}(\bar{\partial}^h_1)$.

From $\bar{\iota}([\langle b, hb \rangle]) = 0$ in $H_1(\operatorname{Tot}(E))$, we have $q(\iota(\langle b, hb \rangle)) = 0$, so $\iota(\langle b, hb \rangle) \in B_1^Z = \operatorname{im}(\partial_2) \cap Z_1$. Unwinding, the element $(0, b) \in \operatorname{im}(\partial_2)$.

By Lemma~\ref{lem:SmallDC.mem_range_barDh1_of_inr_mem_range_d2}, $\langle b, hb \rangle \in \operatorname{im}(\bar{\partial}^h_1)$, as required.
\end{proof}

\begin{lemma}[$\bar{\pi}$ is Surjective]
\label{lem:SmallDC.piBar_surjective}
\lean{SmallDC.piBar_surjective}
\leanok
\uses{def:SmallDC.piBar, def:SmallDC.barDh0}
The map $\bar{\pi} : H_1(\operatorname{Tot}(E)) \to \mathcal{H}_{1,\text{fst}}$ is surjective.
\end{lemma}

\begin{proof}
\leanok
\uses{def:SmallDC.piBar, def:SmallDC.barDh0}
Let $\langle q, hq \rangle \in \ker(\bar{\partial}^h_0) = \mathcal{H}_{1,\text{fst}}$. Lift $q \in H^v_{1,0}$ to a representative $a \in A_{1,0}$. Since $[a] \in \ker(\bar{\partial}^h_0)$, we have $\bar{\partial}^h_0([a]) = 0$ in $H^v_{0,0}$, i.e., $[\partial^h_0(a)] = 0$, so $\partial^h_0(a) \in \operatorname{im}(\partial^v_0)$.

Obtain $b \in A_{0,1}$ with $\partial^h_0(a) = \partial^v_0(b)$. Over $\mathbb{F}_2$, $\partial^h_0(a) + \partial^v_0(b) = 2\partial^v_0(b) = 0$, so $(a,b) \in \ker(\partial_1) = Z_1$.

Consider $[(a,b)] \in H_1(\operatorname{Tot}(E))$. Then $\bar{\pi}([(a,b)]) = [a] \in H^v_{1,0}$, which as an element of $\ker(\bar{\partial}^h_0)$ equals $\langle q, hq \rangle$ since $[a] = q$. Hence $\bar{\pi}$ is surjective.
\end{proof}

\begin{lemma}[$\ker(\bar{\pi}) = \operatorname{im}(\bar{\iota})$]
\label{lem:SmallDC.piBar_ker_eq_iotaBar_range}
\lean{SmallDC.piBar_ker_eq_iotaBar_range}
\leanok
\uses{def:SmallDC.piBar, def:SmallDC.iotaBar, lem:SmallDC.piBar_comp_iotaBar, lem:SmallDC.mem_range_barDh1_of_inr_mem_range_d2}
$\ker(\bar{\pi}) = \operatorname{im}(\bar{\iota})$.
\end{lemma}

\begin{proof}
\leanok
\uses{def:SmallDC.piBar, def:SmallDC.iotaBar, lem:SmallDC.piBar_comp_iotaBar, lem:SmallDC.mem_range_barDh1_of_inr_mem_range_d2}
By extensionality on representatives in $H_1(\operatorname{Tot}(E))$.

$(\subseteq)$: Let $z \in Z_1$ with $\bar{\pi}([z]) = 0$, i.e., $[z] \in \ker(\bar{\pi})$. The condition $\bar{\pi}([z]) = 0$ means $\pi'(z) = 0$ in $\ker(\bar{\partial}^h_0)$, i.e., $z_1 \in \operatorname{im}(\partial^v_1)$. Obtain $w$ with $z_1 = \partial^v_1(w)$.

Set $b' = z_2 + \partial^h_1(w) \in A_{0,1}$. Using the commutativity condition, $\partial^v_0(b') = \partial^v_0(z_2) + \partial^h_0(\partial^v_1(w)) = 0$ (since $z \in \ker(\partial_1)$ gives $\partial^h_0(z_1) + \partial^v_0(z_2) = 0$ and $z_1 = \partial^v_1(w)$, combined with $\partial^v_0 \circ \partial^h_1 = \partial^h_0 \circ \partial^v_1$ over $\mathbb{F}_2$). Thus $\langle b', hb' \rangle \in H^v_{0,1}$.

The difference $z - \iota\langle b', hb' \rangle = (\partial^v_1(w), -\partial^h_1(w)) = \partial_2(w)$ (over $\mathbb{F}_2$, negation is identity) lies in $B_1^Z$. Hence $[z] = [\iota\langle b', hb' \rangle] = \bar{\iota}([\langle b', hb' \rangle]) \in \operatorname{im}(\bar{\iota})$.

$(\supseteq)$: This follows directly from Lemma~\ref{lem:SmallDC.piBar_comp_iotaBar}: for any $y$, $\bar{\pi}(\bar{\iota}(y)) = 0$.
\end{proof}

\begin{theorem}[Isomorphism at Degree 1]
\label{thm:SmallDC.homologyEquiv_n1}
\lean{SmallDC.homologyEquiv_n1}
\leanok
\uses{def:SmallDC.totH1, def:SmallDC.itH1_fst, def:SmallDC.itH1_snd, def:SmallDC.piBar, def:SmallDC.iotaBar, lem:SmallDC.iotaBar_injective, lem:SmallDC.piBar_surjective, lem:SmallDC.piBar_ker_eq_iotaBar_range}
There is a canonical $\mathbb{F}_2$-linear isomorphism
\[
  H_1(\operatorname{Tot}(E)) \cong \mathcal{H}_{1,\text{fst}} \times \mathcal{H}_{1,\text{snd}} = \ker(\bar{\partial}^h_0) \times \bigl(H^v_{0,1}/\operatorname{im}(\bar{\partial}^h_1)\bigr).
\]
\end{theorem}

\begin{proof}
\leanok
\uses{def:SmallDC.totH1, def:SmallDC.itH1_fst, def:SmallDC.itH1_snd, def:SmallDC.piBar, def:SmallDC.iotaBar, lem:SmallDC.iotaBar_injective, lem:SmallDC.piBar_surjective, lem:SmallDC.piBar_ker_eq_iotaBar_range}
We construct a short exact sequence that splits over $\mathbb{F}_2$. By Lemma~\ref{lem:SmallDC.piBar_comp_iotaBar}, Lemma~\ref{lem:SmallDC.iotaBar_injective}, Lemma~\ref{lem:SmallDC.piBar_surjective}, and Lemma~\ref{lem:SmallDC.piBar_ker_eq_iotaBar_range}, the sequence
\[
  0 \to \mathcal{H}_{1,\text{snd}} \xrightarrow{\bar{\iota}} H_1(\operatorname{Tot}(E)) \xrightarrow{\bar{\pi}} \mathcal{H}_{1,\text{fst}} \to 0
\]
is short exact (with $\bar{\iota}$ injective, $\bar{\pi}$ surjective, and $\ker(\bar{\pi}) = \operatorname{im}(\bar{\iota})$).

Over $\mathbb{F}_2$ (a field), every short exact sequence of vector spaces splits. Concretely, we apply the Mathlib lemma \texttt{Submodule.exists\_isCompl} to $\operatorname{im}(\bar{\iota}) \subseteq H_1(\operatorname{Tot}(E))$ to obtain a complement $K$ with $\operatorname{im}(\bar{\iota}) \oplus K = H_1(\operatorname{Tot}(E))$. The isomorphism is assembled as the composition:
\begin{align*}
  H_1(\operatorname{Tot}(E)) &\cong \operatorname{im}(\bar{\iota}) \times K && \text{(split by complement } K\text{)} \\
  &\cong \mathcal{H}_{1,\text{snd}} \times K && \text{(} \bar{\iota} \text{ injective, } \mathcal{H}_{1,\text{snd}} \cong \operatorname{im}(\bar{\iota})\text{)} \\
  &\cong \mathcal{H}_{1,\text{snd}} \times \mathcal{H}_{1,\text{fst}} && \text{(} K \cong H_1/\ker(\bar{\pi}) \cong \operatorname{im}(\bar{\pi}) = \mathcal{H}_{1,\text{fst}}\text{)} \\
  &\cong \mathcal{H}_{1,\text{fst}} \times \mathcal{H}_{1,\text{snd}}. && \text{(swap factors)}
\end{align*}
Each step uses a standard linear equivalence: the complement decomposition, the linear equivalence from injectivity of $\bar{\iota}$, the quotient-kernel equivalence for $\bar{\pi}$, and the product commutativity isomorphism.
\end{proof}

%--- Def_12: FiberBundleDoubleComplex ---
\chapter{Def 12: Fiber Bundle Double Complex}

\begin{definition}[Chain Automorphism]
\label{def:FiberBundle.ChainAutomorphism}
\lean{FiberBundle.ChainAutomorphism}
\leanok
\uses{def:F2}
Let $n_2, m_2 \in \mathbb{N}$ and let $\partial^F : \mathbb{F}_2^{n_2} \to \mathbb{F}_2^{m_2}$ be a linear map (the differential of a 1-complex $F$ over $\mathbb{F}_2$). A \emph{chain automorphism} of $F$ is a structure consisting of:
\begin{itemize}
  \item an invertible linear map $\alpha_1 : \mathbb{F}_2^{n_2} \xrightarrow{\sim} \mathbb{F}_2^{n_2}$ (automorphism on $F_1$),
  \item an invertible linear map $\alpha_0 : \mathbb{F}_2^{m_2} \xrightarrow{\sim} \mathbb{F}_2^{m_2}$ (automorphism on $F_0$),
  \item the \emph{chain map condition}: $\partial^F \circ \alpha_1 = \alpha_0 \circ \partial^F$.
\end{itemize}
\end{definition}

\begin{definition}[Identity Chain Automorphism]
\label{def:FiberBundle.ChainAutomorphism.id}
\lean{FiberBundle.ChainAutomorphism.id}
\leanok
\uses{def:FiberBundle.ChainAutomorphism}
Given a differential $\partial^F : \mathbb{F}_2^{n_2} \to \mathbb{F}_2^{m_2}$, the \emph{identity chain automorphism} $\operatorname{id}_F$ is the chain automorphism with $\alpha_1 = \operatorname{id}_{\mathbb{F}_2^{n_2}}$ and $\alpha_0 = \operatorname{id}_{\mathbb{F}_2^{m_2}}$.
\end{definition}

\begin{proof}
\leanok
\uses{def:FiberBundle.ChainAutomorphism}
We set $\alpha_1 := \operatorname{LinearEquiv.refl}_{\mathbb{F}_2}(\mathbb{F}_2^{n_2})$ and $\alpha_0 := \operatorname{LinearEquiv.refl}_{\mathbb{F}_2}(\mathbb{F}_2^{m_2})$. The chain map condition $\partial^F \circ \alpha_1 = \alpha_0 \circ \partial^F$ holds by simplification using the fact that the linear map underlying the reflexivity equivalence is the identity ($\operatorname{LinearEquiv.refl_toLinearMap}$), so both sides reduce to $\partial^F$.
\end{proof}

\begin{definition}[Connection]
\label{def:FiberBundle.Connection}
\lean{FiberBundle.Connection}
\leanok
\uses{def:FiberBundle.ChainAutomorphism}
A \emph{connection} for a fiber bundle with base dimensions $n_1, m_1$ and fiber differential $\partial^F : \mathbb{F}_2^{n_2} \to \mathbb{F}_2^{m_2}$ is a function
\[
  \varphi : \operatorname{Fin}(n_1) \to \operatorname{Fin}(m_1) \to \operatorname{ChainAutomorphism}(n_2, m_2, \partial^F).
\]
That is, to each pair $(b^1, b^0) \in \operatorname{Fin}(n_1) \times \operatorname{Fin}(m_1)$, the connection assigns a chain automorphism of $F$. Values outside the support of $\partial^B(e_{b^1})$ are irrelevant.
\end{definition}

\begin{definition}[Trivial Connection]
\label{def:FiberBundle.Connection.trivial}
\lean{FiberBundle.Connection.trivial}
\leanok
\uses{def:FiberBundle.Connection, def:FiberBundle.ChainAutomorphism.id}
The \emph{trivial connection} assigns to every pair $(b^1, b^0)$ the identity chain automorphism:
\[
  \varphi_{\operatorname{triv}}(b^1, b^0) := \operatorname{id}_F \quad \text{for all } b^1 \in \operatorname{Fin}(n_1),\ b^0 \in \operatorname{Fin}(m_1).
\]
\end{definition}

\begin{definition}[Twisted Horizontal Differential]
\label{def:FiberBundle.twistedDhLin}
\lean{FiberBundle.twistedDhLin}
\leanok
\uses{def:FiberBundle.Connection}
Let $V$ be an $\mathbb{F}_2$-module, let $\partial^B : \mathbb{F}_2^{n_1} \to \mathbb{F}_2^{m_1}$ be a linear map, and let $\operatorname{autComp} : \operatorname{Fin}(n_1) \to \operatorname{Fin}(m_1) \to (V \to_{\mathbb{F}_2} V)$ be a family of linear endomorphisms. The \emph{twisted horizontal differential}
\[
  \partial^h_\varphi : \mathbb{F}_2^{n_1} \otimes_{\mathbb{F}_2} V \;\longrightarrow\; \mathbb{F}_2^{m_1} \otimes_{\mathbb{F}_2} V
\]
is the linear map defined via the universal property of the tensor product (via $\operatorname{TensorProduct.lift}$): on a pure tensor $b \otimes f$ (with $b : \mathbb{F}_2^{n_1}$, $f : V$), it acts as
\[
  \partial^h_\varphi(b \otimes f) \;=\; \sum_{b_1 \in \operatorname{Fin}(n_1)} \sum_{b_0 \in \operatorname{Fin}(m_1)} b(b_1) \cdot \bigl(\partial^B(e_{b_1})\bigr)(b_0) \cdot \bigl(e_{b_0} \otimes_{\mathbb{F}_2} \operatorname{autComp}(b_1, b_0)(f)\bigr),
\]
where $e_{b_1}$ and $e_{b_0}$ denote the standard basis vectors. This is well-defined as a bilinear map and extends to a linear map on the tensor product.
\end{definition}

\begin{lemma}[Twisted Differential on Pure Tensors]
\label{lem:FiberBundle.twistedDhLin_tmul}
\lean{FiberBundle.twistedDhLin_tmul}
\leanok
\uses{def:FiberBundle.twistedDhLin}
For any $b : \mathbb{F}_2^{n_1}$ and $f : V$,
\[
  \partial^h_\varphi(b \otimes_{\mathbb{F}_2} f) \;=\; \sum_{b_1 \in \operatorname{Fin}(n_1)} \sum_{b_0 \in \operatorname{Fin}(m_1)} b(b_1) \cdot \bigl(\partial^B(e_{b_1})\bigr)(b_0) \cdot \bigl(e_{b_0} \otimes_{\mathbb{F}_2} \operatorname{autComp}(b_1, b_0)(f)\bigr).
\]
\end{lemma}

\begin{proof}
\leanok
\uses{def:FiberBundle.twistedDhLin}
This follows by simplification using the definition of $\operatorname{twistedDhLin}$ and the computation rules for $\operatorname{TensorProduct.lift}$ applied to a pure tensor.
\end{proof}

\begin{lemma}[Commutativity of Twisted Horizontal and Vertical Differentials]
\label{lem:FiberBundle.twistedDh_comm_dv}
\lean{FiberBundle.twistedDh_comm_dv}
\leanok
\uses{def:FiberBundle.twistedDhLin, def:FiberBundle.Connection, def:FiberBundle.ChainAutomorphism}
Let $\partial^B : \mathbb{F}_2^{n_1} \to \mathbb{F}_2^{m_1}$, $\partial^F : \mathbb{F}_2^{n_2} \to \mathbb{F}_2^{m_2}$, and $\varphi$ a connection. The twisted horizontal differential commutes with the vertical differential $\operatorname{id} \otimes \partial^F$:
\[
  \partial^h_{\varphi,\alpha_0} \circ \bigl(\operatorname{id}_{\mathbb{F}_2^{n_1}} \otimes \partial^F\bigr) \;=\; \bigl(\operatorname{id}_{\mathbb{F}_2^{m_1}} \otimes \partial^F\bigr) \circ \partial^h_{\varphi,\alpha_1},
\]
where $\partial^h_{\varphi,\alpha_q}$ denotes the twisted horizontal differential using the automorphism component $\alpha_q$ from the connection.
\end{lemma}

\begin{proof}
\leanok
\uses{def:FiberBundle.twistedDhLin, def:FiberBundle.ChainAutomorphism, lem:FiberBundle.twistedDhLin_tmul}
By tensor product extensionality ($\operatorname{TensorProduct.ext'}$), it suffices to verify equality on all pure tensors $b \otimes f$ with $b : \mathbb{F}_2^{n_1}$ and $f : \mathbb{F}_2^{n_2}$. For such a pure tensor, we simplify both sides using the definition of $\operatorname{twistedDhLin_tmul}$ and $\operatorname{LinearMap.lTensor}$. After unfolding the sums and using congruence, the equality at each summand $(b_1, b_0)$ reduces to showing:
\[
  \partial^F\bigl(\alpha_1(b_1,b_0)(f)\bigr) \;=\; \alpha_0(b_1,b_0)\bigl(\partial^F(f)\bigr).
\]
This is exactly the chain map condition $\varphi(b_1, b_0).\operatorname{comm}$: from $\partial^F \circ (\varphi(b_1,b_0).\alpha_1) = (\varphi(b_1,b_0).\alpha_0) \circ \partial^F$, we extract the value at $f$ via $\operatorname{LinearMap.ext_iff}$ and obtain the required equality by symmetry.
\end{proof}

\begin{definition}[Base Size Function]
\label{def:FiberBundle.baseSize}
\lean{FiberBundle.baseSize}
\leanok

The \emph{base size function} assigns dimensions to degrees:
\[
  \operatorname{baseSize}(n_1, m_1, p) := \begin{cases} n_1 & \text{if } p = 1, \\ m_1 & \text{if } p = 0, \\ 0 & \text{otherwise.} \end{cases}
\]
This encodes the graded structure of the base 1-complex $B = (\mathbb{F}_2^{n_1} \xrightarrow{\partial^B} \mathbb{F}_2^{m_1})$.
\end{definition}

\begin{definition}[Fiber Size Function]
\label{def:FiberBundle.fiberSize}
\lean{FiberBundle.fiberSize}
\leanok

The \emph{fiber size function} assigns dimensions to degrees:
\[
  \operatorname{fiberSize}(n_2, m_2, q) := \begin{cases} n_2 & \text{if } q = 1, \\ m_2 & \text{if } q = 0, \\ 0 & \text{otherwise.} \end{cases}
\]
This encodes the graded structure of the fiber 1-complex $F = (\mathbb{F}_2^{n_2} \xrightarrow{\partial^F} \mathbb{F}_2^{m_2})$.
\end{definition}

\begin{definition}[Graded Object of the Fiber Bundle Double Complex]
\label{def:FiberBundle.fbObj}
\lean{FiberBundle.fbObj}
\leanok
\uses{def:FiberBundle.baseSize, def:FiberBundle.fiberSize, def:F2}
The \emph{graded object} underlying the fiber bundle double complex is the functor
\[
  \operatorname{fbObj}(n_1, m_1, n_2, m_2) : \mathbb{Z} \times \mathbb{Z} \;\longrightarrow\; \operatorname{Mod}_{\mathbb{F}_2},
  \quad (p, q) \;\longmapsto\; \mathbb{F}_2^{\operatorname{baseSize}(p)} \otimes_{\mathbb{F}_2} \mathbb{F}_2^{\operatorname{fiberSize}(q)}.
\]
\end{definition}

\begin{definition}[Fiber Differential]
\label{def:FiberBundle.fiberDiff}
\lean{FiberBundle.fiberDiff}
\leanok
\uses{def:FiberBundle.fiberSize}
Given $\partial^F : \mathbb{F}_2^{n_2} \to \mathbb{F}_2^{m_2}$, the \emph{fiber differential at degrees $(q, q')$} is the linear map
\[
  \operatorname{fiberDiff}(n_2, m_2, \partial^F, q, q') : \mathbb{F}_2^{\operatorname{fiberSize}(q)} \;\longrightarrow\; \mathbb{F}_2^{\operatorname{fiberSize}(q')} \;:=\; \begin{cases} \partial^F & \text{if } q = 1 \text{ and } q' = 0, \\ 0 & \text{otherwise.} \end{cases}
\]
\end{definition}

\begin{definition}[Vertical Differential]
\label{def:FiberBundle.fbDv}
\lean{FiberBundle.fbDv}
\leanok
\uses{def:FiberBundle.fbObj, def:FiberBundle.fiberDiff}
The \emph{vertical differential} $\partial^v = \operatorname{id}_B \otimes \partial^F$ is the morphism in $\operatorname{Mod}_{\mathbb{F}_2}$:
\[
  \operatorname{fbDv}(p, q, q') : \operatorname{fbObj}(p, q) \;\longrightarrow\; \operatorname{fbObj}(p, q'),
  \quad \operatorname{fbDv}(p, q, q') := \operatorname{id}_{\mathbb{F}_2^{\operatorname{baseSize}(p)}} \otimes \operatorname{fiberDiff}(q, q'),
\]
defined via $\operatorname{LinearMap.lTensor}$.
\end{definition}

\begin{definition}[Automorphism at Fiber Degree]
\label{def:FiberBundle.autAtDeg}
\lean{FiberBundle.autAtDeg}
\leanok
\uses{def:FiberBundle.Connection, def:FiberBundle.fiberSize, def:FiberBundle.ChainAutomorphism}
Given a connection $\varphi$ and a fiber degree $q \in \mathbb{Z}$, the \emph{automorphism component at degree $q$} is the family of linear endomorphisms
\[
  \operatorname{autAtDeg}(\partial^F, \varphi, q, b_1, b_0) : \mathbb{F}_2^{\operatorname{fiberSize}(q)} \;\longrightarrow\; \mathbb{F}_2^{\operatorname{fiberSize}(q)} \;:=\; \begin{cases} \varphi(b_1,b_0).\alpha_1 & \text{if } q = 1, \\ \varphi(b_1,b_0).\alpha_0 & \text{if } q = 0, \\ \operatorname{id} & \text{otherwise.} \end{cases}
\]
\end{definition}

\begin{definition}[Twisted Horizontal Differential as Morphism]
\label{def:FiberBundle.fbDh}
\lean{FiberBundle.fbDh}
\leanok
\uses{def:FiberBundle.fbObj, def:FiberBundle.twistedDhLin, def:FiberBundle.autAtDeg, def:FiberBundle.Connection}
The \emph{twisted horizontal differential} is the morphism in $\operatorname{Mod}_{\mathbb{F}_2}$:
\[
  \operatorname{fbDh}(p, p', q) : \operatorname{fbObj}(p, q) \;\longrightarrow\; \operatorname{fbObj}(p', q) \;:=\; \begin{cases} \partial^h_\varphi & \text{if } p = 1 \text{ and } p' = 0, \\ 0 & \text{otherwise,} \end{cases}
\]
where $\partial^h_\varphi = \operatorname{twistedDhLin}(\partial^B, (b_1, b_0) \mapsto \operatorname{autAtDeg}(\partial^F, \varphi, q, b_1, b_0))$ is the connection-twisted horizontal differential.
\end{definition}

\begin{definition}[Fiber Bundle Double Complex]
\label{def:FiberBundle.fiberBundleDoubleComplex}
\lean{FiberBundle.fiberBundleDoubleComplex}
\leanok
\uses{def:DoubleComplexF2, def:FiberBundle.fbObj, def:FiberBundle.fbDh, def:FiberBundle.fbDv, def:FiberBundle.Connection}
Let $\partial^B : \mathbb{F}_2^{n_1} \to \mathbb{F}_2^{m_1}$, $\partial^F : \mathbb{F}_2^{n_2} \to \mathbb{F}_2^{m_2}$, and $\varphi$ a connection. The \emph{fiber bundle double complex} $B \boxtimes_\varphi F$ is the double complex over $\mathbb{F}_2$ with:
\begin{itemize}
  \item \textbf{Objects}: $(p, q) \mapsto \mathbb{F}_2^{\operatorname{baseSize}(p)} \otimes_{\mathbb{F}_2} \mathbb{F}_2^{\operatorname{fiberSize}(q)}$,
  \item \textbf{Horizontal differential}: $\partial^h_\varphi = \operatorname{fbDh}(p, p', q)$ (the connection-twisted map),
  \item \textbf{Vertical differential}: $\partial^v = \operatorname{fbDv}(p, q, q') = \operatorname{id}_B \otimes \partial^F$.
\end{itemize}
This is constructed using $\operatorname{HomologicalComplex}_2.\operatorname{ofGradedObject}$ with the descending complex shape on $\mathbb{Z}$. The double complex axioms are verified by the following conditions:
\begin{enumerate}
  \item $(\partial^h)^2 = 0$: since $\operatorname{fbDh}(p, p', q) \circ \operatorname{fbDh}(p', p'', q) = 0$ (at most one factor can be nonzero, as $p = 1, p' = 0$ and $p' = 1, p'' = 0$ cannot both hold),
  \item $(\partial^v)^2 = 0$: since $\operatorname{fbDv}(p, q', q'') \circ \operatorname{fbDv}(p, q, q') = 0$ (similarly),
  \item Commutativity $\partial^h \circ \partial^v = \partial^v \circ \partial^h$: by Lemma \ref{lem:FiberBundle.twistedDh_comm_dv}, since each $\varphi(b_1, b_0)$ is a chain automorphism.
\end{enumerate}
\end{definition}

\begin{definition}[Fiber Bundle Complex]
\label{def:FiberBundle.fiberBundleComplex}
\lean{FiberBundle.fiberBundleComplex}
\leanok
\uses{def:ChainComplexF2, def:FiberBundle.fiberBundleDoubleComplex, def:DoubleComplexF2.totalComplex}
The \emph{fiber bundle complex} $B \otimes_\varphi F = \operatorname{Tot}(B \boxtimes_\varphi F)$ is the total complex of the fiber bundle double complex:
\[
  B \otimes_\varphi F \;:=\; (B \boxtimes_\varphi F).\operatorname{totalComplex} \;\in\; \operatorname{ChainComplex}_{\mathbb{F}_2}.
\]
It is the chain complex over $\mathbb{F}_2$ obtained by taking the direct sum totalization of the double complex $B \boxtimes_\varphi F$, with differential given by $\partial^h_\varphi + \partial^v$ (summing the horizontal and vertical differentials in each total degree). The existence of the total complex is guaranteed because each graded piece is a finitely-generated $\mathbb{F}_2$-module, so all necessary coproducts exist and the $\operatorname{HasTotal}$ instance holds by type class inference.
\end{definition}

%--- Thm_3: FiberBundleHomology ---
\chapter{Thm 3: Fiber Bundle Homology}

\begin{definition}[Chain Automorphism Acts as Identity on Homology]
\label{def:FiberBundle.ChainAutomorphism.ActsAsIdOnHomology}
\lean{FiberBundle.ChainAutomorphism.ActsAsIdOnHomology}
\leanok
\uses{def:FiberBundle.ChainAutomorphism}
Let $d_F : (\mathbb{F}_2^{n_2}) \to (\mathbb{F}_2^{m_2})$ be a linear map and let $\psi$ be a chain automorphism of the chain complex $F = (F^1 \xrightarrow{d_F} F^0)$.
We say $\psi$ \emph{acts as the identity on homology} if:
\begin{enumerate}
  \item \textbf{Degree 1:} For every $f \in \ker(d_F)$, we have $\alpha_1(f) = f$ (i.e., $\alpha_1$ fixes cycles).
  \item \textbf{Degree 0:} For every $y \in \mathbb{F}_2^{m_2}$, we have $\alpha_0(y) - y \in \operatorname{im}(d_F)$ (i.e., $[\alpha_0(y)] = [y]$ in $H_0(F) = \mathbb{F}_2^{m_2}/\operatorname{im}(d_F)$).
\end{enumerate}
\end{definition}

\begin{definition}[Connection Acts as Identity on Homology]
\label{def:FiberBundle.ActsAsIdentityOnHomology}
\lean{FiberBundle.ActsAsIdentityOnHomology}
\leanok
\uses{def:FiberBundle.Connection, def:FiberBundle.ChainAutomorphism.ActsAsIdOnHomology}
Let $d_B : \mathbb{F}_2^{n_1} \to \mathbb{F}_2^{m_1}$ and $d_F : \mathbb{F}_2^{n_2} \to \mathbb{F}_2^{m_2}$ be linear maps, and let $\varphi$ be a connection of type $\operatorname{Connection}(n_1, m_1, n_2, m_2, d_F)$.

We say $\varphi$ \emph{acts as identity on homology} if for every basis element $b^1 \in \operatorname{Fin}(n_1)$ and every $b^0 \in \operatorname{Fin}(m_1)$ such that $d_B(e_{b^1})(b^0) \neq 0$, the chain automorphism $\varphi(b^1, b^0)$ acts as the identity on the homology of $F$.

Formally:
\[
  \forall\, b^1 \in \operatorname{Fin}(n_1),\; b^0 \in \operatorname{Fin}(m_1),\quad d_B(e_{b^1})(b^0) \neq 0 \;\Rightarrow\; \varphi(b^1, b^0).\operatorname{ActsAsIdOnHomology}(d_F).
\]
\end{definition}

\begin{definition}[Fiber Bundle as Small Double Complex]
\label{def:FiberBundle.fiberBundleSmallDC}
\lean{FiberBundle.fiberBundleSmallDC}
\leanok
\uses{def:FiberBundle.Connection, def:SmallDC, def:FiberBundle.twistedDhLin, def:FiberBundle.autAtDeg, lem:FiberBundle.twistedDh_comm_dv}
Given linear maps $d_B : \mathbb{F}_2^{n_1} \to \mathbb{F}_2^{m_1}$ and $d_F : \mathbb{F}_2^{n_2} \to \mathbb{F}_2^{m_2}$, and a connection $\varphi$, the \emph{fiber bundle small double complex} is the $2\times 2$ double complex $E = B \otimes_\varphi F$ with:
\begin{align*}
  A_{1,1} &= \mathbb{F}_2^{n_1} \otimes_{\mathbb{F}_2} \mathbb{F}_2^{n_2}, & A_{1,0} &= \mathbb{F}_2^{n_1} \otimes_{\mathbb{F}_2} \mathbb{F}_2^{m_2},\\
  A_{0,1} &= \mathbb{F}_2^{m_1} \otimes_{\mathbb{F}_2} \mathbb{F}_2^{n_2}, & A_{0,0} &= \mathbb{F}_2^{m_1} \otimes_{\mathbb{F}_2} \mathbb{F}_2^{m_2},
\end{align*}
with vertical differentials $d_v^i = \operatorname{id} \otimes d_F$ (left-tensored), and horizontal twisted differentials:
\[
  d_h^1 = \widetilde{d_h}(d_B,\, \alpha_1^\varphi), \qquad d_h^0 = \widetilde{d_h}(d_B,\, \alpha_0^\varphi),
\]
where $\alpha_k^\varphi(b^1, b^0) = \operatorname{autAtDeg}(d_F, \varphi, k, b^1, b^0)$. The anticommutativity condition $d_h \circ d_v = d_v \circ d_h$ is supplied by $\operatorname{twistedDh_comm_dv}$.
\end{definition}

\begin{definition}[Kernel Flatness Equivalence for Left Tensor]
\label{def:FiberBundle.lTensor_ker_equiv}
\lean{FiberBundle.lTensor_ker_equiv}
\leanok
\uses{def:FiberBundle.Connection}
Let $M$ be a free $\mathbb{F}_2$-module and $f : \mathbb{F}_2^{n_2} \to \mathbb{F}_2^{m_2}$ a linear map. Since $M$ is free over a field it is flat, so there is a canonical linear equivalence:
\[
  \ker(\operatorname{id}_M \otimes f) \;\simeq\; M \otimes_{\mathbb{F}_2} \ker(f).
\]
\end{definition}

\begin{proof}
\leanok
Since $M$ is free over $\mathbb{F}_2$, it is flat. Flatness implies that $\operatorname{id}_M \otimes \iota_{\ker f}$ is injective, where $\iota_{\ker f} : \ker(f) \hookrightarrow \mathbb{F}_2^{n_2}$ is the inclusion. Moreover, by flatness applied to the exact sequence $\ker(f) \hookrightarrow \mathbb{F}_2^{n_2} \xrightarrow{f} \mathbb{F}_2^{m_2}$, the sequence
\[
  M \otimes \ker(f) \xrightarrow{\operatorname{id}_M \otimes \iota} M \otimes \mathbb{F}_2^{n_2} \xrightarrow{\operatorname{id}_M \otimes f} M \otimes \mathbb{F}_2^{m_2}
\]
is exact, so $\ker(\operatorname{id}_M \otimes f) = \operatorname{im}(\operatorname{id}_M \otimes \iota)$. Combining the equality of submodules with the injectivity of $\operatorname{id}_M \otimes \iota$, we obtain the desired linear equivalence by composing $\operatorname{LinearEquiv.ofEq}$ with $(\operatorname{LinearEquiv.ofInjective})^{-1}$.
\end{proof}

\begin{definition}[Quotient Flatness Equivalence for Left Tensor]
\label{def:FiberBundle.lTensor_quotient_equiv}
\lean{FiberBundle.lTensor_quotient_equiv}
\leanok
\uses{def:FiberBundle.lTensor_ker_equiv}
Let $M$ be an $\mathbb{F}_2$-module and $f : \mathbb{F}_2^{n_2} \to \mathbb{F}_2^{m_2}$. There is a canonical linear equivalence:
\[
  \bigl(M \otimes_{\mathbb{F}_2} \mathbb{F}_2^{m_2}\bigr) \big/ \operatorname{im}(\operatorname{id}_M \otimes f) \;\simeq\; M \otimes_{\mathbb{F}_2} \bigl(\mathbb{F}_2^{m_2} / \operatorname{im}(f)\bigr).
\]
\end{definition}

\begin{proof}
\leanok
This follows from $\operatorname{lTensor.equiv}$ applied to the exact sequence $\mathbb{F}_2^{n_2} \xrightarrow{f} \mathbb{F}_2^{m_2} \xrightarrow{\pi} \mathbb{F}_2^{m_2}/\operatorname{im}(f) \to 0$ (where $\pi$ is surjective), giving the canonical identification of the left-tensored quotient with the tensor product of the quotient.
\end{proof}

\begin{definition}[Kernel Flatness Equivalence for Right Tensor]
\label{def:FiberBundle.rTensor_ker_equiv}
\lean{FiberBundle.rTensor_ker_equiv}
\leanok
\uses{def:FiberBundle.lTensor_ker_equiv}
Let $M$ be a free $\mathbb{F}_2$-module and $f : \mathbb{F}_2^{n_1} \to \mathbb{F}_2^{m_1}$. There is a canonical linear equivalence:
\[
  \ker(f \otimes \operatorname{id}_M) \;\simeq\; \ker(f) \otimes_{\mathbb{F}_2} M.
\]
\end{definition}

\begin{proof}
\leanok
Analogously to the left-tensor case: flatness of $M$ gives injectivity of $\iota_{\ker f} \otimes \operatorname{id}_M$ and the exact sequence $\ker(f) \otimes M \to \mathbb{F}_2^{n_1} \otimes M \xrightarrow{f \otimes \operatorname{id}} \mathbb{F}_2^{m_1} \otimes M$. Combining via $\operatorname{LinearEquiv.ofEq}$ and $(\operatorname{LinearEquiv.ofInjective})^{-1}$ yields the result.
\end{proof}

\begin{definition}[Quotient Flatness Equivalence for Right Tensor]
\label{def:FiberBundle.rTensor_quotient_equiv}
\lean{FiberBundle.rTensor_quotient_equiv}
\leanok
\uses{def:FiberBundle.lTensor_quotient_equiv}
Let $M$ be an $\mathbb{F}_2$-module and $f : \mathbb{F}_2^{n_1} \to \mathbb{F}_2^{m_1}$. There is a canonical linear equivalence:
\[
  \bigl(\mathbb{F}_2^{m_1} \otimes_{\mathbb{F}_2} M\bigr) \big/ \operatorname{im}(f \otimes \operatorname{id}_M) \;\simeq\; \bigl(\mathbb{F}_2^{m_1}/\operatorname{im}(f)\bigr) \otimes_{\mathbb{F}_2} M.
\]
\end{definition}

\begin{proof}
\leanok
This follows from $\operatorname{rTensor.equiv}$ applied to the exact sequence $\mathbb{F}_2^{n_1} \xrightarrow{f} \mathbb{F}_2^{m_1} \xrightarrow{\pi} \mathbb{F}_2^{m_1}/\operatorname{im}(f) \to 0$.
\end{proof}

\begin{lemma}[Tensor Product Basis Decomposition]
\label{lem:FiberBundle.pi_tmul_eq_sum}
\lean{FiberBundle.pi_tmul_eq_sum}
\leanok
\uses{def:FiberBundle.twistedDhLin}
Let $V$ be an $\mathbb{F}_2$-module, $n \in \mathbb{N}$, $g : \operatorname{Fin}(n) \to \mathbb{F}_2$, and $v \in V$. Then:
\[
  g \otimes_{\mathbb{F}_2} v = \sum_{i \in \operatorname{Fin}(n)} g(i) \cdot (e_i \otimes_{\mathbb{F}_2} v),
\]
where $e_i = \operatorname{Pi.single}\,i\,1$.
\end{lemma}

\begin{proof}
\leanok
Writing $g = \sum_i g(i) \cdot e_i$ (by evaluating both sides at each $j$ using $\operatorname{Pi.single}$), we substitute into the left-hand side and use $\operatorname{TensorProduct.sum_tmul}$ to distribute the sum, obtaining the stated equality.
\end{proof}

\begin{lemma}[Twisted $d_h^1$ on Basis Elements in Kernel]
\label{lem:FiberBundle.twistedDh1_single_tmul}
\lean{FiberBundle.twistedDh1_single_tmul}
\leanok
\uses{def:FiberBundle.Connection, def:FiberBundle.twistedDhLin, def:FiberBundle.ActsAsIdentityOnHomology, lem:FiberBundle.twistedDhLin_tmul}
Let $\varphi$ act as identity on homology. For any basis element $e_i \in \mathbb{F}_2^{n_1}$ and any $f \in \ker(d_F)$:
\[
  \widetilde{d_h}(d_B, \alpha_1^\varphi)(e_i \otimes f) = d_B(e_i) \otimes f.
\]
That is, the twisted horizontal differential $d_h^1$ acts on $B_1 \otimes \ker(d_F)$ as the untwisted $d_B \otimes \operatorname{id}$.
\end{lemma}

\begin{proof}
\leanok
\uses{lem:FiberBundle.twistedDhLin_tmul, def:FiberBundle.ActsAsIdentityOnHomology}
Expanding by $\operatorname{twistedDhLin_tmul}$, the left-hand side becomes a double sum over $b^1 \in \operatorname{Fin}(n_1)$ and $b^0 \in \operatorname{Fin}(m_1)$. Since $e_i = \operatorname{Pi.single}\,i\,1$ vanishes at all $b^1 \neq i$, only the $b^1 = i$ term survives. For that term, the inner sum is $\sum_{b^0} d_B(e_i)(b^0) \cdot (e_{b^0} \otimes \alpha_1^\varphi(i, b^0)(f))$. For each $b^0$ with $d_B(e_i)(b^0) \neq 0$, the identity-on-homology hypothesis gives $\alpha_1^\varphi(i, b^0)(f) = f$ (since $f \in \ker(d_F)$). For $b^0$ with $d_B(e_i)(b^0) = 0$, the term vanishes. Hence the sum equals $\sum_{b^0} d_B(e_i)(b^0) \cdot (e_{b^0} \otimes f) = d_B(e_i) \otimes f$ by Lemma~\ref{lem:FiberBundle.pi_tmul_eq_sum} (applied in reverse).
\end{proof}

\begin{lemma}[Twisted $d_h^1$ Agrees with $d_B \otimes \operatorname{id}$ on $B_1 \otimes \ker(d_F)$]
\label{lem:FiberBundle.twistedDh1_comp_lTensor_eq}
\lean{FiberBundle.twistedDh1_comp_lTensor_eq}
\leanok
\uses{def:FiberBundle.twistedDhLin, def:FiberBundle.ActsAsIdentityOnHomology, lem:FiberBundle.twistedDh1_single_tmul}
Assume $\varphi$ acts as identity on homology. Then:
\[
  \widetilde{d_h}(d_B, \alpha_1^\varphi) \circ (\operatorname{id} \otimes \iota_{\ker d_F}) = (\operatorname{id} \otimes \iota_{\ker d_F}) \circ (d_B \otimes \operatorname{id}_{\ker d_F}),
\]
i.e., the following diagram commutes:
\[
\begin{tikzcd}
  \mathbb{F}_2^{n_1} \otimes \ker(d_F) \arrow[r, "\operatorname{id}\otimes\iota"] \arrow[d, "d_B\otimes\operatorname{id}"'] &
  \mathbb{F}_2^{n_1} \otimes \mathbb{F}_2^{n_2} \arrow[d, "\widetilde{d_h^1}"] \\
  \mathbb{F}_2^{m_1} \otimes \ker(d_F) \arrow[r, "\operatorname{id}\otimes\iota"'] &
  \mathbb{F}_2^{m_1} \otimes \mathbb{F}_2^{n_2}
\end{tikzcd}
\]
\end{lemma}

\begin{proof}
\leanok
\uses{lem:FiberBundle.twistedDh1_single_tmul, lem:FiberBundle.pi_tmul_eq_sum}
By $\operatorname{TensorProduct.ext'}$, it suffices to check on pure tensors $b \otimes \langle f, h_f \rangle$ for arbitrary $b \in \mathbb{F}_2^{n_1}$ and $f \in \ker(d_F)$. Writing $b = \sum_i b(i) \cdot e_i$, both sides of the claimed equality reduce (using $\operatorname{TensorProduct.sum_tmul}$ and $\operatorname{map_sum}$) to $\sum_i b(i) \cdot (\text{basis term at }i)$. For each basis term, Lemma~\ref{lem:FiberBundle.twistedDh1_single_tmul} gives $\widetilde{d_h^1}(e_i \otimes f) = d_B(e_i) \otimes f$, and pulling out the scalar $b(i)$ via linearity completes the proof.
\end{proof}

\begin{definition}[Kernel Equivalence from Conjugate Diagram]
\label{def:FiberBundle.kerEquivOfConjugate}
\lean{FiberBundle.kerEquivOfConjugate}
\leanok
\uses{def:FiberBundle.lTensor_ker_equiv}
Suppose $e_1 : M_1 \simeq N_1$ and $e_2 : M_2 \simeq N_2$ are linear equivalences, $f : M_1 \to M_2$ and $g : N_1 \to N_2$ are linear maps, and the diagram
\[
\begin{tikzcd}
  M_1 \arrow[r,"f"] \arrow[d,"e_1"'] & M_2 \arrow[d,"e_2"] \\
  N_1 \arrow[r,"g"'] & N_2
\end{tikzcd}
\]
commutes (i.e., $e_2 \circ f = g \circ e_1$). Then there is a canonical linear equivalence $\ker(f) \simeq \ker(g)$.
\end{definition}

\begin{proof}
\leanok
Define the forward map $\langle x, h_x\rangle \mapsto \langle e_1(x), \cdot\rangle$ where membership in $\ker(g)$ follows from $g(e_1(x)) = e_2(f(x)) = e_2(0) = 0$. The inverse sends $\langle y, h_y\rangle \mapsto \langle e_1^{-1}(y), \cdot\rangle$, with membership verified by applying $e_2$ (which is injective) and using the commutative diagram together with $h_y$. Left and right inverses hold by $\operatorname{simp}$ (using $e_1$'s equivalence properties). Linearity of the map follows from linearity of $e_1$.
\end{proof}

\begin{definition}[Cokernel Equivalence from Conjugate Diagram]
\label{def:FiberBundle.quotientEquivOfConjugate}
\lean{FiberBundle.quotientEquivOfConjugate}
\leanok
\uses{def:FiberBundle.kerEquivOfConjugate}
Under the same hypotheses as Definition~\ref{def:FiberBundle.kerEquivOfConjugate}, there is a canonical linear equivalence $\operatorname{coker}(f) \simeq \operatorname{coker}(g)$, i.e., $(M_2 / \operatorname{im}(f)) \simeq (N_2 / \operatorname{im}(g))$.
\end{definition}

\begin{proof}
\leanok
First, we show $e_2(\operatorname{im}(f)) = \operatorname{im}(g)$. For the forward inclusion: given $e_2(f(x)) \in e_2(\operatorname{im}(f))$, the commutativity $e_2(f(x)) = g(e_1(x))$ places this in $\operatorname{im}(g)$. For the reverse: given $g(w) \in \operatorname{im}(g)$, we have $g(w) = e_2(f(e_1^{-1}(w)))$ by commutativity, so it lies in $e_2(\operatorname{im}(f))$. With $e_2(\operatorname{im}(f)) = \operatorname{im}(g)$ established, the result follows from $\operatorname{Submodule.Quotient.equiv}$.
\end{proof}

\begin{lemma}[Fiber Bundle $\bar{d}_h^1$ Conjugation]
\label{lem:FiberBundle.fiberBundleSmallDC_barDh1_conj}
\lean{FiberBundle.fiberBundleSmallDC_barDh₁_conj}
\leanok
\uses{def:FiberBundle.fiberBundleSmallDC, def:FiberBundle.lTensor_ker_equiv, lem:FiberBundle.twistedDh1_comp_lTensor_eq}
Let $E = \operatorname{fiberBundleSmallDC}(d_B, d_F, \varphi)$ and assume $\varphi$ acts as identity on homology. Then:
\[
  e_2 \circ \bar{d}_h^1 = (d_B \otimes_r \operatorname{id}_{\ker d_F}) \circ e_1,
\]
where $e_k : \mathbb{F}_2^{n_k} \otimes \ker(d_F) \xrightarrow{\sim} \ker(\operatorname{id}_{\mathbb{F}_2^{n_k}} \otimes d_F)$ are the flatness equivalences (Definition~\ref{def:FiberBundle.lTensor_ker_equiv}) for $k=1,2$.

That is, $\bar{d}_h^1$ (the restriction of $d_h^1$ to the kernel of $d_v^1$) is conjugate to $d_B \otimes \operatorname{id}_{\ker(d_F)}$ via these equivalences.
\end{lemma}

\begin{proof}
\leanok
\uses{lem:FiberBundle.twistedDh1_comp_lTensor_eq, def:FiberBundle.lTensor_ker_equiv}
Since $\operatorname{id} \otimes \iota_{\ker d_F}$ is injective (by flatness of $\mathbb{F}_2^{m_1}$), it suffices to show equality after composing with this injection. For any $\langle x, h_x\rangle \in \ker(\operatorname{id} \otimes d_F)$, the left-hand side gives $(\operatorname{id}\otimes\iota)(e_2(\bar{d}_h^1\langle x, h_x\rangle)) = (\bar{d}_h^1\langle x,h_x\rangle).\operatorname{val} = d_h^1(x)$ (using $\operatorname{lTensor_ker_equiv_apply_val}$). On the right, applying Lemma~\ref{lem:FiberBundle.twistedDh1_comp_lTensor_eq} at $e_1\langle x, h_x\rangle$ and using $\operatorname{lTensor_ker_equiv_apply_val}$ again yields the same expression $d_h^1(x)$.
\end{proof}

\begin{lemma}[Twisted $d_h^0$ on Quotient: Basis Elements]
\label{lem:FiberBundle.twistedDh0_single_tmul_quotient}
\lean{FiberBundle.twistedDh0_single_tmul_quotient}
\leanok
\uses{def:FiberBundle.twistedDhLin, def:FiberBundle.ActsAsIdentityOnHomology, lem:FiberBundle.twistedDhLin_tmul}
Assume $\varphi$ acts as identity on homology. For any $e_i \in \mathbb{F}_2^{n_1}$ and $f \in \mathbb{F}_2^{m_2}$:
\[
  (\operatorname{id}_{\mathbb{F}_2^{m_1}} \otimes \pi) \bigl(\widetilde{d_h}(d_B, \alpha_0^\varphi)(e_i \otimes f)\bigr) = d_B(e_i) \otimes [\,f\,],
\]
where $\pi : \mathbb{F}_2^{m_2} \to \mathbb{F}_2^{m_2}/\operatorname{im}(d_F)$ is the quotient map.
\end{lemma}

\begin{proof}
\leanok
\uses{lem:FiberBundle.twistedDhLin_tmul, def:FiberBundle.ActsAsIdentityOnHomology}
Expanding by $\operatorname{twistedDhLin_tmul}$ and distributing the quotient map $\pi$ through the sum and scalar multiples, only the $b^1 = i$ summand survives (since $e_i(b^1) = 0$ for $b^1 \neq i$). In the surviving term, for each $b^0$ with $d_B(e_i)(b^0) \neq 0$, identity-on-homology at degree 0 gives $\pi(\alpha_0^\varphi(i, b^0)(f)) = \pi(f)$ (since $\alpha_0(f) - f \in \operatorname{im}(d_F) = \ker\pi$). For $b^0$ with $d_B(e_i)(b^0) = 0$ the term vanishes. Hence the left side equals $\sum_{b^0} d_B(e_i)(b^0) \cdot (e_{b^0} \otimes [f]) = d_B(e_i) \otimes [f]$.
\end{proof}

\begin{lemma}[Twisted $d_h^0$ Commutes with Left Tensor Quotient]
\label{lem:FiberBundle.twistedDh0_comp_lTensor_quotient_eq}
\lean{FiberBundle.twistedDh0_comp_lTensor_quotient_eq}
\leanok
\uses{def:FiberBundle.twistedDhLin, def:FiberBundle.ActsAsIdentityOnHomology, lem:FiberBundle.twistedDh0_single_tmul_quotient}
Assume $\varphi$ acts as identity on homology. Then:
\[
  (\operatorname{id} \otimes \pi) \circ \widetilde{d_h}(d_B, \alpha_0^\varphi) = (d_B \otimes_r \operatorname{id}) \circ (\operatorname{id} \otimes \pi),
\]
i.e., the following diagram commutes:
\[
\begin{tikzcd}
  \mathbb{F}_2^{n_1} \otimes \mathbb{F}_2^{m_2} \arrow[r,"\widetilde{d_h^0}"] \arrow[d,"\operatorname{id}\otimes\pi"'] &
  \mathbb{F}_2^{m_1} \otimes \mathbb{F}_2^{m_2} \arrow[d,"\operatorname{id}\otimes\pi"] \\
  \mathbb{F}_2^{n_1} \otimes \operatorname{coker}(d_F) \arrow[r,"d_B\otimes\operatorname{id}"'] &
  \mathbb{F}_2^{m_1} \otimes \operatorname{coker}(d_F)
\end{tikzcd}
\]
\end{lemma}

\begin{proof}
\leanok
\uses{lem:FiberBundle.twistedDh0_single_tmul_quotient}
By $\operatorname{TensorProduct.ext'}$ it suffices to check on pure tensors $b \otimes f$. Writing $b = \sum_i b(i) \cdot e_i$, both sides decompose as sums over basis elements. For each $i$, pulling out the scalar $b(i)$ via linearity and applying Lemma~\ref{lem:FiberBundle.twistedDh0_single_tmul_quotient} (using $\operatorname{map_smul}$ and $\operatorname{smul_tmul'}$) gives the equality term by term.
\end{proof}

\begin{lemma}[Fiber Bundle $\bar{d}_h^0$ Quotient Conjugation (Unproven)]
\label{lem:FiberBundle.barDh0_lTensor_quotient_conj}
\lean{FiberBundle.barDh₀_lTensor_quotient_conj}
\leanok
\uses{def:FiberBundle.fiberBundleSmallDC, def:FiberBundle.lTensor_quotient_equiv, lem:FiberBundle.twistedDh0_comp_lTensor_quotient_eq}

\textbf{\color{red}[UNPROVEN -- sorry in Lean source]}

Let $E = \operatorname{fiberBundleSmallDC}(d_B, d_F, \varphi)$ and assume $\varphi$ acts as identity on homology. Then:
\[
  q_2 \circ \bar{d}_h^0 = (d_B \otimes_r \operatorname{id}_{\operatorname{coker}(d_F)}) \circ q_1,
\]
where $q_k : (\mathbb{F}_2^{n_k} \otimes \mathbb{F}_2^{m_2}) / \operatorname{im}(\operatorname{id} \otimes d_F) \xrightarrow{\sim} \mathbb{F}_2^{n_k} \otimes \operatorname{coker}(d_F)$ are the flatness quotient equivalences (Definition~\ref{def:FiberBundle.lTensor_quotient_equiv}).

\textit{Note: This lemma is left with \texttt{sorry} in the formalization. The analogous result for $\bar{d}_h^1$ is fully proved (Lemma~\ref{lem:FiberBundle.fiberBundleSmallDC_barDh1_conj}); the degree-0 case requires additional diagram chasing on the quotient that remains to be completed.}
\end{lemma}

\begin{theorem}[Fiber Bundle Homology at Degree 2]
\label{thm:FiberBundle.fiberBundleHomology_n2}
\lean{FiberBundle.fiberBundleHomology_n2}
\leanok
\uses{def:FiberBundle.Connection, def:FiberBundle.autAtDeg, def:FiberBundle.twistedDhLin, def:SmallDC, lem:FiberBundle.twistedDhLin_tmul, lem:FiberBundle.twistedDh_comm_dv, def:FiberBundle.fiberBundleSmallDC, def:FiberBundle.ActsAsIdentityOnHomology, def:FiberBundle.lTensor_ker_equiv, def:FiberBundle.rTensor_ker_equiv, def:FiberBundle.kerEquivOfConjugate, lem:FiberBundle.fiberBundleSmallDC_barDh1_conj}
Let $d_B : \mathbb{F}_2^{n_1} \to \mathbb{F}_2^{m_1}$, $d_F : \mathbb{F}_2^{n_2} \to \mathbb{F}_2^{m_2}$, and let $\varphi$ be a connection that acts as identity on homology. Then:
\[
  H_2(B \otimes_\varphi F) \;\simeq\; \ker(d_B) \otimes_{\mathbb{F}_2} \ker(d_F).
\]
\end{theorem}

\begin{proof}
\leanok
\uses{def:FiberBundle.fiberBundleSmallDC, def:FiberBundle.ActsAsIdentityOnHomology, def:FiberBundle.lTensor_ker_equiv, def:FiberBundle.rTensor_ker_equiv, def:FiberBundle.kerEquivOfConjugate, lem:FiberBundle.fiberBundleSmallDC_barDh1_conj}
Let $E = \operatorname{fiberBundleSmallDC}(d_B, d_F, \varphi)$.

\textbf{Step 1.} By the SmallDC homology equivalence at degree 2, $\operatorname{totH}_2(E) \simeq \operatorname{itH}_2(E) = \ker(\bar{d}_h^1)$.

\textbf{Step 2.} By Lemma~\ref{lem:FiberBundle.fiberBundleSmallDC_barDh1_conj}, $\bar{d}_h^1$ is conjugate (via the flatness equivalences $e_1, e_2$ of Definition~\ref{def:FiberBundle.lTensor_ker_equiv}) to $d_B \otimes_r \operatorname{id}_{\ker(d_F)}$. Applying Definition~\ref{def:FiberBundle.kerEquivOfConjugate} to this commutative diagram yields $\ker(\bar{d}_h^1) \simeq \ker(d_B \otimes_r \operatorname{id}_{\ker(d_F)})$.

\textbf{Step 3.} By Definition~\ref{def:FiberBundle.rTensor_ker_equiv} (flatness of $\ker(d_F)$ over $\mathbb{F}_2$), $\ker(d_B \otimes_r \operatorname{id}_{\ker(d_F)}) \simeq \ker(d_B) \otimes_{\mathbb{F}_2} \ker(d_F)$.

Composing the three equivalences gives $H_2(B\otimes_\varphi F) \simeq \ker(d_B) \otimes \ker(d_F)$.
\end{proof}

\begin{theorem}[Fiber Bundle Homology at Degree 0]
\label{thm:FiberBundle.fiberBundleHomology_n0}
\lean{FiberBundle.fiberBundleHomology_n0}
\leanok
\uses{def:FiberBundle.Connection, def:FiberBundle.autAtDeg, def:FiberBundle.twistedDhLin, def:SmallDC, lem:FiberBundle.twistedDhLin_tmul, lem:FiberBundle.twistedDh_comm_dv, def:FiberBundle.fiberBundleSmallDC, def:FiberBundle.ActsAsIdentityOnHomology, def:FiberBundle.lTensor_quotient_equiv, def:FiberBundle.rTensor_quotient_equiv, def:FiberBundle.quotientEquivOfConjugate, lem:FiberBundle.barDh0_lTensor_quotient_conj}
Let $d_B$, $d_F$, $\varphi$ be as above with $\varphi$ acting as identity on homology. Then:
\[
  H_0(B \otimes_\varphi F) \;\simeq\; \operatorname{coker}(d_B) \otimes_{\mathbb{F}_2} \operatorname{coker}(d_F),
\]
i.e., $(\mathbb{F}_2^{m_1}/\operatorname{im}(d_B)) \otimes_{\mathbb{F}_2} (\mathbb{F}_2^{m_2}/\operatorname{im}(d_F))$.
\end{theorem}

\begin{proof}
\leanok
\uses{def:FiberBundle.fiberBundleSmallDC, def:FiberBundle.ActsAsIdentityOnHomology, def:FiberBundle.lTensor_quotient_equiv, def:FiberBundle.rTensor_quotient_equiv, def:FiberBundle.quotientEquivOfConjugate, lem:FiberBundle.barDh0_lTensor_quotient_conj}
Let $E = \operatorname{fiberBundleSmallDC}(d_B, d_F, \varphi)$.

\textbf{Step 1.} By the SmallDC homology equivalence at degree 0, $\operatorname{totH}_0(E) \simeq \operatorname{itH}_0(E) = A_{0,0}/\operatorname{im}(\bar{d}_h^0)$.

\textbf{Step 2.} By Lemma~\ref{lem:FiberBundle.barDh0_lTensor_quotient_conj}, $\bar{d}_h^0$ is conjugate (via the flatness quotient equivalences $q_1, q_2$ of Definition~\ref{def:FiberBundle.lTensor_quotient_equiv}) to $d_B \otimes_r \operatorname{id}_{\operatorname{coker}(d_F)}$. Applying Definition~\ref{def:FiberBundle.quotientEquivOfConjugate} gives $\operatorname{coker}(\bar{d}_h^0) \simeq \operatorname{coker}(d_B \otimes_r \operatorname{id}_{\operatorname{coker}(d_F)})$.

\textbf{Step 3.} By Definition~\ref{def:FiberBundle.rTensor_quotient_equiv}, $\operatorname{coker}(d_B \otimes_r \operatorname{id}_{\operatorname{coker}(d_F)}) \simeq \operatorname{coker}(d_B) \otimes_{\mathbb{F}_2} \operatorname{coker}(d_F)$.

Composing yields $H_0(B \otimes_\varphi F) \simeq \operatorname{coker}(d_B) \otimes \operatorname{coker}(d_F)$.
\end{proof}

\begin{theorem}[Fiber Bundle Homology at Degree 1]
\label{thm:FiberBundle.fiberBundleHomology_n1}
\lean{FiberBundle.fiberBundleHomology_n1}
\leanok
\uses{def:FiberBundle.Connection, def:FiberBundle.autAtDeg, def:FiberBundle.twistedDhLin, def:SmallDC, lem:FiberBundle.twistedDhLin_tmul, lem:FiberBundle.twistedDh_comm_dv, def:FiberBundle.fiberBundleSmallDC, def:FiberBundle.ActsAsIdentityOnHomology, def:FiberBundle.lTensor_ker_equiv, def:FiberBundle.lTensor_quotient_equiv, def:FiberBundle.rTensor_ker_equiv, def:FiberBundle.rTensor_quotient_equiv, def:FiberBundle.kerEquivOfConjugate, def:FiberBundle.quotientEquivOfConjugate, lem:FiberBundle.fiberBundleSmallDC_barDh1_conj, lem:FiberBundle.barDh0_lTensor_quotient_conj}
Let $d_B$, $d_F$, $\varphi$ be as above with $\varphi$ acting as identity on homology. Then:
\[
  H_1(B \otimes_\varphi F) \;\simeq\; \bigl(\ker(d_B) \otimes_{\mathbb{F}_2} \operatorname{coker}(d_F)\bigr) \times \bigl(\operatorname{coker}(d_B) \otimes_{\mathbb{F}_2} \ker(d_F)\bigr).
\]
\end{theorem}

\begin{proof}
\leanok
\uses{def:FiberBundle.fiberBundleSmallDC, def:FiberBundle.ActsAsIdentityOnHomology, def:FiberBundle.lTensor_ker_equiv, def:FiberBundle.lTensor_quotient_equiv, def:FiberBundle.rTensor_ker_equiv, def:FiberBundle.rTensor_quotient_equiv, def:FiberBundle.kerEquivOfConjugate, def:FiberBundle.quotientEquivOfConjugate, lem:FiberBundle.fiberBundleSmallDC_barDh1_conj, lem:FiberBundle.barDh0_lTensor_quotient_conj}
Let $E = \operatorname{fiberBundleSmallDC}(d_B, d_F, \varphi)$.

\textbf{Step 1.} By the SmallDC homology equivalence at degree 1, $\operatorname{totH}_1(E) \simeq \operatorname{itH}_1^{\mathrm{fst}}(E) \times \operatorname{itH}_1^{\mathrm{snd}}(E)$, where $\operatorname{itH}_1^{\mathrm{fst}} = \ker(\bar{d}_h^0)$ and $\operatorname{itH}_1^{\mathrm{snd}} = \operatorname{vH}_{0,1}/\operatorname{im}(\bar{d}_h^1)$.

\textbf{Step 2a (first factor).} By Lemma~\ref{lem:FiberBundle.barDh0_lTensor_quotient_conj}, $\bar{d}_h^0$ is conjugate (via the quotient flatness equivalences $q_k$) to $d_B \otimes_r \operatorname{id}_{\operatorname{coker}(d_F)}$. Applying Definition~\ref{def:FiberBundle.kerEquivOfConjugate} gives $\ker(\bar{d}_h^0) \simeq \ker(d_B \otimes_r \operatorname{id}_{\operatorname{coker}(d_F)})$. Then Definition~\ref{def:FiberBundle.rTensor_ker_equiv} gives $\ker(d_B \otimes_r \operatorname{id}_{\operatorname{coker}(d_F)}) \simeq \ker(d_B) \otimes \operatorname{coker}(d_F)$.

\textbf{Step 2b (second factor).} By Lemma~\ref{lem:FiberBundle.fiberBundleSmallDC_barDh1_conj}, $\bar{d}_h^1$ is conjugate (via the kernel flatness equivalences $e_k$) to $d_B \otimes_r \operatorname{id}_{\ker(d_F)}$. Applying Definition~\ref{def:FiberBundle.quotientEquivOfConjugate} gives $\operatorname{coker}(\bar{d}_h^1) \simeq \operatorname{coker}(d_B \otimes_r \operatorname{id}_{\ker(d_F)})$. Then Definition~\ref{def:FiberBundle.rTensor_quotient_equiv} gives $\operatorname{coker}(d_B \otimes_r \operatorname{id}_{\ker(d_F)}) \simeq \operatorname{coker}(d_B) \otimes \ker(d_F)$.

Combining both factors via the product of equivalences ($\operatorname{prodCongr}$) with Step 1 yields:
\[
  H_1(B \otimes_\varphi F) \;\simeq\; \bigl(\ker(d_B)\otimes\operatorname{coker}(d_F)\bigr) \times \bigl(\operatorname{coker}(d_B)\otimes\ker(d_F)\bigr). \qedhere
\]
\end{proof}

%--- Def_13: AugmentedComplex ---
\chapter{Def 13: Augmented Complex}

\begin{definition}[Augmented Complex]
\label{def:FiberBundle.AugmentedComplex}
\lean{FiberBundle.AugmentedComplex}
\leanok

A complex $F = (F_1 \xrightarrow{\partial^F} F_0)$ is \emph{augmented} if there exists a linear map $\varepsilon : F_0 \to \mathbb{F}_2$ such that $\varepsilon \circ \partial^F = 0$. Concretely, for fixed natural numbers $n_2, m_2$ and a linear map $\partial^F : (\operatorname{Fin} n_2 \to \mathbb{F}_2) \to_{\mathbb{F}_2} (\operatorname{Fin} m_2 \to \mathbb{F}_2)$, an \texttt{AugmentedComplex} consists of:
\begin{itemize}
  \item A linear map $\varepsilon : (\operatorname{Fin} m_2 \to \mathbb{F}_2) \to_{\mathbb{F}_2} \mathbb{F}_2$ (the augmentation map),
  \item A proof that $\varepsilon \circ \partial^F = 0$ (the chain map condition).
\end{itemize}
\end{definition}

\begin{definition}[Augmentation Map at Each Degree]
\label{def:FiberBundle.augMap}
\lean{FiberBundle.augMap}
\leanok
\uses{def:FiberBundle.AugmentedComplex, def:FiberBundle.fiberSize}
Given a differential $\partial^F : (\operatorname{Fin} n_2 \to \mathbb{F}_2) \to_{\mathbb{F}_2} (\operatorname{Fin} m_2 \to \mathbb{F}_2)$ and an augmented complex $\mathtt{aug}$, the augmentation map at fiber degree $q \in \mathbb{Z}$ is defined as:
\[
  \operatorname{augMap}(\partial^F, \mathtt{aug}, q) :=
  \begin{cases}
    \varepsilon & \text{if } q = 0,\\
    0 & \text{otherwise.}
  \end{cases}
\]
This defines a linear map $(\operatorname{Fin}(\operatorname{fiberSize}(n_2, m_2, q)) \to \mathbb{F}_2) \to_{\mathbb{F}_2} \mathbb{F}_2$, corresponding to $\pi_0 = \varepsilon$ and $\pi_q = 0$ for $q \ne 0$.
\end{definition}

\begin{definition}[Summand Linear Map for $\pi_*$]
\label{def:FiberBundle.piStarSummandLin}
\lean{FiberBundle.piStarSummandLin}
\leanok
\uses{def:FiberBundle.AugmentedComplex, def:FiberBundle.augMap, def:FiberBundle.baseSize, def:FiberBundle.fiberSize}
For integers $p, q \in \mathbb{Z}$, the linear map underlying the $(p,q)$-summand of $\pi_*$ is:
\[
  \operatorname{piStarSummandLin}(\partial^F, \mathtt{aug}, p, q) :
  \Bigl(\operatorname{Fin}(\operatorname{baseSize}(n_1,m_1,p)) \to \mathbb{F}_2\Bigr)
  \otimes_{\mathbb{F}_2}
  \Bigl(\operatorname{Fin}(\operatorname{fiberSize}(n_2,m_2,q)) \to \mathbb{F}_2\Bigr)
  \;\to_{\mathbb{F}_2}\;
  \Bigl(\operatorname{Fin}(\operatorname{baseSize}(n_1,m_1,p)) \to \mathbb{F}_2\Bigr),
\]
given by $b \otimes f \mapsto \pi_q(f) \cdot b$, where $\pi_q = \operatorname{augMap}(\partial^F, \mathtt{aug}, q)$. Concretely, this is the composition
\[
  (\operatorname{TensorProduct.rid})_{\mathbb{F}_2} \circ (\operatorname{id} \otimes \operatorname{augMap}(\partial^F, \mathtt{aug}, q)).
\]
\end{definition}

\begin{lemma}[Summand Map at $q = 0$]
\label{lem:FiberBundle.piStarSummandLin_tmul_q0}
\lean{FiberBundle.piStarSummandLin_tmul_q0}
\leanok
\uses{def:FiberBundle.piStarSummandLin, def:FiberBundle.augMap, def:FiberBundle.AugmentedComplex, def:FiberBundle.baseSize, def:FiberBundle.fiberSize}
For any $b : \operatorname{Fin}(\operatorname{baseSize}(n_1,m_1,p)) \to \mathbb{F}_2$ and $f : \operatorname{Fin}(\operatorname{fiberSize}(n_2,m_2,0)) \to \mathbb{F}_2$,
\[
  \operatorname{piStarSummandLin}(\partial^F, \mathtt{aug}, p, 0)(b \otimes f) = \varepsilon(f) \cdot b.
\]
\end{lemma}

\begin{proof}
\leanok
\uses{def:FiberBundle.piStarSummandLin, def:FiberBundle.augMap}
By simplification using the definitions of \texttt{piStarSummandLin}, \texttt{augMap\_zero}, and the $\mathbb{F}_2$-module isomorphism $\operatorname{TensorProduct.rid}$, together with the formula for $\operatorname{lTensor}$ on pure tensors, the result follows by reflexivity.
\end{proof}

\begin{lemma}[Summand Map Vanishes for $q \ne 0$]
\label{lem:FiberBundle.piStarSummandLin_q_ne_zero}
\lean{FiberBundle.piStarSummandLin_q_ne_zero}
\leanok
\uses{def:FiberBundle.piStarSummandLin, def:FiberBundle.augMap}
For any $q \ne 0$,
\[
  \operatorname{piStarSummandLin}(\partial^F, \mathtt{aug}, p, q) = 0.
\]
\end{lemma}

\begin{proof}
\leanok
\uses{def:FiberBundle.piStarSummandLin, def:FiberBundle.augMap}
By simplification using the definitions of \texttt{piStarSummandLin} and \texttt{augMap} with the conditional \texttt{dif\_neg hq} (since $q \ne 0$), the augmentation map at degree $q$ is zero, so the composition $\operatorname{TensorProduct.rid} \circ (\operatorname{id} \otimes 0)$ reduces to zero by $\operatorname{LinearMap.comp_zero}$.
\end{proof}

\begin{definition}[Summand Morphism for $\pi_*$]
\label{def:FiberBundle.piStarSummandMor}
\lean{FiberBundle.piStarSummandMor}
\leanok
\uses{def:FiberBundle.fbObj, def:FiberBundle.baseSize, def:FiberBundle.piStarSummandLin, def:FiberBundle.AugmentedComplex}
For $n, p, q \in \mathbb{Z}$ with $p + q = n$, the $(p,q)$-component of $\pi_*$ is a morphism in $\operatorname{ModuleCat}_{\mathbb{F}_2}$:
\[
  \operatorname{piStarSummandMor}(\partial^F, \mathtt{aug}, n, p, q) :
  \operatorname{fbObj}(n_1,m_1,n_2,m_2,(p,q))
  \;\longrightarrow\;
  \mathbb{F}_2^{\operatorname{baseSize}(n_1,m_1,n)},
\]
defined by:
\begin{itemize}
  \item If $q = 0$: the map $b \otimes f \mapsto \varepsilon(f) \cdot b$, composed with a type transport via $\operatorname{Fin.cast}$ using the equality $p = n$ (from $p + 0 = n$),
  \item If $q \ne 0$: the zero morphism.
\end{itemize}
\end{definition}

\begin{definition}[Chain Map $\pi_*$]
\label{def:FiberBundle.piStarMor}
\lean{FiberBundle.piStarMor}
\leanok
\uses{def:FiberBundle.Connection, def:FiberBundle.fiberBundleDoubleComplex, def:FiberBundle.piStarSummandMor, def:FiberBundle.baseSize, def:FiberBundle.AugmentedComplex}
The chain map $\pi_* : \operatorname{Tot}(B \boxtimes_\varphi F)_n \to \mathbb{F}_2^{\operatorname{baseSize}(n_1,m_1,n)}$ at degree $n$ is assembled from the summand maps using $\operatorname{totalDesc}$:
\[
  \operatorname{piStarMor}(\partial^B, \partial^F, \varphi, \mathtt{aug}, n) :
  \bigl(\operatorname{fiberBundleDoubleComplex}(\partial^B, \partial^F, \varphi)\bigr).\operatorname{totalComplex}.X_n
  \;\longrightarrow\;
  \mathbb{F}_2^{\operatorname{baseSize}(n_1,m_1,n)},
\]
given by summing the contributions $\operatorname{piStarSummandMor}(\partial^F, \mathtt{aug}, n, p, q)$ over all $(p, q)$ with $p + q = n$.
\end{definition}

\begin{definition}[Base Differential]
\label{def:FiberBundle.baseDiff}
\lean{FiberBundle.baseDiff}
\leanok
\uses{def:FiberBundle.baseSize}
For $n, n' \in \mathbb{Z}$, the base differential is:
\[
  \operatorname{baseDiff}(n_1, m_1, \partial^B, n, n') :=
  \begin{cases}
    \partial^B & \text{if } n = 1 \text{ and } n' = 0,\\
    0 & \text{otherwise,}
  \end{cases}
\]
as a linear map $(\operatorname{Fin}(\operatorname{baseSize}(n_1,m_1,n)) \to \mathbb{F}_2) \to_{\mathbb{F}_2} (\operatorname{Fin}(\operatorname{baseSize}(n_1,m_1,n')) \to \mathbb{F}_2)$.
\end{definition}

\begin{lemma}[Augmentation Preserved by Chain Automorphism]
\label{lem:FiberBundle.aug_comp__0_eq}
\lean{FiberBundle.aug_comp_α₀_eq}
\leanok
\uses{def:FiberBundle.AugmentedComplex, def:FiberBundle.ChainAutomorphism, def:FiberBundle.ChainAutomorphism.ActsAsIdOnHomology}
Let $\psi$ be a chain automorphism of $F$ and suppose $\psi$ acts as the identity on homology (i.e., $\psi.\operatorname{ActsAsIdOnHomology}$ holds). Then
\[
  \varepsilon \circ \alpha_0 = \varepsilon,
\]
where $\alpha_0 = \psi.\alpha_0$ is the degree-$0$ component of $\psi$.
\end{lemma}

\begin{proof}
\leanok
\uses{def:FiberBundle.AugmentedComplex, def:FiberBundle.ChainAutomorphism, def:FiberBundle.ChainAutomorphism.ActsAsIdOnHomology}
We proceed by extensionality: it suffices to show $\varepsilon(\alpha_0(y)) = \varepsilon(y)$ for an arbitrary $y$. By the degree-$0$ condition of $\operatorname{ActsAsIdOnHomology}$, we have $\alpha_0(y) - y \in \operatorname{im}(\partial^F)$; that is, there exists $z$ such that $\partial^F(z) = \alpha_0(y) - y$. Since $\varepsilon \circ \partial^F = 0$, we rewrite to obtain
\[
  \varepsilon(\alpha_0(y) - y) = \varepsilon(\partial^F(z)) = 0.
\]
Applying $\varepsilon$ to the difference and using linearity of $\varepsilon$:
\[
  \varepsilon(\alpha_0(y)) - \varepsilon(y) = 0,
\]
which gives $\varepsilon(\alpha_0(y)) = \varepsilon(y)$ by $\operatorname{sub_eq_zero}$.
\end{proof}

\begin{lemma}[Augmentation Preserved by Connection]
\label{lem:FiberBundle.aug_preserved_by_connection}
\lean{FiberBundle.aug_preserved_by_connection}
\leanok
\uses{def:FiberBundle.Connection, def:FiberBundle.ActsAsIdentityOnHomology, def:FiberBundle.AugmentedComplex, lem:FiberBundle.aug_comp__0_eq}
Let $\varphi$ be a connection and suppose $\varphi$ acts as the identity on homology (i.e., $\operatorname{ActsAsIdentityOnHomology}(\partial^B, \partial^F, \varphi)$ holds). For any $b_1 : \operatorname{Fin}(n_1)$ and $b_0 : \operatorname{Fin}(m_1)$ with $\partial^B(\delta_{b_1})(b_0) \ne 0$,
\[
  \varepsilon \circ (\varphi\, b_1\, b_0).\alpha_0 = \varepsilon.
\]
\end{lemma}

\begin{proof}
\leanok
\uses{def:FiberBundle.ActsAsIdentityOnHomology, lem:FiberBundle.aug_comp__0_eq}
This is an immediate consequence of \texttt{aug\_comp\_$\alpha_0$\_eq} applied to the chain automorphism $\varphi\, b_1\, b_0$, using the fact that $\operatorname{ActsAsIdentityOnHomology}$ implies $(\varphi\, b_1\, b_0).\operatorname{ActsAsIdOnHomology}$ whenever $\partial^B(\delta_{b_1})(b_0) \ne 0$.
\end{proof}

\begin{lemma}[$\pi_*$ Intertwines Horizontal Differential at $q = 0$]
\label{lem:FiberBundle.piStarChainMap_q0_horizontal}
\lean{FiberBundle.piStarChainMap_q0_horizontal}
\leanok
\uses{def:FiberBundle.ActsAsIdentityOnHomology, def:FiberBundle.Connection, def:FiberBundle.piStarSummandLin, def:FiberBundle.twistedDhLin, def:FiberBundle.autAtDeg, lem:FiberBundle.twistedDhLin_tmul, lem:FiberBundle.aug_preserved_by_connection}
Assume $\operatorname{ActsAsIdentityOnHomology}(\partial^B, \partial^F, \varphi)$. On the $(p,0)$-summand, $\pi_*$ intertwines the horizontal twisted differential with the base differential:
\[
  \operatorname{piStarSummandLin}(\partial^F, \mathtt{aug}, 0, 0)
  \circ
  \operatorname{twistedDhLin}\!\Bigl(\partial^B,\; b_1,b_0 \mapsto \operatorname{autAtDeg}(\partial^F,\varphi,0,b_1,b_0)\Bigr)
  \;=\;
  \partial^B \circ \operatorname{piStarSummandLin}(\partial^F, \mathtt{aug}, 1, 0).
\]
\end{lemma}

\begin{proof}
\leanok
\uses{def:FiberBundle.piStarSummandLin, def:FiberBundle.twistedDhLin, def:FiberBundle.autAtDeg, lem:FiberBundle.twistedDhLin_tmul, lem:FiberBundle.aug_preserved_by_connection, lem:FiberBundle.piStarSummandLin_tmul_q0}
We apply tensor product extensionality: it suffices to verify the equality on pure tensors $b \otimes f$. Let $b : \operatorname{Fin}(\operatorname{baseSize}(n_1,m_1,1)) \to \mathbb{F}_2$ and $f : \operatorname{Fin}(\operatorname{fiberSize}(n_2,m_2,0)) \to \mathbb{F}_2$ be arbitrary.

By \texttt{twistedDhLin\_tmul}, we expand the left-hand side:
\[
  \operatorname{twistedDhLin}(\partial^B, \alpha)(b \otimes f) = \sum_{b_1, b_0} \partial^B(\delta_{b_1})(b_0) \cdot \bigl(\delta_{b_0} \otimes \alpha_{b_1,b_0}(f)\bigr).
\]
Applying $\operatorname{piStarSummandLin}(\partial^F, \mathtt{aug}, 0, 0)$ and using \texttt{piStarSummandLin\_tmul\_q0}, each summand becomes $\varepsilon(\alpha_{b_1,b_0}(f)) \cdot \partial^B(\delta_{b_1})(b_0) \cdot \delta_{b_0}$.

We then simplify using $\operatorname{map_sum}$ and $\operatorname{map_smul}$, and evaluate pointwise at an arbitrary index $j \in \operatorname{Fin}(m_1)$. The sum collapses via $\operatorname{Pi.single_apply}$ and $\operatorname{Finset.sum_ite_eq}$, yielding:
\[
  \sum_{b_1} b(b_1) \cdot \partial^B(\delta_{b_1})(j) \cdot \varepsilon\bigl(\alpha_{b_1,j}(f)\bigr).
\]
For each pair $(b_1, j)$, we distinguish two cases:
\begin{itemize}
  \item If $\partial^B(\delta_{b_1})(j) = 0$, the contribution is zero on both sides.
  \item If $\partial^B(\delta_{b_1})(j) \ne 0$, by \texttt{aug\_preserved\_by\_connection} and \texttt{autAtDeg\_zero}, we have $\varepsilon \circ \alpha_{b_1,j} = \varepsilon$, so $\varepsilon(\alpha_{b_1,j}(f)) = \varepsilon(f)$.
\end{itemize}
Thus each summand equals $b(b_1) \cdot \partial^B(\delta_{b_1})(j) \cdot \varepsilon(f)$, and we can factor out $\varepsilon(f)$:
\[
  \sum_{b_1} b(b_1) \cdot \partial^B(\delta_{b_1})(j) \cdot \varepsilon(f)
  = \varepsilon(f) \cdot \sum_{b_1} (b(b_1) \cdot \partial^B(\delta_{b_1}))(j).
\]
By linearity of $\partial^B$ and the expansion $b = \sum_{b_1} b(b_1) \cdot \delta_{b_1}$, we obtain $(\partial^B(b))(j) \cdot \varepsilon(f)$.

For the right-hand side, $\operatorname{piStarSummandLin}(\partial^F, \mathtt{aug}, 1, 0)(b \otimes f) = \varepsilon(f) \cdot b$ by \texttt{piStarSummandLin\_tmul\_q0}, and then $\partial^B(\varepsilon(f) \cdot b) = \varepsilon(f) \cdot \partial^B(b)$ by linearity. Both sides agree at every index $j$, completing the proof.
\end{proof}

\begin{lemma}[$\pi_*$ Vanishes on Vertical Differential at $q = 1$]
\label{lem:FiberBundle.piStarChainMap_q1_vertical}
\lean{FiberBundle.piStarChainMap_q1_vertical}
\leanok
\uses{def:FiberBundle.fiberDiff, def:FiberBundle.piStarSummandLin, def:FiberBundle.AugmentedComplex}
The composition of $\pi_*$ on the $(p,0)$-summand with $\operatorname{id} \otimes \partial^F$ (the vertical differential from degree $(p,1)$) is zero:
\[
  \operatorname{piStarSummandLin}(\partial^F, \mathtt{aug}, p, 0)
  \circ
  (\operatorname{id} \otimes \operatorname{fiberDiff}(n_2, m_2, \partial^F, 1, 0))
  = 0.
\]
\end{lemma}

\begin{proof}
\leanok
\uses{def:FiberBundle.fiberDiff, def:FiberBundle.piStarSummandLin, def:FiberBundle.AugmentedComplex, lem:FiberBundle.piStarSummandLin_tmul_q0}
We apply tensor product extensionality on pure tensors $b \otimes f$. By \texttt{lTensor\_tmul}, the left-hand side evaluates to:
\[
  \operatorname{piStarSummandLin}(\partial^F, \mathtt{aug}, p, 0)\bigl(b \otimes \operatorname{fiberDiff}(n_2, m_2, \partial^F, 1, 0)(f)\bigr).
\]
Since $\operatorname{fiberDiff}(n_2, m_2, \partial^F, 1, 0)(f) = \partial^F(f)$ by unfolding the definition of \texttt{fiberDiff} and using reflexivity, and by \texttt{piStarSummandLin\_tmul\_q0}, this equals $\varepsilon(\partial^F(f)) \cdot b$. Since $\varepsilon \circ \partial^F = 0$ (the chain map condition of the augmented complex), we have $\varepsilon(\partial^F(f)) = 0$, and therefore $0 \cdot b = 0$.
\end{proof}

\begin{definition}[Cochain Map $\pi^*$]
\label{def:FiberBundle.cochainPiStar}
\lean{FiberBundle.cochainPiStar}
\leanok
\uses{def:FiberBundle.Connection, def:FiberBundle.fiberBundleDoubleComplex, def:FiberBundle.piStarMor, def:FiberBundle.baseSize, def:FiberBundle.AugmentedComplex, def:FiberBundle.ActsAsIdentityOnHomology, def:FiberBundle.ChainAutomorphism, def:FiberBundle.autAtDeg, def:FiberBundle.fiberDiff, def:FiberBundle.twistedDhLin, lem:FiberBundle.twistedDhLin_tmul}
The cochain map $\pi^* : B^*_n \to \operatorname{Tot}(B \boxtimes_\varphi F)^*_n$ at degree $n$ is defined as the $\mathbb{F}_2$-dual (transpose) of the chain map $\pi_*$:
\[
  \operatorname{cochainPiStar}(\partial^B, \partial^F, \varphi, \mathtt{aug}, n)
  :=
  \operatorname{Module.Dual.transpose}_{\mathbb{F}_2}\bigl(\operatorname{piStarMor}(\partial^B, \partial^F, \varphi, \mathtt{aug}, n).\operatorname{hom}\bigr),
\]
which is a linear map
\[
  \operatorname{Module.Dual}_{\mathbb{F}_2}\!\bigl(\operatorname{Fin}(\operatorname{baseSize}(n_1,m_1,n)) \to \mathbb{F}_2\bigr)
  \;\to_{\mathbb{F}_2}\;
  \operatorname{Module.Dual}_{\mathbb{F}_2}\!\bigl(\operatorname{Tot}(B \boxtimes_\varphi F)_n\bigr).
\]
On a cochain $\beta \in B^*_p$, this gives $\pi^*(\beta)(b \otimes f) = \beta(b) \cdot \varepsilon(f)$ for $f \in F_0$, and $\pi^*(\beta)(b \otimes f) = 0$ for $f \in F_1$.
\end{definition}

%--- Thm_4: ProjectionInducesIsomorphism ---
\chapter{Thm 4: Projection Induces Isomorphism}

\begin{definition}[Projection Map $\pi_*$ at Degree 1]
\label{def:FiberBundle.piStarDeg1Lin}
\lean{FiberBundle.piStarDeg1Lin}
\leanok
\uses{def:FiberBundle.fiberBundleSmallDC}
Let $\varepsilon : (\mathbb{F}_2^{m_2}) \to \mathbb{F}_2$ be a linear map (the augmentation of $F_0$).
The projection $\pi_*$ at degree 1 is the linear map
\[
\pi_*^{(1)} : \bigl((\mathbb{F}_2^{n_1} \otimes \mathbb{F}_2^{m_2}) \times (\mathbb{F}_2^{m_1} \otimes \mathbb{F}_2^{n_2})\bigr) \;\longrightarrow\; \mathbb{F}_2^{n_1}
\]
defined by
\[
\pi_*^{(1)}(x_{10},\, x_{01}) \;=\; (\operatorname{rid} \circ (\operatorname{id} \otimes \varepsilon))(x_{10}),
\]
where $\operatorname{rid}$ denotes the right-unit isomorphism $V \otimes \mathbb{F}_2 \simeq V$.
Concretely, on simple tensors $\pi_*^{(1)}(b \otimes f,\, 0) = \varepsilon(f) \cdot b$, and $\pi_*^{(1)}$ is zero on the $\mathbb{F}_2^{m_1} \otimes \mathbb{F}_2^{n_2}$ component.
\end{definition}

\begin{lemma}[Augmentation Is Invariant Under the Connection]
\label{lem:FiberBundle.epsilon_comp_alpha0_eq}
\lean{FiberBundle.epsilon_comp_alpha0_eq}
\leanok
\uses{def:FiberBundle.ChainAutomorphism, def:FiberBundle.ActsAsIdentityOnHomology}
Let $d_F : \mathbb{F}_2^{n_2} \to \mathbb{F}_2^{m_2}$ be a linear map, let $\psi$ be a chain automorphism of $d_F$, and let $\varepsilon : \mathbb{F}_2^{m_2} \to \mathbb{F}_2$ be a linear map satisfying $\varepsilon \circ d_F = 0$. If $\psi$ acts as the identity on the homology of $d_F$ (i.e.\ $\psi.\operatorname{ActsAsIdOnHomology}$), then
\[
\varepsilon \circ \psi.\alpha_0 = \varepsilon.
\]
\end{lemma}

\begin{proof}
\leanok
\uses{def:FiberBundle.ChainAutomorphism, def:FiberBundle.ActsAsIdentityOnHomology}
Let $v$ be arbitrary. By extensionality it suffices to show $\varepsilon(\psi.\alpha_0\, v) = \varepsilon(v)$.

Since $\psi$ acts as the identity on homology at degree 0, we have $\psi.\alpha_0\, v - v \in \operatorname{im}(d_F)$. Obtaining $w$ such that $d_F(w) = \psi.\alpha_0\, v - v$, we compute:
\[
\varepsilon(\psi.\alpha_0\, v - v) = \varepsilon(d_F(w)) = (\varepsilon \circ d_F)(w) = 0,
\]
using the hypothesis $\varepsilon \circ d_F = 0$. Applying $\operatorname{map_sub}$ and $\operatorname{sub_eq_zero}$, we conclude $\varepsilon(\psi.\alpha_0\, v) = \varepsilon(v)$.
\end{proof}

\begin{lemma}[Chain Map Property of $\pi_*$ for the Twisted Differential]
\label{lem:FiberBundle.piStar_chainmap_twistedDh}
\lean{FiberBundle.piStar_chainmap_twistedDh}
\leanok
\uses{def:FiberBundle.ActsAsIdentityOnHomology, def:FiberBundle.fiberBundleSmallDC, def:FiberBundle.Connection}
Let $d_B$, $d_F$, and $\varphi$ be given, and let $\varepsilon$ satisfy $\varepsilon \circ d_F = 0$, with $\varphi$ acting as the identity on homology. For any $x \in \mathbb{F}_2^{n_1} \otimes \mathbb{F}_2^{m_2}$,
\[
(\operatorname{rid} \circ (\operatorname{id} \otimes \varepsilon))\bigl(\widetilde{d}_h^0(x)\bigr) = d_B\bigl((\operatorname{rid} \circ (\operatorname{id} \otimes \varepsilon))(x)\bigr),
\]
where $\widetilde{d}_h^0$ is the twisted horizontal differential at bidegree $(1,0)$.
That is, $\pi_*^{(1)}$ intertwines $\widetilde{d}_h^0$ with $d_B$.
\end{lemma}

\begin{proof}
\leanok
\uses{def:FiberBundle.ActsAsIdentityOnHomology, def:FiberBundle.fiberBundleSmallDC, lem:FiberBundle.epsilon_comp_alpha0_eq}
We proceed by induction on $x$ using the tensor product induction principle.

\textit{Zero case:} Both sides vanish trivially.

\textit{Additivity:} Both sides are linear, so the identity extends from simple tensors.

\textit{Simple tensor case $x = b \otimes f$:} Expanding $\widetilde{d}_h^0(b \otimes f)$ via \texttt{twistedDhLin\_tmul} and simplifying \texttt{autAtDeg}, we obtain
\[
\sum_{b_1, b_0} b(b_1) \cdot d_B(e_{b_1})(b_0) \cdot \varepsilon(\alpha_0^{\varphi(b_1,b_0)} f) \cdot e_{b_0}.
\]
For each term, by \texttt{epsilon\_comp\_alpha0\_eq} (when $d_B(e_{b_1})(b_0) \neq 0$, the chain automorphism $\varphi(b_1,b_0)$ acts as the identity on homology), we replace $\varepsilon(\alpha_0^{\varphi(b_1,b_0)} f)$ by $\varepsilon(f)$. After this substitution the sum becomes
\[
\sum_{b_1, b_0} b(b_1) \cdot d_B(e_{b_1})(b_0) \cdot \varepsilon(f) \cdot e_{b_0}.
\]
Proceeding pointwise and applying the basis expansion $d_B(b) = \sum_{b_1} b(b_1) \cdot d_B(e_{b_1})$, we obtain $\varepsilon(f) \cdot d_B(b) = d_B(\varepsilon(f) \cdot b)$, which equals $d_B((\operatorname{rid} \circ (\operatorname{id} \otimes \varepsilon))(b \otimes f))$.
\end{proof}

\begin{lemma}[$\pi_*$ Maps Cycles to $\ker(d_B)$]
\label{lem:FiberBundle.piStarDeg1_maps_cycles}
\lean{FiberBundle.piStarDeg1_maps_cycles}
\leanok
\uses{def:FiberBundle.ActsAsIdentityOnHomology, def:FiberBundle.fiberBundleSmallDC, lem:FiberBundle.piStar_chainmap_twistedDh}
Under the same hypotheses, for any $z \in \operatorname{totZ}_1$ (the total 1-cycles of $B \otimes_\varphi F$),
\[
\pi_*^{(1)}(z) \;\in\; \ker(d_B).
\]
\end{lemma}

\begin{proof}
\leanok
\uses{def:FiberBundle.ActsAsIdentityOnHomology, def:FiberBundle.fiberBundleSmallDC, lem:FiberBundle.piStar_chainmap_twistedDh}
We must show $d_B(\pi_*^{(1)}(z)) = 0$. Since $z \in \operatorname{totZ}_1 = \ker(d_1)$, we have $d_1(z) = 0$, i.e.\ $\widetilde{d}_h^0(z_1) + d_v^0(z_2) = 0$ (as $d_1 = \operatorname{coprod}(\widetilde{d}_h^0, d_v^0)$).

Applying $\operatorname{rid} \circ (\operatorname{id} \otimes \varepsilon)$ to this equation:
\[
(\operatorname{rid} \circ (\operatorname{id} \otimes \varepsilon))\bigl(\widetilde{d}_h^0(z_1) + d_v^0(z_2)\bigr) = 0.
\]
The $d_v^0$-term vanishes since $d_v^0 = \operatorname{lTensor}(d_F)$, and $(\operatorname{id} \otimes \varepsilon) \circ (\operatorname{id} \otimes d_F) = \operatorname{id} \otimes (\varepsilon \circ d_F) = 0$ by hypothesis. Thus:
\[
(\operatorname{rid} \circ (\operatorname{id} \otimes \varepsilon))\bigl(\widetilde{d}_h^0(z_1)\bigr) = 0.
\]
By the chain map property (Lemma \ref{lem:FiberBundle.piStar_chainmap_twistedDh}), this equals $d_B((\operatorname{rid} \circ (\operatorname{id} \otimes \varepsilon))(z_1)) = d_B(\pi_*^{(1)}(z)) = 0$.
\end{proof}

\begin{definition}[Restriction of $\pi_*$ to Cycles]
\label{def:FiberBundle.piStarOnCycles}
\lean{FiberBundle.piStarOnCycles}
\leanok
\uses{def:FiberBundle.ActsAsIdentityOnHomology, def:FiberBundle.fiberBundleSmallDC, lem:FiberBundle.piStarDeg1_maps_cycles}
The restriction of $\pi_*^{(1)}$ to $\operatorname{totZ}_1$, landing in $\ker(d_B)$:
\[
\pi_*|_{\operatorname{cycles}} : \operatorname{totZ}_1 \;\longrightarrow\; \ker(d_B), \quad z \;\mapsto\; \bigl(\pi_*^{(1)}(z),\; \pi_*^{(1)}(z) \in \ker(d_B)\bigr).
\]
This is a well-defined $\mathbb{F}_2$-linear map by Lemma \ref{lem:FiberBundle.piStarDeg1_maps_cycles}.
\end{definition}

\begin{lemma}[$\pi_*$ Maps Boundaries to Zero]
\label{lem:FiberBundle.piStarOnCycles_maps_boundaries}
\lean{FiberBundle.piStarOnCycles_maps_boundaries}
\leanok
\uses{def:FiberBundle.ActsAsIdentityOnHomology, def:FiberBundle.fiberBundleSmallDC, def:FiberBundle.piStarOnCycles}
Under the same hypotheses,
\[
\operatorname{totB}_1^{Z_1} \;\leq\; \ker(\pi_*|_{\operatorname{cycles}}),
\]
i.e.\ $\pi_*^{(1)}$ vanishes on the total 1-boundaries $\operatorname{totB}_1$ (embedded in $\operatorname{totZ}_1$).
\end{lemma}

\begin{proof}
\leanok
\uses{def:FiberBundle.ActsAsIdentityOnHomology, def:FiberBundle.fiberBundleSmallDC, def:FiberBundle.piStarOnCycles}
Let $z \in \operatorname{totB}_1^{Z_1}$. Rewriting $\operatorname{totB}_1^{Z_1}$ as the preimage under inclusion of $\operatorname{totB}_1 = \operatorname{im}(d_2)$, we obtain $w$ such that $d_2(w) = z.\operatorname{val}$. In particular, $z.\operatorname{val}.1 = d_v^1(w)$ (from the first component of $d_2$).

We must show $\pi_*^{(1)}(z.\operatorname{val}) = 0$. By definition of $\pi_*^{(1)}$ as $\operatorname{coprod}(\operatorname{rid} \circ (\operatorname{id} \otimes \varepsilon),\; 0)$, it suffices to show
\[
(\operatorname{rid} \circ (\operatorname{id} \otimes \varepsilon))(z.\operatorname{val}.1) = 0.
\]
Substituting $z.\operatorname{val}.1 = d_v^1(w) = \operatorname{lTensor}(d_F)(w)$:
\[
(\operatorname{rid} \circ (\operatorname{id} \otimes \varepsilon))(\operatorname{lTensor}(d_F)(w))
= (\operatorname{rid} \circ \operatorname{lTensor}(\varepsilon \circ d_F))(w) = 0,
\]
since $\varepsilon \circ d_F = 0$ by hypothesis.
\end{proof}

\begin{definition}[Induced Map $H_1(\pi_*)$ on Homology]
\label{def:FiberBundle.piStarH1}
\lean{FiberBundle.piStarH1}
\leanok
\uses{def:FiberBundle.ActsAsIdentityOnHomology, def:FiberBundle.fiberBundleSmallDC, def:FiberBundle.piStarOnCycles, lem:FiberBundle.piStarOnCycles_maps_boundaries}
The map induced by $\pi_*$ on degree-1 homology,
\[
H_1(\pi_*) : H_1(B \otimes_\varphi F) = \operatorname{totH}_1 \;\longrightarrow\; \ker(d_B),
\]
defined as the quotient lift
\[
H_1(\pi_*) \;=\; \operatorname{liftQ}\bigl(\operatorname{totB}_1^{Z_1},\; \pi_*|_{\operatorname{cycles}},\; \pi_*|_{\operatorname{cycles}}(\operatorname{totB}_1^{Z_1}) = 0\bigr).
\]
\end{definition}

\begin{lemma}[Surjectivity of $H_1(\pi_*)$]
\label{lem:FiberBundle.piStarH1_surjective}
\lean{FiberBundle.piStarH1_surjective}
\leanok
\uses{def:FiberBundle.ActsAsIdentityOnHomology, def:FiberBundle.fiberBundleSmallDC, def:FiberBundle.piStarH1}
Assume additionally that $\varepsilon : \mathbb{F}_2^{m_2} \to \mathbb{F}_2$ is surjective. Then
\[
H_1(\pi_*) : \operatorname{totH}_1 \;\twoheadrightarrow\; \ker(d_B)
\]
is surjective.
\end{lemma}

\begin{proof}
\leanok
\uses{def:FiberBundle.ActsAsIdentityOnHomology, def:FiberBundle.fiberBundleSmallDC, def:FiberBundle.piStarH1}
Since $\varepsilon$ is surjective, there exists $f_0 \in \mathbb{F}_2^{m_2}$ with $\varepsilon(f_0) = 1$.

Let $v \in \ker(d_B)$. We construct a preimage in $\operatorname{totH}_1$ mapping to $v$.

First, we show $\widetilde{d}_h^0(v \otimes f_0) \in \operatorname{im}(d_v^0)$: since $v \in \ker(d_B)$ and the connection acts as the identity on homology, for each basis vector $e_{b_1}$ with $d_B(e_{b_1})(b_0) \neq 0$, the automorphism $(\varphi(b_1,b_0)).\alpha_0\, f_0 - f_0 \in \operatorname{im}(d_F)$. After expanding via \texttt{twistedDhLin\_tmul}, the sum splits into a term $d_B(v) \otimes f_0 = 0$ (since $v \in \ker(d_B)$) and a sum of terms of the form $e_{b_0} \otimes d_F(g)$, each lying in $\operatorname{im}(\operatorname{lTensor}(d_F))$.

Obtaining $y$ with $d_v^0(y) = \widetilde{d}_h^0(v \otimes f_0)$, we set $z = (v \otimes f_0,\; y)$. Over $\mathbb{F}_2$, $\widetilde{d}_h^0(v \otimes f_0) + d_v^0(y) = 0$, so $z \in \operatorname{totZ}_1$.

The class $[z] \in \operatorname{totH}_1$ satisfies $H_1(\pi_*)([z]) = \pi_*|_{\operatorname{cycles}}(z)$. By definition, this equals $(\operatorname{rid} \circ (\operatorname{id} \otimes \varepsilon))(v \otimes f_0) = \varepsilon(f_0) \cdot v = 1 \cdot v = v$.
\end{proof}

\begin{lemma}[Equal Dimension: $\operatorname{totH}_1$ and $\ker(d_B)$]
\label{lem:FiberBundle.finrank_totH1_eq_finrank_kerDB}
\lean{FiberBundle.finrank_totH1_eq_finrank_kerDB}
\leanok
\uses{def:FiberBundle.ActsAsIdentityOnHomology, def:FiberBundle.fiberBundleSmallDC, thm:FiberBundle.fiberBundleHomology_n1}
Assume:
\begin{enumerate}
  \item $\varphi$ acts as the identity on the homology of $F$,
  \item the augmentation induces an isomorphism $\bar{\varepsilon} : H_0(F) = (\mathbb{F}_2^{m_2} / \operatorname{im}(d_F)) \xrightarrow{\sim} \mathbb{F}_2$,
  \item $H_0(B) = 0$, i.e.\ $\operatorname{im}(d_B) = \mathbb{F}_2^{m_1}$ (surjectivity).
\end{enumerate}
Then
\[
\dim_{\mathbb{F}_2} \operatorname{totH}_1 \;=\; \dim_{\mathbb{F}_2} \ker(d_B).
\]
\end{lemma}

\begin{proof}
\leanok
\uses{def:FiberBundle.ActsAsIdentityOnHomology, def:FiberBundle.fiberBundleSmallDC, thm:FiberBundle.fiberBundleHomology_n1}
By Theorem \ref{thm:FiberBundle.fiberBundleHomology_n1}, there is a linear equivalence
\[
e : \operatorname{totH}_1 \;\simeq\; \bigl(\ker(d_B) \otimes H_0(F)\bigr) \;\times\; \bigl(H_0(B) \otimes \ker(d_F)\bigr).
\]
Hence $\dim \operatorname{totH}_1 = \dim(\ker(d_B) \otimes H_0(F)) + \dim(H_0(B) \otimes \ker(d_F))$.

Since $\operatorname{im}(d_B) = \top$, the cokernel $H_0(B) = \mathbb{F}_2^{m_1} / \operatorname{im}(d_B)$ is the quotient by $\top$, which is subsingleton. Therefore $H_0(B) \otimes \ker(d_F)$ is also subsingleton (as a tensor product with a zero module), and its dimension is $0$.

Since $\bar{\varepsilon} : H_0(F) \simeq \mathbb{F}_2$, we have
\[
\dim(\ker(d_B) \otimes H_0(F)) = \dim(\ker(d_B) \otimes \mathbb{F}_2) = \dim(\ker(d_B)),
\]
using $\operatorname{lTensor}(\bar{\varepsilon})$ and the right-unit isomorphism $\ker(d_B) \otimes \mathbb{F}_2 \simeq \ker(d_B)$.

Combining: $\dim \operatorname{totH}_1 = \dim \ker(d_B) + 0 = \dim \ker(d_B)$.
\end{proof}

\begin{theorem}[$H_1(\pi_*)$ Is an Isomorphism]
\label{thm:FiberBundle.piStarH1_isLinearEquiv}
\lean{FiberBundle.piStarH1_isLinearEquiv}
\leanok
\uses{def:FiberBundle.ActsAsIdentityOnHomology, def:FiberBundle.ChainAutomorphism, def:FiberBundle.fiberBundleSmallDC, thm:FiberBundle.fiberBundleHomology_n1}
Under the following three hypotheses:
\begin{enumerate}
  \item the connection $\varphi$ acts as the identity on the homology of $F$,
  \item $\bar{\varepsilon} : H_0(F) \xrightarrow{\sim} \mathbb{F}_2$ (the augmentation $\varepsilon$ induces an isomorphism on $H_0(F)$),
  \item $\operatorname{im}(d_B) = \top$ (surjectivity of $d_B$, equivalently $H_0(B) = 0$),
\end{enumerate}
and assuming $\varepsilon$ itself is surjective, the map
\[
H_1(\pi_*) : H_1(B \otimes_\varphi F) \;\xrightarrow{\sim}\; H_1(B) = \ker(d_B)
\]
is a bijection (hence a linear isomorphism, since it is a linear map between finite-dimensional $\mathbb{F}_2$-vector spaces of equal dimension).
\end{theorem}

\begin{proof}
\leanok
\uses{def:FiberBundle.ActsAsIdentityOnHomology, def:FiberBundle.fiberBundleSmallDC, thm:FiberBundle.fiberBundleHomology_n1, lem:FiberBundle.piStarH1_surjective, lem:FiberBundle.finrank_totH1_eq_finrank_kerDB}
Surjectivity follows from Lemma \ref{lem:FiberBundle.piStarH1_surjective}. By Lemma \ref{lem:FiberBundle.finrank_totH1_eq_finrank_kerDB}, $\dim \operatorname{totH}_1 = \dim \ker(d_B)$.

We verify that both spaces are finite-dimensional: $\operatorname{totH}_1$ is finite-dimensional because by Theorem \ref{thm:FiberBundle.fiberBundleHomology_n1} it is linearly equivalent to a product of tensor products of finite-dimensional spaces; $\ker(d_B)$ is finite-dimensional as a submodule of a finite-dimensional space.

Since $H_1(\pi_*)$ is a surjective linear map between finite-dimensional spaces of equal dimension, it is also injective by the rank-nullity theorem (as applied via \texttt{LinearMap.injective\_iff\_surjective\_of\_finrank\_eq\_finrank}). Hence $H_1(\pi_*)$ is bijective.
\end{proof}

\begin{definition}[$H_1(\pi_*)$ as a Linear Equivalence]
\label{def:FiberBundle.piStarH1Equiv}
\lean{FiberBundle.piStarH1Equiv}
\leanok
\uses{def:FiberBundle.ActsAsIdentityOnHomology, def:FiberBundle.fiberBundleSmallDC, thm:FiberBundle.piStarH1_isLinearEquiv, def:FiberBundle.piStarH1}
Under the same hypotheses as Theorem \ref{thm:FiberBundle.piStarH1_isLinearEquiv}, the map $H_1(\pi_*)$ is packaged as a linear equivalence
\[
H_1(\pi_*)^{\simeq} : \operatorname{totH}_1 \;\xrightarrow{\;\sim\;} \ker(d_B),
\]
constructed via \texttt{LinearEquiv.ofBijective} applied to the bijection of Theorem \ref{thm:FiberBundle.piStarH1_isLinearEquiv}.
\end{definition}

\begin{definition}[Induced Cohomology Isomorphism $H^1(\pi^*)$]
\label{def:FiberBundle.piUpperStarH1Equiv}
\lean{FiberBundle.piUpperStarH1Equiv}
\leanok
\uses{def:FiberBundle.fiberBundleSmallDC, def:FiberBundle.piStarH1Equiv}
Under the same hypotheses, and assuming $\operatorname{totH}_1$ and $\ker(d_B)$ are finite-dimensional, the dual of the isomorphism $H_1(\pi_*)^{\simeq}$ yields a linear equivalence
\[
H^1(\pi^*) : \operatorname{Hom}(\ker(d_B),\, \mathbb{F}_2) \;\xrightarrow{\;\sim\;} \operatorname{Hom}(\operatorname{totH}_1,\, \mathbb{F}_2).
\]
Explicitly, for $f \in \operatorname{Hom}(\ker(d_B), \mathbb{F}_2)$, we set $H^1(\pi^*)(f) = f \circ H_1(\pi_*)^{\simeq}$, and the inverse sends $g \mapsto g \circ (H_1(\pi_*)^{\simeq})^{-1}$.
\end{definition}

\begin{theorem}[$H^1(\pi^*)$ Is an Isomorphism]
\label{thm:FiberBundle.piUpperStarH1_isLinearEquiv}
\lean{FiberBundle.piUpperStarH1_isLinearEquiv}
\leanok
\uses{def:FiberBundle.ActsAsIdentityOnHomology, def:FiberBundle.ChainAutomorphism, def:FiberBundle.fiberBundleSmallDC, thm:FiberBundle.fiberBundleHomology_n1}
Under the same hypotheses as Theorem \ref{thm:FiberBundle.piStarH1_isLinearEquiv}, the cochain map $\pi^*$ induces an isomorphism on degree-1 cohomology:
\[
H^1(\pi^*) : H^1(B) \;\xrightarrow{\;\sim\;} H^1(B \otimes_\varphi F).
\]
\end{theorem}

\begin{proof}
\leanok
\uses{def:FiberBundle.fiberBundleSmallDC, def:FiberBundle.piUpperStarH1Equiv, thm:FiberBundle.piStarH1_isLinearEquiv}
The map $H^1(\pi^*)$ is defined as $H^1(\pi^*) = H_1(\pi_*)^{\simeq,\operatorname{dual}}$, the dual of the linear equivalence $H_1(\pi_*)^{\simeq}$. Since $H_1(\pi_*)^{\simeq}$ is a linear equivalence (and in particular bijective), its dual $H^1(\pi^*)$ is also a linear equivalence, and its underlying linear map is bijective. This follows directly from \texttt{LinearEquiv.bijective} applied to $H^1(\pi^*)$.
\end{proof}

%--- Def_14: GraphExpansion ---
\chapter{Def 14: Graph Expansion}

\begin{lemma}[Adjacency Matrix is Hermitian]
\label{lem:GraphExpansion.adjMatrix_isHermitian}
\lean{GraphExpansion.adjMatrix_isHermitian}
\leanok

Let $G$ be a simple graph on a finite vertex type $V$ with decidable adjacency. Then the adjacency matrix $G.\operatorname{adjMatrix}(\mathbb{R})$ is Hermitian.
\end{lemma}

\begin{proof}
\leanok

We unfold the definition of \texttt{IsHermitian}. Since the scalar ring is $\mathbb{R}$, the conjugate transpose coincides with the ordinary transpose (by \texttt{Matrix.conjTranspose\_eq\_transpose\_of\_trivial}). Rewriting with this identity and with the symmetry of the adjacency matrix (\texttt{G.isSymm\_adjMatrix}), the goal reduces to reflexivity.
\end{proof}

\begin{definition}[Eigenvalues of the Adjacency Matrix]
\label{def:GraphExpansion.adjEigenvalues}
\lean{GraphExpansion.adjEigenvalues}
\leanok
\uses{lem:GraphExpansion.adjMatrix_isHermitian}
Let $G$ be a simple graph on a finite vertex type $V$ with decidable adjacency, and write $n = |V|$. The \emph{adjacency eigenvalues} of $G$ are the function
\[
  \operatorname{adjEigenvalues}(G) : \operatorname{Fin}(n) \to \mathbb{R}
\]
defined as the eigenvalues $\lambda_0 \geq \lambda_1 \geq \cdots \geq \lambda_{n-1}$ of $G.\operatorname{adjMatrix}(\mathbb{R})$, sorted in decreasing order. Concretely, this is \texttt{(adjMatrix\_isHermitian G).eigenvalues}$_0$.
\end{definition}

\begin{theorem}[Eigenvalues are Antitone]
\label{thm:GraphExpansion.adjEigenvalues_antitone}
\lean{GraphExpansion.adjEigenvalues_antitone}
\leanok
\uses{def:GraphExpansion.adjEigenvalues, lem:GraphExpansion.adjMatrix_isHermitian}
Let $G$ be a simple graph on a finite vertex type $V$ with decidable adjacency. The map $\operatorname{adjEigenvalues}(G) : \operatorname{Fin}(|V|) \to \mathbb{R}$ is antitone: for all $i \leq j$,
\[
  \operatorname{adjEigenvalues}(G)(i) \;\geq\; \operatorname{adjEigenvalues}(G)(j).
\]
\end{theorem}

\begin{proof}
\leanok
\uses{def:GraphExpansion.adjEigenvalues, lem:GraphExpansion.adjMatrix_isHermitian}
This follows directly by applying the field \texttt{eigenvalues}$_0$\texttt{\_antitone} of the Hermitian matrix structure to \texttt{adjMatrix\_isHermitian G}.
\end{proof}

\begin{definition}[Second-Largest Eigenvalue]
\label{def:GraphExpansion.secondLargestEigenvalue}
\lean{GraphExpansion.secondLargestEigenvalue}
\leanok
\uses{def:GraphExpansion.adjEigenvalues}
Let $G$ be a simple graph on a finite vertex type $V$ with decidable adjacency, and suppose $|V| \geq 2$. The \emph{second-largest eigenvalue} $\lambda_2$ of $G$ is
\[
  \lambda_2(G) \;:=\; \operatorname{adjEigenvalues}(G)\!\left(\langle 1, \cdot\rangle\right) \;\in\; \mathbb{R},
\]
i.e.\ the eigenvalue at index $1$ in the decreasing sequence $\lambda_0 \geq \lambda_1 \geq \cdots$.
\end{definition}

\begin{definition}[Spectral Gap]
\label{def:GraphExpansion.spectralGap}
\lean{GraphExpansion.spectralGap}
\leanok
\uses{def:GraphExpansion.secondLargestEigenvalue}
Let $G$ be a finite $s$-regular simple graph on $V$ with $|V| \geq 2$. The \emph{spectral gap} of $G$ is
\[
  \operatorname{spectralGap}(G, s) \;:=\; s - \lambda_2(G) \;\in\; \mathbb{R},
\]
where $\lambda_2(G)$ is the second-largest eigenvalue of $G$. A larger spectral gap indicates better expansion properties.
\end{definition}

\begin{definition}[Expander Family]
\label{def:GraphExpansion.IsExpanderFamily}
\lean{GraphExpansion.IsExpanderFamily}
\leanok
\uses{def:GraphExpansion.secondLargestEigenvalue}
Let $\iota$ be an index type, $n : \iota \to \mathbb{N}$ a family of sizes, $s \in \mathbb{N}$ a uniform regularity degree, and $G_i$ a simple $s$-regular graph on $\operatorname{Fin}(n(i))$ for each $i \in \iota$, with $n(i) \geq 2$ for all $i$. The family $\{G_i\}_{i \in \iota}$ is an \emph{expander family} if there exists $\varepsilon > 0$ such that
\[
  \lambda_2(G_i) \;\leq\; s - \varepsilon \qquad \text{for all } i \in \iota,
\]
i.e.\ the spectral gap of every graph in the family is uniformly bounded below by $\varepsilon$.
\end{definition}

\begin{definition}[Ramanujan Graph]
\label{def:GraphExpansion.IsRamanujan}
\lean{GraphExpansion.IsRamanujan}
\leanok
\uses{def:GraphExpansion.adjEigenvalues}
Let $G$ be a finite $s$-regular simple graph on $V$ with $|V| \geq 2$. The graph $G$ is \emph{Ramanujan} if every eigenvalue $\lambda$ of the adjacency matrix satisfying $|\lambda| < s$ also satisfies
\[
  |\lambda| \;\leq\; 2\sqrt{s - 1}.
\]
Formally, $G$ is Ramanujan if for all $i \in \operatorname{Fin}(|V|)$,
\[
  |\operatorname{adjEigenvalues}(G)(i)| < s \;\Longrightarrow\; |\operatorname{adjEigenvalues}(G)(i)| \leq 2\sqrt{s-1}.
\]
This is the optimal spectral bound for $s$-regular graphs.
\end{definition}

%--- Def_15: TannerCodeLocalSystem ---
\chapter{Def 15: Tanner Code Local System}

\begin{definition}[Labeling]
\label{def:Labeling}
\lean{Labeling}
\leanok

Let $G$ be a simple graph on a finite vertex type $V$, and let $s \in \mathbb{N}$. A \emph{labeling} for the Tanner code on $G$ is a collection
\[
  \Lambda = \{\Lambda_v\}_{v \in V}, \quad \Lambda_v : N(v) \xrightarrow{\;\sim\;} \operatorname{Fin}(s),
\]
where $N(v) = G.\operatorname{neighborSet}(v)$ is the set of neighbors of $v$ and $\Lambda_v$ is a bijection. The neighbor $(\Lambda_v)^{-1}(i)$ corresponds to the $i$-th position in the local code. The existence of such a labeling requires $|N(v)| = s$ for all $v$, which holds when $G$ is $s$-regular.
\end{definition}

\begin{definition}[Local View]
\label{def:localView}
\lean{localView}
\leanok
\uses{def:Labeling, def:F2}
Let $\Lambda$ be a labeling, $v \in V$ a vertex, and $c \in C_1(X) = (G.\operatorname{edgeSet} \to \mathbb{F}_2)$ a 1-chain. The \emph{local view at $v$} is the $\mathbb{F}_2$-linear map
\[
  \operatorname{localView}_\Lambda(v) : (G.\operatorname{edgeSet} \to \mathbb{F}_2) \to (\operatorname{Fin}(s) \to \mathbb{F}_2),
\]
defined by
\[
  \bigl(\operatorname{localView}_\Lambda(v)\,(c)\bigr)(i) = c\bigl(\{v,\, (\Lambda_v)^{-1}(i)\}\bigr),
\]
where $(\Lambda_v)^{-1}(i) \in N(v)$ is the neighbor of $v$ labeled $i$. This gives the restriction of $c$ to the star of $v$, viewed as an element of $\mathbb{F}_2^s$ via the labeling.
\end{definition}

\begin{lemma}[Local View Evaluation]
\label{lem:localView_apply}
\lean{localView_apply}
\leanok
\uses{def:localView, def:Labeling}
For any vertex $v \in V$, 1-chain $c : G.\operatorname{edgeSet} \to \mathbb{F}_2$, and index $i \in \operatorname{Fin}(s)$,
\[
  \operatorname{localView}_\Lambda(v)\,(c)\,(i) = c\bigl(\{v,\, (\Lambda_v)^{-1}(i)\}\bigr).
\]
\end{lemma}

\begin{proof}
\leanok
\uses{def:localView}
This holds by reflexivity, as it is the definitional unfolding of $\operatorname{localView}$.
\end{proof}

\begin{definition}[Differential]
\label{def:differential}
\lean{differential}
\leanok
\uses{def:localView, def:Labeling, def:F2}
Let $m \in \mathbb{N}$ and let $H : (\operatorname{Fin}(s) \to \mathbb{F}_2) \to_{\mathbb{F}_2} (\operatorname{Fin}(m) \to \mathbb{F}_2)$ be a parity-check map. The \emph{differential}
\[
  \partial_\Lambda^H : (G.\operatorname{edgeSet} \to \mathbb{F}_2) \to_{\mathbb{F}_2} (V \to \operatorname{Fin}(m) \to \mathbb{F}_2)
\]
is defined by
\[
  \partial_\Lambda^H(c)(v) = H\bigl(\operatorname{localView}_\Lambda(v)\,(c)\bigr).
\]
This is the component form of the differential $\partial : C_1(X) \to C_0(X) \otimes L_0$ from the paper, using the canonical isomorphism $(V \to \mathbb{F}_2) \otimes (\operatorname{Fin}(m) \to \mathbb{F}_2) \cong V \to \operatorname{Fin}(m) \to \mathbb{F}_2$. For an edge $e = \{v,w\}$, the formula $\partial(e) = e_v \otimes \partial^L(e_{\Lambda_v(w)}) + e_w \otimes \partial^L(e_{\Lambda_w(v)})$ becomes $\partial(e)(v) = \partial^L(e_{\Lambda_v(w)})$ and $\partial(e)(w) = \partial^L(e_{\Lambda_w(v)})$.
\end{definition}

\begin{lemma}[Differential Evaluation]
\label{lem:differential_apply}
\lean{differential_apply}
\leanok
\uses{def:differential, def:localView}
For any parity-check map $H$, 1-chain $c$, and vertex $v \in V$,
\[
  \partial_\Lambda^H(c)(v) = H\bigl(\operatorname{localView}_\Lambda(v)\,(c)\bigr).
\]
\end{lemma}

\begin{proof}
\leanok
\uses{def:differential, def:localView}
This holds by reflexivity, as it is the definitional unfolding of $\partial_\Lambda^H$.
\end{proof}

\begin{definition}[Tanner Code]
\label{def:tannerCode}
\lean{tannerCode}
\leanok
\uses{def:differential, def:Labeling, def:ClassicalCode}
Let $m \in \mathbb{N}$ and $H : (\operatorname{Fin}(s) \to \mathbb{F}_2) \to_{\mathbb{F}_2} (\operatorname{Fin}(m) \to \mathbb{F}_2)$ be a parity-check map. The \emph{Tanner code} is the submodule
\[
  C(X, L, \Lambda) := \ker\bigl(\partial_\Lambda^H\bigr) \subseteq C_1(X),
\]
where $C_1(X) = G.\operatorname{edgeSet} \to \mathbb{F}_2$ is the space of 1-chains. A 1-chain $c = \sum_e a_e \cdot e$ is a codeword if and only if for every vertex $v$, the restriction of the $a_e$ to edges incident to $v$, read via $\Lambda_v$, is a codeword of the local code $L = \ker(H)$.
\end{definition}

\begin{theorem}[Codeword Characterization]
\label{thm:mem_tannerCode_iff}
\lean{mem_tannerCode_iff}
\leanok
\uses{def:tannerCode, def:localView, def:Labeling}
Let $H : (\operatorname{Fin}(s) \to \mathbb{F}_2) \to_{\mathbb{F}_2} (\operatorname{Fin}(m) \to \mathbb{F}_2)$ be a parity-check map. A 1-chain $c \in C_1(X)$ belongs to the Tanner code if and only if for every vertex $v \in V$,
\[
  \operatorname{localView}_\Lambda(v)\,(c) \in \ker(H).
\]
\end{theorem}

\begin{proof}
\leanok
\uses{def:tannerCode, def:localView, def:differential}
We unfold the definitions: $c \in C(X, L, \Lambda)$ means $\partial_\Lambda^H(c) = 0$, i.e., the function $v \mapsto H(\operatorname{localView}_\Lambda(v)\,c)$ is the zero function.

$(\Rightarrow)$: Assume $\partial_\Lambda^H(c) = 0$. For an arbitrary vertex $v$, applying $\operatorname{congr_fun}$ to the equality yields $H(\operatorname{localView}_\Lambda(v)\,c) = 0$, so $\operatorname{localView}_\Lambda(v)\,c \in \ker(H)$.

$(\Leftarrow)$: Assume that for every $v$, $\operatorname{localView}_\Lambda(v)\,c \in \ker(H)$. By function extensionality, for each $v$ the hypothesis gives $H(\operatorname{localView}_\Lambda(v)\,c) = 0$. Therefore $\partial_\Lambda^H(c) = 0$, i.e., $c \in C(X, L, \Lambda)$.

By simplification using the definitions of $\operatorname{tannerCode}$, $\operatorname{LinearMap.mem_ker}$, and $\operatorname{differential}$, both directions are established.
\end{proof}

\begin{theorem}[Codeword Characterization via Classical Code]
\label{thm:mem_tannerCode_iff_localCode}
\lean{mem_tannerCode_iff_localCode}
\leanok
\uses{thm:mem_tannerCode_iff, def:ClassicalCode, def:ClassicalCode.ofParityCheck, thm:ClassicalCode.code_eq_ker}
Let $H : (\operatorname{Fin}(s) \to \mathbb{F}_2) \to_{\mathbb{F}_2} (\operatorname{Fin}(m) \to \mathbb{F}_2)$ be a parity-check map. A 1-chain $c \in C_1(X)$ belongs to the Tanner code if and only if for every vertex $v \in V$,
\[
  \operatorname{localView}_\Lambda(v)\,(c) \in \operatorname{ofParityCheck}(H).\operatorname{code},
\]
where $\operatorname{ofParityCheck}(H)$ is the classical code (Definition~\ref{def:ClassicalCode}) associated to the parity-check matrix $H$.
\end{theorem}

\begin{proof}
\leanok
\uses{thm:mem_tannerCode_iff, def:ClassicalCode.ofParityCheck, thm:ClassicalCode.code_eq_ker}
We rewrite using $\operatorname{mem_tannerCode_iff}$ (Theorem~\ref{thm:mem_tannerCode_iff}). It then suffices to show that $\operatorname{localView}_\Lambda(v)\,c \in \ker(H)$ if and only if $\operatorname{localView}_\Lambda(v)\,c \in \operatorname{ofParityCheck}(H).\operatorname{code}$. By simplification using the definitions of $\operatorname{ClassicalCode.ofParityCheck}$ and $\operatorname{ClassicalCode.code_eq_ker}$ (Theorem~\ref{thm:ClassicalCode.code_eq_ker}), these two conditions are definitionally equal.
\end{proof}

\begin{theorem}[Intersection Characterization of Tanner Code]
\label{thm:tannerCode_eq_iInf}
\lean{tannerCode_eq_iInf}
\leanok
\uses{def:tannerCode, def:localView, thm:mem_tannerCode_iff}
Let $H : (\operatorname{Fin}(s) \to \mathbb{F}_2) \to_{\mathbb{F}_2} (\operatorname{Fin}(m) \to \mathbb{F}_2)$ be a parity-check map. The Tanner code equals the intersection of pullbacks of the local code along all local views:
\[
  C(X, L, \Lambda) = \bigcap_{v \in V} \bigl(\operatorname{localView}_\Lambda(v)\bigr)^{-1}\!\bigl(\ker(H)\bigr).
\]
\end{theorem}

\begin{proof}
\leanok
\uses{def:tannerCode, def:localView, thm:mem_tannerCode_iff}
By extensionality, it suffices to show that for any $c \in C_1(X)$, $c \in C(X, L, \Lambda)$ if and only if $c \in \bigcap_{v \in V} (\operatorname{localView}_\Lambda(v))^{-1}(\ker(H))$.

By simplification using $\operatorname{Submodule.mem_iInf}$, $\operatorname{Submodule.mem_comap}$, and $\operatorname{LinearMap.mem_ker}$, the right-hand side states: for all $v$, $H(\operatorname{localView}_\Lambda(v)\,c) = 0$.

$(\Rightarrow)$: Assume $c \in C(X, L, \Lambda)$. For any $v$, by Theorem~\ref{thm:mem_tannerCode_iff}, $\operatorname{localView}_\Lambda(v)\,c \in \ker(H)$, so applying $\operatorname{LinearMap.mem_ker}$ gives $H(\operatorname{localView}_\Lambda(v)\,c) = 0$.

$(\Leftarrow)$: Assume that for all $v$, $H(\operatorname{localView}_\Lambda(v)\,c) = 0$. Then for each $v$, $\operatorname{localView}_\Lambda(v)\,c \in \ker(H)$ by $\operatorname{LinearMap.mem_ker}$. By Theorem~\ref{thm:mem_tannerCode_iff}, $c \in C(X, L, \Lambda)$.
\end{proof}

%--- Def_16: DualCode ---
\chapter{Def 16: Dual Code}

\begin{definition}[Dual Code]
\label{def:ClassicalCode.dualCode}
\lean{ClassicalCode.dualCode}
\leanok
\uses{def:ClassicalCode, def:ClassicalCode.dimension}
Let $\mathcal{C}$ be a classical code, i.e.\ a subspace $L \subseteq \mathbb{F}_2^n$. The \emph{dual code} is
\[
  L^\perp = \{ w \in \mathbb{F}_2^n : \langle w, c \rangle = 0 \text{ for all } c \in L \},
\]
where $\langle w, c \rangle = \sum_i w_i c_i$ is the standard $\mathbb{F}_2$-dot product. Concretely, $L^\perp$ is defined as the pullback of the dual annihilator $L^{\circ} \subseteq \operatorname{Dual}(\mathbb{F}_2^n)$ through the canonical linear isomorphism
\[
  \operatorname{dotProductEquiv} : \mathbb{F}_2^n \xrightarrow{\;\sim\;} \operatorname{Dual}(\mathbb{F}_2^n), \quad v \mapsto (w \mapsto v \cdot w).
\]
\end{definition}

\begin{theorem}[Membership Characterization of Dual Code]
\label{thm:ClassicalCode.mem_dualCode_iff}
\lean{ClassicalCode.mem_dualCode_iff}
\leanok
\uses{def:ClassicalCode, def:ClassicalCode.dimension, def:ClassicalCode.dualCode}
Let $\mathcal{C}$ be a classical code of block length $n$ and let $w \in \mathbb{F}_2^n$. Then
\[
  w \in L^\perp \iff \forall\, c \in L,\; w \cdot c = 0.
\]
\end{theorem}

\begin{proof}
\leanok
\uses{def:ClassicalCode.dualCode}
We unfold the definition of the dual code: $L^\perp = (L^{\circ}).\operatorname{comap}(\operatorname{dotProductEquiv})$. By simplification using the definitions of \texttt{Submodule.mem\_comap} and \texttt{Submodule.mem\_dualAnnihilator}, membership $w \in L^\perp$ is equivalent to $(\operatorname{dotProductEquiv}\, w)(c) = 0$ for all $c \in L$. We prove both directions separately. In each direction, let the respective hypothesis hold and let $c \in L$ be arbitrary; simplifying the definition of $\operatorname{dotProductEquiv}$, the condition $(\operatorname{dotProductEquiv}\, w)(c) = 0$ is exactly $w \cdot c = 0$, which gives the result.
\end{proof}

\begin{theorem}[Symmetric Membership Characterization]
\label{thm:ClassicalCode.mem_dualCode_iff'}
\lean{ClassicalCode.mem_dualCode_iff'}
\leanok
\uses{def:ClassicalCode.dualCode, thm:ClassicalCode.mem_dualCode_iff}
Let $w \in \mathbb{F}_2^n$. Then
\[
  w \in L^\perp \iff \forall\, c \in L,\; c \cdot w = 0.
\]
Over $\mathbb{F}_2$, the dot product is commutative: $w \cdot c = c \cdot w$.
\end{theorem}

\begin{proof}
\leanok
\uses{thm:ClassicalCode.mem_dualCode_iff}
We rewrite using \texttt{mem\_dualCode\_iff}. It then suffices to show that $(\forall\, c \in L,\, w \cdot c = 0) \Leftrightarrow (\forall\, c \in L,\, c \cdot w = 0)$. We prove both implications simultaneously: given each direction, let $c \in L$ be arbitrary and rewrite using commutativity of the dot product ($w \cdot c = c \cdot w$) to obtain the claim from the hypothesis.
\end{proof}

\begin{definition}[Equivalence with Dual Annihilator]
\label{def:ClassicalCode.dualCodeEquivDualAnnihilator}
\lean{ClassicalCode.dualCodeEquivDualAnnihilator}
\leanok
\uses{def:ClassicalCode, def:ClassicalCode.dualCode}
There is a $\mathbb{F}_2$-linear equivalence
\[
  L^\perp \;\simeq_{\mathbb{F}_2}\; L^{\circ},
\]
between the dual code $L^\perp$ and the dual annihilator $L^{\circ} \subseteq \operatorname{Dual}(\mathbb{F}_2^n)$, induced by the canonical isomorphism $\operatorname{dotProductEquiv}$. Since $(L^\perp).\mathtt{code} = L^{\circ}.\operatorname{comap}(\operatorname{dotProductEquiv})$, the restriction of $\operatorname{dotProductEquiv}^{-1}$ to $L^{\circ}$ gives the desired equivalence.
\end{definition}

\begin{theorem}[Dimension of the Dual Code]
\label{thm:ClassicalCode.dimension_dualCode}
\lean{ClassicalCode.dimension_dualCode}
\leanok
\uses{def:ClassicalCode, def:ClassicalCode.dimension, def:ClassicalCode.dualCode, def:ClassicalCode.dualCodeEquivDualAnnihilator}
Let $\mathcal{C}$ be a classical $[n,k,d]$-code. Then
\[
  \dim L^\perp = n - \dim L.
\]
\end{theorem}

\begin{proof}
\leanok
\uses{def:ClassicalCode.dimension, def:ClassicalCode.dualCode, def:ClassicalCode.dualCodeEquivDualAnnihilator}
We unfold the definition of dimension. First, we transfer the rank of $L^\perp$ to that of the dual annihilator: by the linear equivalence $\operatorname{dualCodeEquivDualAnnihilator}$, we have
\[
  \operatorname{finrank}_{\mathbb{F}_2}(L^\perp) = \operatorname{finrank}_{\mathbb{F}_2}(L^{\circ}).
\]
We rewrite using this equality. Next, we apply the dual annihilator dimension formula
\[
  \operatorname{finrank}_{\mathbb{F}_2}(L) + \operatorname{finrank}_{\mathbb{F}_2}(L^{\circ}) = \operatorname{finrank}_{\mathbb{F}_2}(\mathbb{F}_2^n),
\]
and use the fact that $\operatorname{finrank}_{\mathbb{F}_2}(\mathbb{F}_2^n) = n$ (by \texttt{Module.finrank\_fin\_fun}). The result $\operatorname{finrank}_{\mathbb{F}_2}(L^{\circ}) = n - \operatorname{finrank}_{\mathbb{F}_2}(L)$ then follows by integer arithmetic.
\end{proof}

\begin{theorem}[{Dual Code is an $[n, n-k, d_{L^\perp}]$-Code}]
\label{thm:ClassicalCode.dualCode_isNKDCode}
\lean{ClassicalCode.dualCode_isNKDCode}
\leanok
\uses{def:ClassicalCode, def:ClassicalCode.dimension, def:ClassicalCode.IsNKDCode, def:ClassicalCode.distance, def:ClassicalCode.dualCode, thm:ClassicalCode.dimension_dualCode}
Let $\mathcal{C}$ be a classical $[n,k,d]$-code. Then the dual code $L^\perp$ is an $[n,\, n-k,\, d_{L^\perp}]$-code, where $k = \dim L$ and $d_{L^\perp}$ is the minimum distance of $L^\perp$.
\end{theorem}

\begin{proof}
\leanok
\uses{def:ClassicalCode.IsNKDCode, thm:ClassicalCode.dimension_dualCode}
We verify both components of the \texttt{IsNKDCode} predicate. The first component, namely $\dim(L^\perp) = n - k$, is exactly \texttt{dimension\_dualCode}. The second component, that the distance of $L^\perp$ equals itself, holds by reflexivity.
\end{proof}

%--- Def_17: BinaryEntropyFunction ---
\chapter{Def 17: Binary Entropy Function}

\begin{definition}[Binary Entropy Function]
\label{def:binaryEntropy}
\lean{binaryEntropy}
\leanok
\uses{def:hammingWeight}
The binary entropy function $H_2 : \mathbb{R} \to \mathbb{R}$ is defined by
\[
H_2(\delta) = -\delta \log_2(\delta) - (1-\delta)\log_2(1-\delta).
\]
For $\delta \in (0,1)$, this gives the Shannon entropy of a Bernoulli$(\delta)$ distribution measured in bits. Outside $(0,1)$, the function extends by continuity, with $\log_2(0) = 0$ by Mathlib's convention.
\end{definition}

\begin{lemma}[Key Natural Number Inequality]
\label{lem:key_nat_ineq}
\lean{key_nat_ineq}
\leanok
\uses{def:binaryEntropy}
We have the following natural number inequality:
\[
100^{100} < 2^{50} \cdot 11^{11} \cdot 89^{89}.
\]
\end{lemma}

\begin{proof}
\leanok
This is verified by computation (\texttt{native\_decide}).
\end{proof}

\begin{lemma}[Key Fractional Inequality]
\label{lem:key_frac_ineq}
\lean{key_frac_ineq}
\leanok
\uses{lem:key_nat_ineq}
We have the following inequality in $\mathbb{R}$:
\[
\left(\frac{100}{11}\right)^{11} \cdot \left(\frac{100}{89}\right)^{89} < 2^{50}.
\]
\end{lemma}

\begin{proof}
\leanok
\uses{lem:key_nat_ineq}
Rewriting using $\left(\frac{a}{b}\right)^n = \frac{a^n}{b^n}$ and combining fractions, the inequality $\frac{100^{11} \cdot 100^{89}}{11^{11} \cdot 89^{89}} < 2^{50}$ is equivalent (after clearing the positive denominator $11^{11} \cdot 89^{89}$) to $100^{11} \cdot 100^{89} < 2^{50} \cdot 11^{11} \cdot 89^{89}$. By the ring identity $100^{11} \cdot 100^{89} = 100^{100}$, this reduces to $100^{100} < 2^{50} \cdot 11^{11} \cdot 89^{89}$, which holds by casting \texttt{key\_nat\_ineq} to $\mathbb{R}$.
\end{proof}

\begin{lemma}[$H_2(11/100) < 1/2$]
\label{lem:binaryEntropy_at_eleven_hundredths}
\lean{binaryEntropy_at_eleven_hundredths}
\leanok
\uses{def:binaryEntropy, lem:key_frac_ineq}
The binary entropy function satisfies $H_2(11/100) < 1/2$.
\end{lemma}

\begin{proof}
\leanok
\uses{def:binaryEntropy, lem:key_frac_ineq}
Unfolding the definition of $H_2$, the goal becomes
\[
-\frac{11}{100}\log_2\!\left(\frac{11}{100}\right) - \frac{89}{100}\log_2\!\left(\frac{89}{100}\right) < \frac{1}{2}.
\]
We establish auxiliary facts: $11/100 > 0$, $89/100 > 0$, $\log_2(11/100) < 0$ (since $11/100 < 1$ and the base is $> 1$), and $\log_2(89/100) < 0$ similarly. Using the identity $\log_2(a/b) = \log_2 a - \log_2 b$, we rewrite:
\[
\log_2\!\left(\frac{11}{100}\right) = -\log_2\!\left(\frac{100}{11}\right), \qquad \log_2\!\left(\frac{89}{100}\right) = -\log_2\!\left(\frac{100}{89}\right).
\]
After substitution and ring normalization, the goal becomes
\[
\frac{11}{100}\log_2\!\left(\frac{100}{11}\right) + \frac{89}{100}\log_2\!\left(\frac{100}{89}\right) < \frac{1}{2}.
\]
We establish the key step: $11 \cdot \log_2(100/11) + 89 \cdot \log_2(100/89) < 50$. Indeed, taking $\log_2$ of the inequality $\left(\frac{100}{11}\right)^{11}\cdot\left(\frac{100}{89}\right)^{89} < 2^{50}$ from \texttt{key\_frac\_ineq}, and using $\log_2(x^n) = n\log_2 x$ and $\log_2(2^{50}) = 50$, we obtain:
\[
11\log_2\!\left(\frac{100}{11}\right) + 89\log_2\!\left(\frac{100}{89}\right) < 50.
\]
Since $\log_2(100/11) > 0$ and $\log_2(100/89) > 0$, we conclude by linear arithmetic that $\frac{11}{100}\log_2\!\left(\frac{100}{11}\right) + \frac{89}{100}\log_2\!\left(\frac{100}{89}\right) < \frac{1}{2}$.
\end{proof}

\begin{lemma}[$H_2$ equals $\operatorname{binEntropy} / \log 2$]
\label{lem:binaryEntropy_eq_binEntropy_div}
\lean{binaryEntropy_eq_binEntropy_div}
\leanok
\uses{def:binaryEntropy}
For all $\delta \in \mathbb{R}$,
\[
H_2(\delta) = \frac{\operatorname{binEntropy}(\delta)}{\log 2},
\]
where $\operatorname{binEntropy}$ is Mathlib's natural-logarithm binary entropy function.
\end{lemma}

\begin{proof}
\leanok
\uses{def:binaryEntropy}
Unfolding $H_2(\delta) = -\delta\log_2\delta - (1-\delta)\log_2(1-\delta)$ and $\operatorname{binEntropy}(\delta) = -\delta\log\delta - (1-\delta)\log(1-\delta)$, and using $\log_2 x = \log x / \log 2$ (i.e., $\operatorname{logb}\, 2\, x = \log x / \log 2$), together with $\log(\cdot^{-1}) = -\log(\cdot)$, the equality follows by ring arithmetic.
\end{proof}

\begin{lemma}[Derivative of $H_2$]
\label{lem:hasDerivAt_binaryEntropy}
\lean{hasDerivAt_binaryEntropy}
\leanok
\uses{def:binaryEntropy, lem:binaryEntropy_eq_binEntropy_div}
For $x \in (0, 1)$, the binary entropy function has derivative
\[
H_2'(x) = \log_2\!\left(\frac{1-x}{x}\right).
\]
\end{lemma}

\begin{proof}
\leanok
\uses{def:binaryEntropy, lem:binaryEntropy_eq_binEntropy_div}
Let $x \in (0,1)$. We write $H_2 = \operatorname{binEntropy}(\cdot)/\log 2$ by \texttt{binaryEntropy\_eq\_binEntropy\_div}. Applying Mathlib's \texttt{Real.hasDerivAt\_binEntropy} (which gives derivative $\log((1-x)/x)$ for $\operatorname{binEntropy}$ at $x$, using $x \neq 0$ and $x \neq 1$), and dividing by the constant $\log 2$, we obtain that $H_2$ has a derivative at $x$. Using $\log_2((1-x)/x) = (\log(1-x) - \log x)/\log 2 = \log((1-x)/x)/\log 2$, the result follows.
\end{proof}

\begin{lemma}[$H_2$ is strictly increasing on $(0, 1/2)$]
\label{lem:binaryEntropy_strictMonoOn}
\lean{binaryEntropy_strictMonoOn}
\leanok
\uses{def:binaryEntropy, lem:binaryEntropy_eq_binEntropy_div}
The binary entropy function $H_2$ is strictly increasing on $(0, 1/2)$: for $0 < a < b < 1/2$, we have $H_2(a) < H_2(b)$.
\end{lemma}

\begin{proof}
\leanok
\uses{def:binaryEntropy, lem:binaryEntropy_eq_binEntropy_div}
Let $a, b \in (0, 1/2)$ with $a < b$. Using \texttt{binaryEntropy\_eq\_binEntropy\_div}, we write $H_2(a) = \operatorname{binEntropy}(a)/\log 2$ and $H_2(b) = \operatorname{binEntropy}(b)/\log 2$. Since $\log 2 > 0$ (as $1 < 2$), it suffices to show $\operatorname{binEntropy}(a) < \operatorname{binEntropy}(b)$. Both $a$ and $b$ lie in $[0, 2^{-1}]$ (verified by the hypotheses via $\operatorname{simp}$ and linear arithmetic), and Mathlib's \texttt{Real.binEntropy\_strictMonoOn} asserts strict monotonicity of $\operatorname{binEntropy}$ on $[0, 1/2]$, giving the result.
\end{proof}

\begin{theorem}[$H_2(\delta) < 1/2$ for $\delta \in (0, 11/100)$]
\label{thm:binaryEntropy_lt_half}
\lean{binaryEntropy_lt_half}
\leanok
\uses{def:binaryEntropy, lem:binaryEntropy_strictMonoOn, lem:binaryEntropy_at_eleven_hundredths}
For $\delta \in (0, 11/100)$, we have $H_2(\delta) < 1/2$.
\end{theorem}

\begin{proof}
\leanok
\uses{def:binaryEntropy, lem:binaryEntropy_strictMonoOn, lem:binaryEntropy_at_eleven_hundredths}
Let $\delta$ satisfy $0 < \delta < 11/100$. We note that $\delta < 1/2$ (by linear arithmetic from $11/100 < 1/2$), so $\delta \in (0, 1/2)$. Also $11/100 \in (0, 1/2)$. By \texttt{binaryEntropy\_strictMonoOn}, since $\delta < 11/100$, we have $H_2(\delta) < H_2(11/100)$. By \texttt{binaryEntropy\_at\_eleven\_hundredths}, $H_2(11/100) < 1/2$. Chaining these inequalities gives $H_2(\delta) < 1/2$.
\end{proof}

\begin{definition}[Hamming Ball]
\label{def:hammingBall}
\lean{hammingBall}
\leanok
\uses{def:hammingWeight}
The Hamming ball of radius $r \in \mathbb{R}$ centered at the origin in $\mathbb{F}_2^n$ is the finset
\[
B_0(r) = \{ v \in \mathbb{F}_2^n : \operatorname{wt}(v) \leq r \},
\]
where $\operatorname{wt}(v)$ denotes the Hamming weight of $v$. When $r < 0$, the ball is empty.
\end{definition}

\begin{theorem}[Membership in the Hamming Ball]
\label{thm:mem_hammingBall}
\lean{mem_hammingBall}
\leanok
\uses{def:hammingBall, def:hammingWeight}
For $v : \operatorname{Fin} n \to \mathbb{F}_2$ and $r \in \mathbb{R}$,
\[
v \in B_0(r) \iff \operatorname{wt}(v) \leq r.
\]
\end{theorem}

\begin{proof}
\leanok
\uses{def:hammingBall}
This follows immediately by unfolding the definition of $B_0(r)$ and applying \texttt{simp}.
\end{proof}

\begin{lemma}[Hamming Ball is Empty for Negative Radius]
\label{lem:hammingBall_empty_of_neg}
\lean{hammingBall_empty_of_neg}
\leanok
\uses{def:hammingBall, def:hammingWeight}
For $r < 0$, $B_0(r) = \emptyset$.
\end{lemma}

\begin{proof}
\leanok
\uses{def:hammingBall, thm:mem_hammingBall}
By extensionality, it suffices to show that no $v$ belongs to $B_0(r)$. Suppose $v \in B_0(r)$; then by \texttt{simp} on the definition, $\operatorname{wt}(v) \leq r < 0$. But $\operatorname{wt}(v) \geq 0$ since it is a natural number cast to $\mathbb{R}$, a contradiction by linear arithmetic. The other direction is trivial since $\emptyset$ has no members.
\end{proof}

\begin{lemma}[Hamming Ball is Nonempty for Nonnegative Radius]
\label{lem:hammingBall_nonempty}
\lean{hammingBall_nonempty}
\leanok
\uses{def:hammingBall, def:hammingWeight, thm:mem_hammingBall}
For $r \geq 0$, $B_0(r)$ is nonempty.
\end{lemma}

\begin{proof}
\leanok
\uses{def:hammingBall, thm:mem_hammingBall}
We exhibit the zero vector $0 \in \mathbb{F}_2^n$ as a witness. By \texttt{mem\_hammingBall}, it suffices to show $\operatorname{wt}(0) \leq r$. Using \texttt{simp} with the definitions of \texttt{hammingWeight}, \texttt{hammingNorm}, and \texttt{hammingDist\_self}, we get $\operatorname{wt}(0) = 0$, and $0 \leq r$ holds by assumption.
\end{proof}

\begin{lemma}[Bernoulli Weight Monotonicity]
\label{lem:bernoulli_weight_mono}
\lean{bernoulli_weight_mono}
\leanok
\uses{def:hammingWeight}
Let $0 < \delta \leq 1/2$ and $a \leq k \leq n$. Then
\[
\delta^k (1-\delta)^{n-k} \leq \delta^a (1-\delta)^{n-a}.
\]
\end{lemma}

\begin{proof}
\leanok
We establish $0 < 1 - \delta$ (by linear arithmetic from $\delta \leq 1/2$) and $\delta \leq 1 - \delta$ (since $\delta \leq 1/2$). Writing $k = a + (k - a)$ and $n - a = (k-a) + (n-k)$, and rewriting the powers accordingly, the left-hand side becomes
\[
\delta^a \cdot \delta^{k-a} \cdot (1-\delta)^{n-k} = \delta^a \cdot (1-\delta)^{n-k} \cdot \delta^{k-a}.
\]
Since $\delta \leq 1 - \delta$ and $\delta > 0$, we have $\delta^{k-a} \leq (1-\delta)^{k-a}$ by \texttt{pow\_le\_pow\_left}. Applying \texttt{mul\_le\_mul\_of\_nonneg\_left} twice (first using $(1-\delta)^{n-k} > 0$, then $\delta^a > 0$), we obtain
\[
\delta^a \cdot (1-\delta)^{n-k} \cdot \delta^{k-a} \leq \delta^a \cdot (1-\delta)^{n-k} \cdot (1-\delta)^{k-a} = \delta^a \cdot (1-\delta)^{n-a},
\]
completing the proof by ring arithmetic.
\end{proof}

\begin{lemma}[Bernoulli Weight Sum]
\label{lem:bernoulli_weight_sum}
\lean{bernoulli_weight_sum}
\leanok
\uses{def:hammingWeight, def:F2}
For $0 < \delta \leq 1/2$,
\[
\sum_{v : \mathbb{F}_2^n} \delta^{\operatorname{wt}(v)} (1-\delta)^{n - \operatorname{wt}(v)} = 1.
\]
\end{lemma}

\begin{proof}
\leanok
\uses{def:hammingWeight, def:F2}
We note that every element of $\mathbb{F}_2$ is either $0$ or $1$ (by \texttt{decide}). We define an equivalence $e : (\operatorname{Fin} n \to \mathbb{F}_2) \simeq \operatorname{Finset}(\operatorname{Fin} n)$ by mapping $v$ to its support $\{i \in \operatorname{Fin} n : v_i \neq 0\}$, with inverse sending a subset $s$ to the indicator function $\mathbf{1}_s$. The left inverse is verified case-by-case on $\mathbb{F}_2$, and the right inverse by a boolean case split. Under this equivalence, the cardinality of $e(v)$ equals $\operatorname{wt}(v)$ by definition. Applying Mathlib's \texttt{Fin.sum\_pow\_mul\_eq\_add\_pow}, we get
\[
\sum_{s : \operatorname{Finset}(\operatorname{Fin} n)} \delta^{|s|}(1-\delta)^{n-|s|} = (\delta + (1-\delta))^n = 1^n = 1.
\]
Transporting along $e$ via \texttt{Equiv.sum\_comp} and a congruence argument gives the result.
\end{proof}

\begin{lemma}[Entropy Logarithm Inequality]
\label{lem:entropy_logb_ineq}
\lean{entropy_logb_ineq}
\leanok
\uses{def:binaryEntropy, def:hammingWeight}
Let $0 < \delta \leq 1/2$, $k \leq n$, and $(k : \mathbb{R}) + 1 \leq \delta \cdot n$. Then
\[
-(k \log_2 \delta + (n - k)\log_2(1-\delta)) \leq H_2(\delta) \cdot n.
\]
\end{lemma}

\begin{proof}
\leanok
\uses{def:binaryEntropy}
Unfolding $H_2(\delta) = -\delta\log_2\delta - (1-\delta)\log_2(1-\delta)$, we establish that $\delta n - k > 0$ (by linear arithmetic from $k + 1 \leq \delta n$) and that $\log_2 \delta \leq \log_2(1-\delta)$ (by strict monotonicity of $\log_2$ on $(0,\infty)$, applied at points $\delta < 1 - \delta$ since $\delta \leq 1/2$). Casting $n - k$ appropriately, the inequality then follows by nonlinear arithmetic (\texttt{nlinarith}), using $(\delta n - k) \cdot (\log_2(1-\delta) - \log_2\delta) \geq 0$.
\end{proof}

\begin{theorem}[Hamming Ball Bound]
\label{thm:hammingBallBound}
\lean{hammingBallBound}
\leanok
\uses{def:hammingBall, def:binaryEntropy, def:hammingWeight, lem:bernoulli_weight_mono, lem:bernoulli_weight_sum, lem:entropy_logb_ineq, thm:mem_hammingBall}
Let $n \geq 1$, $0 < \delta \leq 1/2$, and $k + 1 \leq \delta n$. Then
\[
|B_0(k)| \leq 2^{H_2(\delta) \cdot n}.
\]
\end{theorem}

\begin{proof}
\leanok
\uses{def:hammingBall, def:binaryEntropy, lem:bernoulli_weight_mono, lem:bernoulli_weight_sum, lem:entropy_logb_ineq, thm:mem_hammingBall}
We establish $0 < 1 - \delta$ by linear arithmetic. We show $k \leq n$: otherwise $k > n$, so $\delta n < n$, but $k + 1 \leq \delta n < n \leq k$, a contradiction by linear arithmetic. Let $p = \delta^k(1-\delta)^{n-k} > 0$ (by positivity).

\textbf{Step 1:} $|B_0(k)| \cdot \delta^k(1-\delta)^{n-k} \leq 1$. We compute:
\[
|B_0(k)| \cdot p = \sum_{v \in B_0(k)} p \leq \sum_{v \in B_0(k)} \delta^{\operatorname{wt}(v)}(1-\delta)^{n-\operatorname{wt}(v)} \leq \sum_{v : \mathbb{F}_2^n} \delta^{\operatorname{wt}(v)}(1-\delta)^{n-\operatorname{wt}(v)} = 1.
\]
The first inequality uses \texttt{bernoulli\_weight\_mono}: for $v \in B_0(k)$, $\operatorname{wt}(v) \leq k$, so $\delta^k(1-\delta)^{n-k} \leq \delta^{\operatorname{wt}(v)}(1-\delta)^{n-\operatorname{wt}(v)}$. The second uses \texttt{Finset.sum\_le\_univ\_sum\_of\_nonneg} (terms are nonneg). The final equality is \texttt{bernoulli\_weight\_sum}.

\textbf{Step 2:} $p^{-1} \leq 2^{H_2(\delta) n}$. This is equivalent (by \texttt{Real.logb\_le\_iff\_le\_rpow}) to $\log_2(p^{-1}) \leq H_2(\delta) n$. Using $\log_2(p^{-1}) = -\log_2 p$ and $\log_2(\delta^k (1-\delta)^{n-k}) = k\log_2\delta + (n-k)\log_2(1-\delta)$, this becomes $-(k\log_2\delta + (n-k)\log_2(1-\delta)) \leq H_2(\delta)n$, which holds by \texttt{entropy\_logb\_ineq}.

\textbf{Conclusion:} From Step~1, $|B_0(k)| \leq p^{-1}$, and from Step~2, $p^{-1} \leq 2^{H_2(\delta)n}$, so by linear arithmetic $|B_0(k)| \leq 2^{H_2(\delta)n}$.
\end{proof}

%--- Def_18: CheegerConstant ---
\chapter{Def 18: Cheeger Constant}

\begin{definition}[Edge Boundary]
\label{def:CheegerConstant.edgeBoundary}
\lean{CheegerConstant.edgeBoundary}
\leanok

Let $G = (V, E)$ be a finite simple graph with a decidable adjacency relation, and let $S \subseteq V$ be a finite subset of vertices. The \emph{edge boundary} $\delta S$ is the finite set of edges in $G$ that cross the cut determined by $S$:
\[
  \delta S \;=\; \bigl\{\, e \in E(G) \;\big|\; \text{one endpoint of } e \text{ lies in } S \text{ and the other does not}\,\bigr\}.
\]
Formally, using the quotient representation $\operatorname{Sym2}(V)$ for unordered pairs, we set
\[
  \delta S \;=\; \bigl\{\, e \in G.\operatorname{edgeFinset} \;\big|\; (e.\operatorname{out}_1 \in S \land e.\operatorname{out}_2 \notin S) \lor (e.\operatorname{out}_2 \in S \land e.\operatorname{out}_1 \notin S)\,\bigr\}.
\]
\end{definition}

\begin{theorem}[Membership in the Edge Boundary]
\label{thm:CheegerConstant.mem_edgeBoundary}
\lean{CheegerConstant.mem_edgeBoundary}
\leanok
\uses{def:CheegerConstant.edgeBoundary}
Let $G$ be a finite simple graph with decidable adjacency, $S \subseteq V$ a finite subset, and $e \in \operatorname{Sym2}(V)$. Then
\[
  e \in \delta S \;\Longleftrightarrow\; e \in E(G) \;\land\; \exists\, a \in S,\, b \notin S \text{ such that } \{a, b\} = e.
\]
\end{theorem}

\begin{proof}
\leanok
\uses{def:CheegerConstant.edgeBoundary}
We prove both directions. 

$(\Rightarrow)$: Assume $e \in \delta S$. By the definition of $\delta S$ (filtering $G.\operatorname{edgeFinset}$), we obtain $e \in E(G)$ and the disjunction $(e.\operatorname{out}_1 \in S \land e.\operatorname{out}_2 \notin S) \lor (e.\operatorname{out}_2 \in S \land e.\operatorname{out}_1 \notin S)$. In the first case, set $a = e.\operatorname{out}_1$ and $b = e.\operatorname{out}_2$; then $\{a, b\} = e$ by the identity $e = \operatorname{Sym2.mk}(e.\operatorname{out})$, and $a \in S$, $b \notin S$. In the second case, set $a = e.\operatorname{out}_2$ and $b = e.\operatorname{out}_1$; then $\{a, b\} = \{b, a\} = e$ by symmetry of $\operatorname{Sym2}$, and again $a \in S$, $b \notin S$.

$(\Leftarrow)$: Assume $e \in E(G)$ and there exist $a \in S$, $b \notin S$ with $\{a, b\} = e$. We substitute $e = \{a, b\}$. Let $(p, q) = \{a, b\}.\operatorname{out}$; by the characterization $\operatorname{Sym2.mk_eq_mk_iff}$, either $(p, q) = (a, b)$ or $(p, q) = (b, a)$. In the first sub-case, $p = a \in S$ and $q = b \notin S$, giving the left disjunct. In the second sub-case, $p = b$ and $q = a$, so $q = a \in S$ and $p = b \notin S$, giving the right disjunct. Hence $e \in \delta S$.
\end{proof}

\begin{theorem}[Edge Boundary Subset of Edge Set]
\label{thm:CheegerConstant.edgeBoundary_subset_edgeFinset}
\lean{CheegerConstant.edgeBoundary_subset_edgeFinset}
\leanok
\uses{def:CheegerConstant.edgeBoundary}
For any finite simple graph $G$ with decidable adjacency and any $S \subseteq V$, the edge boundary satisfies $\delta S \subseteq E(G)$.
\end{theorem}

\begin{proof}
\leanok
\uses{def:CheegerConstant.edgeBoundary}
This follows immediately from the fact that $\delta S$ is defined as a filter of $G.\operatorname{edgeFinset}$, and any filter of a finset is a subset of it: $\operatorname{Finset.filter_subset}$.
\end{proof}

\begin{theorem}[Adjacent Edge Lies in Boundary]
\label{thm:CheegerConstant.adj_mem_edgeBoundary}
\lean{CheegerConstant.adj_mem_edgeBoundary}
\leanok
\uses{def:CheegerConstant.edgeBoundary, thm:CheegerConstant.mem_edgeBoundary}
Let $G$ be a finite simple graph, $S \subseteq V$ a finite subset, and $u, v \in V$ with $G.\operatorname{Adj}(u, v)$. If $u \in S$ and $v \notin S$, then $\{u, v\} \in \delta S$.
\end{theorem}

\begin{proof}
\leanok
\uses{def:CheegerConstant.edgeBoundary, thm:CheegerConstant.mem_edgeBoundary}
We rewrite using the membership characterization of the edge boundary (Theorem \ref{thm:CheegerConstant.mem_edgeBoundary}). It suffices to provide $\{u, v\} \in E(G)$ (which holds since $G.\operatorname{Adj}(u, v)$ implies $\{u, v\} \in G.\operatorname{edgeFinset}$) and the witnesses $a = u \in S$, $b = v \notin S$ with $\{u, v\} = \{u, v\}$ (by reflexivity).
\end{proof}

\begin{definition}[Eligible Subset]
\label{def:CheegerConstant.EligibleSubset}
\lean{CheegerConstant.EligibleSubset}
\leanok

A finite subset $S \subseteq V$ is called \emph{eligible} (for the Cheeger constant) if:
\begin{enumerate}
  \item $S$ is nonempty: $S \neq \emptyset$, and
  \item $S$ is at most half the vertex set: $2|S| \leq |V|$.
\end{enumerate}
Formally, $\operatorname{EligibleSubset}(S) \;\Longleftrightarrow\; S.\operatorname{Nonempty} \land 2 \cdot |S| \leq |V|$.
\end{definition}

\begin{definition}[Isoperimetric Ratio]
\label{def:CheegerConstant.isoperimetricRatio}
\lean{CheegerConstant.isoperimetricRatio}
\leanok
\uses{def:CheegerConstant.edgeBoundary}
Let $G$ be a finite simple graph with decidable adjacency, and let $S \subseteq V$ be a finite subset. The \emph{isoperimetric ratio} of $S$ is
\[
  h(S) \;=\; \frac{|\delta S|}{|S|} \;\in\; \mathbb{R},
\]
where $|\delta S|$ denotes the cardinality of the edge boundary and $|S|$ the cardinality of $S$.
\end{definition}

\begin{definition}[Cheeger Constant]
\label{def:CheegerConstant.cheegerConstant}
\lean{CheegerConstant.cheegerConstant}
\leanok
\uses{def:CheegerConstant.EligibleSubset, def:CheegerConstant.isoperimetricRatio}
Let $G$ be a finite simple graph with decidable adjacency. The \emph{Cheeger constant} (or \emph{isoperimetric number}) of $G$ is
\[
  h(G) \;=\; \inf_{\substack{S \subseteq V \\ \operatorname{EligibleSubset}(S)}} \frac{|\delta S|}{|S|} \;=\; \operatorname{sInf}\bigl\{\, r \in \mathbb{R} \;\big|\; \exists\, S,\; \operatorname{EligibleSubset}(S) \land h(S) = r\,\bigr\},
\]
where the infimum is taken over all eligible subsets $S$ (nonempty subsets with $2|S| \leq |V|$). When no eligible subset exists (i.e., $|V| < 2$), the set is empty and $h(G) = 0$ by the convention $\operatorname{sInf}\,\emptyset = 0$ for $\mathbb{R}$.
\end{definition}

\begin{theorem}[Eligible Subsets Exist When $|V| \geq 2$]
\label{thm:CheegerConstant.eligibleSubset_nonempty}
\lean{CheegerConstant.eligibleSubset_nonempty}
\leanok
\uses{def:CheegerConstant.EligibleSubset}
If $|V| \geq 2$, then there exists an eligible subset $S \subseteq V$.
\end{theorem}

\begin{proof}
\leanok
\uses{def:CheegerConstant.EligibleSubset}
Since $|V| \geq 2 > 0$, the type $V$ is nonempty. Obtain a vertex $v \in V$. Consider the singleton $S = \{v\}$. Then $S$ is nonempty by $\operatorname{Finset.singleton_nonempty}$, and $|S| = 1$, so $2 \cdot |S| = 2 \leq |V|$ by hypothesis. By integer arithmetic ($\omega$), the condition $2 \cdot 1 \leq |V|$ follows from $|V| \geq 2$. Hence $\{v\}$ is an eligible subset.
\end{proof}

\begin{theorem}[Cheeger Constant is Nonnegative]
\label{thm:CheegerConstant.cheegerConstant_nonneg}
\lean{CheegerConstant.cheegerConstant_nonneg}
\leanok
\uses{def:CheegerConstant.cheegerConstant, def:CheegerConstant.EligibleSubset, def:CheegerConstant.isoperimetricRatio, thm:CheegerConstant.eligibleSubset_nonempty}
For any finite simple graph $G$ with decidable adjacency, $h(G) \geq 0$.
\end{theorem}

\begin{proof}
\leanok
\uses{def:CheegerConstant.cheegerConstant, def:CheegerConstant.EligibleSubset, def:CheegerConstant.isoperimetricRatio, thm:CheegerConstant.eligibleSubset_nonempty}
We unfold the definition of $h(G)$ and consider two cases.

\textbf{Case 1: $|V| \geq 2$.} By Theorem \ref{thm:CheegerConstant.eligibleSubset_nonempty}, obtain an eligible subset $S$. The set $\bigl\{\, r \;\big|\; \exists\, S',\;\operatorname{EligibleSubset}(S') \land h(S') = r\,\bigr\}$ is nonempty (witnessed by $h(S)$). We apply $\operatorname{le_csInf}$: it suffices to show every element $r = h(S')$ in the set satisfies $0 \leq r$. But $r = |\delta S'| / |S'|$, which is nonneg by positivity (both numerator and denominator are nonneg naturals cast to $\mathbb{R}$).

\textbf{Case 2: $|V| < 2$.} Since $|V| < 2$, any eligible subset $S$ would require $S$ nonempty ($|S| \geq 1$) and $2|S| \leq |V| < 2$, giving $|S| < 1$, a contradiction. Hence no eligible subset exists, the defining set is empty, and $h(G) = \operatorname{sInf}\,\emptyset = 0 \geq 0$ by the real conditionally complete lattice convention.
\end{proof}

\begin{theorem}[Axiom: Cheeger Inequalities]
\label{thm:cheegerInequalities}
\lean{cheegerInequalities}
\leanok
\uses{def:GraphExpansion.secondLargestEigenvalue, def:CheegerConstant.cheegerConstant}

\textbf{\color{red}[UNPROVEN AXIOM]}

Let $G$ be a connected $s$-regular finite simple graph on vertex set $V$ with $|V| \geq 2$. Let $\lambda_2 = \lambda_2(G)$ denote the second-largest adjacency eigenvalue of $G$ (Definition \ref{def:GraphExpansion.secondLargestEigenvalue}), and let $h(G)$ be the Cheeger constant of $G$ (Definition \ref{def:CheegerConstant.cheegerConstant}). Then the following \emph{Cheeger inequalities} hold:
\[
  \frac{s - \lambda_2}{2} \;\leq\; h(G) \;\leq\; \sqrt{2s(s - \lambda_2)}.
\]
The lower bound follows from the variational characterization of $\lambda_2$ via the Courant--Fischer min-max theorem, and the upper bound follows from a sweep-cut argument applied to the second eigenvector.

\textit{Note: This is stated as an axiom because the proof requires spectral theory infrastructure (the Courant--Fischer min-max theorem and sweep-cut analysis) that is not available in Mathlib.}
\end{theorem}

%--- Lem_1: RelativeCheeger ---
\chapter{Lem 1: Relative Cheeger Inequality}

\begin{lemma}[Cardinality at Least Two from Subset Bound]
\label{lem:RelativeCheeger.card_ge_two_of_subset}
\lean{RelativeCheeger.card_ge_two_of_subset}
\leanok
\uses{def:GraphExpansion.secondLargestEigenvalue}
Let $V$ be a finite type, $S \subseteq V$ a finset, and $\alpha \in \mathbb{R}$ with $\alpha < 1$.
If $0 < |S|$ and $|S| < \alpha \cdot |V|$, then $|V| \geq 2$.
\end{lemma}

\begin{proof}
\leanok

We argue by contradiction. Suppose $|V| < 2$. Since $|V|$ is a natural number, $|V| \leq 1$.
Since $S \subseteq V$, we have $|S| \leq |V| \leq 1$, so combined with $0 < |S|$ we get $|S| = 1$
and thus $|V| = 1$. Substituting into $|S| < \alpha \cdot |V|$ yields $1 < \alpha \cdot 1 = \alpha$,
but this contradicts $\alpha < 1$. Hence $|V| \geq 2$.
\end{proof}

\begin{theorem}[Axiom: Spectral Laplacian Bound]
\label{thm:RelativeCheeger.spectralLaplacianBound}
\lean{RelativeCheeger.spectralLaplacianBound}
\leanok
\uses{def:GraphExpansion.secondLargestEigenvalue, def:CheegerConstant.edgeBoundary}

\textbf{\color{red}[UNPROVEN AXIOM]}

Let $G = (V, E)$ be a finite, connected, $s$-regular simple graph with $|V| \geq 2$, and let
$\lambda_2$ denote its second-largest adjacency eigenvalue. For any nonempty proper subset
$S \subsetneq V$ (i.e., $0 < |S| < |V|$), the edge boundary $\delta S$ satisfies:
\[
  (s - \lambda_2) \cdot |S| \cdot \left(1 - \frac{|S|}{|V|}\right) \leq |\delta S|.
\]

\textit{Note: This is stated as an axiom because the proof requires eigenvalue decomposition and
the Courant--Fischer variational characterization of $\lambda_2$, which are not currently available
in Mathlib. The proof strategy is: (Step 1) Set $f = \mathbf{1}_S$, decompose $f = \bar{f}\cdot\mathbf{1} + g$
where $\bar{f} = |S|/|V|$; (Step 2) By the spectral theorem, $g^\top(sI - M)g \geq (s - \lambda_2)\,g^\top g$;
(Step 3) Compute $g^\top g = |S|(1 - |S|/|V|)$ and $f^\top(sI-M)f = |\delta S|$;
(Step 4) Since $(sI-M)\cdot\mathbf{1} = 0$, we have $g^\top(sI-M)g = f^\top(sI-M)f$, completing the bound.
This follows the same convention as \texttt{cheegerInequalities} in the Cheeger constant file.}
\end{theorem}

\begin{lemma}[Satisfiability Witness for Spectral Laplacian Bound]
\label{lem:RelativeCheeger.spectralLaplacianBound_satisfiable}
\lean{RelativeCheeger.spectralLaplacianBound_satisfiable}
\leanok
\uses{thm:RelativeCheeger.spectralLaplacianBound}
Let $G = (V, E)$ be a finite, connected simple graph with $|V| \geq 2$. Then there exists a
finset $S \subseteq V$ with $0 < |S|$ and $|S| < |V|$, confirming that the hypotheses of
the spectral Laplacian bound axiom are satisfiable.
\end{lemma}

\begin{proof}
\leanok

Since $|V| \geq 2 > 0$, the type $V$ is nonempty; obtain any vertex $v \in V$. Take $S = \{v\}$.
Then $|S| = 1 > 0$, and $|S| = 1 < |V|$ because $|V| \geq 2$. This provides the required witness.
\end{proof}

\begin{lemma}[Card Ratio Less Than Alpha]
\label{lem:RelativeCheeger.card_div_lt_alpha}
\lean{RelativeCheeger.card_div_lt_alpha}
\leanok

Let $S \subseteq V$ be a finset, $\alpha \in \mathbb{R}$ with $\alpha > 0$, and $|V| > 0$.
If $|S| < \alpha \cdot |V|$, then $|S| / |V| < \alpha$.
\end{lemma}

\begin{proof}
\leanok

Since $|V| > 0$, we have $(|V| : \mathbb{R}) > 0$. Using the equivalence
$|S|/|V| < \alpha \Leftrightarrow |S| < \alpha \cdot |V|$ (valid for positive $|V|$), the
result follows directly from the hypothesis $|S| < \alpha \cdot |V|$ by dividing both sides
by $|V|$.
\end{proof}

\begin{lemma}[Card Strictly Less Than Total From Alpha Bound]
\label{lem:RelativeCheeger.card_lt_of_alpha}
\lean{RelativeCheeger.card_lt_of_alpha}
\leanok

Let $S \subseteq V$ be a finset, $\alpha \in \mathbb{R}$ with $\alpha < 1$.
If $(|S| : \mathbb{R}) < \alpha \cdot |V|$, then $|S| < |V|$ (as natural numbers).
\end{lemma}

\begin{proof}
\leanok

We show $(|S| : \mathbb{R}) < |V|$ and then cast to natural numbers. We have:
\[
  (|S| : \mathbb{R}) < \alpha \cdot |V| < 1 \cdot |V| = |V|,
\]
where the second inequality holds because $\alpha < 1$ and $|V| > 0$ (the latter follows
since $|S| < \alpha \cdot |V|$ with $\alpha > 0$ and $|S| \geq 0$ implies $|V| > 0$; if
$|V| = 0$ then the right-hand side $\alpha \cdot 0 = 0$ and $|S| \geq 0$, which combined with
$|S| < 0$ gives a contradiction). Casting back to natural numbers gives $|S| < |V|$.
\end{proof}

\begin{theorem}[Relative Cheeger Inequality]
\label{thm:RelativeCheeger.relativeCheeger}
\lean{RelativeCheeger.relativeCheeger}
\leanok
\uses{def:CheegerConstant.edgeBoundary, def:GraphExpansion.secondLargestEigenvalue}
Let $G = (V, E)$ be a finite, connected, $s$-regular simple graph, and let $\lambda_2$ denote
its second-largest adjacency eigenvalue (well-defined since $|V| \geq 2$). Let $\alpha \in (0,1)$
and let $S \subseteq V$ with $0 < |S|$ and $|S| < \alpha \cdot |V|$. Then:
\[
  (1 - \alpha)(s - \lambda_2) \leq \frac{|\delta S|}{|S|}.
\]
\end{theorem}

\begin{proof}
\leanok
\uses{thm:RelativeCheeger.spectralLaplacianBound, lem:RelativeCheeger.card_ge_two_of_subset, lem:RelativeCheeger.card_lt_of_alpha, lem:RelativeCheeger.card_div_lt_alpha}
Let $h_{\mathrm{card}} : |V| \geq 2$ be the witness obtained from
$|S| < \alpha \cdot |V|$, $0 < |S|$, $\alpha < 1$ via the lemma
\texttt{card\_ge\_two\_of\_subset}, so that $\lambda_2$ is well-defined.

We establish the following auxiliary facts:
\begin{itemize}
  \item $(|S| : \mathbb{R}) > 0$ by casting $0 < |S|$.
  \item $(|V| : \mathbb{R}) > 0$ since $|V| \geq 2$.
  \item $|S| < |V|$ (as natural numbers) by \texttt{card\_lt\_of\_alpha}.
  \item $|S|/|V| < \alpha$ by \texttt{card\_div\_lt\_alpha}.
  \item $1 - \alpha < 1 - |S|/|V|$ by linear arithmetic from $|S|/|V| < \alpha$.
  \item $1 - \alpha > 0$ by linear arithmetic from $\alpha < 1$.
\end{itemize}

We invoke the spectral Laplacian bound axiom $(\texttt{spectralLaplacianBound})$ to obtain:
\[
  (s - \lambda_2) \cdot |S| \cdot \left(1 - \frac{|S|}{|V|}\right) \leq |\delta S|. \tag{$*$}
\]

We also note $|\delta S| \geq 0$ and $1 - |S|/|V| > 0$ (since $|S| < |V|$ and $|V| > 0$).

\medskip
We now split on the sign of $s - \lambda_2$.

\textbf{Case 1: $s - \lambda_2 \geq 0$.}
We multiply out and rewrite the goal $(1 - \alpha)(s - \lambda_2) \leq |\delta S|/|S|$
equivalently (multiplying both sides by $|S| > 0$) as
$(1 - \alpha)(s - \lambda_2) \cdot |S| \leq |\delta S|$. We compute:
\begin{align*}
  (1 - \alpha)(s - \lambda_2) \cdot |S|
  &= (s - \lambda_2) \cdot |S| \cdot (1 - \alpha) \quad \text{(by ring)} \\
  &\leq (s - \lambda_2) \cdot |S| \cdot \left(1 - \frac{|S|}{|V|}\right)
    \quad \text{(since $1 - \alpha \leq 1 - |S|/|V|$ and $(s-\lambda_2)\cdot|S| \geq 0$)} \\
  &\leq |\delta S| \quad \text{(by $(*)$)}.
\end{align*}

\textbf{Case 2: $s - \lambda_2 < 0$.}
Then $(1-\alpha)(s - \lambda_2) < 0$ (product of positive $1-\alpha > 0$ and negative $s - \lambda_2 < 0$). We have:
\[
  (1 - \alpha)(s - \lambda_2) \leq 0 \leq \frac{|\delta S|}{|S|},
\]
where the second inequality holds since $|\delta S| \geq 0$ and $|S| > 0$.

In both cases the inequality is established.
\end{proof}

%--- Thm_5: AlonBoppanaBound ---
\chapter{Thm 5: Alon--Boppana Bound}

\begin{lemma}[Diameter Implies At Least Two Vertices]
\label{lem:AlonBoppana.card_ge_two_of_diam_ge_two}
\lean{AlonBoppana.card_ge_two_of_diam_ge_two}
\leanok

Let $G$ be a connected simple graph on a finite vertex type $V$ with $\operatorname{diam}(G) \geq 2$. Then $|V| \geq 2$.
\end{lemma}

\begin{proof}
\leanok

By contradiction, assume $|V| \leq 1$. Then $V$ is a subsingleton, so every simple graph on $V$ has extended diameter $\operatorname{ediam}(G) = 0$ (by \texttt{SimpleGraph.ediam\_eq\_zero\_of\_subsingleton}). But then $\operatorname{diam}(G) = 0$, contradicting $\operatorname{diam}(G) \geq 2$.
\end{proof}

\begin{lemma}[Regular Graph Has At Least Two Vertices]
\label{lem:AlonBoppana.card_ge_two_of_regular}
\lean{AlonBoppana.card_ge_two_of_regular}
\leanok

Let $G$ be a connected simple graph on a finite vertex type $V$ that is $s$-regular for some $s \geq 3$. Then $|V| \geq 2$.
\end{lemma}

\begin{proof}
\leanok

By contradiction, assume $|V| \leq 1$. Then $V$ is a subsingleton. Since $G$ is connected, $V$ is nonempty; obtain a vertex $v$. Since $V$ is a subsingleton, the only element is $v$, so the neighborhood of $v$ in the simple graph is empty: for any neighbor $w$ of $v$, adjacency implies $v \neq w$, yet $v = w$ in a subsingleton, a contradiction. Hence $\deg(v) = 0$. But by $s$-regularity, $\deg(v) = s \geq 3$, a contradiction. We conclude by integer arithmetic.
\end{proof}

\begin{theorem}[Axiom: Alon--Boppana Bound (Quantitative Form, Nilli 1991)]
\label{thm:AlonBoppana.alonBoppanaBound}
\lean{AlonBoppana.alonBoppanaBound}
\leanok
\uses{def:GraphExpansion.secondLargestEigenvalue, lem:AlonBoppana.card_ge_two_of_diam_ge_two}

\textbf{\color{red}[UNPROVEN AXIOM]}

Let $G$ be a finite, connected $s$-regular simple graph on a finite vertex type $V$ with $s \geq 3$ and diameter $D = \operatorname{diam}(G) \geq 2$. Let $\lambda_2$ denote the second-largest eigenvalue of the adjacency matrix of $G$. Then
\[
\lambda_2 \geq 2\sqrt{s-1} - \frac{2\sqrt{s-1} - 1}{\lfloor D/2 \rfloor}.
\]

\textit{Note: This is stated as an axiom because the full proof (using Nilli's edge-based Laplacian approach, constructing test vectors supported on BFS layers and bounding the Rayleigh quotient) was not completed in the formalization.}
\end{theorem}

\begin{theorem}[Satisfiability of Alon--Boppana Hypotheses]
\label{thm:AlonBoppana.alonBoppanaBound_satisfiable}
\lean{AlonBoppana.alonBoppanaBound_satisfiable}
\leanok
\uses{thm:AlonBoppana.alonBoppanaBound}
There exist a finite vertex type $V$, a simple graph $G$ on $V$, and $s \geq 3$ such that $G$ is connected, $s$-regular, and $\operatorname{diam}(G) \geq 2$. In particular, the hypotheses of \texttt{alonBoppanaBound} are satisfiable.
\end{theorem}

\begin{proof}
\leanok
\uses{thm:AlonBoppana.alonBoppanaBound}
We exhibit the complete bipartite graph $K_{3,3}$ on vertex set $\operatorname{Fin}(6)$, with parts $\{0,1,2\}$ and $\{3,4,5\}$, where $i$ and $j$ are adjacent iff $i < 3 \leq j$ or $j < 3 \leq i$.

\textbf{Connected:} For any two vertices $u, v$, we find a path by cases. If both $u, v \in \{0,1,2\}$, use the path $u \to 3 \to v$. If both $u, v \in \{3,4,5\}$, use the path $u \to 0 \to v$. If $u \in \{0,1,2\}$ and $v \in \{3,4,5\}$ (or vice versa), then $u$ and $v$ are directly adjacent.

\textbf{3-regular:} For each vertex $v \in \operatorname{Fin}(6)$, $\deg(v) = 3$ is verified by computation (\texttt{native\_decide}).

\textbf{Diameter $\geq 2$:} Let $v_0 = 0$, $v_1 = 1$, $v_3 = 3$. The walk $v_0 \to v_3 \to v_1$ has length 2, so $\operatorname{dist}(v_0, v_1) \leq 2$. Since $v_0 \neq v_1$ (by computation) and $v_0, v_1$ are not adjacent (both in $\{0,1,2\}$, by computation), we have $\operatorname{dist}(v_0, v_1) \neq 0$ and $\neq 1$, hence $\operatorname{dist}(v_0, v_1) = 2$ by integer arithmetic. Since $K_{3,3}$ is connected its extended diameter is finite, and $\operatorname{diam}(K_{3,3}) \geq \operatorname{dist}(v_0, v_1) = 2$.
\end{proof}

\begin{theorem}[Moore Bound]
\label{thm:AlonBoppana.moore_bound}
\lean{AlonBoppana.moore_bound}
\leanok
\uses{lem:AlonBoppana.card_ge_two_of_regular}
Let $G$ be a finite, connected $s$-regular simple graph on vertex type $V$ with $s \geq 3$. Let $D = \operatorname{diam}(G)$. Then
\[
\frac{|V| \cdot (s - 2)}{s} \leq (s-1)^D.
\]
\end{theorem}

\begin{proof}
\leanok
\uses{lem:AlonBoppana.card_ge_two_of_regular}
We use the real-valued bound $|V| \cdot (s-2) / s \leq (s-1)^D$ (rewriting via \texttt{Real.rpow\_natCast}). It suffices to prove the natural number inequality $|V| \cdot (s-2) \leq s \cdot (s-1)^D$ and then cast to $\mathbb{R}$.

\textbf{Geometric sum:} We establish by induction on $D$ that
\[
(s-2) \cdot \sum_{k=0}^{D-1} (s-1)^k = (s-1)^D - 1.
\]
The base case $D = 0$ is trivial. For the inductive step, expanding $\sum_{k=0}^{D} (s-1)^k = \sum_{k=0}^{D-1}(s-1)^k + (s-1)^D$ and using the induction hypothesis, with careful handling of natural number subtraction (using \texttt{zify} and \texttt{ring} after establishing nonnegativity conditions), we verify the telescoping identity.

\textbf{BFS layer decomposition:} Fix a vertex $v_0$. Define layers $L_k = \{w : \operatorname{dist}(v_0, w) = k\}$ for $k \in \{0, \ldots, D\}$. Since $G$ is connected, every vertex lies within distance $D$ of $v_0$, so
\[
|V| = \sum_{k=0}^{D} |L_k|.
\]
The disjointness of distinct layers follows because vertices have unique distances. We have $|L_0| = 1$ (only $v_0$, since $\operatorname{dist}(v_0, v_0) = 0$ and $\operatorname{dist}(v_0, w) = 0$ implies $w = v_0$ for a connected graph).

\textbf{Layer size bound:} We claim $|L_k| \leq s \cdot (s-1)^{k-1}$ for all $1 \leq k \leq D$, proved by induction on $k$:
\begin{itemize}
\item $k = 1$: The vertices at distance 1 are exactly the neighbors of $v_0$, so $|L_1| \leq \deg(v_0) = s$.
\item $k \geq 2$ (inductive step): We show $|L_{k+1}| \leq (s-1) \cdot |L_k|$. For each $w \in L_{k+1}$, let $u$ be the penultimate vertex on a shortest walk from $v_0$ to $w$; then $u \in L_k$ and $G.\text{Adj}(u, w)$. Hence $L_{k+1} \subseteq \bigcup_{u \in L_k} (\mathcal{N}(u) \cap L_{k+1})$.

For each $u \in L_k$, we bound $|\mathcal{N}(u) \cap L_{k+1}| \leq s - 1$. Indeed, $u$ has degree $s$. Since $k \geq 1$, there exists a predecessor vertex $\operatorname{pred}(u)$ at distance $k-1$ from $v_0$ adjacent to $u$ (the penultimate vertex of a shortest walk to $u$). This predecessor is not in $L_{k+1}$, so it lies in $\mathcal{N}(u) \setminus (\mathcal{N}(u) \cap L_{k+1})$, giving $|\mathcal{N}(u) \cap L_{k+1}| < |\mathcal{N}(u)| = s$, i.e., $|\mathcal{N}(u) \cap L_{k+1}| \leq s - 1$.

By $\texttt{Finset.card\_biUnion\_le}$ and summing over $L_k$:
\[
|L_{k+1}| \leq \sum_{u \in L_k} |\mathcal{N}(u) \cap L_{k+1}| \leq |L_k| \cdot (s-1).
\]
Combined with the inductive hypothesis $|L_k| \leq s \cdot (s-1)^{k-1}$, we get $|L_{k+1}| \leq s \cdot (s-1)^k$.
\end{itemize}

\textbf{Combining:} By the Finset sum reindexing (\texttt{Finset.sum\_range\_succ'}),
\[
|V| = 1 + \sum_{k=0}^{D-1} |L_{k+1}| \leq 1 + s \cdot \sum_{k=0}^{D-1} (s-1)^k.
\]
Using the geometric sum identity and natural number arithmetic:
\begin{align*}
|V| \cdot (s-2) &\leq \left(1 + s \cdot \sum_{k=0}^{D-1}(s-1)^k\right)(s-2) \\
&= (s-2) + s \cdot (s-2) \cdot \sum_{k=0}^{D-1}(s-1)^k \\
&= (s-2) + s \cdot ((s-1)^D - 1) \leq s \cdot (s-1)^D,
\end{align*}
where the last step follows since $(s-2) \leq s \cdot 1 \leq s \cdot (s-1)^D$, verified by \texttt{nlinarith} after \texttt{zify}.
\end{proof}

\begin{theorem}[Diameter Lower Bound via Moore Bound]
\label{thm:AlonBoppana.diam_lower_bound}
\lean{AlonBoppana.diam_lower_bound}
\leanok
\uses{thm:AlonBoppana.moore_bound}
Let $G$ be a finite, connected $s$-regular simple graph on vertex type $V$ with $s \geq 3$. Then
\[
\log_{s-1}\!\left(\frac{|V| \cdot (s-2)}{s}\right) \leq \operatorname{diam}(G).
\]
\end{theorem}

\begin{proof}
\leanok
\uses{thm:AlonBoppana.moore_bound}
We have $1 < s - 1$ since $s \geq 3$. We consider two cases.

\textbf{Case 1:} $|V| \cdot (s-2) / s \leq 0$. Then $\log_{s-1}(|V| \cdot (s-2)/s) \leq 0$ (since the argument is $\leq 1$ after noting it is $\geq 0$ by nonnegativity of $|V|$, $s-2$, and $s$). The right-hand side $\operatorname{diam}(G)$ is a natural number, hence $\geq 0$. The inequality follows.

\textbf{Case 2:} $|V| \cdot (s-2)/s > 0$. By \texttt{Real.logb\_le\_iff\_le\_rpow} (valid since $s-1 > 1$ and the argument is positive), the inequality $\log_{s-1}(|V|\cdot(s-2)/s) \leq D$ is equivalent to $|V|\cdot(s-2)/s \leq (s-1)^D$. This follows directly from the Moore bound (\texttt{moore\_bound}).
\end{proof}

\begin{lemma}[Satisfiability of Diameter Lower Bound Hypotheses]
\label{lem:AlonBoppana.diam_lower_bound_satisfiable}
\lean{AlonBoppana.diam_lower_bound_satisfiable}
\leanok
\uses{thm:AlonBoppana.diam_lower_bound}
There exist a finite vertex type $V$, a simple graph $G$ on $V$, and $s \geq 3$ such that $G$ is connected and $s$-regular. In particular, the hypotheses of \texttt{diam\_lower\_bound} are satisfiable.
\end{lemma}

\begin{proof}
\leanok
\uses{thm:AlonBoppana.diam_lower_bound}
We use $V = \operatorname{Fin}(4)$ and $G = \top$ (the complete graph on 4 vertices). The complete graph on $\operatorname{Fin}(4)$ is connected by \texttt{SimpleGraph.connected\_top} and is 3-regular (each vertex has degree $4 - 1 = 3$), verified via \texttt{SimpleGraph.IsRegularOfDegree.top} and simplification with $|\operatorname{Fin}(4)| = 4$.
\end{proof}

\begin{theorem}[Axiom: Asymptotic Alon--Boppana Bound]
\label{thm:AlonBoppana.alonBoppana_liminf}
\lean{AlonBoppana.alonBoppana_liminf}
\leanok
\uses{def:GraphExpansion.secondLargestEigenvalue, thm:AlonBoppana.alonBoppanaBound, thm:AlonBoppana.diam_lower_bound}

\textbf{\color{red}[UNPROVEN AXIOM]}

Let $s \geq 3$, let $n : \mathbb{N} \to \mathbb{N}$, and let $G_i$ be a connected $s$-regular simple graph on $\operatorname{Fin}(n(i))$ for each $i$, with $|\operatorname{Fin}(n(i))| \geq 2$ for all $i$. If $n(i) \to \infty$, then
\[
2\sqrt{s-1} \leq \liminf_{i \to \infty} \lambda_2(G_i),
\]
where $\lambda_2(G_i)$ denotes the second-largest eigenvalue of the adjacency matrix of $G_i$.

\textit{Note: This follows conceptually from \texttt{alonBoppanaBound} and \texttt{diam\_lower\_bound}: since $n(i) \to \infty$, the Moore bound gives $\operatorname{diam}(G_i) \to \infty$, so for large $i$ the correction term $(2\sqrt{s-1} - 1)/\lfloor D_i/2 \rfloor \to 0$ and the bound approaches $2\sqrt{s-1}$. It is stated as an axiom because connecting \texttt{SimpleGraph.diam} growth with \texttt{Filter.liminf} through logarithmic estimates requires infrastructure not directly available for \texttt{SimpleGraph} in Mathlib.}
\end{theorem}

\begin{theorem}[Gershgorin Eigenvalue Bound for Regular Graphs]
\label{thm:AlonBoppana.adjEigenvalue_abs_le}
\lean{AlonBoppana.adjEigenvalue_abs_le}
\leanok
\uses{def:GraphExpansion.adjEigenvalues, lem:GraphExpansion.adjMatrix_isHermitian}
Let $G$ be an $s$-regular simple graph on a finite vertex type $W$. For each $i \in \operatorname{Fin}(|W|)$,
\[
|\lambda_i(G)| \leq s,
\]
where $\lambda_i(G)$ denotes the $i$-th adjacency eigenvalue of $G$.
\end{theorem}

\begin{proof}
\leanok
\uses{def:GraphExpansion.adjEigenvalues, lem:GraphExpansion.adjMatrix_isHermitian}
\textbf{Step 1.} We show the eigenvalue $\lambda_i = (\texttt{adjMatrix\_isHermitian}\ G).\texttt{eigenvalues}_0\ i$ lies in the $\mathbb{R}$-spectrum of the adjacency matrix $A = G.\texttt{adjMatrix}\ \mathbb{R}$, using \texttt{IsHermitian.eigenvalues\_mem\_spectrum\_real}.

\textbf{Step 2.} Since the spectrum of $A$ as a matrix equals the spectrum of the corresponding linear endomorphism $\texttt{toLin}'(A)$ (via \texttt{Matrix.spectrum\_toLin'}), we obtain a module eigenvalue via \texttt{Module.End.HasEigenvalue}.

\textbf{Step 3.} We apply the \textbf{Gershgorin circle theorem} (\texttt{eigenvalue\_mem\_ball}): every eigenvalue lies in a closed disk centered at some diagonal entry $A_{kk}$ with radius $\sum_{j \neq k} |A_{kj}|$ for some row index $k$.

\textbf{Step 4.} The diagonal entry satisfies $A_{kk} = 0$ since $G$ has no self-loops (\texttt{SimpleGraph.adjMatrix\_apply} with \texttt{G.loopless}).

\textbf{Step 5.} We compute the Gershgorin radius of row $k$: since $|A_{kj}| = 1$ if $G.\text{Adj}(k, j)$ and $0$ otherwise, and since $G$ has no self-loops,
\[
\sum_{j \neq k} |A_{kj}| = |\{j : G.\text{Adj}(k, j)\}| = \deg(k) = s.
\]
This uses the identification of the filtered set with the neighbor finset and the regularity hypothesis \texttt{hreg k}.

\textbf{Step 6.} From the Gershgorin bound with center $0$ and radius $s$, we obtain $|\lambda_i| \leq s$ via \texttt{Metric.mem\_closedBall} and \texttt{Real.norm\_eq\_abs}.
\end{proof}

\begin{lemma}[Satisfiability of Eigenvalue Bound Hypotheses]
\label{lem:AlonBoppana.adjEigenvalue_abs_le_satisfiable}
\lean{AlonBoppana.adjEigenvalue_abs_le_satisfiable}
\leanok
\uses{thm:AlonBoppana.adjEigenvalue_abs_le}
There exist a finite vertex type $W$, a simple graph $G$ on $W$, and $s \geq 0$ such that $G$ is $s$-regular. In particular, the hypotheses of \texttt{adjEigenvalue\_abs\_le} are satisfiable.
\end{lemma}

\begin{proof}
\leanok
\uses{thm:AlonBoppana.adjEigenvalue_abs_le}
We use $W = \operatorname{Fin}(4)$, $G = \top$ (the complete graph), and $s = 3$, witnessing $G.\text{IsRegularOfDegree}(3)$ via \texttt{SimpleGraph.IsRegularOfDegree.top}.
\end{proof}

\begin{lemma}[Second-Largest Eigenvalue Lower Bound]
\label{lem:AlonBoppana.secondLargestEigenvalue_ge_neg}
\lean{AlonBoppana.secondLargestEigenvalue_ge_neg}
\leanok
\uses{def:GraphExpansion.secondLargestEigenvalue, thm:AlonBoppana.adjEigenvalue_abs_le}
Let $G$ be an $s$-regular simple graph on a finite vertex type $W$ with $|W| \geq 2$. Then
\[
-s \leq \lambda_2(G),
\]
where $\lambda_2(G)$ denotes the second-largest eigenvalue of the adjacency matrix of $G$.
\end{lemma}

\begin{proof}
\leanok
\uses{def:GraphExpansion.secondLargestEigenvalue, thm:AlonBoppana.adjEigenvalue_abs_le}
Apply \texttt{adjEigenvalue\_abs\_le} at index $i = \langle 1, \cdot \rangle \in \operatorname{Fin}(|W|)$ (valid since $|W| \geq 2$) to obtain $|\lambda_2(G)| \leq s$. By the absolute value bound (\texttt{abs\_le}), this gives both $\lambda_2(G) \leq s$ and $-s \leq \lambda_2(G)$; we use the latter.
\end{proof}

\begin{theorem}[Ramanujan Bound Is Tight]
\label{thm:AlonBoppana.ramanujan_bound_optimal}
\lean{AlonBoppana.ramanujan_bound_optimal}
\leanok
\uses{def:GraphExpansion.secondLargestEigenvalue, thm:AlonBoppana.alonBoppana_liminf, lem:AlonBoppana.secondLargestEigenvalue_ge_neg}
Let $s \geq 3$, let $n : \mathbb{N} \to \mathbb{N}$, and let $G_i$ be a connected $s$-regular simple graph on $\operatorname{Fin}(n(i))$ for each $i$, with $|\operatorname{Fin}(n(i))| \geq 2$ and $n(i) \to \infty$. Suppose that there exists $\varepsilon > 0$ such that
\[
\lambda_2(G_i) \leq s - \varepsilon \quad \text{for all } i.
\]
Then $\varepsilon \leq s - 2\sqrt{s-1}$. In particular, any uniform spectral gap for an expander family cannot exceed $s - 2\sqrt{s-1}$, so the Ramanujan bound $\lambda_2 \leq 2\sqrt{s-1}$ is best possible.
\end{theorem}

\begin{proof}
\leanok
\uses{thm:AlonBoppana.alonBoppana_liminf, lem:AlonBoppana.secondLargestEigenvalue_ge_neg}
\textbf{Step 1.} By the asymptotic Alon--Boppana bound (\texttt{alonBoppana\_liminf}), since $n(i) \to \infty$ and each $G_i$ is connected and $s$-regular with $|V_i| \geq 2$:
\[
2\sqrt{s-1} \leq \liminf_{i \to \infty} \lambda_2(G_i).
\]

\textbf{Step 2.} The uniform bound $\lambda_2(G_i) \leq s - \varepsilon$ holds for all $i$. By \texttt{Filter.Frequently.of\_forall}, this event occurs frequently (in fact, always), so:
\[
\liminf_{i \to \infty} \lambda_2(G_i) \leq s - \varepsilon,
\]
using \texttt{Filter.liminf\_le\_of\_frequently\_le}. The required lower boundedness of $(\lambda_2(G_i))_i$ is provided by $\lambda_2(G_i) \geq -s$ (\texttt{secondLargestEigenvalue\_ge\_neg}), establishing \texttt{Filter.IsBoundedUnder} via \texttt{Filter.Eventually.of\_forall}.

\textbf{Step 3.} Combining Steps 1 and 2:
\[
2\sqrt{s-1} \leq \liminf_{i\to\infty} \lambda_2(G_i) \leq s - \varepsilon,
\]
hence $\varepsilon \leq s - 2\sqrt{s-1}$ by linear arithmetic (\texttt{linarith}).
\end{proof}

%--- Thm_6: AlonChung ---
\chapter{Thm 6: Alon--Chung Bound}

\begin{definition}[Induced Edge Set]
\label{def:AlonChung.inducedEdges}
\lean{AlonChung.inducedEdges}
\leanok

Let $G = (V, E)$ be a simple graph and $S \subseteq V$ a finite subset. The \emph{induced edge set} of $S$ is
\[
  X(S)_1 \;=\; \{ e \in E \mid \forall\, v \in e,\; v \in S \} \;=\; E \cap S^{(2)},
\]
where $S^{(2)}$ denotes the set of unordered pairs of elements of $S$.
\end{definition}

\begin{lemma}[Membership in Induced Edges]
\label{lem:AlonChung.mem_inducedEdges}
\lean{AlonChung.mem_inducedEdges}
\leanok
\uses{def:AlonChung.inducedEdges}
For a simple graph $G$ on $V$, a subset $S \subseteq V$, and an unordered pair $e \in \operatorname{Sym}^2(V)$,
\[
  e \in X(S)_1 \;\Longleftrightarrow\; e \in E(G) \;\land\; \forall\, v \in e,\; v \in S.
\]
\end{lemma}

\begin{proof}
\leanok
\uses{def:AlonChung.inducedEdges}
This follows by unfolding the definition of $X(S)_1 = E(G) \cap S^{(2)}$ and applying the characterisation of membership in $S^{(2)}$: by simplification, $e \in E(G) \cap S^{(2)}$ if and only if $e \in E(G)$ and every vertex of $e$ lies in $S$.
\end{proof}

\begin{lemma}[Nonemptiness of Induced Edges]
\label{lem:AlonChung.inducedEdges_nonempty}
\lean{AlonChung.inducedEdges_nonempty}
\leanok
\uses{def:AlonChung.inducedEdges, lem:AlonChung.mem_inducedEdges}
If there exists an edge $e \in E(G)$ such that every vertex of $e$ lies in $S$, then $X(S)_1$ is nonempty.
\end{lemma}

\begin{proof}
\leanok
\uses{def:AlonChung.inducedEdges, lem:AlonChung.mem_inducedEdges}
From the hypothesis, obtain $e \in E(G)$ with $\forall v \in e,\, v \in S$. Then $e \in X(S)_1$ by the membership characterisation (Lemma~\ref{lem:AlonChung.mem_inducedEdges}), so $X(S)_1$ is nonempty.
\end{proof}

\begin{definition}[Indicator Vector]
\label{def:AlonChung.indicatorVec}
\lean{AlonChung.indicatorVec}
\leanok

For a finite vertex set $V$ and $S \subseteq V$, the \emph{indicator vector} $\mathbf{1}_S : V \to \mathbb{R}$ is defined by
\[
  \mathbf{1}_S(v) \;=\; \begin{cases} 1 & \text{if } v \in S, \\ 0 & \text{if } v \notin S. \end{cases}
\]
\end{definition}

\begin{lemma}[Quadratic Form Equals Twice Induced Edges]
\label{lem:AlonChung.quadratic_form_eq_twice_inducedEdges}
\lean{AlonChung.quadratic_form_eq_twice_inducedEdges}
\leanok
\uses{def:AlonChung.inducedEdges, def:AlonChung.indicatorVec, lem:GraphExpansion.adjMatrix_isHermitian}
Let $G$ be a simple graph on $V$ and $M = A(G)$ its adjacency matrix over $\mathbb{R}$. For any $S \subseteq V$,
\[
  \mathbf{1}_S^\top M\, \mathbf{1}_S \;=\; 2\,|X(S)_1|.
\]
\end{lemma}

\begin{proof}
\leanok
\uses{def:AlonChung.inducedEdges, def:AlonChung.indicatorVec}
We unfold the dot product and matrix-vector multiplication using the definitions of $\mathbf{1}_S$ and $M$:
\[
  \mathbf{1}_S^\top M \mathbf{1}_S = \sum_{v \in V} \mathbf{1}_S(v)\sum_{w \in V} M_{vw}\,\mathbf{1}_S(w).
\]
\textbf{Step 1.} Since $M_{vw} = \mathbf{1}[G.Adj(v,w)]$ and $\mathbf{1}_S(v), \mathbf{1}_S(w) \in \{0,1\}$, the double sum simplifies to
\[
  \sum_{v \in V}\sum_{w \in V} \mathbf{1}[G.\operatorname{Adj}(v,w) \land v \in S \land w \in S].
\]
\textbf{Step 2.} By rewriting via the indicator function of the filtered product set, this equals the cardinality
\[
  \bigl|\{(v,w) \in V \times V \mid G.\operatorname{Adj}(v,w),\, v \in S,\, w \in S\}\bigr|.
\]
\textbf{Step 3.} We exhibit a $2$-to-$1$ map from these ordered pairs to unordered induced edges: by the symmetry of adjacency, each unordered edge $\{v,w\} \in X(S)_1$ corresponds to exactly two ordered pairs $(v,w)$ and $(w,v)$ with $v \neq w$. Concretely, for each $e = s(v,w) \in X(S)_1$ the fibre $\{(a,b) : G.\operatorname{Adj}(a,b) \land a \in S \land b \in S \land s(a,b) = e\}$ equals $\{(v,w),(w,v)\}$, which has cardinality $2$ since $v \neq w$ (graphs have no self-loops). Partitioning the filtered product by edge and applying the constant-fibre sum formula gives the total count $= 2\,|X(S)_1|$. Simplification by ring arithmetic concludes.
\end{proof}

\begin{lemma}[Twice Edge Count for Regular Graphs]
\label{lem:AlonChung.twice_card_edgeFinset_eq}
\lean{AlonChung.twice_card_edgeFinset_eq}
\leanok
\uses{def:AlonChung.indicatorVec}
If $G$ is $s$-regular on $V$, then
\[
  2\,|E(G)| \;=\; s\,|V|.
\]
\end{lemma}

\begin{proof}
\leanok

We use the handshaking lemma $\sum_{v \in V} \deg(v) = 2|E(G)|$. Since $G$ is $s$-regular, $\deg(v) = s$ for all $v$, so $\sum_{v}\deg(v) = s\,|V|$. Substituting and casting to $\mathbb{R}$ gives $2|E(G)| = s|V|$.
\end{proof}

\begin{lemma}[Spectral Decomposition of the Quadratic Form]
\label{lem:AlonChung.spectral_decomp_quadratic}
\lean{AlonChung.spectral_decomp_quadratic}
\leanok
\uses{lem:GraphExpansion.adjMatrix_isHermitian}
Let $A$ be a Hermitian matrix on $V$ with orthonormal eigenvector basis $(v_j)_{j \in V}$ and corresponding real eigenvalues $(\lambda_j)$. For any $x : V \to \mathbb{R}$,
\[
  x^\top A x \;=\; \sum_{j \in V} \lambda_j \,(x \cdot v_j)^2.
\]
\end{lemma}

\begin{proof}
\leanok
\uses{lem:GraphExpansion.adjMatrix_isHermitian}
Let $U$ be the eigenvector unitary (orthogonal matrix) with $A = U D U^\top$ (spectral theorem for real Hermitian matrices), where $D = \operatorname{diag}(\lambda_j)$.

\textbf{Step 1.} We verify $\star U = U^\top$ using that $A$ is real Hermitian.

\textbf{Step 2.} Set $y = U^\top x$. Then $Ax = U(Dy)$ by associativity.

\textbf{Step 3.} Compute $x^\top (Ax) = x^\top U (D y)$. Since $\operatorname{vecMul}(x, U) = U^\top x = y$ (by expanding definitions of vector multiplication and transposition for real matrices), we get $x^\top A x = y^\top D y$.

\textbf{Step 4.} Since $D$ is diagonal, $y^\top D y = \sum_j \lambda_j y_j^2$.

\textbf{Step 5.} Observe that $y_j = (U^\top x)_j = \sum_i x_i U_{ij} = x \cdot v_j$ where $v_j$ is the $j$-th eigenvector. Rewriting using $U_{ij} = (v_j)_i$ and collecting terms gives the result.
\end{proof}

\begin{lemma}[Parseval Identity for Eigenvector Basis]
\label{lem:AlonChung.parseval_eigenvector}
\lean{AlonChung.parseval_eigenvector}
\leanok
\uses{lem:GraphExpansion.adjMatrix_isHermitian}
Let $A$ be Hermitian on $V$ with orthonormal eigenvector basis $(v_j)$. For any $x : V \to \mathbb{R}$,
\[
  \|x\|^2 \;=\; \sum_{j \in V} (x \cdot v_j)^2.
\]
\end{lemma}

\begin{proof}
\leanok
\uses{lem:GraphExpansion.adjMatrix_isHermitian}
Let $U$ be the eigenvector unitary. We first show $\|x\|^2 = \|y\|^2$ where $y = U^\top x$. Indeed, by the orthogonality identity $UU^\top = I$ (from the unitary property and $\star U = U^\top$):
\[
  \|y\|^2 = y^\top D_1 y = (U^\top x)^\top(U^\top x) = x^\top (UU^\top) x = x^\top x = \|x\|^2,
\]
computed via the chain $\|y\|^2 = \operatorname{vecMul}(U^\top x, U^\top) \cdot x = U(U^\top x) \cdot x = (UU^\top x) \cdot x = x \cdot x$. Then $\|y\|^2 = \sum_j y_j^2$, and since $y_j = x \cdot v_j$, this gives $\|x\|^2 = \sum_j (x\cdot v_j)^2$.
\end{proof}

\begin{lemma}[Quadratic Form Bounded by Maximum Eigenvalue]
\label{lem:AlonChung.quadratic_le_bound}
\lean{AlonChung.quadratic_le_bound}
\leanok
\uses{lem:GraphExpansion.adjMatrix_isHermitian, lem:AlonChung.spectral_decomp_quadratic, lem:AlonChung.parseval_eigenvector}
Let $A$ be Hermitian on $V$ and suppose all eigenvalues satisfy $\lambda_j \leq C$. Then for any $x : V \to \mathbb{R}$,
\[
  x^\top A x \;\leq\; C\,\|x\|^2.
\]
\end{lemma}

\begin{proof}
\leanok
\uses{lem:AlonChung.spectral_decomp_quadratic, lem:AlonChung.parseval_eigenvector}
Rewrite using the spectral decomposition (Lemma~\ref{lem:AlonChung.spectral_decomp_quadratic}) and Parseval identity (Lemma~\ref{lem:AlonChung.parseval_eigenvector}):
\[
  x^\top Ax = \sum_j \lambda_j (x\cdot v_j)^2 \;\leq\; \sum_j C(x \cdot v_j)^2 = C\sum_j (x\cdot v_j)^2 = C\,\|x\|^2,
\]
where the inequality uses $\lambda_j \leq C$ and $(x\cdot v_j)^2 \geq 0$ term by term.
\end{proof}

\begin{lemma}[All Eigenvalues of a Regular Graph Are at Most the Degree]
\label{lem:AlonChung.eigenvalue_le_degree}
\lean{AlonChung.eigenvalue_le_degree}
\leanok
\uses{lem:GraphExpansion.adjMatrix_isHermitian, lem:AlonChung.parseval_eigenvector}
If $G$ is $s$-regular and $j \in V$, then $\lambda_j(A(G)) \leq s$.
\end{lemma}

\begin{proof}
\leanok
\uses{lem:GraphExpansion.adjMatrix_isHermitian}
Let $v = v_j$ be the $j$-th orthonormal eigenvector. By the Rayleigh quotient characterisation, $\lambda_j = v^\top A v$ (since $\|v\| = 1$). We bound $v^\top Av$ by the AM--GM inequality.

Since $G$ is $s$-regular, the adjacency matrix satisfies $A_{pq} = \mathbf{1}[G.\operatorname{Adj}(p,q)]$ and $\sum_q \mathbf{1}[G.\operatorname{Adj}(p,q)] = s$ for all $p$.

\textbf{Unit norm.} The orthonormality of the eigenvector basis gives $v^\top v = \sum_i v_i^2 = 1$ (after converting the Euclidean norm condition $\|v\| = 1$ via $\sqrt{\sum_i |v_i|^2} = 1$).

\textbf{AM--GM bound.} For all $a, b \in \mathbb{R}$: $ab \leq (a^2 + b^2)/2$ (since $(a-b)^2 \geq 0$). Therefore
\[
  v^\top Av = \sum_p \sum_q \mathbf{1}[G.\operatorname{Adj}(p,q)]\, v_p v_q \;\leq\; \sum_p\sum_q \mathbf{1}[G.\operatorname{Adj}(p,q)]\,\frac{v_p^2 + v_q^2}{2}.
\]
Splitting the sum: $\sum_p\sum_q \mathbf{1}[G.\operatorname{Adj}(p,q)] v_p^2/2 = (s/2)\sum_p v_p^2$ by regularity, and symmetrically for the $v_q^2$ terms (using $G.\operatorname{Adj}(p,q) \Leftrightarrow G.\operatorname{Adj}(q,p)$). Thus the bound equals $s\sum_p v_p^2 = s$. Since $v^\top v = 1$, we obtain $\lambda_j = v^\top Av \leq s$.
\end{proof}

\begin{lemma}[The Degree is an Eigenvalue of a Regular Graph]
\label{lem:AlonChung.degree_mem_eigenvalues}
\lean{AlonChung.degree_mem_eigenvalues}
\leanok
\uses{lem:GraphExpansion.adjMatrix_isHermitian}
If $G$ is $s$-regular and $|V| \geq 1$, then $s \in \operatorname{Spec}(A(G))$, i.e.\ $s$ equals some eigenvalue.
\end{lemma}

\begin{proof}
\leanok
\uses{lem:GraphExpansion.adjMatrix_isHermitian}
We show $(s : \mathbb{R})$ lies in the spectrum of $A(G)$ by proving it is not in the resolvent. Assume for contradiction that $sI - A$ is invertible. The all-ones vector $\mathbf{1} : V \to \mathbb{R}$ is in the kernel of $sI - A$: for each $v$,
\[
  ((sI - A)\mathbf{1})_v = s - \sum_w A_{vw} = s - \deg(v) = s - s = 0
\]
by regularity. Since $\mathbf{1} \neq 0$ (it evaluates to $1$ at any vertex), injectivity of $sI - A$ is contradicted. Hence $s \in \operatorname{Spec}(A(G))$, and by the spectral theorem for real Hermitian matrices this coincides with the range of the eigenvalue function.
\end{proof}

\begin{lemma}[The Largest Eigenvalue of a Regular Graph Equals the Degree]
\label{lem:AlonChung.max_eigenvalue_eq_degree}
\lean{AlonChung.max_eigenvalue_eq_degree}
\leanok
\uses{def:GraphExpansion.adjEigenvalues, lem:GraphExpansion.adjMatrix_isHermitian, lem:AlonChung.eigenvalue_le_degree, lem:AlonChung.degree_mem_eigenvalues}
If $G$ is $s$-regular and $|V| \geq 2$, then $\lambda_1(A(G)) = s$, i.e.\ the largest eigenvalue equals $s$.
\end{lemma}

\begin{proof}
\leanok
\uses{def:GraphExpansion.adjEigenvalues, lem:GraphExpansion.adjMatrix_isHermitian, lem:AlonChung.eigenvalue_le_degree, lem:AlonChung.degree_mem_eigenvalues}
Let $\lambda_1 = \operatorname{adjEigenvalues}(G)\langle 0, \cdot\rangle$ (the largest eigenvalue in descending order).

\textbf{Upper bound} $\lambda_1 \leq s$: Express $\lambda_1$ in terms of the Hermitian eigenvalue function and apply Lemma~\ref{lem:AlonChung.eigenvalue_le_degree}.

\textbf{Lower bound} $s \leq \lambda_1$: By $|V| \geq 2$ the graph is nonempty, so by Lemma~\ref{lem:AlonChung.degree_mem_eigenvalues} there exists an index $j$ with $\lambda_j = s$. Expressing $\lambda_j$ via the equivalence from $\operatorname{Fin}(|V|)$ and using the antitone property $\lambda_k \leq \lambda_1$ for all $k$ (Lemma~\ref{thm:GraphExpansion.adjEigenvalues_antitone}), we get $s = \lambda_j \leq \lambda_1$.

Combining, $\lambda_1 = s$.
\end{proof}

\begin{lemma}[Norm of the Indicator Vector]
\label{lem:AlonChung.indicatorVec_dotProduct_self}
\lean{AlonChung.indicatorVec_dotProduct_self}
\leanok
\uses{def:AlonChung.indicatorVec}
For any $S \subseteq V$,
\[
  \mathbf{1}_S \cdot \mathbf{1}_S \;=\; |S|.
\]
\end{lemma}

\begin{proof}
\leanok
\uses{def:AlonChung.indicatorVec}
Expanding the dot product and using that $\mathbf{1}_S(v)^2 = \mathbf{1}_S(v)$ (since values are $0$ or $1$), we obtain $\sum_{v \in V} \mathbf{1}_S(v) = |S \cap V| = |S|$ by the boolean sum formula.
\end{proof}

\begin{lemma}[Adjacency Action on Constant Vectors]
\label{lem:AlonChung.regular_mulVec_const}
\lean{AlonChung.regular_mulVec_const}
\leanok
\uses{lem:GraphExpansion.adjMatrix_isHermitian}
If $G$ is $s$-regular and $c \in \mathbb{R}$, then $A(G)\,\mathbf{c} = (sc)\,\mathbf{1}$, i.e.\ the constant vector $\mathbf{c} : V \to \mathbb{R}$ with value $c$ is an eigenvector of $A(G)$ with eigenvalue $s$.
\end{lemma}

\begin{proof}
\leanok

For each $v \in V$: $(A(G)\mathbf{c})_v = \sum_w A_{vw} c = c \sum_w \mathbf{1}[G.\operatorname{Adj}(v,w)] = c\cdot s = sc$ by regularity. Hence $A(G)\mathbf{c} = sc\,\mathbf{1}$.
\end{proof}

\begin{lemma}[Quadratic Form Decomposition]
\label{lem:AlonChung.quadratic_decomp}
\lean{AlonChung.quadratic_decomp}
\leanok
\uses{lem:GraphExpansion.adjMatrix_isHermitian, lem:AlonChung.regular_mulVec_const}
Let $G$ be $s$-regular. Suppose $x = \gamma\mathbf{1} + z$ where $\sum_v z_v = 0$. Then
\[
  x^\top A(G)\,x \;=\; s\gamma^2|V| + z^\top A(G)\,z.
\]
\end{lemma}

\begin{proof}
\leanok
\uses{lem:AlonChung.regular_mulVec_const}
Since $x_v = \gamma + z_v$, we compute $A(G)x$ using linearity and Lemma~\ref{lem:AlonChung.regular_mulVec_const}: $(A(G)x)_v = s\gamma + (A(G)z)_v$. Expanding the quadratic form:
\[
  x^\top A(G)x = \sum_v (\gamma + z_v)(s\gamma + (Az)_v) = \sum_v \bigl[s\gamma^2 + \gamma(Az)_v + s\gamma z_v + z_v(Az)_v\bigr].
\]
Summing term by term:
\begin{itemize}
\item $\sum_v s\gamma^2 = s\gamma^2|V|$.
\item $\sum_v \gamma(Az)_v = \gamma\sum_v(Az)_v = \gamma\cdot s\sum_v z_v = 0$, since $\sum_v (Az)_v = \sum_v\sum_w A_{vw}z_w = \sum_w z_w\sum_v A_{vw} = s\sum_w z_w = 0$ using symmetry of adjacency and regularity, together with $\sum_v z_v = 0$.
\item $\sum_v s\gamma z_v = s\gamma\sum_v z_v = 0$.
\item $\sum_v z_v(Az)_v = z^\top Az$.
\end{itemize}
Therefore $x^\top Ax = s\gamma^2|V| + z^\top Az$.
\end{proof}

\begin{theorem}[Axiom: Regularity Implies Eigenvalue-$s$ Eigenvectors Are Constant]
\label{thm:AlonChung.regular_eigenvector_s_constant}
\lean{AlonChung.regular_eigenvector_s_constant}
\leanok
\uses{def:GraphExpansion.adjEigenvalues, lem:GraphExpansion.adjMatrix_isHermitian}

\textbf{\color{red}[UNPROVEN AXIOM]}

Let $G$ be an $s$-regular graph with $s \geq 1$ and $|V| \geq 2$. Suppose the largest eigenvalue $\lambda_1$ is strictly larger than the second-largest:
\[
  \lambda_1(A(G)) > \lambda_2(A(G)).
\]
If $v : V \to \mathbb{R}$ satisfies $A(G)\,v = s\cdot v$, then $v$ is globally constant, i.e.\ there exists $c \in \mathbb{R}$ such that $v = c\,\mathbf{1}$.

\textit{Note: This is stated as an axiom because the required infrastructure (harmonic function theory on graphs and eigenspace dimension theory) is not currently available in Mathlib. The result combines two classical facts: (1) the discrete maximum principle shows eigenvectors for eigenvalue $s$ are constant on connected components, and (2) a simple top eigenvalue implies the graph is connected (otherwise each component contributes an independent eigenvector, giving multiplicity $\geq 2$).}
\end{theorem}

\begin{lemma}[Spectral Bound on the Quadratic Form]
\label{lem:AlonChung.spectral_bound}
\lean{AlonChung.spectral_bound}
\leanok
\uses{def:GraphExpansion.secondLargestEigenvalue, def:GraphExpansion.adjEigenvalues, lem:GraphExpansion.adjMatrix_isHermitian, def:AlonChung.indicatorVec, lem:AlonChung.quadratic_decomp, lem:AlonChung.spectral_decomp_quadratic, lem:AlonChung.parseval_eigenvector, lem:AlonChung.max_eigenvalue_eq_degree, thm:AlonChung.regular_eigenvector_s_constant}
Let $G$ be $s$-regular with $s \geq 1$, $|V| \geq 2$. Let $n = |V|$, $\gamma = |S|/n$, and $\lambda_2 = \lambda_2(A(G))$. Then
\[
  \mathbf{1}_S^\top A(G)\,\mathbf{1}_S \;\leq\; s\gamma^2 n + \lambda_2 \gamma(1-\gamma)n.
\]
\end{lemma}

\begin{proof}
\leanok
\uses{def:GraphExpansion.secondLargestEigenvalue, lem:GraphExpansion.adjMatrix_isHermitian, def:AlonChung.indicatorVec, lem:AlonChung.quadratic_decomp, lem:AlonChung.spectral_decomp_quadratic, lem:AlonChung.parseval_eigenvector, lem:AlonChung.max_eigenvalue_eq_degree, thm:AlonChung.regular_eigenvector_s_constant}

\textbf{\color{orange}[Uses unproven axiom: thm:AlonChung.regular\_eigenvector\_s\_constant]}

Set $z_v = \mathbf{1}_S(v) - \gamma$. Since $n > 0$ and $\gamma = |S|/n$:
\[
  \sum_v z_v = |S| - \gamma n = |S| - |S| = 0.
\]
Write $\mathbf{1}_S = \gamma\mathbf{1} + z$. By the quadratic decomposition (Lemma~\ref{lem:AlonChung.quadratic_decomp}):
\[
  \mathbf{1}_S^\top A\mathbf{1}_S = s\gamma^2 n + z^\top Az.
\]
It suffices to show $z^\top Az \leq \lambda_2\,\|z\|^2$, since $\|z\|^2 = \gamma(1-\gamma)n$ (computed by expanding $\sum_v(1_S(v)-\gamma)^2 = |S| - 2\gamma|S| + n\gamma^2 = \gamma(1-\gamma)n$).

We use the spectral decomposition (Lemma~\ref{lem:AlonChung.spectral_decomp_quadratic}) and Parseval (Lemma~\ref{lem:AlonChung.parseval_eigenvector}):
\[
  z^\top Az = \sum_j \lambda_j(z\cdot v_j)^2, \qquad \|z\|^2 = \sum_j (z\cdot v_j)^2.
\]
We need $\sum_j \lambda_j c_j^2 \leq \lambda_2 \sum_j c_j^2$ where $c_j = z\cdot v_j$.

\textbf{Case 1}: All eigenvalues $\leq \lambda_2$. Then term by term $\lambda_j c_j^2 \leq \lambda_2 c_j^2$, so the sum inequality holds.

\textbf{Case 2}: Some eigenvalue $> \lambda_2$. This forces a gap $\lambda_1 > \lambda_2$. For each $j$:
\begin{itemize}
\item If $\lambda_j \leq \lambda_2$: $\lambda_j c_j^2 \leq \lambda_2 c_j^2$. \checkmark
\item If $\lambda_j > \lambda_2$: By the antitone ordering of eigenvalues and the gap, $j$ must correspond to index $0$ (the largest), i.e.\ $\lambda_j = \lambda_1 = s$ (Lemma~\ref{lem:AlonChung.max_eigenvalue_eq_degree}). The eigenvector $v_j$ satisfies $A v_j = s\,v_j$, and by the axiom (Theorem~\ref{thm:AlonChung.regular_eigenvector_s_constant}), $v_j = c\mathbf{1}$ for some constant $c$. Then $z\cdot v_j = c\sum_v z_v = 0$, so $c_j = 0$ and $\lambda_j c_j^2 = 0 \leq \lambda_2 c_j^2$. \checkmark
\end{itemize}
Summing over all $j$ gives $z^\top Az \leq \lambda_2\|z\|^2$, completing the proof.
\end{proof}

\begin{theorem}[Alon--Chung Bound]
\label{thm:AlonChung.alonChung}
\lean{AlonChung.alonChung}
\leanok
\uses{def:GraphExpansion.secondLargestEigenvalue, def:GraphExpansion.adjEigenvalues, lem:GraphExpansion.adjMatrix_isHermitian, def:AlonChung.inducedEdges, def:AlonChung.indicatorVec, lem:AlonChung.quadratic_form_eq_twice_inducedEdges, lem:AlonChung.spectral_bound, lem:AlonChung.twice_card_edgeFinset_eq}
Let $G$ be a finite $s$-regular simple graph with $s \geq 1$ and $|V| \geq 2$. Let $\lambda_2 = \lambda_2(A(G))$ be the second-largest eigenvalue of the adjacency matrix. For any nonempty proper subset $S \subsetneq V$, set $n = |V|$, $\gamma = |S|/n$, and
\[
  \alpha \;=\; \gamma^2 + \frac{\lambda_2}{s}\,\gamma(1-\gamma).
\]
Then the number of edges in the induced subgraph on $S$ satisfies
\[
  |X(S)_1| \;\leq\; \alpha\,|E(G)|.
\]
\end{theorem}

\begin{proof}
\leanok
\uses{def:GraphExpansion.secondLargestEigenvalue, def:AlonChung.inducedEdges, def:AlonChung.indicatorVec, lem:AlonChung.quadratic_form_eq_twice_inducedEdges, lem:AlonChung.spectral_bound, lem:AlonChung.twice_card_edgeFinset_eq, thm:AlonChung.regular_eigenvector_s_constant}

\textbf{\color{orange}[Uses unproven axiom: thm:AlonChung.regular\_eigenvector\_s\_constant]}

\textbf{Step 1.} By the quadratic form identity (Lemma~\ref{lem:AlonChung.quadratic_form_eq_twice_inducedEdges}):
\[
  \mathbf{1}_S^\top A(G)\mathbf{1}_S = 2\,|X(S)_1|.
\]

\textbf{Step 2.} By the spectral bound (Lemma~\ref{lem:AlonChung.spectral_bound}):
\[
  \mathbf{1}_S^\top A(G)\mathbf{1}_S \leq s\gamma^2 n + \lambda_2\gamma(1-\gamma)n.
\]
Combining Steps 1 and 2 by linear arithmetic:
\[
  2\,|X(S)_1| \leq s\gamma^2 n + \lambda_2\gamma(1-\gamma)n.
\]

\textbf{Step 3.} By the edge count formula (Lemma~\ref{lem:AlonChung.twice_card_edgeFinset_eq}): $2|E(G)| = sn$, so $|E(G)| = sn/2$.

\textbf{Step 4.} We expand $\alpha\,|E(G)|$:
\[
  \alpha\,|E(G)| = \left(\gamma^2 + \frac{\lambda_2}{s}\gamma(1-\gamma)\right)\cdot\frac{sn}{2} = \frac{(s\gamma^2 + \lambda_2\gamma(1-\gamma))\,n}{2},
\]
where the field simplification uses $s > 0$. Combining with Step 2 and $n > 0$, $s > 0$:
\[
  |X(S)_1| = \frac{2|X(S)_1|}{2} \leq \frac{s\gamma^2 n + \lambda_2\gamma(1-\gamma)n}{2} = \alpha\,|E(G)|. \qedhere
\]
\end{proof}

%--- Cor_1: AlonChungContrapositive ---
\chapter{Cor 1: Alon--Chung Contrapositive}

\begin{definition}[Incident Vertices]
\label{def:AlonChungContrapositive.incidentVertices}
\lean{AlonChungContrapositive.incidentVertices}
\leanok
\uses{def:AlonChung.inducedEdges}
Let $E \subseteq X_1$ be a set of edges in a simple graph $G$ on vertex type $V$. The \emph{set of vertices incident to $E$} is defined as
\[
  \Gamma(E) \;=\; \{v \in V \mid \exists\, e \in E,\; v \in e\}.
\]
That is, $\Gamma(E)$ consists of all vertices that appear as an endpoint of some edge in $E$.
\end{definition}

\begin{lemma}[Membership in Incident Vertices]
\label{lem:AlonChungContrapositive.mem_incidentVertices}
\lean{AlonChungContrapositive.mem_incidentVertices}
\leanok
\uses{def:AlonChungContrapositive.incidentVertices}
For any edge set $E$ and vertex $v$,
\[
  v \in \Gamma(E) \;\Longleftrightarrow\; \exists\, e \in E,\; v \in e.
\]
\end{lemma}

\begin{proof}
\leanok
\uses{def:AlonChungContrapositive.incidentVertices}
This holds by unfolding the definition of $\Gamma(E)$ and simplifying.
\end{proof}

\begin{lemma}[Incident Vertices Nonempty]
\label{lem:AlonChungContrapositive.incidentVertices_nonempty}
\lean{AlonChungContrapositive.incidentVertices_nonempty}
\leanok
\uses{def:AlonChungContrapositive.incidentVertices, lem:AlonChungContrapositive.mem_incidentVertices}
If $E$ is a set of edges such that some edge in $E$ has a vertex, then $\Gamma(E)$ is nonempty.
\end{lemma}

\begin{proof}
\leanok
\uses{def:AlonChungContrapositive.incidentVertices, lem:AlonChungContrapositive.mem_incidentVertices}
Suppose there exists $e \in E$ and $v \in e$. We obtain $e$, $v$, and the membership proof. By the membership characterization of $\Gamma(E)$, $v \in \Gamma(E)$, so $\Gamma(E)$ is nonempty.
\end{proof}

\begin{lemma}[Edges Contained in Induced Edges of Incident Vertices]
\label{lem:AlonChungContrapositive.subset_inducedEdges_incidentVertices}
\lean{AlonChungContrapositive.subset_inducedEdges_incidentVertices}
\leanok
\uses{def:AlonChung.inducedEdges, def:AlonChungContrapositive.incidentVertices, lem:AlonChung.mem_inducedEdges, lem:AlonChungContrapositive.mem_incidentVertices}
Let $G$ be a simple graph and $E \subseteq \operatorname{edgeFinset}(G)$. Then
\[
  E \;\subseteq\; \operatorname{inducedEdges}(G,\, \Gamma(E)).
\]
\end{lemma}

\begin{proof}
\leanok
\uses{def:AlonChung.inducedEdges, def:AlonChungContrapositive.incidentVertices, lem:AlonChung.mem_inducedEdges, lem:AlonChungContrapositive.mem_incidentVertices}
Let $e \in E$ be arbitrary. We must show $e \in \operatorname{inducedEdges}(G, \Gamma(E))$. Rewriting via the membership criterion for $\operatorname{inducedEdges}$, we need: (1) $e \in \operatorname{edgeFinset}(G)$, which holds since $e \in E \subseteq \operatorname{edgeFinset}(G)$; and (2) every vertex $v \in e$ lies in $\Gamma(E)$. For any such $v$, rewriting via the membership criterion for $\Gamma(E)$, we exhibit $e \in E$ and $v \in e$, establishing $v \in \Gamma(E)$.
\end{proof}

\begin{lemma}[Monotonicity of the Alon--Chung Bound Function]
\label{lem:AlonChungContrapositive.alpha_mono}
\lean{AlonChungContrapositive.alpha_mono}
\leanok
\uses{def:AlonChung.inducedEdges}
Let $c \in \mathbb{R}$ with $0 \leq c \leq 1$, and let $0 \leq a \leq b$. Then
\[
  a^2 + c \cdot a \cdot (1 - a) \;\leq\; b^2 + c \cdot b \cdot (1 - b).
\]
\end{lemma}

\begin{proof}
\leanok
\uses{def:AlonChung.inducedEdges}
Define $f(t) = t^2 + c \cdot t(1-t) = (1-c)t^2 + ct$. Then
\[
  f(b) - f(a) = (1-c)(b^2 - a^2) + c(b - a) = (b - a)\bigl[(1-c)(b+a) + c\bigr].
\]
Since $b \geq a \geq 0$ and $0 \leq c \leq 1$, we have $b - a \geq 0$, $(1-c) \geq 0$, $b + a \geq 0$, and $c \geq 0$, so $(b-a)[(1-c)(b+a)+c] \geq 0$. This is verified by nonlinear arithmetic using $a^2 \geq 0$, $b^2 \geq 0$, $(b-a)^2 \geq 0$, and the nonnegativity of $(b-a)\cdot[(1-c)(b+a)+c]$.
\end{proof}

\begin{theorem}[Alon--Chung Contrapositive, Direct Form]
\label{thm:AlonChungContrapositive.alonChungContrapositive_direct}
\lean{AlonChungContrapositive.alonChungContrapositive_direct}
\leanok
\uses{def:AlonChung.inducedEdges, lem:AlonChung.max_eigenvalue_eq_degree, lem:AlonChung.mem_inducedEdges, thm:AlonChung.alonChung}
Let $G$ be a $s$-regular simple graph on $V$ with $|V| \geq 2$ and $s \geq 1$. Let $\lambda_2$ denote the second-largest adjacency eigenvalue of $G$ with $\lambda_2 \geq 0$. Let $E \subseteq \operatorname{edgeFinset}(G)$ and $0 < \gamma < 1$. If
\[
  |\Gamma(E)| \;\leq\; \gamma \cdot |V|,
\]
then
\[
  |E| \;\leq\; \left(\gamma^2 + \frac{\lambda_2}{s}\,\gamma(1-\gamma)\right) |\operatorname{edgeFinset}(G)|.
\]
\end{theorem}

\begin{proof}
\leanok
\uses{def:AlonChung.inducedEdges, lem:AlonChung.max_eigenvalue_eq_degree, lem:AlonChung.mem_inducedEdges, thm:AlonChung.alonChung, lem:AlonChungContrapositive.subset_inducedEdges_incidentVertices, lem:AlonChungContrapositive.alpha_mono}
Let $S = \Gamma(E)$. We consider two cases.

\textbf{Case 1: $S = \emptyset$.} We claim $E = \emptyset$. Indeed, suppose for contradiction that $E$ is nonempty; pick $e \in E$. Since $e \in \operatorname{edgeFinset}(G)$, writing $e = s(v,w)$ via the induction principle for $\operatorname{Sym2}$, we have $v \in S$ by definition of $\Gamma(E)$ (since $v \in e \in E$), contradicting $S = \emptyset$. Thus $|E| = 0 \leq \alpha \cdot |\operatorname{edgeFinset}(G)|$, where $\alpha = \gamma^2 + (\lambda_2/s)\gamma(1-\gamma) \geq 0$ (since $\lambda_2/s \geq 0$, $\gamma > 0$, and $1 - \gamma > 0$, this follows by nonlinear arithmetic).

\textbf{Case 2: $S \neq \emptyset$.} Set $n = |V|$ and $\gamma' = |S|/n$. Then $\gamma' > 0$ (since $S$ is nonempty and $n > 0$) and $\gamma' \leq \gamma$ (from the hypothesis $|S| \leq \gamma n$). Moreover, $|S| < |V|$: if $|S| \geq |V|$, then $|V| \leq |S| \leq \gamma |V|$, giving $1 \leq \gamma$, contradicting $\gamma < 1$.

By Lemma~\ref{lem:AlonChungContrapositive.subset_inducedEdges_incidentVertices}, $E \subseteq \operatorname{inducedEdges}(G, S)$, so $|E| \leq |\operatorname{inducedEdges}(G,S)|$.

Applying the Alon--Chung theorem (Theorem~\ref{thm:AlonChung.alonChung}) to $S$:
\[
  |\operatorname{inducedEdges}(G,S)| \;\leq\; \alpha' \cdot |\operatorname{edgeFinset}(G)|,
\]
where $\alpha' = {\gamma'}^2 + (\lambda_2/s)\,\gamma'(1-\gamma')$.

We verify $\lambda_2/s \leq 1$: by Lemma~\ref{lem:AlonChung.max_eigenvalue_eq_degree}, the largest eigenvalue equals $s$, and by antitonicity of adjacency eigenvalues (Theorem~\ref{thm:GraphExpansion.adjEigenvalues_antitone}), $\lambda_2 \leq s$, so $\lambda_2/s \leq 1$.

By Lemma~\ref{lem:AlonChungContrapositive.alpha_mono} applied with $c = \lambda_2/s \in [0,1]$ and $0 < \gamma' \leq \gamma$, we obtain $\alpha' \leq \alpha = \gamma^2 + (\lambda_2/s)\gamma(1-\gamma)$.

Combining via a calculation:
\[
  |E| \;\leq\; |\operatorname{inducedEdges}(G,S)| \;\leq\; \alpha' \cdot |\operatorname{edgeFinset}(G)| \;\leq\; \alpha \cdot |\operatorname{edgeFinset}(G)|,
\]
where the last step uses $\alpha' \leq \alpha$ and $|\operatorname{edgeFinset}(G)| \geq 0$.
\end{proof}

\begin{theorem}[Alon--Chung Contrapositive]
\label{thm:AlonChungContrapositive.alonChungContrapositive}
\lean{AlonChungContrapositive.alonChungContrapositive}
\leanok
\uses{def:AlonChung.inducedEdges, lem:AlonChung.max_eigenvalue_eq_degree, lem:AlonChung.mem_inducedEdges, thm:AlonChung.alonChung}
Let $G$ be a $s$-regular simple graph on $V$ with $|V| \geq 2$ and $s \geq 1$. Let $\lambda_2 \geq 0$ be the second-largest adjacency eigenvalue. Let $E \subseteq \operatorname{edgeFinset}(G)$ and $0 < \gamma < 1$. Set $\alpha = \gamma^2 + \tfrac{\lambda_2}{s}\,\gamma(1-\gamma)$. If
\[
  |E| \;>\; \alpha \cdot |\operatorname{edgeFinset}(G)|,
\]
then
\[
  |\Gamma(E)| \;>\; \gamma \cdot |V|.
\]
In other words, if $E$ is a large fraction of edges, then $E$ is incident to many vertices.
\end{theorem}

\begin{proof}
\leanok
\uses{def:AlonChung.inducedEdges, lem:AlonChung.max_eigenvalue_eq_degree, lem:AlonChung.mem_inducedEdges, thm:AlonChung.alonChung, thm:AlonChungContrapositive.alonChungContrapositive_direct}
We argue by contradiction. Suppose $|\Gamma(E)| \leq \gamma \cdot |V|$. Then by Theorem~\ref{thm:AlonChungContrapositive.alonChungContrapositive_direct}, we have $|E| \leq \alpha \cdot |\operatorname{edgeFinset}(G)|$. This contradicts the hypothesis $|E| > \alpha \cdot |\operatorname{edgeFinset}(G)|$, so by linear arithmetic we obtain a contradiction. Therefore $|\Gamma(E)| > \gamma \cdot |V|$.
\end{proof}

%--- Def_19: EdgeBoundaryVertex ---
\chapter{Def 19: Edge Boundary and Vertex Incidence for Cell Complexes}

\begin{definition}[Edge Boundary of a Vertex Subset]
\label{def:CellComplex.edgeBoundaryCell}
\lean{CellComplex.edgeBoundaryCell}
\leanok
\uses{def:CellComplex}

Let $X$ be a cell complex and let $S \subseteq X_0$ be a subset of vertices (0-cells). The \emph{edge boundary} of $S$ is the set
\[
  \delta S \;=\; \bigl\{\, e \in X_1 \;\big|\; (\exists\, u \in \partial_1 e,\; u \in S) \;\land\; (\exists\, w \in \partial_1 e,\; w \notin S) \,\bigr\}.
\]
That is, an edge $e \in X_1$ belongs to $\delta S$ if and only if its boundary $\partial_1 e$ intersects both $S$ and its complement $S^c$.
\end{definition}

\begin{theorem}[Membership in the Edge Boundary]
\label{thm:CellComplex.mem_edgeBoundaryCell}
\lean{CellComplex.mem_edgeBoundaryCell}
\leanok
\uses{def:CellComplex, def:CellComplex.edgeBoundaryCell}

Let $X$ be a cell complex, $S \subseteq X_0$, and $e \in X_1$. Then
\[
  e \in \delta S \;\Longleftrightarrow\; \bigl(\exists\, u \in \partial_1 e,\; u \in S\bigr) \;\land\; \bigl(\exists\, w \in \partial_1 e,\; w \notin S\bigr).
\]
\end{theorem}

\begin{proof}
\leanok
\uses{def:CellComplex.edgeBoundaryCell}

This follows immediately by unfolding the definition of $\delta S$ and applying \texttt{simp}.
\end{proof}

\begin{lemma}[Edge Boundary Nonempty]
\label{lem:CellComplex.edgeBoundaryCell_nonempty}
\lean{CellComplex.edgeBoundaryCell_nonempty}
\leanok
\uses{def:CellComplex, def:CellComplex.edgeBoundaryCell, thm:CellComplex.mem_edgeBoundaryCell}

Let $X$ be a cell complex and $S \subseteq X_0$. If there exists an edge $e \in X_1$ such that $\partial_1 e$ contains a vertex in $S$ and a vertex not in $S$, then $\delta S$ is nonempty.
\end{lemma}

\begin{proof}
\leanok
\uses{def:CellComplex.edgeBoundaryCell, thm:CellComplex.mem_edgeBoundaryCell}

Assume there exists $e \in X_1$ with $(\exists\, u \in \partial_1 e,\; u \in S)$ and $(\exists\, w \in \partial_1 e,\; w \notin S)$. From the hypothesis $h$, we obtain such an edge $e$ together with the witness $h_e$ of both conditions. By the membership characterization $\texttt{mem\_edgeBoundaryCell}$, we have $e \in \delta S$, so $\delta S$ is nonempty.
\end{proof}

\begin{definition}[Vertices Incident to a Set of Edges]
\label{def:CellComplex.incidentVerticesCell}
\lean{CellComplex.incidentVerticesCell}
\leanok
\uses{def:CellComplex}

Let $X$ be a cell complex and let $E \subseteq X_1$ be a set of edges (1-cells). The \emph{set of vertices incident to $E$} is
\[
  \Gamma(E) \;=\; \bigl\{\, v \in X_0 \;\big|\; \exists\, e \in E,\; v \in \partial_1 e \,\bigr\}.
\]
\end{definition}

\begin{theorem}[Membership in Incident Vertices]
\label{thm:CellComplex.mem_incidentVerticesCell}
\lean{CellComplex.mem_incidentVerticesCell}
\leanok
\uses{def:CellComplex, def:CellComplex.incidentVerticesCell}

Let $X$ be a cell complex, $E \subseteq X_1$, and $v \in X_0$. Then
\[
  v \in \Gamma(E) \;\Longleftrightarrow\; \exists\, e \in E,\; v \in \partial_1 e.
\]
\end{theorem}

\begin{proof}
\leanok
\uses{def:CellComplex.incidentVerticesCell}

This follows immediately by unfolding the definition of $\Gamma(E)$ and applying \texttt{simp}.
\end{proof}

\begin{lemma}[Incident Vertices Nonempty]
\label{lem:CellComplex.incidentVerticesCell_nonempty}
\lean{CellComplex.incidentVerticesCell_nonempty}
\leanok
\uses{def:CellComplex, def:CellComplex.incidentVerticesCell, thm:CellComplex.mem_incidentVerticesCell}

Let $X$ be a cell complex and $E \subseteq X_1$. If there exists an edge $e \in E$ with a vertex $v \in X_0$ such that $v \in \partial_1 e$, then $\Gamma(E)$ is nonempty.
\end{lemma}

\begin{proof}
\leanok
\uses{def:CellComplex.incidentVerticesCell, thm:CellComplex.mem_incidentVerticesCell}

From the hypothesis $h$, we obtain $e \in E$ and $v \in X_0$ with $v \in \partial_1 e$. By the membership characterization $\texttt{mem\_incidentVerticesCell}$, we conclude $v \in \Gamma(E)$, so $\Gamma(E)$ is nonempty.
\end{proof}

\begin{definition}[Edges Incident to a Vertex (Star)]
\label{def:CellComplex.incidentEdges}
\lean{CellComplex.incidentEdges}
\leanok
\uses{def:CellComplex}

Let $X$ be a cell complex and $v \in X_0$ a vertex. The \emph{set of edges incident to $v$} (the star of $v$) is
\[
  \delta v \;=\; \bigl\{\, e \in X_1 \;\big|\; v \in \partial_1 e \,\bigr\}.
\]
\end{definition}

\begin{theorem}[Membership in Incident Edges]
\label{thm:CellComplex.mem_incidentEdges}
\lean{CellComplex.mem_incidentEdges}
\leanok
\uses{def:CellComplex, def:CellComplex.incidentEdges}

Let $X$ be a cell complex, $v \in X_0$, and $e \in X_1$. Then
\[
  e \in \delta v \;\Longleftrightarrow\; v \in \partial_1 e.
\]
\end{theorem}

\begin{proof}
\leanok
\uses{def:CellComplex.incidentEdges}

This follows immediately by unfolding the definition of $\delta v$ and applying \texttt{simp}.
\end{proof}

\begin{lemma}[Incident Edges Nonempty]
\label{lem:CellComplex.incidentEdges_nonempty}
\lean{CellComplex.incidentEdges_nonempty}
\leanok
\uses{def:CellComplex, def:CellComplex.incidentEdges, thm:CellComplex.mem_incidentEdges}

Let $X$ be a cell complex and $v \in X_0$. If there exists $e \in X_1$ such that $v \in \partial_1 e$, then $\delta v$ is nonempty.
\end{lemma}

\begin{proof}
\leanok
\uses{def:CellComplex.incidentEdges, thm:CellComplex.mem_incidentEdges}

From the hypothesis $h$, we obtain $e \in X_1$ with $v \in \partial_1 e$. By the membership characterization $\texttt{mem\_incidentEdges}$, we conclude $e \in \delta v$, so $\delta v$ is nonempty.
\end{proof}

\begin{definition}[Edge Boundary at a Vertex]
\label{def:CellComplex.edgeBoundaryVertex}
\lean{CellComplex.edgeBoundaryVertex}
\leanok
\uses{def:CellComplex, def:CellComplex.edgeBoundaryCell, def:CellComplex.incidentEdges}

Let $X$ be a cell complex, $S \subseteq X_0$, and $v \in X_0$. The \emph{edge boundary of $S$ at vertex $v$} is the intersection
\[
  (\delta S)_v \;=\; \delta S \cap \delta v \;=\; \bigl\{\, e \in X_1 \;\big|\; e \in \delta S \;\land\; v \in \partial_1 e \,\bigr\}.
\]
\end{definition}

\begin{theorem}[Membership in the Edge Boundary at a Vertex]
\label{thm:CellComplex.mem_edgeBoundaryVertex}
\lean{CellComplex.mem_edgeBoundaryVertex}
\leanok
\uses{def:CellComplex, def:CellComplex.edgeBoundaryVertex, def:CellComplex.edgeBoundaryCell}

Let $X$ be a cell complex, $S \subseteq X_0$, $v \in X_0$, and $e \in X_1$. Then
\[
  e \in (\delta S)_v \;\Longleftrightarrow\; e \in \delta S \;\land\; v \in \partial_1 e.
\]
\end{theorem}

\begin{proof}
\leanok
\uses{def:CellComplex.edgeBoundaryVertex}

This follows immediately by unfolding the definition of $(\delta S)_v$ and applying \texttt{simp}.
\end{proof}

\begin{theorem}[Edge Boundary at Vertex is Subset of Edge Boundary]
\label{thm:CellComplex.edgeBoundaryVertex_subset}
\lean{CellComplex.edgeBoundaryVertex_subset}
\leanok
\uses{def:CellComplex, def:CellComplex.edgeBoundaryVertex, def:CellComplex.edgeBoundaryCell}

Let $X$ be a cell complex, $S \subseteq X_0$, and $v \in X_0$. Then
\[
  (\delta S)_v \;\subseteq\; \delta S.
\]
\end{theorem}

\begin{proof}
\leanok
\uses{def:CellComplex.edgeBoundaryVertex, def:CellComplex.edgeBoundaryCell}

Since $(\delta S)_v$ is defined as the filter of $\delta S$ by the predicate $v \in \partial_1 e$, the inclusion $(\delta S)_v \subseteq \delta S$ holds by the general fact that a filtered finset is a subset of the original finset (\texttt{Finset.filter\_subset}).
\end{proof}

\begin{theorem}[Edge Boundary at Vertex is Subset of Incident Edges]
\label{thm:CellComplex.edgeBoundaryVertex_subset_incidentEdges}
\lean{CellComplex.edgeBoundaryVertex_subset_incidentEdges}
\leanok
\uses{def:CellComplex, def:CellComplex.edgeBoundaryVertex, def:CellComplex.incidentEdges}

Let $X$ be a cell complex, $S \subseteq X_0$, and $v \in X_0$. Then
\[
  (\delta S)_v \;\subseteq\; \delta v.
\]
\end{theorem}

\begin{proof}
\leanok
\uses{def:CellComplex.edgeBoundaryVertex, def:CellComplex.incidentEdges, thm:CellComplex.mem_edgeBoundaryVertex, thm:CellComplex.mem_incidentEdges}

Let $e$ be arbitrary and assume $e \in (\delta S)_v$. Rewriting using the membership characterization $\texttt{mem\_edgeBoundaryVertex}$, we obtain $e \in \delta S$ and $v \in \partial_1 e$. Rewriting the goal using $\texttt{mem\_incidentEdges}$, it suffices to show $v \in \partial_1 e$, which is precisely the second component of the hypothesis.
\end{proof}

\begin{lemma}[Edge Boundary at Vertex Nonempty]
\label{lem:CellComplex.edgeBoundaryVertex_nonempty}
\lean{CellComplex.edgeBoundaryVertex_nonempty}
\leanok
\uses{def:CellComplex, def:CellComplex.edgeBoundaryVertex, def:CellComplex.edgeBoundaryCell, thm:CellComplex.mem_edgeBoundaryVertex}

Let $X$ be a cell complex, $S \subseteq X_0$, and $v \in X_0$. If there exists $e \in X_1$ such that $e \in \delta S$ and $v \in \partial_1 e$, then $(\delta S)_v$ is nonempty.
\end{lemma}

\begin{proof}
\leanok
\uses{def:CellComplex.edgeBoundaryVertex, thm:CellComplex.mem_edgeBoundaryVertex}

From the hypothesis $h$, we obtain $e \in X_1$ with the witness $h_e$ satisfying $e \in \delta S$ and $v \in \partial_1 e$. By the membership characterization $\texttt{mem\_edgeBoundaryVertex}$, we conclude $e \in (\delta S)_v$, so $(\delta S)_v$ is nonempty.
\end{proof}

\begin{theorem}[Alternate Characterization of Edge Boundary at Vertex]
\label{thm:CellComplex.mem_edgeBoundaryVertex_iff}
\lean{CellComplex.mem_edgeBoundaryVertex_iff}
\leanok
\uses{def:CellComplex, def:CellComplex.edgeBoundaryVertex, def:CellComplex.edgeBoundaryCell}

Let $X$ be a cell complex, $S \subseteq X_0$, $v \in X_0$, and $e \in X_1$. Then
\[
  e \in (\delta S)_v \;\Longleftrightarrow\; v \in \partial_1 e \;\land\; \bigl(\exists\, u \in \partial_1 e,\; u \in S\bigr) \;\land\; \bigl(\exists\, w \in \partial_1 e,\; w \notin S\bigr).
\]
\end{theorem}

\begin{proof}
\leanok
\uses{thm:CellComplex.mem_edgeBoundaryVertex, thm:CellComplex.mem_edgeBoundaryCell}

Rewriting with $\texttt{mem\_edgeBoundaryVertex}$, the left-hand side becomes $(e \in \delta S) \land (v \in \partial_1 e)$. Rewriting with $\texttt{mem\_edgeBoundaryCell}$, the condition $e \in \delta S$ expands to $(\exists\, u \in \partial_1 e,\; u \in S) \land (\exists\, w \in \partial_1 e,\; w \notin S)$. The equivalence with the right-hand side then follows by propositional logic (\texttt{tauto}).
\end{proof}

%--- Lem_2: EdgeToVertexExpansion ---
\chapter{Lem 2: Edge-to-Vertex Expansion}

\begin{definition}[Quadratic Inverse $\gamma(\alpha)$]
\label{def:EdgeToVertexExpansion.gammaOfAlpha}
\lean{EdgeToVertexExpansion.gammaOfAlpha}
\leanok
\uses{def:AlonChungContrapositive.incidentVertices}
Let $s, \lambda_2, a \in \mathbb{R}$. The \emph{quadratic inverse} is defined by
\[
\gamma(s, \lambda_2, a) \;:=\; \frac{\sqrt{\lambda_2^2 + 4s(s - \lambda_2)\,a} - \lambda_2}{2(s - \lambda_2)}.
\]
This is the solution $g$ to the equation $a = g^2 + \tfrac{\lambda_2}{s}\,g(1-g)$.
\end{definition}

\begin{definition}[Expansion Parameter $\beta$]
\label{def:EdgeToVertexExpansion.betaParam}
\lean{EdgeToVertexExpansion.betaParam}
\leanok
\uses{def:EdgeToVertexExpansion.gammaOfAlpha}
Let $s, \lambda_2, a \in \mathbb{R}$. The \emph{expansion parameter} is defined by
\[
\beta(s, \lambda_2, a) \;:=\; \frac{\sqrt{\lambda_2^2 + 4s(s - \lambda_2)\,a} - \lambda_2}{s(s-\lambda_2)\,a}.
\]
\end{definition}

\begin{lemma}[$\beta$ Equals $2\gamma / (s\alpha)$]
\label{lem:EdgeToVertexExpansion.betaParam_eq}
\lean{EdgeToVertexExpansion.betaParam_eq}
\leanok
\uses{def:EdgeToVertexExpansion.betaParam, def:EdgeToVertexExpansion.gammaOfAlpha}
Let $s, \lambda_2, a \in \mathbb{R}$ with $s \neq 0$ and $a \neq 0$. Then
\[
\beta(s, \lambda_2, a) \;=\; \frac{2\,\gamma(s, \lambda_2, a)}{s \cdot a}.
\]
\end{lemma}

\begin{proof}
\leanok
\uses{def:EdgeToVertexExpansion.betaParam, def:EdgeToVertexExpansion.gammaOfAlpha}
Unfolding the definitions of $\beta$ and $\gamma$ and simplifying the field expressions, the equality holds by algebraic manipulation.
\end{proof}

\begin{lemma}[Discriminant Nonnegativity]
\label{lem:EdgeToVertexExpansion.discriminant_nonneg}
\lean{EdgeToVertexExpansion.discriminant_nonneg}
\leanok
\uses{def:EdgeToVertexExpansion.gammaOfAlpha}
Let $s, \lambda_2, a \in \mathbb{R}$ with $0 < s$, $\lambda_2 < s$, and $0 \leq a$. Then
\[
0 \;\leq\; \lambda_2^2 + 4s(s - \lambda_2)\,a.
\]
\end{lemma}

\begin{proof}
\leanok
\uses{def:EdgeToVertexExpansion.gammaOfAlpha}
We observe that $0 \leq \lambda_2^2$ (since it is a square), $0 < s - \lambda_2$ (since $\lambda_2 < s$), and hence $0 \leq 4s(s-\lambda_2)a$ by positivity (since $s > 0$, $s - \lambda_2 > 0$, $a \geq 0$). The result follows by linear arithmetic from these three facts.
\end{proof}

\begin{lemma}[$\sqrt{\lambda_2^2 + 4s(s-\lambda_2)a} \geq \lambda_2$]
\label{lem:EdgeToVertexExpansion.sqrt_ge_lam2}
\lean{EdgeToVertexExpansion.sqrt_ge_lam2}
\leanok
\uses{lem:EdgeToVertexExpansion.discriminant_nonneg}
Let $s, \lambda_2, a \in \mathbb{R}$ with $0 < s$, $\lambda_2 < s$, and $0 \leq a$. Then
\[
\lambda_2 \;\leq\; \sqrt{\lambda_2^2 + 4s(s - \lambda_2)\,a}.
\]
\end{lemma}

\begin{proof}
\leanok
\uses{lem:EdgeToVertexExpansion.discriminant_nonneg}
We proceed by cases on whether $\lambda_2 \geq 0$ or $\lambda_2 < 0$.

\textbf{Case 1:} $\lambda_2 \geq 0$. We apply the fact that $\lambda_2 \leq \sqrt{D}$ whenever $\lambda_2^2 \leq D$ and $D \geq 0$. Since $4s(s-\lambda_2)a \geq 0$ (using $s > 0$, $s - \lambda_2 > 0$, $a \geq 0$), we have $\lambda_2^2 \leq \lambda_2^2 + 4s(s-\lambda_2)a$, from which the result follows.

\textbf{Case 2:} $\lambda_2 < 0$. Since $\sqrt{\cdot} \geq 0$, we have $\lambda_2 < 0 \leq \sqrt{\lambda_2^2 + 4s(s-\lambda_2)a}$.
\end{proof}

\begin{lemma}[$\gamma(s, \lambda_2, a) \geq 0$]
\label{lem:EdgeToVertexExpansion.gammaOfAlpha_nonneg}
\lean{EdgeToVertexExpansion.gammaOfAlpha_nonneg}
\leanok
\uses{def:EdgeToVertexExpansion.gammaOfAlpha, lem:EdgeToVertexExpansion.sqrt_ge_lam2}
Let $s, \lambda_2, a \in \mathbb{R}$ with $0 < s$, $\lambda_2 < s$, and $0 \leq a$. Then $\gamma(s, \lambda_2, a) \geq 0$.
\end{lemma}

\begin{proof}
\leanok
\uses{def:EdgeToVertexExpansion.gammaOfAlpha, lem:EdgeToVertexExpansion.sqrt_ge_lam2}
Unfolding the definition, $\gamma(s, \lambda_2, a) = (\sqrt{\lambda_2^2 + 4s(s-\lambda_2)a} - \lambda_2) / (2(s-\lambda_2))$. The numerator is nonneg by the lemma $\sqrt{\lambda_2^2 + 4s(s-\lambda_2)a} \geq \lambda_2$, and the denominator $2(s - \lambda_2) > 0$ since $\lambda_2 < s$. Hence the quotient is nonneg.
\end{proof}

\begin{lemma}[$\gamma(s, \lambda_2, 0) = 0$]
\label{lem:EdgeToVertexExpansion.gammaOfAlpha_zero}
\lean{EdgeToVertexExpansion.gammaOfAlpha_zero}
\leanok
\uses{def:EdgeToVertexExpansion.gammaOfAlpha}
Let $s, \lambda_2 \in \mathbb{R}$ with $0 < s$, $\lambda_2 < s$, and $\lambda_2 \geq 0$. Then $\gamma(s, \lambda_2, 0) = 0$.
\end{lemma}

\begin{proof}
\leanok
\uses{def:EdgeToVertexExpansion.gammaOfAlpha}
Unfolding the definition and substituting $a = 0$, we compute:
\[
\gamma(s, \lambda_2, 0) = \frac{\sqrt{\lambda_2^2 + 4s(s-\lambda_2) \cdot 0} - \lambda_2}{2(s-\lambda_2)} = \frac{\sqrt{\lambda_2^2} - \lambda_2}{2(s-\lambda_2)} = \frac{\lambda_2 - \lambda_2}{2(s-\lambda_2)} = 0,
\]
using $\sqrt{\lambda_2^2} = \lambda_2$ (since $\lambda_2 \geq 0$) by simplification.
\end{proof}

\begin{lemma}[Concavity of $\gamma$]
\label{lem:EdgeToVertexExpansion.gammaOfAlpha_concavity}
\lean{EdgeToVertexExpansion.gammaOfAlpha_concavity}
\leanok
\uses{def:EdgeToVertexExpansion.gammaOfAlpha, lem:EdgeToVertexExpansion.discriminant_nonneg, lem:EdgeToVertexExpansion.sqrt_ge_lam2}
Let $s, \lambda_2, \alpha, \tilde{\alpha} \in \mathbb{R}$ with $0 < s$, $\lambda_2 < s$, $0 < \alpha$, $0 \leq \tilde{\alpha}$, and $\tilde{\alpha} \leq \alpha$. Then
\[
\gamma(s, \lambda_2, \alpha) \cdot \tilde{\alpha} \;\leq\; \gamma(s, \lambda_2, \tilde{\alpha}) \cdot \alpha.
\]
\end{lemma}

\begin{proof}
\leanok
\uses{def:EdgeToVertexExpansion.gammaOfAlpha, lem:EdgeToVertexExpansion.discriminant_nonneg, lem:EdgeToVertexExpansion.sqrt_ge_lam2}
Unfolding the definitions and cross-multiplying by $2(s - \lambda_2) > 0$, it suffices to show
\[
(\sqrt{D} - \lambda_2)\,\tilde{\alpha} \;\leq\; (\sqrt{\tilde{D}} - \lambda_2)\,\alpha,
\]
where $D = \lambda_2^2 + c\alpha$, $\tilde{D} = \lambda_2^2 + c\tilde{\alpha}$, and $c = 4s(s-\lambda_2) > 0$.

Set $u = \sqrt{\tilde{D}}$ and $v = \sqrt{D}$. Both are nonneg, and $\lambda_2 \leq u \leq v$ (by the lemma $\sqrt{\cdot} \geq \lambda_2$ and monotonicity of square root: $\tilde{D} \leq D$ since $c > 0$ and $\tilde{\alpha} \leq \alpha$). In particular, $(u - \lambda_2)(v - \lambda_2) \geq 0$.

It suffices to show $(v - \lambda_2)\tilde{\alpha} \leq (u - \lambda_2)\alpha$, which in turn follows from the intermediate claim
\[
(v - u)\,\tilde{\alpha} \;\leq\; (u - \lambda_2)\,(\alpha - \tilde{\alpha}),
\]
since adding $(u - \lambda_2)\tilde{\alpha}$ to both sides yields the desired inequality.

We prove this intermediate claim by cases on whether $v + u > 0$ or $v + u = 0$.

\textbf{Case 1:} $v + u > 0$. We multiply through by $v + u$ and compute:
\begin{align*}
(v - u)\,\tilde{\alpha}\,(v + u) &= \tilde{\alpha}\,((v-u)(v+u)) = \tilde{\alpha}\,(v^2 - u^2) = \tilde{\alpha}\,c(\alpha - \tilde{\alpha}) \\
&= c\tilde{\alpha}(\alpha - \tilde{\alpha}) = (u^2 - \lambda_2^2)(\alpha - \tilde{\alpha}) = (u-\lambda_2)(u+\lambda_2)(\alpha - \tilde{\alpha}) \\
&\leq (u - \lambda_2)(v + \lambda_2)(\alpha - \tilde{\alpha})
\end{align*}
using $u \leq v$ and $u - \lambda_2 \geq 0$ and $\alpha - \tilde{\alpha} \geq 0$. Then $(u-\lambda_2)(v+\lambda_2)(\alpha-\tilde{\alpha}) \leq (u-\lambda_2)(\alpha-\tilde{\alpha})(v+u)$ using $\lambda_2 \leq u$. Dividing by $v + u > 0$ gives the claim.

\textbf{Case 2:} $v + u = 0$. Since $u, v \geq 0$, we have $u = v = 0$. Then $\lambda_2 \leq u = 0$, so $\lambda_2 \leq 0$. Setting $u = v = 0$ and simplifying, both sides reduce to $0 \leq (-\lambda_2)(\alpha - \tilde{\alpha})$, which holds since $-\lambda_2 \geq 0$ and $\alpha - \tilde{\alpha} \geq 0$.
\end{proof}

\begin{lemma}[Handshaking: $|E(G)| = \tfrac{s}{2}|V|$]
\label{lem:EdgeToVertexExpansion.edgeFinset_eq_half_sn}
\lean{EdgeToVertexExpansion.edgeFinset_eq_half_sn}
\leanok
\uses{lem:AlonChung.twice_card_edgeFinset_eq}
Let $G$ be a finite $s$-regular simple graph on $V$. Then
\[
|E(G)| \;=\; \frac{s}{2}\,|V|.
\]
\end{lemma}

\begin{proof}
\leanok
\uses{lem:AlonChung.twice_card_edgeFinset_eq}
This follows immediately from the handshaking lemma \texttt{twice\_card\_edgeFinset\_eq}, which gives $2|E(G)| = s\,|V|$, by linear arithmetic.
\end{proof}

\begin{lemma}[$|\Gamma(E)| \geq \gamma(\tilde{\alpha})\,|V|$]
\label{lem:EdgeToVertexExpansion.incidentVertices_ge_gamma}
\lean{EdgeToVertexExpansion.incidentVertices_ge_gamma}
\leanok
\uses{def:AlonChungContrapositive.incidentVertices, lem:AlonChungContrapositive.mem_incidentVertices, lem:AlonChungContrapositive.subset_inducedEdges_incidentVertices, def:EdgeToVertexExpansion.gammaOfAlpha, lem:EdgeToVertexExpansion.gammaOfAlpha_concavity, lem:EdgeToVertexExpansion.discriminant_nonneg, lem:EdgeToVertexExpansion.sqrt_ge_lam2, thm:AlonChung.alonChung}
Let $G$ be a finite connected $s$-regular simple graph on $V$ with $|V| \geq 2$ and second-largest eigenvalue $\lambda_2 < s$. Let $E \subseteq E(G)$ be a nonempty set of edges, and set $\tilde{\alpha} = |E|/|E(G)|$. Then
\[
|\Gamma(E)| \;\geq\; \gamma(s, \lambda_2, \tilde{\alpha})\,|V|,
\]
where $\Gamma(E) = \operatorname{incidentVertices}(E)$ is the set of vertices incident to some edge in $E$.
\end{lemma}

\begin{proof}
\leanok
\uses{def:AlonChungContrapositive.incidentVertices, lem:AlonChungContrapositive.mem_incidentVertices, lem:AlonChungContrapositive.subset_inducedEdges_incidentVertices, def:EdgeToVertexExpansion.gammaOfAlpha, lem:EdgeToVertexExpansion.discriminant_nonneg, lem:EdgeToVertexExpansion.sqrt_ge_lam2, thm:AlonChung.alonChung}
Set $S = \Gamma(E)$, $n = |V|$, and $\tilde{\alpha} = |E|/|E(G)|$.

First, $S$ is nonempty: take any edge $e \in E$; writing $e = \{a, b\}$, the vertex $a$ satisfies $a \in \Gamma(E)$ by the membership criterion \texttt{mem\_incidentVertices}.

Moreover, $E \subseteq \operatorname{inducedEdges}(G, S)$ by the lemma \texttt{subset\_inducedEdges\_incidentVertices}.

We consider two cases based on whether $S = V$ or $S \subsetneq V$.

\textbf{Case 1:} $|S| = |V|$ (i.e., $S$ is all of $V$). We must show $\gamma(s, \lambda_2, \tilde{\alpha})\,n \leq |S| = n$, i.e., $\gamma(s,\lambda_2,\tilde{\alpha}) \leq 1$. Unfolding the definition, this is $\sqrt{\lambda_2^2 + 4s(s-\lambda_2)\tilde{\alpha}} - \lambda_2 \leq 2(s-\lambda_2)$, i.e., $\sqrt{\lambda_2^2 + 4s(s-\lambda_2)\tilde{\alpha}} \leq 2s - \lambda_2$. Since $\tilde{\alpha} \leq 1$ (as $|E| \leq |E(G)|$), we have $\lambda_2^2 + 4s(s-\lambda_2)\tilde{\alpha} \leq \lambda_2^2 + 4s(s-\lambda_2) = (2s - \lambda_2)^2$. Taking square roots and using $2s - \lambda_2 \geq 0$, we conclude $\sqrt{\lambda_2^2 + 4s(s-\lambda_2)\tilde{\alpha}} \leq 2s - \lambda_2$.

\textbf{Case 2:} $|S| < |V|$. Let $\gamma' = |S|/n$. Apply the Alon--Chung theorem (\texttt{alonChung}) to $G$ and $S$:
\[
|\operatorname{inducedEdges}(G,S)| \;\leq\; \left(\gamma'^2 + \frac{\lambda_2}{s}\,\gamma'(1-\gamma')\right)|E(G)|.
\]
Since $|E| \leq |\operatorname{inducedEdges}(G,S)|$, we obtain $\tilde{\alpha} \leq \gamma'^2 + \tfrac{\lambda_2}{s}\gamma'(1-\gamma')$.

It remains to show $\gamma(s,\lambda_2,\tilde{\alpha}) \leq \gamma'$. Unfolding, this is
\[
\sqrt{\lambda_2^2 + 4s(s-\lambda_2)\tilde{\alpha}} \;\leq\; \lambda_2 + 2(s-\lambda_2)\gamma'.
\]
Let $R = \lambda_2 + 2(s-\lambda_2)\gamma'$. We verify $R > 0$: since $\tilde{\alpha} > 0$ and $\tilde{\alpha} \leq \alpha(\gamma')$, the quantity $\alpha(\gamma') > 0$, which implies $\gamma' + \tfrac{\lambda_2}{s}(1-\gamma') > 0$, and hence $\lambda_2 + (s-\lambda_2)\gamma' > 0$; adding $(s-\lambda_2)\gamma' > 0$ gives $R > 0$.

Squaring both sides (both nonneg), we need $\lambda_2^2 + 4s(s-\lambda_2)\tilde{\alpha} \leq R^2 = \lambda_2^2 + 4(s-\lambda_2)\lambda_2\gamma' + 4(s-\lambda_2)^2\gamma'^2$, i.e.,
\[
4s(s-\lambda_2)\tilde{\alpha} \;\leq\; 4(s-\lambda_2)\bigl(\lambda_2\gamma' + (s-\lambda_2)\gamma'^2\bigr) = 4(s-\lambda_2)\,s\,\alpha(\gamma').
\]
Dividing by $4(s-\lambda_2) > 0$, this becomes $s\tilde{\alpha} \leq s\,\alpha(\gamma')$, which follows from $\tilde{\alpha} \leq \alpha(\gamma')$ by multiplying by $s > 0$. Taking square roots gives the desired inequality, completing the proof.
\end{proof}

\begin{theorem}[Edge-to-Vertex Expansion]
\label{thm:EdgeToVertexExpansion.edgeToVertexExpansion}
\lean{EdgeToVertexExpansion.edgeToVertexExpansion}
\leanok
\uses{def:AlonChungContrapositive.incidentVertices, lem:AlonChungContrapositive.mem_incidentVertices, lem:AlonChungContrapositive.subset_inducedEdges_incidentVertices, def:EdgeToVertexExpansion.betaParam, def:EdgeToVertexExpansion.gammaOfAlpha, lem:EdgeToVertexExpansion.betaParam_eq, lem:EdgeToVertexExpansion.gammaOfAlpha_concavity, lem:EdgeToVertexExpansion.incidentVertices_ge_gamma, lem:EdgeToVertexExpansion.edgeFinset_eq_half_sn}
Let $G$ be a finite connected $s$-regular simple graph on $V$ with $|V| \geq 2$ and second-largest eigenvalue $\lambda_2 < s$. Let $0 < \alpha \leq 1$ and let $E \subseteq E(G)$ be a set of edges with $|E| \leq \alpha\,|E(G)|$. Then
\[
|\Gamma(E)| \;\geq\; \beta(s, \lambda_2, \alpha)\,|E|,
\]
where $\beta(s, \lambda_2, \alpha) = \dfrac{\sqrt{\lambda_2^2 + 4s(s-\lambda_2)\alpha} - \lambda_2}{s(s-\lambda_2)\alpha}$.
\end{theorem}

\begin{proof}
\leanok
\uses{def:AlonChungContrapositive.incidentVertices, lem:AlonChungContrapositive.mem_incidentVertices, lem:AlonChungContrapositive.subset_inducedEdges_incidentVertices, def:EdgeToVertexExpansion.betaParam, def:EdgeToVertexExpansion.gammaOfAlpha, lem:EdgeToVertexExpansion.betaParam_eq, lem:EdgeToVertexExpansion.gammaOfAlpha_concavity, lem:EdgeToVertexExpansion.incidentVertices_ge_gamma, lem:EdgeToVertexExpansion.edgeFinset_eq_half_sn}
We first handle the edge case $E = \emptyset$: the bound $|\Gamma(\emptyset)| \geq \beta \cdot 0 = 0$ holds trivially by \texttt{simp}.

Assume $E \neq \emptyset$. Let $n = |V|$, $\tilde{\alpha} = |E|/|E(G)|$. We have $0 < \tilde{\alpha} \leq \alpha$ (the upper bound since $|E| \leq \alpha|E(G)|$).

\textbf{Step 1 (lower bound via $\gamma$).} By Lemma~\ref{lem:EdgeToVertexExpansion.incidentVertices_ge_gamma},
\[
|\Gamma(E)| \;\geq\; \gamma(s, \lambda_2, \tilde{\alpha})\,n.
\]

\textbf{Step 2 (concavity).} By Lemma~\ref{lem:EdgeToVertexExpansion.gammaOfAlpha_concavity} (with $\tilde{\alpha} \leq \alpha$),
\[
\gamma(s, \lambda_2, \alpha)\,\tilde{\alpha} \;\leq\; \gamma(s, \lambda_2, \tilde{\alpha})\,\alpha.
\]

\textbf{Step 3 (handshaking).} By Lemma~\ref{lem:EdgeToVertexExpansion.edgeFinset_eq_half_sn}, $|E(G)| = \tfrac{s}{2}n$, so
\[
\tilde{\alpha} \cdot n = \frac{|E|}{|E(G)|}\,n = \frac{|E| \cdot n}{\tfrac{s}{2}n} = \frac{2|E|}{s}.
\]

\textbf{Combining.} Using Lemma~\ref{lem:EdgeToVertexExpansion.betaParam_eq} ($\beta = 2\gamma(\alpha)/(s\alpha)$), we compute:
\begin{align*}
\beta(s,\lambda_2,\alpha)\,|E| &= \frac{2\,\gamma(s,\lambda_2,\alpha)}{s\,\alpha}\,|E| = \frac{\gamma(s,\lambda_2,\alpha)}{\alpha}\cdot\frac{2|E|}{s} = \frac{\gamma(s,\lambda_2,\alpha)}{\alpha}\cdot (\tilde{\alpha}\,n) \\
&= \frac{\gamma(s,\lambda_2,\alpha)\,\tilde{\alpha}}{\alpha}\,n \;\leq\; \frac{\gamma(s,\lambda_2,\tilde{\alpha})\,\alpha}{\alpha}\,n = \gamma(s,\lambda_2,\tilde{\alpha})\,n \;\leq\; |\Gamma(E)|,
\end{align*}
where the first inequality uses Step 2 and $\alpha > 0$, and the last uses Step 1.
\end{proof}

%--- Lem_3: RelativeVertexToEdgeExpansion ---
\chapter{Lem 3: Relative Vertex-to-Edge Expansion}

\begin{definition}[Edge Boundary at Vertex]
\label{def:RelativeVertexToEdgeExpansion.edgeBoundaryAtVertex}
\lean{RelativeVertexToEdgeExpansion.edgeBoundaryAtVertex}
\leanok
\uses{def:CheegerConstant.edgeBoundary}
Let $G = (V, E)$ be a simple graph and $S \subseteq V$. For a vertex $v \in V$, the \emph{edge boundary at $v$} is the set of edges in the edge boundary $\delta S$ that are incident to $v$:
\[
  (\delta S)_v \;:=\; \{ e \in \delta S \mid v \in e \}.
\]
\end{definition}

\begin{definition}[High Expansion Vertices]
\label{def:RelativeVertexToEdgeExpansion.highExpansionVertices}
\lean{RelativeVertexToEdgeExpansion.highExpansionVertices}
\leanok
\uses{def:RelativeVertexToEdgeExpansion.edgeBoundaryAtVertex}
Let $G = (V, E)$ be a simple graph, $S \subseteq V$, $s \in \mathbb{N}$, and $b \in \mathbb{R}$. The set of \emph{high expansion vertices} is
\[
  A \;:=\; \{ v \in S \mid s - b \;\leq\; |(\delta S)_v| \},
\]
i.e., the vertices in $S$ whose local boundary degree is at least $s - b$.
\end{definition}

\begin{lemma}[$A \subseteq S$]
\label{lem:RelativeVertexToEdgeExpansion.highExpansionVertices_subset}
\lean{RelativeVertexToEdgeExpansion.highExpansionVertices_subset}
\leanok
\uses{def:RelativeVertexToEdgeExpansion.highExpansionVertices}
For any graph $G$, subset $S \subseteq V$, and parameters $s \in \mathbb{N}$, $b \in \mathbb{R}$, the set of high expansion vertices satisfies $A \subseteq S$.
\end{lemma}

\begin{proof}
\leanok
\uses{def:RelativeVertexToEdgeExpansion.highExpansionVertices}
Since $A$ is defined as a filtered subset of $S$ (i.e., $A = S.\operatorname{filter}(\cdots)$), the inclusion $A \subseteq S$ follows immediately from the fact that any filtered subset of a finset is contained in that finset.
\end{proof}

\begin{lemma}[Nonemptiness of High Expansion Vertices]
\label{lem:RelativeVertexToEdgeExpansion.highExpansionVertices_nonempty}
\lean{RelativeVertexToEdgeExpansion.highExpansionVertices_nonempty}
\leanok
\uses{def:RelativeVertexToEdgeExpansion.highExpansionVertices}
If there exists a vertex $v \in S$ with $|(\delta S)_v| \geq s - b$, then $A$ is nonempty.
\end{lemma}

\begin{proof}
\leanok
\uses{def:RelativeVertexToEdgeExpansion.highExpansionVertices}
Assume there exist $v \in S$ and a proof $h_{vb}$ that $s - b \leq |(\delta S)_v|$. Decomposing the hypothesis, we obtain the vertex $v$, its membership $v \in S$, and the bound $h_{vb}$. Then $v \in A$ follows directly from the membership criterion $A = S.\operatorname{filter}(s - b \leq |(\delta S)_\cdot|)$, and hence $A$ is nonempty.
\end{proof}

\begin{lemma}[Partition Sum of Edge Boundary]
\label{lem:RelativeVertexToEdgeExpansion.edgeBoundary_partition_sum}
\lean{RelativeVertexToEdgeExpansion.edgeBoundary_partition_sum}
\leanok
\uses{def:CheegerConstant.edgeBoundary, thm:CheegerConstant.mem_edgeBoundary}
For any graph $G$ and subset $S \subseteq V$, the edge boundary decomposes as a disjoint union over vertices in $S$:
\[
  |\delta S| \;=\; \sum_{v \in S} |(\delta S)_v|.
\]
\end{lemma}

\begin{proof}
\leanok
\uses{def:CheegerConstant.edgeBoundary, thm:CheegerConstant.mem_edgeBoundary}
We first show that $\delta S = \bigcup_{v \in S} (\delta S)_v$. For the forward direction: given $e \in \delta S$, by the characterization of the edge boundary (using \texttt{mem\_edgeBoundary}), $e = \{a, b\}$ with $a \in S$ and $b \notin S$; thus $e \in (\delta S)_a$ and $a \in S$, so $e$ lies in the big union. For the reverse direction: if $e \in (\delta S)_v$ for some $v \in S$, then $e \in \delta S$ by the filter condition.

Next, we show the family $\{(\delta S)_v\}_{v \in S}$ is pairwise disjoint. Suppose $e \in (\delta S)_v \cap (\delta S)_w$ with $v \neq w$. Then $e \in \delta S$, and by \texttt{mem\_edgeBoundary}, $e = \{a,b\}$ with $a \in S$, $b \notin S$. Both $v \in e$ and $w \in e$ hold. By the membership condition for $\operatorname{Sym2}$, each of $v, w$ equals either $a$ or $b$. Since $b \notin S$ but $v, w \in S$, both $v = a$ and $w = a$, contradicting $v \neq w$.

Having established the disjoint partition, the cardinality identity follows from the finset big-union cardinality formula for pairwise disjoint families.
\end{proof}

\begin{lemma}[Boundary Degree Bound]
\label{lem:RelativeVertexToEdgeExpansion.edgeBoundaryAtVertex_le_degree}
\lean{RelativeVertexToEdgeExpansion.edgeBoundaryAtVertex_le_degree}
\leanok
\uses{def:RelativeVertexToEdgeExpansion.edgeBoundaryAtVertex, thm:CheegerConstant.edgeBoundary_subset_edgeFinset}
Let $G$ be an $s$-regular graph. For any $S \subseteq V$ and $v \in V$,
\[
  |(\delta S)_v| \;\leq\; s.
\]
\end{lemma}

\begin{proof}
\leanok
\uses{def:RelativeVertexToEdgeExpansion.edgeBoundaryAtVertex, thm:CheegerConstant.edgeBoundary_subset_edgeFinset}
We compute:
\[
  |(\delta S)_v| \;\leq\; |G.\operatorname{incidenceFinset}(v)| \;=\; \deg_G(v) \;=\; s.
\]
The first inequality holds because $(\delta S)_v \subseteq G.\operatorname{incidenceFinset}(v)$: any edge $e \in (\delta S)_v$ lies in the edge boundary $\delta S \subseteq E(G)$ (by \texttt{edgeBoundary\_subset\_edgeFinset}) and satisfies $v \in e$, hence $e \in G.\operatorname{incidenceFinset}(v)$. The equality of the incidence finset cardinality with $\deg_G(v)$ follows from the standard identity, and $\deg_G(v) = s$ by $s$-regularity.
\end{proof}

\begin{lemma}[Relative Cheeger with Non-strict Inequality]
\label{lem:RelativeVertexToEdgeExpansion.relativeCheeger_le}
\lean{RelativeVertexToEdgeExpansion.relativeCheeger_le}
\leanok
\uses{thm:RelativeCheeger.spectralLaplacianBound, def:GraphExpansion.secondLargestEigenvalue}
Let $G$ be a finite, connected, $s$-regular graph on $|V| \geq 2$ vertices with second-largest eigenvalue $\lambda_2$. For $\alpha \in (0,1)$ and $S \subseteq V$ with $0 < |S|$ and $|S| \leq \alpha |V|$,
\[
  (1 - \alpha)(s - \lambda_2)|S| \;\leq\; |\delta S|.
\]
\end{lemma}

\begin{proof}
\leanok
\uses{thm:RelativeCheeger.spectralLaplacianBound, def:GraphExpansion.secondLargestEigenvalue}
We have $|S| > 0$ (given) and $|V| > 0$ (since $|V| \geq 2$). Since $|S| \leq \alpha|V|$ and $\alpha < 1$, we deduce $|S| < |V|$, so $S$ is a proper nonempty subset.

We split into two cases based on the sign of $s - \lambda_2$.

\textbf{Case 1: $s - \lambda_2 \geq 0$.} Apply the spectral Laplacian bound (\texttt{spectralLaplacianBound}):
\[
  (s - \lambda_2)|S|\!\left(1 - \frac{|S|}{|V|}\right) \;\leq\; |\delta S|.
\]
Since $|S|/|V| \leq \alpha$ (as $|S| \leq \alpha|V|$), we have $1 - \alpha \leq 1 - |S|/|V|$. Multiplying both sides of the latter by the nonneg quantity $(s-\lambda_2)|S| \geq 0$ preserves the inequality, yielding
\[
  (1-\alpha)(s-\lambda_2)|S| \;\leq\; (s-\lambda_2)|S|\!\left(1-\frac{|S|}{|V|}\right) \;\leq\; |\delta S|.
\]

\textbf{Case 2: $s - \lambda_2 < 0$.} Then $(1-\alpha)(s-\lambda_2) < 0$ (since $1-\alpha > 0$), so $(1-\alpha)(s-\lambda_2)|S| < 0 \leq |\delta S|$ by linear arithmetic.
\end{proof}

\begin{lemma}[Upper Bound on Edge Boundary]
\label{lem:RelativeVertexToEdgeExpansion.edgeBoundary_upper_bound}
\lean{RelativeVertexToEdgeExpansion.edgeBoundary_upper_bound}
\leanok
\uses{def:RelativeVertexToEdgeExpansion.highExpansionVertices, lem:RelativeVertexToEdgeExpansion.highExpansionVertices_subset, lem:RelativeVertexToEdgeExpansion.edgeBoundary_partition_sum, lem:RelativeVertexToEdgeExpansion.edgeBoundaryAtVertex_le_degree}
Let $G$ be an $s$-regular graph, $S \subseteq V$, $b \in \mathbb{R}$, and $A = \{v \in S \mid |(\delta S)_v| \geq s - b\}$. Then
\[
  |\delta S| \;\leq\; s|A| + (s - b)(|S| - |A|).
\]
\end{lemma}

\begin{proof}
\leanok
\uses{def:RelativeVertexToEdgeExpansion.highExpansionVertices, lem:RelativeVertexToEdgeExpansion.highExpansionVertices_subset, lem:RelativeVertexToEdgeExpansion.edgeBoundary_partition_sum, lem:RelativeVertexToEdgeExpansion.edgeBoundaryAtVertex_le_degree}
Let $B := S \setminus A$. Since $A \subseteq S$, we have the disjoint decomposition $S = A \sqcup B$.

By the partition sum identity (Lemma~\ref{lem:RelativeVertexToEdgeExpansion.edgeBoundary_partition_sum}):
\[
  |\delta S| \;=\; \sum_{v \in S} |(\delta S)_v| \;=\; \sum_{v \in A} |(\delta S)_v| + \sum_{v \in B} |(\delta S)_v|.
\]

For the sum over $A$: since each $|(\delta S)_v| \leq s$ by $s$-regularity (Lemma~\ref{lem:RelativeVertexToEdgeExpansion.edgeBoundaryAtVertex_le_degree}), the bounded sum inequality gives $\sum_{v \in A} |(\delta S)_v| \leq s|A|$.

For the sum over $B$: for any $v \in B = S \setminus A$, we have $v \notin A$, meaning $|(\delta S)_v| < s - b$ (i.e., the defining condition for $A$ fails). Hence $|(\delta S)_v| \leq s - b$ (strict inequality implies $\leq$), and again by the bounded sum inequality, $\sum_{v \in B} |(\delta S)_v| \leq (s-b)|B|$.

Finally, $|B| = |S \setminus A| = |S| - |A|$ (since $A \subseteq S$), obtained by the cardinality identity for set differences. Combining all estimates gives $|\delta S| \leq s|A| + (s-b)(|S| - |A|)$.
\end{proof}

\begin{theorem}[Relative Vertex-to-Edge Expansion]
\label{thm:RelativeVertexToEdgeExpansion.relativeVertexToEdgeExpansion}
\lean{RelativeVertexToEdgeExpansion.relativeVertexToEdgeExpansion}
\leanok
\uses{def:CheegerConstant.edgeBoundary, thm:CheegerConstant.edgeBoundary_subset_edgeFinset, thm:CheegerConstant.mem_edgeBoundary, def:RelativeVertexToEdgeExpansion.highExpansionVertices, lem:RelativeVertexToEdgeExpansion.relativeCheeger_le, lem:RelativeVertexToEdgeExpansion.edgeBoundary_upper_bound}
Let $G$ be a finite, connected, $s$-regular graph on vertex set $V$ with $|V| \geq 2$ and second-largest eigenvalue $\lambda_2$. Let $\alpha \in (0,1)$, $b \in \mathbb{R}$ with $0 < b \leq s$, and $S \subseteq V$ with $|S| \leq \alpha|V|$. Define
\[
  A \;:=\; \bigl\{ v \in S \;\big|\; |(\delta S)_v| \geq s - b \bigr\},
  \qquad
  \beta \;:=\; \frac{(b - \lambda_2) - \alpha(s - \lambda_2)}{b}.
\]
Then
\[
  |A| \;\geq\; \beta \, |S|.
\]
\end{theorem}

\begin{proof}
\leanok
\uses{def:RelativeVertexToEdgeExpansion.highExpansionVertices, lem:RelativeVertexToEdgeExpansion.relativeCheeger_le, lem:RelativeVertexToEdgeExpansion.edgeBoundary_upper_bound}
\textbf{Step 0: Trivial case.} If $|S| = 0$, then both sides are zero and the inequality holds trivially by simplification. Henceforth assume $|S| > 0$.

\textbf{Step 1: Lower bound on $|\delta S|$.} By Lemma~\ref{lem:RelativeVertexToEdgeExpansion.relativeCheeger_le} (relative Cheeger with $\leq$):
\[
  (1-\alpha)(s - \lambda_2)|S| \;\leq\; |\delta S|.
\]

\textbf{Step 2: Upper bound on $|\delta S|$.} By Lemma~\ref{lem:RelativeVertexToEdgeExpansion.edgeBoundary_upper_bound}:
\[
  |\delta S| \;\leq\; s|A| + (s-b)(|S| - |A|).
\]

\textbf{Step 3: Combine.} From Steps 1 and 2 by linear arithmetic:
\[
  (1-\alpha)(s-\lambda_2)|S| \;\leq\; s|A| + (s-b)(|S| - |A|).
\]
Expanding the right side: $s|A| + (s-b)|S| - (s-b)|A| = b|A| + (s-b)|S|$. Hence:
\[
  (1-\alpha)(s-\lambda_2)|S| - (s-b)|S| \;\leq\; b|A|.
\]

\textbf{Step 4: Simplify the coefficient.} We compute by ring arithmetic:
\[
  (1-\alpha)(s-\lambda_2) - (s-b) \;=\; (b - \lambda_2) - \alpha(s-\lambda_2).
\]
Thus:
\[
  \bigl[(b-\lambda_2) - \alpha(s-\lambda_2)\bigr]|S| \;\leq\; b|A|.
\]

\textbf{Step 5: Divide by $b$.} Since $b > 0$, dividing both sides by $b$ gives:
\[
  \frac{(b-\lambda_2) - \alpha(s-\lambda_2)}{b} \cdot |S| \;\leq\; |A|,
\]
which is the desired conclusion. This final step is performed by rewriting with $\operatorname{div_mul_eq_mul_div}$, applying $\operatorname{div_le_iff}$ with $b > 0$, and concluding by linear arithmetic.
\end{proof}

%--- Thm_7: SipserSpielmanExpanderCodeDistance ---
\chapter{Thm 7: Sipser--Spielman Expander Code Distance}

\begin{definition}[Pi Submodule Equivalence]
\label{def:SipserSpielman.piSubmoduleEquiv}
\lean{SipserSpielman.piSubmoduleEquiv}
\leanok
\uses{def:Labeling}
Let $R$ be a commutative semiring, $M$ an $R$-module, and $\iota$ an index type. Given a submodule $p \leq M$, there is a linear equivalence
\[
\operatorname{Submodule.pi}(\operatorname{Set.univ},\, \lambda\,\_\mapsto p) \;\simeq_R\; (\iota \to p).
\]
The forward map sends $x : \iota \to M$ (with $x_i \in p$ for all $i$) to the function $i \mapsto \langle x_i,\, x.\text{prop}\, i\rangle$, and the inverse sends $f : \iota \to p$ to $\langle \lambda\,i \mapsto f_i,\, \lambda\,i\,\_ \mapsto (f_i).\text{prop}\rangle$.
\end{definition}

\begin{lemma}[Finrank of Pi Submodule]
\label{lem:SipserSpielman.finrank_pi_submodule}
\lean{SipserSpielman.finrank_pi_submodule}
\leanok
\uses{def:SipserSpielman.piSubmoduleEquiv}
Let $R$ be a commutative semiring satisfying the strong rank condition, $M$ an $R$-module, $\iota$ a finite index type, and $p \leq M$ a free finite-rank submodule. Then
\[
\operatorname{finrank}_R\bigl(\operatorname{Submodule.pi}(\operatorname{Set.univ},\, \lambda\,\_\mapsto p)\bigr) = |\iota| \cdot \operatorname{finrank}_R(p).
\]
\end{lemma}

\begin{proof}
\leanok
\uses{def:SipserSpielman.piSubmoduleEquiv}
By the linear equivalence $\operatorname{piSubmoduleEquiv}$, the finrank of the pi submodule equals the finrank of $\iota \to p$. Applying the formula $\operatorname{finrank}(\iota \to p) = |\iota| \cdot \operatorname{finrank}(p)$ and simplifying with $|\operatorname{Fin}\,n| = n$ yields the result.
\end{proof}

\begin{lemma}[Range of Differential Contained in Pi of Ranges]
\label{lem:SipserSpielman.range_differential_le}
\lean{SipserSpielman.range_differential_le}
\leanok
\uses{def:Labeling, def:differential, def:localView}
Let $G$ be a simple graph on $V$, $\Lambda$ a labeling of $G$ with degree $s$, and $\operatorname{parityCheck} : (\mathbb{F}_2^s \to_{\mathbb{F}_2} \mathbb{F}_2^m)$ a linear map. Then
\[
\operatorname{range}(\partial_\Lambda^{\operatorname{parityCheck}}) \;\leq\; \operatorname{Submodule.pi}\!\Bigl(\operatorname{Set.univ},\;\lambda\,v \mapsto \operatorname{range}(\operatorname{parityCheck})\Bigr).
\]
\end{lemma}

\begin{proof}
\leanok
\uses{def:differential, def:localView}
Let $f$ be in the range of $\partial_\Lambda^{\operatorname{parityCheck}}$. We must show $f \in \operatorname{Submodule.pi}(\operatorname{Set.univ}, \lambda\,v \mapsto \operatorname{range}(\operatorname{parityCheck}))$, i.e., for every $v \in V$, the value $f(v)$ lies in $\operatorname{range}(\operatorname{parityCheck})$. Since $f$ is in the range, there exists $c$ with $\partial_\Lambda^{\operatorname{parityCheck}}(c) = f$. For any vertex $v$, we have $f(v) = \operatorname{parityCheck}(\operatorname{localView}_\Lambda\,v\,c)$, so $f(v)$ is the image of $\operatorname{localView}_\Lambda\,v\,c$ under $\operatorname{parityCheck}$, hence in its range.
\end{proof}

\begin{lemma}[Upper Bound on Finrank of Range of Differential]
\label{lem:SipserSpielman.finrank_range_differential_le}
\lean{SipserSpielman.finrank_range_differential_le}
\leanok
\uses{def:Labeling, def:differential, lem:SipserSpielman.range_differential_le, lem:SipserSpielman.finrank_pi_submodule}
Let $G$, $\Lambda$, and $\operatorname{parityCheck}$ be as above. Then
\[
\operatorname{finrank}_{\mathbb{F}_2}\bigl(\operatorname{range}(\partial_\Lambda^{\operatorname{parityCheck}})\bigr) \;\leq\; |V| \cdot \operatorname{finrank}_{\mathbb{F}_2}\bigl(\operatorname{range}(\operatorname{parityCheck})\bigr).
\]
\end{lemma}

\begin{proof}
\leanok
\uses{lem:SipserSpielman.range_differential_le, lem:SipserSpielman.finrank_pi_submodule}
We compute:
\begin{align*}
\operatorname{finrank}_{\mathbb{F}_2}\bigl(\operatorname{range}(\partial_\Lambda^{\operatorname{parityCheck}})\bigr)
&\leq \operatorname{finrank}_{\mathbb{F}_2}\Bigl(\operatorname{Submodule.pi}\bigl(\operatorname{Set.univ},\,\lambda\,v\mapsto\operatorname{range}(\operatorname{parityCheck})\bigr)\Bigr) \\
&= |V| \cdot \operatorname{finrank}_{\mathbb{F}_2}\bigl(\operatorname{range}(\operatorname{parityCheck})\bigr).
\end{align*}
The first inequality follows from $\operatorname{Submodule.finrank_mono}$ applied to the containment $\operatorname{range_differential_le}$. The equality follows from $\operatorname{finrank_pi_submodule}$.
\end{proof}

\begin{theorem}[Sipser--Spielman Dimension Bound]
\label{thm:SipserSpielman.sipserSpielman_dimension_bound}
\lean{SipserSpielman.sipserSpielman_dimension_bound}
\leanok
\uses{def:ClassicalCode, def:Labeling, def:differential, def:tannerCode, lem:SipserSpielman.finrank_range_differential_le}
Let $G$ be a simple graph on $V$ that is $s$-regular with $s \geq 1$ and $|V| \geq 2$. Let $\Lambda$ be a labeling of $G$ with degree $s$, and let $\operatorname{parityCheck} : \mathbb{F}_2^s \to \mathbb{F}_2^m$ be a linear map with $k_L = \operatorname{finrank}_{\mathbb{F}_2}(\ker(\operatorname{parityCheck}))$. Then the Tanner code $C(G, \Lambda, \operatorname{parityCheck})$ satisfies
\[
\dim\bigl(C(G,\Lambda,\operatorname{parityCheck})\bigr) \;\geq\; \Bigl(\frac{2k_L}{s} - 1\Bigr)\cdot|X_1|,
\]
where $|X_1| = |E(G)|$ denotes the number of edges.
\end{theorem}

\begin{proof}
\leanok
\uses{def:ClassicalCode, def:Labeling, def:differential, def:tannerCode, lem:SipserSpielman.finrank_range_differential_le, lem:AlonChung.twice_card_edgeFinset_eq}
By definition, $\operatorname{tannerCode}\,\Lambda\,\operatorname{parityCheck} = \ker(\partial_\Lambda^{\operatorname{parityCheck}})$.

\textbf{Step 1 (Rank-nullity for $\partial$).} By rank-nullity for the differential:
\[
\operatorname{finrank}(\operatorname{range}(\partial)) + \operatorname{finrank}(\ker(\partial)) = \operatorname{finrank}(G.\operatorname{edgeSet} \to \mathbb{F}_2) = |X_1|.
\]

\textbf{Step 2 (Rank-nullity for $\operatorname{parityCheck}$).} We have $\operatorname{finrank}(\operatorname{range}(\operatorname{parityCheck})) + k_L = s$, so $\operatorname{finrank}(\operatorname{range}(\operatorname{parityCheck})) = s - k_L$. In particular $k_L \leq s$.

\textbf{Step 3 (Bounding the range).} By $\operatorname{finrank_range_differential_le}$ and Step 2:
\[
\operatorname{finrank}(\operatorname{range}(\partial)) \leq |V|\cdot(s - k_L).
\]

\textbf{Step 4 (Lower bound on dimension).} Combining Steps 1 and 3 by integer arithmetic:
\[
\dim\bigl(\operatorname{tannerCode}\bigr) \geq |X_1| - |V|\cdot(s - k_L).
\]

\textbf{Step 5 (Handshaking).} By the handshaking lemma $\operatorname{twice_card_edgeFinset_eq}$, we have $2|X_1| = s\cdot|V|$, so $|V| = 2|X_1|/s$.

\textbf{Step 6 (Real arithmetic).} Rewriting Step 4 in $\mathbb{R}$ and substituting the handshaking relation, we verify $\dim(\operatorname{tannerCode}) - (2k_L/s - 1)\cdot|X_1| \geq 0$ via the identity
\[
\Bigl(\frac{2k_L}{s}-1\Bigr)|X_1| = |X_1| - (s-k_L)\cdot|V|,
\]
which holds since $(s-k_L)\cdot|V| = (s-k_L)\cdot 2|X_1|/s$. The inequality then follows by nonlinear arithmetic using $s > 0$.
\end{proof}

\begin{definition}[Support of an Edge Function]
\label{def:SipserSpielman.support}
\lean{SipserSpielman.support}
\leanok
\uses{def:Labeling}
Let $G$ be a simple graph on $V$ and $c : G.\operatorname{edgeSet} \to \mathbb{F}_2$. The \emph{support} of $c$ is the finite set
\[
\operatorname{supp}(c) \;:=\; \{e \in G.\operatorname{edgeSet} \mid c(e) \neq 0\}.
\]
\end{definition}

\begin{definition}[Edge Hamming Weight]
\label{def:SipserSpielman.edgeHammingWeight}
\lean{SipserSpielman.edgeHammingWeight}
\leanok
\uses{def:SipserSpielman.support, def:hammingWeight}
The \emph{edge Hamming weight} of $c : G.\operatorname{edgeSet} \to \mathbb{F}_2$ is
\[
\operatorname{wt}(c) \;:=\; |\operatorname{supp}(c)|.
\]
\end{definition}

\begin{definition}[Vertices Incident to Support]
\label{def:SipserSpielman.incidentToSupport}
\lean{SipserSpielman.incidentToSupport}
\leanok
\uses{def:SipserSpielman.support}
Let $c : G.\operatorname{edgeSet} \to \mathbb{F}_2$. The set of vertices \emph{incident to the support} of $c$ is
\[
S(c) \;:=\; \{v \in V \mid \exists\, e \in \operatorname{supp}(c),\; v \in e\}.
\]
\end{definition}

\begin{definition}[Support Degree at a Vertex]
\label{def:SipserSpielman.supportDegree}
\lean{SipserSpielman.supportDegree}
\leanok
\uses{def:Labeling, def:SipserSpielman.support}
Let $\Lambda$ be a labeling of $G$ with degree $s$ and $c : G.\operatorname{edgeSet} \to \mathbb{F}_2$. The \emph{support degree} of $c$ at vertex $v$ is
\[
\deg_{\operatorname{supp}}(v) \;:=\; \bigl|\{i \in \operatorname{Fin}(s) \mid c\bigl(\langle s(v,\,(\Lambda_v)^{-1}(i)),\,\ldots\rangle\bigr) \neq 0\}\bigr|,
\]
the number of edges in $\operatorname{supp}(c)$ incident to $v$.
\end{definition}

\begin{lemma}[Local View Weight Bound]
\label{lem:SipserSpielman.localView_weight_bound}
\lean{SipserSpielman.localView_weight_bound}
\leanok
\uses{def:ClassicalCode, def:Labeling, def:localView, def:tannerCode, def:hammingWeight, lem:localView_apply, thm:mem_tannerCode_iff}
Let $c \in C(G, \Lambda, \operatorname{parityCheck})$ be a codeword, $d_L = \operatorname{dist}(C_L)$ the local code distance, and $v \in V$ a vertex such that $\operatorname{localView}_\Lambda(v, c) \neq 0$. Then
\[
d_L \;\leq\; \operatorname{wt}(\operatorname{localView}_\Lambda(v,c)).
\]
\end{lemma}

\begin{proof}
\leanok
\uses{def:ClassicalCode, def:localView, def:tannerCode, thm:mem_tannerCode_iff}
Rewriting using $d_L = \operatorname{dist}(\operatorname{ClassicalCode.ofParityCheck}(\operatorname{parityCheck}))$. Since $c \in \operatorname{tannerCode}\,\Lambda\,\operatorname{parityCheck}$, by $\operatorname{mem_tannerCode_iff}$, the local view $\operatorname{localView}_\Lambda\,v\,c$ lies in $\ker(\operatorname{parityCheck})$, which equals the local code $C_L$ by $\operatorname{code_eq_ker}$. Since the local view is nonzero and is a codeword of $C_L$, its Hamming weight witnesses that $d_L = \operatorname{sInf}\{\operatorname{wt}(x) \mid x \in C_L,\; x \neq 0\} \leq \operatorname{wt}(\operatorname{localView}_\Lambda\,v\,c)$.
\end{proof}

\begin{lemma}[Nonemptiness of Support]
\label{lem:SipserSpielman.support_nonempty_of_ne_zero}
\lean{SipserSpielman.support_nonempty_of_ne_zero}
\leanok
\uses{def:SipserSpielman.support}
If $c : G.\operatorname{edgeSet} \to \mathbb{F}_2$ is nonzero, then $\operatorname{supp}(c)$ is nonempty.
\end{lemma}

\begin{proof}
\leanok
\uses{def:SipserSpielman.support}
By contraposition: suppose $\operatorname{supp}(c) = \emptyset$. Then for every edge $e$, $e \notin \operatorname{supp}(c)$, which by definition of support and membership in the filter means $c(e) = 0$. By function extensionality, $c = 0$, contradicting the hypothesis $c \neq 0$.
\end{proof}

\begin{lemma}[Positive Edge Hamming Weight of Nonzero Vector]
\label{lem:SipserSpielman.edgeHammingWeight_pos_of_ne_zero}
\lean{SipserSpielman.edgeHammingWeight_pos_of_ne_zero}
\leanok
\uses{def:SipserSpielman.edgeHammingWeight, lem:SipserSpielman.support_nonempty_of_ne_zero}
If $c \neq 0$, then $\operatorname{wt}(c) > 0$.
\end{lemma}

\begin{proof}
\leanok
\uses{lem:SipserSpielman.support_nonempty_of_ne_zero}
By $\operatorname{support_nonempty_of_ne_zero}$, the set $\operatorname{supp}(c)$ is nonempty, so its cardinality is positive.
\end{proof}

\begin{definition}[Support Edges as Elements of $\operatorname{Sym2}(V)$]
\label{def:SipserSpielman.supportEdges}
\lean{SipserSpielman.supportEdges}
\leanok
\uses{def:SipserSpielman.support}
The \emph{support edges} of $c$ viewed in $\operatorname{Sym2}(V)$ are
\[
\operatorname{suppEdges}(c) \;:=\; \{e.\operatorname{val} \mid e \in \operatorname{supp}(c)\} \;\subseteq\; \operatorname{Sym2}(V).
\]
\end{definition}

\begin{lemma}[Support Edges Subset of Edge Finset]
\label{lem:SipserSpielman.supportEdges_subset}
\lean{SipserSpielman.supportEdges_subset}
\leanok
\uses{def:SipserSpielman.supportEdges}
$\operatorname{suppEdges}(c) \subseteq G.\operatorname{edgeFinset}$.
\end{lemma}

\begin{proof}
\leanok
\uses{def:SipserSpielman.supportEdges}
For any $e \in \operatorname{suppEdges}(c)$, by definition there exists a support element $e' \in \operatorname{supp}(c)$ with $e = e'.\operatorname{val}$. Since $e'$ is an element of $G.\operatorname{edgeSet}$, we have $e'.\operatorname{val} \in G.\operatorname{edgeSet}$, so $e \in G.\operatorname{edgeFinset}$.
\end{proof}

\begin{lemma}[Cardinality of Support Edges Equals Edge Hamming Weight]
\label{lem:SipserSpielman.card_supportEdges_eq}
\lean{SipserSpielman.card_supportEdges_eq}
\leanok
\uses{def:SipserSpielman.supportEdges, def:SipserSpielman.edgeHammingWeight}
$|\operatorname{suppEdges}(c)| = \operatorname{wt}(c)$.
\end{lemma}

\begin{proof}
\leanok
\uses{def:SipserSpielman.supportEdges, def:SipserSpielman.edgeHammingWeight}
By definition, $\operatorname{suppEdges}(c)$ is the image of $\operatorname{supp}(c)$ under the map $e \mapsto e.\operatorname{val}$. Since subtype coercion $\operatorname{val}$ is injective on $G.\operatorname{edgeSet}$, this map is injective on $\operatorname{supp}(c)$. Therefore $|\operatorname{suppEdges}(c)| = |\operatorname{supp}(c)| = \operatorname{wt}(c)$.
\end{proof}

\begin{lemma}[Nonzero Local View at Incident Vertices]
\label{lem:SipserSpielman.localView_ne_zero_of_incident}
\lean{SipserSpielman.localView_ne_zero_of_incident}
\leanok
\uses{def:Labeling, def:localView, def:SipserSpielman.incidentToSupport, lem:localView_apply}
If $v \in S(c)$, then $\operatorname{localView}_\Lambda(v,c) \neq 0$.
\end{lemma}

\begin{proof}
\leanok
\uses{def:localView, def:SipserSpielman.incidentToSupport, lem:localView_apply}
Since $v \in S(c)$, by definition of $\operatorname{incidentToSupport}$, there exists an edge $e$ with $c(e) \neq 0$ and $v \in e$. Suppose for contradiction that $\operatorname{localView}_\Lambda\,v\,c = 0$. Then every component $(\operatorname{localView}_\Lambda\,v\,c)(i) = 0$ for all $i \in \operatorname{Fin}(s)$.

Since $e \in G.\operatorname{edgeSet}$ and $v \in e$ (as an element of $\operatorname{Sym2}(V)$), let $w = \operatorname{Sym2.Mem.other}(v \in e)$, so that $s(v,w) = e$ by $\operatorname{Sym2.other_spec}$. Then $G.\operatorname{Adj}\,v\,w$, so $w \in G.\operatorname{neighborSet}(v)$. Let $i = \Lambda_v(w) \in \operatorname{Fin}(s)$.

Applying $\operatorname{localView_apply}$, we get $(\operatorname{localView}_\Lambda\,v\,c)(i) = c(\langle s(v, (\Lambda_v)^{-1}(i)),\ldots\rangle)$. Since $(\Lambda_v)^{-1}(i) = (\Lambda_v)^{-1}(\Lambda_v(w)) = w$ by the invertibility of $\Lambda_v$, this equals $c(\langle s(v,w),\ldots\rangle) = c(e) \neq 0$, contradicting the assumption that the local view is zero.
\end{proof}

\begin{lemma}[Each Incident Vertex Has at Least $d_L$ Support Edges]
\label{lem:SipserSpielman.incident_vertex_degree_ge_dL}
\lean{SipserSpielman.incident_vertex_degree_ge_dL}
\leanok
\uses{def:ClassicalCode, def:Labeling, def:tannerCode, def:SipserSpielman.supportDegree, def:SipserSpielman.incidentToSupport, lem:SipserSpielman.localView_ne_zero_of_incident, lem:SipserSpielman.localView_weight_bound}
For every $v \in S(c)$ and every codeword $c \in C(G, \Lambda, \operatorname{parityCheck})$:
\[
d_L \;\leq\; \deg_{\operatorname{supp}}(v).
\]
\end{lemma}

\begin{proof}
\leanok
\uses{lem:SipserSpielman.localView_ne_zero_of_incident, lem:SipserSpielman.localView_weight_bound}
By $\operatorname{localView_ne_zero_of_incident}$, since $v \in S(c)$ we have $\operatorname{localView}_\Lambda\,v\,c \neq 0$. By $\operatorname{localView_weight_bound}$, $d_L \leq \operatorname{wt}(\operatorname{localView}_\Lambda\,v\,c)$. Converting between the count of nonzero entries of the local view and the support degree at $v$ closes the goal.
\end{proof}

\begin{lemma}[At Most Two Vertices Incident to Each Edge]
\label{lem:SipserSpielman.card_filter_mem_sym2_le_two}
\lean{SipserSpielman.card_filter_mem_sym2_le_two}
\leanok
\uses{def:Labeling}
For any $e \in \operatorname{Sym2}(V)$:
\[
|\{v \in V \mid v \in e\}| \;\leq\; 2.
\]
\end{lemma}

\begin{proof}
\leanok

By induction on $e$ using $\operatorname{Sym2.ind}$, write $e = s(a,b)$. The filter $\{v \in V \mid v \in s(a,b)\}$ is contained in $\{a,b\}$, so its cardinality is at most $|\{a,b\}|$. Using $\operatorname{card_insert_le}$, $|\{a,b\}| \leq |\{b\}| + 1 = 2$. Integer arithmetic closes the goal.
\end{proof}

\begin{lemma}[Double-Counting: $|S| \cdot d_L \leq 2|E|$]
\label{lem:SipserSpielman.incidentCount_le_twice_weight}
\lean{SipserSpielman.incidentCount_le_twice_weight}
\leanok
\uses{def:Labeling, def:tannerCode, def:SipserSpielman.edgeHammingWeight, def:SipserSpielman.incidentToSupport, lem:SipserSpielman.localView_ne_zero_of_incident, lem:SipserSpielman.localView_weight_bound, lem:SipserSpielman.card_filter_mem_sym2_le_two, lem:localView_apply}
For any codeword $c \in C(G, \Lambda, \operatorname{parityCheck})$ with local code distance $d_L$:
\[
|S(c)| \cdot d_L \;\leq\; 2\cdot\operatorname{wt}(c).
\]
\end{lemma}

\begin{proof}
\leanok
\uses{lem:SipserSpielman.localView_ne_zero_of_incident, lem:SipserSpielman.localView_weight_bound, lem:SipserSpielman.card_filter_mem_sym2_le_two, lem:localView_apply}
Define the counting set $T = \{(v, i) \in S(c) \times \operatorname{Fin}(s) \mid (\operatorname{localView}_\Lambda\,v\,c)(i) \neq 0\}$.

\textbf{Lower bound on $|T|$.} For each $v \in S(c)$, by $\operatorname{localView_ne_zero_of_incident}$ we have $\operatorname{localView}_\Lambda\,v\,c \neq 0$, and by $\operatorname{localView_weight_bound}$, the number of nonzero entries of $\operatorname{localView}_\Lambda\,v\,c$ is at least $d_L$. By $\operatorname{card_nsmul_le_sum}$:
\[
|S(c)| \cdot d_L \;\leq\; |T|.
\]

\textbf{Upper bound on $|T|$.} Define the edge map $\varphi : (v, i) \mapsto \langle s(v, (\Lambda_v)^{-1}(i)),\ldots\rangle \in G.\operatorname{edgeSet}$. For each $(v,i) \in T$, the element $\varphi(v,i)$ lies in $\operatorname{supp}(c)$ by $\operatorname{localView_apply}$.

Each support edge $e \in \operatorname{supp}(c)$ has at most 2 preimages under $\varphi$ in $T$: for preimages $(v_1, i_1)$ and $(v_2, i_2)$ with $v_1 = v_2$, equality $\varphi(v_1, i_1) = \varphi(v_1, i_2)$ yields $s(v_1, (\Lambda_{v_1})^{-1}(i_1)) = s(v_1, (\Lambda_{v_1})^{-1}(i_2))$. Applying $\operatorname{Sym2.eq_iff}$ gives either $(\Lambda_{v_1})^{-1}(i_1).\operatorname{val} = (\Lambda_{v_1})^{-1}(i_2).\operatorname{val}$ (so $i_1 = i_2$ by injectivity of $(\Lambda_{v_1})^{-1}$), or $v_1 = (\Lambda_{v_1})^{-1}(i_2).\operatorname{val}$ (impossible since $G$ is loopless). So each fiber has size at most 1 per distinct first coordinate. The first coordinates project into $\{v \mid v \in e\}$, which has at most 2 elements by $\operatorname{card_filter_mem_sym2_le_two}$.

Applying $\operatorname{card_le_mul_card_image}$:
\[
|T| \;\leq\; 2 \cdot |\operatorname{supp}(c)| \;=\; 2\cdot\operatorname{wt}(c).
\]
Combining the bounds gives $|S(c)| \cdot d_L \leq |T| \leq 2\cdot\operatorname{wt}(c)$.
\end{proof}

\begin{lemma}[Hamming Weight at Most Length]
\label{lem:SipserSpielman.hammingWeight_le_card}
\lean{SipserSpielman.hammingWeight_le_card}
\leanok
\uses{def:hammingWeight}
For any $x : \operatorname{Fin}(n) \to \mathbb{F}_2$:
\[
\operatorname{wt}(x) \;\leq\; n.
\]
\end{lemma}

\begin{proof}
\leanok
\uses{def:hammingWeight}
Unfolding the definition, $\operatorname{wt}(x) = |\{i \in \operatorname{Fin}(n) \mid x_i \neq 0\}|$, which is at most $|\operatorname{Fin}(n)| = n$ by $\operatorname{card_filter_le}$ and $\operatorname{Fintype.card_fin}$.
\end{proof}

\begin{lemma}[Code Distance at Most Length]
\label{lem:SipserSpielman.distance_le_length'}
\lean{SipserSpielman.distance_le_length'}
\leanok
\uses{def:ClassicalCode, def:ClassicalCode.distance, lem:SipserSpielman.hammingWeight_le_card}
Let $\operatorname{parityCheck} : \mathbb{F}_2^s \to \mathbb{F}_2^m$ be a linear map whose associated local code has positive distance. Then
\[
\operatorname{dist}(\operatorname{ClassicalCode.ofParityCheck}(\operatorname{parityCheck})) \;\leq\; s.
\]
\end{lemma}

\begin{proof}
\leanok
\uses{def:ClassicalCode, lem:SipserSpielman.hammingWeight_le_card}
Since the distance is positive, the set $\{\operatorname{wt}(x) \mid x \in C_L,\; x \neq 0\}$ is nonempty. We obtain $x \in C_L$ with $x \neq 0$ such that $\operatorname{dist}(C_L) \leq \operatorname{wt}(x)$. By $\operatorname{hammingWeight_le_card}$, $\operatorname{wt}(x) \leq s$. Hence $\operatorname{dist}(C_L) \leq s$.
\end{proof}

\begin{theorem}[Sipser--Spielman Distance Bound]
\label{thm:SipserSpielman.sipserSpielman_distance_bound}
\lean{SipserSpielman.sipserSpielman_distance_bound}
\leanok
\uses{def:ClassicalCode, def:Labeling, def:differential, def:localView, def:tannerCode, lem:localView_apply, thm:mem_tannerCode_iff, lem:SipserSpielman.incidentCount_le_twice_weight, lem:SipserSpielman.distance_le_length', thm:AlonChungContrapositive.alonChungContrapositive_direct, lem:AlonChung.twice_card_edgeFinset_eq}
Let $G$ be a connected $s$-regular simple graph on $V$ with $s \geq 1$ and $|V| \geq 2$. Let $\Lambda$ be a labeling of $G$ with degree $s$, $\operatorname{parityCheck} : \mathbb{F}_2^s \to \mathbb{F}_2^m$ a linear map, $d_L = \operatorname{dist}(C_L)$ the local code distance, and $\lambda_2$ the second largest eigenvalue of $G$ with $\lambda_2 \geq 0$ and $d_L > \lambda_2$. Then every nonzero codeword $c \in C(G, \Lambda, \operatorname{parityCheck})$ satisfies
\[
\operatorname{wt}(c) \;\geq\; \frac{(d_L - \lambda_2)\, d_L}{(s - \lambda_2)\, s} \cdot |X_1|,
\]
where $|X_1| = |E(G)|$.
\end{theorem}

\begin{proof}
\leanok
\uses{def:ClassicalCode, def:Labeling, def:tannerCode, lem:localView_apply, thm:mem_tannerCode_iff, lem:SipserSpielman.incidentCount_le_twice_weight, lem:SipserSpielman.distance_le_length', lem:SipserSpielman.supportEdges_subset, lem:SipserSpielman.card_supportEdges_eq, thm:AlonChungContrapositive.alonChungContrapositive_direct, lem:AlonChung.twice_card_edgeFinset_eq}
Let $c \in C(G, \Lambda, \operatorname{parityCheck})$ be nonzero. Set $S = S(c)$, $E = \operatorname{wt}(c)$, $n = |V|$.

\textbf{Positivity.} Since $d_L > \lambda_2 \geq 0$ and $d_L \in \mathbb{N}$, we have $d_L > 0$. Since $d_L \leq s$ by $\operatorname{distance_le_length'}$, we get $s - \lambda_2 \geq d_L - \lambda_2 > 0$.

\textbf{Key inequality (double counting).} By $\operatorname{incidentCount_le_twice_weight}$:
\[
|S| \cdot d_L \;\leq\; 2E.
\]

\textbf{Nonemptiness of $S$.} Since $c \neq 0$, there exists an edge $e$ with $c(e) \neq 0$; any endpoint of $e$ belongs to $S$, so $S \neq \emptyset$.

\textbf{Handshaking.} By $\operatorname{twice_card_edgeFinset_eq}$: $2|X_1| = s \cdot n$.

\textbf{Case 1: $|S| \geq n$ (all vertices are incident).} Then $n \cdot d_L \leq |S| \cdot d_L \leq 2E$, so $E \geq n d_L/2$. By handshaking, $n d_L / 2 = (d_L/s) \cdot |X_1|$. Since $(d_L - \lambda_2)/(s-\lambda_2) \leq 1$, we have $(d_L/s) \geq (d_L - \lambda_2)d_L/((s-\lambda_2)s)$. Hence $E \geq (d_L - \lambda_2)d_L/((s-\lambda_2)s) \cdot |X_1|$.

\textbf{Case 2: $|S| < n$.} Set $\gamma = |S|/n \in (0,1)$.

By $\operatorname{alonChungContrapositive_direct}$ applied to the support edges $\operatorname{suppEdges}(c) \subseteq G.\operatorname{edgeFinset}$ with expansion parameter $\gamma$ (note $|\operatorname{incidentVertices}(\operatorname{suppEdges}(c))| \leq |S| \leq \gamma n$):
\[
|\operatorname{suppEdges}(c)| \;\leq\; \Bigl(\gamma^2 + \frac{\lambda_2}{s}\gamma(1-\gamma)\Bigr)|X_1|.
\]
By $\operatorname{card_supportEdges_eq}$, $|\operatorname{suppEdges}(c)| = E$, so:
\[
E \;\leq\; \Bigl(\gamma^2 + \frac{\lambda_2}{s}\gamma(1-\gamma)\Bigr)|X_1|.
\]

Combining the double counting bound $\gamma n \cdot d_L \leq 2E$ with $|X_1| = sn/2$:
\[
\gamma \cdot d_L \;\leq\; \Bigl(\gamma^2 + \frac{\lambda_2}{s}\gamma(1-\gamma)\Bigr)\cdot s.
\]
Dividing by $\gamma > 0$:
\[
d_L \;\leq\; \gamma\bigl(s - \lambda_2\bigr) + \lambda_2,
\]
hence $d_L - \lambda_2 \leq \gamma(s - \lambda_2)$, and since $s - \lambda_2 > 0$:
\[
\gamma \;\geq\; \frac{d_L - \lambda_2}{s - \lambda_2}.
\]

Finally:
\begin{align*}
E &\geq \frac{\gamma \cdot n \cdot d_L}{2} \geq \frac{(d_L-\lambda_2)/(s-\lambda_2) \cdot n \cdot d_L}{2} = \frac{(d_L-\lambda_2)\,d_L}{(s-\lambda_2)\,s} \cdot \frac{sn}{2} = \frac{(d_L-\lambda_2)\,d_L}{(s-\lambda_2)\,s} \cdot |X_1|,
\end{align*}
where the last equality uses handshaking. Real arithmetic closes the goal.
\end{proof}

%--- Thm_8: ExpanderViolatedChecks ---
\chapter{Thm 8: Expander Violated Checks}

\begin{definition}[Edge-to-Vertex Expansion Parameter $\beta'$]
\label{def:ExpanderViolatedChecks.beta'}
\lean{ExpanderViolatedChecks.beta'}
\leanok
\uses{def:EdgeToVertexExpansion.betaParam}
The edge-to-vertex expansion parameter $\beta'$ is defined as
\[
\beta'(s, \lambda_2, \alpha) \;=\; \frac{\sqrt{\lambda_2^2 + 4s(s - \lambda_2)\alpha} - \lambda_2}{s(s - \lambda_2)\alpha},
\]
which is equal to $\operatorname{betaParam}(s, \lambda_2, \alpha)$.
\end{definition}

\begin{definition}[Vertex-to-Edge-Boundary Expansion Parameter $\beta''$]
\label{def:ExpanderViolatedChecks.beta''}
\lean{ExpanderViolatedChecks.beta''}
\leanok
\uses{def:RelativeVertexToEdgeExpansion.edgeBoundaryAtVertex}
The vertex-to-edge-boundary expansion parameter $\beta''$ is defined as
\[
\beta''(s, \lambda_2, \alpha, d_L) \;=\; \frac{(d_L - 1 - \lambda_2) - \alpha s(s - \lambda_2)}{d_L - 1}.
\]
This arises from applying the relative vertex-to-edge expansion lemma with $b = d_L - 1$ and $\alpha' = \alpha s$, ensuring that all high-expansion vertices have violated local checks.
\end{definition}

\begin{definition}[Differential Weight]
\label{def:ExpanderViolatedChecks.differentialWeight}
\lean{ExpanderViolatedChecks.differentialWeight}
\leanok
\uses{def:Labeling, def:SipserSpielman.edgeHammingWeight}
Let $G$ be a simple graph, $\Lambda$ a labeling of $G$ of degree $s$, and $\operatorname{parityCheck} : \mathbb{F}_2^s \to \mathbb{F}_2^m$ a linear map. The \emph{differential weight} of a vector $c : G.\operatorname{edgeSet} \to \mathbb{F}_2$ is
\[
\operatorname{differentialWeight}(G, \Lambda, \operatorname{parityCheck}, c)
\;=\;
\bigl|\{v \in V \mid \operatorname{parityCheck}(\operatorname{localView}_\Lambda(v, c)) \neq 0\}\bigr|,
\]
counting the number of vertices at which the local parity check is violated.
\end{definition}

\begin{lemma}[Violated Check from Small Support Weight]
\label{lem:ExpanderViolatedChecks.violated_check_of_small_support}
\lean{ExpanderViolatedChecks.violated_check_of_small_support}
\leanok
\uses{def:SipserSpielman.edgeHammingWeight}
Let $\operatorname{parityCheck} : \mathbb{F}_2^s \to \mathbb{F}_2^m$ be a linear map, $d_L$ the minimum distance of $\operatorname{ClassicalCode.ofParityCheck}(\operatorname{parityCheck})$, with $d_L > 0$. If $x \in \mathbb{F}_2^s$ satisfies $x \neq 0$ and $\operatorname{wt}(x) \leq d_L - 1$, then $\operatorname{parityCheck}(x) \neq 0$.
\end{lemma}

\begin{proof}
\leanok
\uses{def:SipserSpielman.edgeHammingWeight}
Suppose for contradiction that $\operatorname{parityCheck}(x) = 0$. Then $x \in \ker(\operatorname{parityCheck})$, which by \texttt{ClassicalCode.code\_eq\_ker} means $x$ belongs to the classical code. Since $x \neq 0$, the definition of minimum distance gives $d_L \leq \operatorname{wt}(x)$. But $\operatorname{wt}(x) \leq d_L - 1 < d_L$, a contradiction. Hence $\operatorname{parityCheck}(x) \neq 0$.
\end{proof}

\begin{definition}[Support Edges at Vertex]
\label{def:ExpanderViolatedChecks.supportEdgesAtVertex}
\lean{ExpanderViolatedChecks.supportEdgesAtVertex}
\leanok
\uses{def:SipserSpielman.supportEdges}
For a graph $G$, a vector $c : G.\operatorname{edgeSet} \to \mathbb{F}_2$, and a vertex $v \in V$, the \emph{support edges at $v$} is
\[
\operatorname{supportEdgesAtVertex}(G, c, v) \;=\; \bigl|\{e \in \operatorname{supportEdges}(G, c) \mid v \in e\}\bigr|.
\]
\end{definition}

\begin{lemma}[Support Edges Have Both Endpoints in $S$]
\label{lem:ExpanderViolatedChecks.supportEdges_both_endpoints_in_S}
\lean{ExpanderViolatedChecks.supportEdges_both_endpoints_in_S}
\leanok
\uses{def:SipserSpielman.supportEdges, def:SipserSpielman.incidentToSupport}
If $e \in \operatorname{supportEdges}(G, c)$ and $w \in e$, then $w \in \operatorname{incidentToSupport}(G, c)$.
\end{lemma}

\begin{proof}
\leanok
\uses{def:SipserSpielman.supportEdges, def:SipserSpielman.incidentToSupport}
Let $e \in \operatorname{supportEdges}(G, c)$ and $w \in e$. Unfolding definitions, $e$ is the image of some edge $e'$ in the support of $c$ (i.e.\ $c(e') \neq 0$) with $e = e'.val$. Since $w \in e = e'.val$, we have $w \in \operatorname{incidentToSupport}(G, c)$ by definition, as witnessed by the edge $e'$ with $c(e') \neq 0$ and $w \in e'$.
\end{proof}

\begin{lemma}[Support Edges Are Disjoint from the Edge Boundary]
\label{lem:ExpanderViolatedChecks.supportEdge_not_in_edgeBoundary}
\lean{ExpanderViolatedChecks.supportEdge_not_in_edgeBoundary}
\leanok
\uses{def:SipserSpielman.supportEdges, def:SipserSpielman.incidentToSupport, lem:ExpanderViolatedChecks.supportEdges_both_endpoints_in_S}
If $e \in \operatorname{supportEdges}(G, c)$, then $e \notin \partial_E(\operatorname{incidentToSupport}(G, c))$.
\end{lemma}

\begin{proof}
\leanok
\uses{def:SipserSpielman.supportEdges, def:SipserSpielman.incidentToSupport, lem:ExpanderViolatedChecks.supportEdges_both_endpoints_in_S}
Suppose $e \in \partial_E(S)$ where $S = \operatorname{incidentToSupport}(G, c)$. By membership in the edge boundary, there exist $a, b$ with $e = \{a, b\}$, $a \in S$, and $b \notin S$. But since $e \in \operatorname{supportEdges}(G, c)$ and $b \in e$, by Lemma~\ref{lem:ExpanderViolatedChecks.supportEdges_both_endpoints_in_S} we have $b \in S$, a contradiction.
\end{proof}

\begin{lemma}[Within-$S$ Edges Plus Boundary Edges Bounded by Degree]
\label{lem:ExpanderViolatedChecks.withinS_plus_boundary_eq_degree}
\lean{ExpanderViolatedChecks.withinS_plus_boundary_eq_degree}
\leanok
\uses{def:RelativeVertexToEdgeExpansion.edgeBoundaryAtVertex}
Let $G$ be an $s$-regular graph, $S \subseteq V$, and $v \in S$. Then
\[
|(\partial_E S)_v| + |\{e \in G.\operatorname{edgeFinset} \mid v \in e,\; \forall w \in e,\, w \in S\}| \;\leq\; s,
\]
where $(\partial_E S)_v$ denotes the edges in the boundary of $S$ incident to $v$.
\end{lemma}

\begin{proof}
\leanok
\uses{def:RelativeVertexToEdgeExpansion.edgeBoundaryAtVertex}
Both the set of boundary edges at $v$ and the set of within-$S$ edges at $v$ are subsets of the incidence finset $\operatorname{incidenceFinset}(v)$, which has cardinality $\deg(v) = s$. Moreover, these two sets are disjoint: if $e$ is in the edge boundary, then $e$ has one endpoint outside $S$, so it cannot have all endpoints in $S$. By disjointness and the union bound:
\[
|\partial_E(S)_v| + |\text{within-}S\text{ edges at }v| = |\partial_E(S)_v \cup \text{within-}S| \leq |\operatorname{incidenceFinset}(v)| = s. \qedhere
\]
\end{proof}

\begin{lemma}[Local View Weight Bounded by Within-$S$ Edges]
\label{lem:ExpanderViolatedChecks.localView_weight_le_withinS}
\lean{ExpanderViolatedChecks.localView_weight_le_withinS}
\leanok
\uses{def:Labeling, def:SipserSpielman.incidentToSupport, def:SipserSpielman.supportEdges}
Let $G$ be a graph with labeling $\Lambda$ of degree $s$, $c : G.\operatorname{edgeSet} \to \mathbb{F}_2$, and $S \subseteq V$ such that every edge $e$ with $c(e) \neq 0$ has both endpoints in $S$. Then for every $v \in V$,
\[
\operatorname{wt}(\operatorname{localView}_\Lambda(v, c)) \;\leq\; |\{e \in G.\operatorname{edgeFinset} \mid v \in e,\; \forall w \in e,\, w \in S\}|.
\]
\end{lemma}

\begin{proof}
\leanok
\uses{def:Labeling, def:SipserSpielman.incidentToSupport, def:SipserSpielman.supportEdges}
Let $v \in V$. The Hamming weight $\operatorname{wt}(\operatorname{localView}_\Lambda(v, c))$ equals $|\{i \in \operatorname{Fin}(s) \mid \operatorname{localView}_\Lambda(v, c)(i) \neq 0\}|$. Define a map $\operatorname{mapEdge} : \operatorname{Fin}(s) \to \operatorname{Sym2}(V)$ by $i \mapsto \{v, (\Lambda(v))^{-1}(i)\}$.

We show this map sends the support indices into the within-$S$ edges: if $\operatorname{localView}_\Lambda(v, c)(i) \neq 0$, then $c(\{v, (\Lambda(v))^{-1}(i)\}) \neq 0$, so both endpoints lie in $S$ by hypothesis.

We then show $\operatorname{mapEdge}$ is injective on the support indices: if $\{v, w_1\} = \{v, w_2\}$ as elements of $\operatorname{Sym2}(V)$ with $v \neq w_1$, then $w_1 = w_2$, hence the preimages under $\Lambda(v)$ coincide.

By injectivity and inclusion, $|\text{support indices}| = |\operatorname{image}| \leq |\text{within-}S\text{ edges at }v|$.
\end{proof}

\begin{lemma}[Local View Weight Plus Boundary Bounded by Degree]
\label{lem:ExpanderViolatedChecks.localView_weight_add_boundary_le}
\lean{ExpanderViolatedChecks.localView_weight_add_boundary_le}
\leanok
\uses{def:Labeling, def:RelativeVertexToEdgeExpansion.edgeBoundaryAtVertex, lem:ExpanderViolatedChecks.localView_weight_le_withinS, lem:ExpanderViolatedChecks.withinS_plus_boundary_eq_degree}
Under the same hypotheses, for $v \in S$,
\[
\operatorname{wt}(\operatorname{localView}_\Lambda(v, c)) + |(\partial_E S)_v| \;\leq\; s.
\]
\end{lemma}

\begin{proof}
\leanok
\uses{def:Labeling, def:RelativeVertexToEdgeExpansion.edgeBoundaryAtVertex, lem:ExpanderViolatedChecks.localView_weight_le_withinS, lem:ExpanderViolatedChecks.withinS_plus_boundary_eq_degree}
Combine Lemma~\ref{lem:ExpanderViolatedChecks.localView_weight_le_withinS} (bounding the local view weight by within-$S$ edge count) and Lemma~\ref{lem:ExpanderViolatedChecks.withinS_plus_boundary_eq_degree} (bounding the sum of within-$S$ and boundary counts by $s$). The result follows by integer arithmetic.
\end{proof}

\begin{lemma}[Support Edges Cardinality Equals Edge Hamming Weight]
\label{lem:ExpanderViolatedChecks.supportEdges_card_eq}
\lean{ExpanderViolatedChecks.supportEdges_card_eq}
\leanok
\uses{def:SipserSpielman.supportEdges, def:SipserSpielman.edgeHammingWeight}
For any $c : G.\operatorname{edgeSet} \to \mathbb{F}_2$,
\[
|\operatorname{supportEdges}(G, c)| \;=\; \operatorname{edgeHammingWeight}(G, c).
\]
\end{lemma}

\begin{proof}
\leanok
\uses{def:SipserSpielman.supportEdges, def:SipserSpielman.edgeHammingWeight}
Unfolding definitions, $\operatorname{supportEdges}(G, c)$ is the image of the support of $c$ under the injective map $\operatorname{Subtype.val}$. The cardinality of the image under an injective map equals the cardinality of the domain, which is $\operatorname{edgeHammingWeight}(G, c)$.
\end{proof}

\begin{lemma}[Incident-to-Support Equals Incident Vertices of Support Edges]
\label{lem:ExpanderViolatedChecks.incidentToSupport_eq_incidentVertices}
\lean{ExpanderViolatedChecks.incidentToSupport_eq_incidentVertices}
\leanok
\uses{def:SipserSpielman.incidentToSupport, def:SipserSpielman.supportEdges}
For any $c : G.\operatorname{edgeSet} \to \mathbb{F}_2$,
\[
\operatorname{incidentToSupport}(G, c) \;=\; \operatorname{incidentVertices}(\operatorname{supportEdges}(G, c)).
\]
\end{lemma}

\begin{proof}
\leanok
\uses{def:SipserSpielman.incidentToSupport, def:SipserSpielman.supportEdges}
By extensionality: a vertex $v$ belongs to $\operatorname{incidentToSupport}(G, c)$ if and only if there exists an edge $e$ in the support of $c$ with $v \in e$, which holds if and only if there exists $e \in \operatorname{supportEdges}(G, c)$ with $v \in e$, i.e.\ $v \in \operatorname{incidentVertices}(\operatorname{supportEdges}(G, c))$. Both directions follow by unfolding definitions and matching witnesses.
\end{proof}

\begin{lemma}[Handshaking Lemma for Regular Graphs]
\label{lem:ExpanderViolatedChecks.regular_handshaking}
\lean{ExpanderViolatedChecks.regular_handshaking}
\leanok
\uses{def:SipserSpielman.edgeHammingWeight}
If $G$ is $s$-regular, then $2|G.\operatorname{edgeFinset}| = s \cdot |V|$.
\end{lemma}

\begin{proof}
\leanok
\uses{def:SipserSpielman.edgeHammingWeight}
Apply the standard handshaking identity $\sum_{v \in V} \deg(v) = 2|E|$. Since $G$ is $s$-regular, $\deg(v) = s$ for all $v$, so $\sum_{v \in V} s = s \cdot |V| = 2|E|$.
\end{proof}

\begin{lemma}[Incident-to-Support Cardinality at Most Twice Edge Weight]
\label{lem:ExpanderViolatedChecks.incidentToSupport_card_le_twice}
\lean{ExpanderViolatedChecks.incidentToSupport_card_le_twice}
\leanok
\uses{def:SipserSpielman.incidentToSupport, def:SipserSpielman.edgeHammingWeight, lem:ExpanderViolatedChecks.incidentToSupport_eq_incidentVertices, lem:ExpanderViolatedChecks.supportEdges_card_eq}
For any $c : G.\operatorname{edgeSet} \to \mathbb{F}_2$,
\[
|\operatorname{incidentToSupport}(G, c)| \;\leq\; 2 \cdot \operatorname{edgeHammingWeight}(G, c).
\]
\end{lemma}

\begin{proof}
\leanok
\uses{def:SipserSpielman.incidentToSupport, def:SipserSpielman.edgeHammingWeight, lem:ExpanderViolatedChecks.incidentToSupport_eq_incidentVertices, lem:ExpanderViolatedChecks.supportEdges_card_eq}
Let $E = \operatorname{supportEdges}(G, c)$. By Lemma~\ref{lem:ExpanderViolatedChecks.incidentToSupport_eq_incidentVertices}, $|\operatorname{incidentToSupport}| = |\operatorname{incidentVertices}(E)|$. Since each edge $e \in E$ contributes at most 2 vertices (as $|\{v \mid v \in e\}| \leq 2$ for any $e \in \operatorname{Sym2}(V)$), we have:
\[
|\operatorname{incidentVertices}(E)| \leq \sum_{e \in E} |\{v \mid v \in e\}| \leq \sum_{e \in E} 2 = 2|E| = 2\,\operatorname{edgeHammingWeight}(G, c),
\]
where the last equality uses Lemma~\ref{lem:ExpanderViolatedChecks.supportEdges_card_eq}.
\end{proof}

\begin{lemma}[$|S| \leq \alpha s |V|$]
\label{lem:ExpanderViolatedChecks.incidentToSupport_card_le_alphaS}
\lean{ExpanderViolatedChecks.incidentToSupport_card_le_alphaS}
\leanok
\uses{def:SipserSpielman.incidentToSupport, def:SipserSpielman.edgeHammingWeight, lem:ExpanderViolatedChecks.incidentToSupport_card_le_twice, lem:ExpanderViolatedChecks.regular_handshaking}
Let $G$ be $s$-regular, $\alpha > 0$, and suppose $\operatorname{edgeHammingWeight}(G, x) \leq \alpha |G.\operatorname{edgeFinset}|$. Then
\[
|\operatorname{incidentToSupport}(G, x)| \;\leq\; \alpha s |V|.
\]
\end{lemma}

\begin{proof}
\leanok
\uses{def:SipserSpielman.incidentToSupport, def:SipserSpielman.edgeHammingWeight, lem:ExpanderViolatedChecks.incidentToSupport_card_le_twice, lem:ExpanderViolatedChecks.regular_handshaking}
By Lemma~\ref{lem:ExpanderViolatedChecks.incidentToSupport_card_le_twice}, $|S| \leq 2\,\operatorname{edgeHammingWeight}(G, x) \leq 2\alpha|G.\operatorname{edgeFinset}|$. By the handshaking lemma (Lemma~\ref{lem:ExpanderViolatedChecks.regular_handshaking}), $2|G.\operatorname{edgeFinset}| = s|V|$. Therefore:
\[
|S| \leq 2\alpha|G.\operatorname{edgeFinset}| = \alpha \cdot 2|G.\operatorname{edgeFinset}| = \alpha \cdot s|V| = \alpha s |V|. \qedhere
\]
\end{proof}

\begin{lemma}[Local View Nonzero for Vertices Incident to Support]
\label{lem:ExpanderViolatedChecks.localView_ne_zero_of_incidentToSupport}
\lean{ExpanderViolatedChecks.localView_ne_zero_of_incidentToSupport}
\leanok
\uses{def:Labeling, def:SipserSpielman.incidentToSupport}
If $v \in \operatorname{incidentToSupport}(G, c)$, then $\operatorname{localView}_\Lambda(v, c) \neq 0$.
\end{lemma}

\begin{proof}
\leanok
\uses{def:Labeling, def:SipserSpielman.incidentToSupport}
Since $v \in \operatorname{incidentToSupport}(G, c)$, there exists an edge $e$ in the support of $c$ (i.e.\ $c(e) \neq 0$) with $v \in e$. Write $e = \{v, w\}$ for some neighbor $w$ of $v$. Let $i = \Lambda(v)(w, \cdot)$ be the label assigned to $w$. Then by definition of $\operatorname{localView}$, $\operatorname{localView}_\Lambda(v, c)(i) = c(e) \neq 0$. Hence $\operatorname{localView}_\Lambda(v, c) \neq 0$.
\end{proof}

\begin{lemma}[Differential Weight Bounds High-Expansion Vertex Count]
\label{lem:ExpanderViolatedChecks.differentialWeight_ge_highExpansion}
\lean{ExpanderViolatedChecks.differentialWeight_ge_highExpansion}
\leanok
\uses{def:Labeling, def:RelativeVertexToEdgeExpansion.highExpansionVertices, def:SipserSpielman.incidentToSupport, lem:RelativeVertexToEdgeExpansion.highExpansionVertices_subset, lem:ExpanderViolatedChecks.violated_check_of_small_support, lem:ExpanderViolatedChecks.localView_weight_add_boundary_le, lem:ExpanderViolatedChecks.localView_ne_zero_of_incidentToSupport}
Let $G$ be $s$-regular, $\Lambda$ a labeling, $\operatorname{parityCheck}$ a linear map with local code distance $d_L > 0$, and $c : G.\operatorname{edgeSet} \to \mathbb{F}_2$. Setting $S = \operatorname{incidentToSupport}(G, c)$ and $A = \operatorname{highExpansionVertices}(G, S, s, d_L - 1)$, we have
\[
|A| \;\leq\; \operatorname{differentialWeight}(G, \Lambda, \operatorname{parityCheck}, c).
\]
\end{lemma}

\begin{proof}
\leanok
\uses{def:Labeling, def:RelativeVertexToEdgeExpansion.highExpansionVertices, def:SipserSpielman.incidentToSupport, lem:RelativeVertexToEdgeExpansion.highExpansionVertices_subset, lem:ExpanderViolatedChecks.violated_check_of_small_support, lem:ExpanderViolatedChecks.localView_weight_add_boundary_le, lem:ExpanderViolatedChecks.localView_ne_zero_of_incidentToSupport}
It suffices to show $A \subseteq \{v \in V \mid \operatorname{parityCheck}(\operatorname{localView}_\Lambda(v, c)) \neq 0\}$. Let $v \in A$.

By Lemma~\ref{lem:RelativeVertexToEdgeExpansion.highExpansionVertices_subset}, $v \in S$. By the definition of high-expansion vertices, $|(\partial_E S)_v| \geq s - (d_L - 1)$.

Since every edge $e$ with $c(e) \neq 0$ has both endpoints in $S$ (by definition of $S = \operatorname{incidentToSupport}$), Lemma~\ref{lem:ExpanderViolatedChecks.localView_weight_add_boundary_le} gives $\operatorname{wt}(\operatorname{localView}_\Lambda(v, c)) + |(\partial_E S)_v| \leq s$.

Therefore $\operatorname{wt}(\operatorname{localView}_\Lambda(v, c)) \leq s - |(\partial_E S)_v| \leq d_L - 1$.

Since $v \in S = \operatorname{incidentToSupport}(G, c)$, Lemma~\ref{lem:ExpanderViolatedChecks.localView_ne_zero_of_incidentToSupport} gives $\operatorname{localView}_\Lambda(v, c) \neq 0$.

Applying Lemma~\ref{lem:ExpanderViolatedChecks.violated_check_of_small_support} with $x = \operatorname{localView}_\Lambda(v, c)$, $\operatorname{wt}(x) \leq d_L - 1$, and $x \neq 0$, we conclude $\operatorname{parityCheck}(x) \neq 0$.
\end{proof}

\begin{lemma}[$\beta'$ Is Positive]
\label{lem:ExpanderViolatedChecks.beta'_pos}
\lean{ExpanderViolatedChecks.beta'_pos}
\leanok
\uses{def:ExpanderViolatedChecks.beta', def:EdgeToVertexExpansion.betaParam}
If $s \geq 1$, $\lambda_2 < s$, and $\alpha > 0$, then $\beta'(s, \lambda_2, \alpha) > 0$.
\end{lemma}

\begin{proof}
\leanok
\uses{def:ExpanderViolatedChecks.beta', def:EdgeToVertexExpansion.betaParam}
Unfolding $\beta' = \operatorname{betaParam}$, we have
\[
\beta'(s, \lambda_2, \alpha) = \frac{\sqrt{\lambda_2^2 + 4s(s-\lambda_2)\alpha} - \lambda_2}{s(s-\lambda_2)\alpha}.
\]
The denominator satisfies $s(s - \lambda_2)\alpha > 0$ since $s \geq 1$, $s - \lambda_2 > 0$, and $\alpha > 0$.

For the numerator $\sqrt{\lambda_2^2 + 4s(s-\lambda_2)\alpha} - \lambda_2 > 0$, we consider two cases:
\begin{itemize}
\item If $\lambda_2 \geq 0$: since $4s(s-\lambda_2)\alpha > 0$, we have $\lambda_2^2 + 4s(s-\lambda_2)\alpha > \lambda_2^2$, so $\sqrt{\lambda_2^2 + 4s(s-\lambda_2)\alpha} > \sqrt{\lambda_2^2} = \lambda_2$ (using $\lambda_2 \geq 0$). Thus the numerator is positive.
\item If $\lambda_2 < 0$: then $\sqrt{\lambda_2^2 + 4s(s-\lambda_2)\alpha} \geq 0 > \lambda_2$, so the numerator is positive.
\end{itemize}
In both cases $\beta' > 0$.
\end{proof}

\begin{theorem}[Expander Violated Checks]
\label{thm:ExpanderViolatedChecks.expanderViolatedChecks}
\lean{ExpanderViolatedChecks.expanderViolatedChecks}
\leanok
\uses{def:EdgeToVertexExpansion.betaParam, def:RelativeVertexToEdgeExpansion.edgeBoundaryAtVertex, def:RelativeVertexToEdgeExpansion.highExpansionVertices, def:SipserSpielman.edgeHammingWeight, def:SipserSpielman.incidentToSupport, def:SipserSpielman.supportEdges, lem:RelativeVertexToEdgeExpansion.highExpansionVertices_subset, lem:SipserSpielman.supportEdges_subset, thm:EdgeToVertexExpansion.edgeToVertexExpansion, thm:RelativeVertexToEdgeExpansion.relativeVertexToEdgeExpansion}
Let $G$ be a connected $s$-regular graph on $|V| \geq 2$ vertices with second-largest eigenvalue $\lambda_2 < s$. Let $\operatorname{parityCheck} : \mathbb{F}_2^s \to \mathbb{F}_2^m$ be a linear map, $\Lambda$ a labeling of $G$, and $d_L$ the minimum distance of the local code, with $0 < d_L \leq s$. For $\alpha \in (0, 1]$ and $x : G.\operatorname{edgeSet} \to \mathbb{F}_2$ satisfying
\[
\operatorname{edgeHammingWeight}(G, x) \;\leq\; \alpha \cdot |G.\operatorname{edgeFinset}|,
\]
we have
\[
\operatorname{differentialWeight}(G, \Lambda, \operatorname{parityCheck}, x)
\;\geq\;
\beta'(s, \lambda_2, \alpha) \cdot \beta''(s, \lambda_2, \alpha, d_L) \cdot \operatorname{edgeHammingWeight}(G, x),
\]
where
\[
\beta'(s, \lambda_2, \alpha) = \frac{\sqrt{\lambda_2^2 + 4s(s-\lambda_2)\alpha} - \lambda_2}{s(s-\lambda_2)\alpha}, \quad
\beta''(s, \lambda_2, \alpha, d_L) = \frac{(d_L - 1 - \lambda_2) - \alpha s(s - \lambda_2)}{d_L - 1}.
\]
\end{theorem}

\begin{proof}
\leanok
\uses{def:EdgeToVertexExpansion.betaParam, def:RelativeVertexToEdgeExpansion.edgeBoundaryAtVertex, def:RelativeVertexToEdgeExpansion.highExpansionVertices, def:SipserSpielman.edgeHammingWeight, def:SipserSpielman.incidentToSupport, def:SipserSpielman.supportEdges, lem:RelativeVertexToEdgeExpansion.highExpansionVertices_subset, lem:SipserSpielman.supportEdges_subset, thm:EdgeToVertexExpansion.edgeToVertexExpansion, thm:RelativeVertexToEdgeExpansion.relativeVertexToEdgeExpansion, lem:ExpanderViolatedChecks.differentialWeight_ge_highExpansion, lem:ExpanderViolatedChecks.incidentToSupport_card_le_alphaS, lem:ExpanderViolatedChecks.incidentToSupport_eq_incidentVertices, lem:ExpanderViolatedChecks.supportEdges_card_eq, def:ExpanderViolatedChecks.beta', def:ExpanderViolatedChecks.beta'', lem:ExpanderViolatedChecks.beta'_pos}

\textbf{Trivial case.} If $\beta'(s, \lambda_2, \alpha) \cdot \beta''(s, \lambda_2, \alpha, d_L) \leq 0$, then the right-hand side is at most $0 \leq \operatorname{differentialWeight}(\ldots)$, and the bound holds.

\textbf{Main case: $\beta' \cdot \beta'' > 0$.}

By Lemma~\ref{lem:ExpanderViolatedChecks.beta'_pos}, $\beta' > 0$ holds whenever $s \geq 1$, $\lambda_2 < s$, $\alpha > 0$. Since $\beta' \cdot \beta'' > 0$ and $\beta' > 0$, we also have $\beta'' > 0$.

From $\beta'' > 0$, we deduce $d_L \geq 2$ as follows: if $d_L = 1$, then $\beta'' = ({-\lambda_2 - \alpha s(s-\lambda_2)})/0 = 0$ in Lean's division, contradicting $\beta'' > 0$.

From $\beta'' > 0$, we also deduce $\alpha s < 1$: since $d_L \geq 2$, the denominator $d_L - 1 > 0$, so the numerator $(d_L - 1 - \lambda_2) - \alpha s(s - \lambda_2) > 0$. If $\alpha s \geq 1$, then $\alpha s(s - \lambda_2) \geq s - \lambda_2 \geq d_L - \lambda_2 > d_L - 1 - \lambda_2$ (using $d_L \leq s$), contradicting positivity of the numerator. Hence $\alpha s < 1$.

\textbf{Edge case: $x = 0$.} If $\operatorname{edgeHammingWeight}(G, x) = 0$, the right-hand side is $0$ and the bound holds trivially.

\textbf{Main inequality chain.} Assume $\operatorname{edgeHammingWeight}(G, x) > 0$. Set $S = \operatorname{incidentToSupport}(G, x)$, $E = \operatorname{supportEdges}(G, x)$, $A = \operatorname{highExpansionVertices}(G, S, s, d_L - 1)$.

\textbf{Step 1: $|A| \leq \operatorname{differentialWeight}(\ldots)$.} By Lemma~\ref{lem:ExpanderViolatedChecks.differentialWeight_ge_highExpansion}.

\textbf{Step 2: $\beta' \cdot \operatorname{edgeHammingWeight}(G, x) \leq |S|$.} Since $E \subseteq G.\operatorname{edgeFinset}$ (by Lemma~\ref{lem:SipserSpielman.supportEdges_subset}) and $|E| = \operatorname{edgeHammingWeight}(G, x)$ (by Lemma~\ref{lem:ExpanderViolatedChecks.supportEdges_card_eq}) with $|E| \leq \alpha|G.\operatorname{edgeFinset}|$, we may apply the edge-to-vertex expansion theorem (Theorem~\ref{thm:EdgeToVertexExpansion.edgeToVertexExpansion}):
\[
|\operatorname{incidentVertices}(E)| \;\geq\; \operatorname{betaParam}(s, \lambda_2, \alpha) \cdot |E|.
\]
Since $S = \operatorname{incidentVertices}(E)$ (Lemma~\ref{lem:ExpanderViolatedChecks.incidentToSupport_eq_incidentVertices}) and $\beta' = \operatorname{betaParam}$, this yields $|S| \geq \beta' \cdot \operatorname{edgeHammingWeight}(G, x)$.

\textbf{Step 3: $\beta'' \cdot |S| \leq |A|$.} Since $|S| \leq \alpha s |V|$ (Lemma~\ref{lem:ExpanderViolatedChecks.incidentToSupport_card_le_alphaS}) and $\alpha s < 1$, $d_L - 1 \in (0, s]$, we apply the relative vertex-to-edge expansion theorem (Theorem~\ref{thm:RelativeVertexToEdgeExpansion.relativeVertexToEdgeExpansion}) with parameters $\alpha' = \alpha s$ and $b = d_L - 1$:
\[
\frac{(d_L - 1 - \lambda_2) - \alpha s(s - \lambda_2)}{d_L - 1} \cdot |S| \;\leq\; |A|,
\]
which is exactly $\beta'' \cdot |S| \leq |A|$.

\textbf{Combining.} The chain
\[
\beta' \cdot \beta'' \cdot \operatorname{edgeHammingWeight}(G, x)
\;\leq\; \beta'' \cdot |S|
\;\leq\; |A|
\;\leq\; \operatorname{differentialWeight}(G, \Lambda, \operatorname{parityCheck}, x)
\]
follows from Steps 2, 3, and 1 respectively (using $\beta'' > 0$ in Step 2, and applying linear arithmetic to combine the bounds).
\end{proof}

\begin{remark}[Paper Corrections]
\textbf{\color{red}Errata.} The following errors were identified in the original paper and corrected in this formalization:
\begin{itemize}
  \item Paper states |S| ≤ 2α|X₁| = (4α/s)|X₀|, but the correct derivation using the handshaking lemma 2|X₁| = s|X₀| gives 2α|X₁| = αs|X₀|, so α' should be αs, not 4α/s.
  \item Paper uses b = d_L in Lemma 3 (defining A = {v ∈ S | |(δS)_v| ≥ s - d_L}), but the violated-check argument requires the local view weight to be strictly less than d_L (i.e., ≤ d_L - 1). This needs b = d_L - 1 (giving A = {v | |(δS)_v| ≥ s - d_L + 1}), so that within-S edges ≤ d_L - 1 < d_L guarantees the nonzero local view is not a codeword.
  \item Consequently, β'' should be ((d_L - 1 - λ₂) - αs(s - λ₂))/(d_L - 1), not ((d_L - λ₂) - (4α/s)(s - λ₂))/d_L as stated in the paper. The formalization's formula is correct.
\end{itemize}
\end{remark}

%--- Thm_9: ExpanderBitDegree ---
\chapter{Thm 9: Expander Bit Degree}

\begin{definition}[Total Hamming Weight]
\label{def:ExpanderBitDegree.totalWeight}
\lean{ExpanderBitDegree.totalWeight}
\leanok
\uses{def:hammingWeight}
Let $V$ be a finite type and $m \in \mathbb{N}$. For a vector $y : V \to (\operatorname{Fin} m \to \mathbb{F}_2)$, the \emph{total Hamming weight} is
\[
  \operatorname{totalWeight}(y) \;=\; \sum_{v \in V} \operatorname{wt}(y_v) \;\in\; \mathbb{N},
\]
where $\operatorname{wt}(y_v)$ denotes the Hamming weight of the fiber $y_v \in \mathbb{F}_2^m$.
\end{definition}

\begin{definition}[Vertex Support]
\label{def:ExpanderBitDegree.vertexSupport}
\lean{ExpanderBitDegree.vertexSupport}
\leanok
\uses{def:hammingWeight}
For $y : V \to (\operatorname{Fin} m \to \mathbb{F}_2)$, the \emph{vertex support} is
\[
  S(y) \;=\; \{v \in V \mid y_v \neq 0\} \;\subseteq\; V.
\]
\end{definition}

\begin{lemma}[Vertex Support is Nonempty for Nonzero Vectors]
\label{lem:ExpanderBitDegree.vertexSupport_nonempty}
\lean{ExpanderBitDegree.vertexSupport_nonempty}
\leanok
\uses{def:ExpanderBitDegree.vertexSupport}
If $y : V \to (\operatorname{Fin} m \to \mathbb{F}_2)$ satisfies $y \neq 0$, then $S(y) \neq \emptyset$.
\end{lemma}

\begin{proof}
\leanok
\uses{def:ExpanderBitDegree.vertexSupport}
We proceed by contradiction. Suppose $S(y) = \emptyset$. Then for every $v \in V$ we have $y_v = 0$: indeed, if $y_v \neq 0$ then $v \in S(y)$, contradicting $S(y) = \emptyset$. Hence $y(v)(i) = 0$ for all $v$ and $i$, giving $y = 0$, a contradiction.
\end{proof}

\begin{lemma}[Card of Vertex Support Bounds Total Weight from Below]
\label{lem:ExpanderBitDegree.card_vertexSupport_le_totalWeight}
\lean{ExpanderBitDegree.card_vertexSupport_le_totalWeight}
\leanok
\uses{def:ExpanderBitDegree.totalWeight, def:ExpanderBitDegree.vertexSupport, def:hammingWeight}
For any $y : V \to (\operatorname{Fin} m \to \mathbb{F}_2)$,
\[
  |S(y)| \;\leq\; \operatorname{totalWeight}(y).
\]
\end{lemma}

\begin{proof}
\leanok
\uses{def:ExpanderBitDegree.totalWeight, def:ExpanderBitDegree.vertexSupport, def:hammingWeight}
We compute:
\[
  |S(y)| = \sum_{v \in S(y)} 1 \;\leq\; \sum_{v \in S(y)} \operatorname{wt}(y_v) \;\leq\; \sum_{v \in V} \operatorname{wt}(y_v) = \operatorname{totalWeight}(y).
\]
The first inequality holds because for $v \in S(y)$ we have $y_v \neq 0$, so $\operatorname{wt}(y_v) \geq 1$. The second holds by extending the sum to all of $V$ (all terms are non-negative).
\end{proof}

\begin{lemma}[Total Weight Bounded by Degree Times Card]
\label{lem:ExpanderBitDegree.totalWeight_le_mul_card}
\lean{ExpanderBitDegree.totalWeight_le_mul_card}
\leanok
\uses{def:ExpanderBitDegree.totalWeight, def:ExpanderBitDegree.vertexSupport, def:hammingWeight}
For any $y : V \to (\operatorname{Fin} m \to \mathbb{F}_2)$,
\[
  \operatorname{totalWeight}(y) \;\leq\; m \cdot |S(y)|.
\]
\end{lemma}

\begin{proof}
\leanok
\uses{def:ExpanderBitDegree.totalWeight, def:ExpanderBitDegree.vertexSupport, def:hammingWeight}
We split the sum $\sum_{v \in V} \operatorname{wt}(y_v) = \sum_{v \in S(y)} \operatorname{wt}(y_v) + \sum_{v \notin S(y)} \operatorname{wt}(y_v)$. For $v \notin S(y)$ we have $y_v = 0$ so $\operatorname{wt}(y_v) = 0$, hence the second sum vanishes. Then
\[
  \sum_{v \in S(y)} \operatorname{wt}(y_v) \;\leq\; \sum_{v \in S(y)} m \;=\; m \cdot |S(y)|,
\]
since the Hamming weight of any $\mathbb{F}_2^m$-vector is at most $m$.
\end{proof}

\begin{lemma}[Total Weight Divided by $m$ Bounds Vertex Support Card]
\label{lem:ExpanderBitDegree.totalWeight_div_le_card}
\lean{ExpanderBitDegree.totalWeight_div_le_card}
\leanok
\uses{def:ExpanderBitDegree.totalWeight, def:ExpanderBitDegree.vertexSupport, lem:ExpanderBitDegree.totalWeight_le_mul_card}
If $m > 0$, then for any $y : V \to (\operatorname{Fin} m \to \mathbb{F}_2)$,
\[
  \frac{\operatorname{totalWeight}(y)}{m} \;\leq\; |S(y)|.
\]
\end{lemma}

\begin{proof}
\leanok
\uses{def:ExpanderBitDegree.totalWeight, def:ExpanderBitDegree.vertexSupport, lem:ExpanderBitDegree.totalWeight_le_mul_card}
Since $m > 0$, we may divide both sides of $\operatorname{totalWeight}(y) \leq m \cdot |S(y)|$ by $m$ to obtain the result. More precisely, rewriting the inequality as $\operatorname{totalWeight}(y) \leq |S(y)| \cdot m$ and dividing by $m$ (which is positive) yields $\operatorname{totalWeight}(y)/m \leq |S(y)|$ over $\mathbb{R}$.
\end{proof}

\begin{definition}[Coboundary Map]
\label{def:ExpanderBitDegree.coboundary}
\lean{ExpanderBitDegree.coboundary}
\leanok
\uses{def:Labeling, def:differential}
Let $G$ be a simple graph on $V$, $\Lambda$ a labeling of $G$ with parameter $s$, and $H : \mathbb{F}_2^s \to \mathbb{F}_2^m$ a parity-check linear map. The \emph{coboundary map} $\delta$ applied to $y : V \to \mathbb{F}_2^m$ is the function on edge-set of $G$ given by
\[
  \delta(y)(e) \;=\; \sum_{v \in V} \langle y_v,\; \partial_\Lambda^H(\mathbf{1}_e)(v)\rangle \;\in\; \mathbb{F}_2,
\]
where $\mathbf{1}_e$ denotes the indicator of edge $e$ and $\partial_\Lambda^H$ is the Tanner differential.
\end{definition}

\begin{definition}[Coboundary Weight]
\label{def:ExpanderBitDegree.coboundaryWeight}
\lean{ExpanderBitDegree.coboundaryWeight}
\leanok
\uses{def:hammingWeight, def:ExpanderBitDegree.coboundary}
The \emph{coboundary weight} of $y$ is
\[
  \operatorname{wt}(\delta(y)) \;=\; |\{e \in E(G) \mid \delta(y)(e) \neq 0\}|.
\]
\end{definition}

\begin{definition}[Transpose of Parity Check]
\label{def:ExpanderBitDegree.parityCheckTranspose}
\lean{ExpanderBitDegree.parityCheckTranspose}
\leanok
\uses{def:ClassicalCode}
Let $H : \mathbb{F}_2^s \to \mathbb{F}_2^m$ be a linear map. The \emph{transpose} $H^\top : \mathbb{F}_2^m \to \mathbb{F}_2^s$ is defined by
\[
  (H^\top y)(j) \;=\; \langle y,\, H(e_j)\rangle \;\in\; \mathbb{F}_2,
\]
for $j \in \{0,\ldots,s-1\}$, where $e_j = \mathbf{1}_j$ denotes the $j$-th standard basis vector. This is a linear map: additivity follows since the dot product is $\mathbb{F}_2$-bilinear, and scalar compatibility follows since $\mathbb{F}_2$ is a field of characteristic 2.
\end{definition}

\begin{lemma}[Transpose is Injective when Parity Check is Surjective]
\label{lem:ExpanderBitDegree.parityCheckTranspose_injective}
\lean{ExpanderBitDegree.parityCheckTranspose_injective}
\leanok
\uses{def:ExpanderBitDegree.parityCheckTranspose}
If $H : \mathbb{F}_2^s \to \mathbb{F}_2^m$ is surjective, then $H^\top : \mathbb{F}_2^m \to \mathbb{F}_2^s$ is injective.
\end{lemma}

\begin{proof}
\leanok
\uses{def:ExpanderBitDegree.parityCheckTranspose}
We show $\ker(H^\top) = \{0\}$. Suppose $H^\top y = 0$, i.e., $\langle y, H(e_j)\rangle = 0$ for all $j$. We must show $y = 0$. For any index $i$, since $H$ is surjective, there exists $x \in \mathbb{F}_2^s$ with $H(x) = e_i$ (the $i$-th standard basis vector). Writing $x = \sum_j x_j e_j$ and using linearity of $H$,
\[
  \langle y, e_i\rangle = \langle y, H(x)\rangle = \left\langle y, \sum_j x_j H(e_j)\right\rangle = \sum_j x_j \langle y, H(e_j)\rangle = 0,
\]
since each $\langle y, H(e_j)\rangle = (H^\top y)(j) = 0$. But $\langle y, e_i\rangle = y(i)$, so $y(i) = 0$ for all $i$, giving $y = 0$. The converse ($y = 0 \Rightarrow H^\top y = 0$) is immediate.
\end{proof}

\begin{lemma}[Transpose Image Lies in Dual Code]
\label{lem:ExpanderBitDegree.parityCheckTranspose_mem_dualCode}
\lean{ExpanderBitDegree.parityCheckTranspose_mem_dualCode}
\leanok
\uses{def:ExpanderBitDegree.parityCheckTranspose, def:ClassicalCode}
For any $y \in \mathbb{F}_2^m$, we have $H^\top y \in C^\perp$, where $C^\perp$ is the dual of the classical code $C = \ker(H)$ associated to parity-check map $H$.
\end{lemma}

\begin{proof}
\leanok
\uses{def:ExpanderBitDegree.parityCheckTranspose, def:ClassicalCode}
We need to show that for every $c \in C = \ker(H)$, $\langle H^\top y, c\rangle = 0$. We have
\[
  \langle H^\top y, c\rangle = \sum_j (H^\top y)(j) \cdot c_j = \sum_j \langle y, H(e_j)\rangle c_j.
\]
Rewriting the double sum and exchanging order of summation,
\[
  = \sum_i y_i \left(\sum_j H(e_j)_i \cdot c_j\right) = \sum_i y_i \cdot H(c)_i = \langle y, H(c)\rangle.
\]
Since $c \in \ker(H)$, we have $H(c) = 0$, so $\langle y, H(c)\rangle = 0$.
\end{proof}

\begin{lemma}[Transpose Weight Bounds Dual Distance]
\label{lem:ExpanderBitDegree.parityCheckTranspose_weight_ge}
\lean{ExpanderBitDegree.parityCheckTranspose_weight_ge}
\leanok
\uses{def:ClassicalCode, def:hammingWeight, lem:ExpanderBitDegree.parityCheckTranspose_injective, lem:ExpanderBitDegree.parityCheckTranspose_mem_dualCode, def:ExpanderBitDegree.parityCheckTranspose}
Let $H : \mathbb{F}_2^s \to \mathbb{F}_2^m$ be surjective, and let $d_{L^\perp}$ be the minimum distance of the dual code $C^\perp$. For any nonzero $y \in \mathbb{F}_2^m$,
\[
  d_{L^\perp} \;\leq\; \operatorname{wt}(H^\top y).
\]
\end{lemma}

\begin{proof}
\leanok
\uses{def:ClassicalCode, def:hammingWeight, lem:ExpanderBitDegree.parityCheckTranspose_injective, lem:ExpanderBitDegree.parityCheckTranspose_mem_dualCode, def:ExpanderBitDegree.parityCheckTranspose}
Since $H$ is surjective, by Lemma~\ref{lem:ExpanderBitDegree.parityCheckTranspose_injective}, $H^\top$ is injective. Because $y \neq 0$, we have $H^\top y \neq 0$ (otherwise $y = H^\top(0) = 0$ by injectivity). By Lemma~\ref{lem:ExpanderBitDegree.parityCheckTranspose_mem_dualCode}, $H^\top y \in C^\perp$. Therefore $H^\top y$ is a nonzero codeword of $C^\perp$, and so its Hamming weight is at least $d_{L^\perp}$ (by definition of minimum distance as the infimum of weights of nonzero codewords).
\end{proof}

\begin{lemma}[Coboundary as Sum Formula]
\label{lem:ExpanderBitDegree.coboundary_eq_sum}
\lean{ExpanderBitDegree.coboundary_eq_sum}
\leanok
\uses{def:ExpanderBitDegree.coboundary, def:differential}
For any edge $e \in E(G)$,
\[
  \delta(y)(e) \;=\; \sum_{v \in V} \langle y_v,\; \partial_\Lambda^H(\mathbf{1}_e)(v)\rangle.
\]
\end{lemma}

\begin{proof}
\leanok
\uses{def:ExpanderBitDegree.coboundary, def:differential}
This holds by definition (reflexivity).
\end{proof}

\begin{lemma}[Coboundary at Boundary Vertex]
\label{lem:ExpanderBitDegree.coboundary_boundary_vertex}
\lean{ExpanderBitDegree.coboundary_boundary_vertex}
\leanok
\uses{def:ExpanderBitDegree.coboundary, def:differential, def:localView, lem:differential_apply}
Let $e = \{v, w\} \in E(G)$ be an edge such that for every $u \in e$ with $u \neq v$, $y_u = 0$. Then
\[
  \delta(y)(e) \;=\; \langle y_v,\; \partial_\Lambda^H(\mathbf{1}_e)(v)\rangle.
\]
\end{lemma}

\begin{proof}
\leanok
\uses{def:ExpanderBitDegree.coboundary, def:differential, def:localView, lem:differential_apply}
We write $\delta(y)(e) = \langle y_v, \partial_\Lambda^H(\mathbf{1}_e)(v)\rangle + \sum_{u \neq v} \langle y_u, \partial_\Lambda^H(\mathbf{1}_e)(u)\rangle$ by extracting the $v$-term from the full sum. It suffices to show the remaining sum is zero. For $u \neq v$, we consider two cases. If $u \in e$ (i.e., $u$ is the other endpoint), then by hypothesis $y_u = 0$, so $\langle y_u, \partial_\Lambda^H(\mathbf{1}_e)(u)\rangle = 0$. If $u \notin e$, then the local view $\mathcal{L}_\Lambda(u)(\mathbf{1}_e) = 0$, because the edge $\{u, (\Lambda_u)^{-1}(i)\}$ (for any label $i$) is not equal to $e$ (since $u \notin e$); hence $\partial_\Lambda^H(\mathbf{1}_e)(u) = H(0) = 0$ by applying Lemma~\ref{lem:differential_apply}, so the dot product vanishes.
\end{proof}

\begin{lemma}[Non-Boundary Positions Plus Boundary Edges at Most $s$]
\label{lem:ExpanderBitDegree.nonBdyPos_add_boundary_le_degree}
\lean{ExpanderBitDegree.nonBdyPos_add_boundary_le_degree}
\leanok
\uses{def:Labeling, def:RelativeVertexToEdgeExpansion.edgeBoundaryAtVertex, lem:RelativeVertexToEdgeExpansion.edgeBoundaryAtVertex_le_degree}
Let $G$ be $s$-regular with labeling $\Lambda$, and let $S \subseteq V$, $v \in S$. Define the set of label positions corresponding to $S$-neighbors:
\[
  P_S(v) \;=\; \{j \in \operatorname{Fin} s \mid (\Lambda_v)^{-1}(j) \in S\}.
\]
Then
\[
  |P_S(v)| + |(\partial_e S)_v| \;\leq\; s,
\]
where $(\partial_e S)_v$ is the set of edges in $\partial_e S$ incident to $v$.
\end{lemma}

\begin{proof}
\leanok
\uses{def:Labeling, def:RelativeVertexToEdgeExpansion.edgeBoundaryAtVertex, lem:RelativeVertexToEdgeExpansion.edgeBoundaryAtVertex_le_degree}
Define the complementary set $P_{\bar{S}}(v) = \{j \in \operatorname{Fin} s \mid (\Lambda_v)^{-1}(j) \notin S\}$. Since $P_S(v)$ and $P_{\bar{S}}(v)$ partition $\operatorname{Fin} s$, we have $|P_S(v)| + |P_{\bar{S}}(v)| = s$. Thus it suffices to show $|(\partial_e S)_v| \leq |P_{\bar{S}}(v)|$, i.e., there is an injection from boundary edges at $v$ into $P_{\bar{S}}(v)$.

If $s = 0$, then there are no edges and the boundary at $v$ is empty, so the inequality is trivial by Lemma~\ref{lem:RelativeVertexToEdgeExpansion.edgeBoundaryAtVertex_le_degree}.

For $s > 0$, we define a map $f : \partial_e S \to \operatorname{Fin} s$ by sending a boundary edge $e$ (with $v \in e$) to $(\Lambda_v)(\text{other endpoint of } e)$. We show:

\emph{$f$ maps into $P_{\bar{S}}(v)$}: For a boundary edge $e \in \partial_e S$, there exist $a \in S$ and $b \notin S$ with $e = \{a, b\}$. Since $v \in S$, we have $v = a$ and the other endpoint is $b \notin S$, so $f(e)$ corresponds to a neighbor outside $S$.

\emph{$f$ is injective}: If $f(e_1) = f(e_2)$, then both edges have the same other endpoint $w$ at $v$ (by injectivity of $\Lambda_v$), and both edges equal $\{v, w\}$, so $e_1 = e_2$.

Thus $|(\partial_e S)_v| \leq |P_{\bar{S}}(v)|$, giving $|P_S(v)| + |(\partial_e S)_v| \leq |P_S(v)| + |P_{\bar{S}}(v)| = s$.
\end{proof}

\begin{theorem}[Expander Bit Degree]
\label{thm:ExpanderBitDegree.expanderBitDegree}
\lean{ExpanderBitDegree.expanderBitDegree}
\leanok
\uses{def:ClassicalCode, def:GraphExpansion.secondLargestEigenvalue, def:Labeling, def:RelativeVertexToEdgeExpansion.edgeBoundaryAtVertex, def:RelativeVertexToEdgeExpansion.highExpansionVertices, def:differential, def:hammingWeight, def:localView, lem:RelativeVertexToEdgeExpansion.edgeBoundaryAtVertex_le_degree, lem:RelativeVertexToEdgeExpansion.highExpansionVertices_subset, lem:differential_apply, thm:RelativeVertexToEdgeExpansion.relativeVertexToEdgeExpansion}

Let $G$ be a connected $s$-regular simple graph on a finite vertex set $V$ with $|V| \geq 2$ and $s \geq 1$. Let $\lambda_2 = \lambda_2(G)$ be its second largest eigenvalue. Let $H : \mathbb{F}_2^s \to \mathbb{F}_2^m$ be a surjective linear parity-check map, and let $\Lambda$ be a labeling of $G$. Write $k_L + m = s$ (so $k_L$ is the dimension of the local code $\ker(H)$).

Let $d_{L^\perp} \geq 2$ be the minimum distance of the dual local code $C^\perp$, with $d_{L^\perp} \leq s$. Let $\alpha > 0$ with $\alpha m < 1$. Define $b = d_{L^\perp} - 1 > 0$.

For any $y : V \to \mathbb{F}_2^m$ with
\[
  \operatorname{totalWeight}(y) \;\leq\; \alpha \cdot |V| \cdot m,
\]
setting
\[
  \beta \;=\; \frac{(d_{L^\perp} - 1 - \lambda_2) - \alpha m (s - \lambda_2)}{m(d_{L^\perp} - 1)},
\]
we have
\[
  \beta \cdot \operatorname{totalWeight}(y) \;\leq\; \operatorname{wt}(\delta(y)).
\]
\end{theorem}

\begin{proof}
\leanok
\uses{def:ClassicalCode, def:GraphExpansion.secondLargestEigenvalue, def:Labeling, def:RelativeVertexToEdgeExpansion.edgeBoundaryAtVertex, def:RelativeVertexToEdgeExpansion.highExpansionVertices, def:differential, def:hammingWeight, def:localView, lem:RelativeVertexToEdgeExpansion.edgeBoundaryAtVertex_le_degree, lem:RelativeVertexToEdgeExpansion.highExpansionVertices_subset, lem:differential_apply, thm:RelativeVertexToEdgeExpansion.relativeVertexToEdgeExpansion, lem:ExpanderBitDegree.card_vertexSupport_le_totalWeight, lem:ExpanderBitDegree.totalWeight_div_le_card, lem:ExpanderBitDegree.parityCheckTranspose_weight_ge, lem:ExpanderBitDegree.coboundary_boundary_vertex, lem:ExpanderBitDegree.nonBdyPos_add_boundary_le_degree}

\textbf{Base case.} If $y = 0$, then $\operatorname{totalWeight}(y) = 0$ and $\operatorname{wt}(\delta(y)) = 0$, so $\beta \cdot 0 \leq 0$ holds trivially.

\textbf{Setup for $y \neq 0$.} Let $S = S(y)$ be the vertex support of $y$. Since $y \neq 0$, the support $S$ is nonempty, so $|S| > 0$. Since $y \neq 0$ and $y$ takes values in $\mathbb{F}_2^m$, we must have $m > 0$ (otherwise $y$ would be $0$ by vacuity).

Set $\alpha' = \alpha m$ and $b = d_{L^\perp} - 1$. Then $\alpha' > 0$ (since $\alpha > 0$ and $m > 0$), $\alpha' < 1$ (by hypothesis), and $b > 0$ (since $d_{L^\perp} \geq 2$) and $b \leq s$ (since $d_{L^\perp} \leq s$).

By Lemma~\ref{lem:ExpanderBitDegree.card_vertexSupport_le_totalWeight}, $|S| \leq \operatorname{totalWeight}(y)$, and by the hypothesis $\operatorname{totalWeight}(y) \leq \alpha \cdot |V| \cdot m = \alpha' \cdot |V|$, we get $|S| \leq \alpha' \cdot |V|$.

\textbf{Applying the relative vertex-to-edge expansion theorem.} Let $A = A_b(S)$ be the set of high-expansion vertices (those $v \in S$ with $|(\partial_e S)_v| \geq s - b$). By Theorem~\ref{thm:RelativeVertexToEdgeExpansion.relativeVertexToEdgeExpansion} applied with parameters $\alpha'$, $b$, $S$,
\[
  \frac{b - \lambda_2 - \alpha'(s - \lambda_2)}{b} \cdot |S| \;\leq\; |A|.
\]
By Lemma~\ref{lem:ExpanderBitDegree.totalWeight_div_le_card}, $\operatorname{totalWeight}(y)/m \leq |S|$.

\textbf{Reducing to $|A| \leq \operatorname{wt}(\delta(y))$.} We claim it suffices to show $|A| \leq \operatorname{wt}(\delta(y))$. To derive the stated inequality from this:
\begin{itemize}
  \item If $\frac{b - \lambda_2 - \alpha'(s - \lambda_2)}{b} \geq 0$, then
  \[
    \beta \cdot \operatorname{totalWeight}(y) = \frac{b - \lambda_2 - \alpha'(s - \lambda_2)}{mb} \cdot \operatorname{totalWeight}(y)
    = \frac{b - \lambda_2 - \alpha'(s - \lambda_2)}{b} \cdot \frac{\operatorname{totalWeight}(y)}{m}
    \leq \frac{b - \lambda_2 - \alpha'(s - \lambda_2)}{b} \cdot |S|
    \leq |A| \leq \operatorname{wt}(\delta(y)).
  \]
  \item If $\frac{b - \lambda_2 - \alpha'(s - \lambda_2)}{b} < 0$, then $\beta < 0$, so $\beta \cdot \operatorname{totalWeight}(y) \leq 0 \leq \operatorname{wt}(\delta(y))$.
\end{itemize}

\textbf{Proving $|A| \leq \operatorname{wt}(\delta(y))$.} We construct an injection $g : A \to \{e \in E(G) \mid \delta(y)(e) \neq 0\}$.

Note that $G$ is connected with $|V| \geq 2$ and $s$-regular with $s \geq 1$, so $G$ has at least one edge.

For each $v \in A$, we show there exists an edge $e$ incident to $v$ in $\partial_e S$ with $\delta(y)(e) \neq 0$:

By Lemma~\ref{lem:RelativeVertexToEdgeExpansion.highExpansionVertices_subset}, $A \subseteq S$, so $y_v \neq 0$ for all $v \in A$.

Let $v \in A$. By definition of $A$, vertex $v$ has $|(\partial_e S)_v| \geq s - b$. Consider the transpose vector $t_v = H^\top(y_v) \in \mathbb{F}_2^s$. By Lemma~\ref{lem:ExpanderBitDegree.parityCheckTranspose_weight_ge}, $\operatorname{wt}(t_v) \geq d_{L^\perp}$.

Let $\operatorname{supp}(t_v) = \{j \in \operatorname{Fin} s \mid t_v(j) \neq 0\}$ with $|\operatorname{supp}(t_v)| \geq d_{L^\perp} = b + 1$. Define the set of label positions pointing into $S$:
\[
  P_S(v) = \{j \in \operatorname{Fin} s \mid (\Lambda_v)^{-1}(j) \in S\}.
\]
By Lemma~\ref{lem:ExpanderBitDegree.nonBdyPos_add_boundary_le_degree}, $|P_S(v)| \leq s - |(\partial_e S)_v| \leq s - (s - b) = b$. Since $|\operatorname{supp}(t_v)| \geq b + 1 > b \geq |P_S(v)|$, there exists $j \in \operatorname{supp}(t_v) \setminus P_S(v)$.

For this $j$, let $w = (\Lambda_v)^{-1}(j)$ (the neighbor of $v$ at label $j$) and let $e = \{v, w\} \in E(G)$. Since $j \notin P_S(v)$, we have $w \notin S$. Hence $e$ is a boundary edge ($v \in S$, $w \notin S$).

We show $\delta(y)(e) \neq 0$: Since $w \notin S = S(y)$, we have $y_w = 0$. By Lemma~\ref{lem:ExpanderBitDegree.coboundary_boundary_vertex}, $\delta(y)(e) = \langle y_v, \partial_\Lambda^H(\mathbf{1}_e)(v)\rangle$. Applying Lemma~\ref{lem:differential_apply} and computing the local view, the local view of $v$ at edge $e$ is the standard basis vector $e_j$ (since $\Lambda_v$ assigns label $j$ to $w$). Thus
\[
  \delta(y)(e) = \langle y_v, H(e_j)\rangle = (H^\top y_v)(j) = t_v(j) \neq 0
\]
since $j \in \operatorname{supp}(t_v)$.

\textbf{Injectivity of the assignment.} Choose for each $v \in A$ a witness edge $g(v)$ incident to $v$ in $\partial_e S$ with $\delta(y)(g(v)) \neq 0$. We show the map $v \mapsto g(v)$ is injective on $A$. Suppose $g(v_1) = g(v_2) = e$ for $v_1, v_2 \in A$. The edge $e \in \partial_e S$ has one endpoint $a \in S$ and one endpoint $b \notin S$. Both $v_1, v_2 \in A \subseteq S$ are incident to $e$, so both equal $a$, giving $v_1 = v_2$.

\textbf{Conclusion.} By Finset cardinality bounds, $|A| \leq |\{e \in E(G) \mid \delta(y)(e) \neq 0\}| = \operatorname{wt}(\delta(y))$.
\end{proof}

%--- Thm_10: GilbertVarshamovPlus ---
\chapter{Thm 10: Gilbert--Varshamov Plus}

\begin{lemma}[GV Threshold Positivity]
\label{lem:gv_threshold_pos}
\lean{gv_threshold_pos}
\leanok
\uses{def:binaryEntropy, thm:binaryEntropy_lt_half}
Let $\delta \in \mathbb{R}$ with $0 < \delta < 11/100$. Then
\[
0 < \frac{2}{1/2 - h(\delta)},
\]
where $h(\delta)$ denotes the binary entropy function.
\end{lemma}

\begin{proof}
\leanok
\uses{def:binaryEntropy, thm:binaryEntropy_lt_half}
We first obtain $h(\delta) < 1/2$ by applying $\texttt{binaryEntropy\_lt\_half}$ to the hypotheses $0 < \delta$ and $\delta < 11/100$. The result then follows by positivity of the numerator $2 > 0$ (by numerical computation) and positivity of the denominator $1/2 - h(\delta) > 0$ (by linear arithmetic from $h(\delta) < 1/2$).
\end{proof}

\begin{lemma}[GV Lower Bound on $n$]
\label{lem:gv_n_large}
\lean{gv_n_large}
\leanok
\uses{def:binaryEntropy, thm:binaryEntropy_lt_half}
Let $\delta \in \mathbb{R}$ with $0 < \delta < 11/100$, and let $n \in \mathbb{N}$. If
\[
n > \frac{2}{1/2 - h(\delta)},
\]
then $(1/2 - h(\delta)) \cdot n > 2$.
\end{lemma}

\begin{proof}
\leanok
\uses{def:binaryEntropy, thm:binaryEntropy_lt_half}
From $\texttt{binaryEntropy\_lt\_half}$ applied to $0 < \delta$ and $\delta < 11/100$, we have $h(\delta) < 1/2$, so the gap $0 < 1/2 - h(\delta)$ holds by linear arithmetic. Rewriting the hypothesis $n > 2/(1/2 - h(\delta))$ via $\texttt{div\_lt\_iff}$ using this positivity, we obtain $2 < (1/2 - h(\delta)) \cdot n$, and the conclusion follows by linear arithmetic.
\end{proof}

\begin{lemma}[Existence of $k$ in GV Range]
\label{lem:gv_k_exists}
\lean{gv_k_exists}
\leanok
\uses{def:binaryEntropy, thm:binaryEntropy_lt_half}
Let $\delta \in \mathbb{R}$ with $0 < \delta < 11/100$, and let $n \in \mathbb{N}$ with $n > 2/(1/2 - h(\delta))$. Then there exists $k \in \mathbb{N}$ such that:
\[
\frac{n}{2} < k, \quad k + 1 < (1 - h(\delta))\cdot n, \quad k < n.
\]
\end{lemma}

\begin{proof}
\leanok
\uses{def:binaryEntropy, thm:binaryEntropy_lt_half}
We set $k = \lfloor n/2 \rfloor + 1$. From $\texttt{binaryEntropy\_lt\_half}$ we have $h(\delta) < 1/2$, and by Lemma~\ref{lem:gv_n_large} we have $(1/2 - h(\delta)) \cdot n > 2$, which gives $n/2 + 1 < (1 - h(\delta))\cdot n - 1$, hence $n/2 + 1 + 1 < (1 - h(\delta)) \cdot n$.

We verify each condition for $k = \lfloor n/2 \rfloor + 1$:

\textit{Lower bound}: Using $\lfloor n/2 \rfloor \leq n/2$ (from $\texttt{Nat.cast\_div\_le}$) and $n < 2(\lfloor n/2 \rfloor + 1)$ (by integer arithmetic), pushing casts and applying linear arithmetic yields $n/2 < k$.

\textit{Upper bound}: By the same cast inequalities and the gap $n/2 + 1 < (1-h(\delta)) \cdot n - 1$, linear arithmetic gives $k + 1 < (1 - h(\delta)) \cdot n$.

\textit{Strict bound $k < n$}: We argue by contradiction. If $n \leq k$, then integer arithmetic gives $n \leq 2$. One then shows $h(\delta) > 0$ (by expanding the definition: since $0 < \delta < 1/2$, both $\log_2 \delta < 0$ and $\log_2(1-\delta) < 0$, so $h(\delta) = -\delta \log_2 \delta - (1-\delta)\log_2(1-\delta) > 0$ via nonlinear arithmetic). With $n \leq 2$, we get $(1/2 - h(\delta)) \cdot n \leq 1$, contradicting $(1/2 - h(\delta)) \cdot n > 2$.
\end{proof}

\begin{definition}[Evaluation Linear Map]
\label{def:evalMap}
\lean{evalMap}
\leanok
\uses{def:F2}
For $n, r \in \mathbb{N}$ and a vector $v \in \mathbb{F}_2^n$, the \emph{evaluation map} at $v$ is the linear map
\[
\operatorname{eval}_v : \operatorname{Hom}_{\mathbb{F}_2}(\mathbb{F}_2^n, \mathbb{F}_2^r) \to \mathbb{F}_2^r, \quad \phi \mapsto \phi(v).
\]
This is given by $\texttt{LinearMap.apply$_l$}\; v$.
\end{definition}

\begin{lemma}[Evaluation Map Application]
\label{lem:evalMap_apply}
\lean{evalMap_apply}
\leanok
\uses{def:evalMap, def:F2}
For $v \in \mathbb{F}_2^n$ and $\phi \in \operatorname{Hom}_{\mathbb{F}_2}(\mathbb{F}_2^n, \mathbb{F}_2^r)$, we have $\operatorname{eval}_v(\phi) = \phi(v)$.
\end{lemma}

\begin{proof}
\leanok
\uses{def:evalMap, def:F2}
This follows immediately by unfolding the definition of $\operatorname{eval}_v$ as $\texttt{LinearMap.apply$_l$}\; v$ and simplifying.
\end{proof}

\begin{lemma}[Surjectivity of Evaluation Map]
\label{lem:evalMap_surjective}
\lean{evalMap_surjective}
\leanok
\uses{def:evalMap, def:F2}
For $v \in \mathbb{F}_2^n$ with $v \neq 0$, the evaluation map $\operatorname{eval}_v$ is surjective.
\end{lemma}

\begin{proof}
\leanok
\uses{def:evalMap, def:F2, lem:evalMap_apply}
Let $w \in \mathbb{F}_2^r$ be arbitrary. Since $v \neq 0$, there exists an index $j$ with $v_j \neq 0$. By case analysis on $\mathbb{F}_2$ (using $\texttt{decide}$), every element of $\mathbb{F}_2$ is $0$ or $1$, so $v_j = 1$.

We construct $\phi = (\texttt{LinearMap.proj}\; j) \cdot_r w$, i.e.\ the linear map sending $e_j \mapsto w$ and all other standard basis vectors to $0$. Then $\phi(v) = v_j \cdot w = 1 \cdot w = w$, so $\operatorname{eval}_v(\phi) = w$.
\end{proof}

\begin{lemma}[Kernel Rank of Evaluation Map]
\label{lem:evalMap_ker_finrank}
\lean{evalMap_ker_finrank}
\leanok
\uses{def:evalMap, def:F2, lem:evalMap_surjective}
For $v \in \mathbb{F}_2^n$ with $v \neq 0$,
\[
\operatorname{finrank}_{\mathbb{F}_2}(\ker \operatorname{eval}_v) = nr - r.
\]
\end{lemma}

\begin{proof}
\leanok
\uses{def:evalMap, def:F2, lem:evalMap_surjective}
By Lemma~\ref{lem:evalMap_surjective}, $\operatorname{eval}_v$ is surjective. Applying the rank--nullity theorem ($\texttt{LinearMap.finrank\_range\_add\_finrank\_ker}$), we have
\[
\operatorname{finrank}(\operatorname{range}(\operatorname{eval}_v)) + \operatorname{finrank}(\ker \operatorname{eval}_v) = \operatorname{finrank}(\operatorname{Hom}(\mathbb{F}_2^n, \mathbb{F}_2^r)).
\]
Since $\operatorname{eval}_v$ is surjective, $\operatorname{range}(\operatorname{eval}_v) = \mathbb{F}_2^r$ so its rank is $r$. By the formula for the finrank of a Hom space ($\texttt{Module.finrank\_linearMap}$) the total dimension is $nr$. Hence $\operatorname{finrank}(\ker \operatorname{eval}_v) = nr - r$, which follows by integer arithmetic.
\end{proof}

\begin{lemma}[Cardinality of Kernel of Evaluation Map]
\label{lem:card_ker_evalMap}
\lean{card_ker_evalMap}
\leanok
\uses{def:evalMap, def:F2, lem:evalMap_ker_finrank, lem:evalMap_surjective}
For $v \in \mathbb{F}_2^n$ with $v \neq 0$,
\[
|\ker \operatorname{eval}_v| \cdot 2^r = |\operatorname{Hom}_{\mathbb{F}_2}(\mathbb{F}_2^n, \mathbb{F}_2^r)|.
\]
\end{lemma}

\begin{proof}
\leanok
\uses{def:evalMap, def:F2, lem:evalMap_ker_finrank, lem:evalMap_surjective}
We apply $\texttt{Module.card\_eq\_pow\_finrank}$ to both sides. For the kernel, its cardinality is $2^{nr - r}$ by Lemma~\ref{lem:evalMap_ker_finrank}. For the Hom space, its cardinality is $2^{nr}$. Thus the left side is $2^{nr - r} \cdot 2^r = 2^{nr}$. We verify that $r \leq nr$ (since $n > 0$, which follows from $v \neq 0$ implying $n$ is a singleton or larger), and conclude by integer arithmetic that $nr - r + r = nr$.
\end{proof}

\begin{theorem}[Axiom: Existence of Good Linear Map]
\label{thm:exists_good_map}
\lean{exists_good_map}
\leanok
\uses{def:hammingBall, def:hammingWeight, def:F2}

\textbf{\color{red}[UNPROVEN AXIOM]}

Let $n, r \in \mathbb{N}$ with $0 < r \leq n$, and let $t \in \mathbb{R}$. Suppose
\[
\bigl|\operatorname{Ball}(n, t-1)\bigr| < 2^r.
\]
Then there exists a linear map $\phi : \mathbb{F}_2^n \to \mathbb{F}_2^r$ such that for every nonzero $v \in \mathbb{F}_2^n$ with $\operatorname{wt}(v) < t$, we have $\phi(v) \neq 0$.

\textit{Note: This is stated as an axiom because the full probabilistic counting argument was not completed in the formalization. The argument follows the probabilistic method: a uniformly random linear map avoids each bad vector $v$ with probability $1 - 2^{-r}$, and a union bound over the $|\operatorname{Ball}(n,t-1)|$ bad vectors yields existence of a good map when the ball is smaller than $2^r$.}
\end{theorem}

\begin{theorem}[Axiom: Existence of Good Surjective Map with Both Distance Bounds]
\label{thm:exists_good_surjective_map}
\lean{exists_good_surjective_map}
\leanok
\uses{def:hammingBall, def:hammingWeight, def:ClassicalCode, def:ClassicalCode.dualCode, def:ClassicalCode.ofParityCheck, def:F2}

\textbf{\color{red}[UNPROVEN AXIOM]}

Let $n, r \in \mathbb{N}$ with $0 < r \leq n$, and let $t \in \mathbb{R}$. Suppose
\[
\bigl|\operatorname{Ball}(n, t-1)\bigr| < 2^r \quad \text{and} \quad \bigl|\operatorname{Ball}(n, t-1)\bigr| < 2^{n-r}.
\]
Then there exists a surjective linear map $\phi : \mathbb{F}_2^n \to \mathbb{F}_2^r$ such that:
\begin{enumerate}
\item For every nonzero $v \in \mathbb{F}_2^n$ with $\operatorname{wt}(v) < t$, we have $\phi(v) \neq 0$.
\item For every $w \in (\ker \phi)^\perp$ with $w \neq 0$, we have $\operatorname{wt}(w) \geq t$.
\end{enumerate}

\textit{Note: This extends \texttt{exists\_good\_map} by simultaneously bounding both the distance and the dual-distance conditions via a union bound: the probability of a bad distance event plus the probability of a bad dual-distance event is less than $1$, and any map satisfying the dual condition is automatically surjective (otherwise some nonzero $y$ has $\phi^\top(y) = 0$, which has weight $0 < t$).}
\end{theorem}

\begin{lemma}[Distance Lower Bound from Weight Condition]
\label{lem:distance_ge_of_no_lowweight}
\lean{distance_ge_of_no_lowweight}
\leanok
\uses{def:ClassicalCode, def:ClassicalCode.distance, def:hammingWeight}
Let $\mathcal{C}$ be a classical code of length $n$, let $t \in \mathbb{R}$. Suppose every nonzero codeword $v \in \mathcal{C}$ satisfies $\operatorname{wt}(v) \geq t$, and suppose $\mathcal{C}$ contains at least one nonzero element. Then $d(\mathcal{C}) \geq t$.
\end{lemma}

\begin{proof}
\leanok
\uses{def:ClassicalCode, def:ClassicalCode.distance, def:hammingWeight}
Unfolding the definition of distance, let $S = \{w \in \mathbb{N} \mid \exists\, x \in \mathcal{C},\; x \neq 0,\; \operatorname{wt}(x) = w\}$. From the hypothesis there exists a nonzero codeword, so $S$ is nonempty.

For each $w' \in S$, write $w' = \operatorname{wt}(x')$ for some nonzero $x' \in \mathcal{C}$. By the weight hypothesis, $\operatorname{wt}(x') \geq t$, hence $\lceil t \rceil \leq w'$.

Applying $\texttt{le\_csInf}$ to the nonemptiness and this lower bound, we get $\inf S \geq \lceil t \rceil$. Casting to reals:
\[
d(\mathcal{C}) = \inf S \geq \lceil t \rceil \geq t,
\]
where the last inequality is $\texttt{Nat.le\_ceil}$.
\end{proof}

\begin{lemma}[Existence of Submodule of Prescribed Dimension]
\label{lem:exists_submodule_of_finrank}
\lean{exists_submodule_of_finrank}
\leanok
\uses{def:F2}
Let $W$ be a submodule of $\mathbb{F}_2^n$, and let $k \leq \operatorname{finrank}_{\mathbb{F}_2}(W)$. Then there exists a submodule $W' \leq W$ with $\operatorname{finrank}_{\mathbb{F}_2}(W') = k$.
\end{lemma}

\begin{proof}
\leanok
\uses{def:F2}
Since $k \leq \operatorname{finrank}(W) = |\text{ChooseBasisIndex}(W)|$ (by $\texttt{Module.finrank\_eq\_card\_chooseBasisIndex}$), we apply $\texttt{Finset.exists\_subset\_card\_eq}$ to obtain a subset $s$ of the basis index set of cardinality $k$.

Let $B$ be the chosen basis of $W$ and define $W'_\text{sub} = \operatorname{span}_{\mathbb{F}_2}\{B(i) \mid i \in s\}$ as a submodule of $W$. Let $W' = W'_\text{sub}$ mapped along the inclusion $W \hookrightarrow \mathbb{F}_2^n$.

We verify $W' \leq W$ by tracing elements through the map. For the dimension: since $B$ is a basis, the restriction to $s$ gives a linearly independent family (by $\texttt{LinearIndependent.comp}$ with the injective coercion). Thus $\operatorname{finrank}(W'_\text{sub}) = |s| = k$ by $\texttt{finrank\_span\_eq\_card}$. Finally, $\texttt{Submodule.finrank\_map\_subtype\_eq}$ gives $\operatorname{finrank}(W') = \operatorname{finrank}(W'_\text{sub}) = k$.
\end{proof}

\begin{lemma}[Existence of Code with Good Distance]
\label{lem:exists_code_with_good_distance}
\lean{exists_code_with_good_distance}
\leanok
\uses{def:ClassicalCode, def:ClassicalCode.dimension, def:binaryEntropy, def:hammingBall, def:hammingWeight, lem:hammingBall_empty_of_neg, thm:exists_good_map, thm:hammingBallBound, thm:mem_hammingBall}
Let $\delta \in \mathbb{R}$ with $0 < \delta \leq 1/2$, and $n, k \in \mathbb{N}$ with $0 < n$, $0 < k$, $k < n$, and
\[
k + 1 < (1 - h(\delta)) \cdot n.
\]
Then there exists a classical code $\mathcal{C}$ of length $n$ with $\dim(\mathcal{C}) = k$ and $d(\mathcal{C}) \geq \delta n$.
\end{lemma}

\begin{proof}
\leanok
\uses{def:ClassicalCode, def:ClassicalCode.dimension, def:binaryEntropy, def:hammingBall, def:hammingWeight, lem:hammingBall_empty_of_neg, lem:exists_submodule_of_finrank, lem:distance_ge_of_no_lowweight, thm:exists_good_map, thm:hammingBallBound, thm:mem_hammingBall}
Set $r = n - k$. Then $0 < r$ and $r \leq n$. From the hypothesis $k + 1 < (1-h(\delta))n$, we get $h(\delta) \cdot n < n - k = r$.

\textit{Hamming ball bound}: We show $|\operatorname{Ball}(n, \delta n - 1)| < 2^r$.

\textbf{Case} $\delta n < 1$: Then $\delta n - 1 < 0$, so $\operatorname{Ball}(n, \delta n - 1) = \emptyset$ by Lemma~\ref{lem:hammingBall_empty_of_neg}, and the bound is trivial.

\textbf{Case} $\delta n \geq 1$: Set $d' = \lfloor \delta n \rfloor - 1$. We verify $d' + 1 \leq \delta n$ using $\lfloor \delta n \rfloor \geq 1$ and the floor inequality. We also show $\operatorname{Ball}(n, \delta n - 1) \subseteq \operatorname{Ball}(n, d')$ by Theorem~\ref{thm:mem_hammingBall}: for any $v$ with $\operatorname{wt}(v) < \delta n - 1 < \lfloor \delta n \rfloor$, we have $\operatorname{wt}(v) \leq \lfloor \delta n \rfloor - 1 = d'$ by integer arithmetic. Then:
\[
|\operatorname{Ball}(n, \delta n - 1)| \leq |\operatorname{Ball}(n, d')| \leq 2^{h(\delta) \cdot n} < 2^r,
\]
where the middle inequality uses Theorem~\ref{thm:hammingBallBound} with $d' + 1 \leq \delta n$, and the last uses $h(\delta) \cdot n < r$.

\textit{Obtaining the map}: By Theorem~\ref{thm:exists_good_map}, there exists $\phi : \mathbb{F}_2^n \to \mathbb{F}_2^r$ such that every nonzero $v$ with $\operatorname{wt}(v) < \delta n$ satisfies $\phi(v) \neq 0$.

\textit{Dimension of kernel}: By rank--nullity, $\operatorname{finrank}(\operatorname{range}(\phi)) + \operatorname{finrank}(\ker \phi) = n$. Since $\operatorname{finrank}(\operatorname{range}(\phi)) \leq r$, we get $\operatorname{finrank}(\ker \phi) \geq n - r = k$.

\textit{Subcode}: By Lemma~\ref{lem:exists_submodule_of_finrank}, there exists $W \leq \ker \phi$ with $\operatorname{finrank}(W) = k$. The associated code $\mathcal{C} = \langle W \rangle$ has $\dim(\mathcal{C}) = k$.

\textit{Good distance}: For any nonzero $v \in W$, we have $v \in \ker \phi$, so $\phi(v) = 0$. If $\operatorname{wt}(v) < \delta n$, then the property of $\phi$ would require $\phi(v) \neq 0$, a contradiction. Hence $\operatorname{wt}(v) \geq \delta n$.

\textit{Nonemptiness}: Since $\operatorname{finrank}(W) = k \geq 1$, $W$ is nontrivial (as $k > 0$), so it contains a nonzero element.

By Lemma~\ref{lem:distance_ge_of_no_lowweight}, $d(\mathcal{C}) \geq \delta n$.
\end{proof}

\begin{lemma}[Satisfiability of Good Distance Premises]
\label{lem:exists_code_with_good_distance_premises_satisfiable}
\lean{exists_code_with_good_distance_premises_satisfiable}
\leanok
\uses{def:binaryEntropy, thm:binaryEntropy_lt_half}
There exist $\delta \in \mathbb{R}$ and $n, k \in \mathbb{N}$ satisfying all premises of Lemma~\ref{lem:exists_code_with_good_distance}: $0 < \delta$, $\delta \leq 1/2$, $0 < n$, $0 < k$, $k < n$, and $k + 1 < (1 - h(\delta)) \cdot n$.
\end{lemma}

\begin{proof}
\leanok
\uses{def:binaryEntropy, thm:binaryEntropy_lt_half}
We exhibit the witness $\delta = 1/100$, $n = 1000$, $k = 498$. The numerical conditions $0 < 1/100$, $1/100 \leq 1/2$, $0 < 1000$, $0 < 498$, $498 < 1000$ are verified by numerical computation ($\texttt{norm\_num}$). For the bound $499 < (1 - h(1/100)) \cdot 1000$: we first obtain $h(1/100) < 1/2$ from $\texttt{binaryEntropy\_lt\_half}$, then apply nonlinear arithmetic ($\texttt{nlinarith}$) after pushing casts.
\end{proof}

\begin{lemma}[Double Dual Code Identity]
\label{lem:dualCode_dualCode_eq}
\lean{dualCode_dualCode_eq}
\leanok
\uses{def:ClassicalCode, def:ClassicalCode.dualCode, thm:ClassicalCode.mem_dualCode_iff}
For any classical code $\mathcal{C}$ of length $n$,
\[
((\mathcal{C}^\perp)^\perp) = \mathcal{C}.
\]
\end{lemma}

\begin{proof}
\leanok
\uses{def:ClassicalCode, def:ClassicalCode.dualCode, thm:ClassicalCode.mem_dualCode_iff}
We work over the finite-dimensional space $\mathbb{F}_2^n$. By extensionality, it suffices to show that for every $w \in \mathbb{F}_2^n$, $w \in ((\mathcal{C}^\perp)^\perp)$ if and only if $w \in \mathcal{C}$.

($\Leftarrow$): Suppose $w \in \mathcal{C}$. By Theorem~\ref{thm:ClassicalCode.mem_dualCode_iff}, we must show that for every $c \in \mathcal{C}^\perp$, $c \cdot w = 0$. But by definition of the dual code (Theorem~\ref{thm:ClassicalCode.mem_dualCode_iff} applied to $c$), we have $c \cdot w = 0$ for all $w \in \mathcal{C}$; since $w \in \mathcal{C}$ and dot product is commutative, $c \cdot w = w \cdot c = 0$.

($\Rightarrow$): Suppose $w \in ((\mathcal{C}^\perp)^\perp)$, i.e.\ for all $c \in \mathcal{C}^\perp$, $c \cdot w = 0$. We assume for contradiction $w \notin \mathcal{C}$. By $\texttt{Submodule.exists\_dual\_map\_eq\_bot\_of\_notMem}$ (applied to the finite-dimensional setting), there exists a linear functional $f : \mathbb{F}_2^n \to \mathbb{F}_2$ with $f(w) \neq 0$ and $f|_{\mathcal{C}} = 0$. Via the dot-product equivalence $e : \operatorname{Hom}(\mathbb{F}_2^n, \mathbb{F}_2) \cong \mathbb{F}_2^n$, let $c = e^{-1}(f)$.

Since $f|_{\mathcal{C}} = 0$, we have $f(v) = c \cdot v = 0$ for all $v \in \mathcal{C}$, so $c \in \mathcal{C}^\perp$ by Theorem~\ref{thm:ClassicalCode.mem_dualCode_iff}. But then the assumption $w \in ((\mathcal{C}^\perp)^\perp)$ gives $c \cdot w = 0$, which under $e$ reads $f(w) = 0$, contradicting $f(w) \neq 0$.
\end{proof}

\begin{lemma}[Distance of Double Dual]
\label{lem:distance_dualCode_dualCode}
\lean{distance_dualCode_dualCode}
\leanok
\uses{def:ClassicalCode, def:ClassicalCode.dualCode, def:ClassicalCode.distance, lem:dualCode_dualCode_eq}
For any classical code $\mathcal{C}$ of length $n$,
\[
d((\mathcal{C}^\perp)^\perp) = d(\mathcal{C}).
\]
\end{lemma}

\begin{proof}
\leanok
\uses{def:ClassicalCode, def:ClassicalCode.dualCode, def:ClassicalCode.distance, lem:dualCode_dualCode_eq}
Unfolding the definition of distance, both sides are $\inf\{w \in \mathbb{N} \mid \exists\, x, x \in \mathcal{C} \text{ (or } ((\mathcal{C}^\perp)^\perp)\text{)},\; x \neq 0,\; \operatorname{wt}(x) = w\}$. By Lemma~\ref{lem:dualCode_dualCode_eq}, $((\mathcal{C}^\perp)^\perp) = \mathcal{C}$, so the index sets are equal. The result follows by congruence.
\end{proof}

\begin{lemma}[Existence of Code with Good Dual Distance]
\label{lem:exists_code_with_good_dual_distance}
\lean{exists_code_with_good_dual_distance}
\leanok
\uses{def:ClassicalCode, def:ClassicalCode.dimension, def:ClassicalCode.dualCode, def:binaryEntropy, thm:ClassicalCode.dimension_dualCode, lem:exists_code_with_good_distance, lem:distance_dualCode_dualCode}
Let $\delta \in \mathbb{R}$ with $0 < \delta \leq 1/2$, and $n, k \in \mathbb{N}$ with $0 < n$, $0 < k$, $k < n$, and
\[
n - k + 1 < (1 - h(\delta)) \cdot n.
\]
Then there exists a classical code $\mathcal{C}$ of length $n$ with $\dim(\mathcal{C}) = k$ and $d(\mathcal{C}^\perp) \geq \delta n$.
\end{lemma}

\begin{proof}
\leanok
\uses{def:ClassicalCode, def:ClassicalCode.dimension, def:ClassicalCode.dualCode, def:binaryEntropy, thm:ClassicalCode.dimension_dualCode, lem:exists_code_with_good_distance, lem:distance_dualCode_dualCode}
Let $r = n - k$. Then $0 < r$ and $r < n$. We convert the hypothesis: $(n-k) + 1 = (r : \mathbb{R}) + 1 < (1-h(\delta))n$. By Lemma~\ref{lem:exists_code_with_good_distance} applied to $\delta$, $n$, $r$ (with $r$ in place of $k$), there exists a code $D$ of length $n$ with $\dim(D) = r$ and $d(D) \geq \delta n$.

Set $\mathcal{C} = D^\perp$. By Theorem~\ref{thm:ClassicalCode.dimension_dualCode},
\[
\dim(\mathcal{C}) = \dim(D^\perp) = n - \dim(D) = n - r = k.
\]
For the dual distance:
\[
d(\mathcal{C}^\perp) = d((D^\perp)^\perp) = d(D) \geq \delta n,
\]
where we apply Lemma~\ref{lem:distance_dualCode_dualCode}.
\end{proof}

\begin{lemma}[Satisfiability of Good Dual Distance Premises]
\label{lem:exists_code_with_good_dual_distance_premises_satisfiable}
\lean{exists_code_with_good_dual_distance_premises_satisfiable}
\leanok
\uses{def:binaryEntropy, thm:binaryEntropy_lt_half}
There exist $\delta \in \mathbb{R}$ and $n, k \in \mathbb{N}$ satisfying all premises of Lemma~\ref{lem:exists_code_with_good_dual_distance}: $0 < \delta$, $\delta \leq 1/2$, $0 < n$, $0 < k$, $k < n$, and $n - k + 1 < (1-h(\delta)) \cdot n$.
\end{lemma}

\begin{proof}
\leanok
\uses{def:binaryEntropy, thm:binaryEntropy_lt_half}
We exhibit the witness $\delta = 1/100$, $n = 1000$, $k = 502$. All numerical conditions are verified by $\texttt{norm\_num}$. The bound $499 < (1-h(1/100)) \cdot 1000$ follows from $h(1/100) < 1/2$ (via $\texttt{binaryEntropy\_lt\_half}$) and $\texttt{nlinarith}$.
\end{proof}

\begin{lemma}[Key Natural Number Inequality]
\label{lem:key_nat_ineq_hundredth}
\lean{key_nat_ineq_hundredth}
\leanok

$100^{100} < 2^{10} \cdot 99^{99}$.
\end{lemma}

\begin{proof}
\leanok

This is verified by computation ($\texttt{native\_decide}$).
\end{proof}

\begin{lemma}[Key Fractional Inequality]
\label{lem:key_frac_ineq_hundredth}
\lean{key_frac_ineq_hundredth}
\leanok
\uses{lem:key_nat_ineq_hundredth}
$100 \cdot \left(\dfrac{100}{99}\right)^{99} < 2^{10}$.
\end{lemma}

\begin{proof}
\leanok
\uses{lem:key_nat_ineq_hundredth}
We rewrite: $100 \cdot (100/99)^{99} = 100^{100}/99^{99}$, so the inequality is equivalent (after clearing the denominator $99^{99} > 0$) to $100^{100} < 2^{10} \cdot 99^{99}$. We cast to naturals: $100^{100} = \uparrow(100^{100} : \mathbb{N}) < \uparrow(2^{10} \cdot 99^{99})$, which holds by Lemma~\ref{lem:key_nat_ineq_hundredth}.
\end{proof}

\begin{lemma}[Binary Entropy at $1/100$ is Less Than $1/10$]
\label{lem:binaryEntropy_at_one_hundredth}
\lean{binaryEntropy_at_one_hundredth}
\leanok
\uses{def:binaryEntropy, lem:key_frac_ineq_hundredth}
$h(1/100) < 1/10$.
\end{lemma}

\begin{proof}
\leanok
\uses{def:binaryEntropy, lem:key_frac_ineq_hundredth}
We expand $h(1/100) = -(1/100)\log_2(1/100) - (99/100)\log_2(99/100)$. We compute:
\[
\log_2(1/100) = -\log_2(100), \qquad \log_2(99/100) = -\log_2(100/99),
\]
so $h(1/100) = (1/100)\log_2(100) + (99/100)\log_2(100/99)$.

To prove this is less than $1/10$, it suffices to show:
\[
\log_2(100) + 99 \cdot \log_2(100/99) < 10.
\]
We bound the left side: by the logarithm identity,
\[
\log_2\!\bigl(100 \cdot (100/99)^{99}\bigr) = \log_2(100) + 99\log_2(100/99).
\]
Since $100 \cdot (100/99)^{99} < 2^{10}$ by Lemma~\ref{lem:key_frac_ineq_hundredth}, and $\log_2$ is strictly increasing ($\texttt{Real.logb\_lt\_logb}$), we get $\log_2(100 \cdot (100/99)^{99}) < \log_2(2^{10}) = 10$. The conclusion follows by linear arithmetic, using positivity of $\log_2(100) > 0$ and $\log_2(100/99) > 0$.
\end{proof}

\begin{lemma}[Binary Entropy at $1/100$ Tight Bound]
\label{lem:binaryEntropy_at_one_hundredth_tight}
\lean{binaryEntropy_at_one_hundredth_tight}
\leanok
\uses{def:binaryEntropy, lem:binaryEntropy_at_one_hundredth}
$h(1/100) < 499/1000$.
\end{lemma}

\begin{proof}
\leanok
\uses{def:binaryEntropy, lem:binaryEntropy_at_one_hundredth}
This follows immediately from Lemma~\ref{lem:binaryEntropy_at_one_hundredth}: $h(1/100) < 1/10 = 100/1000 < 499/1000$.
\end{proof}

\begin{lemma}[Hamming Ball Bound Below Power of Two]
\label{lem:hammingBall_lt_pow}
\lean{hammingBall_lt_pow}
\leanok
\uses{def:binaryEntropy, def:hammingBall, lem:hammingBall_empty_of_neg, thm:hammingBallBound, thm:mem_hammingBall}
Let $\delta \in \mathbb{R}$ with $0 < \delta \leq 1/2$, $n \in \mathbb{N}$ with $0 < n$, and $m \in \mathbb{N}$ with $h(\delta) \cdot n < m$. Then
\[
\bigl|\operatorname{Ball}(n, \delta n - 1)\bigr| < 2^m.
\]
\end{lemma}

\begin{proof}
\leanok
\uses{def:binaryEntropy, def:hammingBall, lem:hammingBall_empty_of_neg, thm:hammingBallBound, thm:mem_hammingBall}
\textbf{Case} $\delta n < 1$: Then $\delta n - 1 < 0$, so $\operatorname{Ball}(n, \delta n - 1) = \emptyset$ by Lemma~\ref{lem:hammingBall_empty_of_neg}, and $0 < 2^m$.

\textbf{Case} $\delta n \geq 1$: Set $d' = \lfloor \delta n \rfloor - 1$. We verify $d' + 1 \leq \delta n$ using $\lfloor \delta n \rfloor \geq 1$ and the floor estimate. We show $\operatorname{Ball}(n, \delta n - 1) \subseteq \operatorname{Ball}(n, d')$ using Theorem~\ref{thm:mem_hammingBall}: for $v \in \operatorname{Ball}(n, \delta n - 1)$ we have $\operatorname{wt}(v) < \delta n - 1 < \lfloor \delta n \rfloor$, so $\operatorname{wt}(v) \leq d'$. Then:
\[
|\operatorname{Ball}(n, \delta n - 1)| \leq |\operatorname{Ball}(n, d')| \leq 2^{h(\delta) \cdot n} < 2^{(m:\mathbb{R})} = 2^m,
\]
where the second inequality is Theorem~\ref{thm:hammingBallBound} (since $d' + 1 \leq \delta n$), the third uses $h(\delta) \cdot n < m$ and strict monotonicity of $x \mapsto 2^x$.
\end{proof}

\begin{lemma}[Existence of Code with Both Good Distance and Dual Distance]
\label{lem:exists_code_with_both_distances}
\lean{exists_code_with_both_distances}
\leanok
\uses{def:ClassicalCode, def:ClassicalCode.dimension, def:ClassicalCode.dualCode, def:binaryEntropy, def:hammingBall, thm:ClassicalCode.code_eq_ker, thm:ClassicalCode.dimension_dualCode, thm:exists_good_surjective_map, lem:hammingBall_lt_pow, lem:distance_ge_of_no_lowweight}
Let $\delta \in \mathbb{R}$ with $0 < \delta \leq 1/2$, and $n, k \in \mathbb{N}$ with $0 < n$, $0 < k$, $k < n$. Suppose
\[
k + 1 < (1 - h(\delta)) \cdot n \quad \text{and} \quad n - k + 1 < (1 - h(\delta)) \cdot n.
\]
Then there exists a classical code $\mathcal{C}$ of length $n$ with $\dim(\mathcal{C}) = k$, $d(\mathcal{C}) \geq \delta n$, and $d(\mathcal{C}^\perp) \geq \delta n$.
\end{lemma}

\begin{proof}
\leanok
\uses{def:ClassicalCode, def:ClassicalCode.dimension, def:ClassicalCode.dualCode, def:binaryEntropy, def:hammingBall, thm:ClassicalCode.code_eq_ker, thm:ClassicalCode.dimension_dualCode, thm:exists_good_surjective_map, lem:hammingBall_lt_pow, lem:distance_ge_of_no_lowweight}

\textbf{\color{orange}[Uses unproven axiom: thm:exists\_good\_surjective\_map]}

Set $r = n - k$, so $0 < r \leq n$ and $n - r = k$.

\textit{Ball bounds}: From $k + 1 < (1-h(\delta))n$ we get $h(\delta) \cdot n < r$, so by Lemma~\ref{lem:hammingBall_lt_pow}:
\[
|\operatorname{Ball}(n, \delta n - 1)| < 2^r.
\]
From $n - k + 1 < (1-h(\delta))n$ we get $h(\delta) \cdot n < k = n - r$, so similarly:
\[
|\operatorname{Ball}(n, \delta n - 1)| < 2^{n-r}.
\]

\textit{Applying the axiom}: By Theorem~\ref{thm:exists_good_surjective_map}, there exists a surjective $\phi : \mathbb{F}_2^n \to \mathbb{F}_2^r$ with: (i) every nonzero $v$ with $\operatorname{wt}(v) < \delta n$ satisfies $\phi(v) \neq 0$; (ii) every $w \in (\ker\phi)^\perp$ with $w \neq 0$ satisfies $\operatorname{wt}(w) \geq \delta n$.

Let $\mathcal{C} = \texttt{ofParityCheck}(\phi) = \ker \phi$.

\textit{Dimension}: Since $\phi$ is surjective, $\operatorname{finrank}(\operatorname{range}(\phi)) = r$ (the range equals $\mathbb{F}_2^r$). By rank--nullity, $\operatorname{finrank}(\ker\phi) = n - r = k$.

\textit{Distance}: For nonzero $v \in \ker\phi = \mathcal{C}$, we have $\phi(v) = 0$. If $\operatorname{wt}(v) < \delta n$, property (i) would give $\phi(v) \neq 0$, a contradiction. Hence $\operatorname{wt}(v) \geq \delta n$. Since $\operatorname{finrank}(\ker\phi) = k \geq 1$, $\mathcal{C}$ is nontrivial ($\texttt{Module.finrank\_pos\_iff}$), so there exists a nonzero codeword. By Lemma~\ref{lem:distance_ge_of_no_lowweight}, $d(\mathcal{C}) \geq \delta n$.

\textit{Dual distance}: By Theorem~\ref{thm:ClassicalCode.dimension_dualCode} and Theorem~\ref{thm:ClassicalCode.code_eq_ker}, we compute $\dim(\mathcal{C}^\perp) = n - k = r \geq 1$, so $\mathcal{C}^\perp$ is nontrivial. Property (ii) gives that every nonzero $w \in \mathcal{C}^\perp$ has $\operatorname{wt}(w) \geq \delta n$. By Lemma~\ref{lem:distance_ge_of_no_lowweight}, $d(\mathcal{C}^\perp) \geq \delta n$.
\end{proof}

\begin{lemma}[Satisfiability of Both Distance Premises]
\label{lem:exists_code_with_both_distances_premises_satisfiable}
\lean{exists_code_with_both_distances_premises_satisfiable}
\leanok
\uses{def:binaryEntropy, lem:binaryEntropy_at_one_hundredth_tight}
There exist $\delta \in \mathbb{R}$ and $n, k \in \mathbb{N}$ satisfying all premises of Lemma~\ref{lem:exists_code_with_both_distances}: $0 < \delta$, $\delta \leq 1/2$, $0 < n$, $0 < k$, $k < n$, $k + 1 < (1-h(\delta))n$, and $n - k + 1 < (1-h(\delta))n$.
\end{lemma}

\begin{proof}
\leanok
\uses{def:binaryEntropy, lem:binaryEntropy_at_one_hundredth_tight}
We exhibit the witness $\delta = 1/100$, $n = 1000$, $k = 500$. All numerical bounds are checked by $\texttt{norm\_num}$. Both inequalities $501 < (1-h(1/100)) \cdot 1000$ and $501 < (1-h(1/100)) \cdot 1000$ follow from $h(1/100) < 499/1000$ (Lemma~\ref{lem:binaryEntropy_at_one_hundredth_tight}) by linear arithmetic.
\end{proof}

\begin{theorem}[Gilbert--Varshamov Plus]
\label{thm:gilbertVarshamovPlus}
\lean{gilbertVarshamovPlus}
\leanok
\uses{def:ClassicalCode, def:ClassicalCode.dimension, def:ClassicalCode.dualCode, def:binaryEntropy, def:hammingBall, lem:hammingBall_empty_of_neg, thm:ClassicalCode.code_eq_ker, thm:ClassicalCode.dimension_dualCode, thm:ClassicalCode.mem_dualCode_iff, thm:binaryEntropy_lt_half, thm:hammingBallBound, thm:mem_hammingBall}
Let $\delta \in \mathbb{R}$ with $0 < \delta < 11/100$, and let $n \in \mathbb{N}$ with $0 < n$ and
\[
n > \frac{2}{1/2 - h(\delta)}.
\]
Then there exists a classical code $\mathcal{C}$ of length $n$ such that:
\begin{enumerate}
\item $\dim(\mathcal{C}) > n/2$,
\item $d(\mathcal{C}) \geq \delta n$,
\item $d(\mathcal{C}^\perp) \geq \delta n$.
\end{enumerate}
\end{theorem}

\begin{proof}
\leanok
\uses{def:ClassicalCode, def:ClassicalCode.dimension, def:ClassicalCode.dualCode, def:binaryEntropy, def:hammingBall, lem:hammingBall_empty_of_neg, lem:gv_k_exists, lem:gv_n_large, lem:exists_code_with_both_distances, thm:ClassicalCode.code_eq_ker, thm:ClassicalCode.dimension_dualCode, thm:ClassicalCode.mem_dualCode_iff, thm:binaryEntropy_lt_half, thm:hammingBallBound, thm:mem_hammingBall}

\textbf{\color{orange}[Uses unproven axiom: thm:exists\_good\_surjective\_map (via Lemma~\ref{lem:exists_code_with_both_distances})]}

From $\texttt{binaryEntropy\_lt\_half}$ applied to $0 < \delta$ and $\delta < 11/100$, we have $h(\delta) < 1/2$, so $\delta \leq 1/2$ by linear arithmetic.

\textit{Finding $k$}: By Lemma~\ref{lem:gv_k_exists}, there exists $k \in \mathbb{N}$ with
\[
n/2 < k, \quad k + 1 < (1 - h(\delta)) \cdot n, \quad k < n.
\]

\textit{Positivity of $k$}: We show $0 < k$ by contradiction. If $k = 0$, then the lower bound gives $n/2 < 0$, contradicting $n/2 \geq 0$.

\textit{Lower bound for dual condition}: By Lemma~\ref{lem:gv_n_large}, $(1/2 - h(\delta)) \cdot n > 2$. Hence:
\[
n - k + 1 < n - n/2 + 1 = n/2 + 1 < (1 - h(\delta)) \cdot n,
\]
where the last step uses $n/2 + 1 < (1-h(\delta)) \cdot n$, which follows from $(1/2 - h(\delta)) \cdot n > 2$ by linear arithmetic.

\textit{Applying the combined existence lemma}: By Lemma~\ref{lem:exists_code_with_both_distances} (with $\delta$, $n$, $k$ as above), there exists $\mathcal{C}$ with $\dim(\mathcal{C}) = k$, $d(\mathcal{C}) \geq \delta n$, and $d(\mathcal{C}^\perp) \geq \delta n$.

\textit{Conclusion}: We have $\dim(\mathcal{C}) = k > n/2$ (casting to reals from the bound on $k$), together with the two distance bounds.
\end{proof}

%--- Def_20: CayleyGraph ---
\chapter{Def 20: Cayley Graph}

\begin{definition}[Symmetric Generating Set]
\label{def:CayleyGraph.SymmGenSet}
\lean{CayleyGraph.SymmGenSet}
\leanok

A \emph{symmetric generating set} for a group $G$ is a structure consisting of:
\begin{itemize}
  \item A finite set of generators $S \subseteq G$ (the \emph{carrier}),
  \item A proof that the identity $1 \notin S$,
  \item A proof that $S$ is closed under inverses: for all $s \in S$, $s^{-1} \in S$.
\end{itemize}
\end{definition}

\begin{definition}[Cayley Graph]
\label{def:CayleyGraph.cayleyGraph}
\lean{CayleyGraph.cayleyGraph}
\leanok
\uses{def:CayleyGraph.SymmGenSet}
Let $G$ be a finite group with decidable equality, and let $S$ be a symmetric generating set for $G$. The \emph{Cayley graph} $\operatorname{Cay}(G, S)$ is the simple graph with vertex set $G$, where two vertices $g, g' \in G$ are adjacent if and only if there exists $s \in S$ such that $g' = s \cdot g$.

Well-definedness as a simple graph follows from:
\begin{itemize}
  \item \textbf{Symmetry}: If $g' = s \cdot g$ for some $s \in S$, then $g = s^{-1} \cdot g'$ and $s^{-1} \in S$ by the inverse-closure of $S$.
  \item \textbf{Looplessness}: If $g = s \cdot g$ for some $s \in S$, then $s = 1$, contradicting $1 \notin S$.
\end{itemize}
\end{definition}

\begin{theorem}[Cayley Graph is Regular]
\label{thm:CayleyGraph.cayleyGraph_isRegularOfDegree}
\lean{CayleyGraph.cayleyGraph_isRegularOfDegree}
\leanok
\uses{def:CayleyGraph.cayleyGraph, def:CayleyGraph.SymmGenSet}
The Cayley graph $\operatorname{Cay}(G, S)$ is $|S|$-regular: every vertex has degree exactly $|S|$.
\end{theorem}

\begin{proof}
\leanok
\uses{def:CayleyGraph.cayleyGraph, def:CayleyGraph.SymmGenSet, lem:CayleyGraph.cayleyGraph_neighborFinset}
Let $v \in G$ be an arbitrary vertex. We show that the neighbor set of $v$ equals the image $\{s \cdot v \mid s \in S\}$. By extensionality, a vertex $w$ belongs to the neighbor finset of $v$ if and only if there exists $s \in S$ with $w = s \cdot v$, which holds if and only if $w \in S.\operatorname{image}(\cdot \times v)$. This establishes the equality of neighbor finsets.

It then follows that
\[
  \deg(v) = \bigl|\{s \cdot v \mid s \in S\}\bigr| = |S|,
\]
where the last equality holds because the map $s \mapsto s \cdot v$ is injective: if $s \cdot v = s' \cdot v$, then right-cancellation gives $s = s'$. Since $v$ was arbitrary, the graph is $|S|$-regular.
\end{proof}

\begin{lemma}[Neighbor Finset of Cayley Graph]
\label{lem:CayleyGraph.cayleyGraph_neighborFinset}
\lean{CayleyGraph.cayleyGraph_neighborFinset}
\leanok
\uses{def:CayleyGraph.cayleyGraph, def:CayleyGraph.SymmGenSet}
For any vertex $v \in G$, the neighbor finset of $v$ in $\operatorname{Cay}(G, S)$ equals the image of the carrier under left-multiplication by $v$:
\[
  \operatorname{neighborFinset}(v) = \{s \cdot v \mid s \in S\}.
\]
\end{lemma}

\begin{proof}
\leanok
\uses{def:CayleyGraph.cayleyGraph, def:CayleyGraph.SymmGenSet}
By extensionality, a vertex $w$ is in the neighbor finset of $v$ if and only if $w$ and $v$ are adjacent in $\operatorname{Cay}(G, S)$, i.e., there exists $s \in S$ with $w = s \cdot v$. This is precisely the condition for $w \in \{s \cdot v \mid s \in S\}$. Both directions follow immediately by unfolding the definitions.
\end{proof}

\begin{lemma}[Cayley Graph Has Edges When $S$ Is Nonempty]
\label{lem:CayleyGraph.cayleyGraph_nonempty_of_nonempty}
\lean{CayleyGraph.cayleyGraph_nonempty_of_nonempty}
\leanok
\uses{def:CayleyGraph.cayleyGraph, def:CayleyGraph.SymmGenSet}
If the carrier $S$ is nonempty, then the Cayley graph $\operatorname{Cay}(G, S)$ has at least one edge: there exist $g, g' \in G$ such that $g$ and $g'$ are adjacent.
\end{lemma}

\begin{proof}
\leanok
\uses{def:CayleyGraph.cayleyGraph, def:CayleyGraph.SymmGenSet}
Since $S$ is nonempty, we obtain some $s \in S$. We then exhibit the adjacency $1 \sim s \cdot 1$ in $\operatorname{Cay}(G, S)$, witnessed by $s \in S$ and the equality $s \cdot 1 = s \cdot 1$ (by reflexivity). Hence there exist vertices $g = 1$ and $g' = s \cdot 1$ that are adjacent.
\end{proof}

%--- Def_21: LPSExpanderGraphs ---
\chapter{Def 21: LPS Expander Graphs}

\begin{definition}[Projective General Linear Group]
\label{def:LPS.PGL2}
\lean{LPS.PGL2}
\leanok

The \emph{projective general linear group} $\operatorname{PGL}(2, q) = \operatorname{GL}(2, \mathbb{F}_q) / Z(\operatorname{GL}(2, \mathbb{F}_q))$ is defined as the quotient of $\operatorname{GL}(2, \mathbb{Z}/q\mathbb{Z})$ by its center:
\[
  \operatorname{PGL}_2(q) \;:=\; \operatorname{GL}(\operatorname{Fin}~2,\, \mathbb{Z}/q\mathbb{Z}) \,\big/\, Z\!\left(\operatorname{GL}(\operatorname{Fin}~2,\, \mathbb{Z}/q\mathbb{Z})\right).
\]
\end{definition}

\begin{definition}[Projective Special Linear Group]
\label{def:LPS.PSL2}
\lean{LPS.PSL2}
\leanok

The \emph{projective special linear group} $\operatorname{PSL}(2, q) = \operatorname{SL}(2, \mathbb{F}_q) / Z(\operatorname{SL}(2, \mathbb{F}_q))$ is defined as the quotient of $\operatorname{SL}(2, \mathbb{Z}/q\mathbb{Z})$ by its center:
\[
  \operatorname{PSL}_2(q) \;:=\; \operatorname{SL}(\operatorname{Fin}~2,\, \mathbb{Z}/q\mathbb{Z}) \,\big/\, Z\!\left(\operatorname{SL}(\operatorname{Fin}~2,\, \mathbb{Z}/q\mathbb{Z})\right).
\]
\end{definition}

\begin{definition}[The Set $\widetilde{S}_p$]
\label{def:LPS.SpTilde}
\lean{LPS.SpTilde}
\leanok

For a natural number $p$, define the set
\[
  \widetilde{S}_p \;:=\; \bigl\{(a,b,c,d) \in \mathbb{Z}^4 \;\big|\; a^2 + b^2 + c^2 + d^2 = p\bigr\}.
\]
\end{definition}

\begin{definition}[The Normalized Subset $S_p$]
\label{def:LPS.Sp}
\lean{LPS.Sp}
\leanok
\uses{def:LPS.SpTilde}
For an odd prime $p$ (i.e.\ $p \equiv 1$ or $3 \pmod{4}$), the \emph{normalized subset} $S_p \subseteq \widetilde{S}_p$ is the finite set of integer $4$-tuples $(a,b,c,d)$ with $a^2+b^2+c^2+d^2 = p$ satisfying the normalization condition:
\begin{itemize}
  \item If $p \equiv 1 \pmod{4}$: $a > 0$ and $a$ is odd.
  \item If $p \equiv 3 \pmod{4}$: $a$ is even and the first nonzero coordinate among $(a,b,c,d)$ is positive.
\end{itemize}
Concretely, $S_p$ is obtained by filtering $\{-p, \ldots, p\}^4$ by the predicate $\operatorname{SpPred}(p)$.
\end{definition}

\begin{lemma}[Every Element of $S_p$ Lies in $\widetilde{S}_p$]
\label{lem:LPS.Sp_subset_SpTilde}
\lean{LPS.Sp_subset_SpTilde}
\leanok
\uses{def:LPS.Sp, def:LPS.SpTilde}
For every odd prime $p$ and every $t \in S_p$, we have $t \in \widetilde{S}_p$.
\end{lemma}

\begin{proof}
\leanok
\uses{def:LPS.Sp, def:LPS.SpTilde}
Let $t \in S_p$. Unfolding the definition of $S_p$ as a filter, the first component of the predicate $\operatorname{SpPred}(p)$ requires precisely that $t.1^2 + t.2.1^2 + t.2.2.1^2 + t.2.2.2^2 = p$, which is the membership condition for $\widetilde{S}_p$. Hence $t \in \widetilde{S}_p$.
\end{proof}

\begin{theorem}[Axiom: Cardinality of $S_p$]
\label{thm:LPS.Sp_card}
\lean{LPS.Sp_card}
\leanok
\uses{def:LPS.Sp}

\textbf{\color{red}[UNPROVEN AXIOM]}

For every odd prime $p$,
\[
  |S_p| = p + 1.
\]

\textit{Note: This follows from Jacobi's four-square theorem: the number of representations of an odd prime $p$ as a sum of four integer squares is $r_4(p) = 8(p+1)$, and the normalization conditions select exactly $p+1$ representatives. Jacobi's four-square theorem is not available in Mathlib, so this is stated as an axiom.}
\end{theorem}

\begin{lemma}[$\widetilde{S}_p$ is Nonempty]
\label{lem:LPS.SpTilde_nonempty}
\lean{LPS.SpTilde_nonempty}
\leanok
\uses{def:LPS.SpTilde, def:LPS.Sp, thm:LPS.Sp_card, lem:LPS.Sp_subset_SpTilde}
For every odd prime $p$, $\widetilde{S}_p$ is nonempty.
\end{lemma}

\begin{proof}
\leanok
\uses{def:LPS.Sp, def:LPS.SpTilde, thm:LPS.Sp_card, lem:LPS.Sp_subset_SpTilde}
By $\operatorname{Sp_card}$, we have $|S_p| = p+1 \geq 2 > 0$, so $S_p$ is nonempty. Rewriting with the cardinality: if $S_p$ were empty then $|S_p| = 0 \neq p+1$ (since $p \geq 2$), a contradiction by integer arithmetic. Obtaining $t \in S_p$, by $\operatorname{Sp_subset_SpTilde}$, $t \in \widetilde{S}_p$, so $\widetilde{S}_p$ is nonempty.
\end{proof}

\begin{lemma}[$S_p$ is Nonempty]
\label{lem:LPS.Sp_nonempty}
\lean{LPS.Sp_nonempty}
\leanok
\uses{def:LPS.Sp, thm:LPS.Sp_card}
For every odd prime $p$, $S_p$ is nonempty.
\end{lemma}

\begin{proof}
\leanok
\uses{def:LPS.Sp, thm:LPS.Sp_card}
Assume for contradiction that $S_p = \emptyset$. Then $|S_p| = 0$. But by $\operatorname{Sp_card}$, $|S_p| = p+1$. Since $p \geq 2$, we have $p+1 \geq 3 > 0$, contradicting $|S_p| = 0$ by integer arithmetic.
\end{proof}

\begin{definition}[LPS Matrix $M(a,b,c,d)$]
\label{def:LPS.lpsMatrixVal}
\lean{LPS.lpsMatrixVal}
\leanok

Given $x, y \in \mathbb{Z}/q\mathbb{Z}$ and a tuple $t = (a,b,c,d) \in \mathbb{Z}^4$, the \emph{LPS matrix} is the $2 \times 2$ matrix over $\mathbb{Z}/q\mathbb{Z}$ defined by
\[
  M(a,b,c,d) \;:=\;
  \begin{pmatrix}
    \iota(a) + \iota(b)x + \iota(d)y & -\iota(b)y + \iota(c) + \iota(d)x \\
    -\iota(b)y - \iota(c) + \iota(d)x & \iota(a) - \iota(b)x - \iota(d)y
  \end{pmatrix},
\]
where $\iota : \mathbb{Z} \to \mathbb{Z}/q\mathbb{Z}$ denotes the canonical ring homomorphism.
\end{definition}

\begin{lemma}[Determinant of the LPS Matrix]
\label{lem:LPS.lpsMatrixVal_det}
\lean{LPS.lpsMatrixVal_det}
\leanok
\uses{def:LPS.lpsMatrixVal}
For $x, y \in \mathbb{Z}/q\mathbb{Z}$ with $x^2 + y^2 + 1 = 0$ and any $t = (a,b,c,d) \in \mathbb{Z}^4$,
\[
  \det M(a,b,c,d) \;=\; \iota(a)^2 + \iota(b)^2 + \iota(c)^2 + \iota(d)^2.
\]
\end{lemma}

\begin{proof}
\leanok
\uses{def:LPS.lpsMatrixVal}
Unfolding the definition of $M(a,b,c,d)$ and computing the $2\times 2$ determinant using $\det \begin{pmatrix} p & q \\ r & s \end{pmatrix} = ps - qr$, the result follows by a linear combination using the hypothesis $x^2 + y^2 + 1 = 0$: specifically, $\det M = \iota(a)^2 + \iota(b)^2 + \iota(c)^2 + \iota(d)^2$ after simplification via $-(\iota(b)^2 + \iota(d)^2)(x^2 + y^2 + 1) = 0$.
\end{proof}

\begin{lemma}[LPS Matrix is Invertible for $(a,b,c,d) \in \widetilde{S}_p$, $p \neq q$]
\label{lem:LPS.lpsMatrixVal_det_ne_zero}
\lean{LPS.lpsMatrixVal_det_ne_zero}
\leanok
\uses{def:LPS.lpsMatrixVal, def:LPS.SpTilde, lem:LPS.lpsMatrixVal_det}
For distinct odd primes $p, q$ and $x, y \in \mathbb{Z}/q\mathbb{Z}$ with $x^2 + y^2 + 1 = 0$, if $t \in \widetilde{S}_p$ then $\det M(t) \neq 0$.
\end{lemma}

\begin{proof}
\leanok
\uses{def:LPS.lpsMatrixVal, def:LPS.SpTilde, lem:LPS.lpsMatrixVal_det}
By \texttt{lpsMatrixVal\_det}, it suffices to show $\iota(a)^2 + \iota(b)^2 + \iota(c)^2 + \iota(d)^2 \neq 0$ in $\mathbb{Z}/q\mathbb{Z}$. Since $t \in \widetilde{S}_p$, we have $a^2 + b^2 + c^2 + d^2 = p$ in $\mathbb{Z}$. Applying $\iota : \mathbb{Z} \to \mathbb{Z}/q\mathbb{Z}$ and using $\operatorname{push_cast}$, the sum equals $\iota(p) = (p : \mathbb{Z}/q\mathbb{Z})$. We must show $(p : \mathbb{Z}/q\mathbb{Z}) \neq 0$, i.e.\ $q \nmid p$. Since $p$ and $q$ are distinct primes, by primality of $p$: $q \mid p$ implies $q = 1$ or $q = p$, the former contradicting $q$ prime and the latter contradicting $p \neq q$.
\end{proof}

\begin{definition}[Conjugate Tuple]
\label{def:LPS.negConj}
\lean{LPS.negConj}
\leanok

For $t = (a,b,c,d) \in \mathbb{Z}^4$, the \emph{conjugate tuple} is $\bar{t} := (a, -b, -c, -d)$.
\end{definition}

\begin{lemma}[Conjugate Tuple Lies in $\widetilde{S}_p$]
\label{lem:LPS.negConj_mem_SpTilde}
\lean{LPS.negConj_mem_SpTilde}
\leanok
\uses{def:LPS.negConj, def:LPS.SpTilde}
If $t \in \widetilde{S}_p$ then $\bar{t} \in \widetilde{S}_p$.
\end{lemma}

\begin{proof}
\leanok
\uses{def:LPS.negConj, def:LPS.SpTilde}
Since $t \in \widetilde{S}_p$, we have $a^2 + b^2 + c^2 + d^2 = p$. For $\bar{t} = (a,-b,-c,-d)$, we compute $a^2 + (-b)^2 + (-c)^2 + (-d)^2 = a^2 + b^2 + c^2 + d^2 = p$ by the fact that squaring negates signs, which follows by a nonlinear arithmetic argument using $b^2 \geq 0$.
\end{proof}

\begin{lemma}[LPS Matrix of Conjugate Equals Adjugate]
\label{lem:LPS.lpsMatrixVal_negConj}
\lean{LPS.lpsMatrixVal_negConj}
\leanok
\uses{def:LPS.lpsMatrixVal, def:LPS.negConj}
For any $x, y \in \mathbb{Z}/q\mathbb{Z}$ and $t \in \mathbb{Z}^4$,
\[
  M(\bar{t}) \;=\; \operatorname{adj}(M(t)),
\]
where $\operatorname{adj}$ denotes the matrix adjugate.
\end{lemma}

\begin{proof}
\leanok
\uses{def:LPS.lpsMatrixVal, def:LPS.negConj}
We check equality entry by entry using extensionality. For each of the four entries $(i,j) \in \{0,1\}^2$, unfolding the definitions of $M(\bar{t})$ (with $(a,-b,-c,-d)$) and of $\operatorname{adj}(M(t))$ (using the $2\times 2$ adjugate formula $\operatorname{adj}\begin{pmatrix}p&q\\r&s\end{pmatrix} = \begin{pmatrix}s&-q\\-r&p\end{pmatrix}$) and simplifying by ring arithmetic confirms equality.
\end{proof}

\begin{lemma}[LPS Matrix of Conjugate When $a=0$]
\label{lem:LPS.lpsMatrixVal_negConj_zero_fst}
\lean{LPS.lpsMatrixVal_negConj_zero_fst}
\leanok
\uses{def:LPS.lpsMatrixVal, def:LPS.negConj}
If $a = t.1 = 0$, then $M(\bar{t}) = -M(t)$.
\end{lemma}

\begin{proof}
\leanok
\uses{def:LPS.lpsMatrixVal, def:LPS.negConj}
When $a = 0$, entry-by-entry verification using the definitions of $M(\bar{t}) = M(0,-b,-c,-d)$ and $-M(t) = -M(0,b,c,d)$ shows equality via ring arithmetic.
\end{proof}

\begin{lemma}[Conjugate of an Element in $S_p$ with $a > 0$ Stays in $S_p$]
\label{lem:LPS.negConj_mem_Sp_of_fst_pos}
\lean{LPS.negConj_mem_Sp_of_fst_pos}
\leanok
\uses{def:LPS.Sp, def:LPS.negConj}
Let $p$ be an odd prime, $t \in S_p$, and suppose $a = t.1 > 0$. Then $\bar{t} \in S_p$.
\end{lemma}

\begin{proof}
\leanok
\uses{def:LPS.Sp, def:LPS.negConj}
We unfold the membership condition for $S_p$ for both $t$ and $\bar{t} = (a,-b,-c,-d)$, obtaining bounds and the predicate $\operatorname{SpPred}$. The bounds $-p \leq \pm b, \pm c, \pm d \leq p$ are preserved by negation. For the normalization condition: in both the $p \equiv 1 \pmod{4}$ and $p \equiv 3 \pmod{4}$ cases, since $a > 0$ is unchanged and is the first coordinate, the condition $a > 0$ (and $a$ odd if $p \equiv 1$) is preserved, and $\operatorname{firstNonzeroPositive}(\bar{t}) = \operatorname{Or.inl}(a > 0)$ holds.
\end{proof}

\begin{definition}[Canonical Projection to $\operatorname{PGL}_2(q)$]
\label{def:LPS.projPGL}
\lean{LPS.projPGL}
\leanok
\uses{def:LPS.PGL2}
The canonical group homomorphism $\pi_{\operatorname{PGL}} : \operatorname{GL}(2, \mathbb{Z}/q\mathbb{Z}) \to \operatorname{PGL}_2(q)$ is the quotient map $\operatorname{QuotientGroup.mk'}$.
\end{definition}

\begin{definition}[Canonical Projection to $\operatorname{PSL}_2(q)$]
\label{def:LPS.projPSL}
\lean{LPS.projPSL}
\leanok
\uses{def:LPS.PSL2}
The canonical group homomorphism $\pi_{\operatorname{PSL}} : \operatorname{SL}(2, \mathbb{Z}/q\mathbb{Z}) \to \operatorname{PSL}_2(q)$ is the quotient map $\operatorname{QuotientGroup.mk'}$.
\end{definition}

\begin{definition}[Lift Matrix to $\operatorname{GL}$]
\label{def:LPS.matToGL}
\lean{LPS.matToGL}
\leanok

For a prime $q$ and a matrix $M \in M_2(\mathbb{Z}/q\mathbb{Z})$ with $\det M \neq 0$, the element $\operatorname{matToGL}(M) \in \operatorname{GL}(2, \mathbb{Z}/q\mathbb{Z})$ is the unit corresponding to $M$ (using the fact that invertibility of matrices over a field is equivalent to nonzero determinant).
\end{definition}

\begin{definition}[Map Tuple to $\operatorname{GL}$ Element]
\label{def:LPS.tupleToGLElement}
\lean{LPS.tupleToGLElement}
\leanok
\uses{def:LPS.matToGL, def:LPS.lpsMatrixVal, def:LPS.SpTilde, lem:LPS.lpsMatrixVal_det_ne_zero}
For distinct odd primes $p, q$, $x, y \in \mathbb{Z}/q\mathbb{Z}$ with $x^2 + y^2 + 1 = 0$, and $t \in \widetilde{S}_p$, define
\[
  \operatorname{tupleToGL}(t) \;:=\; \operatorname{matToGL}\!\left(M(t)\right) \;\in\; \operatorname{GL}(2, \mathbb{Z}/q\mathbb{Z}),
\]
using $\det M(t) \neq 0$ from \texttt{lpsMatrixVal\_det\_ne\_zero}.
\end{definition}

\begin{lemma}[LPS Matrix Not in Center of $\operatorname{GL}$]
\label{lem:LPS.lpsMatrix_not_in_center}
\lean{LPS.lpsMatrix_not_in_center}
\leanok
\uses{def:LPS.tupleToGLElement, def:LPS.SpTilde}
Let $p, q$ be distinct odd primes with $q > 2\sqrt{p}$, and let $x, y \in \mathbb{Z}/q\mathbb{Z}$ satisfy $x^2 + y^2 + 1 = 0$. For any $t \in \widetilde{S}_p$, the element $\operatorname{tupleToGL}(t)$ does not lie in the center $Z(\operatorname{GL}(2, \mathbb{Z}/q\mathbb{Z}))$.
\end{lemma}

\begin{proof}
\leanok
\uses{def:LPS.tupleToGLElement, def:LPS.SpTilde, def:LPS.lpsMatrixVal}
Suppose for contradiction that $M = \operatorname{tupleToGL}(t) \in Z(\operatorname{GL}(2, \mathbb{Z}/q\mathbb{Z}))$, i.e.\ $M$ commutes with every element. We construct the elementary matrices $E_{12} = \begin{pmatrix}1&1\\0&1\end{pmatrix}$ and $E_{21} = \begin{pmatrix}1&0\\1&1\end{pmatrix}$ as $\operatorname{GL}$ elements (their determinants equal $1 \neq 0$). From commutativity $E_{12} M_{\text{mat}} = M_{\text{mat}} E_{12}$: entry $(1,0)$ gives $M_{10} = 0$ and entries $(0,0),(1,1)$ give $M_{00} = M_{11}$. From $E_{21} M_{\text{mat}} = M_{\text{mat}} E_{21}$: entry $(0,1)$ gives $M_{01} = 0$.

Translating to the LPS matrix entries with $\iota(b) = \bar{b}$, $\iota(c) = \bar{c}$, $\iota(d) = \bar{d}$: the equations yield $\bar{c} = 0$ (from $2\bar{c} = 0$ and $2 \neq 0$ in $\mathbb{Z}/q\mathbb{Z}$ since $q$ is odd), $\bar{d}x = \bar{b}y$, and $\bar{b}x + \bar{d}y = 0$.

Since $x^2 + y^2 = -1$, multiplying: $\bar{b}(x^2+y^2) = 0$ by a linear combination $-y(\bar{d}x - \bar{b}y) + x(\bar{b}x + \bar{d}y) = 0$, giving $-\bar{b} = 0$ so $\bar{b} = 0$. Similarly $\bar{d} = 0$ (if $\bar{d} \neq 0$ then both $x = 0$ and $y = 0$ from the equations, but then $x^2 + y^2 = 0 \neq -1$, contradiction). 

Now $b, c, d$ all have $\mathbb{Z}/q\mathbb{Z}$-images zero, and each satisfies $b^2 \leq p$ (from $a^2 + b^2 + c^2 + d^2 = p$). Since $q > 2\sqrt{p}$, the lemma \texttt{int\_eq\_zero\_of\_zmod\_eq\_zero\_of\_sq\_le} implies $b = c = d = 0$ in $\mathbb{Z}$. Then $a^2 = p$, contradicting the fact that no prime is a perfect square (\texttt{Nat.Prime.not\_perfect\_square}).
\end{proof}

\begin{lemma}[Projection of Conjugate Equals Inverse in $\operatorname{PGL}$]
\label{lem:LPS.projPGL_negConj_eq_inv}
\lean{LPS.projPGL_negConj_eq_inv}
\leanok
\uses{def:LPS.projPGL, def:LPS.tupleToGLElement, def:LPS.negConj, lem:LPS.negConj_mem_SpTilde, lem:LPS.lpsMatrixVal_negConj}
For distinct odd primes $p, q$, $x, y$ with $x^2+y^2+1=0$, and $t \in \widetilde{S}_p$:
\[
  \pi_{\operatorname{PGL}}(\operatorname{tupleToGL}(\bar{t})) \;=\; \pi_{\operatorname{PGL}}(\operatorname{tupleToGL}(t))^{-1}.
\]
\end{lemma}

\begin{proof}
\leanok
\uses{def:LPS.projPGL, def:LPS.tupleToGLElement, def:LPS.negConj, lem:LPS.negConj_mem_SpTilde, lem:LPS.lpsMatrixVal_negConj}
It suffices to show $\pi_{\operatorname{PGL}}(g \cdot g') = 1$ where $g = \operatorname{tupleToGL}(t)$ and $g' = \operatorname{tupleToGL}(\bar{t})$. We compute $g \cdot g'$ using $g'.{\rm val} = \operatorname{adj}(g.{\rm val})$ (from \texttt{lpsMatrixVal\_negConj}): $g.{\rm val} \cdot g'.{\rm val} = M(t) \cdot \operatorname{adj}(M(t)) = \det(M(t)) \cdot I$. Thus $g \cdot g'$ is a scalar matrix $\det(M(t)) \cdot I$. A scalar matrix lies in the center of $\operatorname{GL}(2, \mathbb{Z}/q\mathbb{Z})$, so it maps to $1 \in \operatorname{PGL}_2(q)$ via $\pi_{\operatorname{PGL}} = \operatorname{QuotientGroup.mk'}$.
\end{proof}

\begin{lemma}[Projection of Conjugate Equals Self When $a = 0$ in $\operatorname{PGL}$]
\label{lem:LPS.projPGL_negConj_eq_self_of_fst_zero}
\lean{LPS.projPGL_negConj_eq_self_of_fst_zero}
\leanok
\uses{def:LPS.projPGL, def:LPS.tupleToGLElement, def:LPS.negConj, lem:LPS.lpsMatrixVal_negConj_zero_fst, lem:LPS.projPGL_negConj_eq_inv}
If $t \in \widetilde{S}_p$ and $a = t.1 = 0$, then
\[
  \pi_{\operatorname{PGL}}(\operatorname{tupleToGL}(\bar{t})) \;=\; \pi_{\operatorname{PGL}}(\operatorname{tupleToGL}(t)).
\]
\end{lemma}

\begin{proof}
\leanok
\uses{def:LPS.projPGL, def:LPS.tupleToGLElement, def:LPS.negConj, lem:LPS.lpsMatrixVal_negConj_zero_fst, lem:LPS.projPGL_negConj_eq_inv}
Let $g = \operatorname{tupleToGL}(t)$ and $g' = \operatorname{tupleToGL}(\bar{t})$. We show $\pi_{\operatorname{PGL}}(g') = \pi_{\operatorname{PGL}}(g)$, equivalently $g^{\prime -1} g \in Z(\operatorname{GL}(2, \mathbb{Z}/q\mathbb{Z}))$. Since $a = 0$, by \texttt{lpsMatrixVal\_negConj\_zero\_fst}, $g'.{\rm val} = -g.{\rm val}$. We compute $(g^{\prime -1} g).{\rm val}$: from $(-g.{\rm val})(g^{\prime -1}.{\rm val} \cdot g.{\rm val}) = g.{\rm val}$ (using $g' \cdot g^{\prime -1} = 1$), it follows that $g.{\rm val} \cdot (g^{\prime -1}.{\rm val} \cdot g.{\rm val}) = -g.{\rm val}$. Multiplying on the left by $g^{-1}.{\rm val}$ and using $g^{-1} \cdot g = 1$, we get $(g^{\prime -1}g).{\rm val} = -(1)$. Since $-I$ commutes with every matrix, $g^{\prime -1}g \in Z(\operatorname{GL}(2, \mathbb{Z}/q\mathbb{Z}))$.
\end{proof}

\begin{lemma}[$S_p$ Elements Have Nonnegative First Coordinate]
\label{lem:LPS.Sp_fst_nonneg}
\lean{LPS.Sp_fst_nonneg}
\leanok
\uses{def:LPS.Sp}
For every odd prime $p$ and $t \in S_p$, we have $t.1 \geq 0$.
\end{lemma}

\begin{proof}
\leanok
\uses{def:LPS.Sp}
Unfolding the filter condition for $S_p$: in the case $p \equiv 1 \pmod{4}$, the predicate requires $a > 0$, so $a \geq 0$. In the case $p \equiv 3 \pmod{4}$, the predicate requires $\operatorname{firstNonzeroPositive}(t)$, which by case analysis gives either $a > 0$ (so $a \geq 0$), or $a = 0$ (so $a \geq 0$) in all remaining cases.
\end{proof}

\begin{definition}[LPS Generating Set in $\operatorname{PGL}_2(q)$]
\label{def:LPS.lpsGenSetPGL}
\lean{LPS.lpsGenSetPGL}
\leanok
\uses{def:CayleyGraph.SymmGenSet, def:LPS.PGL2, def:LPS.Sp, def:LPS.tupleToGLElement, def:LPS.projPGL, lem:LPS.lpsMatrix_not_in_center, lem:LPS.negConj_mem_Sp_of_fst_pos, lem:LPS.projPGL_negConj_eq_inv, lem:LPS.Sp_fst_nonneg}
For distinct odd primes $p, q$ with $q > 2\sqrt{p}$ and $x, y \in \mathbb{Z}/q\mathbb{Z}$ with $x^2+y^2+1=0$, the \emph{LPS generating set in $\operatorname{PGL}_2(q)$} is the symmetric generating set
\[
  S_{p,q}^{\operatorname{PGL}} \;:=\; \bigl\{\pi_{\operatorname{PGL}}(\operatorname{tupleToGL}(t)) \;\big|\; t \in S_p\bigr\} \;\subseteq\; \operatorname{PGL}_2(q).
\]
This is a \texttt{SymmGenSet} (i.e.\ closed under inverses and not containing the identity), verified by:
\begin{itemize}
  \item \emph{$1 \notin S_{p,q}^{\operatorname{PGL}}$}: If $\pi_{\operatorname{PGL}}(\operatorname{tupleToGL}(t)) = 1$, then $\operatorname{tupleToGL}(t) \in Z(\operatorname{GL})$, contradicting \texttt{lpsMatrix\_not\_in\_center}.
  \item \emph{Closure under inverses}: For $t \in S_p$ with $a > 0$, use $\bar{t} \in S_p$ and $\pi_{\operatorname{PGL}}(\operatorname{tupleToGL}(\bar{t})) = \pi_{\operatorname{PGL}}(\operatorname{tupleToGL}(t))^{-1}$. For $a = 0$, the element is its own inverse.
\end{itemize}
\end{definition}

\begin{definition}[Scale Matrix to $\operatorname{SL}$]
\label{def:LPS.matToSL}
\lean{LPS.matToSL}
\leanok

For a prime $q$, a matrix $M \in M_2(\mathbb{Z}/q\mathbb{Z})$, and $r \in \mathbb{Z}/q\mathbb{Z}$ with $r \neq 0$ and $r^2 = \det M$, define
\[
  \operatorname{matToSL}(M, r) \;:=\; r^{-1} M \;\in\; \operatorname{SL}(2, \mathbb{Z}/q\mathbb{Z}).
\]
This has determinant $1$ since $\det(r^{-1} M) = (r^{-1})^2 \det M = r^{-2} r^2 = 1$.
\end{definition}

\begin{definition}[Legendre Symbol Square Root]
\label{def:LPS.legendreSqrt}
\lean{LPS.legendreSqrt}
\leanok

For distinct odd primes $p, q$ with Legendre symbol $\left(\frac{p}{q}\right) = 1$, the element $\operatorname{legendreSqrt}(q, p) \in \mathbb{Z}/q\mathbb{Z}$ is the square root of $p$ in $\mathbb{Z}/q\mathbb{Z}$, obtained via \texttt{ZMod.isSquare\_of\_jacobiSym\_eq\_one}.
\end{definition}

\begin{lemma}[Legendre Square Root Satisfies $r^2 = p$]
\label{lem:LPS.legendreSqrt_sq}
\lean{LPS.legendreSqrt_sq}
\leanok
\uses{def:LPS.legendreSqrt}
$\operatorname{legendreSqrt}(q,p)^2 = (p : \mathbb{Z}/q\mathbb{Z})$.
\end{lemma}

\begin{proof}
\leanok
\uses{def:LPS.legendreSqrt}
By the defining property of \texttt{legendreSqrt} via \texttt{ZMod.isSquare\_of\_jacobiSym\_eq\_one.choose\_spec}, the chosen element $r$ satisfies $r \cdot r = (p : \mathbb{Z}/q\mathbb{Z})$, which equals $r^2$ by definition of squaring.
\end{proof}

\begin{lemma}[Legendre Square Root is Nonzero]
\label{lem:LPS.legendreSqrt_ne_zero}
\lean{LPS.legendreSqrt_ne_zero}
\leanok
\uses{def:LPS.legendreSqrt, lem:LPS.legendreSqrt_sq}
$\operatorname{legendreSqrt}(q,p) \neq 0$.
\end{lemma}

\begin{proof}
\leanok
\uses{def:LPS.legendreSqrt, lem:LPS.legendreSqrt_sq}
Suppose $r = 0$. Then $r^2 = 0$. But by \texttt{legendreSqrt\_sq}, $r^2 = (p : \mathbb{Z}/q\mathbb{Z})$. So $(p : \mathbb{Z}/q\mathbb{Z}) = 0$, meaning $q \mid p$. Since $p, q$ are distinct primes, by primality of $p$: either $q = 1$ (contradicting $q$ prime) or $q = p$ (contradicting $p \neq q$). Contradiction.
\end{proof}

\begin{definition}[Map Tuple to $\operatorname{SL}$ Element]
\label{def:LPS.tupleToSLElement}
\lean{LPS.tupleToSLElement}
\leanok
\uses{def:LPS.matToSL, def:LPS.lpsMatrixVal, def:LPS.legendreSqrt, def:LPS.SpTilde, lem:LPS.legendreSqrt_sq, lem:LPS.legendreSqrt_ne_zero}
For distinct odd primes $p, q$ with $\left(\frac{p}{q}\right) = 1$, $x, y \in \mathbb{Z}/q\mathbb{Z}$ with $x^2+y^2+1=0$, and $t \in \widetilde{S}_p$, define
\[
  \operatorname{tupleToSL}(t) \;:=\; \operatorname{matToSL}\!\left(M(t),\, r\right) \;=\; r^{-1} M(t) \;\in\; \operatorname{SL}(2, \mathbb{Z}/q\mathbb{Z}),
\]
where $r = \operatorname{legendreSqrt}(q,p)$ satisfies $r^2 = \det M(t) = (p : \mathbb{Z}/q\mathbb{Z})$.
\end{definition}

\begin{lemma}[$\operatorname{tupleToSL}(\bar{t}) \cdot \operatorname{tupleToSL}(t) = 1$]
\label{lem:LPS.tupleToSL_negConj_mul_self_eq_one}
\lean{LPS.tupleToSL_negConj_mul_self_eq_one}
\leanok
\uses{def:LPS.tupleToSLElement, def:LPS.negConj, lem:LPS.negConj_mem_SpTilde, lem:LPS.lpsMatrixVal_negConj, lem:LPS.legendreSqrt_sq, lem:LPS.legendreSqrt_ne_zero}
For $t \in \widetilde{S}_p$:
\[
  \operatorname{tupleToSL}(\bar{t}) \cdot \operatorname{tupleToSL}(t) \;=\; 1 \;\in\; \operatorname{SL}(2, \mathbb{Z}/q\mathbb{Z}).
\]
\end{lemma}

\begin{proof}
\leanok
\uses{def:LPS.tupleToSLElement, def:LPS.negConj, lem:LPS.negConj_mem_SpTilde, lem:LPS.lpsMatrixVal_negConj, lem:LPS.legendreSqrt_sq, lem:LPS.legendreSqrt_ne_zero}
By definition, $\operatorname{tupleToSL}(\bar{t}).{\rm val} = r^{-1} \operatorname{adj}(M(t))$ (using \texttt{lpsMatrixVal\_negConj}) and $\operatorname{tupleToSL}(t).{\rm val} = r^{-1} M(t)$. Their product is $r^{-1} \operatorname{adj}(M(t)) \cdot r^{-1} M(t) = r^{-1} r^{-1} (\operatorname{adj}(M(t)) \cdot M(t)) = r^{-2} \det(M(t)) \cdot I$. By \texttt{legendreSqrt\_sq}, $\det M(t) = r^2$. So the product equals $r^{-2} r^2 \cdot I = I$. Setting $r^{-1} \cdot r^{-1} \cdot r^2 = (r^{-1} \cdot r^{-1} \cdot r) \cdot r = r^{-1} \cdot 1 \cdot r = 1$ by repeated use of $r^{-1} \cdot r = 1$.
\end{proof}

\begin{lemma}[Projection of Conjugate Equals Self When $a = 0$ in $\operatorname{PSL}$]
\label{lem:LPS.projPSL_negConj_eq_self_of_fst_zero}
\lean{LPS.projPSL_negConj_eq_self_of_fst_zero}
\leanok
\uses{def:LPS.projPSL, def:LPS.tupleToSLElement, def:LPS.negConj, lem:LPS.tupleToSL_negConj_mul_self_eq_one, lem:LPS.lpsMatrixVal_negConj_zero_fst, lem:LPS.legendreSqrt_sq, lem:LPS.legendreSqrt_ne_zero}
If $t \in \widetilde{S}_p$ and $a = t.1 = 0$, then
\[
  \pi_{\operatorname{PSL}}(\operatorname{tupleToSL}(\bar{t})) \;=\; \pi_{\operatorname{PSL}}(\operatorname{tupleToSL}(t)).
\]
\end{lemma}

\begin{proof}
\leanok
\uses{def:LPS.projPSL, def:LPS.tupleToSLElement, def:LPS.negConj, lem:LPS.tupleToSL_negConj_mul_self_eq_one, lem:LPS.lpsMatrixVal_negConj_zero_fst, lem:LPS.legendreSqrt_sq, lem:LPS.legendreSqrt_ne_zero}
Let $\operatorname{nc} = \operatorname{tupleToSL}(\bar{t})$ and $\operatorname{tt} = \operatorname{tupleToSL}(t)$. We show $\operatorname{nc}^{-1} \cdot \operatorname{tt} \in Z(\operatorname{SL}(2, \mathbb{Z}/q\mathbb{Z}))$.

By \texttt{tupleToSL\_negConj\_mul\_self\_eq\_one}, $\operatorname{nc} \cdot \operatorname{tt} = 1$, so $\operatorname{nc}^{-1} = \operatorname{tt}$, hence $\operatorname{nc}^{-1} \cdot \operatorname{tt} = \operatorname{tt}^2$.

We compute $(\operatorname{tt}^2).{\rm val}$: since $\operatorname{tt}.{\rm val} = r^{-1} M(t)$, we get $(\operatorname{tt}^2).{\rm val} = r^{-2} M(t)^2$. Since $a = 0$, $\operatorname{adj}(M(t)) = -M(t)$ by \texttt{lpsMatrixVal\_negConj\_zero\_fst}, so $M(t)^2 = -\det(M(t)) \cdot I$. Using $\det(M(t)) = r^2$: $(\operatorname{tt}^2).{\rm val} = r^{-2} \cdot (-r^2 \cdot I) = -I$.

Since $-I$ commutes with every matrix (as $(-I)g = g(-I)$ for all $g$), $\operatorname{tt}^2 \in Z(\operatorname{SL}(2, \mathbb{Z}/q\mathbb{Z}))$, completing the proof.
\end{proof}

\begin{definition}[LPS Generating Set in $\operatorname{PSL}_2(q)$]
\label{def:LPS.lpsGenSetPSL}
\lean{LPS.lpsGenSetPSL}
\leanok
\uses{def:CayleyGraph.SymmGenSet, def:LPS.PSL2, def:LPS.Sp, def:LPS.tupleToSLElement, def:LPS.projPSL, lem:LPS.tupleToSL_negConj_mul_self_eq_one, lem:LPS.negConj_mem_Sp_of_fst_pos, lem:LPS.Sp_fst_nonneg}
For distinct odd primes $p, q$ with $q > 2\sqrt{p}$ and $\left(\frac{p}{q}\right) = 1$, and $x, y \in \mathbb{Z}/q\mathbb{Z}$ with $x^2+y^2+1=0$, the \emph{LPS generating set in $\operatorname{PSL}_2(q)$} is the symmetric generating set
\[
  S_{p,q}^{\operatorname{PSL}} \;:=\; \bigl\{\pi_{\operatorname{PSL}}(\operatorname{tupleToSL}(t)) \;\big|\; t \in S_p\bigr\} \;\subseteq\; \operatorname{PSL}_2(q).
\]
This is a \texttt{SymmGenSet}, verified by the same argument as for $\operatorname{PGL}_2(q)$ but using the $\operatorname{SL}$ variants of the matrix lemmas.
\end{definition}

\begin{theorem}[Axiom: Cardinality of LPS Generating Set in $\operatorname{PGL}$]
\label{thm:LPS.lpsGenSetPGL_card}
\lean{LPS.lpsGenSetPGL_card}
\leanok
\uses{def:LPS.lpsGenSetPGL}

\textbf{\color{red}[UNPROVEN AXIOM]}

$|S_{p,q}^{\operatorname{PGL}}| = p + 1$.

\textit{Note: The injectivity of the map $S_p \to \operatorname{PGL}_2(q)$ (distinct tuples give distinct projective classes) requires deep number-theoretic arguments about quaternion arithmetic modulo $q$. Together with $|S_p| = p+1$, this gives $|S_{p,q}^{\operatorname{PGL}}| = p+1$. This is stated as an axiom because the required infrastructure is not available in Mathlib.}
\end{theorem}

\begin{theorem}[Axiom: Cardinality of LPS Generating Set in $\operatorname{PSL}$]
\label{thm:LPS.lpsGenSetPSL_card}
\lean{LPS.lpsGenSetPSL_card}
\leanok
\uses{def:LPS.lpsGenSetPSL}

\textbf{\color{red}[UNPROVEN AXIOM]}

$|S_{p,q}^{\operatorname{PSL}}| = p + 1$.

\textit{Note: As with $\operatorname{PGL}$, the injectivity of $S_p \to \operatorname{PSL}_2(q)$ via scaled matrices requires the same quaternion-arithmetic arguments. This is stated as an axiom.}
\end{theorem}

\begin{theorem}[Axiom: $S_{p,q}$ Generates $\operatorname{PGL}_2(q)$]
\label{thm:LPS.lpsGenSetPGL_generates}
\lean{LPS.lpsGenSetPGL_generates}
\leanok
\uses{def:LPS.lpsGenSetPGL, def:LPS.PGL2}

\textbf{\color{red}[UNPROVEN AXIOM]}

The subgroup generated by $S_{p,q}^{\operatorname{PGL}}$ equals all of $\operatorname{PGL}_2(q)$:
\[
  \langle S_{p,q}^{\operatorname{PGL}} \rangle \;=\; \operatorname{PGL}_2(q).
\]

\textit{Note: This deep result follows from the strong approximation theorem for quaternion algebras (Lubotzky--Phillips--Sarnak), which is far beyond current Mathlib infrastructure.}
\end{theorem}

\begin{theorem}[Axiom: $S_{p,q}$ Generates $\operatorname{PSL}_2(q)$]
\label{thm:LPS.lpsGenSetPSL_generates}
\lean{LPS.lpsGenSetPSL_generates}
\leanok
\uses{def:LPS.lpsGenSetPSL, def:LPS.PSL2}

\textbf{\color{red}[UNPROVEN AXIOM]}

The subgroup generated by $S_{p,q}^{\operatorname{PSL}}$ equals all of $\operatorname{PSL}_2(q)$:
\[
  \langle S_{p,q}^{\operatorname{PSL}} \rangle \;=\; \operatorname{PSL}_2(q).
\]

\textit{Note: Follows from the same strong approximation theorem for quaternion algebras (Lubotzky--Phillips--Sarnak).}
\end{theorem}

\begin{definition}[LPS Graph on $\operatorname{PGL}_2(q)$]
\label{def:LPS.lpsGraphPGL}
\lean{LPS.lpsGraphPGL}
\leanok
\uses{def:CayleyGraph.cayleyGraph, def:LPS.lpsGenSetPGL, def:LPS.PGL2}
The \emph{LPS graph on $\operatorname{PGL}_2(q)$} is the Cayley graph
\[
  \operatorname{Cay}(\operatorname{PGL}_2(q),\, S_{p,q}^{\operatorname{PGL}}) \;:=\; \operatorname{cayleyGraph}(S_{p,q}^{\operatorname{PGL}}).
\]
When $\left(\frac{p}{q}\right) = -1$ (i.e.\ the Legendre symbol is $-1$), this is the LPS expander graph on the group $\operatorname{PGL}_2(q)$.
\end{definition}

\begin{definition}[LPS Graph on $\operatorname{PSL}_2(q)$]
\label{def:LPS.lpsGraphPSL}
\lean{LPS.lpsGraphPSL}
\leanok
\uses{def:CayleyGraph.cayleyGraph, def:LPS.lpsGenSetPSL, def:LPS.PSL2}
The \emph{LPS graph on $\operatorname{PSL}_2(q)$} is the Cayley graph
\[
  \operatorname{Cay}(\operatorname{PSL}_2(q),\, S_{p,q}^{\operatorname{PSL}}) \;:=\; \operatorname{cayleyGraph}(S_{p,q}^{\operatorname{PSL}}).
\]
When $\left(\frac{p}{q}\right) = 1$ (i.e.\ the Legendre symbol is $1$), this is the LPS expander graph on the group $\operatorname{PSL}_2(q)$.
\end{definition}

\begin{theorem}[LPS Graph on $\operatorname{PGL}$ is $(p+1)$-Regular]
\label{thm:LPS.lpsGraphPGL_isRegularOfDegree}
\lean{LPS.lpsGraphPGL_isRegularOfDegree}
\leanok
\uses{def:CayleyGraph.SymmGenSet, def:CayleyGraph.cayleyGraph, thm:CayleyGraph.cayleyGraph_isRegularOfDegree, def:LPS.lpsGraphPGL, def:LPS.lpsGenSetPGL, thm:LPS.lpsGenSetPGL_card}
The LPS graph $\operatorname{Cay}(\operatorname{PGL}_2(q), S_{p,q}^{\operatorname{PGL}})$ is $(p+1)$-regular: every vertex has exactly $p+1$ neighbors.
\end{theorem}

\begin{proof}
\leanok
\uses{thm:CayleyGraph.cayleyGraph_isRegularOfDegree, def:LPS.lpsGraphPGL, def:LPS.lpsGenSetPGL, thm:LPS.lpsGenSetPGL_card}
By \texttt{cayleyGraph\_isRegularOfDegree}, the Cayley graph $\operatorname{cayleyGraph}(S)$ is $|S|$-regular for any symmetric generating set $S$. Applying this to $S = S_{p,q}^{\operatorname{PGL}}$ and rewriting with \texttt{lpsGenSetPGL\_card} (which gives $|S_{p,q}^{\operatorname{PGL}}| = p+1$), we obtain the $(p+1)$-regularity.
\end{proof}

\begin{theorem}[LPS Graph on $\operatorname{PSL}$ is $(p+1)$-Regular]
\label{thm:LPS.lpsGraphPSL_isRegularOfDegree}
\lean{LPS.lpsGraphPSL_isRegularOfDegree}
\leanok
\uses{def:CayleyGraph.SymmGenSet, def:CayleyGraph.cayleyGraph, thm:CayleyGraph.cayleyGraph_isRegularOfDegree, def:LPS.lpsGraphPSL, def:LPS.lpsGenSetPSL, thm:LPS.lpsGenSetPSL_card}
The LPS graph $\operatorname{Cay}(\operatorname{PSL}_2(q), S_{p,q}^{\operatorname{PSL}})$ is $(p+1)$-regular.
\end{theorem}

\begin{proof}
\leanok
\uses{thm:CayleyGraph.cayleyGraph_isRegularOfDegree, def:LPS.lpsGraphPSL, def:LPS.lpsGenSetPSL, thm:LPS.lpsGenSetPSL_card}
By \texttt{cayleyGraph\_isRegularOfDegree}, applying it to $S = S_{p,q}^{\operatorname{PSL}}$ and rewriting with \texttt{lpsGenSetPSL\_card} (giving $|S_{p,q}^{\operatorname{PSL}}| = p+1$) yields $(p+1)$-regularity.
\end{proof}

\begin{theorem}[Axiom: Independence of $(x,y)$ for $\operatorname{PGL}$ LPS Graph]
\label{thm:LPS.lpsGraphPGL_independent_of_xy}
\lean{LPS.lpsGraphPGL_independent_of_xy}
\leanok
\uses{def:LPS.lpsGraphPGL}

\textbf{\color{red}[UNPROVEN AXIOM]}

For any two pairs $(x_1, y_1)$ and $(x_2, y_2)$ in $\mathbb{Z}/q\mathbb{Z}$ both satisfying $x_i^2 + y_i^2 + 1 = 0$, the LPS graphs $\operatorname{Cay}(\operatorname{PGL}_2(q), S_{p,q}^{\operatorname{PGL}}(x_1,y_1))$ and $\operatorname{Cay}(\operatorname{PGL}_2(q), S_{p,q}^{\operatorname{PGL}}(x_2,y_2))$ are isomorphic as simple graphs.

\textit{Note: The proof requires constructing an explicit element of the orthogonal group of the norm form $x^2 + y^2 + 1$ over $\mathbb{F}_q$, conjugating one quaternion embedding to the other. This algebraic infrastructure is not available in Mathlib.}
\end{theorem}

\begin{theorem}[Axiom: Independence of $(x,y)$ for $\operatorname{PSL}$ LPS Graph]
\label{thm:LPS.lpsGraphPSL_independent_of_xy}
\lean{LPS.lpsGraphPSL_independent_of_xy}
\leanok
\uses{def:LPS.lpsGraphPSL}

\textbf{\color{red}[UNPROVEN AXIOM]}

For any two pairs $(x_1, y_1)$ and $(x_2, y_2)$ in $\mathbb{Z}/q\mathbb{Z}$ both satisfying $x_i^2 + y_i^2 + 1 = 0$, the LPS graphs $\operatorname{Cay}(\operatorname{PSL}_2(q), S_{p,q}^{\operatorname{PSL}}(x_1,y_1))$ and $\operatorname{Cay}(\operatorname{PSL}_2(q), S_{p,q}^{\operatorname{PSL}}(x_2,y_2))$ are isomorphic as simple graphs.

\textit{Note: The conjugation by $g \in \operatorname{GL}(2, \mathbb{F}_q)$ preserves $\operatorname{SL}(2, \mathbb{F}_q)$ up to scaling, hence descends to $\operatorname{PSL}_2(q)$. Requires the same orthogonal group infrastructure as the $\operatorname{PGL}$ case.}
\end{theorem}

%--- Thm_11: LPSRamanujan ---
\chapter{Thm 11: LPS Graphs are Ramanujan}

\begin{theorem}[$\operatorname{PGL}(2,q)$ Has At Least Six Elements]
\label{thm:LPSRamanujan.pgl2_card_ge_six}
\lean{LPSRamanujan.pgl2_card_ge_six}
\leanok
\uses{def:LPS.PGL2, def:LPS.matToGL, def:LPS.projPGL}

Let $q$ be an odd prime with $q \geq 5$. Then
\[
  6 \;\leq\; |\operatorname{PGL}(2,q)|.
\]
\end{theorem}

\begin{proof}
\leanok
\uses{def:LPS.PGL2, def:LPS.matToGL, def:LPS.projPGL}

We construct an injection from $\mathbb{Z}/q\mathbb{Z} \sqcup \{*\}$ into $\operatorname{PGL}(2,q)$ by sending $a \in \mathbb{Z}/q\mathbb{Z}$ to the class of the upper unipotent matrix
\[
  U(a) = \begin{pmatrix} 1 & a \\ 0 & 1 \end{pmatrix}
\]
and sending the single element of $\{*\}$ to the class of the swap matrix
\[
  S = \begin{pmatrix} 0 & 1 \\ -1 & 0 \end{pmatrix}.
\]

We verify injectivity in four cases. If $\operatorname{projPGL}(U(a)) = \operatorname{projPGL}(U(b))$, then there exists a central element $z$ in $\operatorname{GL}(2,q)$ with $U(a) \cdot z = U(b)$. Since $z$ is central, it commutes with $U(1) = \begin{pmatrix}1&1\\0&1\end{pmatrix}$; by expanding the commutation relation at entry $(1,0)$ we obtain $z_{10} = 0$, and at entry $(0,1)$ we obtain $z_{00} = z_{11}$. Similarly, commuting with $E_{21} = \begin{pmatrix}1&0\\1&1\end{pmatrix}$ yields $z_{01} = 0$. Substituting into the equation $U(a) \cdot z = U(b)$ and reading off entry $(0,0)$ gives $z_{00} = 1$, after which entry $(0,1)$ yields $a \cdot z_{11} = b$. Since $z_{00} = z_{11} = 1$, we conclude $a = b$.

If $\operatorname{projPGL}(S) = \operatorname{projPGL}(U(a))$, then $S \cdot z = U(a)$ for some central $z$. The same argument gives $z_{10} = 0$, $z_{01} = 0$. Reading entry $(1,0)$ of $S \cdot z = U(a)$ yields $-z_{00} = 0$, and entry $(1,1)$ yields $-z_{01} = 1$; but $z_{01} = 0$, giving $0 = 1$, a contradiction. The remaining cases are symmetric or trivial.

Thus the map is injective, and by counting we obtain $|\operatorname{PGL}(2,q)| \geq |\mathbb{Z}/q\mathbb{Z}| + 1 = q + 1 \geq 6$ for $q \geq 5$.
\end{proof}

\begin{theorem}[$\operatorname{PSL}(2,q)$ Has At Least Six Elements]
\label{thm:LPSRamanujan.psl2_card_ge_six}
\lean{LPSRamanujan.psl2_card_ge_six}
\leanok
\uses{def:LPS.PSL2, def:LPS.projPSL}

Let $q$ be an odd prime with $q \geq 5$. Then
\[
  6 \;\leq\; |\operatorname{PSL}(2,q)|.
\]
\end{theorem}

\begin{proof}
\leanok
\uses{def:LPS.PSL2, def:LPS.projPSL}

We construct an injection from $\mathbb{Z}/q\mathbb{Z} \sqcup \{*\}$ into $\operatorname{PSL}(2,q)$ by sending $a$ to the class of the upper unipotent $\operatorname{SL}$ element $U(a) = \bigl(\begin{smallmatrix}1&a\\0&1\end{smallmatrix}\bigr)$ and $*$ to the class of the swap $\operatorname{SL}$ element $S = \bigl(\begin{smallmatrix}0&1\\-1&0\end{smallmatrix}\bigr)$.

Injectivity is established by the same argument as in $\operatorname{PGL}$: if $\operatorname{projPSL}(U(a)) = \operatorname{projPSL}(U(b))$, there is a central $z \in \operatorname{SL}(2,q)$ with $U(a) \cdot z = U(b)$. Commuting $z$ with $E_{12}$ gives $z_{10} = 0$ and $z_{00} = z_{11}$; commuting with $E_{21}$ gives $z_{01} = 0$. Expanding $U(a) \cdot z = U(b)$ then forces $z_{00} = 1$ and $a = b$.

The case $\operatorname{projPSL}(S) = \operatorname{projPSL}(U(a))$ leads to $z_{01} = 0$ and the equation at entry $(1,1)$ of $S \cdot z = U(a)$ gives $-z_{01} = 1$, i.e.\ $0 = 1$, a contradiction. Hence the map is injective, so $|\operatorname{PSL}(2,q)| \geq q + 1 \geq 6$.
\end{proof}

\begin{lemma}[$q \geq 5$ From Spectral Condition]
\label{lem:LPSRamanujan.q_ge_five}
\lean{LPSRamanujan.q_ge_five}
\leanok
\uses{def:LPS.PGL2, def:LPS.PSL2}

Let $p$ and $q$ be odd primes with $q > 2\sqrt{p}$. Then $q \geq 5$.
\end{lemma}

\begin{proof}
\leanok
\uses{def:LPS.PGL2, def:LPS.PSL2}

We argue by contradiction. Suppose $q < 5$, so $q \leq 4$. Since $q$ is an odd prime and $q \leq 4$, we must have $q = 3$.

Since $p$ is an odd prime, $p \geq 3$, hence $(p : \mathbb{R}) \geq 3$. From the hypothesis $q > 2\sqrt{p}$ with $q = 3$, we have $3 > 2\sqrt{p}$, so $9 > 4p$. But $4p \geq 12$, giving $9 > 12$, a contradiction. Formally, let $s = \sqrt{p}$; then $s \geq 0$ and $s^2 = p \geq 3$. The inequality $3 > 2s$ gives $(3 - 2s)^2 > 0$ and $9 > 4s^2 = 4p \geq 12$, which is absurd; this is derived by nonlinear arithmetic from $\operatorname{nlinarith}$ using $(2\sqrt{p} - 3)^2 \geq 0$.
\end{proof}

\begin{theorem}[LPS Graph on $\operatorname{PGL}$ Has At Least Two Vertices]
\label{thm:LPSRamanujan.lpsGraphPGL_card_ge_two}
\lean{LPSRamanujan.lpsGraphPGL_card_ge_two}
\leanok
\uses{def:LPS.PGL2, def:LPS.lpsGraphPGL, def:LPS.projPGL}

Let $p, q$ be distinct odd primes with $q > 2\sqrt{p}$. Then
\[
  2 \;\leq\; |\operatorname{PGL}(2,q)|.
\]
\end{theorem}

\begin{proof}
\leanok
\uses{def:LPS.PGL2, def:LPS.projPGL, thm:LPSRamanujan.pgl2_card_ge_six, lem:LPSRamanujan.q_ge_five}

By Lemma~\ref{lem:LPSRamanujan.q_ge_five}, the hypothesis $q > 2\sqrt{p}$ together with $p, q$ odd primes implies $q \geq 5$. Applying Theorem~\ref{thm:LPSRamanujan.pgl2_card_ge_six} we obtain $|\operatorname{PGL}(2,q)| \geq 6 \geq 2$.
\end{proof}

\begin{theorem}[LPS Graph on $\operatorname{PSL}$ Has At Least Two Vertices]
\label{thm:LPSRamanujan.lpsGraphPSL_card_ge_two}
\lean{LPSRamanujan.lpsGraphPSL_card_ge_two}
\leanok
\uses{def:LPS.PSL2, def:LPS.lpsGraphPSL, def:LPS.projPSL}

Let $p, q$ be distinct odd primes with $q > 2\sqrt{p}$. Then
\[
  2 \;\leq\; |\operatorname{PSL}(2,q)|.
\]
\end{theorem}

\begin{proof}
\leanok
\uses{def:LPS.PSL2, def:LPS.projPSL, thm:LPSRamanujan.psl2_card_ge_six, lem:LPSRamanujan.q_ge_five}

By Lemma~\ref{lem:LPSRamanujan.q_ge_five}, $q \geq 5$. Applying Theorem~\ref{thm:LPSRamanujan.psl2_card_ge_six} gives $|\operatorname{PSL}(2,q)| \geq 6 \geq 2$.
\end{proof}

\begin{theorem}[$\lambda_2 \leq 2\sqrt{p}$ for LPS Graph on $\operatorname{PGL}$]
\label{thm:LPSRamanujan.lpsGraphPGL_lambda2_le}
\lean{LPSRamanujan.lpsGraphPGL_lambda2_le}
\leanok
\uses{def:GraphExpansion.IsRamanujan, def:GraphExpansion.secondLargestEigenvalue, def:LPS.PGL2, def:LPS.lpsGraphPGL, def:LPS.projPGL}

\textbf{\color{blue}[Hypothesis-based]} The following result is proved under explicit assumptions on the Ramanujan property and a spectral bound, which are deep results requiring algebraic geometry not available in Mathlib.

Let $p, q$ be distinct odd primes with $q > 2\sqrt{p}$, and let $x, y \in \mathbb{Z}/q\mathbb{Z}$ satisfy $x^2 + y^2 + 1 = 0$. Let $G = X_{p,q}^{\operatorname{PGL}}$ be the LPS graph on $\operatorname{PGL}(2,q)$. Suppose:
\begin{itemize}
  \item \textbf{(Ramanujan hypothesis)} $G$ is a Ramanujan graph of degree $p+1$, i.e., every eigenvalue $\lambda_i$ with $|\lambda_i| < p+1$ satisfies $|\lambda_i| \leq 2\sqrt{p}$.
  \item \textbf{(Spectral bound hypothesis)} $|\lambda_2(G)| < p+1$, where $\lambda_2$ denotes the second-largest eigenvalue.
\end{itemize}
Then
\[
  \lambda_2(G) \;\leq\; 2\sqrt{p}.
\]
\end{theorem}

\begin{proof}
\leanok
\uses{def:GraphExpansion.IsRamanujan, def:GraphExpansion.secondLargestEigenvalue, def:LPS.lpsGraphPGL, thm:LPSRamanujan.lpsGraphPGL_card_ge_two}

\textbf{\color{blue}[Uses hypothesis-based assumptions: Ramanujan property and spectral bound]}

We unfold the definition of $\operatorname{IsRamanujan}$ in the hypothesis $h_{\mathrm{ram}}$: for every index $i \in \operatorname{Fin}(|\operatorname{PGL}(2,q)|)$ with $|\lambda_i| < p+1$, we have $|\lambda_i| \leq 2\sqrt{p}$. Set the index $i = 1$, which is valid since $|\operatorname{PGL}(2,q)| \geq 2$. Applying $h_{\mathrm{ram}}$ at $i = 1$ and using the hypothesis $h_{\mathrm{spec}} : |\lambda_2| < p+1$ (after casting naturals to reals), we obtain $|\lambda_2| \leq 2\sqrt{p}$. Since $\lambda_2 \leq |\lambda_2|$, we conclude $\lambda_2 \leq 2\sqrt{p}$.
\end{proof}

\begin{remark}[Assumed Hypotheses for $\operatorname{PGL}$ Spectral Bound]
\textbf{\color{blue}Formalization note:} The hypothesis $h_{\mathrm{ram}}$ (the Ramanujan property of LPS graphs) follows from the Ramanujan--Petersson conjecture for $\operatorname{GL}(2)/\mathbb{Q}$, proved by Eichler--Igusa for weight 2 modular forms building on Deligne's proof of the Weil conjectures (1974). The hypothesis $h_{\mathrm{spec}}$ ($|\lambda_2| < p+1$ for connected $(p+1)$-regular graphs) follows from the Perron--Frobenius theorem. Both results require algebraic geometry and spectral theory not currently available in Mathlib, and are therefore taken as explicit hypotheses in this formalization.
\end{remark}

\begin{theorem}[$\lambda_2 \leq 2\sqrt{p}$ for LPS Graph on $\operatorname{PSL}$]
\label{thm:LPSRamanujan.lpsGraphPSL_lambda2_le}
\lean{LPSRamanujan.lpsGraphPSL_lambda2_le}
\leanok
\uses{def:GraphExpansion.IsRamanujan, def:GraphExpansion.secondLargestEigenvalue, def:LPS.PSL2, def:LPS.lpsGraphPSL, def:LPS.projPSL}

\textbf{\color{blue}[Hypothesis-based]} The following result is proved under explicit assumptions on the Ramanujan property and a spectral bound.

Let $p, q$ be distinct odd primes with $q > 2\sqrt{p}$ and $\left(\frac{p}{q}\right) = 1$ (i.e., $p$ is a quadratic residue mod $q$), and let $x, y \in \mathbb{Z}/q\mathbb{Z}$ satisfy $x^2 + y^2 + 1 = 0$. Let $G = X_{p,q}^{\operatorname{PSL}}$ be the LPS graph on $\operatorname{PSL}(2,q)$. Given:
\begin{itemize}
  \item \textbf{(Ramanujan hypothesis)} $G$ is Ramanujan of degree $p+1$.
  \item \textbf{(Spectral bound hypothesis)} $|\lambda_2(G)| < p+1$.
\end{itemize}
Then
\[
  \lambda_2(G) \;\leq\; 2\sqrt{p}.
\]
\end{theorem}

\begin{proof}
\leanok
\uses{def:GraphExpansion.IsRamanujan, def:GraphExpansion.secondLargestEigenvalue, def:LPS.lpsGraphPSL, thm:LPSRamanujan.lpsGraphPSL_card_ge_two}

\textbf{\color{blue}[Uses hypothesis-based assumptions: Ramanujan property and spectral bound]}

We unfold the definition of $\operatorname{IsRamanujan}$ and apply it at index $i = 1 \in \operatorname{Fin}(|\operatorname{PSL}(2,q)|)$, which is valid since $|\operatorname{PSL}(2,q)| \geq 2$. From the spectral bound hypothesis $|\lambda_2| < p+1$ (after casting), the Ramanujan condition yields $|\lambda_2| \leq 2\sqrt{p}$. The conclusion $\lambda_2 \leq 2\sqrt{p}$ follows since $\lambda_2 \leq |\lambda_2|$.
\end{proof}

\begin{remark}[Assumed Hypotheses for $\operatorname{PSL}$ Spectral Bound]
\textbf{\color{blue}Formalization note:} As for the $\operatorname{PGL}$ case, the Ramanujan property of LPS graphs on $\operatorname{PSL}(2,q)$ is a consequence of the Ramanujan--Petersson conjecture, proved via the Weil conjectures. The additional Legendre symbol condition $\left(\frac{p}{q}\right) = 1$ ensures the generators of the LPS graph lift to $\operatorname{PSL}(2,q)$. These deep results are taken as hypotheses.
\end{remark}

\begin{theorem}[LPS Ramanujan Theorem for $\operatorname{PGL}$, Combined]
\label{thm:LPSRamanujan.lpsRamanujan_PGL}
\lean{LPSRamanujan.lpsRamanujan_PGL}
\leanok
\uses{def:GraphExpansion.IsRamanujan, def:GraphExpansion.secondLargestEigenvalue, def:LPS.PGL2, def:LPS.lpsGraphPGL, def:LPS.matToGL, def:LPS.projPGL, thm:LPS.lpsGraphPGL_isRegularOfDegree}

\textbf{\color{blue}[Hypothesis-based]} The following combined result is proved under the hypotheses of connectedness and the Ramanujan property, which are deep results from Lubotzky--Phillips--Sarnak requiring algebraic geometry not available in Mathlib.

Let $p, q$ be distinct odd primes with $q > 2\sqrt{p}$, and let $x, y \in \mathbb{Z}/q\mathbb{Z}$ satisfy $x^2 + y^2 + 1 = 0$. Let $G = X_{p,q}^{\operatorname{PGL}}$ be the LPS graph on $\operatorname{PGL}(2,q)$. Given:
\begin{itemize}
  \item \textbf{(Connectedness hypothesis)} $G$ is connected.
  \item \textbf{(Ramanujan hypothesis)} $G$ is Ramanujan of degree $p+1$.
\end{itemize}
Then simultaneously:
\begin{enumerate}
  \item $G$ is connected,
  \item $G$ is $(p+1)$-regular,
  \item $|\operatorname{PGL}(2,q)| \geq 2$, and
  \item $G$ is a Ramanujan graph of degree $p+1$.
\end{enumerate}
\end{theorem}

\begin{proof}
\leanok
\uses{def:LPS.lpsGraphPGL, def:LPS.PGL2, def:GraphExpansion.IsRamanujan, thm:LPS.lpsGraphPGL_isRegularOfDegree, thm:LPSRamanujan.lpsGraphPGL_card_ge_two}

\textbf{\color{blue}[Uses hypothesis-based assumptions: connectedness and Ramanujan property]}

We combine the hypotheses and previously established results. Connectedness follows directly from $h_{\mathrm{conn}}$. Regularity of degree $p+1$ follows from Theorem~\ref{thm:LPS.lpsGraphPGL_isRegularOfDegree}. The cardinality bound $|\operatorname{PGL}(2,q)| \geq 2$ follows from Theorem~\ref{thm:LPSRamanujan.lpsGraphPGL_card_ge_two}. The Ramanujan property follows from $h_{\mathrm{ram}}$. We package all four conclusions into a conjunction.
\end{proof}

\begin{remark}[Assumed Hypotheses for the Combined PGL Result]
\textbf{\color{blue}Formalization note:} The connectedness of $X_{p,q}^{\operatorname{PGL}}$ follows from the fact that the generating set $S_{p,q}$ generates $\operatorname{PGL}(2,q)$ (Lubotzky--Phillips--Sarnak, 1988). The Ramanujan property follows from the Ramanujan--Petersson conjecture. Both require algebraic number theory and the theory of automorphic forms beyond what Mathlib currently provides; they are therefore stated as explicit hypotheses $h_{\mathrm{conn}}$ and $h_{\mathrm{ram}}$.
\end{remark}

\begin{theorem}[LPS Ramanujan Theorem for $\operatorname{PSL}$, Combined]
\label{thm:LPSRamanujan.lpsRamanujan_PSL}
\lean{LPSRamanujan.lpsRamanujan_PSL}
\leanok
\uses{def:GraphExpansion.IsRamanujan, def:GraphExpansion.secondLargestEigenvalue, def:LPS.PSL2, def:LPS.lpsGraphPSL, def:LPS.matToGL, def:LPS.projPSL, thm:LPS.lpsGraphPSL_isRegularOfDegree}

\textbf{\color{blue}[Hypothesis-based]} Let $p, q$ be distinct odd primes with $q > 2\sqrt{p}$ and $\left(\frac{p}{q}\right) = 1$, and let $x, y \in \mathbb{Z}/q\mathbb{Z}$ satisfy $x^2 + y^2 + 1 = 0$. Let $G = X_{p,q}^{\operatorname{PSL}}$ be the LPS graph on $\operatorname{PSL}(2,q)$. Given:
\begin{itemize}
  \item \textbf{(Connectedness hypothesis)} $G$ is connected.
  \item \textbf{(Ramanujan hypothesis)} $G$ is Ramanujan of degree $p+1$.
\end{itemize}
Then simultaneously:
\begin{enumerate}
  \item $G$ is connected,
  \item $G$ is $(p+1)$-regular,
  \item $|\operatorname{PSL}(2,q)| \geq 2$, and
  \item $G$ is a Ramanujan graph of degree $p+1$.
\end{enumerate}
\end{theorem}

\begin{proof}
\leanok
\uses{def:LPS.lpsGraphPSL, def:LPS.PSL2, def:GraphExpansion.IsRamanujan, thm:LPS.lpsGraphPSL_isRegularOfDegree, thm:LPSRamanujan.lpsGraphPSL_card_ge_two}

\textbf{\color{blue}[Uses hypothesis-based assumptions: connectedness and Ramanujan property]}

We combine the hypotheses and previously established results. Connectedness follows from $h_{\mathrm{conn}}$. Regularity of degree $p+1$ follows from Theorem~\ref{thm:LPS.lpsGraphPSL_isRegularOfDegree}. The cardinality bound $|\operatorname{PSL}(2,q)| \geq 2$ follows from Theorem~\ref{thm:LPSRamanujan.lpsGraphPSL_card_ge_two}. The Ramanujan property is $h_{\mathrm{ram}}$. All four conclusions are packaged into a conjunction.
\end{proof}

\begin{remark}[Assumed Hypotheses for the Combined PSL Result]
\textbf{\color{blue}Formalization note:} As in the $\operatorname{PGL}$ case, connectedness of $X_{p,q}^{\operatorname{PSL}}$ and its Ramanujan property are consequences of the Lubotzky--Phillips--Sarnak theory and the Ramanujan--Petersson conjecture. The Legendre symbol condition $\left(\frac{p}{q}\right) = 1$ ensures the generators lie in $\operatorname{PSL}(2,q)$. These results are taken as hypotheses $h_{\mathrm{conn}}$ and $h_{\mathrm{ram}}$ because the required algebraic geometry is not available in Mathlib.
\end{remark}

%--- Def_22: BalancedProductVectorSpaces ---
\chapter{Def 22: Balanced Product of Vector Spaces}

\begin{definition}[Balanced Product]
\label{def:balancedProduct}
\lean{balancedProduct}
\leanok

Let $H$ be a group, let $V$ and $W$ be $\mathbb{F}_2$-modules equipped with $H$-representations $\rho_V : H \to \operatorname{End}_{\mathbb{F}_2}(V)$ and $\rho_W : H \to \operatorname{End}_{\mathbb{F}_2}(W)$. The \emph{balanced product} $V \otimes_H W$ is defined as the coinvariants of the diagonal $H$-action on $V \otimes_{\mathbb{F}_2} W$:
\[
  V \otimes_H W \;:=\; \operatorname{Coinv}(\rho_V \otimes \rho_W),
\]
where the diagonal action is $h \cdot (v \otimes w) = \rho_V(h)(v) \otimes \rho_W(h)(w)$. Equivalently, this is the tensor product $V \otimes_{\mathbb{F}_2[H]} W$ over the group algebra. Concretely,
\[
  V \otimes_H W \;=\; (V \otimes_{\mathbb{F}_2} W) \,/\, \langle \rho_V(h)(v) \otimes \rho_W(h)(w) - v \otimes w \mid h \in H,\, v \in V,\, w \in W \rangle.
\]
\end{definition}

\begin{definition}[Quotient Map]
\label{def:balancedProduct.mk}
\lean{balancedProduct.mk}
\leanok
\uses{def:balancedProduct}
The \emph{quotient map} $\operatorname{mk} : V \otimes_{\mathbb{F}_2} W \to V \otimes_H W$ is the canonical $\mathbb{F}_2$-linear surjection onto the balanced product, defined as the coinvariants quotient map of the diagonal representation $\rho_V \otimes \rho_W$.
\end{definition}

\begin{theorem}[Diagonal Balance Relation]
\label{thm:balancedProduct.diagonal_balanced}
\lean{balancedProduct.diagonal_balanced}
\leanok
\uses{def:balancedProduct}
For all $h \in H$, $v \in V$, $w \in W$, the following equality holds in $V \otimes_H W$:
\[
  [\rho_V(h)(v) \otimes \rho_W(h)(w)] \;=\; [v \otimes w].
\]
\end{theorem}

\begin{proof}
\leanok
\uses{def:balancedProduct}
Rewriting using the definition of the diagonal tensor product representation, we have
\[
  \rho_V(h)(v) \otimes_{\mathbb{F}_2} \rho_W(h)(w) \;=\; (\rho_V \otimes \rho_W)(h)(v \otimes w),
\]
which follows by simplification using $\operatorname{tprod_apply}$ and $\operatorname{TensorProduct.map_tmul}$. The result then follows immediately from the defining property of coinvariants: $\operatorname{mk}((\rho_V \otimes \rho_W)(h)(x)) = \operatorname{mk}(x)$ for all $h \in H$ and $x \in V \otimes W$.
\end{proof}

\begin{theorem}[Balancing Relation]
\label{thm:balancedProduct.balanced}
\lean{balancedProduct.balanced}
\leanok
\uses{def:balancedProduct, thm:balancedProduct.diagonal_balanced}
For all $h \in H$, $v \in V$, $w \in W$, the following equality holds in $V \otimes_H W$:
\[
  [\rho_V(h^{-1})(v) \otimes w] \;=\; [v \otimes \rho_W(h)(w)].
\]
In right-action notation, this is the familiar relation $[v h \otimes w] = [v \otimes h w]$.
\end{theorem}

\begin{proof}
\leanok
\uses{def:balancedProduct, thm:balancedProduct.diagonal_balanced}
Apply the diagonal balance relation (Theorem~\ref{thm:balancedProduct.diagonal_balanced}) with $h$ and with $v$ replaced by $\rho_V(h^{-1})(v)$ and $w$ replaced by $w$, obtaining:
\[
  [\rho_V(h)(\rho_V(h^{-1})(v)) \otimes \rho_W(h)(w)] \;=\; [\rho_V(h^{-1})(v) \otimes w].
\]
We establish that $\rho_V(h)(\rho_V(h^{-1})(v)) = v$. Indeed, $(\rho_V(h) \circ \rho_V(h^{-1}))(v) = (\rho_V(h \cdot h^{-1}))(v)$ by the homomorphism property, and since $h \cdot h^{-1} = 1$ we get $\rho_V(1)(v) = v$ by the identity axiom. Rewriting with this, the equation becomes $[v \otimes \rho_W(h)(w)] = [\rho_V(h^{-1})(v) \otimes w]$, which gives the result by symmetry.
\end{proof}

\begin{theorem}[Inverse Balancing Relation]
\label{thm:balancedProduct.balanced'}
\lean{balancedProduct.balanced'}
\leanok
\uses{def:balancedProduct, thm:balancedProduct.balanced}
For all $h \in H$, $v \in V$, $w \in W$, the following equality holds in $V \otimes_H W$:
\[
  [\rho_V(h)(v) \otimes w] \;=\; [v \otimes \rho_W(h^{-1})(w)].
\]
\end{theorem}

\begin{proof}
\leanok
\uses{thm:balancedProduct.balanced}
Apply the balancing relation (Theorem~\ref{thm:balancedProduct.balanced}) with $h$ replaced by $h^{-1}$, obtaining $[\rho_V(h^{-1\,-1})(v) \otimes w] = [v \otimes \rho_W(h^{-1})(w)]$. Simplifying $h^{-1\,-1} = h$ by the involution property of group inversion gives the result.
\end{proof}

\begin{theorem}[Extensionality for Maps from the Balanced Product]
\label{thm:balancedProduct.ext'}
\lean{balancedProduct.ext'}
\leanok
\uses{def:balancedProduct}
Let $X$ be an $\mathbb{F}_2$-module and let $f, g : V \otimes_H W \to_{\mathbb{F}_2} X$ be $\mathbb{F}_2$-linear maps. If
\[
  f([\,v \otimes w\,]) \;=\; g([\,v \otimes w\,]) \quad \text{for all } v \in V,\, w \in W,
\]
then $f = g$.
\end{theorem}

\begin{proof}
\leanok
\uses{def:balancedProduct}
We apply the coinvariants extensionality principle $\operatorname{Coinvariants.hom_ext}$ for $\rho_V \otimes \rho_W$, which reduces equality of maps from the quotient to equality on the span of pure tensors. It then suffices to check equality on pure tensors $v \otimes w$, which is exactly the hypothesis. This follows by applying $\operatorname{TensorProduct.ext'}$.
\end{proof}

\begin{theorem}[Surjectivity of the Quotient Map]
\label{thm:balancedProduct.mk_surjective}
\lean{balancedProduct.mk_surjective}
\leanok
\uses{def:balancedProduct.mk}
The quotient map $\operatorname{mk} : V \otimes_{\mathbb{F}_2} W \twoheadrightarrow V \otimes_H W$ is surjective.
\end{theorem}

\begin{proof}
\leanok
\uses{def:balancedProduct.mk}
This follows directly from the general fact that coinvariant quotient maps are surjective: $\operatorname{Coinvariants.mk_surjective}(\rho_V \otimes \rho_W)$.
\end{proof}

\begin{definition}[Descended Averaging Map]
\label{def:avgDescended}
\lean{avgDescended}
\leanok
\uses{def:balancedProduct}
Assume $H$ is a finite group and $|H|$ is invertible in $\mathbb{F}_2$. The \emph{descended averaging map} is the $\mathbb{F}_2$-linear map
\[
  \overline{\operatorname{avg}} : V \otimes_H W \;\longrightarrow\; (V \otimes_{\mathbb{F}_2} W)^H,
\]
defined by lifting the averaging map $\operatorname{avg} : V \otimes_{\mathbb{F}_2} W \to (V \otimes_{\mathbb{F}_2} W)^H$,
\[
  \operatorname{avg}(x) \;=\; \frac{1}{|H|} \sum_{h \in H} (\rho_V \otimes \rho_W)(h)(x),
\]
through the coinvariants quotient. This is well-defined because the kernel of the coinvariants quotient is contained in the kernel of $\operatorname{avg}$: if $x = (\rho_V \otimes \rho_W)(g)(y) - y$ for some $g \in H$, then $\operatorname{avg}(x) = 0$ since the sum $\sum_{h \in H} (\rho_V \otimes \rho_W)(hg)(y) = \sum_{h \in H} (\rho_V \otimes \rho_W)(h)(y)$ by re-indexing.
\end{definition}

\begin{definition}[Inclusion of Invariants into the Balanced Product]
\label{def:invToBalanced}
\lean{invToBalanced}
\leanok
\uses{def:balancedProduct, def:avgDescended}
The \emph{inclusion map} is the $\mathbb{F}_2$-linear map
\[
  \iota : (V \otimes_{\mathbb{F}_2} W)^H \;\longrightarrow\; V \otimes_H W,
\]
defined as the composition of the submodule inclusion $(V \otimes_{\mathbb{F}_2} W)^H \hookrightarrow V \otimes_{\mathbb{F}_2} W$ followed by the quotient map $\operatorname{mk} : V \otimes_{\mathbb{F}_2} W \to V \otimes_H W$.
\end{definition}

\begin{definition}[Isomorphism of Balanced Product with Invariants]
\label{def:balancedProductInvariantsEquiv}
\lean{balancedProductInvariantsEquiv}
\leanok
\uses{def:balancedProduct, def:avgDescended, def:invToBalanced}
Assume $H$ is a finite group and $|H|$ is invertible in $\mathbb{F}_2$ (i.e., $|H|$ is odd). There is a canonical $\mathbb{F}_2$-linear equivalence
\[
  V \otimes_H W \;\cong_{\mathbb{F}_2}\; (V \otimes_{\mathbb{F}_2} W)^H.
\]
The forward map sends $[v \otimes w] \mapsto \frac{1}{|H|}\sum_{h \in H} \rho_V(h)(v) \otimes \rho_W(h)(w)$, and the inverse sends an invariant element $x \in (V \otimes W)^H$ to its class $[x]$ in $V \otimes_H W$.
\end{definition}

\begin{proof}
\leanok
\uses{def:avgDescended, def:invToBalanced}
We verify the two inverse conditions. \textbf{Left inverse:} Let $[x] \in V \otimes_H W$. By induction on coinvariants, it suffices to take $x \in V \otimes_{\mathbb{F}_2} W$. Unfolding definitions, $\iota(\overline{\operatorname{avg}}([x])) = [\operatorname{avg}(x)]$. We must show $[\operatorname{avg}(x)] = [x]$, i.e., $\operatorname{avg}(x) - x \in \ker(\operatorname{mk})$. Expanding the averaging map, $\operatorname{avg}(x) = \frac{1}{|H|}\sum_{h \in H} (\rho_V \otimes \rho_W)(h)(x)$. Using the fact that each $[(\rho_V \otimes \rho_W)(h)(x)] = [x]$ in coinvariants, we get $[\operatorname{avg}(x)] = \frac{1}{|H|} \cdot |H| \cdot [x] = [x]$ (using that $|H|$ is invertible). This uses $\operatorname{invOf_mul_self}$ and the coinvariant quotient property. \textbf{Right inverse:} Let $(x, hx) \in (V \otimes W)^H$ where $hx$ is the invariance proof. Then $\overline{\operatorname{avg}}(\iota(x)) = \overline{\operatorname{avg}}([x]) = \operatorname{avg}(x)$. Since $x$ is invariant, $(\rho_V \otimes \rho_W)(h)(x) = x$ for all $h$, so $\operatorname{avg}(x) = \frac{1}{|H|} \cdot |H| \cdot x = x$, which follows from $\operatorname{averageMap_id}$. The $\mathbb{F}_2$-linearity of the equivalence is inherited from $\overline{\operatorname{avg}}$.
\end{proof}

\begin{definition}[Isomorphism with Coinvariants]
\label{def:balancedProductCoinvariantsEquiv}
\lean{balancedProductCoinvariantsEquiv}
\leanok
\uses{def:balancedProduct}
There is a canonical $\mathbb{F}_2$-linear equivalence
\[
  \operatorname{Coinv}(\rho_V \otimes \mathbf{1}) \;\cong_{\mathbb{F}_2}\; \operatorname{Coinv}(\rho_V),
\]
where $\mathbf{1}$ denotes the trivial $H$-representation on $\mathbb{F}_2$. In other words, $V \otimes_H \mathbb{F}_2 \cong V_H$, identifying the balanced product of $V$ with the trivial module $\mathbb{F}_2$ with the coinvariant module of $\rho_V$. This isomorphism is induced by the canonical $\mathbb{F}_2$-linear isomorphism $\operatorname{rid} : V \otimes_{\mathbb{F}_2} \mathbb{F}_2 \xrightarrow{\sim} V$.
\end{definition}

\begin{proof}
\leanok
\uses{def:balancedProduct}
We show that $\operatorname{rid} : V \otimes_{\mathbb{F}_2} \mathbb{F}_2 \xrightarrow{\sim} V$ maps the coinvariant kernel of $\rho_V \otimes \mathbf{1}$ isomorphically onto the coinvariant kernel of $\rho_V$, establishing
\[
  \operatorname{(Coinv.ker}(\rho_V \otimes \mathbf{1})).\operatorname{map}(\operatorname{rid}) \;=\; \operatorname{Coinv.ker}(\rho_V).
\]
We prove both inclusions. \textbf{Forward inclusion ($\subseteq$):} By induction on the coinvariant kernel, it suffices to handle generators of the form $(\rho_V \otimes \mathbf{1})(g)(t) - t$ for $g \in H$, $t \in V \otimes \mathbb{F}_2$. By further induction on $t$ using $\operatorname{TensorProduct.induction_on}$: for zero tensors the result is trivial; for pure tensors $v \otimes a$, using the trivial action $\mathbf{1}(g) = \operatorname{id}$, we compute $\operatorname{rid}((\rho_V \otimes \mathbf{1})(g)(v \otimes a) - v \otimes a) = a \cdot (\rho_V(g)(v) - v)$, which lies in $\operatorname{Coinv.ker}(\rho_V)$ since $\rho_V(g)(v) - v$ is a generator and the submodule is closed under scalar multiplication; the addition case follows by linearity of $\operatorname{rid}$ and closure under addition. \textbf{Backward inclusion ($\supseteq$):} By induction on $\operatorname{Coinv.ker}(\rho_V)$, generators have the form $\rho_V(g)(v') - v'$. We exhibit a preimage: $\rho_V(g)(v') \otimes 1 - v' \otimes 1 \in \operatorname{Coinv.ker}(\rho_V \otimes \mathbf{1})$ (using the trivial action on $1 \in \mathbb{F}_2$), and $\operatorname{rid}(\rho_V(g)(v') \otimes 1 - v' \otimes 1) = \rho_V(g)(v') - v'$ as required; closures under addition and scalar multiplication are handled by decomposing preimages. The desired linear equivalence then follows from $\operatorname{Submodule.Quotient.equiv}$ applied to $\operatorname{rid}$ with this kernel compatibility.
\end{proof}

\begin{lemma}[Balanced Product is Nonempty]
\label{lem:balancedProduct_nonempty}
\lean{balancedProduct_nonempty}
\leanok
\uses{def:balancedProduct}
The balanced product $V \otimes_H W$ is nonempty. Specifically, $0 \in V \otimes_H W$.
\end{lemma}

\begin{proof}
\leanok
\uses{def:balancedProduct}
The zero element $0$ is an element of $V \otimes_H W$, witnessing nonemptiness.
\end{proof}

%--- Def_23: BalancedProductChainComplex ---
\chapter{Def 23: Balanced Product Chain Complex}

\begin{definition}[H-Equivariant Chain Complex]
\label{def:HEquivariantChainComplex}
\lean{HEquivariantChainComplex}
\leanok
\uses{def:balancedProduct}
Let $H$ be a group. An \emph{$H$-equivariant chain complex} over $\mathbb{F}_2$ consists of:
\begin{itemize}
  \item an underlying chain complex $C_\bullet$ over $\mathbb{F}_2$,
  \item for each degree $i \in \mathbb{Z}$, a representation $\rho_i : H \to \operatorname{GL}_{\mathbb{F}_2}(C_i)$ (i.e., a left $\mathbb{F}_2[H]$-module structure on $C_i$),
  \item an equivariance condition: for all $i, j \in \mathbb{Z}$ and $h \in H$,
        \[
          \partial^{i \to j} \circ \rho_i(h) = \rho_j(h) \circ \partial^{i \to j},
        \]
        i.e., the differential $\partial : C_i \to C_j$ intertwines the $H$-actions.
\end{itemize}
\end{definition}

\begin{definition}[Balanced Product Object]
\label{def:HEquivariantChainComplex.bpObj}
\lean{HEquivariantChainComplex.bpObj}
\leanok
\uses{def:HEquivariantChainComplex, def:balancedProduct}
Let $C$ and $D$ be $H$-equivariant chain complexes. For integers $p, q \in \mathbb{Z}$, the \emph{balanced product object} at $(p,q)$ is
\[
  (C \boxtimes_H D)_{p,q} \;=\; C_p \otimes_H D_q,
\]
the balanced product (coinvariants of $C_p \otimes_{\mathbb{F}_2} D_q$ under the diagonal $H$-action) of the representations $\rho^C_p$ and $\rho^D_q$.
\end{definition}

\begin{definition}[Horizontal Differential of Balanced Product]
\label{def:HEquivariantChainComplex.bpDh}
\lean{HEquivariantChainComplex.bpDh}
\leanok
\uses{def:HEquivariantChainComplex.bpObj, def:HEquivariantChainComplex, def:balancedProduct}
Let $C$ and $D$ be $H$-equivariant chain complexes. The \emph{horizontal differential}
\[
  \partial^h_{p,p',q} : C_p \otimes_H D_q \longrightarrow C_{p'} \otimes_H D_q
\]
is induced by $\partial^C_{p \to p'} \otimes \operatorname{id}_{D_q}$ on the balanced product. This is well-defined on coinvariants by the $H$-equivariance of $\partial^C$.
\end{definition}

\begin{definition}[Vertical Differential of Balanced Product]
\label{def:HEquivariantChainComplex.bpDv}
\lean{HEquivariantChainComplex.bpDv}
\leanok
\uses{def:HEquivariantChainComplex.bpObj, def:HEquivariantChainComplex, def:balancedProduct}
Let $C$ and $D$ be $H$-equivariant chain complexes. The \emph{vertical differential}
\[
  \partial^v_{p,q,q'} : C_p \otimes_H D_q \longrightarrow C_p \otimes_H D_{q'}
\]
is induced by $\operatorname{id}_{C_p} \otimes \partial^D_{q \to q'}$ on the balanced product. This is well-defined on coinvariants by the $H$-equivariance of $\partial^D$.
\end{definition}

\begin{definition}[Balanced Product Double Complex]
\label{def:HEquivariantChainComplex.balancedProductDoubleComplex}
\lean{HEquivariantChainComplex.balancedProductDoubleComplex}
\leanok
\uses{def:HEquivariantChainComplex.bpObj, def:HEquivariantChainComplex.bpDh, def:HEquivariantChainComplex.bpDv, def:HEquivariantChainComplex, def:balancedProduct}
Let $C$ and $D$ be $H$-equivariant chain complexes. The \emph{balanced product double complex} $C \boxtimes_H D$ is the double complex over $\mathbb{F}_2$ with:
\begin{itemize}
  \item objects $(C \boxtimes_H D)_{p,q} = C_p \otimes_H D_q$,
  \item horizontal differential $\partial^h = \partial^C \otimes_H \operatorname{id}_D$ (decreasing $p$),
  \item vertical differential $\partial^v = \operatorname{id}_C \otimes_H \partial^D$ (decreasing $q$).
\end{itemize}
This forms a valid double complex: $(\partial^h)^2 = 0$, $(\partial^v)^2 = 0$, and $\partial^h \circ \partial^v = \partial^v \circ \partial^h$.
\end{definition}

\begin{definition}[Balanced Product Complex]
\label{def:HEquivariantChainComplex.balancedProductComplex}
\lean{HEquivariantChainComplex.balancedProductComplex}
\leanok
\uses{def:HEquivariantChainComplex.balancedProductDoubleComplex, def:HEquivariantChainComplex, def:balancedProduct}
Let $C$ and $D$ be $H$-equivariant chain complexes. The \emph{balanced product complex}
\[
  C \otimes_H D \;=\; \operatorname{Tot}(C \boxtimes_H D)
\]
is the total complex of the balanced product double complex $C \boxtimes_H D$.
\end{definition}

\begin{theorem}[Balanced Product Double Complex Object]
\label{thm:HEquivariantChainComplex.balancedProductDoubleComplex_obj}
\lean{HEquivariantChainComplex.balancedProductDoubleComplex_obj}
\leanok
\uses{def:HEquivariantChainComplex.balancedProductDoubleComplex, def:HEquivariantChainComplex.bpObj}
For all $p, q \in \mathbb{Z}$, the object of the balanced product double complex at $(p,q)$ satisfies
\[
  (C \boxtimes_H D)_{p,q} = C_p \otimes_H D_q.
\]
\end{theorem}

\begin{proof}
\leanok
\uses{def:HEquivariantChainComplex.balancedProductDoubleComplex, def:HEquivariantChainComplex.bpObj}
This holds by reflexivity, as the definition of $\operatorname{balancedProductDoubleComplex}$ sets the object at $(p,q)$ to be $\operatorname{bpObj}(p,q) = C_p \otimes_H D_q$ definitionally.
\end{proof}

\begin{theorem}[Vertical Differential Formula]
\label{thm:HEquivariantChainComplex.balancedProductDoubleComplex_dv_eq}
\lean{HEquivariantChainComplex.balancedProductDoubleComplex_dv_eq}
\leanok
\uses{def:HEquivariantChainComplex.balancedProductDoubleComplex, def:HEquivariantChainComplex.bpDv}
For all $p, q \in \mathbb{Z}$, the vertical differential of the balanced product double complex satisfies
\[
  \partial^v_{p,q} \;=\; \partial^v_{p,q,q-1} : C_p \otimes_H D_q \to C_p \otimes_H D_{q-1},
\]
i.e., it is induced by $\operatorname{id}_{C_p} \otimes \partial^D_{q \to q-1}$.
\end{theorem}

\begin{proof}
\leanok
\uses{def:HEquivariantChainComplex.balancedProductDoubleComplex, def:HEquivariantChainComplex.bpDv}
This holds by reflexivity, as $\operatorname{dv}$ at $(p,q)$ in the balanced product double complex is defined as the linear map underlying $\operatorname{bpDv}(p, q, q-1)$.
\end{proof}

\begin{theorem}[Horizontal Differential Formula]
\label{thm:HEquivariantChainComplex.balancedProductDoubleComplex_dh_eq}
\lean{HEquivariantChainComplex.balancedProductDoubleComplex_dh_eq}
\leanok
\uses{def:HEquivariantChainComplex.balancedProductDoubleComplex, def:HEquivariantChainComplex.bpDh}
For all $p, q \in \mathbb{Z}$, the horizontal differential of the balanced product double complex satisfies
\[
  \partial^h_{p,q} \;=\; \partial^h_{p,p-1,q} : C_p \otimes_H D_q \to C_{p-1} \otimes_H D_q,
\]
i.e., it is induced by $\partial^C_{p \to p-1} \otimes \operatorname{id}_{D_q}$.
\end{theorem}

\begin{proof}
\leanok
\uses{def:HEquivariantChainComplex.balancedProductDoubleComplex, def:HEquivariantChainComplex.bpDh}
This holds by reflexivity, as $\operatorname{dh}$ at $(p,q)$ in the balanced product double complex is defined as the linear map underlying $\operatorname{bpDh}(p, p-1, q)$.
\end{proof}

\begin{lemma}[Balanced Product Complex is Nonempty]
\label{lem:HEquivariantChainComplex.balancedProductComplex_nonempty}
\lean{HEquivariantChainComplex.balancedProductComplex_nonempty}
\leanok
\uses{def:HEquivariantChainComplex.balancedProductComplex}
For every $n \in \mathbb{Z}$, the degree-$n$ component of the balanced product complex $C \otimes_H D$ is nonempty (as it contains the zero element).
\end{lemma}

\begin{proof}
\leanok
\uses{def:HEquivariantChainComplex.balancedProductComplex}
We exhibit the zero element $0 \in (C \otimes_H D)_n$, which witnesses that the type is nonempty.
\end{proof}

%--- Lem_4: KunnethBalancedProduct ---
\chapter{Lem 4: Künneth Formula for Balanced Products}

\begin{definition}[Cycles Action]
\label{def:HEquivariantChainComplex.cyclesAction}
\lean{HEquivariantChainComplex.cyclesAction}
\leanok
\uses{def:HEquivariantChainComplex}
Let $H$ be a finite group and $C$ an $H$-equivariant chain complex over $\mathbb{F}_2$. For each $p \in \mathbb{Z}$, the group $H$ acts on the cycles $Z_p(C) = \ker \partial_p$ by the \emph{cycles action}
\[\rho^Z_p : H \to \operatorname{End}_{\mathbb{F}_2}(Z_p(C)),\quad h \cdot [z] = [\rho_p(h)(z)],\]
where $\rho_p$ is the given $H$-action on $C_p$. This is well-defined: if $z \in \ker \partial_p$ then $\partial_p(\rho_p(h)(z)) = \rho_{p-1}(h)(\partial_p(z)) = \rho_{p-1}(h)(0) = 0$, using the $H$-equivariance of $\partial_p$.
\end{definition}

\begin{proof}
\leanok
\uses{def:HEquivariantChainComplex}
Let $z \in Z_p(C)$, so $\partial_p(z) = 0$. We must show $\rho_p(h)(z) \in Z_p(C)$. By the equivariance hypothesis $\partial_p \circ \rho_p(h) = \rho_{p-1}(h) \circ \partial_p$, evaluating at $z$ gives $\partial_p(\rho_p(h)(z)) = \rho_{p-1}(h)(\partial_p(z)) = \rho_{p-1}(h)(0) = 0$. The map $h \mapsto \rho^Z_p(h)$ is a group homomorphism: $\rho^Z_p(1) = \mathrm{id}$ follows from $\rho_p(1) = \mathrm{id}$, and $\rho^Z_p(g_1 g_2) = \rho^Z_p(g_1) \circ \rho^Z_p(g_2)$ follows from $\rho_p(g_1 g_2) = \rho_p(g_1) \circ \rho_p(g_2)$.
\end{proof}

\begin{definition}[Boundaries Action]
\label{def:HEquivariantChainComplex.boundariesAction}
\lean{HEquivariantChainComplex.boundariesAction}
\leanok
\uses{def:HEquivariantChainComplex}
Let $H$ be a finite group and $C$ an $H$-equivariant chain complex over $\mathbb{F}_2$. For each $p \in \mathbb{Z}$, the group $H$ acts on the boundaries $B_p(C) = \operatorname{im} \partial_{p+1}$ by the \emph{boundaries action}
\[\rho^B_p : H \to \operatorname{End}_{\mathbb{F}_2}(B_p(C)),\quad h \cdot [\partial_{p+1}(x)] = [\partial_{p+1}(\rho_{p+1}(h)(x))],\]
using $H$-equivariance of $\partial_{p+1}$.
\end{definition}

\begin{proof}
\leanok
\uses{def:HEquivariantChainComplex}
Let $b \in B_p(C)$, so there exists $x \in C_{p+1}$ with $b = \partial_{p+1}(x)$. We must show $\rho_p(h)(b) \in B_p(C)$. We claim $\rho_p(h)(b) = \partial_{p+1}(\rho_{p+1}(h)(x))$. By the equivariance condition $\partial_{p+1} \circ \rho_{p+1}(h) = \rho_p(h) \circ \partial_{p+1}$, evaluating at $x$ gives $\partial_{p+1}(\rho_{p+1}(h)(x)) = \rho_p(h)(\partial_{p+1}(x)) = \rho_p(h)(b)$. The homomorphism property follows from that of $\rho_p$.
\end{proof}

\begin{definition}[Boundaries-in-Cycles Action]
\label{def:HEquivariantChainComplex.boundariesSubCyclesAction}
\lean{HEquivariantChainComplex.boundariesSubCyclesAction}
\leanok
\uses{def:HEquivariantChainComplex}
Let $H$ be a finite group and $C$ an $H$-equivariant chain complex over $\mathbb{F}_2$. For each $p \in \mathbb{Z}$, the group $H$ acts on the submodule $B_p(C) \subseteq Z_p(C)$ (boundaries viewed inside cycles) by the \emph{boundaries-sub-cycles action}, induced from $\rho^Z_p$ and $\rho^B_p$ compatibly: for $z \in Z_p(C)$ with $z \in B_p(C)$, one sets $h \cdot z = \rho_p(h)(z)$, which lies in both $Z_p(C)$ and $B_p(C)$.
\end{definition}

\begin{proof}
\leanok
\uses{def:HEquivariantChainComplex}
Let $z$ be an element of $C.complex.boundariesSubCycles\ p$, meaning $z$ is a cycle and a boundary. The image $\rho_p(h)(z)$ is a cycle by the cycles action, and it is a boundary by the boundaries action. The $\mathbb{F}_2$-linearity and group homomorphism property follow from those of $\rho_p$.
\end{proof}

\begin{definition}[Homology Action]
\label{def:HEquivariantChainComplex.homologyAction}
\lean{HEquivariantChainComplex.homologyAction}
\leanok
\uses{def:HEquivariantChainComplex}
Let $H$ be a finite group and $C$ an $H$-equivariant chain complex over $\mathbb{F}_2$. For each $p \in \mathbb{Z}$, the \emph{induced $H$-action on homology} $H_p(C) = Z_p(C)/B_p(C)$ is the representation
\[\rho^{H_p} : H \to \operatorname{End}_{\mathbb{F}_2}(H_p(C)),\quad h \cdot [z] = [\rho_p(h)(z)],\]
defined via the quotient map $\operatorname{mapQ}$ applied to the cycles action $\rho^Z_p(h)$, using the fact that $\rho_p(h)$ preserves $B_p(C)$.
\end{definition}

\begin{proof}
\leanok
\uses{def:HEquivariantChainComplex}
We define the action using $\operatorname{mapQ}$: for each $h \in H$, the linear endomorphism $\rho^Z_p(h)$ of $Z_p(C)$ sends $B_p(C)$ into itself (since $\rho^B_p(h)$ is well-defined), so it descends to a well-defined linear map on the quotient $H_p(C) = Z_p(C)/B_p(C)$. For the homomorphism property, $\rho^{H_p}(1) = \mathrm{id}$: on a quotient representative $[z]$, $[\rho_p(1)(z)] = [z]$ since $\rho_p(1) = \mathrm{id}$. For $g_1, g_2 \in H$, $[\rho_p(g_1 g_2)(z)] = [\rho_p(g_1)(\rho_p(g_2)(z))]$ by $\rho_p(g_1 g_2) = \rho_p(g_1) \circ \rho_p(g_2)$, verified on each quotient element by surjectivity.
\end{proof}

\begin{definition}[Balanced Künneth Summand]
\label{def:HEquivariantChainComplex.balancedKunnethSummand}
\lean{HEquivariantChainComplex.balancedKunnethSummand}
\leanok
\uses{def:HEquivariantChainComplex}
Let $H$ be a finite group and $C$, $D$ $H$-equivariant chain complexes over $\mathbb{F}_2$. For $n, p \in \mathbb{Z}$, the \emph{balanced Künneth summand} at index $p$ is
\[K_{n,p}(C,D) := H_p(C) \otimes_H H_{n-p}(D),\]
defined as the coinvariants of the tensor product representation $\rho^{H_p}_C \otimes \rho^{H_{n-p}}_D$ of $H$ on $H_p(C) \otimes_{\mathbb{F}_2} H_{n-p}(D)$.
\end{definition}

\begin{definition}[Balanced Künneth Sum]
\label{def:HEquivariantChainComplex.balancedKunnethSum}
\lean{HEquivariantChainComplex.balancedKunnethSum}
\leanok
\uses{def:HEquivariantChainComplex.balancedKunnethSummand}
Let $H$ be a finite group and $C$, $D$ $H$-equivariant chain complexes over $\mathbb{F}_2$. For $n \in \mathbb{Z}$, the \emph{balanced Künneth sum} is
\[K_n(C,D) := \bigoplus_{p \in \mathbb{Z}} H_p(C) \otimes_H H_{n-p}(D),\]
the direct sum of all balanced Künneth summands.
\end{definition}

\begin{definition}[Invertible Cardinality in Odd Order]
\label{def:HEquivariantChainComplex.invertible_card_of_odd_order}
\lean{HEquivariantChainComplex.invertible_card_of_odd_order}
\leanok
\uses{def:HEquivariantChainComplex}
Let $H$ be a finite group of odd order. Then the cardinality $|H|$ is invertible in $\mathbb{F}_2$, i.e.\ $|H| \equiv 1 \pmod{2}$ means $(|H| : \mathbb{F}_2) = 1$.
\end{definition}

\begin{proof}
\leanok
\uses{def:HEquivariantChainComplex}
Since $|H|$ is odd, write $|H| = 2k + 1$ for some $k \in \mathbb{N}$. Then in $\mathbb{F}_2$:
$(|H| : \mathbb{F}_2) = (2k + 1 : \mathbb{F}_2) = 2k + 1 = 0 \cdot k + 1 = 1,$
using $(2 : \mathbb{F}_2) = 0$ (verified by computation). Hence $(|H| : \mathbb{F}_2) = 1$, which is invertible (with $\operatorname{Invertible} 1$ given by the standard instance for units).
\end{proof}

\begin{theorem}[Axiom: Künneth Formula for Balanced Products]
\label{thm:HEquivariantChainComplex.kunnethBalancedProduct}
\lean{HEquivariantChainComplex.kunnethBalancedProduct}
\leanok
\uses{def:HEquivariantChainComplex, def:HEquivariantChainComplex.balancedKunnethSum, def:HEquivariantChainComplex.homologyAction}

\textbf{\color{red}[UNPROVEN AXIOM]}

Let $H$ be a finite group of odd order, and let $C$, $D$ be $H$-equivariant chain complexes over $\mathbb{F}_2$. Then for every $n \in \mathbb{Z}$ there is a linear isomorphism
\[H_n(C \otimes_H D) \;\cong_{\mathbb{F}_2}\; \bigoplus_{p \in \mathbb{Z}} H_p(C) \otimes_H H_{n-p}(D).\]

\textit{Note: This is stated as an axiom because the full proof requires infrastructure not available in Mathlib: specifically, (1) identifying the balanced product complex $C \otimes_H D$ (from $\operatorname{balancedProductComplex}$, Def.\ 23) with the coinvariants of the tensor product complex $(C \otimes D)_H$ at the chain complex level; (2) showing that over $\mathbb{F}_2$ with $|H|$ odd the coinvariants functor $(-)_H$ is exact (so it commutes with homology); and (3) combining with the ordinary Künneth formula (Thm.\ 1) and the fact that $(-)_H$ commutes with finite direct sums. The proof sketch is: $H_n(C \otimes_H D) \cong H_n((C \otimes D)_H) \cong (H_n(C \otimes D))_H \cong \bigoplus_{p+q=n}(H_p(C) \otimes H_q(D))_H = \bigoplus_{p+q=n} H_p(C) \otimes_H H_q(D)$.}
\end{theorem}

\begin{lemma}[Satisfiability Witness for Künneth Balanced Product]
\label{lem:HEquivariantChainComplex.kunnethBalancedProduct_satisfiable}
\lean{HEquivariantChainComplex.kunnethBalancedProduct_satisfiable}
\leanok
\uses{def:HEquivariantChainComplex, thm:HEquivariantChainComplex.kunnethBalancedProduct}
The premises of \texttt{kunnethBalancedProduct} are satisfiable: there exist a finite group $H$, $H$-equivariant chain complexes $C$ and $D$ over $\mathbb{F}_2$, and a proof that $|H|$ is odd.
\end{lemma}

\begin{proof}
\leanok
\uses{def:HEquivariantChainComplex}
Take $H = \mathbf{1}$ (the trivial group). Then $|H| = 1$, which is odd. Let $C = D$ be the trivial $H$-equivariant chain complex given by the single-object chain complex concentrated at degree $0$ with value $\mathbb{F}_2$, equipped with the trivial $H$-action (the only possible action since $H$ is trivial). This is equivariant vacuously. Thus the witness is $(H, C, D)$ with $\operatorname{Odd}(|\mathbf{1}|)$ witnessed by $1 = 2 \cdot 0 + 1$.
\end{proof}

\begin{lemma}[Balanced Künneth Summand is Nonempty]
\label{lem:HEquivariantChainComplex.balancedKunnethSummand_nonempty}
\lean{HEquivariantChainComplex.balancedKunnethSummand_nonempty}
\leanok
\uses{def:HEquivariantChainComplex.balancedKunnethSummand}
For any $n, p \in \mathbb{Z}$, the balanced Künneth summand $H_p(C) \otimes_H H_{n-p}(D)$ is nonempty.
\end{lemma}

\begin{proof}
\leanok
\uses{def:HEquivariantChainComplex.balancedKunnethSummand}
The zero element $0$ belongs to every module, so $H_p(C) \otimes_H H_{n-p}(D)$ contains at least $0$.
\end{proof}

\begin{lemma}[Balanced Künneth Sum is Nonempty]
\label{lem:HEquivariantChainComplex.balancedKunnethSum_nonempty}
\lean{HEquivariantChainComplex.balancedKunnethSum_nonempty}
\leanok
\uses{def:HEquivariantChainComplex.balancedKunnethSum}
For any $n \in \mathbb{Z}$, the balanced Künneth sum $\bigoplus_{p \in \mathbb{Z}} H_p(C) \otimes_H H_{n-p}(D)$ is nonempty.
\end{lemma}

\begin{proof}
\leanok
\uses{def:HEquivariantChainComplex.balancedKunnethSum}
The zero element $0$ belongs to every direct sum of modules, so the balanced Künneth sum contains at least $0$.
\end{proof}

%--- Def_24: QuotientGraphTrivialization ---
\chapter{Def 24: Quotient Graph Trivialization}

\begin{definition}[Graph with Group Action]
\label{def:GraphWithGroupAction}
\lean{GraphWithGroupAction}
\leanok
\uses{def:CellComplex}
A \emph{graph with group action} consists of a finite graph $X$ (given as a 1-dimensional cell complex, i.e., a \texttt{CellComplex}) equipped with a free right action of a finite group $H$ on both vertices $X_0$ and edges $X_1$, compatible with the boundary map, and satisfying a quotient condition.

Formally, it is a structure $\mathtt{GraphWithGroupAction}(H)$ for a group $H$ with $|H| < \infty$, comprising:
\begin{itemize}
  \item A cell complex $\mathtt{graph} : \mathtt{CellComplex}$ (the underlying graph).
  \item A right $H$-action $\mathtt{vertexAction}$ on the vertices $X.\mathtt{graph}.\mathtt{cells}(0)$, modeled as a left \texttt{MulAction} via $h \cdot v = v \cdot h^{-1}$.
  \item A right $H$-action $\mathtt{edgeAction}$ on the edges $X.\mathtt{graph}.\mathtt{cells}(1)$.
  \item Freeness of the vertex action: for all $h \in H$ and $v \in X_0$, if $h \cdot v = v$ then $h = 1$.
  \item Freeness of the edge action: for all $h \in H$ and $e \in X_1$, if $h \cdot e = e$ then $h = 1$.
  \item Boundary equivariance: for all $e \in X_1$, $h \in H$, and $\tau \in X_0$,
    \[
      \tau \in \partial(h \cdot e) \iff h^{-1} \cdot \tau \in \partial e.
    \]
  \item Quotient condition: for all $e \in X_1$, $v \in X_0$, and $h \in H$, if $v \in \partial e$ and $h \cdot v \in \partial e$, then $h = 1$. (No edge connects a vertex to a nontrivially translated copy of itself.)
\end{itemize}
\end{definition}

\begin{definition}[Quotient Vertices]
\label{def:GraphWithGroupAction.quotientVertices}
\lean{GraphWithGroupAction.quotientVertices}
\leanok
\uses{def:CellComplex}
Let $X : \mathtt{GraphWithGroupAction}(H)$. The \emph{quotient vertices} of $X$ by $H$ are the orbits of the $H$-action on $X_0$:
\[
  (X/H)_0 \;:=\; X_0 \,/\, H \;=\; \mathtt{MulAction.orbitRel.Quotient}\, H\, (X.\mathtt{graph}.\mathtt{cells}\, 0).
\]
\end{definition}

\begin{definition}[Quotient Edges]
\label{def:GraphWithGroupAction.quotientEdges}
\lean{GraphWithGroupAction.quotientEdges}
\leanok
\uses{def:CellComplex}
Let $X : \mathtt{GraphWithGroupAction}(H)$. The \emph{quotient edges} of $X$ by $H$ are the orbits of the $H$-action on $X_1$:
\[
  (X/H)_1 \;:=\; X_1 \,/\, H \;=\; \mathtt{MulAction.orbitRel.Quotient}\, H\, (X.\mathtt{graph}.\mathtt{cells}\, 1).
\]
\end{definition}

\begin{definition}[Vertex Quotient Map]
\label{def:GraphWithGroupAction.vertexQuot}
\lean{GraphWithGroupAction.vertexQuot}
\leanok
\uses{def:CellComplex}
The \emph{vertex quotient map} $\pi_0 : X_0 \to (X/H)_0$ sends a vertex $v$ to its $H$-orbit:
\[
  \pi_0(v) \;:=\; [v] \;\in\; X_0/H.
\]
\end{definition}

\begin{definition}[Edge Quotient Map]
\label{def:GraphWithGroupAction.edgeQuot}
\lean{GraphWithGroupAction.edgeQuot}
\leanok
\uses{def:CellComplex}
The \emph{edge quotient map} $\pi_1 : X_1 \to (X/H)_1$ sends an edge $e$ to its $H$-orbit:
\[
  \pi_1(e) \;:=\; [e] \;\in\; X_1/H.
\]
\end{definition}

\begin{lemma}[Boundary Orbit Compatibility]
\label{lem:GraphWithGroupAction.bdry_orbit_compat}
\lean{GraphWithGroupAction.bdry_orbit_compat}
\leanok
\uses{def:CellComplex}
Let $X : \mathtt{GraphWithGroupAction}(H)$, $e \in X_1$, $h \in H$, and $\tau \in X_0$. If $\tau \in \partial e$, then $h \cdot \tau \in \partial(h \cdot e)$.
\end{lemma}

\begin{proof}
\leanok
\uses{def:CellComplex}
We must show $h \cdot \tau \in \partial(h \cdot e)$. By boundary equivariance, $h \cdot \tau \in \partial(h \cdot e)$ if and only if $h^{-1} \cdot (h \cdot \tau) \in \partial e$. Since $h^{-1} \cdot (h \cdot \tau) = \tau$ and $\tau \in \partial e$ by hypothesis, the result follows by simplification.
\end{proof}

\begin{definition}[Quotient Boundary Map]
\label{def:GraphWithGroupAction.quotientBdry}
\lean{GraphWithGroupAction.quotientBdry}
\leanok
\uses{def:CellComplex}
The \emph{quotient boundary map} on the quotient graph assigns to each edge orbit $E \in (X/H)_1$ a finite set of vertex orbits:
\[
  \partial_{X/H}(E) \;:=\; \bigl\{\, \pi_0(\tau) \;\mid\; \tau \in \partial e \,\bigr\}
\]
where $e$ is any representative of $E$. This is well-defined: if $e_1, e_2$ represent the same orbit, so $\exists\, g \in H,\; g \cdot e_2 = e_1$, then for any $\tau \in \partial e_1$, we have $g^{-1} \cdot \tau \in \partial e_2$ by boundary equivariance, and $\pi_0(g^{-1} \cdot \tau) = \pi_0(\tau)$ since they are in the same $H$-orbit; and vice versa.
\end{definition}

\begin{lemma}[Quotient Boundary is Nonempty]
\label{lem:GraphWithGroupAction.quotientBdry_nonempty}
\lean{GraphWithGroupAction.quotientBdry_nonempty}
\leanok
\uses{def:CellComplex}
Let $X : \mathtt{GraphWithGroupAction}(H)$, and suppose every edge has at least one boundary vertex: $\forall\, e \in X_1,\; \partial e \neq \emptyset$. Then for every edge orbit $E \in (X/H)_1$, the quotient boundary $\partial_{X/H}(E)$ is nonempty.
\end{lemma}

\begin{proof}
\leanok
\uses{def:CellComplex}
By induction on $E$ using a representative $e$. By the unfolding of $\mathtt{quotientBdry}$, we have $\partial_{X/H}([e]) = \{\pi_0(\tau) \mid \tau \in \partial e\}$. Since $\partial e$ is nonempty by hypothesis, choose $\tau \in \partial e$. Then $\pi_0(\tau) \in \partial_{X/H}([e])$, so the set is nonempty.
\end{proof}

\begin{definition}[Trivialization]
\label{def:GraphWithGroupAction.Trivialization}
\lean{GraphWithGroupAction.Trivialization}
\leanok
\uses{def:CellComplex}
A \emph{trivialization} of the vertex quotient map $\pi_0 : X_0 \to (X/H)_0$ is a section, i.e., a function
\[
  R : (X/H)_0 \;\longrightarrow\; X_0
\]
such that for every vertex orbit $V \in (X/H)_0$, the representative $R(V)$ lies in $V$:
\[
  \pi_0(R(V)) = V \quad \forall\, V \in (X/H)_0.
\]
The set $\{R(V) \mid V \in (X/H)_0\} \subseteq X_0$ contains exactly one representative from each $H$-orbit of vertices.
\end{definition}

\begin{theorem}[Trivialization Exists]
\label{thm:GraphWithGroupAction.trivialization_exists}
\lean{GraphWithGroupAction.trivialization_exists}
\leanok
\uses{def:CellComplex}
For any $X : \mathtt{GraphWithGroupAction}(H)$, a trivialization exists: $\mathtt{Nonempty}(X.\mathtt{Trivialization})$.
\end{theorem}

\begin{proof}
\leanok
\uses{def:CellComplex}
We first observe that for every vertex orbit $V \in (X/H)_0$, there exists a vertex $v \in X_0$ with $\pi_0(v) = V$: this follows by induction on $V$ using a representative. Specifically, if $V = [v]$ for some $v \in X_0$, then $\pi_0(v) = [v] = V$ by reflexivity.

Having shown existence for each orbit, we apply the axiom of choice to obtain a selection function $V \mapsto \varepsilon(V)$ where $\varepsilon(V)$ is a chosen representative with $\pi_0(\varepsilon(V)) = V$. This selection defines a trivialization $R = \langle \varepsilon, \text{(the proof of representativeness)} \rangle$, so $X.\mathtt{Trivialization}$ is nonempty.
\end{proof}

\begin{lemma}[Trivialization is Nonempty]
\label{lem:GraphWithGroupAction.trivialization_nonempty}
\lean{GraphWithGroupAction.trivialization_nonempty}
\leanok
\uses{def:CellComplex}
For any $X : \mathtt{GraphWithGroupAction}(H)$, $\mathtt{Nonempty}(X.\mathtt{Trivialization})$. This follows immediately from \texttt{trivialization\_exists}.
\end{lemma}

\begin{proof}
\leanok
\uses{def:CellComplex}
This is an immediate consequence of $X.\mathtt{trivialization\_exists}$.
\end{proof}

\begin{definition}[Vertex Group Element]
\label{def:GraphWithGroupAction.vertexGroupElement}
\lean{GraphWithGroupAction.vertexGroupElement}
\leanok
\uses{def:CellComplex}
Given a trivialization $R$ of $X$ and a vertex $v \in X_0$, the \emph{vertex group element} $g_v \in H$ is the unique group element such that
\[
  g_v \cdot R(\pi_0(v)) = v.
\]
Since $\pi_0(R(\pi_0(v))) = \pi_0(v)$, the vertices $R(\pi_0(v))$ and $v$ lie in the same $H$-orbit, so such a $g_v$ exists (chosen via the axiom of choice).
\end{definition}

\begin{lemma}[Vertex Group Element Specification]
\label{lem:GraphWithGroupAction.vertexGroupElement_spec}
\lean{GraphWithGroupAction.vertexGroupElement_spec}
\leanok
\uses{def:CellComplex}
For any trivialization $R$ and vertex $v \in X_0$:
\[
  g_v \cdot R(\pi_0(v)) = v,
\]
where $g_v = \mathtt{vertexGroupElement}(R, v)$.
\end{lemma}

\begin{proof}
\leanok
\uses{def:CellComplex}
This follows directly from the definition of $\mathtt{vertexGroupElement}$ as the chosen witness of the existential statement, whose proof gives $g_v \cdot R(\pi_0(v)) = v$.
\end{proof}

\begin{definition}[Connection on Edge]
\label{def:GraphWithGroupAction.connectionOnEdge}
\lean{GraphWithGroupAction.connectionOnEdge}
\leanok
\uses{def:CellComplex}
Given a trivialization $R$, an edge $e \in X_1$, and two boundary vertices $\tau_1, \tau_2 \in \partial e$ lying in distinct vertex orbits ($\pi_0(\tau_1) \neq \pi_0(\tau_2)$), the \emph{connection} of $e$ relative to $\tau_1$ and $\tau_2$ is the group element
\[
  \phi_R(e;\,\tau_1,\tau_2) \;:=\; g_{\tau_1}^{-1} \cdot g_{\tau_2} \;\in\; H,
\]
where $g_{\tau_i} = \mathtt{vertexGroupElement}(R, \tau_i)$ is the unique element satisfying $g_{\tau_i} \cdot R(\pi_0(\tau_i)) = \tau_i$.
\end{definition}

\begin{lemma}[Connection is Constant on Orbits]
\label{lem:GraphWithGroupAction.connectionOnEdge_orbit_compat}
\lean{GraphWithGroupAction.connectionOnEdge_orbit_compat}
\leanok
\uses{def:CellComplex}
Let $R$ be a trivialization, $e \in X_1$, $g \in H$, and $\tau_1, \tau_2 \in \partial e$ with $\pi_0(\tau_1) \neq \pi_0(\tau_2)$. Then the connection is invariant under the $H$-action:
\[
  \phi_R(g \cdot e;\; g \cdot \tau_1,\; g \cdot \tau_2) \;=\; \phi_R(e;\; \tau_1,\; \tau_2).
\]
\end{lemma}

\begin{proof}
\leanok
\uses{def:CellComplex}
We unfold $\mathtt{connectionOnEdge}$ on both sides. It suffices to show that for any $v \in X_0$,
\[
  \mathtt{vertexGroupElement}(R, g \cdot v) \;=\; g \cdot \mathtt{vertexGroupElement}(R, v).
\]
Let $h = \mathtt{vertexGroupElement}(R, v)$, so $h \cdot R(\pi_0(v)) = v$.

First, note $\pi_0(g \cdot v) = \pi_0(v)$, since $g \cdot v$ and $v$ are in the same $H$-orbit (with $\pi_0$ sending each element to its orbit, which is a sound quotient relation via $\exists g, g \cdot v' = v$). Hence $R(\pi_0(g \cdot v)) = R(\pi_0(v))$.

Let $h' = \mathtt{vertexGroupElement}(R, g \cdot v)$, so $h' \cdot R(\pi_0(v)) = g \cdot v$. Also $(g \cdot h) \cdot R(\pi_0(v)) = g \cdot (h \cdot R(\pi_0(v))) = g \cdot v$. Thus
\[
  h' \cdot R(\pi_0(v)) = (g \cdot h) \cdot R(\pi_0(v)).
\]
By freeness of the vertex action applied to $(h')^{-1} \cdot (g \cdot h) \in H$ acting on $R(\pi_0(v))$, we get $(h')^{-1} \cdot (g \cdot h) = 1$, i.e., $h' = g \cdot h$.

Substituting back, $\phi_R(g \cdot e; g \cdot \tau_1, g \cdot \tau_2) = (g \cdot g_{\tau_1})^{-1} \cdot (g \cdot g_{\tau_2}) = g_{\tau_1}^{-1} \cdot g^{-1} \cdot g \cdot g_{\tau_2} = g_{\tau_1}^{-1} \cdot g_{\tau_2} = \phi_R(e; \tau_1, \tau_2)$, using commutativity of the inverse and simplification by linear arithmetic.
\end{proof}

\begin{definition}[Induced Connection]
\label{def:GraphWithGroupAction.inducedConnection}
\lean{GraphWithGroupAction.inducedConnection}
\leanok
\uses{def:CellComplex}
Given a trivialization $R$ of $X$, the \emph{induced connection} is a map
\[
  \phi_R : (X/H)_1 \;\longrightarrow\; H
\]
defined as follows. For an edge orbit $E \in (X/H)_1$, let $e = \mathtt{out}(E)$ be a chosen representative edge with boundary $\partial e$. If there exist $\tau_1, \tau_2 \in \partial e$ with $\pi_0(\tau_1) \neq \pi_0(\tau_2)$, we set
\[
  \phi_R(E) \;:=\; g_{\tau_1}^{-1} \cdot g_{\tau_2},
\]
where $\tau_1, \tau_2$ are chosen witnesses and $g_{\tau_i} = \mathtt{vertexGroupElement}(R,\tau_i)$. Otherwise (all boundary vertices are in the same orbit, or the boundary is small), we set $\phi_R(E) := 1$.

By the orbit compatibility lemma, this is independent of the choice of representative edge.
\end{definition}

\begin{theorem}[Induced Connection Specification]
\label{thm:GraphWithGroupAction.inducedConnection_spec}
\lean{GraphWithGroupAction.inducedConnection_spec}
\leanok
\uses{def:CellComplex}
Let $R$ be a trivialization of $X$, let $e \in X_1$, and let $\tau_1, \tau_2 \in \partial e$ with $\pi_0(\tau_1) \neq \pi_0(\tau_2)$. Then there exists an edge $e' \in X_1$ such that:
\begin{enumerate}
  \item $\pi_1(e') = \pi_1(e)$ (i.e., $e'$ is in the same orbit as $e$),
  \item $R(\pi_0(\tau_1)) \in \partial e'$,
  \item $\phi_R(e;\tau_1,\tau_2) \cdot R(\pi_0(\tau_2)) \in \partial e'$.
\end{enumerate}
\end{theorem}

\begin{proof}
\leanok
\uses{def:CellComplex}
Let $g_1 = \mathtt{vertexGroupElement}(R, \tau_1)$, so that $g_1 \cdot R(\pi_0(\tau_1)) = \tau_1$. Set $e' := g_1^{-1} \cdot e$.

\textbf{(1) Same orbit:} We have $\pi_1(e') = \pi_1(g_1^{-1} \cdot e) = \pi_1(e)$, since $e'$ and $e$ are related by $g_1^{-1} \in H$, giving a sound quotient relation (i.e., $\exists\, h,\, h \cdot e = e'$ with $h = g_1^{-1}$).

\textbf{(2) First boundary vertex:} We must show $R(\pi_0(\tau_1)) \in \partial(g_1^{-1} \cdot e)$. By boundary equivariance, $R(\pi_0(\tau_1)) \in \partial(g_1^{-1} \cdot e)$ iff $(g_1^{-1})^{-1} \cdot R(\pi_0(\tau_1)) = g_1 \cdot R(\pi_0(\tau_1)) \in \partial e$. Since $g_1 \cdot R(\pi_0(\tau_1)) = \tau_1$ by the specification of $\mathtt{vertexGroupElement}$, and $\tau_1 \in \partial e$ by hypothesis, this holds.

\textbf{(3) Second boundary vertex:} We must show $\phi_R(e;\tau_1,\tau_2) \cdot R(\pi_0(\tau_2)) \in \partial(g_1^{-1} \cdot e)$. Unfolding $\mathtt{connectionOnEdge}$, we have $\phi_R(e;\tau_1,\tau_2) = g_1^{-1} \cdot g_2$ where $g_2 = \mathtt{vertexGroupElement}(R,\tau_2)$. By boundary equivariance, $(g_1^{-1} \cdot g_2) \cdot R(\pi_0(\tau_2)) \in \partial(g_1^{-1} \cdot e)$ iff $g_1 \cdot (g_1^{-1} \cdot g_2) \cdot R(\pi_0(\tau_2)) = g_2 \cdot R(\pi_0(\tau_2)) = \tau_2 \in \partial e$, which holds by hypothesis. Rewriting using $g_1 \cdot g_1^{-1} \cdot \tau_2 = \tau_2$ completes the argument.
\end{proof}

\begin{lemma}[Quotient Vertices Nonempty]
\label{lem:GraphWithGroupAction.quotientVertices_nonempty}
\lean{GraphWithGroupAction.quotientVertices_nonempty}
\leanok
\uses{def:CellComplex}
If $X_0$ is nonempty, then $(X/H)_0$ is nonempty.
\end{lemma}

\begin{proof}
\leanok
\uses{def:CellComplex}
Since $X_0$ is nonempty, choose an arbitrary vertex $v \in X_0$ by $\mathtt{Classical.arbitrary}$. Then $[v] \in X_0/H$ witnesses that $(X/H)_0$ is nonempty.
\end{proof}

\begin{lemma}[Quotient Edges Nonempty]
\label{lem:GraphWithGroupAction.quotientEdges_nonempty}
\lean{GraphWithGroupAction.quotientEdges_nonempty}
\leanok
\uses{def:CellComplex}
If $X_1$ is nonempty, then $(X/H)_1$ is nonempty.
\end{lemma}

\begin{proof}
\leanok
\uses{def:CellComplex}
Since $X_1$ is nonempty, choose an arbitrary edge $e \in X_1$ by $\mathtt{Classical.arbitrary}$. Then $[e] \in X_1/H$ witnesses that $(X/H)_1$ is nonempty.
\end{proof}

%--- Def_25: InvariantLabeling ---
\chapter{Def 25: Invariant Labeling}

\begin{definition}[Cell Incident Edges]
\label{def:GraphWithGroupAction.cellIncidentEdges}
\lean{GraphWithGroupAction.cellIncidentEdges}
\leanok
\uses{def:GraphWithGroupAction}
Let $X$ be a graph with a group action by $H$. The set of edges incident to a vertex $v \in X_0$ in the cell complex is defined as
\[
\delta v = \{ e \in X_1 \mid v \in \partial e \}.
\]
Formally, $\operatorname{cellIncidentEdges}(v) = \{ e \in X_1 \mid v \in X.\operatorname{graph}.\operatorname{bdry}(e) \}$.
\end{definition}

\begin{theorem}[Membership in Cell Incident Edges]
\label{thm:GraphWithGroupAction.mem_cellIncidentEdges}
\lean{GraphWithGroupAction.mem_cellIncidentEdges}
\leanok
\uses{def:GraphWithGroupAction.cellIncidentEdges}
For a vertex $v \in X_0$ and an edge $e \in X_1$, we have
\[
e \in \delta v \iff v \in \partial e.
\]
\end{theorem}

\begin{proof}
\leanok
\uses{def:GraphWithGroupAction.cellIncidentEdges}
This follows immediately by simplification using the definition of $\operatorname{cellIncidentEdges}$.
\end{proof}

\begin{lemma}[Cell Incident Edges Nonempty]
\label{lem:GraphWithGroupAction.cellIncidentEdges_nonempty}
\lean{GraphWithGroupAction.cellIncidentEdges_nonempty}
\leanok
\uses{def:GraphWithGroupAction.cellIncidentEdges, thm:GraphWithGroupAction.mem_cellIncidentEdges}
For a vertex $v \in X_0$ and an edge $e \in X_1$ with $v \in \partial e$, the set $\delta v$ is non-empty.
\end{lemma}

\begin{proof}
\leanok
\uses{def:GraphWithGroupAction.cellIncidentEdges, thm:GraphWithGroupAction.mem_cellIncidentEdges}
The element $e$ witnesses non-emptiness: since $v \in \partial e$, by the membership characterization we have $e \in \delta v$, so $\langle e, \cdot \rangle$ provides the required element.
\end{proof}

\begin{definition}[Cell Complex Labeling]
\label{def:GraphWithGroupAction.CellLabeling}
\lean{GraphWithGroupAction.CellLabeling}
\leanok
\uses{def:GraphWithGroupAction, def:GraphWithGroupAction.cellIncidentEdges, def:Labeling}
A \emph{cell complex labeling of degree $s$} on $X$ is a family $\Lambda = \{\Lambda_v\}_{v \in X_0}$ assigning to each vertex $v \in X_0$ a bijection
\[
\Lambda_v : \delta v \xrightarrow{\sim} \operatorname{Fin}(s).
\]
Formally, $\operatorname{CellLabeling}(s) = \prod_{v : X_0} (\delta v \simeq \operatorname{Fin}(s))$.
\end{definition}

\begin{lemma}[Group Action Preserves Incidence]
\label{lem:GraphWithGroupAction.smul_incident_of_incident}
\lean{GraphWithGroupAction.smul_incident_of_incident}
\leanok
\uses{def:GraphWithGroupAction, def:GraphWithGroupAction.cellIncidentEdges, lem:GraphWithGroupAction.bdry_orbit_compat}
If $e \in \delta v$ (i.e., $e$ is incident to $v$), then $h \cdot e \in \delta(h \cdot v)$ for all $h \in H$. That is, the $H$-action preserves incidence.
\end{lemma}

\begin{proof}
\leanok
\uses{def:GraphWithGroupAction.cellIncidentEdges, thm:GraphWithGroupAction.mem_cellIncidentEdges, lem:GraphWithGroupAction.bdry_orbit_compat}
Rewriting using the membership characterization of $\operatorname{cellIncidentEdges}$ on both sides, it suffices to show $h \cdot v \in \partial(h \cdot e)$. This follows directly from the boundary orbit compatibility lemma $\operatorname{bdry_orbit_compat}$, which states that $v \in \partial e \Rightarrow h \cdot v \in \partial(h \cdot e)$.
\end{proof}

\begin{definition}[Transported Incident Edge]
\label{def:GraphWithGroupAction.smulIncidentEdge}
\lean{GraphWithGroupAction.smulIncidentEdge}
\leanok
\uses{def:GraphWithGroupAction, def:GraphWithGroupAction.cellIncidentEdges, lem:GraphWithGroupAction.smul_incident_of_incident}
For a vertex $v \in X_0$, a group element $h \in H$, and an incident edge $e \in \delta v$, define the transported incident edge
\[
h \cdot e \in \delta(h \cdot v)
\]
as the pair $\langle h \cdot e.{\operatorname{val}},\, \operatorname{smul_incident_of_incident}(v, e.{\operatorname{val}}, h, e.{\operatorname{prop}})\rangle$.
\end{definition}

\begin{definition}[$H$-Invariant Labeling]
\label{def:GraphWithGroupAction.IsInvariantLabeling}
\lean{GraphWithGroupAction.IsInvariantLabeling}
\leanok
\uses{def:GraphWithGroupAction, def:GraphWithGroupAction.CellLabeling, def:GraphWithGroupAction.smulIncidentEdge}
A cell complex labeling $\Lambda$ is \emph{$H$-invariant} if for all vertices $v \in X_0$, all group elements $h \in H$, and all incident edges $e \in \delta v$,
\[
\Lambda_{h \cdot v}(h \cdot e) = \Lambda_v(e).
\]
Formally, $\operatorname{IsInvariantLabeling}(\Lambda)$ is the proposition
\[
\forall\, v \in X_0,\; \forall\, h \in H,\; \forall\, e \in \delta v,\quad \Lambda(h \cdot v)(\operatorname{smulIncidentEdge}(v, h, e)) = \Lambda(v)(e).
\]
\end{definition}

\begin{lemma}[Invariance Condition is Satisfiable]
\label{lem:GraphWithGroupAction.isInvariantLabeling_trivial_satisfiable}
\lean{GraphWithGroupAction.isInvariantLabeling_trivial_satisfiable}
\leanok
\uses{def:GraphWithGroupAction, def:GraphWithGroupAction.CellLabeling, def:GraphWithGroupAction.IsInvariantLabeling}
There exist a group $H$, a graph with $H$-action $X$, a degree $s \in \mathbb{N}$, a cell labeling $\Lambda : X.\operatorname{CellLabeling}(s)$, and a proof that $\Lambda$ is $H$-invariant. In other words, the invariance premises are satisfiable.
\end{lemma}

\begin{proof}
\leanok
\uses{def:GraphWithGroupAction, def:GraphWithGroupAction.CellLabeling, def:GraphWithGroupAction.IsInvariantLabeling}
We take $H = \operatorname{Unit}$ (the trivial group) with $s = 0$. We construct $X$ as the graph with trivial action in which all cell types are $\operatorname{PEmpty}$ (no vertices or edges), all structure maps are vacuously defined, and the vertex/edge actions are trivial. The boundary map sends every edge (of which there are none) to the empty set, and the boundary-of-boundary condition holds vacuously by simplification.

We provide the labeling $\Lambda$ and the invariance proof vacuously: since there are no vertices ($X_0 = \operatorname{PEmpty}$), the universal quantifiers in both the labeling and invariance condition hold trivially via $\operatorname{PEmpty.elim}$.
\end{proof}

\begin{definition}[Quotient Incident Edges]
\label{def:GraphWithGroupAction.quotientIncidentEdges}
\lean{GraphWithGroupAction.quotientIncidentEdges}
\leanok
\uses{def:GraphWithGroupAction, def:GraphWithGroupAction.edgeQuot, def:GraphWithGroupAction.vertexQuot, def:GraphWithGroupAction.quotientEdges, def:GraphWithGroupAction.quotientVertices}
For a vertex orbit $[v] \in X/H$, the \emph{quotient incident edges} are the edge orbits $[e] \in X_1/H$ such that some representative edge has a boundary vertex in the orbit of $[v]$:
\[
\delta[v] = \bigl\{ [e] \in X_1/H \mid \exists\, e \in [e],\; \exists\, v' \in [v],\; v' \in \partial e \bigr\}.
\]
\end{definition}

\begin{theorem}[Membership in Quotient Incident Edges]
\label{thm:GraphWithGroupAction.mem_quotientIncidentEdges}
\lean{GraphWithGroupAction.mem_quotientIncidentEdges}
\leanok
\uses{def:GraphWithGroupAction.quotientIncidentEdges}
For a vertex orbit $[v] \in X/H$ and an edge orbit $[e] \in X_1/H$,
\[
[e] \in \delta[v] \iff \exists\, e' \in [e],\; \exists\, v' \in [v],\; v' \in \partial e'.
\]
\end{theorem}

\begin{proof}
\leanok
\uses{def:GraphWithGroupAction.quotientIncidentEdges}
This follows immediately by simplification using the definition of $\operatorname{quotientIncidentEdges}$.
\end{proof}

\begin{lemma}[Quotient Incidence from Incidence]
\label{lem:GraphWithGroupAction.quotientIncident_of_incident}
\lean{GraphWithGroupAction.quotientIncident_of_incident}
\leanok
\uses{def:GraphWithGroupAction.quotientIncidentEdges, def:GraphWithGroupAction.cellIncidentEdges, thm:GraphWithGroupAction.mem_quotientIncidentEdges, thm:GraphWithGroupAction.mem_cellIncidentEdges}
If $e \in \delta v$ (i.e., $e$ is incident to $v$ in the cell complex), then the edge orbit $[e]$ is incident to the vertex orbit $[v]$ in the quotient graph:
\[
[e] \in \delta[v].
\]
\end{lemma}

\begin{proof}
\leanok
\uses{def:GraphWithGroupAction.quotientIncidentEdges, thm:GraphWithGroupAction.mem_quotientIncidentEdges, thm:GraphWithGroupAction.mem_cellIncidentEdges}
Rewriting using the membership characterization of $\operatorname{quotientIncidentEdges}$, we provide the witnesses: $e$ itself as the edge representative, $v$ as the vertex representative, the reflexivity proofs $[e] = [e]$ and $[v] = [v]$, and the boundary membership $v \in \partial e$ extracted from $e \in \delta v$ via the membership characterization of $\operatorname{cellIncidentEdges}$.
\end{proof}

\begin{lemma}[Quotient Incident Edges Nonempty]
\label{lem:GraphWithGroupAction.quotientIncidentEdges_nonempty}
\lean{GraphWithGroupAction.quotientIncidentEdges_nonempty}
\leanok
\uses{def:GraphWithGroupAction.quotientIncidentEdges, lem:GraphWithGroupAction.quotientIncident_of_incident, thm:GraphWithGroupAction.mem_cellIncidentEdges}
For any vertex $v \in X_0$ and edge $e \in X_1$ with $v \in \partial e$, the quotient incident edge set $\delta[v]$ is non-empty.
\end{lemma}

\begin{proof}
\leanok
\uses{def:GraphWithGroupAction.quotientIncidentEdges, lem:GraphWithGroupAction.quotientIncident_of_incident, thm:GraphWithGroupAction.mem_cellIncidentEdges}
The edge orbit $[e]$ witnesses non-emptiness: since $v \in \partial e$, by the membership characterization of $\operatorname{cellIncidentEdges}$ we have $e \in \delta v$, and then by $\operatorname{quotientIncident_of_incident}$ we obtain $[e] \in \delta[v]$.
\end{proof}

\begin{definition}[Quotient Labeling]
\label{def:GraphWithGroupAction.quotientLabeling}
\lean{GraphWithGroupAction.quotientLabeling}
\leanok
\uses{def:GraphWithGroupAction, def:GraphWithGroupAction.CellLabeling, def:GraphWithGroupAction.IsInvariantLabeling, def:GraphWithGroupAction.quotientIncidentEdges, def:GraphWithGroupAction.smulIncidentEdge, lem:GraphWithGroupAction.quotientIncident_of_incident}
Given an $H$-invariant labeling $\Lambda$ (i.e., $\operatorname{IsInvariantLabeling}(\Lambda)$ holds), the \emph{quotient labeling} $\bar{\Lambda}$ is a well-defined bijection
\[
\bar{\Lambda}_{[v]} : \delta[v] \xrightarrow{\sim} \operatorname{Fin}(s)
\]
for each vertex orbit $[v] \in X/H$, defined by $\bar{\Lambda}_{[v]}([e]) = \Lambda_v(e)$ for any representatives $v \in [v]$ and $e \in [e]$ with $e \in \delta v$.
\end{definition}

\begin{proof}
\leanok
\uses{def:GraphWithGroupAction.CellLabeling, def:GraphWithGroupAction.IsInvariantLabeling, def:GraphWithGroupAction.quotientIncidentEdges, def:GraphWithGroupAction.cellIncidentEdges, def:GraphWithGroupAction.smulIncidentEdge, lem:GraphWithGroupAction.quotientIncident_of_incident, thm:GraphWithGroupAction.mem_quotientIncidentEdges, thm:GraphWithGroupAction.mem_cellIncidentEdges, lem:GraphWithGroupAction.bdry_orbit_compat}
We construct the bijection as follows. Given a vertex orbit $[v]$, we choose a representative vertex $v$ using $\operatorname{Quotient.exists_rep}$, obtaining $v$ and the proof $[v] = [v]$.

We define a translation map: for each quotient incident edge $[e] \in \delta[v]$, we produce an incident edge in $\delta v$. By the membership characterization of $\operatorname{quotientIncidentEdges}$, there exists a representative edge $e$ with $[e'] = [e]$ and a vertex $v'$ with $[v'] = [v]$ and $v' \in \partial e'$. Since $[v'] = [v]$, by $\operatorname{Quotient.exact}$ there exists $g \in H$ with $g \cdot v = v'$. Then $g^{-1} \cdot e'$ is incident to $v$: indeed $v = g^{-1} \cdot v' \in \partial(g^{-1} \cdot e')$ by the boundary equivariance of $X$. We thus obtain the element $\langle g^{-1} \cdot e', \cdot \rangle \in \delta v$, and it lies in the orbit $[e]$ since $g \cdot (g^{-1} \cdot e') = e'$ has the same quotient.

We define $\operatorname{pickEdge}([e])$ to be this translated incident edge. We verify:

\textbf{Injectivity of pickEdge:} If $\operatorname{pickEdge}([e_1]) = \operatorname{pickEdge}([e_2])$, then both have the same underlying edge, hence the same image under $\operatorname{edgeQuot}$ (by $\operatorname{pickEdge_spec}$), so $[e_1] = [e_2]$ as elements of $\delta[v]$.

\textbf{Surjectivity of pickEdge:} For any $e \in \delta v$, the quotient edge $[e] \in \delta[v]$ (by $\operatorname{quotientIncident_of_incident}$). We have $\operatorname{pickEdge}([e]) = $ some $e' \in \delta v$ with $[e'] = [e]$. By $\operatorname{Quotient.exact}$, there exists $g$ with $g \cdot e' = e$ (or $g \cdot e = e'$). Both $e$ and $e'$ are in $\delta v$, so $v \in \partial e$ and $v \in \partial e'$. From $g \cdot e' = e$ and equivariance, $g^{-1} \cdot v \in \partial e'$, and by the quotient condition $\operatorname{quotientCondition}$ (which states that if $v$ and $g^{-1} \cdot v$ both lie in $\partial e'$ then $g^{-1} = 1$), we get $g = 1$, hence $e = e'$.

The composition $\Lambda_v \circ \operatorname{pickEdge}$ is therefore a bijection $\delta[v] \to \operatorname{Fin}(s)$: it is injective as a composition of two injections ($\operatorname{pickEdge}$ injective and $\Lambda_v$ injective), and surjective as a composition of two surjections ($\operatorname{pickEdge}$ surjective and $\Lambda_v$ surjective). We conclude by $\operatorname{Equiv.ofBijective}$.
\end{proof}

\begin{theorem}[Quotient Labeling is Well-Defined]
\label{thm:GraphWithGroupAction.quotientLabeling_well_defined}
\lean{GraphWithGroupAction.quotientLabeling_well_defined}
\leanok
\uses{def:GraphWithGroupAction, def:GraphWithGroupAction.CellLabeling, def:GraphWithGroupAction.IsInvariantLabeling, def:GraphWithGroupAction.cellIncidentEdges, thm:GraphWithGroupAction.mem_cellIncidentEdges}
Let $\Lambda$ be an $H$-invariant labeling. For any vertex $v \in X_0$, edge $e \in X_1$ with $v \in \partial e$, and $h \in H$ (so that $h \cdot v \in \partial(h \cdot e)$), the labeling value is independent of the representative:
\[
\Lambda_{h \cdot v}\bigl(\langle h \cdot e,\, \cdot \rangle\bigr) = \Lambda_v\bigl(\langle e,\, \cdot \rangle\bigr).
\]
\end{theorem}

\begin{proof}
\leanok
\uses{def:GraphWithGroupAction.IsInvariantLabeling, def:GraphWithGroupAction.smulIncidentEdge, thm:GraphWithGroupAction.mem_cellIncidentEdges}
We apply the $H$-invariance condition to $v$, $h$, and the incident edge $\langle e, \cdot \rangle \in \delta v$, obtaining
\[
\Lambda(h \cdot v)\bigl(\operatorname{smulIncidentEdge}(v, h, \langle e, \cdot \rangle)\bigr) = \Lambda(v)\bigl(\langle e, \cdot \rangle\bigr).
\]
We observe that the underlying cell of $\operatorname{smulIncidentEdge}(v, h, \langle e, \cdot \rangle)$ is definitionally equal to $h \cdot e$. Rewriting with the $H$-invariance equation and using congruence of the subtype (since the values agree), we conclude the result.
\end{proof}

%--- Def_26: BalancedProductTannerCycleCode ---
\chapter{Def 26: Balanced Product Tanner Cycle Code}

\begin{definition}[Cell Local View]
\label{def:BalancedProductTannerCycleCode.cellLocalView}
\lean{BalancedProductTannerCycleCode.cellLocalView}
\leanok
\uses{def:CSSCode.complex, def:HEquivariantChainComplex}
Let $H$ be a finite group acting on a graph with group action $X$, and let $\Lambda$ be a cell labeling of $X$ with $s$ labels. For a vertex $v \in C_0(X)$ (a 0-cell of $X$), the \emph{cell local view} at $v$ is the linear map
\[
\operatorname{cellLocalView}(v) : (C_1(X) \to \mathbb{F}_2) \to (\operatorname{Fin}(s) \to \mathbb{F}_2)
\]
defined by $\operatorname{cellLocalView}(v)(c)(i) = c\!\left(\Lambda(v)^{-1}(i)\right)$, where $\Lambda(v)^{-1}(i)$ denotes the preimage of $i$ under the labeling bijection at $v$.
\end{definition}

\begin{definition}[Tanner Differential]
\label{def:BalancedProductTannerCycleCode.tannerDifferential}
\lean{BalancedProductTannerCycleCode.tannerDifferential}
\leanok
\uses{def:CSSCode.complex, def:HEquivariantChainComplex, def:BalancedProductTannerCycleCode.cellLocalView}
The \emph{Tanner differential} is the linear map
\[
\partial^{\mathrm{Tan}} : (C_1(X) \to \mathbb{F}_2) \to (C_0(X) \to \operatorname{Fin}(s) \to \mathbb{F}_2)
\]
defined by $\partial^{\mathrm{Tan}}(c)(v) = \operatorname{cellLocalView}(v)(c)$, i.e.,
\[
\partial^{\mathrm{Tan}}(c)(v)(i) = c\!\left(\Lambda(v)^{-1}(i)\right).
\]
This map assembles the local views at all vertices.
\end{definition}

\begin{definition}[Tanner Code Chain Complex]
\label{def:BalancedProductTannerCycleCode.tannerCodeComplex}
\lean{BalancedProductTannerCycleCode.tannerCodeComplex}
\leanok
\uses{def:CSSCode.complex, def:HEquivariantChainComplex, def:BalancedProductTannerCycleCode.tannerDifferential}
The \emph{Tanner code chain complex} $C(X, \Lambda)$ is the chain complex over $\mathbb{F}_2$ defined by:
\[
\cdots \to 0 \to C_1(X,\Lambda) \xrightarrow{\partial^{\mathrm{Tan}}} C_0(X,\Lambda) \to 0 \to \cdots
\]
where
\begin{itemize}
  \item $C_1(X,\Lambda) = (C_1(X) \to \mathbb{F}_2)$, the $\mathbb{F}_2$-module of functions on 1-cells of $X$,
  \item $C_0(X,\Lambda) = (C_0(X) \to \operatorname{Fin}(s) \to \mathbb{F}_2)$, the $\mathbb{F}_2$-module of labeled functions on 0-cells,
  \item the differential at degrees $(1,0)$ is $\partial^{\mathrm{Tan}}$, and all other differentials are zero,
  \item all modules outside degrees $0$ and $1$ are the trivial module $\{*\}$.
\end{itemize}
The chain complex condition $\partial^{\mathrm{Tan}} \circ \partial^{\mathrm{Tan}} = 0$ holds trivially since the complex is concentrated in two consecutive degrees.
\end{definition}

\begin{definition}[$H$-Action on $C_1$]
\label{def:BalancedProductTannerCycleCode.actionC1}
\lean{BalancedProductTannerCycleCode.actionC1}
\leanok
\uses{def:HEquivariantChainComplex}
The \emph{$H$-representation on $C_1(X,\Lambda)$} is the $\mathbb{F}_2$-linear group action
\[
\rho_1 : H \to \operatorname{GL}(C_1(X) \to \mathbb{F}_2)
\]
defined by $(\rho_1(h) \cdot c)(e) = c(h^{-1} \cdot e)$ for $h \in H$ and $e \in C_1(X)$. This is the standard permutation representation by precomposition with the inverse group action.
\end{definition}

\begin{definition}[$H$-Action on $C_0$]
\label{def:BalancedProductTannerCycleCode.actionC0}
\lean{BalancedProductTannerCycleCode.actionC0}
\leanok
\uses{def:HEquivariantChainComplex}
The \emph{$H$-representation on $C_0(X,\Lambda)$} is the $\mathbb{F}_2$-linear group action
\[
\rho_0 : H \to \operatorname{GL}(C_0(X) \to \operatorname{Fin}(s) \to \mathbb{F}_2)
\]
defined by $(\rho_0(h) \cdot f)(v) = f(h^{-1} \cdot v)$ for $h \in H$ and $v \in C_0(X)$.
\end{definition}

\begin{definition}[Tanner Code $H$-Equivariant Complex]
\label{def:BalancedProductTannerCycleCode.tannerCodeHEquivariant}
\lean{BalancedProductTannerCycleCode.tannerCodeHEquivariant}
\leanok
\uses{def:CSSCode.complex, def:HEquivariantChainComplex, def:BalancedProductTannerCycleCode.tannerCodeComplex, def:BalancedProductTannerCycleCode.actionC1, def:BalancedProductTannerCycleCode.actionC0}
Given an $H$-invariant labeling $\Lambda$ (i.e., $\Lambda$ satisfies the invariance condition $\operatorname{IsInvariantLabeling}(\Lambda)$), the Tanner code chain complex $C(X, \Lambda)$ equipped with the $H$-representations $\rho_1$ on $C_1$ and $\rho_0$ on $C_0$ forms an \emph{$H$-equivariant chain complex} $\widetilde{C}(X,\Lambda) \in \operatorname{HEquivariantChainComplex}(H)$.

The equivariance condition requires that $\partial^{\mathrm{Tan}} \circ \rho_1(h) = \rho_0(h) \circ \partial^{\mathrm{Tan}}$ for all $h \in H$, i.e., the Tanner differential commutes with the $H$-action.
\end{definition}

\begin{definition}[Cycle Complex]
\label{def:BalancedProductTannerCycleCode.cycleComplex}
\lean{BalancedProductTannerCycleCode.cycleComplex}
\leanok
\uses{def:CSSCode.complex, def:HEquivariantChainComplex, def:CycleGraph.differential}
For $\ell \geq 1$, the \emph{cycle graph chain complex} $C(C_\ell)$ is the chain complex over $\mathbb{F}_2$ defined by:
\[
\cdots \to 0 \to (\operatorname{Fin}(\ell) \to \mathbb{F}_2) \xrightarrow{\partial_{C_\ell}} (\operatorname{Fin}(\ell) \to \mathbb{F}_2) \to 0 \to \cdots
\]
where both degree-$1$ and degree-$0$ components equal $(\operatorname{Fin}(\ell) \to \mathbb{F}_2)$, the differential at $(1,0)$ is the cycle graph differential $\partial_{C_\ell}$, and all modules outside degrees $0$ and $1$ are the trivial module.
\end{definition}

\begin{definition}[Cycle Compatible Action]
\label{def:BalancedProductTannerCycleCode.CycleCompatibleAction}
\lean{BalancedProductTannerCycleCode.CycleCompatibleAction}
\leanok
\uses{def:HEquivariantChainComplex}
A \emph{cycle compatible $H$-action} on $\operatorname{Fin}(\ell)$ is an $H$-action satisfying the compatibility condition: for all $h \in H$ and $i \in \operatorname{Fin}(\ell)$,
\[
h \cdot \left\langle (i + \ell - 1) \bmod \ell \right\rangle = \left\langle (h \cdot i + \ell - 1) \bmod \ell \right\rangle,
\]
where $\langle k \rangle$ denotes the element of $\operatorname{Fin}(\ell)$ with value $k$. In other words, the $H$-action commutes with taking the predecessor modulo $\ell$.
\end{definition}

\begin{definition}[Cycle Chain Action]
\label{def:BalancedProductTannerCycleCode.cycleChainAction}
\lean{BalancedProductTannerCycleCode.cycleChainAction}
\leanok
\uses{def:HEquivariantChainComplex}
The \emph{cycle chain $H$-representation} is the $\mathbb{F}_2$-linear group action
\[
\rho_C : H \to \operatorname{GL}(\operatorname{Fin}(\ell) \to \mathbb{F}_2)
\]
defined by $(\rho_C(h) \cdot f)(i) = f(h^{-1} \cdot i)$ for $h \in H$ and $i \in \operatorname{Fin}(\ell)$.
\end{definition}

\begin{definition}[Cycle Graph $H$-Equivariant Complex]
\label{def:BalancedProductTannerCycleCode.cycleGraphHEquivariant}
\lean{BalancedProductTannerCycleCode.cycleGraphHEquivariant}
\leanok
\uses{def:CSSCode.complex, def:HEquivariantChainComplex, def:BalancedProductTannerCycleCode.cycleComplex, def:BalancedProductTannerCycleCode.cycleChainAction, def:BalancedProductTannerCycleCode.CycleCompatibleAction}
Given a cycle compatible $H$-action on $\operatorname{Fin}(\ell)$, the cycle graph chain complex $C(C_\ell)$ equipped with the $H$-representation $\rho_C$ on both degrees $0$ and $1$ forms an \emph{$H$-equivariant chain complex} $\widetilde{C}(C_\ell) \in \operatorname{HEquivariantChainComplex}(H)$.

The equivariance condition requires $\partial_{C_\ell} \circ \rho_C(h) = \rho_C(h) \circ \partial_{C_\ell}$ for all $h \in H$, which follows from the cycle compatible action hypothesis.
\end{definition}

\begin{definition}[Balanced Product Tanner Cycle Code]
\label{def:BalancedProductTannerCycleCode.balancedProductTannerCycleCode}
\lean{BalancedProductTannerCycleCode.balancedProductTannerCycleCode}
\leanok
\uses{def:CSSCode.complex, def:HEquivariantChainComplex, def:BalancedProductTannerCycleCode.tannerCodeHEquivariant, def:BalancedProductTannerCycleCode.cycleGraphHEquivariant, def:HEquivariantChainComplex.balancedProductComplex}
The \emph{balanced product Tanner cycle code} is the chain complex
\[
C(X,\Lambda) \otimes_{\mathbb{Z}_\ell} C(C_\ell)
\]
defined as the balanced product complex of the $H$-equivariant Tanner code complex and the $H$-equivariant cycle graph complex:
\[
\operatorname{balancedProductTannerCycleCode}(X, \Lambda, \ell) = \widetilde{C}(X,\Lambda) \otimes_H \widetilde{C}(C_\ell).
\]
This requires $\ell \geq 3$, $\ell$ odd, an $H$-invariant labeling $\Lambda$, and a cycle compatible $H$-action on $\operatorname{Fin}(\ell)$.

The total complex has three non-trivial degrees:
\begin{align*}
  C_2 &= C_1(X,\Lambda) \otimes_H C_1(C_\ell), \\
  C_1 &= \bigl(C_1(X,\Lambda) \otimes_H C_0(C_\ell)\bigr) \oplus \bigl(C_0(X,\Lambda) \otimes_H C_1(C_\ell)\bigr), \\
  C_0 &= C_0(X,\Lambda) \otimes_H C_0(C_\ell),
\end{align*}
where physical qubits correspond to $C_1$, Z-checks are given by $\partial_2 : C_2 \to C_1$, and X-checks are given by $\partial_1 : C_1 \to C_0$.
\end{definition}

\begin{definition}[Inclusion of Tanner Component into $C_1$]
\label{def:BalancedProductTannerCycleCode.iotaC1_tanner_component}
\lean{BalancedProductTannerCycleCode.ιC1_tanner_component}
\leanok
\uses{def:CSSCode.complex, def:HEquivariantChainComplex, def:BalancedProductTannerCycleCode.balancedProductTannerCycleCode}
The canonical inclusion map
\[
\iota_{\mathrm{Tan}} : C_1(X,\Lambda) \otimes_H C_0(C_\ell) \hookrightarrow \bigl(\text{balancedProductTannerCycleCode}\bigr)_1
\]
is the map into the total complex at degree $1$ corresponding to the $(1,0)$-summand of the balanced product double complex.
\end{definition}

\begin{definition}[Inclusion of Cycle Component into $C_1$]
\label{def:BalancedProductTannerCycleCode.iotaC1_cycle_component}
\lean{BalancedProductTannerCycleCode.ιC1_cycle_component}
\leanok
\uses{def:CSSCode.complex, def:HEquivariantChainComplex, def:BalancedProductTannerCycleCode.balancedProductTannerCycleCode}
The canonical inclusion map
\[
\iota_{\mathrm{Cyc}} : C_0(X,\Lambda) \otimes_H C_1(C_\ell) \hookrightarrow \bigl(\text{balancedProductTannerCycleCode}\bigr)_1
\]
is the map into the total complex at degree $1$ corresponding to the $(0,1)$-summand of the balanced product double complex.
\end{definition}

\begin{definition}[Inclusion into $C_2$]
\label{def:BalancedProductTannerCycleCode.iotaC2_component}
\lean{BalancedProductTannerCycleCode.ιC2_component}
\leanok
\uses{def:CSSCode.complex, def:HEquivariantChainComplex, def:BalancedProductTannerCycleCode.balancedProductTannerCycleCode}
The canonical inclusion map
\[
\iota_2 : C_1(X,\Lambda) \otimes_H C_1(C_\ell) \hookrightarrow \bigl(\text{balancedProductTannerCycleCode}\bigr)_2
\]
is the map into the total complex at degree $2$ corresponding to the $(1,1)$-summand of the balanced product double complex.
\end{definition}

\begin{definition}[Inclusion into $C_0$]
\label{def:BalancedProductTannerCycleCode.iotaC0_component}
\lean{BalancedProductTannerCycleCode.ιC0_component}
\leanok
\uses{def:CSSCode.complex, def:HEquivariantChainComplex, def:BalancedProductTannerCycleCode.balancedProductTannerCycleCode}
The canonical inclusion map
\[
\iota_0 : C_0(X,\Lambda) \otimes_H C_0(C_\ell) \hookrightarrow \bigl(\text{balancedProductTannerCycleCode}\bigr)_0
\]
is the map into the total complex at degree $0$ corresponding to the $(0,0)$-summand of the balanced product double complex.
\end{definition}

\begin{definition}[Z-Check Map]
\label{def:BalancedProductTannerCycleCode.zCheckMap}
\lean{BalancedProductTannerCycleCode.zCheckMap}
\leanok
\uses{def:CSSCode.complex, def:HEquivariantChainComplex, def:BalancedProductTannerCycleCode.balancedProductTannerCycleCode}
The \emph{Z-check map} of the balanced product Tanner cycle code is the differential
\[
\partial_2 : \bigl(\text{balancedProductTannerCycleCode}\bigr)_2 \to \bigl(\text{balancedProductTannerCycleCode}\bigr)_1,
\]
i.e., the chain complex differential from degree $2$ to degree $1$. Z-errors are detected by this map.
\end{definition}

\begin{definition}[X-Check Map]
\label{def:BalancedProductTannerCycleCode.xCheckMap}
\lean{BalancedProductTannerCycleCode.xCheckMap}
\leanok
\uses{def:CSSCode.complex, def:HEquivariantChainComplex, def:BalancedProductTannerCycleCode.balancedProductTannerCycleCode}
The \emph{X-check map} of the balanced product Tanner cycle code is the differential
\[
\partial_1 : \bigl(\text{balancedProductTannerCycleCode}\bigr)_1 \to \bigl(\text{balancedProductTannerCycleCode}\bigr)_0,
\]
i.e., the chain complex differential from degree $1$ to degree $0$. X-errors are detected by this map.
\end{definition}

\begin{theorem}[Tensor Product Identification at $C_2$]
\label{thm:BalancedProductTannerCycleCode.balancedProductTannerCycleCode_C2_eq}
\lean{BalancedProductTannerCycleCode.balancedProductTannerCycleCode_C2_eq}
\leanok
\uses{def:CSSCode.complex, def:HEquivariantChainComplex, def:BalancedProductTannerCycleCode.balancedProductTannerCycleCode, def:HEquivariantChainComplex.bpObj}
The degree-$(1,1)$ entry of the balanced product double complex equals the balanced product object at $(1,1)$:
\[
\bigl(\widetilde{C}(X,\Lambda) \boxtimes_H \widetilde{C}(C_\ell)\bigr)_{1,1} = \operatorname{bpObj}(1, 1) = C_1(X,\Lambda) \otimes_H C_1(C_\ell).
\]
\end{theorem}

\begin{proof}
\leanok
\uses{def:CSSCode.complex, def:HEquivariantChainComplex, def:HEquivariantChainComplex.bpObj}
This holds by reflexivity, i.e., definitional equality of the two expressions.
\end{proof}

\begin{theorem}[Tensor Product Identification at $C_0$]
\label{thm:BalancedProductTannerCycleCode.balancedProductTannerCycleCode_C0_eq}
\lean{BalancedProductTannerCycleCode.balancedProductTannerCycleCode_C0_eq}
\leanok
\uses{def:CSSCode.complex, def:HEquivariantChainComplex, def:BalancedProductTannerCycleCode.balancedProductTannerCycleCode, def:HEquivariantChainComplex.bpObj}
The degree-$(0,0)$ entry of the balanced product double complex equals the balanced product object at $(0,0)$:
\[
\bigl(\widetilde{C}(X,\Lambda) \boxtimes_H \widetilde{C}(C_\ell)\bigr)_{0,0} = \operatorname{bpObj}(0, 0) = C_0(X,\Lambda) \otimes_H C_0(C_\ell).
\]
\end{theorem}

\begin{proof}
\leanok
\uses{def:CSSCode.complex, def:HEquivariantChainComplex, def:HEquivariantChainComplex.bpObj}
This holds by reflexivity, i.e., definitional equality of the two expressions.
\end{proof}

\begin{theorem}[CSS Condition]
\label{thm:BalancedProductTannerCycleCode.css_condition}
\lean{BalancedProductTannerCycleCode.css_condition}
\leanok
\uses{def:CSSCode.complex, def:HEquivariantChainComplex, def:BalancedProductTannerCycleCode.balancedProductTannerCycleCode, def:BalancedProductTannerCycleCode.zCheckMap, def:BalancedProductTannerCycleCode.xCheckMap}
The composition of the Z-check map and the X-check map is zero:
\[
\partial_1 \circ \partial_2 = 0 : C_2 \to C_1 \to C_0.
\]
This is the fundamental CSS condition ensuring that Z-stabilizers are orthogonal to X-stabilizers, making the balanced product Tanner cycle code a valid quantum CSS code.
\end{theorem}

\begin{proof}
\leanok
\uses{def:BalancedProductTannerCycleCode.balancedProductTannerCycleCode, def:BalancedProductTannerCycleCode.zCheckMap, def:BalancedProductTannerCycleCode.xCheckMap}
By the chain complex condition applied to the balanced product Tanner cycle code, the composition of consecutive differentials vanishes. Specifically, we apply $\texttt{d\_comp\_d}$ at degrees $2$, $1$, $0$ to the chain complex $\operatorname{balancedProductTannerCycleCode}(X, \Lambda, \ell)$, yielding $\partial_1 \circ \partial_2 = 0$ directly.
\end{proof}

%--- Def_27: HorizontalVerticalHomologySplitting ---
\chapter{Def 27: Horizontal/Vertical Homology Splitting}

\begin{definition}[Tanner Code H-Equivariant Chain Complex (local abbreviation)]
\label{def:HorizontalVerticalHomologySplitting.tannerHEq}
\lean{HorizontalVerticalHomologySplitting.tannerHEq}
\leanok
\uses{def:BalancedProductTannerCycleCode.tannerCodeHEquivariant, def:GraphWithGroupAction.IsInvariantLabeling, def:GraphWithGroupAction.CellLabeling}

Let $H$ be a finite group acting on a graph $X$ with cell labeling $\Lambda$ that is invariant under the action. The \emph{Tanner code $H$-equivariant chain complex} is defined as
\[
  \operatorname{tannerHEq}(X, \Lambda, h_\Lambda) \;:=\; \operatorname{tannerCodeHEquivariant}(X, \Lambda, h_\Lambda),
\]
the $H$-equivariant chain complex associated to the Tanner code on the graph with group action $X$ and labeling $\Lambda$.
\end{definition}

\begin{definition}[Cycle Graph H-Equivariant Chain Complex (local abbreviation)]
\label{def:HorizontalVerticalHomologySplitting.cycleHEq}
\lean{HorizontalVerticalHomologySplitting.cycleHEq}
\leanok
\uses{def:BalancedProductTannerCycleCode.cycleGraphHEquivariant, def:BalancedProductTannerCycleCode.CycleCompatibleAction}

Let $\ell \geq 3$ be an odd natural number with $H$ acting on $\operatorname{Fin}(\ell)$ via a cycle-compatible action $h_{\mathrm{compat}}$. The \emph{cycle graph $H$-equivariant chain complex} is defined as
\[
  \operatorname{cycleHEq}(\ell, h_{\mathrm{compat}}) \;:=\; \operatorname{cycleGraphHEquivariant}(\ell, h_{\mathrm{compat}}),
\]
the $H$-equivariant chain complex associated to the cycle graph $C_\ell$.
\end{definition}

\begin{definition}[Balanced Product Complex]
\label{def:HorizontalVerticalHomologySplitting.bpComplex}
\lean{HorizontalVerticalHomologySplitting.bpComplex}
\leanok
\uses{def:HEquivariantChainComplex.balancedProductComplex, def:HorizontalVerticalHomologySplitting.tannerHEq, def:HorizontalVerticalHomologySplitting.cycleHEq}

The \emph{balanced product complex} $C(X, \Lambda) \otimes_H C(C_\ell)$ is defined as
\[
  \operatorname{bpComplex}(X, \Lambda, \ell, h_\Lambda, h_{\mathrm{compat}}) \;:=\;
  \operatorname{tannerHEq}(X, \Lambda, h_\Lambda).\operatorname{balancedProductComplex}\!\bigl(\operatorname{cycleHEq}(\ell, h_{\mathrm{compat}})\bigr),
\]
the balanced product of the Tanner code complex and the cycle graph complex over $H$.
\end{definition}

\begin{definition}[Degree-1 Homology of the Balanced Product]
\label{def:HorizontalVerticalHomologySplitting.H1}
\lean{HorizontalVerticalHomologySplitting.H1}
\leanok
\uses{def:HorizontalVerticalHomologySplitting.bpComplex, def:ChainComplexF2.homologyPrime}

The \emph{degree-1 homology} of the balanced product complex is defined as
\[
  H_1 \;:=\; \operatorname{bpComplex}(X, \Lambda, \ell, h_\Lambda, h_{\mathrm{compat}}).\operatorname{homology}'(1),
\]
the first homology group of $C(X, \Lambda) \otimes_H C(C_\ell)$ over $\mathbb{F}_2$.
\end{definition}

\begin{definition}[Horizontal Homology $H_1^h$]
\label{def:HorizontalVerticalHomologySplitting.H1h}
\lean{HorizontalVerticalHomologySplitting.H1h}
\leanok
\uses{def:HEquivariantChainComplex.balancedKunnethSummand, def:HorizontalVerticalHomologySplitting.tannerHEq, def:HorizontalVerticalHomologySplitting.cycleHEq}

The \emph{horizontal homology} is the balanced K\"{u}nneth summand at indices $p = 1$, $q = 0$:
\[
  H_1^h \;:=\; \operatorname{balancedKunnethSummand}\!\bigl(\operatorname{tannerHEq}(X, \Lambda, h_\Lambda),\; \operatorname{cycleHEq}(\ell, h_{\mathrm{compat}}),\; 1,\; 1\bigr)
  \;=\; H_1\!\bigl(C(X, \Lambda)\bigr) \otimes_H H_0\!\bigl(C_\ell\bigr).
\]
\end{definition}

\begin{definition}[Vertical Homology $H_1^v$]
\label{def:HorizontalVerticalHomologySplitting.H1v}
\lean{HorizontalVerticalHomologySplitting.H1v}
\leanok
\uses{def:HEquivariantChainComplex.balancedKunnethSummand, def:HorizontalVerticalHomologySplitting.tannerHEq, def:HorizontalVerticalHomologySplitting.cycleHEq}

The \emph{vertical homology} is the balanced K\"{u}nneth summand at indices $p = 0$, $q = 1$:
\[
  H_1^v \;:=\; \operatorname{balancedKunnethSummand}\!\bigl(\operatorname{tannerHEq}(X, \Lambda, h_\Lambda),\; \operatorname{cycleHEq}(\ell, h_{\mathrm{compat}}),\; 1,\; 0\bigr)
  \;=\; H_0\!\bigl(C(X, \Lambda)\bigr) \otimes_H H_1\!\bigl(C_\ell\bigr).
\]
\end{definition}

\begin{definition}[K\"{u}nneth Isomorphism at Degree 1]
\label{def:HorizontalVerticalHomologySplitting.kunnethIso}
\lean{HorizontalVerticalHomologySplitting.kunnethIso}
\leanok
\uses{thm:HEquivariantChainComplex.kunnethBalancedProduct, def:HEquivariantChainComplex.balancedKunnethSum, def:HorizontalVerticalHomologySplitting.H1}

Assuming $|H|$ is odd, the \emph{K\"{u}nneth isomorphism at degree $1$} is the $\mathbb{F}_2$-linear equivalence
\[
  \phi \;:\; H_1\!\bigl(C(X,\Lambda) \otimes_H C(C_\ell)\bigr)
  \;\xrightarrow{\;\sim\;}
  \bigoplus_{p \in \mathbb{Z}} H_p\!\bigl(C(X,\Lambda)\bigr) \otimes_H H_{1-p}\!\bigl(C_\ell\bigr),
\]
defined as $\phi := \operatorname{kunnethBalancedProduct}\!\bigl(\operatorname{tannerHEq}(X, \Lambda, h_\Lambda),\; \operatorname{cycleHEq}(\ell, h_{\mathrm{compat}}),\; h_{\mathrm{odd}},\; 1\bigr)$.
\end{definition}

\begin{definition}[Direct Sum Projection onto $p$-th Summand]
\label{def:HorizontalVerticalHomologySplitting.directSumProj}
\lean{HorizontalVerticalHomologySplitting.directSumProj}
\leanok
\uses{def:HEquivariantChainComplex.balancedKunnethSum, def:HEquivariantChainComplex.balancedKunnethSummand, def:HorizontalVerticalHomologySplitting.tannerHEq, def:HorizontalVerticalHomologySplitting.cycleHEq}

For each integer $p \in \mathbb{Z}$, the \emph{canonical projection onto the $p$-th summand} is the $\mathbb{F}_2$-linear map
\[
  \pi_p \;:\; \bigoplus_{q} H_q\!\bigl(C(X,\Lambda)\bigr) \otimes_H H_{1-q}\!\bigl(C_\ell\bigr)
  \;\longrightarrow\; H_p\!\bigl(C(X,\Lambda)\bigr) \otimes_H H_{1-p}\!\bigl(C_\ell\bigr),
\]
defined as $\pi_p := \operatorname{DirectSum.component}_{\mathbb{F}_2}(-)(p)$.
\end{definition}

\begin{definition}[Direct Sum Inclusion from $p$-th Summand]
\label{def:HorizontalVerticalHomologySplitting.directSumInc}
\lean{HorizontalVerticalHomologySplitting.directSumInc}
\leanok
\uses{def:HEquivariantChainComplex.balancedKunnethSum, def:HEquivariantChainComplex.balancedKunnethSummand, def:HorizontalVerticalHomologySplitting.tannerHEq, def:HorizontalVerticalHomologySplitting.cycleHEq}

For each integer $p \in \mathbb{Z}$, the \emph{canonical inclusion of the $p$-th summand} into the direct sum is the $\mathbb{F}_2$-linear map
\[
  \iota_p \;:\; H_p\!\bigl(C(X,\Lambda)\bigr) \otimes_H H_{1-p}\!\bigl(C_\ell\bigr)
  \;\longrightarrow\; \bigoplus_{q} H_q\!\bigl(C(X,\Lambda)\bigr) \otimes_H H_{1-q}\!\bigl(C_\ell\bigr),
\]
defined as $\iota_p := \operatorname{DirectSum.lof}_{\mathbb{F}_2}(-)(p)$.
\end{definition}

\begin{definition}[Horizontal Projection $p^h$]
\label{def:HorizontalVerticalHomologySplitting.projH}
\lean{HorizontalVerticalHomologySplitting.projH}
\leanok
\uses{def:HorizontalVerticalHomologySplitting.directSumProj, def:HorizontalVerticalHomologySplitting.kunnethIso, def:HorizontalVerticalHomologySplitting.H1, def:HorizontalVerticalHomologySplitting.H1h}

The \emph{horizontal projection} $p^h : H_1 \to H_1^h$ is the $\mathbb{F}_2$-linear map defined as the composition of the K\"{u}nneth isomorphism with the canonical projection onto the $p = 1$ summand:
\[
  p^h \;:=\; \pi_1 \circ \phi,
\]
where $\phi : H_1 \xrightarrow{\sim} \bigoplus_p H_p \otimes_H H_{1-p}$ is the K\"{u}nneth isomorphism and $\pi_1$ is the projection onto the $p = 1$ component.
\end{definition}

\begin{definition}[Vertical Projection $p^v$]
\label{def:HorizontalVerticalHomologySplitting.projV}
\lean{HorizontalVerticalHomologySplitting.projV}
\leanok
\uses{def:HorizontalVerticalHomologySplitting.directSumProj, def:HorizontalVerticalHomologySplitting.kunnethIso, def:HorizontalVerticalHomologySplitting.H1, def:HorizontalVerticalHomologySplitting.H1v}

The \emph{vertical projection} $p^v : H_1 \to H_1^v$ is the $\mathbb{F}_2$-linear map defined as the composition of the K\"{u}nneth isomorphism with the canonical projection onto the $p = 0$ summand:
\[
  p^v \;:=\; \pi_0 \circ \phi,
\]
where $\phi : H_1 \xrightarrow{\sim} \bigoplus_p H_p \otimes_H H_{1-p}$ is the K\"{u}nneth isomorphism and $\pi_0$ is the projection onto the $p = 0$ component.
\end{definition}

\begin{definition}[Horizontal Homology Class]
\label{def:HorizontalVerticalHomologySplitting.IsHorizontal}
\lean{HorizontalVerticalHomologySplitting.IsHorizontal}
\leanok
\uses{def:HorizontalVerticalHomologySplitting.projV, def:HorizontalVerticalHomologySplitting.H1}

A homology class $x \in H_1$ is \emph{horizontal} if its vertical projection vanishes:
\[
  \operatorname{IsHorizontal}(x) \;:\Leftrightarrow\; p^v(x) = 0.
\]
Equivalently, under the K\"{u}nneth isomorphism, $x$ lies entirely in the horizontal summand $H_1^h \oplus \{0\}$.
\end{definition}

\begin{definition}[Vertical Homology Class]
\label{def:HorizontalVerticalHomologySplitting.IsVertical}
\lean{HorizontalVerticalHomologySplitting.IsVertical}
\leanok
\uses{def:HorizontalVerticalHomologySplitting.projH, def:HorizontalVerticalHomologySplitting.H1}

A homology class $x \in H_1$ is \emph{vertical} if its horizontal projection vanishes:
\[
  \operatorname{IsVertical}(x) \;:\Leftrightarrow\; p^h(x) = 0.
\]
Equivalently, under the K\"{u}nneth isomorphism, $x$ lies entirely in the vertical summand $\{0\} \oplus H_1^v$.
\end{definition}

\begin{definition}[Horizontal Cohomology $H^1_h$]
\label{def:HorizontalVerticalHomologySplitting.coH1h}
\lean{HorizontalVerticalHomologySplitting.coH1h}
\leanok
\uses{def:HEquivariantChainComplex.balancedKunnethSummand, def:HorizontalVerticalHomologySplitting.tannerHEq, def:HorizontalVerticalHomologySplitting.cycleHEq}

The \emph{horizontal cohomology} is the balanced K\"{u}nneth summand at indices $p = 1$, $q = 0$:
\[
  H^1_h \;:=\; \operatorname{balancedKunnethSummand}\!\bigl(\operatorname{tannerHEq}(X, \Lambda, h_\Lambda),\; \operatorname{cycleHEq}(\ell, h_{\mathrm{compat}}),\; 1,\; 1\bigr)
  \;=\; H^1\!\bigl(C(X, \Lambda)\bigr) \otimes_H H^0\!\bigl(C_\ell\bigr).
\]
By the homology-cohomology equivalence over $\mathbb{F}_2$, this coincides with $H_1^h$.
\end{definition}

\begin{definition}[Vertical Cohomology $H^1_v$]
\label{def:HorizontalVerticalHomologySplitting.coH1v}
\lean{HorizontalVerticalHomologySplitting.coH1v}
\leanok
\uses{def:HEquivariantChainComplex.balancedKunnethSummand, def:HorizontalVerticalHomologySplitting.tannerHEq, def:HorizontalVerticalHomologySplitting.cycleHEq}

The \emph{vertical cohomology} is the balanced K\"{u}nneth summand at indices $p = 0$, $q = 1$:
\[
  H^1_v \;:=\; \operatorname{balancedKunnethSummand}\!\bigl(\operatorname{tannerHEq}(X, \Lambda, h_\Lambda),\; \operatorname{cycleHEq}(\ell, h_{\mathrm{compat}}),\; 1,\; 0\bigr)
  \;=\; H^0\!\bigl(C(X, \Lambda)\bigr) \otimes_H H^1\!\bigl(C_\ell\bigr).
\]
By the homology-cohomology equivalence over $\mathbb{F}_2$, this coincides with $H_1^v$.
\end{definition}

\begin{definition}[Cohomology Horizontal Projection $p_h$]
\label{def:HorizontalVerticalHomologySplitting.coprojH}
\lean{HorizontalVerticalHomologySplitting.coprojH}
\leanok
\uses{def:HorizontalVerticalHomologySplitting.projH, def:HorizontalVerticalHomologySplitting.coH1h}

The \emph{cohomology horizontal projection} $p_h : H^1 \to H^1_h$ is defined as
\[
  p_h \;:=\; p^h,
\]
identifying $\operatorname{coprojH}$ with $\operatorname{projH}$ via the homology-cohomology equivalence over $\mathbb{F}_2$.
\end{definition}

\begin{definition}[Cohomology Vertical Projection $p_v$]
\label{def:HorizontalVerticalHomologySplitting.coprojV}
\lean{HorizontalVerticalHomologySplitting.coprojV}
\leanok
\uses{def:HorizontalVerticalHomologySplitting.projV, def:HorizontalVerticalHomologySplitting.coH1v}

The \emph{cohomology vertical projection} $p_v : H^1 \to H^1_v$ is defined as
\[
  p_v \;:=\; p^v,
\]
identifying $\operatorname{coprojV}$ with $\operatorname{projV}$ via the homology-cohomology equivalence over $\mathbb{F}_2$.
\end{definition}

\begin{lemma}[Horizontal and Vertical Projections are Jointly Surjective]
\label{lem:HorizontalVerticalHomologySplitting.projH_projV_jointly_surjective}
\lean{HorizontalVerticalHomologySplitting.projH_projV_jointly_surjective}
\leanok
\uses{def:HorizontalVerticalHomologySplitting.projH, def:HorizontalVerticalHomologySplitting.projV, def:HorizontalVerticalHomologySplitting.kunnethIso, def:HorizontalVerticalHomologySplitting.directSumInc, def:HorizontalVerticalHomologySplitting.H1, def:HorizontalVerticalHomologySplitting.H1h, def:HorizontalVerticalHomologySplitting.H1v}

For every pair of elements $y_h \in H_1^h$ and $y_v \in H_1^v$, there exists $x \in H_1$ such that
\[
  p^h(x) = y_h \quad \text{and} \quad p^v(x) = y_v.
\]
\end{lemma}

\begin{proof}
\leanok
\uses{def:HorizontalVerticalHomologySplitting.kunnethIso, def:HorizontalVerticalHomologySplitting.directSumInc, def:HorizontalVerticalHomologySplitting.directSumProj}

Let $\phi : H_1 \xrightarrow{\sim} \bigoplus_p H_p \otimes_H H_{1-p}$ be the K\"{u}nneth isomorphism. We use the inverse $\phi^{-1}$.

Construct the element
\[
  z \;:=\; \iota_1(y_h) + \iota_0(y_v) \;\in\; \bigoplus_p H_p \otimes_H H_{1-p},
\]
where $\iota_1$ and $\iota_0$ are the canonical inclusions of the $p = 1$ and $p = 0$ summands respectively. Set $x := \phi^{-1}(z)$.

We verify both projections. For the horizontal projection, simplifying using the definitions of $\operatorname{projH}$, $\pi_1$, $\iota_1$, $\iota_0$, and the fact that $\phi \circ \phi^{-1} = \operatorname{id}$:
\begin{align*}
  p^h(x) &= \pi_1\!\bigl(\phi(\phi^{-1}(z))\bigr) = \pi_1(z) \\
  &= \pi_1\!\bigl(\iota_1(y_h) + \iota_0(y_v)\bigr) \\
  &= \pi_1(\iota_1(y_h)) + \pi_1(\iota_0(y_v)).
\end{align*}
By the defining properties of direct sum components, $\pi_1(\iota_1(y_h)) = y_h$ and $\pi_1(\iota_0(y_v)) = 0$ since $0 \neq 1$ (verified by integer arithmetic). Hence $p^h(x) = y_h$.

For the vertical projection, by the same reasoning with $\pi_0$:
\begin{align*}
  p^v(x) &= \pi_0(z) = \pi_0(\iota_1(y_h)) + \pi_0(\iota_0(y_v)) = 0 + y_v = y_v,
\end{align*}
since $1 \neq 0$ (verified by integer arithmetic) implies $\pi_0(\iota_1(y_h)) = 0$, and $\pi_0(\iota_0(y_v)) = y_v$. Thus $p^v(x) = y_v$.
\end{proof}

\begin{lemma}[Horizontal Homology is Nonempty]
\label{lem:HorizontalVerticalHomologySplitting.H1h_nonempty}
\lean{HorizontalVerticalHomologySplitting.H1h_nonempty}
\leanok
\uses{def:HorizontalVerticalHomologySplitting.H1h}

The type $H_1^h$ is nonempty; it contains the zero element.
\end{lemma}

\begin{proof}
\leanok
\uses{def:HorizontalVerticalHomologySplitting.H1h}

The zero element $0 \in H_1^h$ witnesses nonemptiness.
\end{proof}

\begin{lemma}[Vertical Homology is Nonempty]
\label{lem:HorizontalVerticalHomologySplitting.H1v_nonempty}
\lean{HorizontalVerticalHomologySplitting.H1v_nonempty}
\leanok
\uses{def:HorizontalVerticalHomologySplitting.H1v}

The type $H_1^v$ is nonempty; it contains the zero element.
\end{lemma}

\begin{proof}
\leanok
\uses{def:HorizontalVerticalHomologySplitting.H1v}

The zero element $0 \in H_1^v$ witnesses nonemptiness.
\end{proof}

\begin{lemma}[Degree-1 Homology is Nonempty]
\label{lem:HorizontalVerticalHomologySplitting.H1_nonempty}
\lean{HorizontalVerticalHomologySplitting.H1_nonempty}
\leanok
\uses{def:HorizontalVerticalHomologySplitting.H1}

The type $H_1$ is nonempty; it contains the zero element.
\end{lemma}

\begin{proof}
\leanok
\uses{def:HorizontalVerticalHomologySplitting.H1}

The zero element $0 \in H_1$ witnesses nonemptiness.
\end{proof}

%--- Def_28: IotaPiMaps ---
\chapter{Def 28: Iota and Pi Maps}

\begin{definition}[Quotient Local View]
\label{def:IotaPiMaps.quotientLocalView}
\lean{IotaPiMaps.quotientLocalView}
\leanok
\uses{def:GraphWithGroupAction, def:GraphWithGroupAction.IsInvariantLabeling, def:GraphWithGroupAction.quotientEdges, def:GraphWithGroupAction.quotientVertices, def:GraphWithGroupAction.quotientLabeling}

Let $X$ be a finite graph with a free right $H$-action, $\Lambda$ a cell labeling with $s$ local
code bits, and $h_\Lambda$ a proof that $\Lambda$ is $H$-invariant. For a vertex orbit
$V \in X_0/H$, the \emph{quotient local view} is the $\mathbb{F}_2$-linear map
\[
  \operatorname{quotientLocalView}(V) : (X_1/H \to \mathbb{F}_2) \to (\operatorname{Fin}(s) \to \mathbb{F}_2)
\]
defined by $c \mapsto \bigl(i \mapsto c\bigl((\Lambda^{X/H}_V)^{-1}(i)\bigr)\bigr)$, where
$\Lambda^{X/H}_V$ is the quotient labeling bijection at $V$.

Linearity follows: for $c_1, c_2 : X_1/H \to \mathbb{F}_2$ and $r \in \mathbb{F}_2$,
\[
  \operatorname{quotientLocalView}(V)(c_1 + c_2) = \operatorname{quotientLocalView}(V)(c_1) + \operatorname{quotientLocalView}(V)(c_2),
\]
\[
  \operatorname{quotientLocalView}(V)(r \cdot c) = r \cdot \operatorname{quotientLocalView}(V)(c).
\]
Both properties follow by extensionality and simplification of the pointwise operations.
\end{definition}

\begin{definition}[Quotient Tanner Differential]
\label{def:IotaPiMaps.quotientTannerDifferential}
\lean{IotaPiMaps.quotientTannerDifferential}
\leanok
\uses{def:GraphWithGroupAction, def:GraphWithGroupAction.IsInvariantLabeling, def:GraphWithGroupAction.quotientEdges, def:GraphWithGroupAction.quotientVertices, def:IotaPiMaps.quotientLocalView}

The \emph{quotient Tanner differential} is the $\mathbb{F}_2$-linear map
\[
  \partial : (X_1/H \to \mathbb{F}_2) \to (X_0/H \to \operatorname{Fin}(s) \to \mathbb{F}_2)
\]
defined by $c \mapsto (V \mapsto \operatorname{quotientLocalView}(V)(c))$.

Linearity: for $c_1, c_2 : X_1/H \to \mathbb{F}_2$ and $r \in \mathbb{F}_2$, the result
follows by extensionality over vertex orbits $V$ and local indices $i$, together with
simplification of the pointwise addition and scalar multiplication.
\end{definition}

\begin{definition}[Quotient Tanner Code Complex]
\label{def:IotaPiMaps.quotientTannerComplex}
\lean{IotaPiMaps.quotientTannerComplex}
\leanok
\uses{def:GraphWithGroupAction, def:GraphWithGroupAction.IsInvariantLabeling, def:GraphWithGroupAction.quotientEdges, def:GraphWithGroupAction.quotientVertices, def:IotaPiMaps.quotientTannerDifferential}

The \emph{quotient Tanner code complex} $C(X/H, L)$ is the chain complex over $\mathbb{F}_2$
with chain spaces
\[
  C_1(X/H) = (X_1/H \to \mathbb{F}_2), \qquad C_0(X/H) = (X_0/H \to \operatorname{Fin}(s) \to \mathbb{F}_2),
\]
and $C_i = 0$ for $i \geq 2$ or $i < 0$. The differential is
\[
  d_{1,0} = \partial : C_1(X/H) \to C_0(X/H)
\]
(the quotient Tanner differential), and all other differentials are zero.

The chain complex condition $d \circ d = 0$ is verified:
\begin{itemize}
  \item For the pair $(i,j,k) = (1,0,k)$: we have $d_{1,0} \circ 0 = 0$ by the zero
    postcomposition rule.
  \item For $i = 0$: the outgoing differential is zero, so the composition vanishes.
  \item For $i \geq 2$ or $i < 0$: the incoming differential is zero, so the composition
    vanishes by the zero precomposition rule.
  \item All cases with non-adjacent indices satisfy the shape condition vacuously (the shape
    predicate fails, verified by positivity or simplification).
\end{itemize}

The homology in degree 1 is denoted $H_1(C(X/H, L)) := \ker(d_{1,0}) / \operatorname{im}(0)
= \ker(\partial)$.
\end{definition}

\begin{definition}[Quotient Tanner Homology]
\label{def:IotaPiMaps.quotientTannerHomology}
\lean{IotaPiMaps.quotientTannerHomology}
\leanok
\uses{def:GraphWithGroupAction, def:GraphWithGroupAction.IsInvariantLabeling, def:IotaPiMaps.quotientTannerComplex}

The \emph{quotient Tanner homology} is the degree-1 homology of the quotient Tanner code complex:
\[
  H_1(C(X/H, L)) := \bigl(\operatorname{quotientTannerComplex}(X, \Lambda, h_\Lambda)\bigr).H_1.
\]
\end{definition}

\begin{theorem}[Axiom: Iota Map]
\label{thm:IotaPiMaps.iotaMap}
\lean{IotaPiMaps.iotaMap}
\leanok
\uses{def:GraphWithGroupAction, def:GraphWithGroupAction.IsInvariantLabeling, def:BalancedProductTannerCycleCode.CycleCompatibleAction, def:HorizontalVerticalHomologySplitting.H1h, def:IotaPiMaps.quotientTannerHomology}

\textbf{\color{red}[UNPROVEN AXIOM]}

There exists an $\mathbb{F}_2$-linear map
\[
  \iota : H_1(C(X/H, L)) \to H_1^h(C(X,L) \otimes_H C(C_\ell)),
\]
where $H_1^h = H_1^h(X, \Lambda, \ell, h_\Lambda, h_{\mathrm{compat}})$ is the horizontal
homology of the balanced product Tanner cycle code from Definition 27.

The map $\iota$ is defined on homology classes as follows: given
$[\sum_{\mathcal{E}} a_{\mathcal{E}} \mathcal{E}] \in H_1(C(X/H,L))$
(where $\mathcal{E} = eH$ ranges over edge orbits), set
\[
  \iota\!\left[\sum_{\mathcal{E}} a_{\mathcal{E}} \mathcal{E}\right]
  = \left[\left(\sum_{\mathcal{E}} a_{\mathcal{E}} \sum_{e \in \mathcal{E}} e \otimes y_0,\; 0\right)\right],
\]
where $y_0$ is a fixed $0$-cell of the cycle graph $C_\ell$.

Well-definedness: the orbit indicator $\sum_{e \in \mathcal{E}} \delta_e$ is $H$-invariant in
$C_1(X,L)$, so tensoring with $\delta_{y_0} \in C_0(C_\ell)$ gives a well-defined element
of $C_1(X,L) \otimes_H C_0(C_\ell)$. The map sends cycles to cycles (the quotient differential
is compatible with the balanced product differential) and boundaries to boundaries (well-defined on homology).

\textit{Note: This is stated as an axiom because the full proof was not completed in the formalization.}
\end{theorem}

\begin{theorem}[Axiom: Pi Map]
\label{thm:IotaPiMaps.piMap}
\lean{IotaPiMaps.piMap}
\leanok
\uses{def:GraphWithGroupAction, def:GraphWithGroupAction.IsInvariantLabeling, def:BalancedProductTannerCycleCode.CycleCompatibleAction, def:HorizontalVerticalHomologySplitting.H1h, def:IotaPiMaps.quotientTannerHomology}

\textbf{\color{red}[UNPROVEN AXIOM]}

There exists an $\mathbb{F}_2$-linear map
\[
  \pi : H_1^h(C(X,L) \otimes_H C(C_\ell)) \to H_1(C(X/H, L)).
\]
The map $\pi$ projects from horizontal homology to quotient Tanner code homology by collapsing
orbits. By the K\"{u}nneth decomposition, elements of
$H_1^h = H_1(C(X,L)) \otimes_H H_0(C_\ell)$ are represented by cycles of the form
$(\sum_e a_e\, e \otimes y_0, 0)$ with coefficients $a_e$ constant on $H$-orbits. The map sends
this class to
\[
  \pi\!\left[\left(\sum_e a_e\, e \otimes y_0, 0\right)\right] = \left[\sum_e a_e\, eH\right]
  \in H_1(C(X/H,L)).
\]
Well-definedness: since $a_e$ is constant on orbits, each orbit $\mathcal{E} = eH$ contributes
$a_{\mathcal{E}} \cdot \mathcal{E}$, and the quotient differential is compatible with the balanced
product differential.

\textit{Note: This is stated as an axiom because the full proof was not completed in the formalization.}
\end{theorem}

\begin{theorem}[Axiom: $\pi \circ \iota = \mathrm{id}$]
\label{thm:IotaPiMaps.piMap_comp_iotaMap}
\lean{IotaPiMaps.piMap_comp_iotaMap}
\leanok
\uses{def:GraphWithGroupAction, def:GraphWithGroupAction.IsInvariantLabeling, def:BalancedProductTannerCycleCode.CycleCompatibleAction, def:HorizontalVerticalHomologySplitting.H1h, def:IotaPiMaps.quotientTannerHomology}

\textbf{\color{red}[UNPROVEN AXIOM]}

Assuming $\ell \equiv 1 \pmod{2}$ and $|H|$ is odd, we have
\[
  \pi \circ \iota = \mathrm{id}_{H_1(C(X/H,L))}.
\]
Each orbit $\mathcal{E} = eH$ has $|H| = \ell$ elements. Applying $\iota$ to $[\mathcal{E}]$
gives $[\sum_{e \in \mathcal{E}} e \otimes y_0]$, and then $\pi$ maps each $e \otimes y_0$ back
to $\mathcal{E}$, yielding $\ell \cdot [\mathcal{E}]$. Since $\ell$ is odd, $\ell \equiv 1
\pmod{2}$, so $\ell \cdot [\mathcal{E}] = [\mathcal{E}]$ in $\mathbb{F}_2$.

\textit{Note: This is stated as an axiom because the full proof was not completed in the formalization.}
\end{theorem}

\begin{theorem}[Axiom: $\iota \circ \pi = \mathrm{id}$]
\label{thm:IotaPiMaps.iotaMap_comp_piMap}
\lean{IotaPiMaps.iotaMap_comp_piMap}
\leanok
\uses{def:GraphWithGroupAction, def:GraphWithGroupAction.IsInvariantLabeling, def:BalancedProductTannerCycleCode.CycleCompatibleAction, def:HorizontalVerticalHomologySplitting.H1h, def:IotaPiMaps.quotientTannerHomology}

\textbf{\color{red}[UNPROVEN AXIOM]}

Assuming $\ell \equiv 1 \pmod{2}$ and $|H|$ is odd, we have
\[
  \iota \circ \pi = \mathrm{id}_{H_1^h}.
\]
For a horizontal class $[(\sum_e a_e\, e \otimes y_0, 0)]$ with $a_e$ constant on orbits,
$\pi$ sends it to $[\sum_e a_e\, eH]$, and $\iota$ lifts back to
$[(\sum_e a_e \sum_{e' \in eH} e' \otimes y_0, 0)]$. Since each orbit of size $\ell$
contributes $\ell$ copies and $\ell \equiv 1 \pmod{2}$, this equals the original class.

\textit{Note: This is stated as an axiom because the full proof was not completed in the formalization.}
\end{theorem}

\begin{definition}[Witness Graph]
\label{def:IotaPiMaps.witnessGraph}
\lean{IotaPiMaps.witnessGraph}
\leanok
\uses{def:GraphWithGroupAction}

The \emph{witness graph} is a trivial graph with a single vertex and a single edge, equipped
with the trivial $\mathrm{Unit}$-action: the group $H = \mathrm{Unit}$ acts on both vertices
and edges as the identity. The cell sets are $\mathrm{cells}(n) = \mathrm{Unit}$ for all $n$,
with empty boundary maps (satisfying $\partial \circ \partial = 0$ vacuously). Both the vertex
and edge actions are free (since the only element of $\mathrm{Unit}$ fixing a cell is the
identity). This provides a non-trivial graph with at least one vertex and one edge on which all
required hypotheses hold.
\end{definition}

\begin{definition}[Witness Labeling]
\label{def:IotaPiMaps.witnessLabeling}
\lean{IotaPiMaps.witnessLabeling}
\leanok
\uses{def:GraphWithGroupAction, def:GraphWithGroupAction.CellLabeling, def:IotaPiMaps.witnessGraph}

The \emph{witness labeling} is a cell labeling $\Lambda : \operatorname{witnessGraph}.\operatorname{CellLabeling}(0)$
with $s = 0$ local code bits. Since $s = 0$, the labeling map at any vertex is a bijection
from the empty set of incident edges to $\operatorname{Fin}(0) = \emptyset$. It is defined by:
for each vertex, the labeling bijection sends any element $e$ with $e \in \emptyset$ to its
absurdity (by \texttt{Finset.notMem\_empty}), and the inverse sends any element of
$\operatorname{Fin}(0)$ to its absurdity (by \texttt{Fin.elim0}).
\end{definition}

\begin{lemma}[Witness Labeling is Invariant]
\label{lem:IotaPiMaps.witnessLabeling_invariant}
\lean{IotaPiMaps.witnessLabeling_invariant}
\leanok
\uses{def:GraphWithGroupAction.IsInvariantLabeling, def:IotaPiMaps.witnessLabeling}

The witness labeling is $H$-invariant:
\[
  \operatorname{GraphWithGroupAction.IsInvariantLabeling}(\operatorname{witnessLabeling}).
\]
\end{lemma}

\begin{proof}
\leanok
\uses{def:IotaPiMaps.witnessLabeling}

For any vertex $v$, group element $h$, and incident edge $e$, we need $\Lambda(v)(e) = \Lambda(h \cdot v)(h \cdot e)$. Since the labeling has $s = 0$ and the edge $e$ must satisfy $e \in \emptyset$ (the empty boundary), the hypothesis $e.2 \in \operatorname{Finset.empty}$ is absurd, so the conclusion holds vacuously by \texttt{Finset.notMem\_empty}.
\end{proof}

\begin{definition}[Witness Unit Mul Action]
\label{def:IotaPiMaps.witnessUnitMulAction}
\lean{IotaPiMaps.witnessUnitMulAction}
\leanok
\uses{def:GraphWithGroupAction}

The \emph{witness $\mathrm{Unit}$ action on $\operatorname{Fin}(3)$} is the trivial $\mathrm{MulAction}$
where the unique group element $() \in \mathrm{Unit}$ acts on every $i \in \operatorname{Fin}(3)$
as the identity: $() \cdot i = i$. The axioms $\mathtt{one\_smul}$ and $\mathtt{mul\_smul}$
hold by reflexivity.
\end{definition}

\begin{lemma}[Witness Cycle Compatibility]
\label{lem:IotaPiMaps.witnessCycleCompat}
\lean{IotaPiMaps.witnessCycleCompat}
\leanok
\uses{def:BalancedProductTannerCycleCode.CycleCompatibleAction, def:IotaPiMaps.witnessUnitMulAction}

With the witness $\mathrm{Unit}$-action on $\operatorname{Fin}(3)$, the cycle-compatible action
condition holds for $H = \mathrm{Unit}$ and $\ell = 3$:
\[
  \operatorname{BalancedProductTannerCycleCode.CycleCompatibleAction}(\mathrm{Unit}, 3).
\]
\end{lemma}

\begin{proof}
\leanok
\uses{def:IotaPiMaps.witnessUnitMulAction}

For any $h : \mathrm{Unit}$ and $i : \operatorname{Fin}(3)$, the condition requires a compatibility
between the group action and the cycle graph structure. By the definition of
$\operatorname{witnessUnitMulAction}$, the action is trivial: $h \cdot i = i$.
The required identity holds after unfolding the definition and simplifying.
\end{proof}

\begin{lemma}[Iota Map Premises Satisfiable]
\label{lem:IotaPiMaps.iotaMap_satisfiable}
\lean{IotaPiMaps.iotaMap_satisfiable}
\leanok
\uses{def:GraphWithGroupAction, def:GraphWithGroupAction.IsInvariantLabeling, def:GraphWithGroupAction.CellLabeling, def:BalancedProductTannerCycleCode.CycleCompatibleAction, def:IotaPiMaps.witnessGraph, def:IotaPiMaps.witnessLabeling, def:IotaPiMaps.witnessUnitMulAction}

The premises of the iota map are satisfiable: there exist a group $H$, a graph with $H$-action
$X$, a cell labeling $\Lambda$, a natural number $\ell$, a group action of $H$ on
$\operatorname{Fin}(\ell)$, invariance $h_\Lambda$, and cycle-compatibility $h_{\mathrm{compat}}$,
such that the underlying graph has at least one vertex and one edge.
\end{lemma}

\begin{proof}
\leanok
\uses{def:IotaPiMaps.witnessGraph, def:IotaPiMaps.witnessLabeling, def:IotaPiMaps.witnessUnitMulAction, lem:IotaPiMaps.witnessLabeling_invariant, lem:IotaPiMaps.witnessCycleCompat}

We take $H = \mathrm{Unit}$, $X = \operatorname{witnessGraph}$, $s = 0$,
$\Lambda = \operatorname{witnessLabeling}$, $\ell = 3$, the action from
$\operatorname{witnessUnitMulAction}$, invariance from $\operatorname{witnessLabeling_invariant}$,
and compatibility from $\operatorname{witnessCycleCompat}$. The witness graph has
$\operatorname{cells}(0) = \mathrm{Unit}$ and $\operatorname{cells}(1) = \mathrm{Unit}$,
so $(() , \operatorname{trivial})$ witnesses both a vertex and an edge.
\end{proof}

\begin{lemma}[$\pi \circ \iota$ Premises Satisfiable]
\label{lem:IotaPiMaps.piMap_comp_iotaMap_satisfiable}
\lean{IotaPiMaps.piMap_comp_iotaMap_satisfiable}
\leanok
\uses{def:GraphWithGroupAction, def:GraphWithGroupAction.IsInvariantLabeling, def:BalancedProductTannerCycleCode.CycleCompatibleAction, def:IotaPiMaps.witnessGraph, def:IotaPiMaps.witnessLabeling, def:IotaPiMaps.witnessUnitMulAction}

The premises of $\pi \circ \iota = \mathrm{id}$ (including $\ell \geq 3$, $\ell$ odd, and $|H|$
odd) are satisfiable with a non-trivial graph.
\end{lemma}

\begin{proof}
\leanok
\uses{def:IotaPiMaps.witnessGraph, def:IotaPiMaps.witnessLabeling, def:IotaPiMaps.witnessUnitMulAction, lem:IotaPiMaps.witnessLabeling_invariant, lem:IotaPiMaps.witnessCycleCompat}

We take the same witnesses as in $\operatorname{iotaMap_satisfiable}$, additionally verifying:
$\ell = 3 \geq 3$ (by reflexivity of $\leq$); $3 \bmod 2 = 1$ (by reflexivity); and
$\operatorname{Fintype.card}(\mathrm{Unit}) = 1$ is odd, witnessed by $(0, \operatorname{rfl})$
since $1 = 2 \cdot 0 + 1$.
\end{proof}

\begin{lemma}[$\iota \circ \pi$ Premises Satisfiable]
\label{lem:IotaPiMaps.iotaMap_comp_piMap_satisfiable}
\lean{IotaPiMaps.iotaMap_comp_piMap_satisfiable}
\leanok
\uses{def:GraphWithGroupAction, def:GraphWithGroupAction.IsInvariantLabeling, def:BalancedProductTannerCycleCode.CycleCompatibleAction, def:IotaPiMaps.witnessGraph, def:IotaPiMaps.witnessLabeling, def:IotaPiMaps.witnessUnitMulAction}

The premises of $\iota \circ \pi = \mathrm{id}$ (including $\ell \geq 3$, $\ell$ odd, and $|H|$
odd) are satisfiable with a non-trivial graph.
\end{lemma}

\begin{proof}
\leanok
\uses{def:IotaPiMaps.witnessGraph, def:IotaPiMaps.witnessLabeling, def:IotaPiMaps.witnessUnitMulAction, lem:IotaPiMaps.witnessLabeling_invariant, lem:IotaPiMaps.witnessCycleCompat}

We use the same witnesses as in $\operatorname{piMap_comp_iotaMap_satisfiable}$. The oddness
of $|H| = |\mathrm{Unit}| = 1$ is witnessed by the pair $(0, \operatorname{rfl})$.
\end{proof}

\begin{theorem}[$\iota$ is Injective]
\label{thm:IotaPiMaps.iotaMap_injective}
\lean{IotaPiMaps.iotaMap_injective}
\leanok
\uses{def:GraphWithGroupAction, def:GraphWithGroupAction.IsInvariantLabeling, def:BalancedProductTannerCycleCode.CycleCompatibleAction, def:HorizontalVerticalHomologySplitting.H1h, def:IotaPiMaps.quotientTannerHomology}

Assuming $\ell \equiv 1 \pmod{2}$ and $|H|$ is odd, the map $\iota$ is injective:
\[
  \operatorname{Function.Injective}(\iota).
\]
\end{theorem}

\begin{proof}
\leanok
\uses{def:GraphWithGroupAction, def:GraphWithGroupAction.IsInvariantLabeling, def:BalancedProductTannerCycleCode.CycleCompatibleAction, def:HorizontalVerticalHomologySplitting.H1h, def:IotaPiMaps.quotientTannerHomology}

We apply $\pi \circ \iota = \mathrm{id}$ (from \texttt{piMap\_comp\_iotaMap}). Let $a, b$ be such
that $\iota(a) = \iota(b)$. We have
\[
  \pi(\iota(a)) = \pi(\iota(b)).
\]
Rewriting using $\pi \circ \iota = \mathrm{id}$, we obtain $\mathrm{id}(a) = \mathrm{id}(b)$,
i.e., $a = b$. Concretely, using \texttt{congr\_arg} with $\pi$ applied to the hypothesis
$\iota(a) = \iota(b)$, then rewriting via \texttt{LinearMap.comp\_apply} in both directions
and applying the composite identity, we get $\mathrm{id}(a) = \mathrm{id}(b)$, hence $a = b$.
\end{proof}

\begin{theorem}[$\pi$ is Surjective]
\label{thm:IotaPiMaps.piMap_surjective}
\lean{IotaPiMaps.piMap_surjective}
\leanok
\uses{def:GraphWithGroupAction, def:GraphWithGroupAction.IsInvariantLabeling, def:BalancedProductTannerCycleCode.CycleCompatibleAction, def:HorizontalVerticalHomologySplitting.H1h, def:IotaPiMaps.quotientTannerHomology}

Assuming $\ell \equiv 1 \pmod{2}$ and $|H|$ is odd, the map $\pi$ is surjective:
\[
  \operatorname{Function.Surjective}(\pi).
\]
\end{theorem}

\begin{proof}
\leanok
\uses{def:GraphWithGroupAction, def:GraphWithGroupAction.IsInvariantLabeling, def:BalancedProductTannerCycleCode.CycleCompatibleAction, def:HorizontalVerticalHomologySplitting.H1h, def:IotaPiMaps.quotientTannerHomology}

We apply $\pi \circ \iota = \mathrm{id}$. Given any $y \in H_1(C(X/H,L))$, we claim
$\iota(y)$ is a preimage of $y$ under $\pi$. By \texttt{LinearMap.ext\_iff} applied to the
identity $\pi \circ \iota = \mathrm{id}$ at $y$, we get $\pi(\iota(y)) = \mathrm{id}(y) = y$,
after simplifying via \texttt{LinearMap.comp\_apply} and \texttt{LinearMap.id\_apply}.
\end{proof}

\begin{theorem}[$\pi$ is Injective]
\label{thm:IotaPiMaps.piMap_injective}
\lean{IotaPiMaps.piMap_injective}
\leanok
\uses{def:GraphWithGroupAction, def:GraphWithGroupAction.IsInvariantLabeling, def:BalancedProductTannerCycleCode.CycleCompatibleAction, def:HorizontalVerticalHomologySplitting.H1h, def:IotaPiMaps.quotientTannerHomology}

Assuming $\ell \equiv 1 \pmod{2}$ and $|H|$ is odd, the map $\pi$ is injective:
\[
  \operatorname{Function.Injective}(\pi).
\]
\end{theorem}

\begin{proof}
\leanok
\uses{def:GraphWithGroupAction, def:GraphWithGroupAction.IsInvariantLabeling, def:BalancedProductTannerCycleCode.CycleCompatibleAction, def:HorizontalVerticalHomologySplitting.H1h, def:IotaPiMaps.quotientTannerHomology}

We apply $\iota \circ \pi = \mathrm{id}$ (from \texttt{iotaMap\_comp\_piMap}). Let $a, b$ satisfy
$\pi(a) = \pi(b)$. Applying $\iota$ to both sides via \texttt{congr\_arg}, we get
$\iota(\pi(a)) = \iota(\pi(b))$. Rewriting via \texttt{LinearMap.comp\_apply} and the identity
$\iota \circ \pi = \mathrm{id}$, we obtain $\mathrm{id}(a) = \mathrm{id}(b)$, i.e., $a = b$.
\end{proof}

\begin{theorem}[$\iota$ is Surjective]
\label{thm:IotaPiMaps.iotaMap_surjective}
\lean{IotaPiMaps.iotaMap_surjective}
\leanok
\uses{def:GraphWithGroupAction, def:GraphWithGroupAction.IsInvariantLabeling, def:BalancedProductTannerCycleCode.CycleCompatibleAction, def:HorizontalVerticalHomologySplitting.H1h, def:IotaPiMaps.quotientTannerHomology}

Assuming $\ell \equiv 1 \pmod{2}$ and $|H|$ is odd, the map $\iota$ is surjective:
\[
  \operatorname{Function.Surjective}(\iota).
\]
\end{theorem}

\begin{proof}
\leanok
\uses{def:GraphWithGroupAction, def:GraphWithGroupAction.IsInvariantLabeling, def:BalancedProductTannerCycleCode.CycleCompatibleAction, def:HorizontalVerticalHomologySplitting.H1h, def:IotaPiMaps.quotientTannerHomology}

We apply $\iota \circ \pi = \mathrm{id}$. Given any $y \in H_1^h$, we exhibit $\pi(y)$ as a
preimage under $\iota$. By \texttt{LinearMap.ext\_iff} applied to $\iota \circ \pi = \mathrm{id}$
at $y$, we get $\iota(\pi(y)) = \mathrm{id}(y) = y$, after simplifying via
\texttt{LinearMap.comp\_apply} and \texttt{LinearMap.id\_apply}.
\end{proof}

\begin{definition}[Iota Linear Equivalence]
\label{def:IotaPiMaps.iotaEquiv}
\lean{IotaPiMaps.iotaEquiv}
\leanok
\uses{def:GraphWithGroupAction, def:GraphWithGroupAction.IsInvariantLabeling, def:BalancedProductTannerCycleCode.CycleCompatibleAction, def:HorizontalVerticalHomologySplitting.H1h, def:IotaPiMaps.quotientTannerHomology}

Assuming $\ell \equiv 1 \pmod{2}$ and $|H|$ is odd, the map $\iota$ is a linear isomorphism:
\[
  \iota : H_1(C(X/H, L)) \xrightarrow{\;\sim\;} H_1^h.
\]
This equivalence $\iota^{\equiv}$ is constructed via
$\operatorname{LinearEquiv.ofBijective}(\iota, \langle h_{\mathrm{inj}}, h_{\mathrm{surj}}\rangle)$,
where $h_{\mathrm{inj}}$ is from \texttt{iotaMap\_injective} and $h_{\mathrm{surj}}$ from
\texttt{iotaMap\_surjective}.
\end{definition}

\begin{definition}[Pi Linear Equivalence]
\label{def:IotaPiMaps.piEquiv}
\lean{IotaPiMaps.piEquiv}
\leanok
\uses{def:GraphWithGroupAction, def:GraphWithGroupAction.IsInvariantLabeling, def:BalancedProductTannerCycleCode.CycleCompatibleAction, def:HorizontalVerticalHomologySplitting.H1h, def:IotaPiMaps.quotientTannerHomology}

Assuming $\ell \equiv 1 \pmod{2}$ and $|H|$ is odd, the map $\pi$ is a linear isomorphism:
\[
  \pi : H_1^h \xrightarrow{\;\sim\;} H_1(C(X/H, L)).
\]
This equivalence $\pi^{\equiv}$ is constructed via
$\operatorname{LinearEquiv.ofBijective}(\pi, \langle h_{\mathrm{inj}}, h_{\mathrm{surj}}\rangle)$,
where $h_{\mathrm{inj}}$ is from \texttt{piMap\_injective} and $h_{\mathrm{surj}}$ from
\texttt{piMap\_surjective}.
\end{definition}

\begin{lemma}[Quotient Tanner Homology is Nonempty]
\label{lem:IotaPiMaps.quotientTannerHomology_nonempty}
\lean{IotaPiMaps.quotientTannerHomology_nonempty}
\leanok
\uses{def:GraphWithGroupAction, def:GraphWithGroupAction.IsInvariantLabeling, def:IotaPiMaps.quotientTannerHomology}

The quotient Tanner homology $H_1(C(X/H, L))$ is nonempty.
\end{lemma}

\begin{proof}
\leanok
\uses{def:IotaPiMaps.quotientTannerHomology}

The zero element $0 \in H_1(C(X/H, L))$ witnesses non-emptiness: $\langle 0 \rangle$.
\end{proof}

\begin{lemma}[Iota Equivalence is Nonempty]
\label{lem:IotaPiMaps.iotaEquiv_nonempty}
\lean{IotaPiMaps.iotaEquiv_nonempty}
\leanok
\uses{def:GraphWithGroupAction, def:GraphWithGroupAction.IsInvariantLabeling, def:BalancedProductTannerCycleCode.CycleCompatibleAction, def:HorizontalVerticalHomologySplitting.H1h, def:IotaPiMaps.quotientTannerHomology, def:IotaPiMaps.iotaEquiv}

Assuming $\ell \equiv 1 \pmod{2}$ and $|H|$ is odd, the type of linear isomorphisms
$H_1(C(X/H, L)) \simeq_{\mathbb{F}_2} H_1^h$ is nonempty.
\end{lemma}

\begin{proof}
\leanok
\uses{def:IotaPiMaps.iotaEquiv}

The equivalence $\operatorname{iotaEquiv}(X, \Lambda, \ell, h_\Lambda, h_{\mathrm{compat}}, h_{\ell\mathrm{odd}}, h_{\mathrm{odd}})$
witnesses non-emptiness.
\end{proof}

%--- Thm_12: EncodingRateCircle ---
\chapter{Thm 12: Encoding Rate (Circle Case)}

\begin{theorem}[Encoding Rate (Circle Case)]
\label{thm:EncodingRateCircle.encodingRateCircle}
\lean{EncodingRateCircle.encodingRateCircle}
\leanok
\uses{def:HorizontalVerticalHomologySplitting.H1h, def:IotaPiMaps.quotientTannerHomology, def:IotaPiMaps.iotaEquiv, def:GraphWithGroupAction, def:GraphWithGroupAction.CellLabeling, def:GraphWithGroupAction.IsInvariantLabeling, def:BalancedProductTannerCycleCode.CycleCompatibleAction}

Let $H$ be a finite group with $\operatorname{DecidableEq}$ instance, let $X$ be a graph with $H$-action, let $\Lambda$ be an $H$-invariant cell labeling of $X$ of width $s$, let $\ell \in \mathbb{N}$ with $\ell \neq 0$ and $H \curvearrowright \operatorname{Fin}(\ell)$. Assume $\ell$ is odd (i.e., $\ell \equiv 1 \pmod{2}$) and $|H|$ is odd. Then the horizontal homology of the balanced product Tanner cycle code has the same $\mathbb{F}_2$-dimension as the quotient Tanner code homology:
\[
\dim_{\mathbb{F}_2} H_1^h\!\left(C(X, \Lambda) \otimes_{\mathbb{Z}_\ell} C(C_\ell)\right)
\;=\;
\dim_{\mathbb{F}_2} H_1\!\left(C(X/H, \Lambda)\right).
\]
\end{theorem}

\begin{proof}
\leanok
\uses{def:IotaPiMaps.iotaEquiv, def:HorizontalVerticalHomologySplitting.H1h, def:IotaPiMaps.quotientTannerHomology}

By the linear equivalence $\iota\colon H_1(C(X/H,\Lambda)) \xrightarrow{\;\simeq\;} H_1^h$ constructed in \texttt{IotaPiMaps.iotaEquiv}, which is a $\mathbb{F}_2$-linear isomorphism under the assumptions that $\ell$ is odd and $|H|$ is odd, linear equivalences preserve $\operatorname{finrank}$. Applying \texttt{LinearEquiv.finrank\_eq} to this isomorphism and taking the symmetric equality, we obtain the desired dimension equality.
\end{proof}

\begin{lemma}[Encoding Rate Satisfiability Witness]
\label{lem:EncodingRateCircle.encodingRateCircle_satisfiable}
\lean{EncodingRateCircle.encodingRateCircle_satisfiable}
\leanok
\uses{def:GraphWithGroupAction, def:GraphWithGroupAction.CellLabeling, def:GraphWithGroupAction.IsInvariantLabeling, def:BalancedProductTannerCycleCode.CycleCompatibleAction}

The premises of \texttt{encodingRateCircle} are satisfiable: there exist a finite group $H$, a graph with $H$-action $X$, an $H$-invariant cell labeling $\Lambda$, a positive natural number $\ell$ with an $H$-action on $\operatorname{Fin}(\ell)$, a compatible cycle action, $\ell$ odd, and $|H|$ odd, such that $\top$ holds.
\end{lemma}

\begin{proof}
\leanok
\uses{lem:IotaPiMaps.piMap_comp_iotaMap_satisfiable}

We apply \texttt{IotaPiMaps.piMap\_comp\_iotaMap\_satisfiable}, which provides a witness tuple $\langle H, \text{hG}, \text{hF}, \text{hD}, X, s, \Lambda, \ell, \text{hNZ}, \text{hMA}, h\Lambda, \text{hcompat}, -, h_{\ell\text{\_odd}}, \text{hodd}, -\rangle$. Decomposing this existential, we extract all required components and assemble the existential for the satisfiability statement, concluding with \texttt{trivial}.
\end{proof}

%--- Def_29: HorizontalSubsystemBalancedProductCode ---
\chapter{Def 29: Horizontal Subsystem Balanced Product Code}

\begin{definition}[Horizontal Künneth Inclusion]
\label{def:HorizontalSubsystemBalancedProductCode.incH}
\lean{HorizontalSubsystemBalancedProductCode.incH}
\leanok
\uses{def:HorizontalVerticalHomologySplitting.H1h, def:HorizontalVerticalHomologySplitting.directSumInc, def:HEquivariantChainComplex.balancedKunnethSum}
Let $H$ be a finite group acting on a graph-with-group-action $X$ with cell labeling $\Lambda$, and let $\ell \geq 3$ be an odd integer with a compatible cycle action. The \emph{horizontal Künneth inclusion} is the $\mathbb{F}_2$-linear map
\[
  \operatorname{incH} : H_1^h \;\longrightarrow\; \bigoplus_{p} H_p(\mathcal{C}(X,\Lambda)) \otimes_H H_{1-p}(\mathcal{C}_\ell),
\]
defined as the direct sum inclusion at index $p = 1$. Here $H_1^h$ denotes the horizontal homology summand of the balanced product complex.
\end{definition}

\begin{definition}[Vertical Künneth Inclusion]
\label{def:HorizontalSubsystemBalancedProductCode.incV}
\lean{HorizontalSubsystemBalancedProductCode.incV}
\leanok
\uses{def:HorizontalVerticalHomologySplitting.H1v, def:HorizontalVerticalHomologySplitting.directSumInc, def:HEquivariantChainComplex.balancedKunnethSum}
The \emph{vertical Künneth inclusion} is the $\mathbb{F}_2$-linear map
\[
  \operatorname{incV} : H_1^v \;\longrightarrow\; \bigoplus_{p} H_p(\mathcal{C}(X,\Lambda)) \otimes_H H_{1-p}(\mathcal{C}_\ell),
\]
defined as the direct sum inclusion at index $p = 0$. Here $H_1^v$ denotes the vertical homology summand of the balanced product complex.
\end{definition}

\begin{definition}[Horizontal Embedding into $H_1$]
\label{def:HorizontalSubsystemBalancedProductCode.embH}
\lean{HorizontalSubsystemBalancedProductCode.embH}
\leanok
\uses{def:HorizontalSubsystemBalancedProductCode.incH, def:HorizontalVerticalHomologySplitting.kunnethIso, def:HorizontalVerticalHomologySplitting.H1, def:HorizontalVerticalHomologySplitting.H1h}
The \emph{horizontal embedding} is the $\mathbb{F}_2$-linear map
\[
  \operatorname{embH} : H_1^h \;\longrightarrow\; H_1,
\]
defined as the composition
\[
  \operatorname{embH} = \kappa^{-1} \circ \operatorname{incH},
\]
where $\kappa : H_1 \xrightarrow{\sim} \bigoplus_p H_p(\mathcal{C}(X,\Lambda)) \otimes_H H_{1-p}(\mathcal{C}_\ell)$ is the Künneth isomorphism (Def.\ 27).
\end{definition}

\begin{definition}[Vertical Embedding into $H_1$]
\label{def:HorizontalSubsystemBalancedProductCode.embV}
\lean{HorizontalSubsystemBalancedProductCode.embV}
\leanok
\uses{def:HorizontalSubsystemBalancedProductCode.incV, def:HorizontalVerticalHomologySplitting.kunnethIso, def:HorizontalVerticalHomologySplitting.H1, def:HorizontalVerticalHomologySplitting.H1v}
The \emph{vertical embedding} is the $\mathbb{F}_2$-linear map
\[
  \operatorname{embV} : H_1^v \;\longrightarrow\; H_1,
\]
defined as the composition
\[
  \operatorname{embV} = \kappa^{-1} \circ \operatorname{incV},
\]
where $\kappa$ is the Künneth isomorphism.
\end{definition}

\begin{definition}[Logical Submodule $H_1^L$]
\label{def:HorizontalSubsystemBalancedProductCode.HL}
\lean{HorizontalSubsystemBalancedProductCode.HL}
\leanok
\uses{def:HorizontalSubsystemBalancedProductCode.embH}
The \emph{logical submodule} $H_1^L \subseteq H_1$ is defined as the image of the horizontal embedding:
\[
  H_1^L := \operatorname{im}(\operatorname{embH}) = \operatorname{embH}(H_1^h) \subseteq H_1.
\]
Logical $Z$-operators of the subsystem code correspond to nontrivial elements of the horizontal homology $H_1^h$.
\end{definition}

\begin{definition}[Gauge Submodule $H_1^G$]
\label{def:HorizontalSubsystemBalancedProductCode.HG}
\lean{HorizontalSubsystemBalancedProductCode.HG}
\leanok
\uses{def:HorizontalSubsystemBalancedProductCode.embV}
The \emph{gauge submodule} $H_1^G \subseteq H_1$ is defined as the image of the vertical embedding:
\[
  H_1^G := \operatorname{im}(\operatorname{embV}) = \operatorname{embV}(H_1^v) \subseteq H_1.
\]
Gauge $Z$-operators correspond to elements of the vertical homology $H_1^v$.
\end{definition}

\begin{definition}[Number of Logical Qubits]
\label{def:HorizontalSubsystemBalancedProductCode.logicalQubits}
\lean{HorizontalSubsystemBalancedProductCode.logicalQubits}
\leanok
\uses{def:HorizontalSubsystemBalancedProductCode.HL}
The \emph{number of logical qubits} of the horizontal subsystem balanced product code is the $\mathbb{F}_2$-rank of the logical submodule:
\[
  k := \dim_{\mathbb{F}_2} H_1^L.
\]
\end{definition}

\begin{theorem}[Complementarity of Logical and Gauge Submodules]
\label{thm:HorizontalSubsystemBalancedProductCode.hcompl}
\lean{HorizontalSubsystemBalancedProductCode.hcompl}
\leanok
\uses{def:HorizontalSubsystemBalancedProductCode.HL, def:HorizontalSubsystemBalancedProductCode.HG, def:HorizontalVerticalHomologySplitting.kunnethIso, def:HorizontalVerticalHomologySplitting.directSumInc, def:HEquivariantChainComplex.balancedKunnethSum}
The logical and gauge submodules form a complementary pair in $H_1$:
\[
  H_1 = H_1^L \oplus H_1^G,
\]
i.e., $H_1^L$ and $H_1^G$ are complementary submodules ($\operatorname{IsCompl}(H_1^L, H_1^G)$).
\end{theorem}

\begin{proof}
\leanok
\uses{def:HorizontalSubsystemBalancedProductCode.HL, def:HorizontalSubsystemBalancedProductCode.HG, def:HorizontalSubsystemBalancedProductCode.embH, def:HorizontalSubsystemBalancedProductCode.embV, def:HorizontalSubsystemBalancedProductCode.incH, def:HorizontalSubsystemBalancedProductCode.incV, def:HorizontalVerticalHomologySplitting.kunnethIso, def:HorizontalVerticalHomologySplitting.directSumInc}
We verify both conditions for complementarity.

\medskip
\noindent\textbf{Disjointness} ($H_1^L \cap H_1^G = \{0\}$): We use the submodule disjointness criterion. Let $x \in H_1^L \cap H_1^G$. Then by definition of $H_1^L = \operatorname{im}(\operatorname{embH})$ and $H_1^G = \operatorname{im}(\operatorname{embV})$, there exist $a \in H_1^h$ and $b \in H_1^v$ such that $x = \operatorname{embH}(a) = \operatorname{embV}(b)$.

Let $\kappa$ denote the Künneth isomorphism. Applying $\kappa$ to both sides, we obtain
\[
  \kappa(\operatorname{embH}(a)) = \kappa(\operatorname{embV}(b)).
\]
Since $\operatorname{embH} = \kappa^{-1} \circ \operatorname{incH}$ and $\operatorname{embV} = \kappa^{-1} \circ \operatorname{incV}$, applying $\kappa$ gives $\kappa \circ \kappa^{-1} = \operatorname{id}$, so:
\[
  \operatorname{incH}(a) = \operatorname{incV}(b).
\]
Now $\operatorname{incH}$ includes into the direct summand at index $p = 1$, and $\operatorname{incV}$ includes into the summand at $p = 0$. Applying the component projection at $p = 1$ to both sides:
\[
  a = \pi_1(\operatorname{incH}(a)) = \pi_1(\operatorname{incV}(b)) = 0,
\]
since $\pi_1(\operatorname{incV}(b)) = 0$ as $\operatorname{incV}$ maps into the $p = 0$ summand. Thus $a = 0$, and by simplification $x = \operatorname{embH}(0) = 0$.

\medskip
\noindent\textbf{Codisjointness} ($H_1^L + H_1^G = H_1$): We show every $x \in H_1$ can be written as a sum of an element of $H_1^L$ and an element of $H_1^G$. Let $z = \kappa(x)$ be the image of $x$ under the Künneth isomorphism. Since the balanced Künneth sum at degree 1 decomposes over the summands at $p = 0$ and $p = 1$ (all other summands are trivial, as the Tanner and cycle complexes have PUnit modules outside degrees $\{0,1\}$), we have
\[
  z = \iota_1(z_h) + \iota_0(z_v),
\]
where $z_h = \pi_1(z) \in H_1^h$ and $z_v = \pi_0(z) \in H_1^v$ are the components at $p = 1$ and $p = 0$ respectively, and $\iota_p = \operatorname{directSumInc}(-,p)$ are the canonical inclusions.

Set $x_h = \operatorname{embH}(z_h) \in H_1^L$ and $x_v = \operatorname{embV}(z_v) \in H_1^G$. We claim $x = x_h + x_v$. Indeed, applying $\kappa$ (which is injective):
\[
  \kappa(x_h + x_v) = \kappa(\operatorname{embH}(z_h)) + \kappa(\operatorname{embV}(z_v)) = \operatorname{incH}(z_h) + \operatorname{incV}(z_v) = z = \kappa(x).
\]
By injectivity of $\kappa$, we conclude $x = x_h + x_v$, completing the proof.
\end{proof}

\begin{theorem}[Cohomology Complementarity]
\label{thm:HorizontalSubsystemBalancedProductCode.cohomology_iscompl}
\lean{HorizontalSubsystemBalancedProductCode.cohomology_iscompl}
\leanok
\uses{thm:HorizontalSubsystemBalancedProductCode.hcompl, def:HorizontalVerticalHomologySplitting.coH1h, def:HorizontalVerticalHomologySplitting.coH1v}
The cohomological logical submodule $H^1_L$ and gauge submodule $H^1_G$ are complementary:
\[
  H^1 = H^1_L \oplus H^1_G.
\]
Over $\mathbb{F}_2$, the cohomology and homology splitting are identified via Def.\ 2 and Def.\ 27, so $H^1_L = H_1^L$ and $H^1_G = H_1^G$ definitionally.
\end{theorem}

\begin{proof}
\leanok
\uses{thm:HorizontalSubsystemBalancedProductCode.hcompl}
Since over $\mathbb{F}_2$ the cohomology summands $\operatorname{coH}1h$ and $\operatorname{coH}1v$ defined in Def.\ 27 coincide definitionally with the homology summands $H_1^h$ and $H_1^v$ respectively, the cohomology logical and gauge submodules $\operatorname{coHL}$ and $\operatorname{coHG}$ are definitionally equal to $H_1^L$ and $H_1^G$. The result therefore follows immediately from the homology complementarity theorem $\operatorname{hcompl}$.
\end{proof}

\begin{theorem}[Logical Submodule Vanishes When Quotient Differential is Injective]
\label{thm:HorizontalSubsystemBalancedProductCode.HL_eq_bot_of_quotient_ker_trivial}
\lean{HorizontalSubsystemBalancedProductCode.HL_eq_bot_of_quotient_ker_trivial}
\leanok
\uses{def:HorizontalSubsystemBalancedProductCode.HL, def:IotaPiMaps.quotientTannerDifferential, def:IotaPiMaps.quotientTannerHomology, thm:IotaPiMaps.iotaMap_comp_piMap, thm:IotaPiMaps.iotaMap, thm:IotaPiMaps.piMap}
Suppose $\ell$ is odd ($\ell \equiv 1 \pmod{2}$) and $|H|$ is odd. If the kernel of the quotient Tanner differential $\partial : C_1(X/H, \Lambda) \to C_0(X/H, \Lambda)$ is trivial, i.e.,
\[
  \ker(\partial_{\mathrm{quot}}) = \{0\},
\]
then the logical submodule vanishes: $H_1^L = \{0\}$.
\end{theorem}

\begin{proof}
\leanok
\uses{def:HorizontalSubsystemBalancedProductCode.HL, def:HorizontalSubsystemBalancedProductCode.embH, def:IotaPiMaps.quotientTannerDifferential, def:IotaPiMaps.quotientTannerHomology, def:IotaPiMaps.quotientTannerComplex, thm:IotaPiMaps.iotaMap_comp_piMap, thm:IotaPiMaps.piMap, thm:IotaPiMaps.iotaMap}
We show that $H_1^L = \operatorname{im}(\operatorname{embH}) = \{0\}$ by proving every element in the image is zero.

Let $x \in H_1^L$, so there exists $a \in H_1^h$ with $x = \operatorname{embH}(a)$. We claim $a = 0$.

First, note that the cycles at degree 1 of the quotient Tanner complex equal the kernel of the quotient differential: $Z_1(\mathcal{C}(X/H,\Lambda)) = \ker(\partial_{\mathrm{quot}})$. Since $\ker(\partial_{\mathrm{quot}}) = \{0\}$ by hypothesis, the cycles submodule is trivial, and consequently the quotient Tanner homology $H_1(\mathcal{C}(X/H, \Lambda))$ is trivial (every element equals zero, since the homology is a quotient of a trivial submodule).

By the iota/pi isomorphism of Def.\ 28 ($\ell$ odd, $|H|$ odd), we have $\operatorname{iota} \circ \operatorname{pi} = \operatorname{id}$ on $H_1^h$. Applying this identity to $a$:
\[
  a = \operatorname{iota}\!\left(\operatorname{pi}(a)\right).
\]
Since $H_1(\mathcal{C}(X/H,\Lambda))$ is trivial, $\operatorname{pi}(a) = 0$, and therefore
\[
  a = \operatorname{iota}(0) = 0.
\]
Thus $x = \operatorname{embH}(0) = 0$, completing the proof that $H_1^L = \{0\}$.
\end{proof}

\begin{theorem}[Gauge Submodule Vanishes When $H_0$ of Tanner Complex is Trivial]
\label{thm:HorizontalSubsystemBalancedProductCode.HG_eq_bot_of_tanner_H0_subsingleton}
\lean{HorizontalSubsystemBalancedProductCode.HG_eq_bot_of_tanner_H0_subsingleton}
\leanok
\uses{def:HorizontalSubsystemBalancedProductCode.HG, def:HorizontalVerticalHomologySplitting.H1v, def:HEquivariantChainComplex.balancedKunnethSummand, def:HorizontalVerticalHomologySplitting.tannerHEq, def:HorizontalVerticalHomologySplitting.cycleHEq}
Suppose $H_0(\mathcal{C}(X,\Lambda))$ is trivial (i.e., a subsingleton). Then the gauge submodule vanishes: $H_1^G = \{0\}$.

\textit{Note:} The paper derives this from the weaker condition $H_0(\mathcal{C}(X/H,\Lambda)) = 0$ (surjectivity of the quotient differential). The present formalization uses the stronger hypothesis $H_0(\mathcal{C}(X,\Lambda)) = 0$ because the degree-0 iota/pi maps connecting quotient $H_0$ to equivariant $H_0$ coinvariants are not yet formalized.
\end{theorem}

\begin{proof}
\leanok
\uses{def:HorizontalSubsystemBalancedProductCode.HG, def:HorizontalSubsystemBalancedProductCode.embV, def:HorizontalVerticalHomologySplitting.H1v}
We show $H_1^G = \operatorname{im}(\operatorname{embV}) = \{0\}$ by proving every element in the image is zero.

Let $x \in H_1^G$, so there exists $a \in H_1^v$ with $x = \operatorname{embV}(a)$. We claim $a = 0$.

Recall that $H_1^v$ is the balanced Künneth summand at $p = 0$:
\[
  H_1^v = \operatorname{Coinv}_H\!\left(H_0(\mathcal{C}(X,\Lambda)) \otimes_{\mathbb{F}_2} H_1(\mathcal{C}_\ell)\right).
\]
Since $H_0(\mathcal{C}(X,\Lambda))$ is subsingleton by hypothesis, the tensor product $H_0(\mathcal{C}(X,\Lambda)) \otimes_{\mathbb{F}_2} H_1(\mathcal{C}_\ell)$ is also subsingleton. Indeed, any element of the tensor product can be written as a sum of pure tensors $m \otimes n$; since $m \in H_0(\mathcal{C}(X,\Lambda))$ is forced to be zero (the module is subsingleton), each pure tensor equals $0 \otimes n = 0$, so the entire element is zero.

Since the coinvariants module is a quotient of a subsingleton module, it is itself subsingleton. Therefore $H_1^v$ is subsingleton, which forces $a = 0$.

Hence $x = \operatorname{embV}(0) = 0$, so $H_1^G = \{0\}$.
\end{proof}

\begin{lemma}[Logical Submodule is Nonempty]
\label{lem:HorizontalSubsystemBalancedProductCode.HL_nonempty}
\lean{HorizontalSubsystemBalancedProductCode.HL_nonempty}
\leanok
\uses{def:HorizontalSubsystemBalancedProductCode.HL}
The carrier of the logical submodule $H_1^L$ is nonempty: $H_1^L \neq \emptyset$. In particular, $0 \in H_1^L$.
\end{lemma}

\begin{proof}
\leanok
\uses{def:HorizontalSubsystemBalancedProductCode.HL}
Since $H_1^L$ is a submodule of $H_1$, it contains the zero element. Therefore its carrier is nonempty, witnessed by $0 \in H_1^L$.
\end{proof}

\begin{lemma}[Gauge Submodule is Nonempty]
\label{lem:HorizontalSubsystemBalancedProductCode.HG_nonempty}
\lean{HorizontalSubsystemBalancedProductCode.HG_nonempty}
\leanok
\uses{def:HorizontalSubsystemBalancedProductCode.HG}
The carrier of the gauge submodule $H_1^G$ is nonempty: $H_1^G \neq \emptyset$. In particular, $0 \in H_1^G$.
\end{lemma}

\begin{proof}
\leanok
\uses{def:HorizontalSubsystemBalancedProductCode.HG}
Since $H_1^G$ is a submodule of $H_1$, it contains the zero element. Therefore its carrier is nonempty, witnessed by $0 \in H_1^G$.
\end{proof}

%--- Thm_13: HomologicalDistanceBound ---
\chapter{Thm 13: Homological Distance Bound}

\begin{definition}[Hamming Weight]
\label{def:HomologicalDistanceBound.hammingWeight}
\lean{HomologicalDistanceBound.hammingWeight}
\leanok
\uses{def:F2}
Let $A$ be a finite type with decidable equality. The \emph{Hamming weight} of a vector $x : A \to \mathbb{F}_2$ is the number of coordinates where $x$ is nonzero:
\[
  \operatorname{wt}(x) := \bigl|\{ a \in A \mid x(a) \neq 0 \}\bigr|.
\]
\end{definition}

\begin{lemma}[Nonzero Vectors Have Positive Hamming Weight]
\label{lem:HomologicalDistanceBound.hammingWeight_pos_of_ne_zero}
\lean{HomologicalDistanceBound.hammingWeight_pos_of_ne_zero}
\leanok
\uses{def:HomologicalDistanceBound.hammingWeight, def:F2}
Let $A$ be a finite type with decidable equality. If $x : A \to \mathbb{F}_2$ satisfies $x \neq 0$, then $\operatorname{wt}(x) > 0$.
\end{lemma}

\begin{proof}
\leanok
\uses{def:HomologicalDistanceBound.hammingWeight}
We proceed by contradiction. Suppose $\operatorname{wt}(x) = 0$, i.e., the filter set $\{a \in A \mid x(a) \neq 0\}$ has cardinality zero. By pushing negations and rewriting via $\operatorname{Finset.card_pos}$, the emptiness of the filter set means that for all $a \in A$, $x(a) = 0$. By function extensionality, this gives $x = 0$, contradicting $x \neq 0$.
\end{proof}

\begin{definition}[$(\alpha, \beta)$-Expanding Linear Map]
\label{def:HomologicalDistanceBound.IsExpandingLinMap}
\lean{HomologicalDistanceBound.IsExpandingLinMap}
\leanok
\uses{def:F2, def:HomologicalDistanceBound.hammingWeight}
Let $A$, $B$, $C$ be finite types with decidable equality. A linear map $f : (A \to \mathbb{F}_2) \to_{\mathbb{F}_2} (B \to C \to \mathbb{F}_2)$ is \emph{$(\alpha, \beta)$-expanding} (written $\operatorname{IsExpandingLinMap}(f, \alpha, \beta)$) if:
\begin{enumerate}
  \item $0 < \alpha \leq 1$ and $0 < \beta$;
  \item For every $x : A \to \mathbb{F}_2$ with $\operatorname{wt}(x) \leq \alpha \cdot |A|$, we have
    \[
      \bigl|\{(b,c) \in B \times C \mid f(x)(b,c) \neq 0\}\bigr| \;\geq\; \beta \cdot \operatorname{wt}(x).
    \]
\end{enumerate}
\end{definition}

\begin{lemma}[Classical Expander Distance Bound]
\label{lem:HomologicalDistanceBound.expanding_ker_weight_bound}
\lean{HomologicalDistanceBound.expanding_ker_weight_bound}
\leanok
\uses{def:HomologicalDistanceBound.IsExpandingLinMap, def:HomologicalDistanceBound.hammingWeight, def:F2}
Let $f : (A \to \mathbb{F}_2) \to_{\mathbb{F}_2} (B \to C \to \mathbb{F}_2)$ be an $(\alpha, \beta)$-expanding linear map. If $x \neq 0$ and $f(x) = 0$, then
\[
  \operatorname{wt}(x) > \alpha \cdot |A|.
\]
\end{lemma}

\begin{proof}
\leanok
\uses{def:HomologicalDistanceBound.IsExpandingLinMap, def:HomologicalDistanceBound.hammingWeight}
We proceed by contradiction. Suppose $\operatorname{wt}(x) \leq \alpha \cdot |A|$. From the $(\alpha, \beta)$-expanding property (decomposing $\operatorname{hexp}$ into its components $0 < \alpha$, $\alpha \leq 1$, $0 < \beta$, and the expansion inequality), we apply the expansion to $x$ and $h$ to obtain:
\[
  \bigl|\{(b,c) \mid f(x)(b,c) \neq 0\}\bigr| \;\geq\; \beta \cdot \operatorname{wt}(x).
\]
Since $f(x) = 0$, the output filter is empty: $|\{(b,c) \mid f(x)(b,c) \neq 0\}| = 0$. Indeed, for any $(b,c)$, we have $f(x)(b,c) = (f(x))(b)(c)$, and rewriting using $f(x) = 0$ gives $f(x)(b)(c) = 0$. Thus after rewriting with $\operatorname{hout}$, we obtain $0 \geq \beta \cdot \operatorname{wt}(x)$.

We then establish that $\operatorname{wt}(x) > 0$ as a real number: since $\operatorname{wt}(x) = |\{a \mid x(a) \neq 0\}|$ and the filter set is nonempty (if it were empty, every $x(a) = 0$ and function extensionality would yield $x = 0$, contradicting $x \neq 0$), we have $\operatorname{wt}(x) \geq 1 > 0$. Together with $\beta > 0$, this gives $\beta \cdot \operatorname{wt}(x) > 0$, contradicting $0 \geq \beta \cdot \operatorname{wt}(x)$ by nonlinear arithmetic.
\end{proof}

\begin{theorem}[Axiom: Hamming Weight on Cycle Representatives]
\label{thm:HomologicalDistanceBound.cycleRepWeight}
\lean{HomologicalDistanceBound.cycleRepWeight}
\leanok
\uses{def:HorizontalVerticalHomologySplitting.H1, def:GraphWithGroupAction, def:GraphWithGroupAction.IsInvariantLabeling, def:BalancedProductTannerCycleCode.CycleCompatibleAction}

\textbf{\color{red}[UNPROVEN AXIOM]}

Let $H$ be a finite group, $X$ a graph with $H$-action, $\Lambda$ a cell labeling, $\ell \geq 1$ an integer, $\mu_H$ an $H$-action on $\operatorname{Fin}(\ell)$, with $\Lambda$ invariant and the cycle-compatible action condition satisfied, and $|H|$ odd. There exists a function
\[
  \operatorname{cycleRepWeight} : H_1(C(X,\Lambda) \otimes_H C(C_\ell)) \to \mathbb{N}
\]
assigning to each homology class $[x] \in H_1$ the minimum Hamming weight $|u| + |v|$ over all cycle representatives $x = (u, v) \in \operatorname{Tot}_1 = E_{1,0} \oplus E_{0,1}$ of that class, where the basis is the coinvariant basis of the balanced product modules.

\textit{Note: This is stated as an axiom because the full proof—constructing the coinvariant basis and transporting Hamming weight through balanced product quotients—requires infrastructure beyond Mathlib's current scope. See van Lint, \emph{Introduction to Coding Theory}, Ch.\ 1; MacWilliams–Sloane, \emph{The Theory of Error-Correcting Codes}, \S1.1; and Panteleev–Kalachev [PK22].}
\end{theorem}

\begin{theorem}[Axiom: Hamming Weight of Zero is Zero]
\label{thm:HomologicalDistanceBound.cycleRepWeight_zero}
\lean{HomologicalDistanceBound.cycleRepWeight_zero}
\leanok
\uses{thm:HomologicalDistanceBound.cycleRepWeight, def:HorizontalVerticalHomologySplitting.H1, def:GraphWithGroupAction.IsInvariantLabeling, def:BalancedProductTannerCycleCode.CycleCompatibleAction}

\textbf{\color{red}[UNPROVEN AXIOM]}

Under the same hypotheses, $\operatorname{cycleRepWeight}(0) = 0$. The zero vector has no nonzero coordinates, so its Hamming weight is zero.

\textit{Note: This is stated as an axiom because it depends on the coinvariant basis infrastructure axiomatized in $\operatorname{cycleRepWeight}$.}
\end{theorem}

\begin{theorem}[Axiom: Nonzero Classes Have Positive Weight]
\label{thm:HomologicalDistanceBound.cycleRepWeight_pos_of_ne_zero}
\lean{HomologicalDistanceBound.cycleRepWeight_pos_of_ne_zero}
\leanok
\uses{thm:HomologicalDistanceBound.cycleRepWeight, def:HorizontalVerticalHomologySplitting.H1, def:GraphWithGroupAction.IsInvariantLabeling, def:BalancedProductTannerCycleCode.CycleCompatibleAction}

\textbf{\color{red}[UNPROVEN AXIOM]}

Under the same hypotheses, for any $x \in H_1$ with $x \neq 0$, we have $\operatorname{cycleRepWeight}(x) > 0$. If $x \neq 0$, then at least one coordinate is nonzero, so $\operatorname{wt}(x) \geq 1 > 0$.

\textit{Note: This is stated as an axiom because it depends on the coinvariant basis infrastructure axiomatized in $\operatorname{cycleRepWeight}$.}
\end{theorem}

\begin{lemma}[$H_1$ is Inhabited]
\label{lem:HomologicalDistanceBound.H1_inhabited}
\lean{HomologicalDistanceBound.H1_inhabited}
\leanok
\uses{def:HorizontalVerticalHomologySplitting.H1, def:GraphWithGroupAction.IsInvariantLabeling, def:BalancedProductTannerCycleCode.CycleCompatibleAction}
The type $H_1(C(X,\Lambda) \otimes_H C(C_\ell))$ is inhabited: there exists an element $x$ (namely $x = 0$) with $x = x$.
\end{lemma}

\begin{proof}
\leanok
\uses{def:HorizontalVerticalHomologySplitting.H1}
We exhibit the witness $x := 0$ and note that $0 = 0$ holds by reflexivity.
\end{proof}

\begin{theorem}[Axiom: Homological Distance Bound, Case 1 (Horizontal)]
\label{thm:HomologicalDistanceBound.homologicalDistanceBound_horizontal}
\lean{HomologicalDistanceBound.homologicalDistanceBound_horizontal}
\leanok
\uses{def:HomologicalDistanceBound.IsExpandingLinMap, thm:HomologicalDistanceBound.cycleRepWeight, def:HorizontalVerticalHomologySplitting.H1, def:HorizontalVerticalHomologySplitting.projH, def:BalancedProductTannerCycleCode.tannerDifferential, def:GraphWithGroupAction, def:GraphWithGroupAction.IsInvariantLabeling, def:BalancedProductTannerCycleCode.CycleCompatibleAction}

\textbf{\color{red}[UNPROVEN AXIOM]}

Let $H$ be a finite group of odd order, $X$ a graph with $H$-action, $\Lambda$ an $H$-invariant cell labeling, $\ell \geq 3$ an odd integer with $H$ acting on $\operatorname{Fin}(\ell)$ compatibly, and suppose the Tanner differential $\partial : C_1(X, \Lambda) \to C_0(X, \Lambda)$ is $(\alpha_{\mathrm{ho}}, \beta_{\mathrm{ho}})$-expanding. If $x \in H_1$ is a nonzero homology class with nontrivial horizontal projection $p^h([x]) \neq 0$, then the minimum Hamming weight of a cycle representative satisfies:
\[
  \operatorname{wt}(x) \;\geq\; |X_1| \cdot \min\!\left(\frac{\alpha_{\mathrm{ho}}}{2},\; \frac{\alpha_{\mathrm{ho}} \beta_{\mathrm{ho}}}{4}\right).
\]

\textit{Note: This is Theorem 5, Case 1 of Panteleev--Kalachev [PK22]. The proof proceeds by case analysis on $|u|$ relative to $\alpha_{\mathrm{ho}}|X_1|/2$: if $|u| \geq \alpha_{\mathrm{ho}}|X_1|/2$ the bound follows directly; if $|u| < \alpha_{\mathrm{ho}}|X_1|/2$ the expansion of $\partial$ applied fiber-by-fiber gives $|\partial u_{\mathrm{fiber}}| \geq \beta_{\mathrm{ho}}|u_{\mathrm{fiber}}|$, and the cycle condition $\partial^h u + \partial^v v = 0$ forces $|v| \geq \beta_{\mathrm{ho}}|u|/s$. It is axiomatized because the coinvariant basis and fiber decomposition infrastructure is beyond Mathlib's current scope.}
\end{theorem}

\begin{theorem}[Axiom: Homological Distance Bound, Case 2 (Vertical)]
\label{thm:HomologicalDistanceBound.homologicalDistanceBound_vertical}
\lean{HomologicalDistanceBound.homologicalDistanceBound_vertical}
\leanok
\uses{def:HomologicalDistanceBound.IsExpandingLinMap, thm:HomologicalDistanceBound.cycleRepWeight, def:HorizontalVerticalHomologySplitting.H1, def:HorizontalVerticalHomologySplitting.projH, def:BalancedProductTannerCycleCode.tannerDifferential, def:GraphWithGroupAction, def:GraphWithGroupAction.IsInvariantLabeling, def:BalancedProductTannerCycleCode.CycleCompatibleAction}

\textbf{\color{red}[UNPROVEN AXIOM]}

Under the same hypotheses as the horizontal bound (with additionally $s \geq 1$), if $x \in H_1$ is a nonzero homology class with trivial horizontal projection $p^h([x]) = 0$ (purely vertical class), then:
\[
  \operatorname{wt}(x) \;\geq\; \ell \cdot \min\!\left(\frac{\alpha_{\mathrm{ho}}}{4s},\; \frac{\alpha_{\mathrm{ho}} \beta_{\mathrm{ho}}}{4s}\right).
\]

\textit{Note: This is Theorem 5, Case 2 of Panteleev--Kalachev [PK22]. The proof decomposes $v$ into $\ell$ slices along the cycle graph edges. The cycle condition forces $\partial(v_j - v_{j+1}) = 0$, and expansion applied to consecutive differences with a pigeonhole argument over slices yields the bound. It is axiomatized for the same reasons as Case 1.}
\end{theorem}

\begin{lemma}[Horizontal Bound Multiplier is at Most One]
\label{lem:HomologicalDistanceBound.horizontal_bound_le_edges}
\lean{HomologicalDistanceBound.horizontal_bound_le_edges}
\leanok
\uses{def:HomologicalDistanceBound.IsExpandingLinMap}
For valid expansion parameters $0 < \alpha \leq 1$ and $0 < \beta$:
\[
  \min\!\left(\frac{\alpha}{2},\; \frac{\alpha\beta}{4}\right) \leq 1.
\]
\end{lemma}

\begin{proof}
\leanok
\uses{def:HomologicalDistanceBound.IsExpandingLinMap}
We apply $\min\_le\_of\_left\_le$ to reduce to showing $\alpha/2 \leq 1$. Since $\alpha \leq 1$, we have $\alpha/2 \leq 1/2 \leq 1$ by linear arithmetic.
\end{proof}

\begin{lemma}[Vertical Bound Multiplier is at Most One]
\label{lem:HomologicalDistanceBound.vertical_bound_le_one}
\lean{HomologicalDistanceBound.vertical_bound_le_one}
\leanok
\uses{def:HomologicalDistanceBound.IsExpandingLinMap}
For valid expansion parameters $0 < \alpha \leq 1$, $0 < \beta$, and integer $s \geq 1$:
\[
  \min\!\left(\frac{\alpha}{4s},\; \frac{\alpha\beta}{4s}\right) \leq 1.
\]
\end{lemma}

\begin{proof}
\leanok
\uses{def:HomologicalDistanceBound.IsExpandingLinMap}
We apply $\min\_le\_of\_left\_le$ to reduce to $\alpha/(4s) \leq 1$. We establish $0 < s$ as a real number via $\operatorname{Nat.cast_pos}$ (since $s \geq 1$ by $\omega$), and $(1 : \mathbb{R}) \leq s$ via exact casting of the natural number hypothesis. Rewriting $\alpha/(4s) \leq 1$ as $\alpha \leq 4s$ via $\operatorname{div_le_one}$ (using positivity of $4s$), the conclusion follows by linear arithmetic from $\alpha \leq 1 \leq s$.
\end{proof}

\begin{lemma}[Horizontal Bound is Positive]
\label{lem:HomologicalDistanceBound.horizontal_bound_pos}
\lean{HomologicalDistanceBound.horizontal_bound_pos}
\leanok
\uses{def:HomologicalDistanceBound.IsExpandingLinMap}
For valid expansion parameters $0 < \alpha$, $0 < \beta$, and any edge count $n_E \geq 1$:
\[
  0 \;<\; n_E \cdot \min\!\left(\frac{\alpha}{2},\; \frac{\alpha\beta}{4}\right).
\]
\end{lemma}

\begin{proof}
\leanok
\uses{def:HomologicalDistanceBound.IsExpandingLinMap}
We apply $\operatorname{mul_pos}$. The factor $n_E > 0$ follows from $\operatorname{Nat.cast_pos}$ applied to $hn_E$. For the minimum factor, we apply $\operatorname{lt_min}$ and use positivity to establish $0 < \alpha/2$ and $0 < \alpha\beta/4$ from $\alpha > 0$ and $\beta > 0$.
\end{proof}

\begin{lemma}[Vertical Bound is Positive]
\label{lem:HomologicalDistanceBound.vertical_bound_pos}
\lean{HomologicalDistanceBound.vertical_bound_pos}
\leanok
\uses{def:HomologicalDistanceBound.IsExpandingLinMap}
For valid expansion parameters $0 < \alpha$, $0 < \beta$, integer $s \geq 1$, and cycle length $0 < \ell$:
\[
  0 \;<\; \ell \cdot \min\!\left(\frac{\alpha}{4s},\; \frac{\alpha\beta}{4s}\right).
\]
\end{lemma}

\begin{proof}
\leanok
\uses{def:HomologicalDistanceBound.IsExpandingLinMap}
We apply $\operatorname{mul_pos}$. The factor $\ell > 0$ as a real number follows from $\operatorname{Nat.cast_pos}$ applied to $h\ell$. For the minimum factor, we establish $0 < s$ as a real number via $\operatorname{Nat.cast_pos}$ (using $s \geq 1$ by $\omega$), then apply $\operatorname{lt_min}$ and use positivity to establish $0 < \alpha/(4s)$ and $0 < \alpha\beta/(4s)$ from $\alpha > 0$, $\beta > 0$, and $s > 0$.
\end{proof}

%--- Thm_14: CohomologicalDistanceBound ---
\chapter{Thm 14: Cohomological Distance Bound}

\begin{definition}[Coboundary Map]
\label{def:CohomologicalDistanceBound.coboundaryMap}
\lean{CohomologicalDistanceBound.coboundaryMap}
\leanok
\uses{def:GraphWithGroupAction, def:GraphWithGroupAction.CellLabeling, def:GraphWithGroupAction.cellIncidentEdges, def:F2}

Let $H$ be a finite group, $X$ a graph with $H$-action, and $\Lambda$ a cell labeling of $X$ with $s$ labels. The \emph{coboundary map} (or transpose Tanner differential)
\[
\delta = \partial^T : C_0(X, \Lambda) \longrightarrow C_1(X, \Lambda)
\]
is the linear map $\delta : (X_0 \to \operatorname{Fin}(s) \to \mathbb{F}_2) \to (X_1 \to \mathbb{F}_2)$ defined by
\[
\delta(y)(e) = \sum_{v \in X_0} \begin{cases} y\bigl(v,\, \Lambda_v(e)\bigr) & \text{if } e \in \operatorname{cellIncidentEdges}(v), \\ 0 & \text{otherwise.} \end{cases}
\]
Here $X_0$ denotes the $0$-cells (vertices) and $X_1$ the $1$-cells (edges) of $X$. This map is $\mathbb{F}_2$-linear, being the transpose of the Tanner differential $\partial : C_1(X,\Lambda) \to C_0(X,\Lambda)$.
\end{definition}

\begin{definition}[Expanding Coboundary Property]
\label{def:CohomologicalDistanceBound.IsExpandingCoboundary}
\lean{CohomologicalDistanceBound.IsExpandingCoboundary}
\leanok
\uses{def:HomologicalDistanceBound.IsExpandingLinMap, def:HomologicalDistanceBound.hammingWeight}

Let $A$, $B$, $C$ be finite types with decidable equality, and let $f : (A \to B \to \mathbb{F}_2) \to_{\mathbb{F}_2} (C \to \mathbb{F}_2)$ be a linear map over $\mathbb{F}_2$. We say $f$ is \emph{$(\alpha, \beta)$-expanding} (written $\operatorname{IsExpandingCoboundary}(f, \alpha, \beta)$) if:
\begin{enumerate}
  \item $0 < \alpha \leq 1$ and $0 < \beta$,
  \item For every $x : A \to B \to \mathbb{F}_2$ with
    \[
    \bigl|\{(a,b) \in A \times B \mid x(a,b) \neq 0\}\bigr| \leq \alpha \cdot |A| \cdot |B|,
    \]
    we have
    \[
    \bigl|\{c \in C \mid f(x)(c) \neq 0\}\bigr| \geq \beta \cdot \bigl|\{(a,b) \in A \times B \mid x(a,b) \neq 0\}\bigr|.
    \]
\end{enumerate}
This is the expanding property for a map whose domain is a matrix-shaped space $(A \to B \to \mathbb{F}_2)$, dual to the $\operatorname{IsExpandingLinMap}$ condition of Definition~\ref{def:HomologicalDistanceBound.IsExpandingLinMap}.
\end{definition}

\begin{definition}[Number of Vertices]
\label{def:CohomologicalDistanceBound.numVertices}
\lean{CohomologicalDistanceBound.numVertices}
\leanok
\uses{def:GraphWithGroupAction}

Let $H$ be a finite group and $X$ a graph with $H$-action. The \emph{number of $0$-cells (vertices)} of $X$ is
\[
|X_0| := \bigl|\operatorname{Fintype.card}(X\text{.graph.cells}(0))\bigr|.
\]
\end{definition}

\begin{theorem}[Axiom: Cohomological Distance Bound, Case 1 (Vertical)]
\label{thm:CohomologicalDistanceBound.cohomologicalDistanceBound_vertical}
\lean{CohomologicalDistanceBound.cohomologicalDistanceBound_vertical}
\leanok
\uses{def:GraphWithGroupAction, def:GraphWithGroupAction.CellLabeling, def:GraphWithGroupAction.IsInvariantLabeling, def:BalancedProductTannerCycleCode.CycleCompatibleAction, def:CohomologicalDistanceBound.coboundaryMap, def:CohomologicalDistanceBound.IsExpandingCoboundary, def:HorizontalVerticalHomologySplitting.H1, def:HorizontalVerticalHomologySplitting.coprojV, thm:HomologicalDistanceBound.cycleRepWeight}

\textbf{\color{red}[UNPROVEN AXIOM]}

Let $H$ be a finite group, $X$ a graph with $H$-action, $\Lambda$ a cell labeling with $s$ labels, $\ell \geq 3$ an odd integer with $H \curvearrowright \operatorname{Fin}(\ell)$, $\Lambda$ invariant, the action cycle-compatible, and $|H|$ odd. Let $\alpha_{\mathrm{co}}, \beta_{\mathrm{co}} \in \mathbb{R}$ be such that the coboundary map $\delta = \partial^T$ is $(\alpha_{\mathrm{co}}, \beta_{\mathrm{co}})$-expanding. Let $[x] \in H^1$ be a nontrivial cohomology class (i.e.\ $x \neq 0$) with nontrivial vertical cohomological projection: $p^v([x]) \neq 0$.

Then the minimum Hamming weight of a cycle representative satisfies
\[
|x| \;\geq\; |X_0| \cdot s \cdot \min\!\Bigl(\tfrac{\alpha_{\mathrm{co}}}{2},\; \tfrac{\alpha_{\mathrm{co}} \cdot \beta_{\mathrm{co}}}{4}\Bigr).
\]

\textit{Note: This is stated as an axiom because the full proof was not completed in the formalization. It is the cohomological dual of Theorem~\ref{thm:HomologicalDistanceBound.homologicalDistanceBound_vertical}, obtained by transposing the roles of boundary and coboundary. The factor $|X_0| \cdot s$ arises because the domain of $\delta$ is $C_0(X,\Lambda) \cong \mathbb{F}_2^{X_0 \times [s]}$.}
\end{theorem}

\begin{theorem}[Axiom: Cohomological Distance Bound, Case 2 (Horizontal)]
\label{thm:CohomologicalDistanceBound.cohomologicalDistanceBound_horizontal}
\lean{CohomologicalDistanceBound.cohomologicalDistanceBound_horizontal}
\leanok
\uses{def:GraphWithGroupAction, def:GraphWithGroupAction.CellLabeling, def:GraphWithGroupAction.IsInvariantLabeling, def:BalancedProductTannerCycleCode.CycleCompatibleAction, def:CohomologicalDistanceBound.coboundaryMap, def:CohomologicalDistanceBound.IsExpandingCoboundary, def:HorizontalVerticalHomologySplitting.H1, def:HorizontalVerticalHomologySplitting.coprojV, thm:HomologicalDistanceBound.cycleRepWeight}

\textbf{\color{red}[UNPROVEN AXIOM]}

Let $H$, $X$, $\Lambda$, $\ell$, $\alpha_{\mathrm{co}}$, $\beta_{\mathrm{co}}$ be as in Theorem~\ref{thm:CohomologicalDistanceBound.cohomologicalDistanceBound_vertical}, with the additional assumption $s \geq 1$. Let $[x] \in H^1$ be a nontrivial cohomology class with \emph{trivial} vertical projection: $p^v([x]) = 0$ (i.e.\ $[x]$ is purely horizontal).

Then
\[
|x| \;\geq\; \ell \cdot \min\!\Bigl(\tfrac{\alpha_{\mathrm{co}}}{4s},\; \tfrac{\alpha_{\mathrm{co}} \cdot \beta_{\mathrm{co}}}{4s}\Bigr).
\]

\textit{Note: This is stated as an axiom because the full proof was not completed in the formalization. It is the cohomological dual of Theorem~\ref{thm:HomologicalDistanceBound.homologicalDistanceBound_horizontal}, with the coboundary expansion applied fiber-by-fiber along the $\ell$ edges of $C_\ell$.}
\end{theorem}

\begin{theorem}[Axiom: Satisfiability Witness for Vertical Cohomological Distance Bound]
\label{thm:CohomologicalDistanceBound.cohomologicalDistanceBound_vertical_satisfiable}
\lean{CohomologicalDistanceBound.cohomologicalDistanceBound_vertical_satisfiable}
\leanok
\uses{def:GraphWithGroupAction, def:GraphWithGroupAction.CellLabeling, def:GraphWithGroupAction.IsInvariantLabeling, def:BalancedProductTannerCycleCode.CycleCompatibleAction, def:CohomologicalDistanceBound.coboundaryMap, def:CohomologicalDistanceBound.IsExpandingCoboundary, def:HorizontalVerticalHomologySplitting.H1, def:HorizontalVerticalHomologySplitting.coprojV, thm:HomologicalDistanceBound.cycleRepWeight, thm:CohomologicalDistanceBound.cohomologicalDistanceBound_vertical}

\textbf{\color{red}[UNPROVEN AXIOM]}

The hypotheses of Theorem~\ref{thm:CohomologicalDistanceBound.cohomologicalDistanceBound_vertical} are jointly satisfiable. That is, there exist a finite group $H$, a graph $X$ with $H$-action, a cell labeling $\Lambda$ with $s$ labels, an odd $\ell \geq 3$ with cycle-compatible $H$-action on $\operatorname{Fin}(\ell)$, invariant labeling, odd $|H|$, parameters $\alpha_{\mathrm{co}}, \beta_{\mathrm{co}} > 0$ with the coboundary $\delta$ being $(\alpha_{\mathrm{co}}, \beta_{\mathrm{co}})$-expanding, and a nonzero cohomology class $x \in H^1$ with $p^v(x) \neq 0$, such that
\[
|x| \;\geq\; |X_0| \cdot s \cdot \min\!\Bigl(\tfrac{\alpha_{\mathrm{co}}}{2},\; \tfrac{\alpha_{\mathrm{co}} \cdot \beta_{\mathrm{co}}}{4}\Bigr).
\]

\textit{Note: This axiom serves as a satisfiability witness confirming the premises of the vertical cohomological distance bound are consistent.}
\end{theorem}

\begin{theorem}[Axiom: Satisfiability Witness for Horizontal Cohomological Distance Bound]
\label{thm:CohomologicalDistanceBound.cohomologicalDistanceBound_horizontal_satisfiable}
\lean{CohomologicalDistanceBound.cohomologicalDistanceBound_horizontal_satisfiable}
\leanok
\uses{def:GraphWithGroupAction, def:GraphWithGroupAction.CellLabeling, def:GraphWithGroupAction.IsInvariantLabeling, def:BalancedProductTannerCycleCode.CycleCompatibleAction, def:CohomologicalDistanceBound.coboundaryMap, def:CohomologicalDistanceBound.IsExpandingCoboundary, def:HorizontalVerticalHomologySplitting.H1, def:HorizontalVerticalHomologySplitting.coprojV, thm:HomologicalDistanceBound.cycleRepWeight, thm:CohomologicalDistanceBound.cohomologicalDistanceBound_horizontal}

\textbf{\color{red}[UNPROVEN AXIOM]}

The hypotheses of Theorem~\ref{thm:CohomologicalDistanceBound.cohomologicalDistanceBound_horizontal} are jointly satisfiable. That is, there exist a finite group $H$, a graph $X$ with $H$-action, integers $s \geq 1$ and $\ell \geq 3$ odd, a cell labeling $\Lambda$, cycle-compatible $H$-action on $\operatorname{Fin}(\ell)$, invariant labeling, odd $|H|$, parameters $\alpha_{\mathrm{co}}, \beta_{\mathrm{co}} > 0$ with the coboundary $\delta$ being $(\alpha_{\mathrm{co}}, \beta_{\mathrm{co}})$-expanding, and a nonzero cohomology class $x \in H^1$ with $p^v(x) = 0$ (purely horizontal), such that
\[
|x| \;\geq\; \ell \cdot \min\!\Bigl(\tfrac{\alpha_{\mathrm{co}}}{4s},\; \tfrac{\alpha_{\mathrm{co}} \cdot \beta_{\mathrm{co}}}{4s}\Bigr).
\]

\textit{Note: This axiom serves as a satisfiability witness confirming the premises of the horizontal cohomological distance bound are consistent, including the requirement $s \geq 1$ needed for the $\frac{1}{4s}$ factors.}
\end{theorem}

%--- Cor_2: SubsystemCodeParameters ---
\chapter{Cor 2: Subsystem Code Parameters}

\begin{theorem}[Number of Physical Qubits]
\label{thm:SubsystemCodeParameters.numQubits_eq}
\lean{SubsystemCodeParameters.numQubits_eq}
\leanok
\uses{def:BalancedProductTannerCycleCode.balancedProductTannerCycleCode, def:GraphWithGroupAction.IsInvariantLabeling, def:BalancedProductTannerCycleCode.CycleCompatibleAction}

Let $H$ be a finite group acting on $\mathbb{F}_\ell$ for some $\ell \geq 3$ odd, let $X$ be a graph with $H$-action, let $\Lambda$ be an $s$-regular $H$-invariant cell labeling, and let the action be cycle-compatible. Suppose:
\begin{itemize}
  \item $\operatorname{hdim_components}$: $\dim_{\mathbb{F}_2}(\operatorname{Tot}_1) = |X_1| + s \cdot |X_0|$, where $\operatorname{Tot}_1$ denotes the degree-$1$ component of the balanced product Tanner cycle complex;
  \item $\operatorname{hreg}$: $2|X_1| = s \cdot |X_0|$ (i.e.\ the graph $X$ is $s$-regular).
\end{itemize}
Then
\[
  \dim_{\mathbb{F}_2}(\operatorname{Tot}_1) = 3|X_1|.
\]
\end{theorem}

\begin{proof}
\leanok
\uses{def:BalancedProductTannerCycleCode.balancedProductTannerCycleCode}

Rewriting with the hypothesis $\operatorname{hdim_components}$, the goal becomes
$|X_1| + s \cdot |X_0| = 3|X_1|$.
By the regularity hypothesis $\operatorname{hreg}$, we have $s \cdot |X_0| = 2|X_1|$, so $|X_1| + s \cdot |X_0| = |X_1| + 2|X_1| = 3|X_1|$. This follows by linear integer arithmetic.
\end{proof}

\begin{theorem}[Logical Qubit Count]
\label{thm:SubsystemCodeParameters.subsystemCodeParameters_K}
\lean{SubsystemCodeParameters.subsystemCodeParameters_K}
\leanok
\uses{def:HorizontalSubsystemBalancedProductCode.HL, def:HorizontalVerticalHomologySplitting.H1h, def:HorizontalSubsystemBalancedProductCode.embH, def:HorizontalVerticalHomologySplitting.kunnethIso, def:HorizontalSubsystemBalancedProductCode.incH, def:HorizontalVerticalHomologySplitting.directSumInc}

The horizontal subsystem logical qubit space satisfies
\[
  \dim_{\mathbb{F}_2}\bigl(\operatorname{HL}(X, \Lambda, \ell)\bigr)
  = \dim_{\mathbb{F}_2}\bigl(H_1^h(X, \Lambda, \ell)\bigr).
\]
\end{theorem}

\begin{proof}
\leanok
\uses{def:HorizontalSubsystemBalancedProductCode.HL, def:HorizontalVerticalHomologySplitting.H1h, def:HorizontalSubsystemBalancedProductCode.embH, def:HorizontalVerticalHomologySplitting.kunnethIso, def:HorizontalSubsystemBalancedProductCode.incH, def:HorizontalVerticalHomologySplitting.directSumInc}

Unfolding the definition of $\operatorname{HL}$ as the range of the embedding map $\operatorname{embH}$, it suffices to show that $\operatorname{embH}$ is injective, and then apply the fact that the rank of the range of an injective linear map equals the rank of the domain.

We show injectivity of $\operatorname{embH} = \operatorname{kunnethIso}^{-1} \circ \operatorname{incH}$. Suppose $\operatorname{embH}(a) = \operatorname{embH}(b)$. Unfolding, we have $\operatorname{kunnethIso}^{-1}(\operatorname{incH}(a)) = \operatorname{kunnethIso}^{-1}(\operatorname{incH}(b))$. By injectivity of $\operatorname{kunnethIso}^{-1}$ (as a linear equivalence), we obtain $\operatorname{incH}(a) = \operatorname{incH}(b)$. Unfolding $\operatorname{incH}$ as the direct sum inclusion $\operatorname{DirectSum.lof}(\cdot, 1)$, injectivity of direct sum inclusions then gives $a = b$.
\end{proof}

\begin{theorem}[Logical Qubit Lower Bound]
\label{thm:SubsystemCodeParameters.subsystemCodeParameters_K_lowerBound}
\lean{SubsystemCodeParameters.subsystemCodeParameters_K_lowerBound}
\leanok
\uses{thm:SubsystemCodeParameters.subsystemCodeParameters_K, thm:EncodingRateCircle.encodingRateCircle, def:IotaPiMaps.quotientTannerHomology, def:IotaPiMaps.quotientTannerDifferential}

Let $k_L = \dim_{\mathbb{F}_2}(\ker(\partial^{\mathrm{quot}}))$ be the local code dimension, and let $m = |(X/H)_1|$ be the number of edges of the quotient graph, satisfying $m = |X_1|/\ell$ (the free action hypothesis). Assume $s \geq 1$ and that the quotient Tanner code homology satisfies the Sipser--Spielman bound
\[
  \dim_{\mathbb{F}_2}\bigl(H_1(C(X/H, L))\bigr) \geq \Bigl(\frac{2k_L}{s} - 1\Bigr) \cdot m.
\]
Then
\[
  \dim_{\mathbb{F}_2}\bigl(\operatorname{HL}(X, \Lambda, \ell)\bigr) \geq \Bigl(\frac{2k_L}{s} - 1\Bigr) \cdot \frac{|X_1|}{\ell}.
\]
\end{theorem}

\begin{proof}
\leanok
\uses{thm:SubsystemCodeParameters.subsystemCodeParameters_K, thm:EncodingRateCircle.encodingRateCircle, def:IotaPiMaps.quotientTannerHomology, def:IotaPiMaps.quotientTannerDifferential}

We proceed in three steps.

\textbf{Step 1.} By Theorem~\ref{thm:SubsystemCodeParameters.subsystemCodeParameters_K},
\[
  \dim_{\mathbb{F}_2}(\operatorname{HL}) = \dim_{\mathbb{F}_2}(H_1^h).
\]

\textbf{Step 2.} By Theorem~\ref{thm:EncodingRateCircle.encodingRateCircle} (encoding rate for the circle complex),
\[
  \dim_{\mathbb{F}_2}(H_1^h) = \dim_{\mathbb{F}_2}(H_1(C(X/H, L))).
\]
Combining steps 1 and 2, we obtain
\[
  \dim_{\mathbb{F}_2}(\operatorname{HL}) = \dim_{\mathbb{F}_2}(H_1(C(X/H, L))).
\]

\textbf{Step 3.} Substituting the hypothesis $m = |X_1|/\ell$ (i.e.\ $\operatorname{hQuotEdges}$) into the assumed Sipser--Spielman lower bound $\operatorname{hQuotHomology}$, and combining with step 2 via the chain of equalities (cast to $\mathbb{R}$), the conclusion follows directly.
\end{proof}

\begin{theorem}[Z-Distance Bound]
\label{thm:SubsystemCodeParameters.subsystemCodeParameters_DZ}
\lean{SubsystemCodeParameters.subsystemCodeParameters_DZ}
\leanok
\uses{def:HomologicalDistanceBound.IsExpandingLinMap, def:BalancedProductTannerCycleCode.tannerDifferential, def:HorizontalVerticalHomologySplitting.H1, def:HorizontalVerticalHomologySplitting.projH, thm:HomologicalDistanceBound.cycleRepWeight, thm:HomologicalDistanceBound.homologicalDistanceBound_horizontal}

Let $\alpha_{\mathrm{ho}}, \beta_{\mathrm{ho}} > 0$ and suppose the Tanner differential $\partial^{\mathrm{Tanner}}$ is $(\alpha_{\mathrm{ho}}, \beta_{\mathrm{ho}})$-expanding. For every nontrivial homology class $x \in H_1 \setminus \{0\}$ whose horizontal projection $\pi_H(x) \neq 0$,
\[
  \operatorname{wt}(x) \geq |X_1| \cdot \min\!\Bigl(\frac{\alpha_{\mathrm{ho}}}{2},\, \frac{\alpha_{\mathrm{ho}} \beta_{\mathrm{ho}}}{4}\Bigr).
\]
\end{theorem}

\begin{proof}
\leanok
\uses{thm:HomologicalDistanceBound.homologicalDistanceBound_horizontal}

This is a direct application of Theorem~\ref{thm:HomologicalDistanceBound.homologicalDistanceBound_horizontal} (homological distance bound, horizontal case) with the given expansion parameters and the hypotheses $x \neq 0$ and $\pi_H(x) \neq 0$.
\end{proof}

\begin{theorem}[X-Distance Bound, Vertical Case]
\label{thm:SubsystemCodeParameters.subsystemCodeParameters_DX_vertical}
\lean{SubsystemCodeParameters.subsystemCodeParameters_DX_vertical}
\leanok
\uses{def:CohomologicalDistanceBound.IsExpandingCoboundary, def:CohomologicalDistanceBound.coboundaryMap, def:HorizontalVerticalHomologySplitting.H1, def:HorizontalVerticalHomologySplitting.coprojV, thm:HomologicalDistanceBound.cycleRepWeight, thm:CohomologicalDistanceBound.cohomologicalDistanceBound_vertical}

Let $\alpha_{\mathrm{co}}, \beta_{\mathrm{co}} > 0$ and suppose the coboundary map is $(\alpha_{\mathrm{co}}, \beta_{\mathrm{co}})$-expanding. For every nontrivial cohomology class $x \in H_1 \setminus \{0\}$ whose vertical co-projection $\tilde{\pi}_V(x) \neq 0$,
\[
  \operatorname{wt}(x) \geq (|X_0| \cdot s) \cdot \min\!\Bigl(\frac{\alpha_{\mathrm{co}}}{2},\, \frac{\alpha_{\mathrm{co}} \beta_{\mathrm{co}}}{4}\Bigr).
\]
Since $s$-regularity gives $s \cdot |X_0| = 2|X_1|$, this is equivalent to
\[
  \operatorname{wt}(x) \geq \min\!\Bigl(\alpha_{\mathrm{co}} |X_1|,\, \frac{\alpha_{\mathrm{co}} \beta_{\mathrm{co}} |X_1|}{2}\Bigr).
\]
\end{theorem}

\begin{proof}
\leanok
\uses{thm:CohomologicalDistanceBound.cohomologicalDistanceBound_vertical}

This is a direct application of Theorem~\ref{thm:CohomologicalDistanceBound.cohomologicalDistanceBound_vertical} (cohomological distance bound, vertical case) with the given expansion parameters and the hypotheses $x \neq 0$ and $\tilde{\pi}_V(x) \neq 0$.
\end{proof}

\begin{theorem}[X-Distance Bound, Horizontal Case]
\label{thm:SubsystemCodeParameters.subsystemCodeParameters_DX_horizontal}
\lean{SubsystemCodeParameters.subsystemCodeParameters_DX_horizontal}
\leanok
\uses{def:CohomologicalDistanceBound.IsExpandingCoboundary, def:CohomologicalDistanceBound.coboundaryMap, def:HorizontalVerticalHomologySplitting.H1, def:HorizontalVerticalHomologySplitting.coprojV, thm:HomologicalDistanceBound.cycleRepWeight, thm:CohomologicalDistanceBound.cohomologicalDistanceBound_horizontal}

Let $\alpha_{\mathrm{co}}, \beta_{\mathrm{co}} > 0$, $s \geq 1$, and suppose the coboundary map is $(\alpha_{\mathrm{co}}, \beta_{\mathrm{co}})$-expanding. For every nontrivial cohomology class $x \in H_1 \setminus \{0\}$ whose vertical co-projection satisfies $\tilde{\pi}_V(x) = 0$ (purely horizontal class),
\[
  \operatorname{wt}(x) \geq \ell \cdot \min\!\Bigl(\frac{\alpha_{\mathrm{co}}}{4s},\, \frac{\alpha_{\mathrm{co}} \beta_{\mathrm{co}}}{4s}\Bigr).
\]
\end{theorem}

\begin{proof}
\leanok
\uses{thm:CohomologicalDistanceBound.cohomologicalDistanceBound_horizontal}

This is a direct application of Theorem~\ref{thm:CohomologicalDistanceBound.cohomologicalDistanceBound_horizontal} (cohomological distance bound, horizontal case) with the given expansion parameters and the hypotheses $x \neq 0$ and $\tilde{\pi}_V(x) = 0$.
\end{proof}

\begin{theorem}[Combined X-Distance Bound]
\label{thm:SubsystemCodeParameters.subsystemCodeParameters_DX_combined}
\lean{SubsystemCodeParameters.subsystemCodeParameters_DX_combined}
\leanok
\uses{thm:SubsystemCodeParameters.subsystemCodeParameters_DX_vertical, thm:SubsystemCodeParameters.subsystemCodeParameters_DX_horizontal, def:CohomologicalDistanceBound.IsExpandingCoboundary, def:CohomologicalDistanceBound.coboundaryMap, def:HorizontalVerticalHomologySplitting.H1, thm:HomologicalDistanceBound.cycleRepWeight}

Let $\alpha_{\mathrm{co}}, \beta_{\mathrm{co}} > 0$, $s \geq 1$, and suppose the coboundary map is $(\alpha_{\mathrm{co}}, \beta_{\mathrm{co}})$-expanding. For every nontrivial homology class $x \in H_1 \setminus \{0\}$,
\[
  \operatorname{wt}(x) \geq \min\!\left(
    (|X_0| \cdot s) \cdot \min\!\Bigl(\frac{\alpha_{\mathrm{co}}}{2},\, \frac{\alpha_{\mathrm{co}} \beta_{\mathrm{co}}}{4}\Bigr),\;
    \ell \cdot \min\!\Bigl(\frac{\alpha_{\mathrm{co}}}{4s},\, \frac{\alpha_{\mathrm{co}} \beta_{\mathrm{co}}}{4s}\Bigr)
  \right).
\]
Using $s$-regularity $s \cdot |X_0| = 2|X_1|$, this gives
\[
  D_X \geq \min\!\Bigl(\alpha_{\mathrm{co}} |X_1|,\; \tfrac{\alpha_{\mathrm{co}} \beta_{\mathrm{co}} |X_1|}{2},\; \tfrac{\ell \alpha_{\mathrm{co}}}{4s},\; \tfrac{\ell \alpha_{\mathrm{co}} \beta_{\mathrm{co}}}{4s}\Bigr).
\]
\end{theorem}

\begin{proof}
\leanok
\uses{thm:SubsystemCodeParameters.subsystemCodeParameters_DX_vertical, thm:SubsystemCodeParameters.subsystemCodeParameters_DX_horizontal}

We perform a case analysis on whether the vertical co-projection $\tilde{\pi}_V(x)$ is zero or not.

\textbf{Case 1: $\tilde{\pi}_V(x) = 0$ (purely horizontal).} We apply Theorem~\ref{thm:SubsystemCodeParameters.subsystemCodeParameters_DX_horizontal} to obtain
\[
  \operatorname{wt}(x) \geq \ell \cdot \min\!\Bigl(\frac{\alpha_{\mathrm{co}}}{4s},\, \frac{\alpha_{\mathrm{co}} \beta_{\mathrm{co}}}{4s}\Bigr).
\]
Since the right-hand side equals the second term of the outer minimum, the conclusion follows from $\min(A, B) \leq B$.

\textbf{Case 2: $\tilde{\pi}_V(x) \neq 0$ (nontrivial vertical component).} We apply Theorem~\ref{thm:SubsystemCodeParameters.subsystemCodeParameters_DX_vertical} to obtain
\[
  \operatorname{wt}(x) \geq (|X_0| \cdot s) \cdot \min\!\Bigl(\frac{\alpha_{\mathrm{co}}}{2},\, \frac{\alpha_{\mathrm{co}} \beta_{\mathrm{co}}}{4}\Bigr).
\]
Since the right-hand side equals the first term of the outer minimum, the conclusion follows from $\min(A, B) \leq A$.
\end{proof}

\begin{theorem}[Combined X-Distance Bound for Nontrivial Classes]
\label{thm:SubsystemCodeParameters.subsystemCodeParameters_DX_for_nontrivial}
\lean{SubsystemCodeParameters.subsystemCodeParameters_DX_for_nontrivial}
\leanok
\uses{thm:SubsystemCodeParameters.subsystemCodeParameters_DX_combined, def:CohomologicalDistanceBound.IsExpandingCoboundary, def:CohomologicalDistanceBound.coboundaryMap, def:HorizontalVerticalHomologySplitting.H1, thm:HomologicalDistanceBound.cycleRepWeight}

Under the same hypotheses as Theorem~\ref{thm:SubsystemCodeParameters.subsystemCodeParameters_DX_combined}, for every nontrivial homology class $x \in H_1 \setminus \{0\}$,
\[
  \operatorname{wt}(x) \geq \min\!\left(
    (|X_0| \cdot s) \cdot \min\!\Bigl(\frac{\alpha_{\mathrm{co}}}{2},\, \frac{\alpha_{\mathrm{co}} \beta_{\mathrm{co}}}{4}\Bigr),\;
    \ell \cdot \min\!\Bigl(\frac{\alpha_{\mathrm{co}}}{4s},\, \frac{\alpha_{\mathrm{co}} \beta_{\mathrm{co}}}{4s}\Bigr)
  \right).
\]

\textit{Note: The existence of nontrivial homology classes for specific graph families (Cayley expanders) is a constructive combinatorial result from the balanced product code construction \cite{PK22}.}
\end{theorem}

\begin{proof}
\leanok
\uses{thm:SubsystemCodeParameters.subsystemCodeParameters_DX_combined}

This is a direct application of Theorem~\ref{thm:SubsystemCodeParameters.subsystemCodeParameters_DX_combined} with the given expansion parameters and the hypothesis $x \neq 0$.
\end{proof}

\begin{remark}[Paper Corrections]
\textbf{\color{red}Errata.} The following errors were identified in the original paper and corrected in this formalization:
\begin{itemize}
  \item Paper Corollary (cor:distanceboundssybsystemcode) states D_X ≥ min{α_co|X₁|, α_co|X₁|/2, ℓα_co/(4s), ℓα_coβ_co/(4s)} but the second term is missing β_co. From Theorem thm:distco Case 1, |x| ≥ |X₀|s·min{α_co/2, α_coβ_co/4}. Substituting |X₀|s = 2|X₁| gives min{α_co|X₁|, α_coβ_co|X₁|/2}, so the correct second term should be α_coβ_co|X₁|/2, not α_co|X₁|/2.
\end{itemize}
\end{remark}

%--- Def_30: UnipotentSubgroupForLPS ---
\chapter{Def 30: Unipotent Subgroup for LPS Graphs}

\begin{definition}[Unipotent GL Element]
\label{def:LPS.unipotentGL}
\lean{LPS.unipotentGL}
\leanok
\uses{def:LPS.matToGL}
For a prime $q$ and $x \in \mathbb{Z}/q\mathbb{Z}$, the \emph{unipotent GL element} is the upper triangular unipotent matrix
\[
  \operatorname{unipotentGL}(q, x) = \begin{pmatrix} 1 & x \\ 0 & 1 \end{pmatrix} \in \operatorname{GL}(2, \mathbb{F}_q),
\]
with inverse $\begin{pmatrix} 1 & -x \\ 0 & 1 \end{pmatrix}$.
\end{definition}

\begin{lemma}[Unipotent GL Sends Zero to Identity]
\label{lem:LPS.unipotentGL_zero}
\lean{LPS.unipotentGL_zero}
\leanok
\uses{def:LPS.unipotentGL}
$\operatorname{unipotentGL}(q, 0) = 1 \in \operatorname{GL}(2, \mathbb{F}_q)$.
\end{lemma}

\begin{proof}
\leanok
\uses{def:LPS.unipotentGL}
By extensionality on $\operatorname{GL}$ elements, it suffices to check all entries $(i,j)$ with $i,j \in \{0,1\}$. Each entry is verified by simplification using the definition of $\operatorname{unipotentGL}$.
\end{proof}

\begin{lemma}[Unipotent GL is Additive]
\label{lem:LPS.unipotentGL_add}
\lean{LPS.unipotentGL_add}
\leanok
\uses{def:LPS.unipotentGL}
For $a, b \in \mathbb{Z}/q\mathbb{Z}$,
\[
  \operatorname{unipotentGL}(q, a + b) = \operatorname{unipotentGL}(q, a) \cdot \operatorname{unipotentGL}(q, b).
\]
\end{lemma}

\begin{proof}
\leanok
\uses{def:LPS.unipotentGL}
By extensionality on $\operatorname{GL}$ elements, it suffices to check all entries $(i,j)$ with $i,j \in \{0,1\}$. Each entry equality is verified by simplification using the definition of $\operatorname{unipotentGL}$, the matrix multiplication formula, and ring arithmetic.
\end{proof}

\begin{definition}[Unipotent PGL Element]
\label{def:LPS.unipotentPGL}
\lean{LPS.unipotentPGL}
\leanok
\uses{def:LPS.PGL2, def:LPS.projPGL, def:LPS.unipotentGL}
For $x \in \mathbb{Z}/q\mathbb{Z}$, the \emph{unipotent PGL element} is the image of $\operatorname{unipotentGL}(q,x)$ under the projection:
\[
  \operatorname{unipotentPGL}(q, x) = \operatorname{projPGL}(q)(\operatorname{unipotentGL}(q, x)) \in \operatorname{PGL}(2, \mathbb{F}_q).
\]
\end{definition}

\begin{definition}[Unipotent PGL Group Homomorphism]
\label{def:LPS.unipotentPGLHom}
\lean{LPS.unipotentPGLHom}
\leanok
\uses{def:LPS.PGL2, def:LPS.projPGL, def:LPS.unipotentPGL, lem:LPS.unipotentGL_zero, lem:LPS.unipotentGL_add}
The \emph{unipotent homomorphism} is the group homomorphism
\[
  \operatorname{unipotentPGLHom}(q) : \operatorname{Multiplicative}(\mathbb{Z}/q\mathbb{Z}) \to \operatorname{PGL}(2, \mathbb{F}_q),
\]
defined by $x \mapsto \operatorname{unipotentPGL}(q, \operatorname{toAdd}(x))$.

The map-one property holds because $\operatorname{unipotentGL}(q, 0) = 1$, and the map-mul property holds because $\operatorname{unipotentGL}(q, a+b) = \operatorname{unipotentGL}(q,a) \cdot \operatorname{unipotentGL}(q,b)$ and $\operatorname{projPGL}$ is a group homomorphism.
\end{definition}

\begin{lemma}[Injectivity of the Unipotent Map into PGL]
\label{lem:LPS.unipotentPGL_injective}
\lean{LPS.unipotentPGL_injective}
\leanok
\uses{def:LPS.PGL2, def:LPS.projPGL, def:LPS.matToGL, def:LPS.unipotentPGL, def:LPS.unipotentGL}
The map $x \mapsto \operatorname{unipotentPGL}(q, x)$ is injective.
\end{lemma}

\begin{proof}
\leanok
\uses{def:LPS.projPGL, def:LPS.matToGL, def:LPS.unipotentPGL, def:LPS.unipotentGL}
Let $a, b \in \mathbb{Z}/q\mathbb{Z}$ with $\operatorname{unipotentPGL}(q,a) = \operatorname{unipotentPGL}(q,b)$. Unfolding the definition of $\operatorname{projPGL}$ and using the quotient group equality, we obtain that
\[
  z := \operatorname{unipotentGL}(q,a)^{-1} \cdot \operatorname{unipotentGL}(q,b)
\]
lies in the center of $\operatorname{GL}(2, \mathbb{F}_q)$, i.e., $k \cdot z = z \cdot k$ for all $k$.

We compute $z.\operatorname{val} = \begin{pmatrix} 1 & b-a \\ 0 & 1 \end{pmatrix}$ by multiplying $\begin{pmatrix}1&-a\\0&1\end{pmatrix}\begin{pmatrix}1&b\\0&1\end{pmatrix}$ and simplifying.

Since $z$ is in the center, it commutes with $E_{21} = \operatorname{matToGL}\!\left(\begin{pmatrix}1&0\\1&1\end{pmatrix}\right)$. From the $(0,0)$-entry of the commutation identity $E_{21} \cdot z.\operatorname{val} = z.\operatorname{val} \cdot E_{21}$, we extract by simplification that $b - a = 0$, hence $a = b$.
\end{proof}

\begin{definition}[Unipotent Subgroup]
\label{def:LPS.unipotentSubgroup}
\lean{LPS.unipotentSubgroup}
\leanok
\uses{def:LPS.PGL2, def:LPS.unipotentPGLHom}
The \emph{unipotent subgroup} is the image of the unipotent homomorphism:
\[
  H := \operatorname{unipotentSubgroup}(q) = \operatorname{range}(\operatorname{unipotentPGLHom}(q)) \leq \operatorname{PGL}(2,\mathbb{F}_q).
\]
It consists of all elements of the form $\begin{pmatrix}1&x\\0&1\end{pmatrix}$ for $x \in \mathbb{F}_q$, and is isomorphic to $(\mathbb{F}_q, +) \cong \mathbb{Z}/q\mathbb{Z}$.
\end{definition}

\begin{lemma}[Membership in Unipotent Subgroup]
\label{lem:LPS.mem_unipotentSubgroup}
\lean{LPS.mem_unipotentSubgroup}
\leanok
\uses{def:LPS.unipotentSubgroup, def:LPS.unipotentPGL}
For $g \in \operatorname{PGL}(2,\mathbb{F}_q)$,
\[
  g \in \operatorname{unipotentSubgroup}(q) \iff \exists\, x \in \mathbb{Z}/q\mathbb{Z},\; \operatorname{unipotentPGL}(q,x) = g.
\]
\end{lemma}

\begin{proof}
\leanok
\uses{def:LPS.unipotentSubgroup, def:LPS.unipotentPGL, def:LPS.unipotentPGLHom}
This follows directly by simplification using the definitions of $\operatorname{unipotentSubgroup}$, $\operatorname{MonoidHom.mem_range}$, and $\operatorname{unipotentPGLHom_apply}$.
\end{proof}

\begin{theorem}[Cardinality of the Unipotent Subgroup]
\label{thm:LPS.unipotentSubgroup_card}
\lean{LPS.unipotentSubgroup_card}
\leanok
\uses{def:LPS.unipotentSubgroup, lem:LPS.unipotentPGL_injective, def:LPS.unipotentPGLHom}
$|\operatorname{unipotentSubgroup}(q)| = q$.
\end{theorem}

\begin{proof}
\leanok
\uses{def:LPS.unipotentSubgroup, lem:LPS.unipotentPGL_injective, def:LPS.unipotentPGLHom}
We first show $\operatorname{unipotentPGLHom}(q)$ is injective. Let $a, b$ with $\operatorname{unipotentPGLHom}(q)(a) = \operatorname{unipotentPGLHom}(q)(b)$. Unfolding the definition of $\operatorname{unipotentPGLHom}$, this gives $\operatorname{unipotentPGL}(q, \operatorname{toAdd}(a)) = \operatorname{unipotentPGL}(q, \operatorname{toAdd}(b))$, and by $\operatorname{unipotentPGL_injective}$ we get $\operatorname{toAdd}(a) = \operatorname{toAdd}(b)$, hence $a = b$.

Using injectivity, we construct a bijection
\[
  \operatorname{Multiplicative}(\mathbb{Z}/q\mathbb{Z}) \simeq \operatorname{unipotentSubgroup}(q)
\]
via $x \mapsto (\operatorname{unipotentPGLHom}(q)(x), \langle x, \operatorname{rfl}\rangle)$. Rewriting with this cardinality equivalence, and using $|\operatorname{Multiplicative}(\mathbb{Z}/q\mathbb{Z})| = |\mathbb{Z}/q\mathbb{Z}| = q$, we obtain $|\operatorname{unipotentSubgroup}(q)| = q$.
\end{proof}

\begin{theorem}[Odd Cardinality of the Unipotent Subgroup]
\label{thm:LPS.unipotentSubgroup_card_odd}
\lean{LPS.unipotentSubgroup_card_odd}
\leanok
\uses{def:LPS.unipotentSubgroup, thm:LPS.unipotentSubgroup_card}
If $q$ is an odd prime (i.e., $q \equiv 1 \pmod{2}$), then $|\operatorname{unipotentSubgroup}(q)|$ is odd.
\end{theorem}

\begin{proof}
\leanok
\uses{thm:LPS.unipotentSubgroup_card}
Rewriting with $\operatorname{unipotentSubgroup_card}$, the goal becomes showing $q$ is odd. Rewriting with $\operatorname{Nat.odd_iff}$, this is exactly the hypothesis $q \% 2 = 1$.
\end{proof}

\begin{lemma}[Determinant of a Unipotent GL Element]
\label{lem:LPS.unipotentGL_det}
\lean{LPS.unipotentGL_det}
\leanok
\uses{def:LPS.unipotentGL}
For $x \in \mathbb{Z}/q\mathbb{Z}$, $\det(\operatorname{unipotentGL}(q,x).\operatorname{val}) = 1$.
\end{lemma}

\begin{proof}
\leanok
\uses{def:LPS.unipotentGL}
By simplification using the $2 \times 2$ determinant formula and the definition of $\operatorname{unipotentGL}$: $\det\begin{pmatrix}1&x\\0&1\end{pmatrix} = 1\cdot 1 - x \cdot 0 = 1$.
\end{proof}

\begin{lemma}[LPS Matrix Determinant Equals $p$]
\label{lem:LPS.lpsMatrixVal_det_eq_p}
\lean{LPS.lpsMatrixVal_det_eq_p}
\leanok
\uses{def:LPS.SpTilde, def:LPS.lpsMatrixVal, lem:LPS.lpsMatrixVal_det}
Let $p \neq q$ be primes, $x, y \in \mathbb{Z}/q\mathbb{Z}$ with $x^2 + y^2 + 1 = 0$, and $t \in \widetilde{S}_p$. Then
\[
  \det(\operatorname{lpsMatrixVal}(q, x, y, t)) = p \in \mathbb{Z}/q\mathbb{Z}.
\]
\end{lemma}

\begin{proof}
\leanok
\uses{def:LPS.SpTilde, def:LPS.lpsMatrixVal, lem:LPS.lpsMatrixVal_det}
Rewriting with $\operatorname{lpsMatrixVal_det}$, the determinant equals the sum of squares of the components of $t$. Since $t \in \widetilde{S}_p$, we have $t_0^2 + t_1^2 + t_2^2 + t_3^2 = p$ as integers (from the definition of $\widetilde{S}_p$). Applying the ring homomorphism $\mathbb{Z} \to \mathbb{Z}/q\mathbb{Z}$ and pushing the cast through, we get $\det = p \in \mathbb{Z}/q\mathbb{Z}$.
\end{proof}

\begin{lemma}[Determinant of a Tuple-to-GL Element]
\label{lem:LPS.tupleToGLElement_det}
\lean{LPS.tupleToGLElement_det}
\leanok
\uses{def:LPS.tupleToGLElement, def:LPS.SpTilde, lem:LPS.lpsMatrixVal_det_eq_p}
Under the same hypotheses as above, $\det(\operatorname{tupleToGLElement}(q, p, \ldots, t, \cdot).\operatorname{val}) = p \in \mathbb{Z}/q\mathbb{Z}$.
\end{lemma}

\begin{proof}
\leanok
\uses{def:LPS.tupleToGLElement, lem:LPS.lpsMatrixVal_det_eq_p}
Rewriting with $\operatorname{tupleToGLElement_val}$, the value equals $\operatorname{lpsMatrixVal}(q,x,y,t)$, and the result follows from $\operatorname{lpsMatrixVal_det_eq_p}$.
\end{proof}

\begin{lemma}[Unipotent Subgroup Elements Have Square Determinant Lift]
\label{lem:LPS.unipotentSubgroup_det_is_square}
\lean{LPS.unipotentSubgroup_det_is_square}
\leanok
\uses{def:LPS.unipotentSubgroup, def:LPS.unipotentPGL, lem:LPS.mem_unipotentSubgroup}
If $g \in \operatorname{unipotentSubgroup}(q)$, then there exists $x \in \mathbb{Z}/q\mathbb{Z}$ such that $\operatorname{unipotentPGL}(q,x) = g$.
\end{lemma}

\begin{proof}
\leanok
\uses{lem:LPS.mem_unipotentSubgroup}
This follows directly by rewriting with $\operatorname{mem_unipotentSubgroup}$: the membership condition $g \in \operatorname{unipotentSubgroup}(q)$ is equivalent to the existence of such $x$.
\end{proof}

\begin{lemma}[$p$ Does Not Vanish Mod $q$]
\label{lem:LPS.natCast_prime_ne_zero}
\lean{LPS.natCast_prime_ne_zero}
\leanok
\uses{def:LPS.PGL2}
If $p \neq q$ are primes, then $(p : \mathbb{Z}/q\mathbb{Z}) \neq 0$.
\end{lemma}

\begin{proof}
\leanok
\uses{def:LPS.PGL2}
Suppose $(p : \mathbb{Z}/q\mathbb{Z}) = 0$. By $\operatorname{ZMod.natCast_eq_zero_iff}$, this means $q \mid p$. Since $q$ is prime and $p$ is prime, the divisibility $q \mid p$ implies $q = p$ (as $p$ is prime, its only divisors are $1$ and $p$, and $q > 1$), contradicting $p \neq q$.
\end{proof}

\begin{lemma}[$p$ is Not a Square When Legendre Symbol is $-1$]
\label{lem:LPS.not_isSquare_of_legendreSym_neg_one}
\lean{LPS.not_isSquare_of_legendreSym_neg_one}
\leanok
\uses{def:LPS.PGL2}
If $\operatorname{legendreSym}(q, p) = -1$, then $(p : \mathbb{Z}/q\mathbb{Z})$ is not a perfect square.
\end{lemma}

\begin{proof}
\leanok
\uses{def:LPS.PGL2}
Rewriting with $\operatorname{jacobiSym.legendreSym.to_jacobiSym}$, the condition $\operatorname{legendreSym}(q, p) = -1$ becomes $\operatorname{jacobiSym}(p, q) = -1$. The result then follows from $\operatorname{ZMod.nonsquare_of_jacobiSym_eq_neg_one}$.
\end{proof}

\begin{lemma}[Center of $\operatorname{GL}(2, \mathbb{F}_q)$ Has Square Determinant]
\label{lem:LPS.center_GL2_det_isSquare}
\lean{LPS.center_GL2_det_isSquare}
\leanok
\uses{def:LPS.matToGL}
If $z \in \operatorname{GL}(2, \mathbb{F}_q)$ lies in the center of $\operatorname{GL}(2, \mathbb{F}_q)$, then $\det(z.\operatorname{val})$ is a perfect square in $\mathbb{F}_q$.
\end{lemma}

\begin{proof}
\leanok
\uses{def:LPS.matToGL}
Since $z$ lies in the center, it commutes with all elements of $\operatorname{GL}(2,\mathbb{F}_q)$. Let
\[
  E_{12} = \operatorname{matToGL}\!\left(\begin{pmatrix}1&1\\0&1\end{pmatrix}\right), \quad E_{21} = \operatorname{matToGL}\!\left(\begin{pmatrix}1&0\\1&1\end{pmatrix}\right).
\]
From $E_{12} \cdot z.\operatorname{val} = z.\operatorname{val} \cdot E_{12}$, we extract by simplification the $(1,1)$-entry equation, giving $z_{10} = 0$, and the $(0,1)$-entry equation, giving $z_{00} = z_{11}$. From $E_{21} \cdot z.\operatorname{val} = z.\operatorname{val} \cdot E_{21}$, we extract the $(0,0)$-entry equation giving $z_{01} = 0$. Thus $z.\operatorname{val}$ is a scalar matrix $\begin{pmatrix}c&0\\0&c\end{pmatrix}$, and $\det(z.\operatorname{val}) = c^2 = (z_{00})^2$, which is a perfect square.
\end{proof}

\begin{theorem}[Generating Set Disjoint from Unipotent Subgroup]
\label{thm:LPS.lpsGenSet_disjoint_unipotentSubgroup}
\lean{LPS.lpsGenSet_disjoint_unipotentSubgroup}
\leanok
\uses{def:CayleyGraph.SymmGenSet, def:LPS.PGL2, def:LPS.SpTilde, def:LPS.lpsGenSetPGL, def:LPS.lpsMatrixVal, def:LPS.matToGL, def:LPS.projPGL, def:LPS.tupleToGLElement, lem:LPS.Sp_subset_SpTilde, lem:LPS.lpsMatrixVal_det, def:LPS.unipotentSubgroup, lem:LPS.natCast_prime_ne_zero, lem:LPS.not_isSquare_of_legendreSym_neg_one}
Let $p, q$ be distinct odd primes with $(q : \mathbb{R}) > 2\sqrt{p}$, $x,y \in \mathbb{F}_q$ satisfying $x^2 + y^2 + 1 = 0$, and $\operatorname{legendreSym}(q, p) = -1$. Then $S_{p,q} \cap H = \emptyset$, i.e., no generator lies in the unipotent subgroup $H$.
\end{theorem}

\begin{proof}
\leanok
\uses{def:CayleyGraph.SymmGenSet, def:LPS.lpsGenSetPGL, def:LPS.tupleToGLElement, def:LPS.matToGL, lem:LPS.Sp_subset_SpTilde, lem:LPS.natCast_prime_ne_zero, lem:LPS.not_isSquare_of_legendreSym_neg_one, lem:LPS.unipotentPGL_injective, lem:LPS.lpsMatrixVal_det_eq_p, lem:LPS.unipotentGL_det}
Suppose for contradiction that $s \in S_{p,q}$ and $s \in H$. From $s \in H$, by $\operatorname{mem_unipotentSubgroup}$, we obtain $a \in \mathbb{F}_q$ with $\operatorname{unipotentPGL}(q,a) = s$. From $s \in S_{p,q}$, we obtain $t \in S_p$ (a tuple) such that $\operatorname{projPGL}(\operatorname{tupleToGLElement}(q,p,\ldots,t)) = s$.

By $\operatorname{Sp_subset_SpTilde}$, $t \in \widetilde{S}_p$. Set
\[
  z := \operatorname{tupleToGLElement}(\ldots, t)^{-1} \cdot \operatorname{unipotentGL}(q, a).
\]
From the equality in $\operatorname{PGL}$, $z$ lies in the center of $\operatorname{GL}(2, \mathbb{F}_q)$. Let $E_{12}$ and $E_{21}$ be the elementary matrices from $\operatorname{matToGL}$. The commutation of $z$ with $E_{12}$ and $E_{21}$ (extracted via the $(1,1)$-, $(0,1)$- and $(0,0)$-entries respectively) shows $z_{10} = 0$, $z_{00} = z_{11}$, and $z_{01} = 0$, so $\det(z.\operatorname{val}) = (z_{00})^2$.

For the determinant of $z$: by $\operatorname{lpsMatrixVal_det_eq_p}$ and $\operatorname{unipotentGL_det}$, we compute $\det(z.\operatorname{val}) = (p : \mathbb{F}_q)^{-1}$. Thus $(p : \mathbb{F}_q)^{-1}$ is a perfect square, so by $\operatorname{isSquare_inv}$, $(p : \mathbb{F}_q)$ is a perfect square. However, by $\operatorname{not_isSquare_of_legendreSym_neg_one}$, $(p : \mathbb{F}_q)$ is not a square --- contradiction.
\end{proof}

\begin{theorem}[Generating Set Disjoint from All Conjugates of Unipotent Subgroup]
\label{thm:LPS.lpsGenSet_disjoint_conjugate}
\lean{LPS.lpsGenSet_disjoint_conjugate}
\leanok
\uses{def:CayleyGraph.SymmGenSet, def:LPS.PGL2, def:LPS.SpTilde, def:LPS.lpsGenSetPGL, def:LPS.lpsMatrixVal, def:LPS.matToGL, def:LPS.projPGL, def:LPS.tupleToGLElement, lem:LPS.Sp_subset_SpTilde, lem:LPS.lpsMatrixVal_det, def:LPS.unipotentSubgroup, lem:LPS.natCast_prime_ne_zero, lem:LPS.not_isSquare_of_legendreSym_neg_one, lem:LPS.center_GL2_det_isSquare}
Under the same hypotheses as the previous theorem, for any $g \in \operatorname{PGL}(2,\mathbb{F}_q)$, $h \in H$, and $s \in S_{p,q}$, we have $s \neq g h g^{-1}$.
\end{theorem}

\begin{proof}
\leanok
\uses{def:LPS.lpsGenSetPGL, def:LPS.tupleToGLElement, def:LPS.matToGL, lem:LPS.Sp_subset_SpTilde, lem:LPS.natCast_prime_ne_zero, lem:LPS.not_isSquare_of_legendreSym_neg_one, lem:LPS.center_GL2_det_isSquare, lem:LPS.lpsMatrixVal_det_eq_p, lem:LPS.unipotentGL_det}
Suppose $s = g h g^{-1}$. From $h \in H$, write $h = \operatorname{unipotentPGL}(q,a)$ for some $a$. From $s \in S_{p,q}$, obtain $t \in S_p$ with $\operatorname{projPGL}(\operatorname{tupleToGLElement}(\ldots,t)) = s$.

Lift $g$ to a GL element $G$ via the surjectivity of $\operatorname{projPGL}$. Then the equality becomes $\operatorname{projPGL}(\operatorname{tupleToGLElement}(\ldots,t)) = \operatorname{projPGL}(G \cdot \operatorname{unipotentGL}(q,a) \cdot G^{-1})$ in $\operatorname{PGL}$. This places $z := \operatorname{tupleToGLElement}(\ldots,t)^{-1} \cdot (G \cdot \operatorname{unipotentGL}(q,a) \cdot G^{-1})$ in the center of $\operatorname{GL}(2, \mathbb{F}_q)$.

By $\operatorname{center_GL2_det_isSquare}$, $\det(z.\operatorname{val})$ is a perfect square. We compute $\det(z.\operatorname{val})$: since $\det(G \cdot \operatorname{unipotentGL}(q,a) \cdot G^{-1}) = \det(G) \cdot 1 \cdot \det(G^{-1}) = 1$, and $\det(\operatorname{tupleToGLElement}(\ldots,t)) = p$ by $\operatorname{lpsMatrixVal_det_eq_p}$, we get $\det(z.\operatorname{val}) = p^{-1}$. By $\operatorname{isSquare_inv}$, $(p : \mathbb{F}_q)$ is a perfect square, contradicting $\operatorname{not_isSquare_of_legendreSym_neg_one}$.
\end{proof}

\begin{lemma}[Satisfiability of Disjoint Conjugate Premises]
\label{lem:LPS.lpsGenSet_disjoint_conjugate_satisfiable}
\lean{LPS.lpsGenSet_disjoint_conjugate_satisfiable}
\leanok
\uses{def:LPS.lpsGenSetPGL}
There exist distinct odd primes $p, q$ satisfying the premises of $\operatorname{lpsGenSet_disjoint_conjugate}$.
\end{lemma}

\begin{proof}
\leanok
\uses{def:LPS.lpsGenSetPGL}
This follows directly from $\operatorname{lpsGenSetPGL_satisfiable}$, which provides a witness.
\end{proof}

\begin{lemma}[Unipotent Subgroup is Nonempty]
\label{lem:LPS.unipotentSubgroup_nonempty}
\lean{LPS.unipotentSubgroup_nonempty}
\leanok
\uses{def:LPS.unipotentSubgroup}
The unipotent subgroup $\operatorname{unipotentSubgroup}(q)$ is nonempty as a subset of $\operatorname{PGL}(2, \mathbb{F}_q)$.
\end{lemma}

\begin{proof}
\leanok
\uses{def:LPS.unipotentSubgroup}
The identity element $1 \in \operatorname{PGL}(2, \mathbb{F}_q)$ belongs to $\operatorname{unipotentSubgroup}(q)$ since every subgroup contains the identity.
\end{proof}

%--- Thm_15: ExplicitFamilyQuantumCodes ---
\chapter{Thm 15: Explicit Family of Quantum Codes}

\begin{definition}[Valid Prime]
\label{def:ExplicitFamilyQuantumCodes.ValidPrime}
\lean{ExplicitFamilyQuantumCodes.ValidPrime}
\leanok
\uses{def:HasBoundedWeight}
A \emph{valid prime} for the LPS construction with $p = 401$ is a structure consisting of:
\begin{itemize}
  \item A prime $q \in \mathbb{N}$ with $q$ odd, $q \neq 401$, and $q > 2\sqrt{401}$,
  \item Elements $x, y \in \mathbb{Z}/q\mathbb{Z}$ satisfying $x^2 + y^2 + 1 = 0$.
\end{itemize}
The condition $x^2 + y^2 + 1 = 0$ witnesses that $-1$ is not a quadratic residue mod $q$, i.e., the Legendre symbol $(p/q) = -1$.
\end{definition}

\begin{definition}[Code Length]
\label{def:ExplicitFamilyQuantumCodes.codeLength}
\lean{ExplicitFamilyQuantumCodes.codeLength}
\leanok
\uses{def:ExplicitFamilyQuantumCodes.ValidPrime}
For a valid prime $q$, the code length is
\[
  N = 3 \cdot 201 \cdot q(q^2 - 1).
\]
This equals $3|X_1|$ where $|X_1| = \tfrac{s}{2} \cdot |\operatorname{PGL}(2,q)| = 201 \cdot q(q^2-1)$ is the number of edges of the LPS graph, obtained by the handshaking lemma for $s$-regular graphs with $|X_0| = |\operatorname{PGL}(2,q)| = q(q^2-1)$ vertices and $s = 402$.
\end{definition}

\begin{definition}[LDPC Weight Bound]
\label{def:ExplicitFamilyQuantumCodes.ldpcWeightBound}
\lean{ExplicitFamilyQuantumCodes.ldpcWeightBound}
\leanok
\uses{def:HasBoundedWeight}
The LDPC weight bound is $w = 2s = 804$.
\end{definition}

\begin{definition}[LPS Code Data]
\label{def:ExplicitFamilyQuantumCodes.LPSCodeData}
\lean{ExplicitFamilyQuantumCodes.LPSCodeData}
\leanok
\uses{def:ExplicitFamilyQuantumCodes.ValidPrime, def:ExplicitFamilyQuantumCodes.codeLength, def:SubsystemCSSCode}
For a valid prime $q$, the \emph{LPS code data} is a structure consisting of:
\begin{itemize}
  \item The number of $X$-type check rows $r_X \in \mathbb{N}$,
  \item The number of $Z$-type check rows $r_Z \in \mathbb{N}$,
  \item A subsystem CSS code of length $N = \operatorname{codeLength}(q)$.
\end{itemize}
\end{definition}

\begin{definition}[LPS Construction Data]
\label{def:ExplicitFamilyQuantumCodes.LPSConstructionData}
\lean{ExplicitFamilyQuantumCodes.LPSConstructionData}
\leanok
\uses{def:ExplicitFamilyQuantumCodes.ValidPrime, def:ExplicitFamilyQuantumCodes.codeLength, def:ExplicitFamilyQuantumCodes.ldpcWeightBound, def:SubsystemCSSCode, def:HasBoundedWeight}
For a valid prime $q$, the \emph{LPS construction data} is a structure consisting of:
\begin{itemize}
  \item The number of $X$-type check rows $r_X \in \mathbb{N}$,
  \item The number of $Z$-type check rows $r_Z \in \mathbb{N}$,
  \item A subsystem CSS code of length $N = \operatorname{codeLength}(q)$,
  \item A proof that the $X$-type parity check matrix $H_X$ has bounded weight $\leq 2s = 804$ (LDPC property),
  \item A proof that the $Z$-type parity check matrix $H_{ZT}$ has bounded weight $\leq 2s = 804$ (LDPC property).
\end{itemize}
The bound $2s = 804$ arises from the Kronecker product structure: the Tanner code differential has weight $\leq s$ (from $s$-regularity) and the cycle graph differential has weight $2$, giving total weight $\leq 2s$ after the balanced product quotient.
\end{definition}

\begin{theorem}[Axiom: Balanced Product Subsystem CSS Code]
\label{thm:ExplicitFamilyQuantumCodes.balanced_product_code}
\lean{ExplicitFamilyQuantumCodes.balanced_product_code}
\leanok
\uses{def:ExplicitFamilyQuantumCodes.ValidPrime, def:ExplicitFamilyQuantumCodes.LPSCodeData}

\textbf{\color{red}[UNPROVEN AXIOM]}

For every valid prime $q$, the balanced product $C(X, L) \otimes_{\mathbb{Z}_q} C(C_q)$ yields a subsystem CSS code of length $N = \operatorname{codeLength}(q)$, with the horizontal/vertical homology splitting providing the logical/gauge decomposition.

\textit{Note: This is stated as an axiom citing Panteleev--Kalachev \cite{PK22}, Def.~26 and Def.~29. The full formalization of the balanced product construction was not completed.}
\end{theorem}

\begin{theorem}[Axiom: LDPC Property of Balanced Product Codes]
\label{thm:ExplicitFamilyQuantumCodes.balanced_product_ldpc}
\lean{ExplicitFamilyQuantumCodes.balanced_product_ldpc}
\leanok
\uses{def:ExplicitFamilyQuantumCodes.ValidPrime, def:ExplicitFamilyQuantumCodes.ldpcWeightBound, def:HasBoundedWeight, thm:ExplicitFamilyQuantumCodes.balanced_product_code}

\textbf{\color{red}[UNPROVEN AXIOM]}

For every valid prime $q$, the parity check matrices of the balanced product code have bounded weight:
\[
  H_X \text{ has bounded weight} \leq 2s = 804, \quad H_{ZT} \text{ has bounded weight} \leq 2s = 804.
\]

\textit{Note: This is stated as an axiom citing Panteleev--Kalachev \cite{PK22}, \S4.3, Step~8. The bound $2s = 804$ follows from the Kronecker product structure.}
\end{theorem}

\begin{theorem}[Axiom: Infinitely Many Valid Primes]
\label{thm:ExplicitFamilyQuantumCodes.infinitelyManyValidPrimes}
\lean{ExplicitFamilyQuantumCodes.infinitelyManyValidPrimes}
\leanok
\uses{def:ExplicitFamilyQuantumCodes.ValidPrime}

\textbf{\color{red}[UNPROVEN AXIOM]}

For any $M \in \mathbb{N}$, there exists a valid prime $q > M$.

\textit{Note: This is stated as an axiom appealing to Dirichlet's theorem on primes in arithmetic progressions (Dirichlet 1837) and quadratic reciprocity for the Legendre symbol condition. The required Dirichlet theorem infrastructure was not formalized from scratch.}
\end{theorem}

\begin{theorem}[Axiom: Logical Qubit Lower Bound]
\label{thm:ExplicitFamilyQuantumCodes.lps_K_lower_bound}
\lean{ExplicitFamilyQuantumCodes.lps_K_lower_bound}
\leanok
\uses{def:ExplicitFamilyQuantumCodes.ValidPrime, def:ExplicitFamilyQuantumCodes.LPSConstructionData, def:SubsystemCSSCode.logicalQubits}

\textbf{\color{red}[UNPROVEN AXIOM]}

There exists a constant $c > 0$ such that for every valid prime $q$, letting $D$ denote the balanced product construction data,
\[
  c \cdot (q^2 - 1) \leq K,
\]
where $K = \operatorname{logicalQubits}(D.\operatorname{code})$ is the number of logical qubits.

\textit{Note: This is stated as an axiom citing Sipser--Spielman 1996 (Theorem~7 in this work). Since $k_L > s/2 = 201$ and $|X_1|/\ell = 201(q^2-1)$, one obtains $K \geq c(q^2-1)$ for a constant $c > 0$.}
\end{theorem}

\begin{theorem}[Axiom: Logical Qubit Upper Bound]
\label{thm:ExplicitFamilyQuantumCodes.lps_K_upper_bound}
\lean{ExplicitFamilyQuantumCodes.lps_K_upper_bound}
\leanok
\uses{def:ExplicitFamilyQuantumCodes.ValidPrime, def:ExplicitFamilyQuantumCodes.LPSConstructionData, def:SubsystemCSSCode.logicalQubits}

\textbf{\color{red}[UNPROVEN AXIOM]}

There exists a constant $C > 0$ such that for every valid prime $q$, letting $D$ denote the balanced product construction data,
\[
  K \leq C \cdot (q^2 - 1),
\]
where $K = \operatorname{logicalQubits}(D.\operatorname{code})$.

\textit{Note: This is stated as an axiom appealing to the Rank--Nullity Theorem (Axler, \textit{Linear Algebra Done Right}, Thm.~3.22). The total chain module $\operatorname{Tot}_1$ of the balanced product has dimension $\in \Theta(q^3)$, giving $K = \dim(H_1) \leq \dim(\operatorname{Tot}_1) \in O(q^2)$ after dividing by $\ell = q$.}
\end{theorem}

\begin{theorem}[Axiom: Z-Distance Lower Bound]
\label{thm:ExplicitFamilyQuantumCodes.lps_DZ_lower_bound}
\lean{ExplicitFamilyQuantumCodes.lps_DZ_lower_bound}
\leanok
\uses{def:ExplicitFamilyQuantumCodes.ValidPrime, def:ExplicitFamilyQuantumCodes.LPSConstructionData, def:SubsystemCSSCode.dZ}

\textbf{\color{red}[UNPROVEN AXIOM]}

There exists a constant $c > 0$ such that for every valid prime $q$, letting $D$ denote the balanced product construction data,
\[
  c \cdot q(q^2 - 1) \leq D_Z,
\]
where $D_Z = D.\operatorname{code}.\operatorname{dZ}$ is the $Z$-distance of the subsystem code.

\textit{Note: This is stated as an axiom citing Panteleev--Kalachev \cite{PK22}, Theorem~13 (homological distance bound). Since $|X_1| = 201 \cdot q(q^2-1)$, the bound gives $D_Z \geq c \cdot q(q^2-1)$ for a constant $c > 0$ determined by the expansion parameters $\alpha_{ho} = 10^{-3}$ and $\beta_{ho} > 0$ from Theorem~8.}
\end{theorem}

\begin{theorem}[Axiom: X-Distance Lower Bound]
\label{thm:ExplicitFamilyQuantumCodes.lps_DX_lower_bound}
\lean{ExplicitFamilyQuantumCodes.lps_DX_lower_bound}
\leanok
\uses{def:ExplicitFamilyQuantumCodes.ValidPrime, def:ExplicitFamilyQuantumCodes.LPSConstructionData, def:SubsystemCSSCode.dX}

\textbf{\color{red}[UNPROVEN AXIOM]}

There exists a constant $c > 0$ such that for every valid prime $q$, letting $D$ denote the balanced product construction data,
\[
  D_X \geq c \cdot q,
\]
where $D_X = D.\operatorname{code}.\operatorname{dX}$ is the $X$-distance of the subsystem code.

\textit{Note: This is stated as an axiom citing Panteleev--Kalachev \cite{PK22}, Theorem~14 (cohomological distance bound). Using $|X_0|s = 2|X_1| \in \Theta(q^3)$ and $\ell = q$, the minimum over horizontal and vertical contributions is $\Theta(q)$. Expansion parameters $\alpha_{co} = 10^{-5}$ and $\beta_{co} > 0$ come from Theorem~9.}
\end{theorem}

\begin{theorem}[Explicit Family of Quantum Codes]
\label{thm:ExplicitFamilyQuantumCodes.explicitFamilyQuantumCodes}
\lean{ExplicitFamilyQuantumCodes.explicitFamilyQuantumCodes}
\leanok
\uses{def:HasBoundedWeight}

\textbf{\color{orange}[Uses unproven axioms: thm:ExplicitFamilyQuantumCodes.balanced\_product\_code, thm:ExplicitFamilyQuantumCodes.balanced\_product\_ldpc, thm:ExplicitFamilyQuantumCodes.infinitelyManyValidPrimes, thm:ExplicitFamilyQuantumCodes.lps\_K\_lower\_bound, thm:ExplicitFamilyQuantumCodes.lps\_K\_upper\_bound, thm:ExplicitFamilyQuantumCodes.lps\_DZ\_lower\_bound, thm:ExplicitFamilyQuantumCodes.lps\_DX\_lower\_bound]}

There exists an explicit construction of an infinite family of $[\![N, K, D_X, D_Z]\!]$ LDPC subsystem CSS quantum codes satisfying the following simultaneously:
\begin{enumerate}
  \item \textbf{(Code existence)} For each valid prime $q$, there exist row counts $r_X, r_Z \in \mathbb{N}$ and a subsystem CSS code of length $N = 3 \cdot 201 \cdot q(q^2-1)$ whose parity check matrices $H_X$ and $H_{ZT}$ both have bounded weight $\leq 804$.
  \item \textbf{(Infinite family)} For any $M \in \mathbb{N}$, there exists a valid prime $q$ with $N > M$.
  \item \textbf{(Logical qubit count)} There exist constants $0 < c \leq C$ such that for every valid prime $q$,
  \[
    c(q^2 - 1) \leq K \leq C(q^2 - 1),
  \]
  i.e., $K \in \Theta(q^2) = \Theta(N^{2/3})$.
  \item \textbf{(X-distance)} There exists a constant $c > 0$ such that for every valid prime $q$,
  \[
    D_X \geq c \cdot q = \Omega(N^{1/3}).
  \]
  \item \textbf{(Z-distance)} There exist constants $0 < c \leq C$ such that for every valid prime $q$,
  \[
    c \cdot N \leq D_Z \leq C \cdot N,
  \]
  i.e., $D_Z \in \Theta(N)$.
\end{enumerate}
\end{theorem}

\begin{proof}
\leanok
\uses{def:HasBoundedWeight, thm:ExplicitFamilyQuantumCodes.balanced_product_code, thm:ExplicitFamilyQuantumCodes.balanced_product_ldpc, thm:ExplicitFamilyQuantumCodes.infinitelyManyValidPrimes, thm:ExplicitFamilyQuantumCodes.lps_K_lower_bound, thm:ExplicitFamilyQuantumCodes.lps_K_upper_bound, thm:ExplicitFamilyQuantumCodes.lps_DZ_lower_bound, thm:ExplicitFamilyQuantumCodes.lps_DX_lower_bound}

\textbf{\color{orange}[Uses unproven axioms as listed above.]}

We prove each part in turn using the \texttt{refine} tactic to split the conjunction.

\textbf{Part 1 (LDPC subsystem CSS code).} Let $q$ be any valid prime. Define $D = \operatorname{balanced_product_construction}(q)$, which packages the axioms \texttt{balanced\_product\_code} and \texttt{balanced\_product\_ldpc} into a single structure. We extract $D.r_X$, $D.r_Z$, $D.\operatorname{code}$, $D.\operatorname{isLDPC_HX}$, and $D.\operatorname{isLDPC_HZT}$ to satisfy Part~1.

\textbf{Part 2 (Infinite family).} This is exactly the content of the lemma \texttt{codeLength\_unbounded}, which is proved as follows. Let $M \in \mathbb{N}$ be arbitrary. By the axiom \texttt{infinitelyManyValidPrimes}, there exists a valid prime $q > M$. Since $q \geq 41$ (proved from the lower bound $q > 2\sqrt{401} \geq 40$), and $q^2 - 1 \geq 1$, we compute:
\[
  N = 3 \cdot 201 \cdot q(q^2-1) \geq 3 \cdot 201 \cdot q \cdot 1 \geq q > M.
\]

\textbf{Part 3 ($K \in \Theta(q^2)$).} From the axiom \texttt{lps\_K\_lower\_bound}, obtain $c > 0$ and the bound $c(q^2-1) \leq K$ for all valid primes. From the axiom \texttt{lps\_K\_upper\_bound}, obtain $C > 0$ and the bound $K \leq C(q^2-1)$ for all valid primes. Taking both together gives Part~3.

\textbf{Part 4 ($D_X \in \Omega(q)$).} This follows directly from the axiom \texttt{lps\_DX\_lower\_bound}, which provides $c > 0$ with $D_X \geq c \cdot q$ for all valid primes.

\textbf{Part 5 ($D_Z \in \Theta(N)$).} From the axiom \texttt{lps\_DZ\_lower\_bound}, obtain $c > 0$ and the bound $c \cdot q(q^2-1) \leq D_Z$ for all valid primes. We set $c' = c/(3 \cdot 201)$ and $C' = 1$.

For the lower bound: since $N = 3 \cdot 201 \cdot q(q^2-1)$ (the lemma \texttt{codeLength\_eq\_real}), we compute
\[
  \frac{c}{3 \cdot 201} \cdot N = \frac{c}{3 \cdot 201} \cdot 3 \cdot 201 \cdot q(q^2-1) = c \cdot q(q^2-1) \leq D_Z.
\]
This is verified by simplifying with \texttt{field\_simp} and applying linear arithmetic.

For the upper bound: $D_Z \leq N$ is the content of the lemma \texttt{dZ\_le\_codeLength}. Indeed, $D_Z$ is the infimum of a set of Hamming weights of nonzero logical $Z$ codewords; if this set is empty, then $D_Z = 0 \leq N$; otherwise any element $w$ of the set satisfies $w = \operatorname{hammingWeight}(x) \leq n = N$ (by \texttt{SipserSpielman.hammingWeight\_le\_card}), so $D_Z \leq w \leq N$. Hence $D_Z \leq 1 \cdot N$.
\end{proof}

\begin{remark}[Paper Corrections]
\textbf{\color{red}Errata.} The following errors were identified in the original paper and corrected in this formalization:
\begin{itemize}
  \item In Corollary (cor:distanceboundssybsystemcode), the D_X formula has second term α_co|X_1|/2 but should be α_co·β_co·|X_1|/2, derived from Theorem distco Case 1: |X_0|s·α_co·β_co/4 = 2|X_1|·α_co·β_co/4 = α_co·β_co·|X_1|/2. The β_co factor was dropped. This is a typo that does not affect the asymptotic result since the minimum is dominated by the ℓ·α_co·β_co/(4s) = Θ(q) terms, not the Θ(q³) terms.
\end{itemize}
\end{remark}

%--- Cor_3: DistanceBalancedFamily ---
\chapter{Cor 3: Distance-Balanced Family of Quantum Codes}

\begin{definition}[Classical Code Length]
\label{def:DistanceBalancedFamily.classicalCodeLength}
\lean{DistanceBalancedFamily.classicalCodeLength}
\leanok
\uses{def:ExplicitFamilyQuantumCodes.ValidPrime}

For a valid prime $q$, the classical code length is $n_c = q^2$.
\end{definition}

\begin{definition}[Balanced Code Length]
\label{def:DistanceBalancedFamily.balancedCodeLength}
\lean{DistanceBalancedFamily.balancedCodeLength}
\leanok
\uses{def:ExplicitFamilyQuantumCodes.ValidPrime, def:ExplicitFamilyQuantumCodes.codeLength, def:DistanceBalancedFamily.classicalCodeLength}

For a valid prime $q$, the balanced code length is $N = N_0 \cdot n_c = \operatorname{codeLength}(vp) \cdot q^2$, where $N_0 = \operatorname{codeLength}(vp)$ is the subsystem code length from Theorem~15.
\end{definition}

\begin{definition}[Balanced Code Data]
\label{def:DistanceBalancedFamily.BalancedCodeData}
\lean{DistanceBalancedFamily.BalancedCodeData}
\leanok
\uses{def:ExplicitFamilyQuantumCodes.ValidPrime, def:DistanceBalancedFamily.balancedCodeLength, def:CSSCode, def:HasBoundedWeight}

For a valid prime $vp$, a \emph{balanced code datum} is a structure bundling:
\begin{itemize}
  \item row counts $r_X, r_Z \in \mathbb{N}$ for the X- and Z-type parity checks,
  \item a CSS code $\mathcal{C}$ of block length $N = \operatorname{balancedCodeLength}(vp)$,
  \item a weight bound $w \in \mathbb{N}$ such that the X-parity check matrix $H_X$ and the transposed Z-parity check matrix $H_{ZT}$ both have bounded row and column weight at most $w$ (LDPC property),
  \item a lower bound on logical qubits: $K \geq K_0 \cdot \lfloor n_c/2 \rfloor + 1$, where $K_0$ is the logical qubit count of the subsystem code from Theorem~15 and $n_c = q^2$ is the classical code length,
  \item a lower bound on X-distance: $D_X \geq D_{X,0} \cdot \lfloor n_c / 10 \rfloor$, where $D_{X,0}$ is the X-distance of the subsystem code,
  \item an upper bound on logical qubits: $K \leq K_0 \cdot n_c$,
  \item a lower bound on Z-distance: $D_Z \geq D_{Z,0}$, where $D_{Z,0}$ is the Z-distance of the subsystem code.
\end{itemize}
\end{definition}

\begin{theorem}[Axiom: Evra--Kaufman--Zemor Distance Balancing]
\label{thm:DistanceBalancedFamily.ekz_distance_balancing}
\lean{DistanceBalancedFamily.ekz_distance_balancing}
\leanok
\uses{def:ExplicitFamilyQuantumCodes.ValidPrime, def:DistanceBalancedFamily.BalancedCodeData}

\textbf{\color{red}[UNPROVEN AXIOM]}

For every valid prime $vp$, there exists a balanced code datum $\operatorname{BalancedCodeData}(vp)$.

That is, given the subsystem CSS code from Theorem~15 (with parameters $[[N_0, K_0, D_{X,0}, D_{Z,0}]]$) and the Gilbert--Varshamov classical code of length $n_c = q^2$ (Theorem~10), the Evra--Kaufman--Zemor distance balancing procedure (EKZ22, Theorem~1) produces a (non-subsystem) CSS code with:
\begin{itemize}
  \item Block length $N = N_0 \cdot n_c$,
  \item Logical qubits $K \geq K_0 \cdot k_c$ where $k_c \geq n_c/2$,
  \item X-distance $D_X \geq D_{X,0} \cdot d_c$ where $d_c \geq n_c/10$,
  \item Z-distance $D_Z \geq D_{Z,0}$,
  \item LDPC property preserved (row and column weights bounded by a constant).
\end{itemize}

\textit{Note: This is stated as an axiom because the full proof was not completed in the formalization. The EKZ result is from: Evra, Kaufman, Zemor, ``Decodable quantum LDPC codes beyond the $\sqrt{n}$ distance barrier using high-dimensional expanders'', FOCS 2022.}
\end{theorem}

\begin{lemma}[EKZ Satisfiability]
\label{lem:DistanceBalancedFamily.ekz_satisfiable}
\lean{DistanceBalancedFamily.ekz_satisfiable}
\leanok
\uses{thm:ExplicitFamilyQuantumCodes.infinitelyManyValidPrimes, thm:DistanceBalancedFamily.ekz_distance_balancing}

There exists a valid prime $vp$ such that $\operatorname{BalancedCodeData}(vp)$ is nonempty. In other words, the EKZ axiom's hypothesis is satisfiable.
\end{lemma}

\begin{proof}
\leanok
\uses{thm:ExplicitFamilyQuantumCodes.infinitelyManyValidPrimes, thm:DistanceBalancedFamily.ekz_distance_balancing}

From \texttt{infinitelyManyValidPrimes} applied at bound $0$, we obtain a valid prime $vp$. Then the axiom \texttt{ekz\_distance\_balancing} applied to $vp$ yields an element of $\operatorname{BalancedCodeData}(vp)$, establishing nonemptiness.
\end{proof}

\begin{definition}[Distance Balancing]
\label{def:DistanceBalancedFamily.distanceBalancing}
\lean{DistanceBalancedFamily.distanceBalancing}
\leanok
\uses{def:ExplicitFamilyQuantumCodes.ValidPrime, def:DistanceBalancedFamily.BalancedCodeData, thm:DistanceBalancedFamily.ekz_distance_balancing}

For a valid prime $vp$, $\operatorname{distanceBalancing}(vp)$ is the balanced code datum obtained by applying the EKZ distance balancing axiom: $\operatorname{distanceBalancing}(vp) := \operatorname{ekz_distance_balancing}(vp)$.
\end{definition}

\begin{lemma}[Classical $k_c$ Lower Bound (I)]
\label{lem:DistanceBalancedFamily.classicalK_ge}
\lean{DistanceBalancedFamily.classicalK_ge}
\leanok
\uses{def:ExplicitFamilyQuantumCodes.ValidPrime, def:DistanceBalancedFamily.classicalCodeLength}

For any valid prime $vp$, we have $\lfloor n_c / 2 \rfloor + 1 \geq 1$, where $n_c = \operatorname{classicalCodeLength}(vp) = q^2$.
\end{lemma}

\begin{proof}
\leanok
\uses{def:DistanceBalancedFamily.classicalCodeLength}

This follows immediately by integer arithmetic (\texttt{omega}).
\end{proof}

\begin{lemma}[Classical $k_c$ Lower Bound (II)]
\label{lem:DistanceBalancedFamily.classicalK_lower}
\lean{DistanceBalancedFamily.classicalK_lower}
\leanok
\uses{def:ExplicitFamilyQuantumCodes.ValidPrime, def:DistanceBalancedFamily.classicalCodeLength}

For any valid prime $vp$ (so $q \geq 41$), the real-valued inequality
\[
  \lfloor n_c / 2 \rfloor + 1 \;\geq\; \frac{q^2}{4}
\]
holds, where $n_c = q^2$.
\end{lemma}

\begin{proof}
\leanok
\uses{def:DistanceBalancedFamily.classicalCodeLength}

Unfolding $\operatorname{classicalCodeLength}$ gives $n_c = q^2$. Casting to $\mathbb{R}$ via \texttt{push\_cast} and using the fact that $q^2 \geq 0$ by \texttt{positivity}, the bound $\lfloor q^2/2 \rfloor + 1 \geq q^2/4$ follows from linear arithmetic (\texttt{linarith}).
\end{proof}

\begin{lemma}[Classical $d_c$ Lower Bound]
\label{lem:DistanceBalancedFamily.classicalD_lower}
\lean{DistanceBalancedFamily.classicalD_lower}
\leanok
\uses{def:ExplicitFamilyQuantumCodes.ValidPrime, def:DistanceBalancedFamily.classicalCodeLength}

For any valid prime $vp$, the real-valued inequality
\[
  \lfloor n_c / 10 \rfloor \;\geq\; \frac{q^2}{20}
\]
holds, where $n_c = q^2$.
\end{lemma}

\begin{proof}
\leanok
\uses{def:DistanceBalancedFamily.classicalCodeLength}

Unfolding $\operatorname{classicalCodeLength}$ gives $n_c = q^2$. Casting to $\mathbb{R}$ and using $q^2 \geq 0$, the bound follows from linear arithmetic.
\end{proof}

\begin{lemma}[Balanced Code Length Equation]
\label{lem:DistanceBalancedFamily.balancedCodeLength_eq}
\lean{DistanceBalancedFamily.balancedCodeLength_eq}
\leanok
\uses{def:DistanceBalancedFamily.balancedCodeLength, def:DistanceBalancedFamily.classicalCodeLength}

For any valid prime $vp$,
\[
  \operatorname{balancedCodeLength}(vp) \;=\; \operatorname{codeLength}(vp) \cdot \operatorname{classicalCodeLength}(vp).
\]
\end{lemma}

\begin{proof}
\leanok
\uses{def:DistanceBalancedFamily.balancedCodeLength, def:DistanceBalancedFamily.classicalCodeLength}

This holds by reflexivity of the definition.
\end{proof}

\begin{lemma}[$N$ Lower Bound]
\label{lem:DistanceBalancedFamily.balancedN_lower}
\lean{DistanceBalancedFamily.balancedN_lower}
\leanok
\uses{def:ExplicitFamilyQuantumCodes.ValidPrime, def:DistanceBalancedFamily.balancedCodeLength, def:ExplicitFamilyQuantumCodes.codeLength}

There exists a constant $c > 0$ such that for every valid prime $vp$,
\[
  c \cdot q^5 \;\leq\; N,
\]
where $N = \operatorname{balancedCodeLength}(vp)$. Concretely, $c = \frac{3 \cdot 201}{2}$ works.
\end{lemma}

\begin{proof}
\leanok
\uses{def:DistanceBalancedFamily.balancedCodeLength, def:ExplicitFamilyQuantumCodes.codeLength}

We take $c = 3 \cdot 201 / 2$, which is positive by \texttt{positivity}. For any valid prime $vp$, we simplify $\operatorname{balancedCodeLength}(vp) = \operatorname{codeLength}(vp) \cdot q^2$ and cast to $\mathbb{R}$. We have $q \geq 41$ from \texttt{validPrime\_q\_ge}. Using the explicit formula $\operatorname{codeLength}(vp) = 3 \cdot q \cdot (q^2 - 1)$ (from \texttt{codeLength\_eq\_real}), and the bounds $q^2 - 1 \geq q^2/2$ (since $q \geq 41$ implies $q^2 \geq 41^2$), the result follows by nonlinear arithmetic (\texttt{nlinarith}).
\end{proof}

\begin{lemma}[$N$ Upper Bound]
\label{lem:DistanceBalancedFamily.balancedN_upper}
\lean{DistanceBalancedFamily.balancedN_upper}
\leanok
\uses{def:ExplicitFamilyQuantumCodes.ValidPrime, def:DistanceBalancedFamily.balancedCodeLength, def:ExplicitFamilyQuantumCodes.codeLength}

There exists a constant $C > 0$ such that for every valid prime $vp$,
\[
  N \;\leq\; C \cdot q^5,
\]
where $N = \operatorname{balancedCodeLength}(vp)$. Concretely, $C = 3 \cdot 201$ works.
\end{lemma}

\begin{proof}
\leanok
\uses{def:DistanceBalancedFamily.balancedCodeLength, def:ExplicitFamilyQuantumCodes.codeLength}

We take $C = 3 \cdot 201$, which is positive by \texttt{positivity}. For any valid prime $vp$, we simplify $N = \operatorname{codeLength}(vp) \cdot q^2$ and cast to $\mathbb{R}$. Using the formula $\operatorname{codeLength}(vp) = 3q(q^2-1)$ and the bound $q^2 - 1 \leq q^2$, the result $3q(q^2-1) \cdot q^2 \leq 3 \cdot 201 \cdot q^5$ follows by nonlinear arithmetic.
\end{proof}

\begin{lemma}[$K$ Lower Bound]
\label{lem:DistanceBalancedFamily.balancedK_lower}
\lean{DistanceBalancedFamily.balancedK_lower}
\leanok
\uses{def:ExplicitFamilyQuantumCodes.ValidPrime, def:DistanceBalancedFamily.distanceBalancing, thm:ExplicitFamilyQuantumCodes.lps_K_lower_bound, lem:DistanceBalancedFamily.classicalK_lower}

There exists a constant $c > 0$ such that for every valid prime $vp$,
\[
  c \cdot q^4 \;\leq\; K,
\]
where $K$ is the number of logical qubits of $\operatorname{distanceBalancing}(vp).\mathcal{C}$.
\end{lemma}

\begin{proof}
\leanok
\uses{thm:ExplicitFamilyQuantumCodes.lps_K_lower_bound, def:DistanceBalancedFamily.distanceBalancing, lem:DistanceBalancedFamily.classicalK_lower}

From \texttt{lps\_K\_lower\_bound}, we obtain $c_0 > 0$ such that $K_0 \geq c_0 (q^2 - 1)$ for all valid primes. We take $c = c_0/8$. For any valid prime $vp$:
\begin{itemize}
  \item From the EKZ axiom (\texttt{hK} field of $\operatorname{distanceBalancing}(vp)$), we have $K \geq K_0 \cdot (\lfloor n_c/2 \rfloor + 1)$.
  \item From \texttt{lps\_K\_lower\_bound}: $K_0 \geq c_0(q^2 - 1)$.
  \item From \texttt{classicalK\_lower}: $\lfloor n_c/2 \rfloor + 1 \geq q^2/4$.
\end{itemize}
Since $q \geq 41$, we have $q^2 - 1 \geq q^2/2$. Combining:
\[
  K \;\geq\; K_0 \cdot \frac{q^2}{4} \;\geq\; c_0(q^2-1) \cdot \frac{q^2}{4} \;\geq\; c_0 \cdot \frac{q^2}{2} \cdot \frac{q^2}{4} \;=\; \frac{c_0}{8} q^4.
\]
This chain of inequalities is verified by \texttt{nlinarith} and \texttt{linarith}.
\end{proof}

\begin{lemma}[$K$ Upper Bound]
\label{lem:DistanceBalancedFamily.balancedK_upper}
\lean{DistanceBalancedFamily.balancedK_upper}
\leanok
\uses{def:ExplicitFamilyQuantumCodes.ValidPrime, def:DistanceBalancedFamily.distanceBalancing, thm:ExplicitFamilyQuantumCodes.lps_K_upper_bound, def:DistanceBalancedFamily.classicalCodeLength}

There exists a constant $C > 0$ such that for every valid prime $vp$,
\[
  K \;\leq\; C \cdot q^4,
\]
where $K$ is the number of logical qubits of $\operatorname{distanceBalancing}(vp).\mathcal{C}$.
\end{lemma}

\begin{proof}
\leanok
\uses{thm:ExplicitFamilyQuantumCodes.lps_K_upper_bound, def:DistanceBalancedFamily.distanceBalancing, def:DistanceBalancedFamily.classicalCodeLength}

From \texttt{lps\_K\_upper\_bound}, we obtain $C_0 > 0$ such that $K_0 \leq C_0(q^2 - 1)$ for all valid primes. We take $C = C_0$. For any valid prime $vp$:
\begin{itemize}
  \item From the EKZ axiom (\texttt{hK\_upper} field): $K \leq K_0 \cdot n_c$.
  \item From \texttt{lps\_K\_upper\_bound}: $K_0 \leq C_0(q^2 - 1)$.
  \item By definition: $n_c = q^2$.
\end{itemize}
Combining: $K \leq K_0 \cdot q^2 \leq C_0(q^2-1) \cdot q^2 \leq C_0 q^4$, where the last step uses $q^2 - 1 \leq q^2$, verified by \texttt{nlinarith}.
\end{proof}

\begin{lemma}[$D_X$ Lower Bound]
\label{lem:DistanceBalancedFamily.balancedDX_lower}
\lean{DistanceBalancedFamily.balancedDX_lower}
\leanok
\uses{def:ExplicitFamilyQuantumCodes.ValidPrime, def:DistanceBalancedFamily.distanceBalancing, thm:ExplicitFamilyQuantumCodes.lps_DX_lower_bound, lem:DistanceBalancedFamily.classicalD_lower}

There exists a constant $c > 0$ such that for every valid prime $vp$,
\[
  c \cdot q^3 \;\leq\; D_X,
\]
where $D_X$ is the X-distance of $\operatorname{distanceBalancing}(vp).\mathcal{C}$.
\end{lemma}

\begin{proof}
\leanok
\uses{thm:ExplicitFamilyQuantumCodes.lps_DX_lower_bound, def:DistanceBalancedFamily.distanceBalancing, lem:DistanceBalancedFamily.classicalD_lower}

From \texttt{lps\_DX\_lower\_bound}, we obtain $c_0 > 0$ such that $D_{X,0} \geq c_0 \cdot q$ for all valid primes. We take $c = c_0/20$. For any valid prime $vp$:
\begin{itemize}
  \item From the EKZ axiom (\texttt{hDX} field): $D_X \geq D_{X,0} \cdot \lfloor n_c/10 \rfloor$.
  \item From \texttt{lps\_DX\_lower\_bound}: $D_{X,0} \geq c_0 q$.
  \item From \texttt{classicalD\_lower}: $\lfloor n_c/10 \rfloor \geq q^2/20$.
\end{itemize}
Combining:
\[
  D_X \;\geq\; D_{X,0} \cdot \frac{q^2}{20} \;\geq\; c_0 q \cdot \frac{q^2}{20} \;=\; \frac{c_0}{20} q^3.
\]
The chain is assembled by \texttt{calc} steps using \texttt{linarith} and exact casts.
\end{proof}

\begin{lemma}[$D_Z$ Lower Bound]
\label{lem:DistanceBalancedFamily.balancedDZ_lower}
\lean{DistanceBalancedFamily.balancedDZ_lower}
\leanok
\uses{def:ExplicitFamilyQuantumCodes.ValidPrime, def:DistanceBalancedFamily.distanceBalancing, thm:ExplicitFamilyQuantumCodes.lps_DZ_lower_bound}

There exists a constant $c > 0$ such that for every valid prime $vp$,
\[
  c \cdot q^3 \;\leq\; D_Z,
\]
where $D_Z$ is the Z-distance of $\operatorname{distanceBalancing}(vp).\mathcal{C}$.
\end{lemma}

\begin{proof}
\leanok
\uses{thm:ExplicitFamilyQuantumCodes.lps_DZ_lower_bound, def:DistanceBalancedFamily.distanceBalancing}

From \texttt{lps\_DZ\_lower\_bound}, we obtain $c_0 > 0$ such that $D_{Z,0} \geq c_0 \cdot q(q^2-1)$ for all valid primes. We take $c = c_0/2$. For any valid prime $vp$:
\begin{itemize}
  \item From the EKZ axiom (\texttt{hDZ} field): $D_Z \geq D_{Z,0}$.
  \item From \texttt{lps\_DZ\_lower\_bound}: $D_{Z,0} \geq c_0 q(q^2-1)$.
  \item Since $q \geq 41$, we have $q^2 \geq 41^2$, hence $q^2 - 1 \geq q^2/2$.
\end{itemize}
Therefore $q(q^2-1) \geq q \cdot q^2/2 = q^3/2$ (verified by \texttt{nlinarith}), and:
\[
  D_Z \;\geq\; D_{Z,0} \;\geq\; c_0 q(q^2-1) \;\geq\; c_0 \cdot \frac{q^3}{2} \;=\; \frac{c_0}{2} q^3.
\]
\end{proof}

\begin{lemma}[$N$ Unbounded]
\label{lem:DistanceBalancedFamily.balancedN_unbounded}
\lean{DistanceBalancedFamily.balancedN_unbounded}
\leanok
\uses{def:DistanceBalancedFamily.balancedCodeLength, lem:DistanceBalancedFamily.balancedN_lower, thm:ExplicitFamilyQuantumCodes.infinitelyManyValidPrimes}

The family is infinite: for every $M \in \mathbb{N}$, there exists a valid prime $vp$ such that $\operatorname{balancedCodeLength}(vp) > M$.
\end{lemma}

\begin{proof}
\leanok
\uses{lem:DistanceBalancedFamily.balancedN_lower, thm:ExplicitFamilyQuantumCodes.infinitelyManyValidPrimes}

Let $M \in \mathbb{N}$ be arbitrary. From \texttt{balancedN\_lower}, we obtain $c > 0$ such that $c \cdot q^5 \leq N$ for all valid primes. Applying \texttt{infinitelyManyValidPrimes} at the bound $\lceil (M+1)/c \rceil + 1$, we obtain a valid prime $vp$ with $q > \lceil (M+1)/c \rceil + 1$. Since $\lceil (M+1)/c \rceil \geq (M+1)/c$, we have $q > (M+1)/c$, i.e., $cq > M+1$. Since $q \geq 41 \geq 1$, we get $q^5 \geq q^1 \geq q$, hence $c q^5 \geq cq > M$. Therefore $N \geq c q^5 > M$, establishing $N > M$ (cast back to $\mathbb{N}$).
\end{proof}

\begin{theorem}[Corollary 3: Distance-Balanced Family of Quantum Codes]
\label{thm:DistanceBalancedFamily.distanceBalancedFamily}
\lean{DistanceBalancedFamily.distanceBalancedFamily}
\leanok
\uses{def:ExplicitFamilyQuantumCodes.ValidPrime, def:ExplicitFamilyQuantumCodes.codeLength, thm:ExplicitFamilyQuantumCodes.infinitelyManyValidPrimes, thm:ExplicitFamilyQuantumCodes.lps_DX_lower_bound, thm:ExplicitFamilyQuantumCodes.lps_DZ_lower_bound, thm:ExplicitFamilyQuantumCodes.lps_K_lower_bound, thm:ExplicitFamilyQuantumCodes.lps_K_upper_bound}

There exists an explicit construction of an infinite family of $[[N, K, D]]$ LDPC quantum CSS codes such that:
\begin{enumerate}
  \item \textbf{(LDPC existence)} For each valid prime $vp$, there exists an LDPC CSS code of block length $N = \operatorname{balancedCodeLength}(vp)$ with parity check matrices of bounded weight.
  \item \textbf{(Infinitely many)} For every $M \in \mathbb{N}$, there exists a valid prime $vp$ with $\operatorname{balancedCodeLength}(vp) > M$.
  \item \textbf{($K \in \Theta(N^{4/5})$)} There exist constants $0 < c \leq C$ such that for every valid prime $vp$,
    \[
      c \cdot q^4 \;\leq\; K \;\leq\; C \cdot q^4,
    \]
    where $K$ is the logical qubit count of the balanced code and $N \in \Theta(q^5)$ gives $q \in \Theta(N^{1/5})$.
  \item \textbf{($D \in \Omega(N^{3/5})$)} There exists a constant $c > 0$ such that for every valid prime $vp$,
    \[
      c \cdot q^3 \;\leq\; \min(D_X, D_Z),
    \]
    so $D = \min(D_X, D_Z) \in \Omega(q^3) = \Omega(N^{3/5})$.
  \item \textbf{($N \in \Theta(q^5)$)} There exist constants $0 < c \leq C$ such that for every valid prime $vp$,
    \[
      c \cdot q^5 \;\leq\; N \;\leq\; C \cdot q^5.
    \]
\end{enumerate}
\end{theorem}

\begin{proof}
\leanok
\uses{lem:DistanceBalancedFamily.balancedN_unbounded, lem:DistanceBalancedFamily.balancedK_lower, lem:DistanceBalancedFamily.balancedK_upper, lem:DistanceBalancedFamily.balancedDX_lower, lem:DistanceBalancedFamily.balancedDZ_lower, lem:DistanceBalancedFamily.balancedN_lower, lem:DistanceBalancedFamily.balancedN_upper, def:DistanceBalancedFamily.distanceBalancing, thm:ExplicitFamilyQuantumCodes.infinitelyManyValidPrimes, thm:ExplicitFamilyQuantumCodes.lps_K_lower_bound, thm:ExplicitFamilyQuantumCodes.lps_K_upper_bound, thm:ExplicitFamilyQuantumCodes.lps_DX_lower_bound, thm:ExplicitFamilyQuantumCodes.lps_DZ_lower_bound}

\textbf{\color{orange}[Uses unproven axiom: thm:DistanceBalancedFamily.ekz\_distance\_balancing (EKZ22, Theorem 1)]}

We verify each part separately.

\textbf{Part 1 (LDPC existence).} For each valid prime $vp$, let $D = \operatorname{distanceBalancing}(vp)$. The fields $D.r_X$, $D.r_Z$, $D.\mathcal{C}$, $D.w$, $D.\texttt{isLDPC\_HX}$, and $D.\texttt{isLDPC\_HZT}$ directly provide the required LDPC CSS code.

\textbf{Part 2 (Infinitely many).} This is exactly the statement of \texttt{balancedN\_unbounded}.

\textbf{Part 3 ($K \in \Theta(q^4)$).} From \texttt{balancedK\_lower}, we obtain $c > 0$ with $c q^4 \leq K$ for all valid primes. From \texttt{balancedK\_upper}, we obtain $C > 0$ with $K \leq C q^4$ for all valid primes. We then set $c, C$ and verify $c q^4 \leq K \leq C q^4$ pointwise.

\textbf{Part 4 ($D \in \Omega(q^3)$).} From \texttt{balancedDX\_lower}, we obtain $c_1 > 0$ with $c_1 q^3 \leq D_X$. From \texttt{balancedDZ\_lower}, we obtain $c_2 > 0$ with $c_2 q^3 \leq D_Z$. Setting $c = \min(c_1, c_2) > 0$ (since $c > 0$ iff $c_1 > 0$ and $c_2 > 0$), we verify:
\begin{itemize}
  \item $\min(c_1,c_2) q^3 \leq c_1 q^3 \leq D_X$, using $\min(c_1,c_2) \leq c_1$ and $q^3 \geq 0$.
  \item $\min(c_1,c_2) q^3 \leq c_2 q^3 \leq D_Z$, using $\min(c_1,c_2) \leq c_2$ and $q^3 \geq 0$.
\end{itemize}
Both bounds are verified by \texttt{le\_min\_iff} decomposition and \texttt{calc}.

\textbf{Part 5 ($N \in \Theta(q^5)$).} From \texttt{balancedN\_lower} and \texttt{balancedN\_upper}, we obtain constants $c, C > 0$ with $c q^5 \leq N \leq C q^5$ for all valid primes.
\end{proof}

\chapter{Overview: Balanced Product Codes Formalization}

This chapter provides the root module of the balanced product codes formalization. The file \texttt{balanced\_product\_codes.lean} contains only import declarations, assembling all definitions, lemmas, theorems, corollaries, and remarks from the submodules into a single coherent library. No new declarations are made at this level; all mathematical content is introduced in the imported files listed below.

The formalization is organized into the following components:

\begin{enumerate}
\item \textbf{Definitions (Def\_1 through Def\_30):} Chain complexes and cohomology, classical codes, CSS codes, LDPC codes, subsystem CSS codes, cell complexes, cycle graphs, double complexes, total complexes, tensor product double complexes, fiber bundle double complexes, augmented complexes, graph expansion, Tanner code local systems, dual codes, binary entropy, Cheeger constants, edge boundary vertices, Cayley graphs, LPS expander graphs, balanced product vector spaces, balanced product chain complexes, quotient graph trivializations, invariant labelings, balanced product Tanner cycle codes, horizontal/vertical homology splittings, iota/pi maps, horizontal subsystem balanced product codes, and unipotent subgroups for LPS.

\item \textbf{Theorems and Corollaries (Thm\_1 through Thm\_15, Cor\_1 through Cor\_3):} The K\"{u}nneth formula, small double complex homology, fiber bundle homology, the projection-induces-isomorphism theorem, the Alon--Boppana bound, the Alon--Chung theorem, the Alon--Chung contrapositive, the Sipser--Spielman expander code distance bound, expander violated checks, expander bit degree bounds, the Gilbert--Varshamov-plus theorem, LPS Ramanujan graphs, encoding rate for circle codes, homological and cohomological distance bounds, subsystem code parameters, explicit families of quantum codes, and the distance-balanced family theorem.

\item \textbf{Lemmas (Lem\_1 through Lem\_4):} Relative Cheeger inequalities, edge-to-vertex expansion, relative vertex-to-edge expansion, and the K\"{u}nneth balanced product.

\item \textbf{Remarks (Rem\_1 through Rem\_3):} Base field conventions, notation conventions, and the expanding matrix definition.
\end{enumerate}

